% =============================================================================
%  Cuerpo del Tema 4. Materiales dieléctricos
% =============================================================================
\hacerportada

% ------------------------------------------------------------- Página de resumen
\thispagestyle{empty}
\begin{center}
  {\small\color{tenue}\espaciado{Electromagnetismo I}\quad$\cdot$\quad\espaciado{Tema 4}}\\[1.6em]
  {\fuentetitulo\fontsize{21}{25}\selectfont Materiales dieléctricos}\\[0.9em]
  {\color{acentoA}\rule{2.6cm}{0.6pt}}\\[1.1em]
  {\large Daniel Martín Guerrero}\\[0.3em]
  {\small\itshape\color{tenue} Grado en Física, Facultad de Ciencias, Universidad de Salamanca}
\end{center}
\vspace{1.6em}

\noindent
\begin{minipage}[t]{0.63\textwidth}
\small
{\color{acentoA}\espaciado{Resumen.}}\enspace
Estudiamos la electrostática en presencia de aislantes. Partiendo de un modelo atómico sencillo
justificamos la \emph{aproximación dipolar}: un dieléctrico polarizado es una distribución continua de
dipolos, descrita por el vector polarización $\vP$. Demostramos que su campo es el de unas densidades de
carga ligada $\rhoP=-\divg\vP$ y $\sigmaP=\vP\cdot\nn$, de carga total nula, e introducimos el vector
desplazamiento $\vD=\eo\vE+\vP$, cuyo flujo solo ve la carga libre; su rotacional, $\rot\vD=\rot\vP$,
explica cuándo puede anularse. Para medios lineales, homogéneos e isótropos definimos la susceptibilidad
y la permitividad, y deducimos las condiciones de frontera y la ley de refracción de las líneas de campo.
Cerramos la teoría con el campo en cavidades aciculares y lenticulares, que da sentido operacional a $\vE$
y $\vD$. Los seminarios resuelven los doce problemas del tema.
\end{minipage}\hfill
\begin{minipage}[t]{0.33\textwidth}\centering\vspace{-0.6em}
\begin{tikzpicture}
\begin{axis}[marco, width=5.6cm, xmin=-3, xmax=3, ymin=-2.2, ymax=2.2]
  \fill[dielec] (axis cs:0,0) circle[radius=1];
  %% Generado por scripts/lineas_campo.py -- no editar a mano
\addplot[equipot] coordinates {(2.781,-2.200) (2.782,-2.165) (2.782,-2.130) (2.783,-2.094) (2.784,-2.059) (2.784,-2.024) (2.785,-1.989) (2.785,-1.954) (2.786,-1.918) (2.787,-1.883) (2.787,-1.848) (2.788,-1.813) (2.789,-1.778) (2.789,-1.742) (2.790,-1.707) (2.791,-1.672) (2.791,-1.637) (2.792,-1.602) (2.793,-1.566) (2.793,-1.531) (2.794,-1.496) (2.795,-1.461) (2.795,-1.426) (2.796,-1.390) (2.797,-1.355) (2.797,-1.320) (2.798,-1.285) (2.799,-1.250) (2.799,-1.214) (2.800,-1.179) (2.801,-1.144) (2.801,-1.109) (2.802,-1.074) (2.802,-1.038) (2.803,-1.003) (2.804,-0.968) (2.804,-0.933) (2.805,-0.898) (2.805,-0.862) (2.806,-0.827) (2.807,-0.792) (2.807,-0.757) (2.808,-0.722) (2.808,-0.686) (2.809,-0.651) (2.809,-0.616) (2.809,-0.581) (2.810,-0.546) (2.810,-0.510) (2.811,-0.475) (2.811,-0.440) (2.811,-0.405) (2.812,-0.370) (2.812,-0.334) (2.812,-0.299) (2.812,-0.264) (2.813,-0.229) (2.813,-0.194) (2.813,-0.158) (2.813,-0.123) (2.813,-0.088) (2.813,-0.053) (2.813,-0.018) (2.813,0.018) (2.813,0.053) (2.813,0.088) (2.813,0.123) (2.813,0.158) (2.813,0.194) (2.813,0.229) (2.812,0.264) (2.812,0.299) (2.812,0.334) (2.812,0.370) (2.811,0.405) (2.811,0.440) (2.811,0.475) (2.810,0.510) (2.810,0.546) (2.809,0.581) (2.809,0.616) (2.809,0.651) (2.808,0.686) (2.808,0.722) (2.807,0.757) (2.807,0.792) (2.806,0.827) (2.805,0.862) (2.805,0.898) (2.804,0.933) (2.804,0.968) (2.803,1.003) (2.802,1.038) (2.802,1.074) (2.801,1.109) (2.801,1.144) (2.800,1.179) (2.799,1.214) (2.799,1.250) (2.798,1.285) (2.797,1.320) (2.797,1.355) (2.796,1.390) (2.795,1.426) (2.795,1.461) (2.794,1.496) (2.793,1.531) (2.793,1.566) (2.792,1.602) (2.791,1.637) (2.791,1.672) (2.790,1.707) (2.789,1.742) (2.789,1.778) (2.788,1.813) (2.787,1.848) (2.787,1.883) (2.786,1.918) (2.785,1.954) (2.785,1.989) (2.784,2.024) (2.784,2.059) (2.783,2.094) (2.782,2.130) (2.782,2.165) (2.781,2.200)};
\addplot[equipot] coordinates {(2.286,-2.200) (2.287,-2.165) (2.287,-2.130) (2.288,-2.094) (2.289,-2.059) (2.290,-2.024) (2.291,-1.989) (2.292,-1.954) (2.293,-1.918) (2.294,-1.883) (2.295,-1.848) (2.296,-1.813) (2.297,-1.778) (2.298,-1.742) (2.299,-1.707) (2.300,-1.672) (2.301,-1.637) (2.302,-1.602) (2.303,-1.566) (2.304,-1.531) (2.306,-1.496) (2.307,-1.461) (2.308,-1.426) (2.309,-1.390) (2.310,-1.355) (2.311,-1.320) (2.312,-1.285) (2.314,-1.250) (2.315,-1.214) (2.316,-1.179) (2.317,-1.144) (2.318,-1.109) (2.319,-1.074) (2.321,-1.039) (2.322,-1.003) (2.323,-0.968) (2.324,-0.933) (2.325,-0.898) (2.326,-0.863) (2.327,-0.827) (2.328,-0.792) (2.329,-0.757) (2.330,-0.722) (2.331,-0.687) (2.332,-0.651) (2.333,-0.616) (2.334,-0.581) (2.335,-0.546) (2.335,-0.511) (2.336,-0.475) (2.337,-0.440) (2.338,-0.405) (2.338,-0.370) (2.339,-0.335) (2.339,-0.299) (2.340,-0.264) (2.340,-0.229) (2.340,-0.194) (2.341,-0.158) (2.341,-0.123) (2.341,-0.088) (2.341,-0.053) (2.341,-0.018) (2.341,0.018) (2.341,0.053) (2.341,0.088) (2.341,0.123) (2.341,0.158) (2.340,0.194) (2.340,0.229) (2.340,0.264) (2.339,0.299) (2.339,0.335) (2.338,0.370) (2.338,0.405) (2.337,0.440) (2.336,0.475) (2.335,0.511) (2.335,0.546) (2.334,0.581) (2.333,0.616) (2.332,0.651) (2.331,0.687) (2.330,0.722) (2.329,0.757) (2.328,0.792) (2.327,0.827) (2.326,0.863) (2.325,0.898) (2.324,0.933) (2.323,0.968) (2.322,1.003) (2.321,1.039) (2.319,1.074) (2.318,1.109) (2.317,1.144) (2.316,1.179) (2.315,1.214) (2.314,1.250) (2.312,1.285) (2.311,1.320) (2.310,1.355) (2.309,1.390) (2.308,1.426) (2.307,1.461) (2.306,1.496) (2.304,1.531) (2.303,1.566) (2.302,1.602) (2.301,1.637) (2.300,1.672) (2.299,1.707) (2.298,1.742) (2.297,1.778) (2.296,1.813) (2.295,1.848) (2.294,1.883) (2.293,1.918) (2.292,1.954) (2.291,1.989) (2.290,2.024) (2.289,2.059) (2.288,2.094) (2.287,2.130) (2.287,2.165) (2.286,2.200)};
\addplot[equipot] coordinates {(1.789,-2.200) (1.790,-2.165) (1.792,-2.130) (1.793,-2.094) (1.794,-2.059) (1.795,-2.024) (1.797,-1.989) (1.798,-1.953) (1.799,-1.918) (1.801,-1.883) (1.802,-1.848) (1.804,-1.813) (1.806,-1.777) (1.807,-1.742) (1.809,-1.707) (1.810,-1.672) (1.812,-1.637) (1.814,-1.601) (1.816,-1.566) (1.818,-1.531) (1.820,-1.496) (1.822,-1.461) (1.824,-1.425) (1.826,-1.390) (1.828,-1.355) (1.830,-1.320) (1.832,-1.285) (1.834,-1.250) (1.836,-1.214) (1.838,-1.179) (1.840,-1.144) (1.843,-1.109) (1.845,-1.074) (1.847,-1.039) (1.849,-1.003) (1.851,-0.968) (1.854,-0.933) (1.856,-0.898) (1.858,-0.863) (1.860,-0.827) (1.862,-0.792) (1.864,-0.757) (1.866,-0.722) (1.868,-0.687) (1.870,-0.652) (1.872,-0.616) (1.874,-0.581) (1.876,-0.546) (1.877,-0.511) (1.879,-0.476) (1.881,-0.440) (1.882,-0.405) (1.883,-0.370) (1.884,-0.335) (1.885,-0.300) (1.886,-0.264) (1.887,-0.229) (1.888,-0.194) (1.889,-0.159) (1.889,-0.123) (1.890,-0.088) (1.890,-0.053) (1.890,-0.018) (1.890,0.018) (1.890,0.053) (1.890,0.088) (1.889,0.123) (1.889,0.159) (1.888,0.194) (1.887,0.229) (1.886,0.264) (1.885,0.300) (1.884,0.335) (1.883,0.370) (1.882,0.405) (1.881,0.440) (1.879,0.476) (1.877,0.511) (1.876,0.546) (1.874,0.581) (1.872,0.616) (1.870,0.652) (1.868,0.687) (1.866,0.722) (1.864,0.757) (1.862,0.792) (1.860,0.827) (1.858,0.863) (1.856,0.898) (1.854,0.933) (1.851,0.968) (1.849,1.003) (1.847,1.039) (1.845,1.074) (1.843,1.109) (1.840,1.144) (1.838,1.179) (1.836,1.214) (1.834,1.250) (1.832,1.285) (1.830,1.320) (1.828,1.355) (1.826,1.390) (1.824,1.425) (1.822,1.461) (1.820,1.496) (1.818,1.531) (1.816,1.566) (1.814,1.601) (1.812,1.637) (1.810,1.672) (1.809,1.707) (1.807,1.742) (1.806,1.777) (1.804,1.813) (1.802,1.848) (1.801,1.883) (1.799,1.918) (1.798,1.953) (1.797,1.989) (1.795,2.024) (1.794,2.059) (1.793,2.094) (1.792,2.130) (1.790,2.165) (1.789,2.200)};
\addplot[equipot] coordinates {(1.289,-2.200) (1.290,-2.165) (1.292,-2.130) (1.293,-2.095) (1.295,-2.060) (1.297,-2.025) (1.298,-1.990) (1.300,-1.955) (1.302,-1.920) (1.304,-1.885) (1.306,-1.850) (1.308,-1.815) (1.311,-1.780) (1.313,-1.745) (1.316,-1.710) (1.318,-1.675) (1.321,-1.640) (1.324,-1.605) (1.326,-1.570) (1.329,-1.535) (1.332,-1.500) (1.336,-1.465) (1.339,-1.430) (1.342,-1.395) (1.346,-1.360) (1.350,-1.326) (1.353,-1.291) (1.357,-1.256) (1.361,-1.221) (1.365,-1.186) (1.370,-1.151) (1.374,-1.116) (1.378,-1.082) (1.383,-1.047) (1.387,-1.012) (1.392,-0.977) (1.396,-0.943) (1.401,-0.908) (1.405,-0.873) (1.410,-0.838) (1.414,-0.803) (1.419,-0.769) (1.423,-0.734) (1.428,-0.699) (1.432,-0.664) (1.436,-0.629) (1.440,-0.595) (1.444,-0.560) (1.448,-0.525) (1.452,-0.490) (1.455,-0.455) (1.459,-0.420) (1.462,-0.385) (1.464,-0.350) (1.467,-0.315) (1.469,-0.280) (1.472,-0.245) (1.473,-0.210) (1.475,-0.175) (1.476,-0.140) (1.477,-0.105) (1.478,-0.070) (1.479,-0.035) (1.479,0.000) (1.479,0.035) (1.478,0.070) (1.477,0.105) (1.476,0.140) (1.475,0.175) (1.473,0.210) (1.472,0.245) (1.469,0.280) (1.467,0.315) (1.464,0.350) (1.462,0.385) (1.459,0.420) (1.455,0.455) (1.452,0.490) (1.448,0.525) (1.444,0.560) (1.440,0.595) (1.436,0.629) (1.432,0.664) (1.428,0.699) (1.423,0.734) (1.419,0.769) (1.414,0.803) (1.410,0.838) (1.405,0.873) (1.401,0.908) (1.396,0.943) (1.392,0.977) (1.387,1.012) (1.383,1.047) (1.378,1.082) (1.374,1.116) (1.370,1.151) (1.365,1.186) (1.361,1.221) (1.357,1.256) (1.353,1.291) (1.350,1.326) (1.346,1.360) (1.342,1.395) (1.339,1.430) (1.336,1.465) (1.332,1.500) (1.329,1.535) (1.326,1.570) (1.324,1.605) (1.321,1.640) (1.318,1.675) (1.316,1.710) (1.313,1.745) (1.311,1.780) (1.308,1.815) (1.306,1.850) (1.304,1.885) (1.302,1.920) (1.300,1.955) (1.298,1.990) (1.297,2.025) (1.295,2.060) (1.293,2.095) (1.292,2.130) (1.290,2.165) (1.289,2.200)};
\addplot[equipot] coordinates {(0.781,-2.200) (0.782,-2.165) (0.784,-2.130) (0.785,-2.094) (0.787,-2.059) (0.788,-2.024) (0.790,-1.989) (0.792,-1.954) (0.794,-1.918) (0.797,-1.883) (0.799,-1.848) (0.801,-1.813) (0.804,-1.778) (0.807,-1.743) (0.810,-1.708) (0.813,-1.673) (0.817,-1.638) (0.820,-1.602) (0.824,-1.567) (0.828,-1.532) (0.833,-1.498) (0.837,-1.463) (0.842,-1.428) (0.848,-1.393) (0.853,-1.358) (0.859,-1.323) (0.866,-1.289) (0.872,-1.254) (0.879,-1.220) (0.887,-1.185) (0.894,-1.151) (0.903,-1.117) (0.911,-1.082) (0.920,-1.048) (0.929,-1.014) (0.938,-0.980) (0.947,-0.946) (0.957,-0.912) (0.967,-0.878) (0.977,-0.845) (0.987,-0.811) (0.997,-0.777) (1.007,-0.743) (1.017,-0.710) (1.027,-0.676) (1.036,-0.642) (1.046,-0.608) (1.055,-0.574) (1.063,-0.540) (1.072,-0.505) (1.080,-0.471) (1.088,-0.437) (1.095,-0.402) (1.102,-0.368) (1.108,-0.333) (1.113,-0.298) (1.119,-0.263) (1.123,-0.228) (1.127,-0.193) (1.130,-0.158) (1.133,-0.123) (1.135,-0.088) (1.136,-0.053) (1.137,-0.018) (1.137,0.018) (1.136,0.053) (1.135,0.088) (1.133,0.123) (1.130,0.158) (1.127,0.193) (1.123,0.228) (1.119,0.263) (1.113,0.298) (1.108,0.333) (1.102,0.368) (1.095,0.402) (1.088,0.437) (1.080,0.471) (1.072,0.505) (1.063,0.540) (1.055,0.574) (1.046,0.608) (1.036,0.642) (1.027,0.676) (1.017,0.710) (1.007,0.743) (0.997,0.777) (0.987,0.811) (0.977,0.845) (0.967,0.878) (0.957,0.912) (0.947,0.946) (0.938,0.980) (0.929,1.014) (0.920,1.048) (0.911,1.082) (0.903,1.117) (0.894,1.151) (0.887,1.185) (0.879,1.220) (0.872,1.254) (0.866,1.289) (0.859,1.323) (0.853,1.358) (0.848,1.393) (0.842,1.428) (0.837,1.463) (0.833,1.498) (0.828,1.532) (0.824,1.567) (0.820,1.602) (0.817,1.638) (0.813,1.673) (0.810,1.708) (0.807,1.743) (0.804,1.778) (0.801,1.813) (0.799,1.848) (0.797,1.883) (0.794,1.918) (0.792,1.954) (0.790,1.989) (0.788,2.024) (0.787,2.059) (0.785,2.094) (0.784,2.130) (0.782,2.165) (0.781,2.200)};
\addplot[equipot] coordinates {(0.262,-2.200) (0.263,-2.165) (0.263,-2.130) (0.264,-2.095) (0.265,-2.060) (0.266,-2.025) (0.266,-1.990) (0.267,-1.955) (0.268,-1.919) (0.270,-1.884) (0.271,-1.849) (0.272,-1.814) (0.273,-1.779) (0.275,-1.744) (0.277,-1.709) (0.278,-1.674) (0.280,-1.639) (0.283,-1.604) (0.285,-1.569) (0.288,-1.534) (0.291,-1.499) (0.294,-1.464) (0.298,-1.429) (0.302,-1.395) (0.307,-1.360) (0.312,-1.325) (0.318,-1.291) (0.324,-1.256) (0.332,-1.222) (0.340,-1.188) (0.350,-1.154) (0.360,-1.121) (0.373,-1.088) (0.386,-1.055) (0.401,-1.024) (0.417,-0.993) (0.435,-0.962) (0.453,-0.933) (0.473,-0.904) (0.493,-0.875) (0.500,-0.842) (0.500,-0.807) (0.500,-0.772) (0.500,-0.737) (0.500,-0.701) (0.500,-0.666) (0.500,-0.631) (0.500,-0.596) (0.500,-0.561) (0.500,-0.526) (0.500,-0.491) (0.500,-0.456) (0.500,-0.421) (0.500,-0.386) (0.500,-0.351) (0.500,-0.316) (0.500,-0.281) (0.500,-0.246) (0.500,-0.210) (0.500,-0.175) (0.500,-0.140) (0.500,-0.105) (0.500,-0.070) (0.500,-0.035) (0.500,0.000) (0.500,0.035) (0.500,0.070) (0.500,0.105) (0.500,0.140) (0.500,0.175) (0.500,0.210) (0.500,0.246) (0.500,0.281) (0.500,0.316) (0.500,0.351) (0.500,0.386) (0.500,0.421) (0.500,0.456) (0.500,0.491) (0.500,0.526) (0.500,0.561) (0.500,0.596) (0.500,0.631) (0.500,0.666) (0.500,0.701) (0.500,0.737) (0.500,0.772) (0.500,0.807) (0.500,0.842) (0.493,0.875) (0.473,0.904) (0.453,0.933) (0.435,0.962) (0.417,0.993) (0.401,1.024) (0.386,1.055) (0.373,1.088) (0.360,1.121) (0.350,1.154) (0.340,1.188) (0.332,1.222) (0.324,1.256) (0.318,1.291) (0.312,1.325) (0.307,1.360) (0.302,1.395) (0.298,1.429) (0.294,1.464) (0.291,1.499) (0.288,1.534) (0.285,1.569) (0.283,1.604) (0.280,1.639) (0.278,1.674) (0.277,1.709) (0.275,1.744) (0.273,1.779) (0.272,1.814) (0.271,1.849) (0.270,1.884) (0.268,1.919) (0.267,1.955) (0.266,1.990) (0.266,2.025) (0.265,2.060) (0.264,2.095) (0.263,2.130) (0.263,2.165) (0.262,2.200)};
\addplot[equipot] coordinates {(-0.262,-2.200) (-0.263,-2.165) (-0.263,-2.130) (-0.264,-2.095) (-0.265,-2.060) (-0.266,-2.025) (-0.266,-1.990) (-0.267,-1.955) (-0.268,-1.919) (-0.270,-1.884) (-0.271,-1.849) (-0.272,-1.814) (-0.273,-1.779) (-0.275,-1.744) (-0.277,-1.709) (-0.278,-1.674) (-0.280,-1.639) (-0.283,-1.604) (-0.285,-1.569) (-0.288,-1.534) (-0.291,-1.499) (-0.294,-1.464) (-0.298,-1.429) (-0.302,-1.395) (-0.307,-1.360) (-0.312,-1.325) (-0.318,-1.291) (-0.324,-1.256) (-0.332,-1.222) (-0.340,-1.188) (-0.350,-1.154) (-0.360,-1.121) (-0.373,-1.088) (-0.386,-1.055) (-0.401,-1.024) (-0.417,-0.993) (-0.435,-0.962) (-0.453,-0.933) (-0.473,-0.904) (-0.493,-0.875) (-0.500,-0.842) (-0.500,-0.807) (-0.500,-0.772) (-0.500,-0.737) (-0.500,-0.701) (-0.500,-0.666) (-0.500,-0.631) (-0.500,-0.596) (-0.500,-0.561) (-0.500,-0.526) (-0.500,-0.491) (-0.500,-0.456) (-0.500,-0.421) (-0.500,-0.386) (-0.500,-0.351) (-0.500,-0.316) (-0.500,-0.281) (-0.500,-0.246) (-0.500,-0.210) (-0.500,-0.175) (-0.500,-0.140) (-0.500,-0.105) (-0.500,-0.070) (-0.500,-0.035) (-0.500,-0.000) (-0.500,0.035) (-0.500,0.070) (-0.500,0.105) (-0.500,0.140) (-0.500,0.175) (-0.500,0.210) (-0.500,0.246) (-0.500,0.281) (-0.500,0.316) (-0.500,0.351) (-0.500,0.386) (-0.500,0.421) (-0.500,0.456) (-0.500,0.491) (-0.500,0.526) (-0.500,0.561) (-0.500,0.596) (-0.500,0.631) (-0.500,0.666) (-0.500,0.701) (-0.500,0.737) (-0.500,0.772) (-0.500,0.807) (-0.500,0.842) (-0.493,0.875) (-0.473,0.904) (-0.453,0.933) (-0.435,0.962) (-0.417,0.993) (-0.401,1.024) (-0.386,1.055) (-0.373,1.088) (-0.360,1.121) (-0.350,1.154) (-0.340,1.188) (-0.332,1.222) (-0.324,1.256) (-0.318,1.291) (-0.312,1.325) (-0.307,1.360) (-0.302,1.395) (-0.298,1.429) (-0.294,1.464) (-0.291,1.499) (-0.288,1.534) (-0.285,1.569) (-0.283,1.604) (-0.280,1.639) (-0.278,1.674) (-0.277,1.709) (-0.275,1.744) (-0.273,1.779) (-0.272,1.814) (-0.271,1.849) (-0.270,1.884) (-0.268,1.919) (-0.267,1.955) (-0.266,1.990) (-0.266,2.025) (-0.265,2.060) (-0.264,2.095) (-0.263,2.130) (-0.263,2.165) (-0.262,2.200)};
\addplot[equipot] coordinates {(-0.781,-2.200) (-0.782,-2.165) (-0.784,-2.130) (-0.785,-2.094) (-0.787,-2.059) (-0.788,-2.024) (-0.790,-1.989) (-0.792,-1.954) (-0.794,-1.918) (-0.797,-1.883) (-0.799,-1.848) (-0.801,-1.813) (-0.804,-1.778) (-0.807,-1.743) (-0.810,-1.708) (-0.813,-1.673) (-0.817,-1.638) (-0.820,-1.602) (-0.824,-1.567) (-0.828,-1.532) (-0.833,-1.498) (-0.837,-1.463) (-0.842,-1.428) (-0.848,-1.393) (-0.853,-1.358) (-0.859,-1.323) (-0.866,-1.289) (-0.872,-1.254) (-0.879,-1.220) (-0.887,-1.185) (-0.894,-1.151) (-0.903,-1.117) (-0.911,-1.082) (-0.920,-1.048) (-0.929,-1.014) (-0.938,-0.980) (-0.947,-0.946) (-0.957,-0.912) (-0.967,-0.878) (-0.977,-0.845) (-0.987,-0.811) (-0.997,-0.777) (-1.007,-0.743) (-1.017,-0.710) (-1.027,-0.676) (-1.036,-0.642) (-1.046,-0.608) (-1.055,-0.574) (-1.063,-0.540) (-1.072,-0.505) (-1.080,-0.471) (-1.088,-0.437) (-1.095,-0.402) (-1.102,-0.368) (-1.108,-0.333) (-1.113,-0.298) (-1.119,-0.263) (-1.123,-0.228) (-1.127,-0.193) (-1.130,-0.158) (-1.133,-0.123) (-1.135,-0.088) (-1.136,-0.053) (-1.137,-0.018) (-1.137,0.018) (-1.136,0.053) (-1.135,0.088) (-1.133,0.123) (-1.130,0.158) (-1.127,0.193) (-1.123,0.228) (-1.119,0.263) (-1.113,0.298) (-1.108,0.333) (-1.102,0.368) (-1.095,0.402) (-1.088,0.437) (-1.080,0.471) (-1.072,0.505) (-1.063,0.540) (-1.055,0.574) (-1.046,0.608) (-1.036,0.642) (-1.027,0.676) (-1.017,0.710) (-1.007,0.743) (-0.997,0.777) (-0.987,0.811) (-0.977,0.845) (-0.967,0.878) (-0.957,0.912) (-0.947,0.946) (-0.938,0.980) (-0.929,1.014) (-0.920,1.048) (-0.911,1.082) (-0.903,1.117) (-0.894,1.151) (-0.887,1.185) (-0.879,1.220) (-0.872,1.254) (-0.866,1.289) (-0.859,1.323) (-0.853,1.358) (-0.848,1.393) (-0.842,1.428) (-0.837,1.463) (-0.833,1.498) (-0.828,1.532) (-0.824,1.567) (-0.820,1.602) (-0.817,1.638) (-0.813,1.673) (-0.810,1.708) (-0.807,1.743) (-0.804,1.778) (-0.801,1.813) (-0.799,1.848) (-0.797,1.883) (-0.794,1.918) (-0.792,1.954) (-0.790,1.989) (-0.788,2.024) (-0.787,2.059) (-0.785,2.094) (-0.784,2.130) (-0.782,2.165) (-0.781,2.200)};
\addplot[equipot] coordinates {(-1.289,-2.200) (-1.290,-2.165) (-1.292,-2.130) (-1.293,-2.095) (-1.295,-2.060) (-1.297,-2.025) (-1.298,-1.990) (-1.300,-1.955) (-1.302,-1.920) (-1.304,-1.885) (-1.306,-1.850) (-1.308,-1.815) (-1.311,-1.780) (-1.313,-1.745) (-1.316,-1.710) (-1.318,-1.675) (-1.321,-1.640) (-1.324,-1.605) (-1.326,-1.570) (-1.329,-1.535) (-1.332,-1.500) (-1.336,-1.465) (-1.339,-1.430) (-1.342,-1.395) (-1.346,-1.360) (-1.350,-1.326) (-1.353,-1.291) (-1.357,-1.256) (-1.361,-1.221) (-1.365,-1.186) (-1.370,-1.151) (-1.374,-1.116) (-1.378,-1.082) (-1.383,-1.047) (-1.387,-1.012) (-1.392,-0.977) (-1.396,-0.943) (-1.401,-0.908) (-1.405,-0.873) (-1.410,-0.838) (-1.414,-0.803) (-1.419,-0.769) (-1.423,-0.734) (-1.428,-0.699) (-1.432,-0.664) (-1.436,-0.629) (-1.440,-0.595) (-1.444,-0.560) (-1.448,-0.525) (-1.452,-0.490) (-1.455,-0.455) (-1.459,-0.420) (-1.462,-0.385) (-1.464,-0.350) (-1.467,-0.315) (-1.469,-0.280) (-1.472,-0.245) (-1.473,-0.210) (-1.475,-0.175) (-1.476,-0.140) (-1.477,-0.105) (-1.478,-0.070) (-1.479,-0.035) (-1.479,0.000) (-1.479,0.035) (-1.478,0.070) (-1.477,0.105) (-1.476,0.140) (-1.475,0.175) (-1.473,0.210) (-1.472,0.245) (-1.469,0.280) (-1.467,0.315) (-1.464,0.350) (-1.462,0.385) (-1.459,0.420) (-1.455,0.455) (-1.452,0.490) (-1.448,0.525) (-1.444,0.560) (-1.440,0.595) (-1.436,0.629) (-1.432,0.664) (-1.428,0.699) (-1.423,0.734) (-1.419,0.769) (-1.414,0.803) (-1.410,0.838) (-1.405,0.873) (-1.401,0.908) (-1.396,0.943) (-1.392,0.977) (-1.387,1.012) (-1.383,1.047) (-1.378,1.082) (-1.374,1.116) (-1.370,1.151) (-1.365,1.186) (-1.361,1.221) (-1.357,1.256) (-1.353,1.291) (-1.350,1.326) (-1.346,1.360) (-1.342,1.395) (-1.339,1.430) (-1.336,1.465) (-1.332,1.500) (-1.329,1.535) (-1.326,1.570) (-1.324,1.605) (-1.321,1.640) (-1.318,1.675) (-1.316,1.710) (-1.313,1.745) (-1.311,1.780) (-1.308,1.815) (-1.306,1.850) (-1.304,1.885) (-1.302,1.920) (-1.300,1.955) (-1.298,1.990) (-1.297,2.025) (-1.295,2.060) (-1.293,2.095) (-1.292,2.130) (-1.290,2.165) (-1.289,2.200)};
\addplot[equipot] coordinates {(-1.789,-2.200) (-1.790,-2.165) (-1.792,-2.130) (-1.793,-2.094) (-1.794,-2.059) (-1.795,-2.024) (-1.797,-1.989) (-1.798,-1.953) (-1.799,-1.918) (-1.801,-1.883) (-1.802,-1.848) (-1.804,-1.813) (-1.806,-1.777) (-1.807,-1.742) (-1.809,-1.707) (-1.810,-1.672) (-1.812,-1.637) (-1.814,-1.601) (-1.816,-1.566) (-1.818,-1.531) (-1.820,-1.496) (-1.822,-1.461) (-1.824,-1.425) (-1.826,-1.390) (-1.828,-1.355) (-1.830,-1.320) (-1.832,-1.285) (-1.834,-1.250) (-1.836,-1.214) (-1.838,-1.179) (-1.840,-1.144) (-1.843,-1.109) (-1.845,-1.074) (-1.847,-1.039) (-1.849,-1.003) (-1.851,-0.968) (-1.854,-0.933) (-1.856,-0.898) (-1.858,-0.863) (-1.860,-0.827) (-1.862,-0.792) (-1.864,-0.757) (-1.866,-0.722) (-1.868,-0.687) (-1.870,-0.652) (-1.872,-0.616) (-1.874,-0.581) (-1.876,-0.546) (-1.877,-0.511) (-1.879,-0.476) (-1.881,-0.440) (-1.882,-0.405) (-1.883,-0.370) (-1.884,-0.335) (-1.885,-0.300) (-1.886,-0.264) (-1.887,-0.229) (-1.888,-0.194) (-1.889,-0.159) (-1.889,-0.123) (-1.890,-0.088) (-1.890,-0.053) (-1.890,-0.018) (-1.890,0.018) (-1.890,0.053) (-1.890,0.088) (-1.889,0.123) (-1.889,0.159) (-1.888,0.194) (-1.887,0.229) (-1.886,0.264) (-1.885,0.300) (-1.884,0.335) (-1.883,0.370) (-1.882,0.405) (-1.881,0.440) (-1.879,0.476) (-1.877,0.511) (-1.876,0.546) (-1.874,0.581) (-1.872,0.616) (-1.870,0.652) (-1.868,0.687) (-1.866,0.722) (-1.864,0.757) (-1.862,0.792) (-1.860,0.827) (-1.858,0.863) (-1.856,0.898) (-1.854,0.933) (-1.851,0.968) (-1.849,1.003) (-1.847,1.039) (-1.845,1.074) (-1.843,1.109) (-1.840,1.144) (-1.838,1.179) (-1.836,1.214) (-1.834,1.250) (-1.832,1.285) (-1.830,1.320) (-1.828,1.355) (-1.826,1.390) (-1.824,1.425) (-1.822,1.461) (-1.820,1.496) (-1.818,1.531) (-1.816,1.566) (-1.814,1.601) (-1.812,1.637) (-1.810,1.672) (-1.809,1.707) (-1.807,1.742) (-1.806,1.777) (-1.804,1.813) (-1.802,1.848) (-1.801,1.883) (-1.799,1.918) (-1.798,1.953) (-1.797,1.989) (-1.795,2.024) (-1.794,2.059) (-1.793,2.094) (-1.792,2.130) (-1.790,2.165) (-1.789,2.200)};
\addplot[equipot] coordinates {(-2.286,-2.200) (-2.287,-2.165) (-2.287,-2.130) (-2.288,-2.094) (-2.289,-2.059) (-2.290,-2.024) (-2.291,-1.989) (-2.292,-1.954) (-2.293,-1.918) (-2.294,-1.883) (-2.295,-1.848) (-2.296,-1.813) (-2.297,-1.778) (-2.298,-1.742) (-2.299,-1.707) (-2.300,-1.672) (-2.301,-1.637) (-2.302,-1.602) (-2.303,-1.566) (-2.304,-1.531) (-2.306,-1.496) (-2.307,-1.461) (-2.308,-1.426) (-2.309,-1.390) (-2.310,-1.355) (-2.311,-1.320) (-2.312,-1.285) (-2.314,-1.250) (-2.315,-1.214) (-2.316,-1.179) (-2.317,-1.144) (-2.318,-1.109) (-2.319,-1.074) (-2.321,-1.039) (-2.322,-1.003) (-2.323,-0.968) (-2.324,-0.933) (-2.325,-0.898) (-2.326,-0.863) (-2.327,-0.827) (-2.328,-0.792) (-2.329,-0.757) (-2.330,-0.722) (-2.331,-0.687) (-2.332,-0.651) (-2.333,-0.616) (-2.334,-0.581) (-2.335,-0.546) (-2.335,-0.511) (-2.336,-0.475) (-2.337,-0.440) (-2.338,-0.405) (-2.338,-0.370) (-2.339,-0.335) (-2.339,-0.299) (-2.340,-0.264) (-2.340,-0.229) (-2.340,-0.194) (-2.341,-0.158) (-2.341,-0.123) (-2.341,-0.088) (-2.341,-0.053) (-2.341,-0.018) (-2.341,0.018) (-2.341,0.053) (-2.341,0.088) (-2.341,0.123) (-2.341,0.158) (-2.340,0.194) (-2.340,0.229) (-2.340,0.264) (-2.339,0.299) (-2.339,0.335) (-2.338,0.370) (-2.338,0.405) (-2.337,0.440) (-2.336,0.475) (-2.335,0.511) (-2.335,0.546) (-2.334,0.581) (-2.333,0.616) (-2.332,0.651) (-2.331,0.687) (-2.330,0.722) (-2.329,0.757) (-2.328,0.792) (-2.327,0.827) (-2.326,0.863) (-2.325,0.898) (-2.324,0.933) (-2.323,0.968) (-2.322,1.003) (-2.321,1.039) (-2.319,1.074) (-2.318,1.109) (-2.317,1.144) (-2.316,1.179) (-2.315,1.214) (-2.314,1.250) (-2.312,1.285) (-2.311,1.320) (-2.310,1.355) (-2.309,1.390) (-2.308,1.426) (-2.307,1.461) (-2.306,1.496) (-2.304,1.531) (-2.303,1.566) (-2.302,1.602) (-2.301,1.637) (-2.300,1.672) (-2.299,1.707) (-2.298,1.742) (-2.297,1.778) (-2.296,1.813) (-2.295,1.848) (-2.294,1.883) (-2.293,1.918) (-2.292,1.954) (-2.291,1.989) (-2.290,2.024) (-2.289,2.059) (-2.288,2.094) (-2.287,2.130) (-2.287,2.165) (-2.286,2.200)};
\addplot[equipot] coordinates {(-2.781,-2.200) (-2.782,-2.165) (-2.782,-2.130) (-2.783,-2.094) (-2.784,-2.059) (-2.784,-2.024) (-2.785,-1.989) (-2.785,-1.954) (-2.786,-1.918) (-2.787,-1.883) (-2.787,-1.848) (-2.788,-1.813) (-2.789,-1.778) (-2.789,-1.742) (-2.790,-1.707) (-2.791,-1.672) (-2.791,-1.637) (-2.792,-1.602) (-2.793,-1.566) (-2.793,-1.531) (-2.794,-1.496) (-2.795,-1.461) (-2.795,-1.426) (-2.796,-1.390) (-2.797,-1.355) (-2.797,-1.320) (-2.798,-1.285) (-2.799,-1.250) (-2.799,-1.214) (-2.800,-1.179) (-2.801,-1.144) (-2.801,-1.109) (-2.802,-1.074) (-2.802,-1.038) (-2.803,-1.003) (-2.804,-0.968) (-2.804,-0.933) (-2.805,-0.898) (-2.805,-0.862) (-2.806,-0.827) (-2.807,-0.792) (-2.807,-0.757) (-2.808,-0.722) (-2.808,-0.686) (-2.809,-0.651) (-2.809,-0.616) (-2.809,-0.581) (-2.810,-0.546) (-2.810,-0.510) (-2.811,-0.475) (-2.811,-0.440) (-2.811,-0.405) (-2.812,-0.370) (-2.812,-0.334) (-2.812,-0.299) (-2.812,-0.264) (-2.813,-0.229) (-2.813,-0.194) (-2.813,-0.158) (-2.813,-0.123) (-2.813,-0.088) (-2.813,-0.053) (-2.813,-0.018) (-2.813,0.018) (-2.813,0.053) (-2.813,0.088) (-2.813,0.123) (-2.813,0.158) (-2.813,0.194) (-2.813,0.229) (-2.812,0.264) (-2.812,0.299) (-2.812,0.334) (-2.812,0.370) (-2.811,0.405) (-2.811,0.440) (-2.811,0.475) (-2.810,0.510) (-2.810,0.546) (-2.809,0.581) (-2.809,0.616) (-2.809,0.651) (-2.808,0.686) (-2.808,0.722) (-2.807,0.757) (-2.807,0.792) (-2.806,0.827) (-2.805,0.862) (-2.805,0.898) (-2.804,0.933) (-2.804,0.968) (-2.803,1.003) (-2.802,1.038) (-2.802,1.074) (-2.801,1.109) (-2.801,1.144) (-2.800,1.179) (-2.799,1.214) (-2.799,1.250) (-2.798,1.285) (-2.797,1.320) (-2.797,1.355) (-2.796,1.390) (-2.795,1.426) (-2.795,1.461) (-2.794,1.496) (-2.793,1.531) (-2.793,1.566) (-2.792,1.602) (-2.791,1.637) (-2.791,1.672) (-2.790,1.707) (-2.789,1.742) (-2.789,1.778) (-2.788,1.813) (-2.787,1.848) (-2.787,1.883) (-2.786,1.918) (-2.785,1.954) (-2.785,1.989) (-2.784,2.024) (-2.784,2.059) (-2.783,2.094) (-2.782,2.130) (-2.782,2.165) (-2.781,2.200)};
\addplot[lineaD, flechaen=0.142, flechaen=0.858] coordinates {(-3.000,0.196) (-2.965,0.196) (-2.930,0.196) (-2.895,0.196) (-2.860,0.196) (-2.824,0.196) (-2.789,0.196) (-2.754,0.195) (-2.719,0.195) (-2.684,0.195) (-2.649,0.195) (-2.614,0.195) (-2.579,0.194) (-2.544,0.194) (-2.509,0.194) (-2.473,0.194) (-2.438,0.193) (-2.403,0.193) (-2.368,0.193) (-2.333,0.193) (-2.298,0.192) (-2.263,0.192) (-2.228,0.192) (-2.193,0.191) (-2.158,0.191) (-2.122,0.190) (-2.087,0.190) (-2.052,0.189) (-2.017,0.189) (-1.982,0.188) (-1.947,0.188) (-1.912,0.187) (-1.877,0.187) (-1.842,0.186) (-1.806,0.185) (-1.771,0.184) (-1.736,0.183) (-1.701,0.183) (-1.666,0.182) (-1.631,0.181) (-1.596,0.179) (-1.561,0.178) (-1.526,0.177) (-1.491,0.176) (-1.456,0.174) (-1.421,0.173) (-1.385,0.171) (-1.350,0.169) (-1.315,0.167) (-1.280,0.165) (-1.245,0.163) (-1.210,0.161) (-1.175,0.158) (-1.140,0.155) (-1.105,0.152) (-1.070,0.149) (-1.035,0.146) (-1.000,0.143) (-0.965,0.141) (-0.930,0.141) (-0.895,0.141) (-0.860,0.141) (-0.825,0.141) (-0.790,0.141) (-0.755,0.141) (-0.720,0.141) (-0.685,0.141) (-0.649,0.141) (-0.614,0.141) (-0.579,0.141) (-0.544,0.141) (-0.509,0.141) (-0.474,0.141) (-0.439,0.141) (-0.404,0.141) (-0.369,0.141) (-0.333,0.141) (-0.298,0.141) (-0.263,0.141) (-0.228,0.141) (-0.193,0.141) (-0.158,0.141) (-0.123,0.141) (-0.088,0.141) (-0.053,0.141) (-0.018,0.141) (0.018,0.141) (0.053,0.141) (0.088,0.141) (0.123,0.141) (0.158,0.141) (0.193,0.141) (0.228,0.141) (0.263,0.141) (0.298,0.141) (0.333,0.141) (0.369,0.141) (0.404,0.141) (0.439,0.141) (0.474,0.141) (0.509,0.141) (0.544,0.141) (0.579,0.141) (0.614,0.141) (0.649,0.141) (0.685,0.141) (0.720,0.141) (0.755,0.141) (0.790,0.141) (0.825,0.141) (0.860,0.141) (0.895,0.141) (0.930,0.141) (0.965,0.141) (1.000,0.143) (1.035,0.146) (1.070,0.149) (1.105,0.152) (1.140,0.155) (1.175,0.158) (1.210,0.161) (1.245,0.163) (1.280,0.165) (1.315,0.167) (1.350,0.169) (1.385,0.171) (1.421,0.173) (1.456,0.174) (1.491,0.176) (1.526,0.177) (1.561,0.178) (1.596,0.179) (1.631,0.181) (1.666,0.182) (1.701,0.183) (1.736,0.183) (1.771,0.184) (1.806,0.185) (1.842,0.186) (1.877,0.187) (1.912,0.187) (1.947,0.188) (1.982,0.188) (2.017,0.189) (2.052,0.189) (2.087,0.190) (2.122,0.190) (2.158,0.191) (2.193,0.191) (2.228,0.192) (2.263,0.192) (2.298,0.192) (2.333,0.193) (2.368,0.193) (2.403,0.193) (2.438,0.193) (2.473,0.194) (2.509,0.194) (2.544,0.194) (2.579,0.194) (2.614,0.195) (2.649,0.195) (2.684,0.195) (2.719,0.195) (2.754,0.195) (2.789,0.196) (2.824,0.196) (2.860,0.196) (2.895,0.196) (2.930,0.196) (2.965,0.196) (3.000,0.196)};
\addplot[lineaD, flechaen=0.141, flechaen=0.859] coordinates {(-3.000,0.393) (-2.965,0.393) (-2.930,0.392) (-2.895,0.392) (-2.859,0.392) (-2.824,0.392) (-2.789,0.391) (-2.754,0.391) (-2.719,0.391) (-2.684,0.390) (-2.648,0.390) (-2.613,0.390) (-2.578,0.389) (-2.543,0.389) (-2.508,0.388) (-2.473,0.388) (-2.438,0.387) (-2.402,0.387) (-2.367,0.386) (-2.332,0.386) (-2.297,0.385) (-2.262,0.384) (-2.227,0.384) (-2.191,0.383) (-2.156,0.382) (-2.121,0.381) (-2.086,0.381) (-2.051,0.380) (-2.016,0.379) (-1.981,0.378) (-1.945,0.377) (-1.910,0.375) (-1.875,0.374) (-1.840,0.373) (-1.805,0.371) (-1.770,0.370) (-1.735,0.368) (-1.700,0.367) (-1.664,0.365) (-1.629,0.363) (-1.594,0.361) (-1.559,0.359) (-1.524,0.356) (-1.489,0.354) (-1.454,0.351) (-1.419,0.348) (-1.384,0.345) (-1.349,0.342) (-1.314,0.338) (-1.279,0.334) (-1.244,0.330) (-1.209,0.326) (-1.174,0.321) (-1.140,0.316) (-1.105,0.310) (-1.070,0.304) (-1.036,0.298) (-1.001,0.292) (-0.967,0.284) (-0.932,0.283) (-0.896,0.283) (-0.861,0.283) (-0.826,0.283) (-0.791,0.283) (-0.756,0.283) (-0.721,0.283) (-0.686,0.283) (-0.650,0.283) (-0.615,0.283) (-0.580,0.283) (-0.545,0.283) (-0.510,0.283) (-0.475,0.283) (-0.439,0.283) (-0.404,0.283) (-0.369,0.283) (-0.334,0.283) (-0.299,0.283) (-0.264,0.283) (-0.229,0.283) (-0.193,0.283) (-0.158,0.283) (-0.123,0.283) (-0.088,0.283) (-0.053,0.283) (-0.018,0.283) (0.018,0.283) (0.053,0.283) (0.088,0.283) (0.123,0.283) (0.158,0.283) (0.193,0.283) (0.229,0.283) (0.264,0.283) (0.299,0.283) (0.334,0.283) (0.369,0.283) (0.404,0.283) (0.439,0.283) (0.475,0.283) (0.510,0.283) (0.545,0.283) (0.580,0.283) (0.615,0.283) (0.650,0.283) (0.686,0.283) (0.721,0.283) (0.756,0.283) (0.791,0.283) (0.826,0.283) (0.861,0.283) (0.896,0.283) (0.932,0.283) (0.967,0.284) (1.001,0.292) (1.036,0.298) (1.070,0.304) (1.105,0.310) (1.140,0.316) (1.174,0.321) (1.209,0.326) (1.244,0.330) (1.279,0.334) (1.314,0.338) (1.349,0.342) (1.384,0.345) (1.419,0.348) (1.454,0.351) (1.489,0.354) (1.524,0.356) (1.559,0.359) (1.594,0.361) (1.629,0.363) (1.664,0.365) (1.700,0.367) (1.735,0.368) (1.770,0.370) (1.805,0.371) (1.840,0.373) (1.875,0.374) (1.910,0.375) (1.945,0.377) (1.981,0.378) (2.016,0.379) (2.051,0.380) (2.086,0.381) (2.121,0.381) (2.156,0.382) (2.191,0.383) (2.227,0.384) (2.262,0.384) (2.297,0.385) (2.332,0.386) (2.367,0.386) (2.402,0.387) (2.438,0.387) (2.473,0.388) (2.508,0.388) (2.543,0.389) (2.578,0.389) (2.613,0.390) (2.648,0.390) (2.684,0.390) (2.719,0.391) (2.754,0.391) (2.789,0.391) (2.824,0.392) (2.859,0.392) (2.895,0.392) (2.930,0.392) (2.965,0.393) (3.000,0.393)};
\addplot[lineaD, flechaen=0.141, flechaen=0.859] coordinates {(-3.000,0.590) (-2.965,0.589) (-2.930,0.589) (-2.895,0.589) (-2.860,0.588) (-2.825,0.588) (-2.790,0.587) (-2.755,0.587) (-2.720,0.587) (-2.685,0.586) (-2.650,0.586) (-2.615,0.585) (-2.580,0.584) (-2.545,0.584) (-2.510,0.583) (-2.474,0.583) (-2.439,0.582) (-2.404,0.581) (-2.369,0.580) (-2.334,0.580) (-2.299,0.579) (-2.264,0.578) (-2.229,0.577) (-2.194,0.576) (-2.159,0.575) (-2.124,0.574) (-2.089,0.572) (-2.054,0.571) (-2.019,0.570) (-1.984,0.569) (-1.949,0.567) (-1.914,0.565) (-1.879,0.564) (-1.844,0.562) (-1.809,0.560) (-1.774,0.558) (-1.739,0.556) (-1.704,0.554) (-1.669,0.551) (-1.634,0.549) (-1.599,0.546) (-1.565,0.543) (-1.530,0.540) (-1.495,0.537) (-1.460,0.533) (-1.425,0.529) (-1.390,0.525) (-1.356,0.521) (-1.321,0.516) (-1.286,0.511) (-1.251,0.506) (-1.217,0.500) (-1.182,0.494) (-1.148,0.487) (-1.114,0.480) (-1.080,0.473) (-1.045,0.465) (-1.011,0.456) (-0.978,0.447) (-0.944,0.437) (-0.911,0.426) (-0.876,0.424) (-0.841,0.424) (-0.806,0.424) (-0.771,0.424) (-0.736,0.424) (-0.701,0.424) (-0.666,0.424) (-0.631,0.424) (-0.596,0.424) (-0.561,0.424) (-0.526,0.424) (-0.491,0.424) (-0.455,0.424) (-0.420,0.424) (-0.385,0.424) (-0.350,0.424) (-0.315,0.424) (-0.280,0.424) (-0.245,0.424) (-0.210,0.424) (-0.175,0.424) (-0.140,0.424) (-0.105,0.424) (-0.070,0.424) (-0.035,0.424) (0.000,0.424) (0.035,0.424) (0.070,0.424) (0.105,0.424) (0.140,0.424) (0.175,0.424) (0.210,0.424) (0.245,0.424) (0.280,0.424) (0.315,0.424) (0.350,0.424) (0.385,0.424) (0.420,0.424) (0.455,0.424) (0.491,0.424) (0.526,0.424) (0.561,0.424) (0.596,0.424) (0.631,0.424) (0.666,0.424) (0.701,0.424) (0.736,0.424) (0.771,0.424) (0.806,0.424) (0.841,0.424) (0.876,0.424) (0.911,0.426) (0.944,0.437) (0.978,0.447) (1.011,0.456) (1.045,0.465) (1.080,0.473) (1.114,0.480) (1.148,0.487) (1.182,0.494) (1.217,0.500) (1.251,0.506) (1.286,0.511) (1.321,0.516) (1.356,0.521) (1.390,0.525) (1.425,0.529) (1.460,0.533) (1.495,0.537) (1.530,0.540) (1.565,0.543) (1.599,0.546) (1.634,0.549) (1.669,0.551) (1.704,0.554) (1.739,0.556) (1.774,0.558) (1.809,0.560) (1.844,0.562) (1.879,0.564) (1.914,0.565) (1.949,0.567) (1.984,0.569) (2.019,0.570) (2.054,0.571) (2.089,0.572) (2.124,0.574) (2.159,0.575) (2.194,0.576) (2.229,0.577) (2.264,0.578) (2.299,0.579) (2.334,0.580) (2.369,0.580) (2.404,0.581) (2.439,0.582) (2.474,0.583) (2.510,0.583) (2.545,0.584) (2.580,0.584) (2.615,0.585) (2.650,0.586) (2.685,0.586) (2.720,0.587) (2.755,0.587) (2.790,0.587) (2.825,0.588) (2.860,0.588) (2.895,0.589) (2.930,0.589) (2.965,0.589) (3.000,0.590)};
\addplot[lineaD, flechaen=0.141, flechaen=0.859] coordinates {(-3.000,0.787) (-2.965,0.786) (-2.930,0.786) (-2.895,0.786) (-2.859,0.785) (-2.824,0.785) (-2.789,0.784) (-2.754,0.783) (-2.719,0.783) (-2.684,0.782) (-2.649,0.782) (-2.613,0.781) (-2.578,0.780) (-2.543,0.780) (-2.508,0.779) (-2.473,0.778) (-2.438,0.777) (-2.402,0.776) (-2.367,0.775) (-2.332,0.774) (-2.297,0.773) (-2.262,0.772) (-2.227,0.771) (-2.192,0.770) (-2.157,0.769) (-2.121,0.767) (-2.086,0.766) (-2.051,0.764) (-2.016,0.763) (-1.981,0.761) (-1.946,0.759) (-1.911,0.758) (-1.876,0.756) (-1.841,0.754) (-1.805,0.751) (-1.770,0.749) (-1.735,0.747) (-1.700,0.744) (-1.665,0.741) (-1.630,0.738) (-1.595,0.735) (-1.560,0.732) (-1.525,0.728) (-1.490,0.724) (-1.455,0.720) (-1.420,0.716) (-1.386,0.711) (-1.351,0.706) (-1.316,0.701) (-1.281,0.696) (-1.247,0.690) (-1.212,0.683) (-1.178,0.676) (-1.143,0.669) (-1.109,0.661) (-1.075,0.653) (-1.041,0.643) (-1.007,0.634) (-0.974,0.623) (-0.940,0.612) (-0.907,0.600) (-0.874,0.587) (-0.842,0.574) (-0.809,0.566) (-0.773,0.566) (-0.738,0.566) (-0.703,0.566) (-0.668,0.566) (-0.633,0.566) (-0.598,0.566) (-0.562,0.566) (-0.527,0.566) (-0.492,0.566) (-0.457,0.566) (-0.422,0.566) (-0.387,0.566) (-0.352,0.566) (-0.316,0.566) (-0.281,0.566) (-0.246,0.566) (-0.211,0.566) (-0.176,0.566) (-0.141,0.566) (-0.105,0.566) (-0.070,0.566) (-0.035,0.566) (0.000,0.566) (0.035,0.566) (0.070,0.566) (0.105,0.566) (0.141,0.566) (0.176,0.566) (0.211,0.566) (0.246,0.566) (0.281,0.566) (0.316,0.566) (0.352,0.566) (0.387,0.566) (0.422,0.566) (0.457,0.566) (0.492,0.566) (0.527,0.566) (0.562,0.566) (0.598,0.566) (0.633,0.566) (0.668,0.566) (0.703,0.566) (0.738,0.566) (0.773,0.566) (0.809,0.566) (0.842,0.574) (0.874,0.587) (0.907,0.600) (0.940,0.612) (0.974,0.623) (1.007,0.634) (1.041,0.643) (1.075,0.653) (1.109,0.661) (1.143,0.669) (1.178,0.676) (1.212,0.683) (1.247,0.690) (1.281,0.696) (1.316,0.701) (1.351,0.706) (1.386,0.711) (1.420,0.716) (1.455,0.720) (1.490,0.724) (1.525,0.728) (1.560,0.732) (1.595,0.735) (1.630,0.738) (1.665,0.741) (1.700,0.744) (1.735,0.747) (1.770,0.749) (1.805,0.751) (1.841,0.754) (1.876,0.756) (1.911,0.758) (1.946,0.759) (1.981,0.761) (2.016,0.763) (2.051,0.764) (2.086,0.766) (2.121,0.767) (2.157,0.769) (2.192,0.770) (2.227,0.771) (2.262,0.772) (2.297,0.773) (2.332,0.774) (2.367,0.775) (2.402,0.776) (2.438,0.777) (2.473,0.778) (2.508,0.779) (2.543,0.780) (2.578,0.780) (2.613,0.781) (2.649,0.782) (2.684,0.782) (2.719,0.783) (2.754,0.783) (2.789,0.784) (2.824,0.785) (2.859,0.785) (2.895,0.786) (2.930,0.786) (2.965,0.786) (3.000,0.787)};
\addplot[lineaD, flechaen=0.140, flechaen=0.860] coordinates {(-3.000,0.984) (-2.965,0.984) (-2.930,0.983) (-2.895,0.983) (-2.860,0.982) (-2.825,0.982) (-2.789,0.981) (-2.754,0.981) (-2.719,0.980) (-2.684,0.979) (-2.649,0.979) (-2.614,0.978) (-2.579,0.977) (-2.544,0.976) (-2.509,0.975) (-2.474,0.974) (-2.439,0.973) (-2.404,0.972) (-2.368,0.971) (-2.333,0.970) (-2.298,0.969) (-2.263,0.968) (-2.228,0.967) (-2.193,0.966) (-2.158,0.964) (-2.123,0.963) (-2.088,0.961) (-2.053,0.960) (-2.018,0.958) (-1.983,0.956) (-1.948,0.954) (-1.913,0.952) (-1.878,0.950) (-1.843,0.948) (-1.808,0.946) (-1.773,0.943) (-1.738,0.941) (-1.703,0.938) (-1.668,0.935) (-1.633,0.932) (-1.598,0.929) (-1.563,0.926) (-1.528,0.922) (-1.493,0.918) (-1.458,0.914) (-1.423,0.910) (-1.389,0.906) (-1.354,0.901) (-1.319,0.896) (-1.284,0.890) (-1.250,0.884) (-1.215,0.878) (-1.181,0.872) (-1.146,0.865) (-1.112,0.857) (-1.078,0.849) (-1.044,0.841) (-1.010,0.831) (-0.976,0.822) (-0.943,0.811) (-0.910,0.800) (-0.877,0.788) (-0.844,0.775) (-0.812,0.761) (-0.780,0.746) (-0.749,0.731) (-0.718,0.713) (-0.684,0.707) (-0.649,0.707) (-0.614,0.707) (-0.579,0.707) (-0.544,0.707) (-0.509,0.707) (-0.474,0.707) (-0.439,0.707) (-0.404,0.707) (-0.368,0.707) (-0.333,0.707) (-0.298,0.707) (-0.263,0.707) (-0.228,0.707) (-0.193,0.707) (-0.158,0.707) (-0.123,0.707) (-0.088,0.707) (-0.053,0.707) (-0.018,0.707) (0.018,0.707) (0.053,0.707) (0.088,0.707) (0.123,0.707) (0.158,0.707) (0.193,0.707) (0.228,0.707) (0.263,0.707) (0.298,0.707) (0.333,0.707) (0.368,0.707) (0.404,0.707) (0.439,0.707) (0.474,0.707) (0.509,0.707) (0.544,0.707) (0.579,0.707) (0.614,0.707) (0.649,0.707) (0.684,0.707) (0.718,0.713) (0.749,0.731) (0.780,0.746) (0.812,0.761) (0.844,0.775) (0.877,0.788) (0.910,0.800) (0.943,0.811) (0.976,0.822) (1.010,0.831) (1.044,0.841) (1.078,0.849) (1.112,0.857) (1.146,0.865) (1.181,0.872) (1.215,0.878) (1.250,0.884) (1.284,0.890) (1.319,0.896) (1.354,0.901) (1.389,0.906) (1.423,0.910) (1.458,0.914) (1.493,0.918) (1.528,0.922) (1.563,0.926) (1.598,0.929) (1.633,0.932) (1.668,0.935) (1.703,0.938) (1.738,0.941) (1.773,0.943) (1.808,0.946) (1.843,0.948) (1.878,0.950) (1.913,0.952) (1.948,0.954) (1.983,0.956) (2.018,0.958) (2.053,0.960) (2.088,0.961) (2.123,0.963) (2.158,0.964) (2.193,0.966) (2.228,0.967) (2.263,0.968) (2.298,0.969) (2.333,0.970) (2.368,0.971) (2.404,0.972) (2.439,0.973) (2.474,0.974) (2.509,0.975) (2.544,0.976) (2.579,0.977) (2.614,0.978) (2.649,0.979) (2.684,0.979) (2.719,0.980) (2.754,0.981) (2.789,0.981) (2.825,0.982) (2.860,0.982) (2.895,0.983) (2.930,0.983) (2.965,0.984) (3.000,0.984)};
\addplot[lineaD, flechaen=0.139, flechaen=0.861] coordinates {(-3.000,1.182) (-2.965,1.182) (-2.930,1.181) (-2.895,1.181) (-2.860,1.180) (-2.825,1.180) (-2.790,1.179) (-2.755,1.178) (-2.720,1.178) (-2.685,1.177) (-2.650,1.176) (-2.615,1.175) (-2.580,1.174) (-2.545,1.174) (-2.510,1.173) (-2.475,1.172) (-2.440,1.171) (-2.404,1.170) (-2.369,1.169) (-2.334,1.168) (-2.299,1.166) (-2.264,1.165) (-2.229,1.164) (-2.194,1.163) (-2.159,1.161) (-2.124,1.160) (-2.089,1.158) (-2.054,1.157) (-2.019,1.155) (-1.984,1.153) (-1.949,1.151) (-1.914,1.149) (-1.879,1.147) (-1.844,1.145) (-1.809,1.143) (-1.774,1.141) (-1.740,1.138) (-1.705,1.136) (-1.670,1.133) (-1.635,1.130) (-1.600,1.127) (-1.565,1.124) (-1.530,1.121) (-1.495,1.117) (-1.460,1.113) (-1.426,1.110) (-1.391,1.105) (-1.356,1.101) (-1.321,1.097) (-1.287,1.092) (-1.252,1.087) (-1.217,1.081) (-1.183,1.076) (-1.148,1.070) (-1.114,1.063) (-1.079,1.056) (-1.045,1.049) (-1.011,1.042) (-0.977,1.034) (-0.943,1.025) (-0.909,1.016) (-0.875,1.006) (-0.842,0.996) (-0.808,0.985) (-0.775,0.974) (-0.742,0.961) (-0.710,0.948) (-0.678,0.934) (-0.646,0.919) (-0.615,0.903) (-0.585,0.885) (-0.555,0.867) (-0.526,0.849) (-0.491,0.849) (-0.456,0.849) (-0.420,0.849) (-0.385,0.849) (-0.350,0.849) (-0.315,0.849) (-0.280,0.849) (-0.245,0.849) (-0.210,0.849) (-0.175,0.849) (-0.140,0.849) (-0.105,0.849) (-0.070,0.849) (-0.035,0.849) (0.000,0.849) (0.035,0.849) (0.070,0.849) (0.105,0.849) (0.140,0.849) (0.175,0.849) (0.210,0.849) (0.245,0.849) (0.280,0.849) (0.315,0.849) (0.350,0.849) (0.385,0.849) (0.420,0.849) (0.456,0.849) (0.491,0.849) (0.526,0.849) (0.555,0.867) (0.585,0.885) (0.615,0.903) (0.646,0.919) (0.678,0.934) (0.710,0.948) (0.742,0.961) (0.775,0.974) (0.808,0.985) (0.842,0.996) (0.875,1.006) (0.909,1.016) (0.943,1.025) (0.977,1.034) (1.011,1.042) (1.045,1.049) (1.079,1.056) (1.114,1.063) (1.148,1.070) (1.183,1.076) (1.217,1.081) (1.252,1.087) (1.287,1.092) (1.321,1.097) (1.356,1.101) (1.391,1.105) (1.426,1.110) (1.460,1.113) (1.495,1.117) (1.530,1.121) (1.565,1.124) (1.600,1.127) (1.635,1.130) (1.670,1.133) (1.705,1.136) (1.740,1.138) (1.774,1.141) (1.809,1.143) (1.844,1.145) (1.879,1.147) (1.914,1.149) (1.949,1.151) (1.984,1.153) (2.019,1.155) (2.054,1.157) (2.089,1.158) (2.124,1.160) (2.159,1.161) (2.194,1.163) (2.229,1.164) (2.264,1.165) (2.299,1.166) (2.334,1.168) (2.369,1.169) (2.404,1.170) (2.440,1.171) (2.475,1.172) (2.510,1.173) (2.545,1.174) (2.580,1.174) (2.615,1.175) (2.650,1.176) (2.685,1.177) (2.720,1.178) (2.755,1.178) (2.790,1.179) (2.825,1.180) (2.860,1.180) (2.895,1.181) (2.930,1.181) (2.965,1.182) (3.000,1.182)};
\addplot[lineaD, flechaen=0.139, flechaen=0.861] coordinates {(-3.000,1.381) (-2.965,1.380) (-2.930,1.380) (-2.895,1.379) (-2.860,1.379) (-2.825,1.378) (-2.790,1.377) (-2.755,1.377) (-2.720,1.376) (-2.685,1.375) (-2.650,1.374) (-2.615,1.374) (-2.580,1.373) (-2.545,1.372) (-2.510,1.371) (-2.475,1.370) (-2.439,1.369) (-2.404,1.368) (-2.369,1.367) (-2.334,1.366) (-2.299,1.365) (-2.264,1.364) (-2.229,1.362) (-2.194,1.361) (-2.159,1.360) (-2.124,1.358) (-2.089,1.357) (-2.054,1.355) (-2.019,1.354) (-1.984,1.352) (-1.949,1.350) (-1.914,1.348) (-1.879,1.347) (-1.844,1.345) (-1.809,1.343) (-1.774,1.340) (-1.739,1.338) (-1.704,1.336) (-1.669,1.333) (-1.635,1.331) (-1.600,1.328) (-1.565,1.325) (-1.530,1.322) (-1.495,1.319) (-1.460,1.316) (-1.425,1.313) (-1.390,1.309) (-1.355,1.305) (-1.321,1.302) (-1.286,1.298) (-1.251,1.293) (-1.216,1.289) (-1.181,1.284) (-1.147,1.280) (-1.112,1.274) (-1.077,1.269) (-1.043,1.264) (-1.008,1.258) (-0.974,1.252) (-0.939,1.246) (-0.905,1.239) (-0.871,1.232) (-0.836,1.225) (-0.802,1.217) (-0.768,1.209) (-0.734,1.201) (-0.700,1.193) (-0.666,1.184) (-0.632,1.174) (-0.599,1.164) (-0.565,1.154) (-0.532,1.144) (-0.498,1.133) (-0.465,1.121) (-0.432,1.109) (-0.399,1.097) (-0.367,1.085) (-0.334,1.071) (-0.302,1.058) (-0.270,1.044) (-0.237,1.031) (-0.205,1.017) (-0.173,1.003) (-0.140,0.991) (-0.105,0.990) (-0.070,0.990) (-0.035,0.990) (-0.000,0.990) (0.035,0.990) (0.070,0.990) (0.105,0.990) (0.140,0.991) (0.173,1.003) (0.205,1.017) (0.237,1.031) (0.270,1.044) (0.302,1.058) (0.334,1.071) (0.367,1.085) (0.399,1.097) (0.432,1.109) (0.465,1.121) (0.498,1.133) (0.532,1.144) (0.565,1.154) (0.599,1.164) (0.632,1.174) (0.666,1.184) (0.700,1.193) (0.734,1.201) (0.768,1.209) (0.802,1.217) (0.836,1.225) (0.871,1.232) (0.905,1.239) (0.939,1.246) (0.974,1.252) (1.008,1.258) (1.043,1.264) (1.077,1.269) (1.112,1.274) (1.147,1.280) (1.181,1.284) (1.216,1.289) (1.251,1.293) (1.286,1.298) (1.321,1.302) (1.355,1.305) (1.390,1.309) (1.425,1.313) (1.460,1.316) (1.495,1.319) (1.530,1.322) (1.565,1.325) (1.600,1.328) (1.635,1.331) (1.669,1.333) (1.704,1.336) (1.739,1.338) (1.774,1.340) (1.809,1.343) (1.844,1.345) (1.879,1.347) (1.914,1.348) (1.949,1.350) (1.984,1.352) (2.019,1.354) (2.054,1.355) (2.089,1.357) (2.124,1.358) (2.159,1.360) (2.194,1.361) (2.229,1.362) (2.264,1.364) (2.299,1.365) (2.334,1.366) (2.369,1.367) (2.404,1.368) (2.439,1.369) (2.475,1.370) (2.510,1.371) (2.545,1.372) (2.580,1.373) (2.615,1.374) (2.650,1.374) (2.685,1.375) (2.720,1.376) (2.755,1.377) (2.790,1.377) (2.825,1.378) (2.860,1.379) (2.895,1.379) (2.930,1.380) (2.965,1.380) (3.000,1.381)};
\addplot[lineaD, flechaen=0.141, flechaen=0.859] coordinates {(-3.000,1.580) (-2.965,1.579) (-2.930,1.579) (-2.895,1.578) (-2.860,1.578) (-2.825,1.577) (-2.790,1.576) (-2.755,1.576) (-2.720,1.575) (-2.685,1.574) (-2.650,1.573) (-2.615,1.573) (-2.580,1.572) (-2.545,1.571) (-2.510,1.570) (-2.475,1.569) (-2.440,1.568) (-2.405,1.567) (-2.370,1.566) (-2.335,1.565) (-2.300,1.564) (-2.265,1.563) (-2.230,1.562) (-2.195,1.561) (-2.160,1.559) (-2.125,1.558) (-2.090,1.557) (-2.055,1.555) (-2.020,1.554) (-1.985,1.552) (-1.950,1.551) (-1.915,1.549) (-1.880,1.547) (-1.845,1.546) (-1.810,1.544) (-1.775,1.542) (-1.740,1.540) (-1.705,1.538) (-1.670,1.536) (-1.635,1.533) (-1.600,1.531) (-1.565,1.529) (-1.530,1.526) (-1.495,1.524) (-1.460,1.521) (-1.425,1.518) (-1.390,1.515) (-1.355,1.512) (-1.320,1.509) (-1.286,1.506) (-1.251,1.503) (-1.216,1.499) (-1.181,1.496) (-1.146,1.492) (-1.111,1.488) (-1.077,1.484) (-1.042,1.480) (-1.007,1.476) (-0.972,1.471) (-0.937,1.467) (-0.903,1.462) (-0.868,1.458) (-0.833,1.453) (-0.799,1.448) (-0.764,1.443) (-0.729,1.438) (-0.695,1.432) (-0.660,1.427) (-0.626,1.421) (-0.591,1.416) (-0.556,1.410) (-0.522,1.405) (-0.487,1.399) (-0.453,1.394) (-0.418,1.388) (-0.383,1.383) (-0.349,1.378) (-0.314,1.373) (-0.279,1.368) (-0.245,1.364) (-0.210,1.360) (-0.175,1.357) (-0.140,1.354) (-0.105,1.352) (-0.070,1.350) (-0.035,1.349) (0.000,1.348) (0.035,1.349) (0.070,1.350) (0.105,1.352) (0.140,1.354) (0.175,1.357) (0.210,1.360) (0.245,1.364) (0.279,1.368) (0.314,1.373) (0.349,1.378) (0.383,1.383) (0.418,1.388) (0.453,1.394) (0.487,1.399) (0.522,1.405) (0.556,1.410) (0.591,1.416) (0.626,1.421) (0.660,1.427) (0.695,1.432) (0.729,1.438) (0.764,1.443) (0.799,1.448) (0.833,1.453) (0.868,1.458) (0.903,1.462) (0.937,1.467) (0.972,1.471) (1.007,1.476) (1.042,1.480) (1.077,1.484) (1.111,1.488) (1.146,1.492) (1.181,1.496) (1.216,1.499) (1.251,1.503) (1.286,1.506) (1.320,1.509) (1.355,1.512) (1.390,1.515) (1.425,1.518) (1.460,1.521) (1.495,1.524) (1.530,1.526) (1.565,1.529) (1.600,1.531) (1.635,1.533) (1.670,1.536) (1.705,1.538) (1.740,1.540) (1.775,1.542) (1.810,1.544) (1.845,1.546) (1.880,1.547) (1.915,1.549) (1.950,1.551) (1.985,1.552) (2.020,1.554) (2.055,1.555) (2.090,1.557) (2.125,1.558) (2.160,1.559) (2.195,1.561) (2.230,1.562) (2.265,1.563) (2.300,1.564) (2.335,1.565) (2.370,1.566) (2.405,1.567) (2.440,1.568) (2.475,1.569) (2.510,1.570) (2.545,1.571) (2.580,1.572) (2.615,1.573) (2.650,1.573) (2.685,1.574) (2.720,1.575) (2.755,1.576) (2.790,1.576) (2.825,1.577) (2.860,1.578) (2.895,1.578) (2.930,1.579) (2.965,1.579) (3.000,1.580)};
\addplot[lineaD, flechaen=0.141, flechaen=0.859] coordinates {(-3.000,1.779) (-2.965,1.779) (-2.930,1.778) (-2.895,1.777) (-2.859,1.777) (-2.824,1.776) (-2.789,1.776) (-2.754,1.775) (-2.719,1.774) (-2.684,1.774) (-2.649,1.773) (-2.613,1.772) (-2.578,1.771) (-2.543,1.770) (-2.508,1.770) (-2.473,1.769) (-2.438,1.768) (-2.403,1.767) (-2.367,1.766) (-2.332,1.765) (-2.297,1.764) (-2.262,1.763) (-2.227,1.762) (-2.192,1.761) (-2.157,1.760) (-2.121,1.758) (-2.086,1.757) (-2.051,1.756) (-2.016,1.755) (-1.981,1.753) (-1.946,1.752) (-1.911,1.750) (-1.876,1.749) (-1.841,1.747) (-1.805,1.746) (-1.770,1.744) (-1.735,1.742) (-1.700,1.741) (-1.665,1.739) (-1.630,1.737) (-1.595,1.735) (-1.560,1.733) (-1.525,1.731) (-1.489,1.729) (-1.454,1.727) (-1.419,1.724) (-1.384,1.722) (-1.349,1.720) (-1.314,1.717) (-1.279,1.715) (-1.244,1.712) (-1.209,1.709) (-1.174,1.707) (-1.139,1.704) (-1.104,1.701) (-1.069,1.698) (-1.034,1.695) (-0.999,1.692) (-0.964,1.689) (-0.929,1.686) (-0.894,1.682) (-0.859,1.679) (-0.824,1.676) (-0.789,1.673) (-0.754,1.669) (-0.719,1.666) (-0.684,1.663) (-0.649,1.659) (-0.614,1.656) (-0.579,1.653) (-0.544,1.649) (-0.509,1.646) (-0.474,1.643) (-0.439,1.640) (-0.404,1.637) (-0.369,1.635) (-0.334,1.632) (-0.299,1.630) (-0.263,1.628) (-0.228,1.626) (-0.193,1.624) (-0.158,1.622) (-0.123,1.621) (-0.088,1.620) (-0.053,1.620) (-0.018,1.619) (0.018,1.619) (0.053,1.620) (0.088,1.620) (0.123,1.621) (0.158,1.622) (0.193,1.624) (0.228,1.626) (0.263,1.628) (0.299,1.630) (0.334,1.632) (0.369,1.635) (0.404,1.637) (0.439,1.640) (0.474,1.643) (0.509,1.646) (0.544,1.649) (0.579,1.653) (0.614,1.656) (0.649,1.659) (0.684,1.663) (0.719,1.666) (0.754,1.669) (0.789,1.673) (0.824,1.676) (0.859,1.679) (0.894,1.682) (0.929,1.686) (0.964,1.689) (0.999,1.692) (1.034,1.695) (1.069,1.698) (1.104,1.701) (1.139,1.704) (1.174,1.707) (1.209,1.709) (1.244,1.712) (1.279,1.715) (1.314,1.717) (1.349,1.720) (1.384,1.722) (1.419,1.724) (1.454,1.727) (1.489,1.729) (1.525,1.731) (1.560,1.733) (1.595,1.735) (1.630,1.737) (1.665,1.739) (1.700,1.741) (1.735,1.742) (1.770,1.744) (1.805,1.746) (1.841,1.747) (1.876,1.749) (1.911,1.750) (1.946,1.752) (1.981,1.753) (2.016,1.755) (2.051,1.756) (2.086,1.757) (2.121,1.758) (2.157,1.760) (2.192,1.761) (2.227,1.762) (2.262,1.763) (2.297,1.764) (2.332,1.765) (2.367,1.766) (2.403,1.767) (2.438,1.768) (2.473,1.769) (2.508,1.770) (2.543,1.770) (2.578,1.771) (2.613,1.772) (2.649,1.773) (2.684,1.774) (2.719,1.774) (2.754,1.775) (2.789,1.776) (2.824,1.776) (2.859,1.777) (2.895,1.777) (2.930,1.778) (2.965,1.779) (3.000,1.779)};
\addplot[lineaD, flechaen=0.142, flechaen=0.858] coordinates {(-3.000,1.979) (-2.965,1.978) (-2.930,1.978) (-2.895,1.977) (-2.860,1.977) (-2.824,1.976) (-2.789,1.975) (-2.754,1.975) (-2.719,1.974) (-2.684,1.973) (-2.649,1.973) (-2.614,1.972) (-2.579,1.971) (-2.544,1.971) (-2.508,1.970) (-2.473,1.969) (-2.438,1.968) (-2.403,1.967) (-2.368,1.967) (-2.333,1.966) (-2.298,1.965) (-2.263,1.964) (-2.228,1.963) (-2.192,1.962) (-2.157,1.961) (-2.122,1.960) (-2.087,1.959) (-2.052,1.958) (-2.017,1.956) (-1.982,1.955) (-1.947,1.954) (-1.912,1.953) (-1.877,1.951) (-1.841,1.950) (-1.806,1.949) (-1.771,1.947) (-1.736,1.946) (-1.701,1.944) (-1.666,1.943) (-1.631,1.941) (-1.596,1.940) (-1.561,1.938) (-1.526,1.936) (-1.491,1.935) (-1.455,1.933) (-1.420,1.931) (-1.385,1.929) (-1.350,1.927) (-1.315,1.926) (-1.280,1.924) (-1.245,1.922) (-1.210,1.920) (-1.175,1.917) (-1.140,1.915) (-1.105,1.913) (-1.070,1.911) (-1.035,1.909) (-1.000,1.907) (-0.965,1.905) (-0.930,1.902) (-0.895,1.900) (-0.859,1.898) (-0.824,1.896) (-0.789,1.893) (-0.754,1.891) (-0.719,1.889) (-0.684,1.887) (-0.649,1.885) (-0.614,1.882) (-0.579,1.880) (-0.544,1.878) (-0.509,1.876) (-0.474,1.875) (-0.439,1.873) (-0.404,1.871) (-0.369,1.869) (-0.334,1.868) (-0.298,1.867) (-0.263,1.865) (-0.228,1.864) (-0.193,1.863) (-0.158,1.862) (-0.123,1.862) (-0.088,1.861) (-0.053,1.861) (-0.018,1.861) (0.018,1.861) (0.053,1.861) (0.088,1.861) (0.123,1.862) (0.158,1.862) (0.193,1.863) (0.228,1.864) (0.263,1.865) (0.298,1.867) (0.334,1.868) (0.369,1.869) (0.404,1.871) (0.439,1.873) (0.474,1.875) (0.509,1.876) (0.544,1.878) (0.579,1.880) (0.614,1.882) (0.649,1.885) (0.684,1.887) (0.719,1.889) (0.754,1.891) (0.789,1.893) (0.824,1.896) (0.859,1.898) (0.895,1.900) (0.930,1.902) (0.965,1.905) (1.000,1.907) (1.035,1.909) (1.070,1.911) (1.105,1.913) (1.140,1.915) (1.175,1.917) (1.210,1.920) (1.245,1.922) (1.280,1.924) (1.315,1.926) (1.350,1.927) (1.385,1.929) (1.420,1.931) (1.455,1.933) (1.491,1.935) (1.526,1.936) (1.561,1.938) (1.596,1.940) (1.631,1.941) (1.666,1.943) (1.701,1.944) (1.736,1.946) (1.771,1.947) (1.806,1.949) (1.841,1.950) (1.877,1.951) (1.912,1.953) (1.947,1.954) (1.982,1.955) (2.017,1.956) (2.052,1.958) (2.087,1.959) (2.122,1.960) (2.157,1.961) (2.192,1.962) (2.228,1.963) (2.263,1.964) (2.298,1.965) (2.333,1.966) (2.368,1.967) (2.403,1.967) (2.438,1.968) (2.473,1.969) (2.508,1.970) (2.544,1.971) (2.579,1.971) (2.614,1.972) (2.649,1.973) (2.684,1.973) (2.719,1.974) (2.754,1.975) (2.789,1.975) (2.824,1.976) (2.860,1.977) (2.895,1.977) (2.930,1.978) (2.965,1.978) (3.000,1.979)};
\addplot[lineaD, flechaen=0.142, flechaen=0.858] coordinates {(-3.000,2.179) (-2.965,2.178) (-2.930,2.178) (-2.895,2.177) (-2.860,2.177) (-2.824,2.176) (-2.789,2.176) (-2.754,2.175) (-2.719,2.174) (-2.684,2.174) (-2.649,2.173) (-2.614,2.173) (-2.579,2.172) (-2.544,2.171) (-2.509,2.170) (-2.473,2.170) (-2.438,2.169) (-2.403,2.168) (-2.368,2.167) (-2.333,2.167) (-2.298,2.166) (-2.263,2.165) (-2.228,2.164) (-2.193,2.163) (-2.158,2.162) (-2.123,2.161) (-2.087,2.161) (-2.052,2.160) (-2.017,2.159) (-1.982,2.158) (-1.947,2.156) (-1.912,2.155) (-1.877,2.154) (-1.842,2.153) (-1.807,2.152) (-1.772,2.151) (-1.737,2.150) (-1.701,2.148) (-1.666,2.147) (-1.631,2.146) (-1.596,2.145) (-1.561,2.143) (-1.526,2.142) (-1.491,2.141) (-1.456,2.139) (-1.421,2.138) (-1.386,2.136) (-1.351,2.135) (-1.316,2.133) (-1.281,2.132) (-1.245,2.130) (-1.210,2.129) (-1.175,2.127) (-1.140,2.126) (-1.105,2.124) (-1.070,2.122) (-1.035,2.121) (-1.000,2.119) (-0.965,2.118) (-0.930,2.116) (-0.895,2.114) (-0.860,2.113) (-0.825,2.111) (-0.790,2.110) (-0.754,2.108) (-0.719,2.107) (-0.684,2.105) (-0.649,2.104) (-0.614,2.102) (-0.579,2.101) (-0.544,2.100) (-0.509,2.098) (-0.474,2.097) (-0.439,2.096) (-0.404,2.095) (-0.369,2.094) (-0.333,2.093) (-0.298,2.092) (-0.263,2.091) (-0.228,2.090) (-0.193,2.090) (-0.158,2.089) (-0.123,2.089) (-0.088,2.089) (-0.053,2.088) (-0.018,2.088) (0.018,2.088) (0.053,2.088) (0.088,2.089) (0.123,2.089) (0.158,2.089) (0.193,2.090) (0.228,2.090) (0.263,2.091) (0.298,2.092) (0.333,2.093) (0.369,2.094) (0.404,2.095) (0.439,2.096) (0.474,2.097) (0.509,2.098) (0.544,2.100) (0.579,2.101) (0.614,2.102) (0.649,2.104) (0.684,2.105) (0.719,2.107) (0.754,2.108) (0.790,2.110) (0.825,2.111) (0.860,2.113) (0.895,2.114) (0.930,2.116) (0.965,2.118) (1.000,2.119) (1.035,2.121) (1.070,2.122) (1.105,2.124) (1.140,2.126) (1.175,2.127) (1.210,2.129) (1.245,2.130) (1.281,2.132) (1.316,2.133) (1.351,2.135) (1.386,2.136) (1.421,2.138) (1.456,2.139) (1.491,2.141) (1.526,2.142) (1.561,2.143) (1.596,2.145) (1.631,2.146) (1.666,2.147) (1.701,2.148) (1.737,2.150) (1.772,2.151) (1.807,2.152) (1.842,2.153) (1.877,2.154) (1.912,2.155) (1.947,2.156) (1.982,2.158) (2.017,2.159) (2.052,2.160) (2.087,2.161) (2.123,2.161) (2.158,2.162) (2.193,2.163) (2.228,2.164) (2.263,2.165) (2.298,2.166) (2.333,2.167) (2.368,2.167) (2.403,2.168) (2.438,2.169) (2.473,2.170) (2.509,2.170) (2.544,2.171) (2.579,2.172) (2.614,2.173) (2.649,2.173) (2.684,2.174) (2.719,2.174) (2.754,2.175) (2.789,2.176) (2.824,2.176) (2.860,2.177) (2.895,2.177) (2.930,2.178) (2.965,2.178) (3.000,2.179)};
\addplot[lineaD, flechaen=0.142, flechaen=0.858] coordinates {(-3.000,-0.196) (-2.965,-0.196) (-2.930,-0.196) (-2.895,-0.196) (-2.860,-0.196) (-2.824,-0.196) (-2.789,-0.196) (-2.754,-0.195) (-2.719,-0.195) (-2.684,-0.195) (-2.649,-0.195) (-2.614,-0.195) (-2.579,-0.194) (-2.544,-0.194) (-2.509,-0.194) (-2.473,-0.194) (-2.438,-0.193) (-2.403,-0.193) (-2.368,-0.193) (-2.333,-0.193) (-2.298,-0.192) (-2.263,-0.192) (-2.228,-0.192) (-2.193,-0.191) (-2.158,-0.191) (-2.122,-0.190) (-2.087,-0.190) (-2.052,-0.189) (-2.017,-0.189) (-1.982,-0.188) (-1.947,-0.188) (-1.912,-0.187) (-1.877,-0.187) (-1.842,-0.186) (-1.806,-0.185) (-1.771,-0.184) (-1.736,-0.183) (-1.701,-0.183) (-1.666,-0.182) (-1.631,-0.181) (-1.596,-0.179) (-1.561,-0.178) (-1.526,-0.177) (-1.491,-0.176) (-1.456,-0.174) (-1.421,-0.173) (-1.385,-0.171) (-1.350,-0.169) (-1.315,-0.167) (-1.280,-0.165) (-1.245,-0.163) (-1.210,-0.161) (-1.175,-0.158) (-1.140,-0.155) (-1.105,-0.152) (-1.070,-0.149) (-1.035,-0.146) (-1.000,-0.143) (-0.965,-0.141) (-0.930,-0.141) (-0.895,-0.141) (-0.860,-0.141) (-0.825,-0.141) (-0.790,-0.141) (-0.755,-0.141) (-0.720,-0.141) (-0.685,-0.141) (-0.649,-0.141) (-0.614,-0.141) (-0.579,-0.141) (-0.544,-0.141) (-0.509,-0.141) (-0.474,-0.141) (-0.439,-0.141) (-0.404,-0.141) (-0.369,-0.141) (-0.333,-0.141) (-0.298,-0.141) (-0.263,-0.141) (-0.228,-0.141) (-0.193,-0.141) (-0.158,-0.141) (-0.123,-0.141) (-0.088,-0.141) (-0.053,-0.141) (-0.018,-0.141) (0.018,-0.141) (0.053,-0.141) (0.088,-0.141) (0.123,-0.141) (0.158,-0.141) (0.193,-0.141) (0.228,-0.141) (0.263,-0.141) (0.298,-0.141) (0.333,-0.141) (0.369,-0.141) (0.404,-0.141) (0.439,-0.141) (0.474,-0.141) (0.509,-0.141) (0.544,-0.141) (0.579,-0.141) (0.614,-0.141) (0.649,-0.141) (0.685,-0.141) (0.720,-0.141) (0.755,-0.141) (0.790,-0.141) (0.825,-0.141) (0.860,-0.141) (0.895,-0.141) (0.930,-0.141) (0.965,-0.141) (1.000,-0.143) (1.035,-0.146) (1.070,-0.149) (1.105,-0.152) (1.140,-0.155) (1.175,-0.158) (1.210,-0.161) (1.245,-0.163) (1.280,-0.165) (1.315,-0.167) (1.350,-0.169) (1.385,-0.171) (1.421,-0.173) (1.456,-0.174) (1.491,-0.176) (1.526,-0.177) (1.561,-0.178) (1.596,-0.179) (1.631,-0.181) (1.666,-0.182) (1.701,-0.183) (1.736,-0.183) (1.771,-0.184) (1.806,-0.185) (1.842,-0.186) (1.877,-0.187) (1.912,-0.187) (1.947,-0.188) (1.982,-0.188) (2.017,-0.189) (2.052,-0.189) (2.087,-0.190) (2.122,-0.190) (2.158,-0.191) (2.193,-0.191) (2.228,-0.192) (2.263,-0.192) (2.298,-0.192) (2.333,-0.193) (2.368,-0.193) (2.403,-0.193) (2.438,-0.193) (2.473,-0.194) (2.509,-0.194) (2.544,-0.194) (2.579,-0.194) (2.614,-0.195) (2.649,-0.195) (2.684,-0.195) (2.719,-0.195) (2.754,-0.195) (2.789,-0.196) (2.824,-0.196) (2.860,-0.196) (2.895,-0.196) (2.930,-0.196) (2.965,-0.196) (3.000,-0.196)};
\addplot[lineaD, flechaen=0.141, flechaen=0.859] coordinates {(-3.000,-0.393) (-2.965,-0.393) (-2.930,-0.392) (-2.895,-0.392) (-2.859,-0.392) (-2.824,-0.392) (-2.789,-0.391) (-2.754,-0.391) (-2.719,-0.391) (-2.684,-0.390) (-2.648,-0.390) (-2.613,-0.390) (-2.578,-0.389) (-2.543,-0.389) (-2.508,-0.388) (-2.473,-0.388) (-2.438,-0.387) (-2.402,-0.387) (-2.367,-0.386) (-2.332,-0.386) (-2.297,-0.385) (-2.262,-0.384) (-2.227,-0.384) (-2.191,-0.383) (-2.156,-0.382) (-2.121,-0.381) (-2.086,-0.381) (-2.051,-0.380) (-2.016,-0.379) (-1.981,-0.378) (-1.945,-0.377) (-1.910,-0.375) (-1.875,-0.374) (-1.840,-0.373) (-1.805,-0.371) (-1.770,-0.370) (-1.735,-0.368) (-1.700,-0.367) (-1.664,-0.365) (-1.629,-0.363) (-1.594,-0.361) (-1.559,-0.359) (-1.524,-0.356) (-1.489,-0.354) (-1.454,-0.351) (-1.419,-0.348) (-1.384,-0.345) (-1.349,-0.342) (-1.314,-0.338) (-1.279,-0.334) (-1.244,-0.330) (-1.209,-0.326) (-1.174,-0.321) (-1.140,-0.316) (-1.105,-0.310) (-1.070,-0.304) (-1.036,-0.298) (-1.001,-0.292) (-0.967,-0.284) (-0.932,-0.283) (-0.896,-0.283) (-0.861,-0.283) (-0.826,-0.283) (-0.791,-0.283) (-0.756,-0.283) (-0.721,-0.283) (-0.686,-0.283) (-0.650,-0.283) (-0.615,-0.283) (-0.580,-0.283) (-0.545,-0.283) (-0.510,-0.283) (-0.475,-0.283) (-0.439,-0.283) (-0.404,-0.283) (-0.369,-0.283) (-0.334,-0.283) (-0.299,-0.283) (-0.264,-0.283) (-0.229,-0.283) (-0.193,-0.283) (-0.158,-0.283) (-0.123,-0.283) (-0.088,-0.283) (-0.053,-0.283) (-0.018,-0.283) (0.018,-0.283) (0.053,-0.283) (0.088,-0.283) (0.123,-0.283) (0.158,-0.283) (0.193,-0.283) (0.229,-0.283) (0.264,-0.283) (0.299,-0.283) (0.334,-0.283) (0.369,-0.283) (0.404,-0.283) (0.439,-0.283) (0.475,-0.283) (0.510,-0.283) (0.545,-0.283) (0.580,-0.283) (0.615,-0.283) (0.650,-0.283) (0.686,-0.283) (0.721,-0.283) (0.756,-0.283) (0.791,-0.283) (0.826,-0.283) (0.861,-0.283) (0.896,-0.283) (0.932,-0.283) (0.967,-0.284) (1.001,-0.292) (1.036,-0.298) (1.070,-0.304) (1.105,-0.310) (1.140,-0.316) (1.174,-0.321) (1.209,-0.326) (1.244,-0.330) (1.279,-0.334) (1.314,-0.338) (1.349,-0.342) (1.384,-0.345) (1.419,-0.348) (1.454,-0.351) (1.489,-0.354) (1.524,-0.356) (1.559,-0.359) (1.594,-0.361) (1.629,-0.363) (1.664,-0.365) (1.700,-0.367) (1.735,-0.368) (1.770,-0.370) (1.805,-0.371) (1.840,-0.373) (1.875,-0.374) (1.910,-0.375) (1.945,-0.377) (1.981,-0.378) (2.016,-0.379) (2.051,-0.380) (2.086,-0.381) (2.121,-0.381) (2.156,-0.382) (2.191,-0.383) (2.227,-0.384) (2.262,-0.384) (2.297,-0.385) (2.332,-0.386) (2.367,-0.386) (2.402,-0.387) (2.438,-0.387) (2.473,-0.388) (2.508,-0.388) (2.543,-0.389) (2.578,-0.389) (2.613,-0.390) (2.648,-0.390) (2.684,-0.390) (2.719,-0.391) (2.754,-0.391) (2.789,-0.391) (2.824,-0.392) (2.859,-0.392) (2.895,-0.392) (2.930,-0.392) (2.965,-0.393) (3.000,-0.393)};
\addplot[lineaD, flechaen=0.141, flechaen=0.859] coordinates {(-3.000,-0.590) (-2.965,-0.589) (-2.930,-0.589) (-2.895,-0.589) (-2.860,-0.588) (-2.825,-0.588) (-2.790,-0.587) (-2.755,-0.587) (-2.720,-0.587) (-2.685,-0.586) (-2.650,-0.586) (-2.615,-0.585) (-2.580,-0.584) (-2.545,-0.584) (-2.510,-0.583) (-2.474,-0.583) (-2.439,-0.582) (-2.404,-0.581) (-2.369,-0.580) (-2.334,-0.580) (-2.299,-0.579) (-2.264,-0.578) (-2.229,-0.577) (-2.194,-0.576) (-2.159,-0.575) (-2.124,-0.574) (-2.089,-0.572) (-2.054,-0.571) (-2.019,-0.570) (-1.984,-0.569) (-1.949,-0.567) (-1.914,-0.565) (-1.879,-0.564) (-1.844,-0.562) (-1.809,-0.560) (-1.774,-0.558) (-1.739,-0.556) (-1.704,-0.554) (-1.669,-0.551) (-1.634,-0.549) (-1.599,-0.546) (-1.565,-0.543) (-1.530,-0.540) (-1.495,-0.537) (-1.460,-0.533) (-1.425,-0.529) (-1.390,-0.525) (-1.356,-0.521) (-1.321,-0.516) (-1.286,-0.511) (-1.251,-0.506) (-1.217,-0.500) (-1.182,-0.494) (-1.148,-0.487) (-1.114,-0.480) (-1.080,-0.473) (-1.045,-0.465) (-1.011,-0.456) (-0.978,-0.447) (-0.944,-0.437) (-0.911,-0.426) (-0.876,-0.424) (-0.841,-0.424) (-0.806,-0.424) (-0.771,-0.424) (-0.736,-0.424) (-0.701,-0.424) (-0.666,-0.424) (-0.631,-0.424) (-0.596,-0.424) (-0.561,-0.424) (-0.526,-0.424) (-0.491,-0.424) (-0.455,-0.424) (-0.420,-0.424) (-0.385,-0.424) (-0.350,-0.424) (-0.315,-0.424) (-0.280,-0.424) (-0.245,-0.424) (-0.210,-0.424) (-0.175,-0.424) (-0.140,-0.424) (-0.105,-0.424) (-0.070,-0.424) (-0.035,-0.424) (0.000,-0.424) (0.035,-0.424) (0.070,-0.424) (0.105,-0.424) (0.140,-0.424) (0.175,-0.424) (0.210,-0.424) (0.245,-0.424) (0.280,-0.424) (0.315,-0.424) (0.350,-0.424) (0.385,-0.424) (0.420,-0.424) (0.455,-0.424) (0.491,-0.424) (0.526,-0.424) (0.561,-0.424) (0.596,-0.424) (0.631,-0.424) (0.666,-0.424) (0.701,-0.424) (0.736,-0.424) (0.771,-0.424) (0.806,-0.424) (0.841,-0.424) (0.876,-0.424) (0.911,-0.426) (0.944,-0.437) (0.978,-0.447) (1.011,-0.456) (1.045,-0.465) (1.080,-0.473) (1.114,-0.480) (1.148,-0.487) (1.182,-0.494) (1.217,-0.500) (1.251,-0.506) (1.286,-0.511) (1.321,-0.516) (1.356,-0.521) (1.390,-0.525) (1.425,-0.529) (1.460,-0.533) (1.495,-0.537) (1.530,-0.540) (1.565,-0.543) (1.599,-0.546) (1.634,-0.549) (1.669,-0.551) (1.704,-0.554) (1.739,-0.556) (1.774,-0.558) (1.809,-0.560) (1.844,-0.562) (1.879,-0.564) (1.914,-0.565) (1.949,-0.567) (1.984,-0.569) (2.019,-0.570) (2.054,-0.571) (2.089,-0.572) (2.124,-0.574) (2.159,-0.575) (2.194,-0.576) (2.229,-0.577) (2.264,-0.578) (2.299,-0.579) (2.334,-0.580) (2.369,-0.580) (2.404,-0.581) (2.439,-0.582) (2.474,-0.583) (2.510,-0.583) (2.545,-0.584) (2.580,-0.584) (2.615,-0.585) (2.650,-0.586) (2.685,-0.586) (2.720,-0.587) (2.755,-0.587) (2.790,-0.587) (2.825,-0.588) (2.860,-0.588) (2.895,-0.589) (2.930,-0.589) (2.965,-0.589) (3.000,-0.590)};
\addplot[lineaD, flechaen=0.141, flechaen=0.859] coordinates {(-3.000,-0.787) (-2.965,-0.786) (-2.930,-0.786) (-2.895,-0.786) (-2.859,-0.785) (-2.824,-0.785) (-2.789,-0.784) (-2.754,-0.783) (-2.719,-0.783) (-2.684,-0.782) (-2.649,-0.782) (-2.613,-0.781) (-2.578,-0.780) (-2.543,-0.780) (-2.508,-0.779) (-2.473,-0.778) (-2.438,-0.777) (-2.402,-0.776) (-2.367,-0.775) (-2.332,-0.774) (-2.297,-0.773) (-2.262,-0.772) (-2.227,-0.771) (-2.192,-0.770) (-2.157,-0.769) (-2.121,-0.767) (-2.086,-0.766) (-2.051,-0.764) (-2.016,-0.763) (-1.981,-0.761) (-1.946,-0.759) (-1.911,-0.758) (-1.876,-0.756) (-1.841,-0.754) (-1.805,-0.751) (-1.770,-0.749) (-1.735,-0.747) (-1.700,-0.744) (-1.665,-0.741) (-1.630,-0.738) (-1.595,-0.735) (-1.560,-0.732) (-1.525,-0.728) (-1.490,-0.724) (-1.455,-0.720) (-1.420,-0.716) (-1.386,-0.711) (-1.351,-0.706) (-1.316,-0.701) (-1.281,-0.696) (-1.247,-0.690) (-1.212,-0.683) (-1.178,-0.676) (-1.143,-0.669) (-1.109,-0.661) (-1.075,-0.653) (-1.041,-0.643) (-1.007,-0.634) (-0.974,-0.623) (-0.940,-0.612) (-0.907,-0.600) (-0.874,-0.587) (-0.842,-0.574) (-0.809,-0.566) (-0.773,-0.566) (-0.738,-0.566) (-0.703,-0.566) (-0.668,-0.566) (-0.633,-0.566) (-0.598,-0.566) (-0.562,-0.566) (-0.527,-0.566) (-0.492,-0.566) (-0.457,-0.566) (-0.422,-0.566) (-0.387,-0.566) (-0.352,-0.566) (-0.316,-0.566) (-0.281,-0.566) (-0.246,-0.566) (-0.211,-0.566) (-0.176,-0.566) (-0.141,-0.566) (-0.105,-0.566) (-0.070,-0.566) (-0.035,-0.566) (0.000,-0.566) (0.035,-0.566) (0.070,-0.566) (0.105,-0.566) (0.141,-0.566) (0.176,-0.566) (0.211,-0.566) (0.246,-0.566) (0.281,-0.566) (0.316,-0.566) (0.352,-0.566) (0.387,-0.566) (0.422,-0.566) (0.457,-0.566) (0.492,-0.566) (0.527,-0.566) (0.562,-0.566) (0.598,-0.566) (0.633,-0.566) (0.668,-0.566) (0.703,-0.566) (0.738,-0.566) (0.773,-0.566) (0.809,-0.566) (0.842,-0.574) (0.874,-0.587) (0.907,-0.600) (0.940,-0.612) (0.974,-0.623) (1.007,-0.634) (1.041,-0.643) (1.075,-0.653) (1.109,-0.661) (1.143,-0.669) (1.178,-0.676) (1.212,-0.683) (1.247,-0.690) (1.281,-0.696) (1.316,-0.701) (1.351,-0.706) (1.386,-0.711) (1.420,-0.716) (1.455,-0.720) (1.490,-0.724) (1.525,-0.728) (1.560,-0.732) (1.595,-0.735) (1.630,-0.738) (1.665,-0.741) (1.700,-0.744) (1.735,-0.747) (1.770,-0.749) (1.805,-0.751) (1.841,-0.754) (1.876,-0.756) (1.911,-0.758) (1.946,-0.759) (1.981,-0.761) (2.016,-0.763) (2.051,-0.764) (2.086,-0.766) (2.121,-0.767) (2.157,-0.769) (2.192,-0.770) (2.227,-0.771) (2.262,-0.772) (2.297,-0.773) (2.332,-0.774) (2.367,-0.775) (2.402,-0.776) (2.438,-0.777) (2.473,-0.778) (2.508,-0.779) (2.543,-0.780) (2.578,-0.780) (2.613,-0.781) (2.649,-0.782) (2.684,-0.782) (2.719,-0.783) (2.754,-0.783) (2.789,-0.784) (2.824,-0.785) (2.859,-0.785) (2.895,-0.786) (2.930,-0.786) (2.965,-0.786) (3.000,-0.787)};
\addplot[lineaD, flechaen=0.140, flechaen=0.860] coordinates {(-3.000,-0.984) (-2.965,-0.984) (-2.930,-0.983) (-2.895,-0.983) (-2.860,-0.982) (-2.825,-0.982) (-2.789,-0.981) (-2.754,-0.981) (-2.719,-0.980) (-2.684,-0.979) (-2.649,-0.979) (-2.614,-0.978) (-2.579,-0.977) (-2.544,-0.976) (-2.509,-0.975) (-2.474,-0.974) (-2.439,-0.973) (-2.404,-0.972) (-2.368,-0.971) (-2.333,-0.970) (-2.298,-0.969) (-2.263,-0.968) (-2.228,-0.967) (-2.193,-0.966) (-2.158,-0.964) (-2.123,-0.963) (-2.088,-0.961) (-2.053,-0.960) (-2.018,-0.958) (-1.983,-0.956) (-1.948,-0.954) (-1.913,-0.952) (-1.878,-0.950) (-1.843,-0.948) (-1.808,-0.946) (-1.773,-0.943) (-1.738,-0.941) (-1.703,-0.938) (-1.668,-0.935) (-1.633,-0.932) (-1.598,-0.929) (-1.563,-0.926) (-1.528,-0.922) (-1.493,-0.918) (-1.458,-0.914) (-1.423,-0.910) (-1.389,-0.906) (-1.354,-0.901) (-1.319,-0.896) (-1.284,-0.890) (-1.250,-0.884) (-1.215,-0.878) (-1.181,-0.872) (-1.146,-0.865) (-1.112,-0.857) (-1.078,-0.849) (-1.044,-0.841) (-1.010,-0.831) (-0.976,-0.822) (-0.943,-0.811) (-0.910,-0.800) (-0.877,-0.788) (-0.844,-0.775) (-0.812,-0.761) (-0.780,-0.746) (-0.749,-0.731) (-0.718,-0.713) (-0.684,-0.707) (-0.649,-0.707) (-0.614,-0.707) (-0.579,-0.707) (-0.544,-0.707) (-0.509,-0.707) (-0.474,-0.707) (-0.439,-0.707) (-0.404,-0.707) (-0.368,-0.707) (-0.333,-0.707) (-0.298,-0.707) (-0.263,-0.707) (-0.228,-0.707) (-0.193,-0.707) (-0.158,-0.707) (-0.123,-0.707) (-0.088,-0.707) (-0.053,-0.707) (-0.018,-0.707) (0.018,-0.707) (0.053,-0.707) (0.088,-0.707) (0.123,-0.707) (0.158,-0.707) (0.193,-0.707) (0.228,-0.707) (0.263,-0.707) (0.298,-0.707) (0.333,-0.707) (0.368,-0.707) (0.404,-0.707) (0.439,-0.707) (0.474,-0.707) (0.509,-0.707) (0.544,-0.707) (0.579,-0.707) (0.614,-0.707) (0.649,-0.707) (0.684,-0.707) (0.718,-0.713) (0.749,-0.731) (0.780,-0.746) (0.812,-0.761) (0.844,-0.775) (0.877,-0.788) (0.910,-0.800) (0.943,-0.811) (0.976,-0.822) (1.010,-0.831) (1.044,-0.841) (1.078,-0.849) (1.112,-0.857) (1.146,-0.865) (1.181,-0.872) (1.215,-0.878) (1.250,-0.884) (1.284,-0.890) (1.319,-0.896) (1.354,-0.901) (1.389,-0.906) (1.423,-0.910) (1.458,-0.914) (1.493,-0.918) (1.528,-0.922) (1.563,-0.926) (1.598,-0.929) (1.633,-0.932) (1.668,-0.935) (1.703,-0.938) (1.738,-0.941) (1.773,-0.943) (1.808,-0.946) (1.843,-0.948) (1.878,-0.950) (1.913,-0.952) (1.948,-0.954) (1.983,-0.956) (2.018,-0.958) (2.053,-0.960) (2.088,-0.961) (2.123,-0.963) (2.158,-0.964) (2.193,-0.966) (2.228,-0.967) (2.263,-0.968) (2.298,-0.969) (2.333,-0.970) (2.368,-0.971) (2.404,-0.972) (2.439,-0.973) (2.474,-0.974) (2.509,-0.975) (2.544,-0.976) (2.579,-0.977) (2.614,-0.978) (2.649,-0.979) (2.684,-0.979) (2.719,-0.980) (2.754,-0.981) (2.789,-0.981) (2.825,-0.982) (2.860,-0.982) (2.895,-0.983) (2.930,-0.983) (2.965,-0.984) (3.000,-0.984)};
\addplot[lineaD, flechaen=0.139, flechaen=0.861] coordinates {(-3.000,-1.182) (-2.965,-1.182) (-2.930,-1.181) (-2.895,-1.181) (-2.860,-1.180) (-2.825,-1.180) (-2.790,-1.179) (-2.755,-1.178) (-2.720,-1.178) (-2.685,-1.177) (-2.650,-1.176) (-2.615,-1.175) (-2.580,-1.174) (-2.545,-1.174) (-2.510,-1.173) (-2.475,-1.172) (-2.440,-1.171) (-2.404,-1.170) (-2.369,-1.169) (-2.334,-1.168) (-2.299,-1.166) (-2.264,-1.165) (-2.229,-1.164) (-2.194,-1.163) (-2.159,-1.161) (-2.124,-1.160) (-2.089,-1.158) (-2.054,-1.157) (-2.019,-1.155) (-1.984,-1.153) (-1.949,-1.151) (-1.914,-1.149) (-1.879,-1.147) (-1.844,-1.145) (-1.809,-1.143) (-1.774,-1.141) (-1.740,-1.138) (-1.705,-1.136) (-1.670,-1.133) (-1.635,-1.130) (-1.600,-1.127) (-1.565,-1.124) (-1.530,-1.121) (-1.495,-1.117) (-1.460,-1.113) (-1.426,-1.110) (-1.391,-1.105) (-1.356,-1.101) (-1.321,-1.097) (-1.287,-1.092) (-1.252,-1.087) (-1.217,-1.081) (-1.183,-1.076) (-1.148,-1.070) (-1.114,-1.063) (-1.079,-1.056) (-1.045,-1.049) (-1.011,-1.042) (-0.977,-1.034) (-0.943,-1.025) (-0.909,-1.016) (-0.875,-1.006) (-0.842,-0.996) (-0.808,-0.985) (-0.775,-0.974) (-0.742,-0.961) (-0.710,-0.948) (-0.678,-0.934) (-0.646,-0.919) (-0.615,-0.903) (-0.585,-0.885) (-0.555,-0.867) (-0.526,-0.849) (-0.491,-0.849) (-0.456,-0.849) (-0.420,-0.849) (-0.385,-0.849) (-0.350,-0.849) (-0.315,-0.849) (-0.280,-0.849) (-0.245,-0.849) (-0.210,-0.849) (-0.175,-0.849) (-0.140,-0.849) (-0.105,-0.849) (-0.070,-0.849) (-0.035,-0.849) (0.000,-0.849) (0.035,-0.849) (0.070,-0.849) (0.105,-0.849) (0.140,-0.849) (0.175,-0.849) (0.210,-0.849) (0.245,-0.849) (0.280,-0.849) (0.315,-0.849) (0.350,-0.849) (0.385,-0.849) (0.420,-0.849) (0.456,-0.849) (0.491,-0.849) (0.526,-0.849) (0.555,-0.867) (0.585,-0.885) (0.615,-0.903) (0.646,-0.919) (0.678,-0.934) (0.710,-0.948) (0.742,-0.961) (0.775,-0.974) (0.808,-0.985) (0.842,-0.996) (0.875,-1.006) (0.909,-1.016) (0.943,-1.025) (0.977,-1.034) (1.011,-1.042) (1.045,-1.049) (1.079,-1.056) (1.114,-1.063) (1.148,-1.070) (1.183,-1.076) (1.217,-1.081) (1.252,-1.087) (1.287,-1.092) (1.321,-1.097) (1.356,-1.101) (1.391,-1.105) (1.426,-1.110) (1.460,-1.113) (1.495,-1.117) (1.530,-1.121) (1.565,-1.124) (1.600,-1.127) (1.635,-1.130) (1.670,-1.133) (1.705,-1.136) (1.740,-1.138) (1.774,-1.141) (1.809,-1.143) (1.844,-1.145) (1.879,-1.147) (1.914,-1.149) (1.949,-1.151) (1.984,-1.153) (2.019,-1.155) (2.054,-1.157) (2.089,-1.158) (2.124,-1.160) (2.159,-1.161) (2.194,-1.163) (2.229,-1.164) (2.264,-1.165) (2.299,-1.166) (2.334,-1.168) (2.369,-1.169) (2.404,-1.170) (2.440,-1.171) (2.475,-1.172) (2.510,-1.173) (2.545,-1.174) (2.580,-1.174) (2.615,-1.175) (2.650,-1.176) (2.685,-1.177) (2.720,-1.178) (2.755,-1.178) (2.790,-1.179) (2.825,-1.180) (2.860,-1.180) (2.895,-1.181) (2.930,-1.181) (2.965,-1.182) (3.000,-1.182)};
\addplot[lineaD, flechaen=0.139, flechaen=0.861] coordinates {(-3.000,-1.381) (-2.965,-1.380) (-2.930,-1.380) (-2.895,-1.379) (-2.860,-1.379) (-2.825,-1.378) (-2.790,-1.377) (-2.755,-1.377) (-2.720,-1.376) (-2.685,-1.375) (-2.650,-1.374) (-2.615,-1.374) (-2.580,-1.373) (-2.545,-1.372) (-2.510,-1.371) (-2.475,-1.370) (-2.439,-1.369) (-2.404,-1.368) (-2.369,-1.367) (-2.334,-1.366) (-2.299,-1.365) (-2.264,-1.364) (-2.229,-1.362) (-2.194,-1.361) (-2.159,-1.360) (-2.124,-1.358) (-2.089,-1.357) (-2.054,-1.355) (-2.019,-1.354) (-1.984,-1.352) (-1.949,-1.350) (-1.914,-1.348) (-1.879,-1.347) (-1.844,-1.345) (-1.809,-1.343) (-1.774,-1.340) (-1.739,-1.338) (-1.704,-1.336) (-1.669,-1.333) (-1.635,-1.331) (-1.600,-1.328) (-1.565,-1.325) (-1.530,-1.322) (-1.495,-1.319) (-1.460,-1.316) (-1.425,-1.313) (-1.390,-1.309) (-1.355,-1.305) (-1.321,-1.302) (-1.286,-1.298) (-1.251,-1.293) (-1.216,-1.289) (-1.181,-1.284) (-1.147,-1.280) (-1.112,-1.274) (-1.077,-1.269) (-1.043,-1.264) (-1.008,-1.258) (-0.974,-1.252) (-0.939,-1.246) (-0.905,-1.239) (-0.871,-1.232) (-0.836,-1.225) (-0.802,-1.217) (-0.768,-1.209) (-0.734,-1.201) (-0.700,-1.193) (-0.666,-1.184) (-0.632,-1.174) (-0.599,-1.164) (-0.565,-1.154) (-0.532,-1.144) (-0.498,-1.133) (-0.465,-1.121) (-0.432,-1.109) (-0.399,-1.097) (-0.367,-1.085) (-0.334,-1.071) (-0.302,-1.058) (-0.270,-1.044) (-0.237,-1.031) (-0.205,-1.017) (-0.173,-1.003) (-0.140,-0.991) (-0.105,-0.990) (-0.070,-0.990) (-0.035,-0.990) (-0.000,-0.990) (0.035,-0.990) (0.070,-0.990) (0.105,-0.990) (0.140,-0.991) (0.173,-1.003) (0.205,-1.017) (0.237,-1.031) (0.270,-1.044) (0.302,-1.058) (0.334,-1.071) (0.367,-1.085) (0.399,-1.097) (0.432,-1.109) (0.465,-1.121) (0.498,-1.133) (0.532,-1.144) (0.565,-1.154) (0.599,-1.164) (0.632,-1.174) (0.666,-1.184) (0.700,-1.193) (0.734,-1.201) (0.768,-1.209) (0.802,-1.217) (0.836,-1.225) (0.871,-1.232) (0.905,-1.239) (0.939,-1.246) (0.974,-1.252) (1.008,-1.258) (1.043,-1.264) (1.077,-1.269) (1.112,-1.274) (1.147,-1.280) (1.181,-1.284) (1.216,-1.289) (1.251,-1.293) (1.286,-1.298) (1.321,-1.302) (1.355,-1.305) (1.390,-1.309) (1.425,-1.313) (1.460,-1.316) (1.495,-1.319) (1.530,-1.322) (1.565,-1.325) (1.600,-1.328) (1.635,-1.331) (1.669,-1.333) (1.704,-1.336) (1.739,-1.338) (1.774,-1.340) (1.809,-1.343) (1.844,-1.345) (1.879,-1.347) (1.914,-1.348) (1.949,-1.350) (1.984,-1.352) (2.019,-1.354) (2.054,-1.355) (2.089,-1.357) (2.124,-1.358) (2.159,-1.360) (2.194,-1.361) (2.229,-1.362) (2.264,-1.364) (2.299,-1.365) (2.334,-1.366) (2.369,-1.367) (2.404,-1.368) (2.439,-1.369) (2.475,-1.370) (2.510,-1.371) (2.545,-1.372) (2.580,-1.373) (2.615,-1.374) (2.650,-1.374) (2.685,-1.375) (2.720,-1.376) (2.755,-1.377) (2.790,-1.377) (2.825,-1.378) (2.860,-1.379) (2.895,-1.379) (2.930,-1.380) (2.965,-1.380) (3.000,-1.381)};
\addplot[lineaD, flechaen=0.141, flechaen=0.859] coordinates {(-3.000,-1.580) (-2.965,-1.579) (-2.930,-1.579) (-2.895,-1.578) (-2.860,-1.578) (-2.825,-1.577) (-2.790,-1.576) (-2.755,-1.576) (-2.720,-1.575) (-2.685,-1.574) (-2.650,-1.573) (-2.615,-1.573) (-2.580,-1.572) (-2.545,-1.571) (-2.510,-1.570) (-2.475,-1.569) (-2.440,-1.568) (-2.405,-1.567) (-2.370,-1.566) (-2.335,-1.565) (-2.300,-1.564) (-2.265,-1.563) (-2.230,-1.562) (-2.195,-1.561) (-2.160,-1.559) (-2.125,-1.558) (-2.090,-1.557) (-2.055,-1.555) (-2.020,-1.554) (-1.985,-1.552) (-1.950,-1.551) (-1.915,-1.549) (-1.880,-1.547) (-1.845,-1.546) (-1.810,-1.544) (-1.775,-1.542) (-1.740,-1.540) (-1.705,-1.538) (-1.670,-1.536) (-1.635,-1.533) (-1.600,-1.531) (-1.565,-1.529) (-1.530,-1.526) (-1.495,-1.524) (-1.460,-1.521) (-1.425,-1.518) (-1.390,-1.515) (-1.355,-1.512) (-1.320,-1.509) (-1.286,-1.506) (-1.251,-1.503) (-1.216,-1.499) (-1.181,-1.496) (-1.146,-1.492) (-1.111,-1.488) (-1.077,-1.484) (-1.042,-1.480) (-1.007,-1.476) (-0.972,-1.471) (-0.937,-1.467) (-0.903,-1.462) (-0.868,-1.458) (-0.833,-1.453) (-0.799,-1.448) (-0.764,-1.443) (-0.729,-1.438) (-0.695,-1.432) (-0.660,-1.427) (-0.626,-1.421) (-0.591,-1.416) (-0.556,-1.410) (-0.522,-1.405) (-0.487,-1.399) (-0.453,-1.394) (-0.418,-1.388) (-0.383,-1.383) (-0.349,-1.378) (-0.314,-1.373) (-0.279,-1.368) (-0.245,-1.364) (-0.210,-1.360) (-0.175,-1.357) (-0.140,-1.354) (-0.105,-1.352) (-0.070,-1.350) (-0.035,-1.349) (0.000,-1.348) (0.035,-1.349) (0.070,-1.350) (0.105,-1.352) (0.140,-1.354) (0.175,-1.357) (0.210,-1.360) (0.245,-1.364) (0.279,-1.368) (0.314,-1.373) (0.349,-1.378) (0.383,-1.383) (0.418,-1.388) (0.453,-1.394) (0.487,-1.399) (0.522,-1.405) (0.556,-1.410) (0.591,-1.416) (0.626,-1.421) (0.660,-1.427) (0.695,-1.432) (0.729,-1.438) (0.764,-1.443) (0.799,-1.448) (0.833,-1.453) (0.868,-1.458) (0.903,-1.462) (0.937,-1.467) (0.972,-1.471) (1.007,-1.476) (1.042,-1.480) (1.077,-1.484) (1.111,-1.488) (1.146,-1.492) (1.181,-1.496) (1.216,-1.499) (1.251,-1.503) (1.286,-1.506) (1.320,-1.509) (1.355,-1.512) (1.390,-1.515) (1.425,-1.518) (1.460,-1.521) (1.495,-1.524) (1.530,-1.526) (1.565,-1.529) (1.600,-1.531) (1.635,-1.533) (1.670,-1.536) (1.705,-1.538) (1.740,-1.540) (1.775,-1.542) (1.810,-1.544) (1.845,-1.546) (1.880,-1.547) (1.915,-1.549) (1.950,-1.551) (1.985,-1.552) (2.020,-1.554) (2.055,-1.555) (2.090,-1.557) (2.125,-1.558) (2.160,-1.559) (2.195,-1.561) (2.230,-1.562) (2.265,-1.563) (2.300,-1.564) (2.335,-1.565) (2.370,-1.566) (2.405,-1.567) (2.440,-1.568) (2.475,-1.569) (2.510,-1.570) (2.545,-1.571) (2.580,-1.572) (2.615,-1.573) (2.650,-1.573) (2.685,-1.574) (2.720,-1.575) (2.755,-1.576) (2.790,-1.576) (2.825,-1.577) (2.860,-1.578) (2.895,-1.578) (2.930,-1.579) (2.965,-1.579) (3.000,-1.580)};
\addplot[lineaD, flechaen=0.141, flechaen=0.859] coordinates {(-3.000,-1.779) (-2.965,-1.779) (-2.930,-1.778) (-2.895,-1.777) (-2.859,-1.777) (-2.824,-1.776) (-2.789,-1.776) (-2.754,-1.775) (-2.719,-1.774) (-2.684,-1.774) (-2.649,-1.773) (-2.613,-1.772) (-2.578,-1.771) (-2.543,-1.770) (-2.508,-1.770) (-2.473,-1.769) (-2.438,-1.768) (-2.403,-1.767) (-2.367,-1.766) (-2.332,-1.765) (-2.297,-1.764) (-2.262,-1.763) (-2.227,-1.762) (-2.192,-1.761) (-2.157,-1.760) (-2.121,-1.758) (-2.086,-1.757) (-2.051,-1.756) (-2.016,-1.755) (-1.981,-1.753) (-1.946,-1.752) (-1.911,-1.750) (-1.876,-1.749) (-1.841,-1.747) (-1.805,-1.746) (-1.770,-1.744) (-1.735,-1.742) (-1.700,-1.741) (-1.665,-1.739) (-1.630,-1.737) (-1.595,-1.735) (-1.560,-1.733) (-1.525,-1.731) (-1.489,-1.729) (-1.454,-1.727) (-1.419,-1.724) (-1.384,-1.722) (-1.349,-1.720) (-1.314,-1.717) (-1.279,-1.715) (-1.244,-1.712) (-1.209,-1.709) (-1.174,-1.707) (-1.139,-1.704) (-1.104,-1.701) (-1.069,-1.698) (-1.034,-1.695) (-0.999,-1.692) (-0.964,-1.689) (-0.929,-1.686) (-0.894,-1.682) (-0.859,-1.679) (-0.824,-1.676) (-0.789,-1.673) (-0.754,-1.669) (-0.719,-1.666) (-0.684,-1.663) (-0.649,-1.659) (-0.614,-1.656) (-0.579,-1.653) (-0.544,-1.649) (-0.509,-1.646) (-0.474,-1.643) (-0.439,-1.640) (-0.404,-1.637) (-0.369,-1.635) (-0.334,-1.632) (-0.299,-1.630) (-0.263,-1.628) (-0.228,-1.626) (-0.193,-1.624) (-0.158,-1.622) (-0.123,-1.621) (-0.088,-1.620) (-0.053,-1.620) (-0.018,-1.619) (0.018,-1.619) (0.053,-1.620) (0.088,-1.620) (0.123,-1.621) (0.158,-1.622) (0.193,-1.624) (0.228,-1.626) (0.263,-1.628) (0.299,-1.630) (0.334,-1.632) (0.369,-1.635) (0.404,-1.637) (0.439,-1.640) (0.474,-1.643) (0.509,-1.646) (0.544,-1.649) (0.579,-1.653) (0.614,-1.656) (0.649,-1.659) (0.684,-1.663) (0.719,-1.666) (0.754,-1.669) (0.789,-1.673) (0.824,-1.676) (0.859,-1.679) (0.894,-1.682) (0.929,-1.686) (0.964,-1.689) (0.999,-1.692) (1.034,-1.695) (1.069,-1.698) (1.104,-1.701) (1.139,-1.704) (1.174,-1.707) (1.209,-1.709) (1.244,-1.712) (1.279,-1.715) (1.314,-1.717) (1.349,-1.720) (1.384,-1.722) (1.419,-1.724) (1.454,-1.727) (1.489,-1.729) (1.525,-1.731) (1.560,-1.733) (1.595,-1.735) (1.630,-1.737) (1.665,-1.739) (1.700,-1.741) (1.735,-1.742) (1.770,-1.744) (1.805,-1.746) (1.841,-1.747) (1.876,-1.749) (1.911,-1.750) (1.946,-1.752) (1.981,-1.753) (2.016,-1.755) (2.051,-1.756) (2.086,-1.757) (2.121,-1.758) (2.157,-1.760) (2.192,-1.761) (2.227,-1.762) (2.262,-1.763) (2.297,-1.764) (2.332,-1.765) (2.367,-1.766) (2.403,-1.767) (2.438,-1.768) (2.473,-1.769) (2.508,-1.770) (2.543,-1.770) (2.578,-1.771) (2.613,-1.772) (2.649,-1.773) (2.684,-1.774) (2.719,-1.774) (2.754,-1.775) (2.789,-1.776) (2.824,-1.776) (2.859,-1.777) (2.895,-1.777) (2.930,-1.778) (2.965,-1.779) (3.000,-1.779)};
\addplot[lineaD, flechaen=0.142, flechaen=0.858] coordinates {(-3.000,-1.979) (-2.965,-1.978) (-2.930,-1.978) (-2.895,-1.977) (-2.860,-1.977) (-2.824,-1.976) (-2.789,-1.975) (-2.754,-1.975) (-2.719,-1.974) (-2.684,-1.973) (-2.649,-1.973) (-2.614,-1.972) (-2.579,-1.971) (-2.544,-1.971) (-2.508,-1.970) (-2.473,-1.969) (-2.438,-1.968) (-2.403,-1.967) (-2.368,-1.967) (-2.333,-1.966) (-2.298,-1.965) (-2.263,-1.964) (-2.228,-1.963) (-2.192,-1.962) (-2.157,-1.961) (-2.122,-1.960) (-2.087,-1.959) (-2.052,-1.958) (-2.017,-1.956) (-1.982,-1.955) (-1.947,-1.954) (-1.912,-1.953) (-1.877,-1.951) (-1.841,-1.950) (-1.806,-1.949) (-1.771,-1.947) (-1.736,-1.946) (-1.701,-1.944) (-1.666,-1.943) (-1.631,-1.941) (-1.596,-1.940) (-1.561,-1.938) (-1.526,-1.936) (-1.491,-1.935) (-1.455,-1.933) (-1.420,-1.931) (-1.385,-1.929) (-1.350,-1.927) (-1.315,-1.926) (-1.280,-1.924) (-1.245,-1.922) (-1.210,-1.920) (-1.175,-1.917) (-1.140,-1.915) (-1.105,-1.913) (-1.070,-1.911) (-1.035,-1.909) (-1.000,-1.907) (-0.965,-1.905) (-0.930,-1.902) (-0.895,-1.900) (-0.859,-1.898) (-0.824,-1.896) (-0.789,-1.893) (-0.754,-1.891) (-0.719,-1.889) (-0.684,-1.887) (-0.649,-1.885) (-0.614,-1.882) (-0.579,-1.880) (-0.544,-1.878) (-0.509,-1.876) (-0.474,-1.875) (-0.439,-1.873) (-0.404,-1.871) (-0.369,-1.869) (-0.334,-1.868) (-0.298,-1.867) (-0.263,-1.865) (-0.228,-1.864) (-0.193,-1.863) (-0.158,-1.862) (-0.123,-1.862) (-0.088,-1.861) (-0.053,-1.861) (-0.018,-1.861) (0.018,-1.861) (0.053,-1.861) (0.088,-1.861) (0.123,-1.862) (0.158,-1.862) (0.193,-1.863) (0.228,-1.864) (0.263,-1.865) (0.298,-1.867) (0.334,-1.868) (0.369,-1.869) (0.404,-1.871) (0.439,-1.873) (0.474,-1.875) (0.509,-1.876) (0.544,-1.878) (0.579,-1.880) (0.614,-1.882) (0.649,-1.885) (0.684,-1.887) (0.719,-1.889) (0.754,-1.891) (0.789,-1.893) (0.824,-1.896) (0.859,-1.898) (0.895,-1.900) (0.930,-1.902) (0.965,-1.905) (1.000,-1.907) (1.035,-1.909) (1.070,-1.911) (1.105,-1.913) (1.140,-1.915) (1.175,-1.917) (1.210,-1.920) (1.245,-1.922) (1.280,-1.924) (1.315,-1.926) (1.350,-1.927) (1.385,-1.929) (1.420,-1.931) (1.455,-1.933) (1.491,-1.935) (1.526,-1.936) (1.561,-1.938) (1.596,-1.940) (1.631,-1.941) (1.666,-1.943) (1.701,-1.944) (1.736,-1.946) (1.771,-1.947) (1.806,-1.949) (1.841,-1.950) (1.877,-1.951) (1.912,-1.953) (1.947,-1.954) (1.982,-1.955) (2.017,-1.956) (2.052,-1.958) (2.087,-1.959) (2.122,-1.960) (2.157,-1.961) (2.192,-1.962) (2.228,-1.963) (2.263,-1.964) (2.298,-1.965) (2.333,-1.966) (2.368,-1.967) (2.403,-1.967) (2.438,-1.968) (2.473,-1.969) (2.508,-1.970) (2.544,-1.971) (2.579,-1.971) (2.614,-1.972) (2.649,-1.973) (2.684,-1.973) (2.719,-1.974) (2.754,-1.975) (2.789,-1.975) (2.824,-1.976) (2.860,-1.977) (2.895,-1.977) (2.930,-1.978) (2.965,-1.978) (3.000,-1.979)};
\addplot[lineaD, flechaen=0.142, flechaen=0.858] coordinates {(-3.000,-2.179) (-2.965,-2.178) (-2.930,-2.178) (-2.895,-2.177) (-2.860,-2.177) (-2.824,-2.176) (-2.789,-2.176) (-2.754,-2.175) (-2.719,-2.174) (-2.684,-2.174) (-2.649,-2.173) (-2.614,-2.173) (-2.579,-2.172) (-2.544,-2.171) (-2.509,-2.170) (-2.473,-2.170) (-2.438,-2.169) (-2.403,-2.168) (-2.368,-2.167) (-2.333,-2.167) (-2.298,-2.166) (-2.263,-2.165) (-2.228,-2.164) (-2.193,-2.163) (-2.158,-2.162) (-2.123,-2.161) (-2.087,-2.161) (-2.052,-2.160) (-2.017,-2.159) (-1.982,-2.158) (-1.947,-2.156) (-1.912,-2.155) (-1.877,-2.154) (-1.842,-2.153) (-1.807,-2.152) (-1.772,-2.151) (-1.737,-2.150) (-1.701,-2.148) (-1.666,-2.147) (-1.631,-2.146) (-1.596,-2.145) (-1.561,-2.143) (-1.526,-2.142) (-1.491,-2.141) (-1.456,-2.139) (-1.421,-2.138) (-1.386,-2.136) (-1.351,-2.135) (-1.316,-2.133) (-1.281,-2.132) (-1.245,-2.130) (-1.210,-2.129) (-1.175,-2.127) (-1.140,-2.126) (-1.105,-2.124) (-1.070,-2.122) (-1.035,-2.121) (-1.000,-2.119) (-0.965,-2.118) (-0.930,-2.116) (-0.895,-2.114) (-0.860,-2.113) (-0.825,-2.111) (-0.790,-2.110) (-0.754,-2.108) (-0.719,-2.107) (-0.684,-2.105) (-0.649,-2.104) (-0.614,-2.102) (-0.579,-2.101) (-0.544,-2.100) (-0.509,-2.098) (-0.474,-2.097) (-0.439,-2.096) (-0.404,-2.095) (-0.369,-2.094) (-0.333,-2.093) (-0.298,-2.092) (-0.263,-2.091) (-0.228,-2.090) (-0.193,-2.090) (-0.158,-2.089) (-0.123,-2.089) (-0.088,-2.089) (-0.053,-2.088) (-0.018,-2.088) (0.018,-2.088) (0.053,-2.088) (0.088,-2.089) (0.123,-2.089) (0.158,-2.089) (0.193,-2.090) (0.228,-2.090) (0.263,-2.091) (0.298,-2.092) (0.333,-2.093) (0.369,-2.094) (0.404,-2.095) (0.439,-2.096) (0.474,-2.097) (0.509,-2.098) (0.544,-2.100) (0.579,-2.101) (0.614,-2.102) (0.649,-2.104) (0.684,-2.105) (0.719,-2.107) (0.754,-2.108) (0.790,-2.110) (0.825,-2.111) (0.860,-2.113) (0.895,-2.114) (0.930,-2.116) (0.965,-2.118) (1.000,-2.119) (1.035,-2.121) (1.070,-2.122) (1.105,-2.124) (1.140,-2.126) (1.175,-2.127) (1.210,-2.129) (1.245,-2.130) (1.281,-2.132) (1.316,-2.133) (1.351,-2.135) (1.386,-2.136) (1.421,-2.138) (1.456,-2.139) (1.491,-2.141) (1.526,-2.142) (1.561,-2.143) (1.596,-2.145) (1.631,-2.146) (1.666,-2.147) (1.701,-2.148) (1.737,-2.150) (1.772,-2.151) (1.807,-2.152) (1.842,-2.153) (1.877,-2.154) (1.912,-2.155) (1.947,-2.156) (1.982,-2.158) (2.017,-2.159) (2.052,-2.160) (2.087,-2.161) (2.123,-2.161) (2.158,-2.162) (2.193,-2.163) (2.228,-2.164) (2.263,-2.165) (2.298,-2.166) (2.333,-2.167) (2.368,-2.167) (2.403,-2.168) (2.438,-2.169) (2.473,-2.170) (2.509,-2.170) (2.544,-2.171) (2.579,-2.172) (2.614,-2.173) (2.649,-2.173) (2.684,-2.174) (2.719,-2.174) (2.754,-2.175) (2.789,-2.176) (2.824,-2.176) (2.860,-2.177) (2.895,-2.177) (2.930,-2.178) (2.965,-2.178) (3.000,-2.179)};
\addplot[lineaD, flechaen=0.142, flechaen=0.858] coordinates {(-3.000,0.000) (3.000,0.000)};

  \draw[texto, line width=0.5pt] (axis cs:0,0) circle[radius=1];
\end{axis}
\end{tikzpicture}\\[-0.2em]
{\scriptsize\color{tenue}\itshape Líneas de $\vD$ en torno a una esfera dieléctrica\\ en un campo uniforme ($\er=4$).}
\end{minipage}

\vspace{1.4em}
\noindent{\small{\color{acentoA}\espaciado{Palabras clave.}}\enspace dieléctrico $\cdot$ polarizabilidad
$\cdot$ vector polarización $\cdot$ carga ligada $\cdot$ vector desplazamiento $\cdot$ electrete $\cdot$
susceptibilidad $\cdot$ permitividad $\cdot$ condiciones de frontera $\cdot$ refracción de las líneas de
campo $\cdot$ cavidades}

\vspace{0.6em}
\noindent{\small{\color{acentoA}\espaciado{Clasificación PACS.}}\enspace 41.20.Cv (electrostática;
ecuaciones de Poisson y Laplace, problemas de contorno) $\cdot$ 77.22.$-$d (propiedades dieléctricas de
sólidos y líquidos) $\cdot$ 77.22.Ej (polarización y despolarización)}

\vfill
\begin{center}
\small
$\displaystyle \phi(\vr)=\frac{1}{4\pi\eo}\!\int_{V'}\!\frac{\rhoP}{R}\dd V'+\frac{1}{4\pi\eo}\oiint_{S'}\!\frac{\sigmaP}{R}\dd S'
\qquad \oiint_S\vD\cdot\dS=\qfenc
\qquad \frac{\tan\theta_2}{\varepsilon_2}=\frac{\tan\theta_1}{\varepsilon_1}$
\end{center}
\vspace{1em}
\newpage

% ------------------------------------------------------------- Índice
\thispagestyle{empty}
\enlargethispage{3\baselineskip}
\tableofcontents
\thispagestyle{empty}
\newpage
\setcounter{page}{1}

% =============================================================================
\section{Introducción: del conductor al dieléctrico}\label{sec:introduccion}
% =============================================================================

En el tema anterior dimos el primer paso hacia la materia con el conductor perfecto, un cuerpo con una
reserva ilimitada de cargas que se mueven sin oposición~\cite{T3}. En el otro extremo está el
\emph{aislante perfecto}, o \emph{dieléctrico}: un material que no admite cargas libres. Sus cargas
---núcleos y electrones--- están ligadas a los átomos o moléculas y no pueden abandonar su posición,
aunque sí desplazarse ligeramente respecto de ella. Ambas idealizaciones funcionan muy bien en
electrostática, y entre las dos cubren la mayor parte de los materiales de interés.

La pregunta de este tema es qué hace un campo eléctrico con un dieléctrico. La respuesta es que lo
\emph{polariza}: cada átomo o molécula se convierte en un pequeño dipolo, o alinea parcialmente el que
ya tenía. Un dieléctrico polarizado será, a escala macroscópica, una distribución continua de dipolos,
y para describirlo bastarán la teoría del dipolo y el desarrollo multipolar de los temas
anteriores~\cite{T1,T2}. La novedad es que, como en el conductor, el campo y la polarización se
determinan mutuamente: el campo polariza el medio, y los dipolos crean a su vez un campo que se suma al
aplicado.

Seguimos el orden de los apuntes manuscritos de la asignatura y nos apoyamos en los apuntes de teoría y
en sus figuras~\cite{DFA-teoria,DFA-figuras}, que hemos recalculado y redibujado; para ampliar pueden
consultarse los textos clásicos~\cite{Griffiths,Wangsness,Purcell,Feynman,Jackson}.

\subsection{El dieléctrico ideal}\label{sec:dielectrico-ideal}

\begin{definicion}[Dieléctrico]\label{def:dielectrico}
Un \emph{dieléctrico} (o aislante perfecto) es un cuerpo material, que ocupa una región $V\subset\mathbb{R}^3$
de frontera $S=\partial V$, tal que:
\begin{enumerate}
  \item no contiene cargas libres: todas sus cargas están ligadas a átomos o moléculas, alrededor de cuya
        posición de equilibrio solo pueden desplazarse distancias del orden atómico o menores;
  \item es globalmente neutro en ausencia de cargas externas depositadas en él.
\end{enumerate}
Llamaremos \emph{cargas libres} a todas las demás: las de los conductores, los iones o electrones que
podamos depositar en el material, y en general cualquier carga cuyo movimiento no esté restringido a la
escala atómica.
\end{definicion}

\begin{observacion}[Libre y ligada no es una distinción microscópica]\label{obs:libre-ligada}
Un electrón es un electrón, y crea el mismo campo esté ligado o no. La distinción es de
\emph{control}: la carga libre es la que fijamos nosotros (depositándola, conectando una batería), mientras
que la ligada es la que el material reorganiza por sí solo en respuesta al campo. Esta asimetría es la que
hará útil el vector desplazamiento de la \cref{sec:desplazamiento}.
\end{observacion}

\subsection{Notación y convenios}\label{sec:notacion}

\begin{notacion}\label{not:tema4}
Usamos la notación de los temas anteriores~\cite{T0,T1,T2,T3} y, además:
\begin{itemize}
  \item $\vr$ es el punto de campo y $\vr\,'$ el punto fuente; $\vR=\vr-\vr\,'$, $R=\abs{\vR}$ y
        $\uv R=\vR/R$. El operador $\grad'$ actúa sobre las coordenadas primadas.
  \item $\rhof$ y $\sigmaf$ son las densidades de carga libre, volúmica y superficial; $\rhoP$ y $\sigmaP$, las de
        carga ligada o \emph{de polarización}. $\qfenc$ y $\qPenc$ son las cargas libre y ligada
        encerradas por una superficie.
  \item $\vp$ es el momento dipolar de un átomo o molécula; $\vP$, el vector polarización; $\vD$, el vector
        desplazamiento; $\chi$, $\varepsilon$ y $\er$, la susceptibilidad, la permitividad y la permitividad relativa.
  \item $\nn$ es la normal unitaria exterior al dieléctrico. En una superficie que separa dos medios, $\nn$ apunta
        del medio~1 al medio~2, y los subíndices $n$ y $t$ indican las componentes normal y tangencial.
  \item \LHI{} abrevia «lineal, homogéneo e isótropo». Mientras no se diga lo contrario,
        $\phi(\infty)=0$ y $\eo\simeq\SI{8.854e-12}{F/m}$~\cite{CODATA}.
\end{itemize}
\end{notacion}

% =============================================================================
\section{Modelo microscópico de la polarización}\label{sec:microscopico}
% =============================================================================

Vamos a describir los dieléctricos mediante la \emph{aproximación dipolar}: un dieléctrico polarizado es,
a escala macroscópica, una distribución de dipolos. Antes de adoptarla conviene justificarla con un
modelo microscópico sencillo, que además nos dará los órdenes de magnitud. Hay que distinguir dos clases
de sustancias según que sus moléculas tengan o no momento dipolar en ausencia de campo.

\subsection{Sustancias no polares: el átomo de nube uniforme}\label{sec:no-polares}

En una sustancia no polar el momento dipolar de cada átomo o molécula es nulo, $\vp=\vcero$, a escala
atómica. El modelo más simple de átomo es un núcleo puntual de carga $+q$ rodeado de una nube electrónica
esférica de carga $-q$ uniformemente repartida en un radio $a$ (\cref{fig:atomo}). Sin campo aplicado,
el centro de la nube coincide con el núcleo y el momento dipolar es nulo.

% Figura: modelo de nube electrónica uniforme (átomo sin polarizar / polarizado)
% y campo de la nube en función de la distancia a su centro.
\begin{figure}[t]
\centering
\begin{subfigure}[b]{0.3\textwidth}\centering
\begin{tikzpicture}[scale=1.0]
  \shade[inner color=acentoA!38!fondo, outer color=acentoA!14!fondo] (0,0) circle (1.15);
  \draw[acentoA!80!fondo, line width=0.6pt] (0,0) circle (1.15);
  \draw[tenue, dashed, line width=0.4pt] (0,0) -- (-40:1.15) node[midway, below left=-1pt, etiq] {$a$};
  \fill[cargapos] (0,0) circle (0.1);
  \node[etiq, above=2pt] at (0,0.05) {$+q$};
  \node[etiq] at (0.72,0.72) {$-q$};
  \node[etiq] at (0,1.55) {$\vE=\vcero$};
\end{tikzpicture}
\caption{Átomo sin polarizar}
\end{subfigure}\hfill
\begin{subfigure}[b]{0.3\textwidth}\centering
\begin{tikzpicture}[scale=1.0]
  \shade[inner color=acentoA!38!fondo, outer color=acentoA!14!fondo] (0,0) circle (1.15);
  \draw[acentoA!80!fondo, line width=0.6pt] (0,0) circle (1.15);
  \fill[texto] (0,0) circle (0.035);
  \draw[campoP, line width=0.9pt] (0,0) -- (0.42,0);
  \node[etiq, below=1pt] at (0.21,0) {$d$};
  \fill[cargapos] (0.42,0) circle (0.1);
  \node[etiq, above=2pt] at (0.42,0.05) {$+q$};
  \node[etiq] at (-0.62,0.72) {$-q$};
  \draw[vecE] (-0.7,1.55) -- (0.7,1.55) node[right, etiq] {$\vE$};
\end{tikzpicture}
\caption{Átomo polarizado}
\end{subfigure}\hfill
\begin{subfigure}[b]{0.36\textwidth}\centering
\begin{tikzpicture}
\begin{axis}[width=5.6cm, height=4.1cm, axis lines=left, xmin=0, xmax=3.2, ymin=0, ymax=1.18,
  xlabel={$r/a$}, ylabel={$E_{\text{nube}}\big/\tfrac{q}{4\pi\eo a^{2}}$},
  xtick={0,1,2,3}, ytick={0,0.5,1}, ylabel style={yshift=-3pt}]
  \addplot[curvaE, domain=0:1, samples=2] {x};
  \addplot[curvaE, domain=1:3.2, samples=80] {1/x^2};
  \draw[equipot] (axis cs:1,0) -- (axis cs:1,1.18);
  \node[etiqt, anchor=west] at (axis cs:0.08,0.86) {$\propto r$};
  \node[etiqt, anchor=west] at (axis cs:1.75,0.52) {$\propto 1/r^{2}$};
\end{axis}
\end{tikzpicture}
\caption{Campo de la nube}
\end{subfigure}
\caption{Modelo atómico de nube uniforme. (a) Sin campo aplicado, el centro de la nube electrónica
(carga $-q$, radio $a$) coincide con el núcleo y el momento dipolar es nulo. (b) Un campo externo
desplaza el núcleo una distancia $d\ll a$ respecto al centro de la nube, y aparece $\vp=q\vec d$.
(c) Campo de la propia nube: dentro crece linealmente con $r$ (ley de Hooke electrostática),
fuera decae como el de una carga puntual; el equilibrio se alcanza en la rama lineal.}
\label{fig:atomo}
\end{figure}


Si aplicamos un campo externo $\vE$, el núcleo es empujado en el sentido del campo y la nube en el
contrario. Supondremos que la nube se desplaza rígidamente, sin deformarse, y que el desplazamiento
relativo $d$ es pequeño frente a $a$.

\begin{proposicion}[Polarizabilidad del átomo de nube uniforme]\label{prop:polarizabilidad}
En el modelo de nube uniforme, un campo externo $\vE$ induce el momento dipolar
\begin{equation}\label{eq:p-atomo}
  \vp=q\,\vec d=4\pi\eo a^3\,\vE ,
\end{equation}
donde $\vec d$ es el desplazamiento del núcleo respecto al centro de la nube.
\end{proposicion}

\begin{proof}
El campo de la nube en un punto interior a distancia $r$ de su centro se obtiene con la ley de Gauss: una
esfera concéntrica de radio $r<a$ encierra la fracción $(r/a)^3$ de la carga $-q$, luego
\[
  E_{\text{nube}}(r)\,4\pi r^2=\frac{-q}{\eo}\frac{r^3}{a^3}
  \quad\Longrightarrow\quad
  \vE_{\text{nube}}(\vr)=-\frac{q}{4\pi\eo a^3}\,\vr .
\]
Es un campo lineal en $\vr$ y dirigido hacia el centro (\cref{fig:atomo}\,(c)). Sobre el núcleo, situado en
$\vec d$, actúan el campo externo y el de la nube; el equilibrio exige
\[
  q\vE+q\,\vE_{\text{nube}}(\vec d)=\vcero
  \quad\Longrightarrow\quad
  \vE=\frac{q}{4\pi\eo a^3}\,\vec d ,
\]
y multiplicando por $4\pi\eo a^3$ se obtiene \eqref{eq:p-atomo}. La fuerza de la nube es recuperadora y
proporcional al desplazamiento, como la de un muelle: es la «ley de Hooke» del átomo.
\end{proof}

El resultado tiene la forma general de la respuesta lineal de un átomo isótropo.

\begin{definicion}[Polarizabilidad]\label{def:polarizabilidad}
Si el momento dipolar inducido en un átomo o molécula es proporcional al campo que actúa sobre él,
\begin{equation}\label{eq:alfa}
  \vp=\alpha\,\vE ,
\end{equation}
la constante $\alpha>0$ se llama \emph{polarizabilidad} (atómica o molecular). Sus unidades son
\[
  [\alpha]=\frac{\si{C\,m}}{\si{V/m}}=\frac{\si{C^2\,m}}{\si{N}}=\si{F\,m^2},
\]
pues $\si{C/V}=\si{F}$. En el modelo de nube uniforme, $\alpha=4\pi\eo a^3$: es proporcional al volumen atómico.
\end{definicion}

\begin{ejemplo}[El átomo de hidrógeno]\label{ej:hidrogeno}
Con el radio de Bohr, $a_0\simeq\SI{5.29e-11}{m}$, el modelo predice
\[
  \alpha_{\text{mod}}=4\pi\eo a_0^3\simeq 4\pi\,(\SI{8.854e-12}{F/m})(\SI{5.29e-11}{m})^3\simeq\SI{1.65e-41}{F\,m^2},
\]
mientras que el valor experimental es $\alpha_{\mathrm H}\simeq\SI{7.42e-41}{F\,m^2}$~\cite{Griffiths}: el modelo da
$\alpha_{\text{mod}}\simeq0{,}22\,\alpha_{\mathrm H}$. Para un modelo tan tosco, el orden de magnitud y la dependencia con
$a^3$ son correctos.
\end{ejemplo}

\begin{observacion}[El resultado cuántico]\label{obs:hidrogeno-cuantico}
La mecánica cuántica da para el estado fundamental del hidrógeno, de forma exacta,
$\alpha=\tfrac92\,4\pi\eo a_0^3$~\cite{Griffiths}. El factor $9/2$ es justamente el cociente $1/0{,}22$ del
\cref{ej:hidrogeno}: la nube real no es una esfera uniforme y rígida, sino que se deforma y es más
«blanda» que en nuestro modelo.
\end{observacion}

\begin{ejemplo}[Cuánto se desplaza la nube]\label{ej:desplazamiento}
Un campo de $E=\SI{5e5}{V/m}$ es muy intenso en el laboratorio, aunque inferior al de ruptura dieléctrica del
aire ($\sim\SI{3e6}{V/m}$~\cite{T3}). Sobre un átomo de hidrógeno induce
\[
  p=\alpha_{\mathrm H}E\simeq(\SI{7.42e-41}{F\,m^2})(\SI{5e5}{V/m})\simeq\SI{3.7e-35}{C\,m},
\]
y el núcleo se separa del centro de la nube una distancia
\[
  d=\frac{p}{e}\simeq\frac{\SI{3.7e-35}{C\,m}}{\SI{1.602e-19}{C}}\simeq\SI{2.3e-16}{m}=\SI{0.23}{fm}.
\]
El desplazamiento es del orden del tamaño de un núcleo y cinco órdenes de magnitud menor que el átomo: la
hipótesis $d\ll a$ está más que justificada, y con ella la respuesta lineal.
\end{ejemplo}

\subsection{Sustancias polares: orientación de dipolos permanentes}\label{sec:polares}

En una sustancia polar cada molécula tiene un momento dipolar permanente, $\vp\neq\vcero$, porque los
centros de carga positiva y negativa no coinciden. El ejemplo típico es el agua
(\cref{fig:agua}), con $p_{\mathrm{H_2O}}\simeq\SI{6.1e-30}{C\,m}$ \cite{Griffiths}, unas $10^5$ veces el
dipolo inducido del \cref{ej:desplazamiento}.

% Figura: la molécula de agua como dipolo permanente.
\begin{figure}[t]
\centering
\begin{subfigure}[b]{0.44\textwidth}\centering
\begin{tikzpicture}[scale=1.0]
  % geometría: ángulo H-O-H de 104,5 grados, dipolo según el bisector (eje +x)
  \coordinate (O) at (0,0);
  \coordinate (H1) at (52.25:1.25);
  \coordinate (H2) at (-52.25:1.25);
  \shade[ball color=carganeg!85!fondo] (H1) circle (0.52);
  \shade[ball color=carganeg!85!fondo] (H2) circle (0.52);
  \shade[ball color=cargapos!85!fondo] (O) circle (0.95);
  \node[font=\bfseries, text=white] at (-0.15,0) {$-2\delta$};
  \node[font=\bfseries, text=white] at ($(H1)+(0.12,0.08)$) {$+\delta$};
  \node[font=\bfseries, text=white] at ($(H2)+(0.12,-0.08)$) {$+\delta$};
  \draw[vecP, line width=1.3pt] (0.35,0) -- (2.15,0) node[right, etiq] {$\vp$};
  \node[etiqt] at ($(H1)+(0.95,0.25)$) {H};
  \node[etiqt] at ($(H2)+(0.95,-0.25)$) {H};
  \node[etiqt] at (-1.25,0.0) {O};
  \draw[tenue, line width=0.4pt] (52.25:0.55) arc (52.25:-52.25:0.55);
\end{tikzpicture}
\caption{Distribución de carga y dipolo}
\end{subfigure}\hfill
\begin{subfigure}[b]{0.52\textwidth}\centering
\begin{tikzpicture}[scale=0.95]
  \tikzset{electron/.style={circle, fill=acentoC, draw=fondo, line width=0.3pt, inner sep=1.5pt}}
  % oxígeno: núcleo y dos capas
  \draw[tenue, line width=0.5pt] (0,0) circle (0.62);
  \draw[tenue, line width=0.5pt] (0,0) circle (1.1);
  \shade[ball color=cargapos!85!fondo] (0,0) circle (0.36);
  \node[font=\tiny, text=white, align=center] at (0,0) {$8p$\\[-3pt]$8n$};
  \foreach \t in {90,270} \node[electron] at (\t:0.62) {};
  \foreach \t in {30,90,150,210,-30,-90} \node[electron] at (\t:1.1) {};
  % hidrógenos (enlaces covalentes: capas solapadas)
  \foreach \h/\ang in {A/-35,B/-145} {
    \coordinate (\h) at (\ang:1.82);
    \draw[tenue, line width=0.5pt] (\h) circle (0.62);
    \shade[ball color=carganeg!85!fondo] (\h) circle (0.2);
    \node[font=\tiny, text=white] at (\h) {$+$};
  }
  \node[electron] at ($(A)+(150:0.62)$) {};
  \node[electron] at ($(B)+(30:0.62)$) {};
  \node[etiqt] at (0,1.42) {oxígeno};
  \node[etiqt] at ($(A)+(0.1,-0.86)$) {H};
  \node[etiqt] at ($(B)+(-0.1,-0.86)$) {H};
\end{tikzpicture}
\caption{Estructura electrónica (esquema)}
\end{subfigure}
\caption{La molécula de agua, ejemplo de sustancia polar. (a) El oxígeno atrae con más fuerza los
electrones de los enlaces: queda una carga parcial $-2\delta$ en el oxígeno y $+\delta$ en cada hidrógeno.
Como la molécula no es lineal (ángulo de $104{,}5^\circ$), los centros de carga positiva y negativa no
coinciden y existe un momento dipolar permanente $p_{\mathrm{H_2O}}\simeq\SI{6.1e-30}{C\,m}$, dirigido
del oxígeno hacia los hidrógenos. (b) Esquema de capas: cada hidrógeno comparte su electrón con la
capa externa del oxígeno.}
\label{fig:agua}
\end{figure}


\begin{proposicion}[Dipolo en un campo uniforme]\label{prop:par}
Sobre un dipolo rígido $\vp=q\vec L$ en un campo uniforme $\vE$ la fuerza neta es nula y el momento de las
fuerzas respecto a cualquier punto es
\begin{equation}\label{eq:par}
  \vtau=\vp\times\vE ,
\end{equation}
que tiende a alinear $\vp$ con $\vE$. Su energía potencial es $U=-\vp\cdot\vE$, mínima cuando $\vp\parallel\vE$.
\end{proposicion}

\begin{proof}
Con las cargas $\pm q$ en $\vr_\pm$, la fuerza neta es $q\vE-q\vE=\vcero$ y el momento
$\vtau=\vr_+\times q\vE+\vr_-\times(-q\vE)=q(\vr_+-\vr_-)\times\vE=\vp\times\vE$, independiente del origen.
Si $\theta$ es el ángulo entre $\vp$ y $\vE$, $\abs{\vtau}=pE\sen\theta$ y el trabajo para girar el dipolo desde
$\theta=\pi/2$ es $\int_{\pi/2}^{\theta}pE\sen\theta'\,\dd\theta'=-pE\cos\theta=-\vp\cdot\vE$.
\end{proof}

Sin campo aplicado, la agitación térmica reparte las orientaciones al azar y el momento dipolar medio
de cualquier volumen macroscópico es nulo. Con campo, el par \eqref{eq:par} compite con la agitación
térmica y la orientación es solo parcial, pero el momento medio deja de anularse y tiene la dirección de
$\vE$ (\cref{fig:orientacion}).

% Figura: (a) par sobre un dipolo en un campo uniforme; (b) y (c) orientación
% térmica de dipolos permanentes sin campo y con campo.
\begin{figure}[t]
\centering
\begin{subfigure}[b]{0.38\textwidth}\centering
\begin{tikzpicture}[scale=1.0]
  \foreach \y in {-1.6,-1.07,...,1.62}
    \draw[campoE!55!fondo, line width=0.45pt, flechamedio=0.96] (-2.4,\y) -- (2.4,\y);
  \node[etiq, text=campoE] at (2.62,1.6) {$\vE$};
  \fill[dielec, draw=texto!60!fondo, line width=0.4pt, rotate=40] (0,0) ellipse (0.6 and 1.45);
  \coordinate (m) at (220:1.08); \coordinate (p) at (40:1.08);
  \draw[texto, line width=0.6pt] (m) -- (p);
  \draw[vecP, line width=1.2pt] (0,0) -- (40:0.7);
  \node[etiq] at (128:0.36) {$\vp$};
  \draw[texto!60!fondo, dashed, line width=0.35pt] (0,0) -- (1.05,0);
  \draw[tenue, line width=0.4pt] (0.5,0) arc (0:40:0.5);
  \node[etiqm, font=\footnotesize] at (17:0.72) {$\theta$};
  \fill[cargapos] (p) circle (0.085); \fill[carganeg] (m) circle (0.085);
  \draw[draw=cargapos, line width=1.0pt, -{Stealth[length=2.2mm]}] (p) -- ++(1.05,0);
  \node[etiqf, anchor=south] at ($(p)+(0.8,0.14)$) {$\vF_{+}=q\vE$};
  \draw[draw=carganeg, line width=1.0pt, -{Stealth[length=2.2mm]}] (m) -- ++(-1.05,0);
  \node[etiqf, anchor=north] at ($(m)+(-0.8,-0.14)$) {$\vF_{-}=-q\vE$};
  \node[etiqf, font=\footnotesize, anchor=north west] at ($(p)+(0.06,-0.1)$) {$+q$};
  \node[etiqf, font=\footnotesize, anchor=south east] at ($(m)+(-0.06,0.1)$) {$-q$};
  \node[etiq] at (0,-2.15) {$\vtau=\vp\times\vE$\quad($\otimes$, entrante)};
\end{tikzpicture}
\caption{Par de fuerzas sobre $\vp$}
\end{subfigure}\hfill
\begin{subfigure}[b]{0.29\textwidth}\centering
\begin{tikzpicture}[scale=0.82]
  \fill[dielec, rounded corners=2pt] (0,0) rectangle (3.2,3.2);
  \pgfmathsetseed{17}
  \foreach \i in {0,...,5} \foreach \j in {0,...,5} {
    \pgfmathsetmacro{\ang}{360*rnd}
    \draw[campoP, line width=0.75pt, -{Stealth[length=1.5mm,width=1.2mm]}]
      ({0.32+0.51*\i-0.17*cos(\ang)},{0.32+0.51*\j-0.17*sin(\ang)}) --
      ({0.32+0.51*\i+0.17*cos(\ang)},{0.32+0.51*\j+0.17*sin(\ang)});
  }
  \node[etiq] at (1.6,3.55) {$\vE=\vcero$};
\end{tikzpicture}
\caption{$\langle\vp\rangle=\vcero$}
\end{subfigure}\hfill
\begin{subfigure}[b]{0.29\textwidth}\centering
\begin{tikzpicture}[scale=0.82]
  \fill[dielec, rounded corners=2pt] (0,0) rectangle (3.2,3.2);
  \pgfmathsetseed{17}
  \foreach \i in {0,...,5} \foreach \j in {0,...,5} {
    \pgfmathsetmacro{\ang}{110*(rnd-0.5)}
    \draw[campoP, line width=0.75pt, -{Stealth[length=1.5mm,width=1.2mm]}]
      ({0.32+0.51*\i-0.17*cos(\ang)},{0.32+0.51*\j-0.17*sin(\ang)}) --
      ({0.32+0.51*\i+0.17*cos(\ang)},{0.32+0.51*\j+0.17*sin(\ang)});
  }
  \draw[vecE] (0.7,3.55) -- (2.5,3.55) node[right, etiq] {$\vE$};
\end{tikzpicture}
\caption{$\langle\vp\rangle\neq\vcero$}
\end{subfigure}
\caption{Sustancias polares. (a) En un campo uniforme la fuerza neta sobre un dipolo es nula, pero
el par $\vtau=\vp\times\vE$ tiende a alinearlo con el campo. (b) Sin campo, la agitación térmica
orienta los dipolos al azar y el momento medio es nulo. (c) Con campo, la orientación es parcial:
el par compite con la agitación térmica, y el momento medio, no nulo, tiene la dirección de $\vE$.}
\label{fig:orientacion}
\end{figure}


\begin{observacion}[Polarizabilidad de orientación]\label{obs:langevin}
La mecánica estadística (teoría de Langevin y Debye) da, para $pE\ll k_BT$, un momento medio
$\langle\vp\rangle=\dfrac{p^2}{3k_BT}\,\vE$~\cite{Purcell,Jackson}. Para el agua a temperatura ambiente,
$p^2/3k_BT\simeq\SI{3e-39}{F\,m^2}$, unas cuarenta veces la polarizabilidad electrónica del hidrógeno: por
eso las sustancias polares responden mucho más intensamente al campo, y por eso la permitividad relativa
del agua es tan grande ($\er\simeq80$; \cref{tab:permitividades}).
\end{observacion}

\subsection{La aproximación dipolar}\label{sec:aprox-dipolar}

Tanto si la sustancia es polar como si no, el resultado es el mismo: un dieléctrico en un campo eléctrico
adquiere, en cada elemento de volumen macroscópico, un momento dipolar neto en general no nulo.

\begin{definicion}[Aproximación dipolar]\label{def:aprox-dipolar}
En la \emph{aproximación dipolar}, un dieléctrico polarizado se describe como una distribución continua de
dipolos puntuales: cada elemento de volumen $\dd V'$ centrado en $\vr\,'$, pequeño a escala macroscópica pero
con muchas moléculas, posee un momento dipolar $\dd\vp$, y su contribución al potencial en $\vr$ es la de un
dipolo puntual~\cite{T2},
\begin{equation}\label{eq:dphi-dipolo}
  \dd\phi(\vr)=\frac{1}{4\pi\eo}\frac{\dd\vp\cdot\uv R}{R^2}.
\end{equation}
Se desprecian los multipolos de orden superior de cada elemento, cuya contribución decae más deprisa con la
distancia y es irrelevante a escala macroscópica.
\end{definicion}

% =============================================================================
\section{Polarización y densidades de carga ligada}\label{sec:polarizacion}
% =============================================================================

\subsection{El vector polarización}\label{sec:vector-P}

\begin{definicion}[Vector polarización]\label{def:polarizacion}
Se llama \emph{vector polarización} del dieléctrico en el punto $\vr\,'$ al momento dipolar por unidad de
volumen,
\begin{equation}\label{eq:P}
  \vP(\vr\,')\defeq\frac{\dd\vp}{\dd V'} ,
\end{equation}
donde $\dd\vp$ es el momento dipolar neto del elemento $\dd V'$. Sus unidades son $\si{C\,m/m^3}=\si{C/m^2}$. Si hay
$N$ moléculas por unidad de volumen con momento medio $\langle\vp\rangle$, $\vP=N\langle\vp\rangle$. Se dice que el
dieléctrico está \emph{uniformemente polarizado} si $\vP$ no depende del punto.
\end{definicion}

Con esta definición, el momento dipolar de cada elemento es $\dd\vp=\vP\,\dd V'$, y el potencial del
dieléctrico completo se obtiene sumando \eqref{eq:dphi-dipolo}:
\begin{equation}\label{eq:phi-P}
  \phi(\vr)=\frac{1}{4\pi\eo}\int_{V'}\frac{\vP(\vr\,')\cdot\uv R}{R^2}\,\dd V' .
\end{equation}
Esta expresión es correcta, pero poco transparente. El resultado central del tema es que puede
reescribirse como el potencial de una distribución de carga ordinaria.

\subsection{Densidades de carga de polarización}\label{sec:densidades}

\begin{lema}[Gradiente de $1/R$]\label{lem:grad-1R}
Para $\vr\neq\vr\,'$,
\begin{equation}\label{eq:grad-1R}
  \grad\Big(\frac1R\Big)=-\frac{\uv R}{R^2},\qquad
  \grad'\Big(\frac1R\Big)=+\frac{\uv R}{R^2}.
\end{equation}
\end{lema}

\begin{proof}
Para cualquier $n$, $\grad R^n=nR^{n-1}\grad R=nR^{n-1}\uv R$, pues $\grad R=\vR/R$~\cite{T0}. Con $n=-1$ se
obtiene la primera igualdad. Como $R$ depende de $\vr-\vr\,'$, derivar respecto a $\vr\,'$ cambia el signo,
$\grad'f(\vr-\vr\,')=-\grad f(\vr-\vr\,')$, de donde la segunda.
\end{proof}

\begin{teorema}[Densidades de carga de polarización]\label{teo:densidades}
El potencial \eqref{eq:phi-P} de un dieléctrico polarizado que ocupa la región $V'$, de frontera
$S'=\partial V'$ con normal exterior $\nn$, es el de una densidad volúmica $\rhoP$ y una densidad superficial
$\sigmaP$ de carga,
\begin{equation}\label{eq:phi-densidades}
  \phi(\vr)=\frac{1}{4\pi\eo}\int_{V'}\frac{\rhoP(\vr\,')}{R}\,\dd V'
            +\frac{1}{4\pi\eo}\oiint_{S'}\frac{\sigmaP(\vr\,')}{R}\,\dd S' ,
\end{equation}
con
\begin{equation}\label{eq:rhoP-sigmaP}
  \boxed{\;\rhoP(\vr\,')=-\grad'\!\cdot\vP(\vr\,'),\qquad \sigmaP(\vr\,')=\vP(\vr\,')\cdot\nn .\;}
\end{equation}
Se llaman \emph{densidades de carga de polarización} (o de carga ligada).
\end{teorema}

\begin{proof}
Por el \cref{lem:grad-1R}, el integrando de \eqref{eq:phi-P} es $\vP\cdot\grad'(1/R)$. La identidad del producto
$\grad'\!\cdot(f\vA)=f\,\grad'\!\cdot\vA+\vA\cdot\grad'f$, con $f=1/R$ y $\vA=\vP$, da
\[
  \frac{\vP\cdot\uv R}{R^2}=\vP\cdot\grad'\Big(\frac1R\Big)
  =-\frac{\grad'\!\cdot\vP}{R}+\grad'\!\cdot\Big(\frac{\vP}{R}\Big).
\]
Integrando en $V'$ y usando la linealidad de la integral,
\[
  \phi(\vr)=\frac{1}{4\pi\eo}\left[\int_{V'}\frac{-\grad'\!\cdot\vP}{R}\,\dd V'
            +\int_{V'}\grad'\!\cdot\Big(\frac{\vP}{R}\Big)\dd V'\right].
\]
A la segunda integral se le aplica el teorema de la divergencia de Gauss,
$\iiint_{V}\grad\cdot\vA\,\dd V=\oiint_{\partial V}\vA\cdot\dS$, con $\dS'=\nn\,\dd S'$:
\[
  \iiint_{V'}\grad'\!\cdot\Big(\frac{\vP}{R}\Big)\dd V'=\oiint_{S'}\frac{\vP\cdot\nn}{R}\,\dd S' .
\]
Basta llamar $\rhoP$ y $\sigmaP$ a los numeradores para obtener \eqref{eq:phi-densidades}.
\end{proof}

\begin{observacion}[Sobre el rigor de la deducción]\label{obs:singular}
Si $\vr$ está dentro del dieléctrico, el integrando de \eqref{eq:phi-P} es singular en $\vr\,'=\vr$, pero solo como
$1/R^2$, que es integrable en tres dimensiones; basta excluir una bola de radio $\delta$ alrededor de $\vr$ y hacer
$\delta\to0$, pues la contribución de su superficie es $\mathcal O(\delta)$. Por otra parte, la aproximación dipolar
promedia sobre volúmenes que contienen muchas moléculas: los campos que obtenemos son \emph{macroscópicos}, y en
particular \eqref{eq:phi-densidades} no describe las fluctuaciones a escala atómica.
\end{observacion}

\begin{corolario}[Las cargas ligadas son cargas de verdad]\label{cor:campo-ligadas}
El campo de un dieléctrico polarizado es $\vE=-\grad\phi$, es decir,
\begin{equation}\label{eq:E-densidades}
  \vE(\vr)=\frac{1}{4\pi\eo}\int_{V'}\frac{\rhoP\,\uv R}{R^2}\,\dd V'
           +\frac{1}{4\pi\eo}\oiint_{S'}\frac{\sigmaP\,\uv R}{R^2}\,\dd S' ,
\end{equation}
el mismo que crearían cargas libres con esas densidades. Toda la electrostática del vacío se aplica a ellas.
\end{corolario}

\begin{proof}
Se deriva \eqref{eq:phi-densidades} bajo el signo integral usando la primera igualdad de \eqref{eq:grad-1R}.
\end{proof}

\subsection{Interpretación física}\label{sec:interpretacion}

¿De dónde salen estas cargas, si el dieléctrico es neutro? Pensemos en una cadena de dipolos idénticos
(\cref{fig:cadena}\,(a)): la carga positiva de cada dipolo compensa la negativa del siguiente, y solo
«sobran» la carga negativa del primero y la positiva del último. La cadena equivale a dos cargas $\mp q$ en
sus extremos.

% Figura: interpretación de las densidades de polarización. Cadena de dipolos,
% polarización uniforme (solo carga superficial) y no uniforme (carga de volumen).
\begin{figure}[t]
\centering
\begin{subfigure}[b]{\textwidth}\centering
\begin{tikzpicture}[scale=1.0]
  \foreach \i in {0,...,4} {
    \pgfmathsetmacro{\x}{1.05*\i}
    \draw[campoP, line width=0.9pt, -{Stealth[length=1.6mm]}] (\x,0) -- (\x+0.7,0);
    \node[cneg] at (\x+0.02,-0.27) {$-$};
    \node[cpos] at (\x+0.68,-0.27) {$+$};
  }
  \draw[gauss, rounded corners=3pt] (1.65,-0.48) rectangle (2.20,0.3);
  \node[etiqt, anchor=north] at (1.93,-0.55) {neutro};
  \node[etiq] at (5.75,-0.05) {$=$};
  \draw[texto, line width=0.6pt] (6.4,0) -- (10.5,0);
  \fill[carganeg] (6.4,0) circle (0.08); \fill[cargapos] (10.5,0) circle (0.08);
  \node[etiq, anchor=north] at (6.4,-0.12) {$-q$};
  \node[etiq, anchor=north] at (10.5,-0.12) {$+q$};
\end{tikzpicture}
\caption{Cadena de dipolos idénticos: el interior es localmente neutro y solo «sobran» las cargas de los extremos}
\end{subfigure}\\[0.9em]
\begin{subfigure}[b]{0.48\textwidth}\centering
\begin{tikzpicture}[scale=0.95]
  \fill[dielec] (0,0) rectangle (3.4,2.0);
  \draw[borde] (0,0) rectangle (3.4,2.0);
  \foreach \i in {0,...,4} \foreach \j in {0,...,4}
    \draw[campoP, line width=0.75pt, -{Stealth[length=1.5mm,width=1.2mm]}]
      ({0.2+0.65*\i},{0.22+0.39*\j}) -- ({0.62+0.65*\i},{0.22+0.39*\j});
  \foreach \j in {0,...,4} {
    \node[cneg] at (-0.18,{0.22+0.39*\j}) {$-$};
    \node[cpos] at (3.58,{0.22+0.39*\j}) {$+$};
  }
  \draw[cargapos, dashed, line width=0.7pt] (1.72,1.0) ellipse (0.35 and 0.8);
  \draw[vecn] (0,1.0) -- (-0.75,1.0) node[above, etiq] {$\nn$};
  \draw[vecn] (3.4,1.0) -- (4.15,1.0) node[above, etiq] {$\nn$};
\end{tikzpicture}
\caption{$\vP$ uniforme: $\rhoP=0$, $\sigmaP=\vP\cdot\nn$}
\end{subfigure}\hfill
\begin{subfigure}[b]{0.48\textwidth}\centering
\begin{tikzpicture}[scale=0.95]
  \fill[dielec] (0,0) rectangle (3.4,2.0);
  \draw[borde] (0,0) rectangle (3.4,2.0);
  \foreach \i/\L in {0/0.18,1/0.28,2/0.38,3/0.48,4/0.58} \foreach \j in {0,...,4}
    \draw[campoP, line width=0.75pt, -{Stealth[length=1.5mm,width=1.2mm]}]
      ({0.42+0.65*\i-\L/2},{0.22+0.39*\j}) -- ({0.42+0.65*\i+\L/2},{0.22+0.39*\j});
  \foreach \j in {0,...,4} {
    \node[cneg] at (-0.18,{0.22+0.39*\j}) {$-$};
    \node[cpos] at (3.58,{0.22+0.39*\j}) {$+$};
    \node[cpos] at (3.78,{0.22+0.39*\j}) {$+$};
  }
  \draw[cargapos, dashed, line width=0.7pt] (1.72,1.0) ellipse (0.35 and 0.8);
  \draw[vecn] (0,1.0) -- (-0.75,1.0) node[above, etiq] {$\nn$};
  \draw[vecn] (3.9,1.0) -- (4.55,1.0) node[above, etiq] {$\nn$};
\end{tikzpicture}
\caption{$\vP$ no uniforme: $\rhoP=-\divg\vP\neq0$}
\end{subfigure}
\caption{Origen de las cargas de polarización. (a) En una cadena de dipolos la carga positiva de cada
uno compensa la negativa del siguiente. (b) Si $\vP$ es uniforme, un volumen interior cualquiera (trazos)
contiene tantas cabezas como colas de dipolos y es neutro: la carga ligada aparece solo en la superficie.
(c) Si $\vP$ crece hacia la derecha, en el mismo volumen entran más extremos negativos que positivos
salen y queda una carga neta: es la densidad volúmica $\rhoP=-\divg\vP$.}
\label{fig:cadena}
\end{figure}


Lo mismo ocurre en un bloque uniformemente polarizado (\cref{fig:cadena}\,(b)): cualquier volumen interior
contiene tantas cabezas como colas de dipolos y es neutro, y solo aparece carga en las caras por las que
$\vP$ entra (carga negativa) o sale (positiva). Si la polarización no es uniforme
(\cref{fig:cadena}\,(c)), en un volumen interior entran y salen distintos números de extremos, y queda una
carga neta: esa es la densidad $\rhoP$.

\begin{ejemplo}[Cálculo directo de $\sigmaP$]\label{ej:sigmaP-directo}
Sea un cilindro de longitud $\ell$ y sección $A$ uniformemente polarizado según su eje
(\cref{fig:cilindro-uniforme}\,(a)). Su momento dipolar es $p=P\,A\ell$, y el de un par de cargas $\mp q$ en las
caras es $p=q\ell$. Igualando, $q=PA$, de modo que la densidad en cada cara es $\sigmaP=q/A=P$. Si la cara
se corta oblicuamente, con normal $\nn$ formando un ángulo $\theta$ con $\vP$ (\cref{fig:cilindro-uniforme}\,(b)), la
misma carga se reparte en un área $A'=A/\cos\theta$, y
\[
  \sigmaP=\frac{q}{A'}=\frac{PA}{A/\cos\theta}=P\cos\theta=\vP\cdot\nn ,
\]
que es exactamente la expresión del \cref{teo:densidades}. En la superficie lateral $\vP\perp\nn$ y $\sigmaP=0$.
\end{ejemplo}

% Figura: cálculo directo de sigma_P en un cilindro uniformemente polarizado,
% corte oblicuo (sigma_P = P cos theta) y polarización radial (rho_P < 0).
\begin{figure}[t]
\centering
\begin{subfigure}[b]{0.46\textwidth}\centering
\begin{tikzpicture}[scale=0.9]
  % cilindro (vista lateral con perspectiva caballera)
  \fill[dielec] (0,-0.6) -- (2.6,-0.6) arc (-90:90:0.24 and 0.6) -- (0,0.6) arc (90:270:0.24 and 0.6);
  \draw[borde] (0,-0.6) -- (2.6,-0.6) (0,0.6) -- (2.6,0.6);
  \draw[borde] (0,0) ellipse (0.24 and 0.6);
  \fill[dielecB] (2.6,0) ellipse (0.24 and 0.6);
  \draw[borde] (2.6,0) ellipse (0.24 and 0.6);
  \node[etiq] at (2.6,0) {$A$};
  \draw[vecP] (0.6,0) -- (1.9,0) node[above, midway, etiq] {$\vP$};
  \draw[cota] (0,0.92) -- (2.6,0.92) node[midway, above, etiq] {$\ell$};
  % equivalente
  \node[etiq] at (3.45,0) {$\equiv$};
  \draw[borde] (4.3,0) ellipse (0.24 and 0.6);
  \draw[borde] (6.6,0) ellipse (0.24 and 0.6);
  \draw[texto!70!fondo, line width=0.5pt] (4.3,0) -- (6.6,0);
  \node[cneg, font=\small] at (4.3,0.85) {$-q$};
  \node[cpos, font=\small] at (6.6,0.85) {$+q$};
  \draw[cota] (4.3,-0.92) -- (6.6,-0.92) node[midway, below, etiq] {$\ell$};
\end{tikzpicture}
\caption{$p=q\ell=P\,A\ell\ \Rightarrow\ \sigmaP=q/A=P$}
\end{subfigure}\hfill
\begin{subfigure}[b]{0.27\textwidth}\centering
\begin{tikzpicture}[scale=0.9]
  \fill[dielec] (0,-0.6) -- (2.2,-0.6) -- (2.75,0.6) -- (0,0.6) arc (90:270:0.24 and 0.6);
  \draw[borde] (0,-0.6) -- (2.2,-0.6) (0,0.6) -- (2.75,0.6);
  \draw[borde] (0,0) ellipse (0.24 and 0.6);
  \fill[dielecB] (2.475,0) ellipse [x radius=0.24, y radius=0.66, rotate=-24.6];
  \draw[borde] (2.475,0) ellipse [x radius=0.24, y radius=0.66, rotate=-24.6];
  \node[etiq] at (3.05,0.82) {$A'$};
  \draw[vecP] (0.45,0) -- (1.65,0) node[above, midway, etiq] {$\vP$};
  \draw[vecn] (2.475,0) -- ++(-24.6:1.05) node[right, etiq] {$\nn$};
  \draw[texto!60!fondo, dashed, line width=0.35pt] (2.475,0) -- (3.75,0);
  \draw[tenue, line width=0.4pt] (3.375,0) arc (0:-24.6:0.9);
  \node[etiqt] at ($(2.475,0)+(-11:1.1)$) {$\theta$};
\end{tikzpicture}
\caption{$\sigmaP=\dfrac{q}{A'}=P\cos\theta$}
\end{subfigure}\hfill
\begin{subfigure}[b]{0.23\textwidth}\centering
\begin{tikzpicture}[scale=0.9]
  \fill[dielec] (0,0) circle (1.05);
  \draw[equipot] (0,0) circle (1.05);
  \foreach \t in {0,45,...,315} {
    \draw[campoP, line width=0.8pt, -{Stealth[length=1.6mm]}] (\t:0.3) -- (\t:0.95);
    \node[cpos] at (\t:1.28) {$+$};
  }
  \foreach \t in {22.5,112.5,202.5,292.5} \node[cneg, font=\small\bfseries] at (\t:0.15) {$-$};
\end{tikzpicture}
\caption{$\divg\vP>0$}
\end{subfigure}
\caption{(a) Un cilindro uniformemente polarizado equivale a dos discos con cargas $\mp q$: de
$P=p/V=q\ell/(A\ell)$ se sigue $\sigmaP=P$. (b) Si la cara se corta oblicuamente, la misma carga se
reparte sobre un área mayor $A'=A/\cos\theta$, y $\sigmaP=P\cos\theta=\vP\cdot\nn$. (c) Una polarización
divergente arrastra carga positiva hacia la superficie y deja un defecto negativo en el interior:
$\rhoP=-\divg\vP<0$ y $\sigmaP>0$, con carga ligada total nula.}
\label{fig:cilindro-uniforme}
\end{figure}


\subsection{La carga de polarización total es nula}\label{sec:carga-total}

Una polarización divergente arrastra carga positiva hacia la superficie y deja un defecto negativo en el
interior (\cref{fig:cilindro-uniforme}\,(c)): la carga que falta en el volumen aparece en la superficie,
$\int_{V'}\rhoP\,\dd V'=-\oiint_{S'}\sigmaP\,\dd S'$. Es una consecuencia general.

\begin{proposicion}[Neutralidad de la carga ligada]\label{prop:QP-nula}
La carga de polarización total de un dieléctrico es nula:
\begin{equation}\label{eq:QP}
  \QP=\int_{V'}\rhoP\,\dd V'+\oiint_{S'}\sigmaP\,\dd S'=0 .
\end{equation}
\end{proposicion}

\begin{proof}
Por el teorema de la divergencia,
$\int_{V'}\rhoP\,\dd V'=-\int_{V'}\grad'\!\cdot\vP\,\dd V'=-\oiint_{S'}\vP\cdot\nn\,\dd S'=-\oiint_{S'}\sigmaP\,\dd S'$.
\end{proof}

\begin{observacion}[El momento dipolar de las cargas ligadas]\label{obs:momento-ligado}
Como el monopolo de las cargas ligadas es nulo, el término dominante de su desarrollo multipolar~\cite{T2} es
el dipolar, y su momento es el esperado:
\[
  \int_{V'}\vr\,'\rhoP\,\dd V'+\oiint_{S'}\vr\,'\sigmaP\,\dd S'=\int_{V'}\vP\,\dd V' .
\]
Para la componente $x$, por ejemplo, $\grad'\!\cdot(x'\vP)=x'\grad'\!\cdot\vP+P_x$, y el teorema de la divergencia da
$\oiint_{S'}x'\vP\cdot\nn\,\dd S'-\int_{V'}x'\grad'\!\cdot\vP\,\dd V'=\int_{V'}P_x\,\dd V'$. Lejos del dieléctrico, su campo es el
de un dipolo $\vp_{\text{tot}}=\int\vP\,\dd V'$, como debe ser.
\end{observacion}

La comprobación explícita de \eqref{eq:QP} en un cubo con polarización radial es el \cref{prob:cubo} de
los seminarios, y en una esfera uniformemente polarizada, el \cref{prob:esfera-polarizada}.

% =============================================================================
\section{Ley de Gauss en dieléctricos: el vector desplazamiento}\label{sec:desplazamiento}
% =============================================================================

\subsection{Carga libre y carga ligada}\label{sec:gauss-D}

En presencia de dieléctricos, el campo electrostático lo crean dos tipos de carga: la libre, que
controlamos, y la ligada, que el material dispone por sí solo y que depende a su vez del campo. La ley de
Gauss del vacío sigue siendo válida si se cuentan ambas,
\begin{equation}\label{eq:gauss-total}
  \divg\vE=\frac{\rhof+\rhoP}{\eo},\qquad
  \oiint_S\vE\cdot\dS=\frac{\qfenc+\qPenc}{\eo},
\end{equation}
pero en la práctica no sirve de mucho: para usarla habría que conocer la carga de polarización, que es
precisamente lo que no sabemos hasta haber resuelto el problema. Conviene disponer de una ley que solo
involucre la carga libre.

\begin{lema}[Carga ligada encerrada]\label{lem:qP-enc}
Para cualquier superficie cerrada $S$, la carga de polarización que encierra es menos el flujo de $\vP$:
\begin{equation}\label{eq:qP-enc}
  \qPenc=-\oiint_S\vP\cdot\dS .
\end{equation}
\end{lema}

\begin{proof}
Sea $V_S$ la región encerrada por $S$ y $V$ la del dieléctrico. La carga ligada encerrada consta de la
volúmica de $V\cap V_S$ y de la superficial de la parte de $\partial V$ contenida en $V_S$. La frontera de
$V\cap V_S$ está formada por $\partial V\cap V_S$ y $S\cap V$; aplicando el teorema de la divergencia en esa región,
\[
  \int_{V\cap V_S}\rhoP\,\dd V=-\oiint_{\partial(V\cap V_S)}\vP\cdot\dS
  =-\int_{\partial V\cap V_S}\sigmaP\,\dd S-\int_{S\cap V}\vP\cdot\dS .
\]
Pasando el primer término de la derecha al otro miembro se obtiene $\qPenc=-\int_{S\cap V}\vP\cdot\dS$, que es
\eqref{eq:qP-enc} porque $\vP=\vcero$ fuera del dieléctrico.
\end{proof}

Llevando \eqref{eq:qP-enc}, o su forma local $\rhoP=-\divg\vP$, a \eqref{eq:gauss-total}, el flujo de $\vP$ pasa al
primer miembro:
\[
  \divg\vE=\frac{\rhof-\divg\vP}{\eo}
  \quad\Longrightarrow\quad
  \divg\big(\eo\vE+\vP\big)=\rhof .
\]
Esto sugiere agrupar los dos campos en uno solo.

\begin{definicion}[Vector desplazamiento eléctrico]\label{def:D}
Se llama \emph{vector desplazamiento eléctrico} (o inducción eléctrica) al campo
\begin{equation}\label{eq:D}
  \boxed{\;\vD(\vr)\defeq\eo\vE(\vr)+\vP(\vr)\; .}
\end{equation}
Sus unidades son las de $\vP$, $\si{C/m^2}$. El nombre se debe a Maxwell, que lo interpretaba como un
desplazamiento de las cargas del éter; hoy es simplemente una herramienta de cálculo.
\end{definicion}

\begin{teorema}[Ley de Gauss para $\vD$]\label{teo:gauss-D}
El vector desplazamiento satisface, en forma diferencial e integral,
\begin{equation}\label{eq:gauss-D}
  \divg\vD=\rhof,\qquad
  \oiint_S\vD\cdot\dS=\qfenc ,
\end{equation}
para cualquier superficie cerrada $S$: su flujo solo «ve» la carga libre.
\end{teorema}

\begin{proof}
La forma diferencial se acaba de deducir. Para la integral, por \eqref{eq:gauss-total} y el \cref{lem:qP-enc},
$\oiint_S\eo\vE\cdot\dS=\qfenc+\qPenc=\qfenc-\oiint_S\vP\cdot\dS$, que es \eqref{eq:gauss-D}. Ambas formas son equivalentes por el
teorema de la divergencia.
\end{proof}

\begin{observacion}[Lo que el teorema no dice]\label{obs:D-no-determinado}
La ley de Gauss para $\vD$ es tan útil como la del vacío \emph{cuando hay simetría suficiente}: plana, cilíndrica
o esférica. En general, conocer la divergencia de un campo no basta para determinarlo; hace falta también su
rotacional, y el de $\vD$ no es nulo (\cref{sec:rot-D}). En particular, que la carga libre encerrada sea nula
\emph{no} implica que $\vD$ sea nulo.
\end{observacion}

\subsection{Electretes}\label{sec:electretes}

\begin{definicion}[Electrete]\label{def:electrete}
Un \emph{electrete} es un material con una polarización permanente, que no se debe a un campo aplicado y no
desaparece al retirarlo. Los materiales ferroeléctricos son el ejemplo más importante; son el análogo eléctrico
de los imanes permanentes (ferromagnéticos) del magnetismo.
\end{definicion}

En un electrete $\vP$ es un dato del problema, y basta el \cref{teo:densidades} para calcular su campo. Los
problemas~\ref{prob:cilindro-polarizado}, \ref{prob:electrete-esferico}, \ref{prob:electrete-coaxial}
y~\ref{prob:electrete-condensador} son de este tipo.

\subsection{El rotacional de $\vD$}\label{sec:rot-D}

El campo electrostático es conservativo, $\rot\vE=\vcero$. Tomando el rotacional de \eqref{eq:D},
\begin{equation}\label{eq:rot-D}
  \rot\vD=\rot\vP .
\end{equation}
El desplazamiento no es, en general, conservativo: su circulación no se anula allí donde la tiene $\vP$.
Recordemos el teorema de Helmholtz~\cite{T0}.

\begin{lema}[Teorema de Helmholtz]\label{lem:helmholtz}
Un campo vectorial $\vF$ de clase $\mathcal C^1$ a trozos que tiende a cero en el infinito más deprisa que $1/r$
queda determinado de forma única por su divergencia y su rotacional (entendidos en sentido distribucional,
de modo que incluyen las discontinuidades de sus componentes normal y tangencial en las superficies).
\end{lema}

En una superficie en que $\vP$ salta de $\vP_1$ a $\vP_2$, la parte superficial del rotacional es
$\nn\times(\vP_2-\vP_1)\,\delta_S$: es nula si y solo si se conserva la componente tangencial de $\vP$. Con esto
podemos responder a la pregunta de la \cref{obs:D-no-determinado}.

\begin{proposicion}[Cuándo se anula $\vD$]\label{prop:D-nulo}
Si en todo el espacio no hay carga libre y la polarización es irrotacional, incluidas sus superficies de
discontinuidad,
\[
  \rhof\equiv0,\qquad \rot\vP=\vcero\ \text{en cada medio},\qquad \nn\times(\vP_2-\vP_1)=\vcero\ \text{en cada interfase},
\]
entonces $\vD\equiv\vcero$ y, por tanto, $\vE=-\vP/\eo$. Recíprocamente, si $\vD\equiv\vcero$, no hay carga libre y $\vP$
es irrotacional.
\end{proposicion}

\begin{proof}
Por \eqref{eq:gauss-D} y \eqref{eq:rot-D}, $\vD$ tiene divergencia y rotacional nulos en todo el espacio, incluidos
los términos superficiales. Si el dieléctrico es acotado, fuera de él $\vD=\eo\vE$ decae como el campo de una
distribución neutra (\cref{prop:QP-nula}), al menos como $1/r^3$. Por el \cref{lem:helmholtz}, el único campo con
estas propiedades es el nulo. El recíproco es inmediato: si $\vD=\vcero$, $\rhof=\divg\vD=0$ y $\rot\vP=\rot\vD=\vcero$.
\end{proof}

\begin{ejemplo}[Polarizaciones radiales]\label{ej:radial}
Una polarización $\vP=f(r)\,\uv r$ (esférica) o $\vP=f(\rho)\,\uv\rho$ (cilíndrica) es irrotacional en el interior del
material, y en sus superficies, que son esferas o cilindros coaxiales, $\vP$ es normal y $\nn\times\vP=\vcero$. Sin
carga libre, $\vD=\vcero$ en todo el espacio y $\vE=-\vP/\eo$ dentro del material: así se resuelven los
problemas~\ref{prob:electrete-esferico} y~\ref{prob:electrete-coaxial}.
\end{ejemplo}

\begin{ejemplo}[El cilindro uniformemente polarizado]\label{ej:stokes}
En un cilindro con $\vP$ uniforme y axial, $\rot\vP=\vcero$ dentro, pero en la pared lateral $\vP$ es tangente y salta a
cero. Un pequeño rectángulo que atraviese la pared (\cref{fig:p44}\,(a)) tiene circulación $\oint\vP\cdot\dl=P\,\Delta\ell\neq0$ y,
por el teorema de Stokes, $\rot\vP\neq\vcero$ en la pared. Aunque no haya carga libre, $\vD\neq\vcero$
(\cref{prob:cilindro-polarizado}): sus líneas se cierran como las del campo de un solenoide
(\cref{fig:cilindro-lineas}).
\end{ejemplo}

\begin{trampa}[Trampa habitual]
«Como no hay carga libre, $\vD=\vcero$» solo es cierto si además $\vP$ es irrotacional (\cref{prop:D-nulo}). En el
cilindro, la esfera o el cubo uniformemente polarizados, $\vD\neq\vcero$ aunque $\rhof=0$.
\end{trampa}

% =============================================================================
\section{Dieléctricos lineales, homogéneos e isótropos}\label{sec:lineales}
% =============================================================================

Pasamos ahora a los dieléctricos ordinarios, en los que la polarización no es permanente sino la respuesta
del material al campo que actúa sobre él. Para cerrar el sistema formado por \eqref{eq:gauss-D} y
$\rot\vE=\vcero$ hace falta una relación entre $\vP$ y $\vE$, que depende del material y debe obtenerse de su
estructura microscópica o de la experiencia.

\subsection{La relación constitutiva}\label{sec:constitutiva}

Desde el punto de vista macroscópico, si $\vP$ se anula con el campo y este es débil frente a los campos
atómicos ($\sim\SI{e11}{V/m}$), cada componente de la polarización puede desarrollarse en potencias de las del campo,
\begin{equation}\label{eq:constitutiva-general}
  P_i(\vr)=\sum_{j}\alpha_{ij}(\vr)\,E_j(\vr)+\sum_{j,k}\beta_{ijk}(\vr)\,E_j(\vr)E_k(\vr)+\cdots,
  \qquad i,j,k\in\{x,y,z\},
\end{equation}
donde los coeficientes son propiedades del material en el punto $\vr$. Es una relación \emph{local}: la
polarización en $\vr$ solo depende del campo en $\vr$.

\begin{definicion}[Medio lineal, homogéneo e isótropo]\label{def:LHI}
Un dieléctrico es
\begin{enumerate}
  \item \emph{lineal} si en \eqref{eq:constitutiva-general} basta el primer término; se escribe entonces
        \begin{equation}\label{eq:chi-tensor}
          P_i(\vr)=\eo\sum_j\chi_{ij}(\vr)\,E_j(\vr),
        \end{equation}
        y la matriz adimensional $(\chi_{ij})$ es el \emph{tensor de susceptibilidad eléctrica};
  \item \emph{homogéneo} si $\chi_{ij}$ no depende de $\vr$: es una propiedad intrínseca del material, la misma en
        todos sus puntos;
  \item \emph{isótropo} si no tiene direcciones privilegiadas: $\chi_{ij}=\chi\,\delta_{ij}$, es decir,
        $\chi_{xx}=\chi_{yy}=\chi_{zz}=\chi$ y $\chi_{ij}=0$ para $i\neq j$.
\end{enumerate}
En un medio \LHI{}, por tanto,
\begin{equation}\label{eq:chi}
  \boxed{\;\vP(\vr)=\eo\chi\,\vE(\vr)\;},
\end{equation}
con $\chi>0$ una constante, la \emph{susceptibilidad eléctrica} del material.
\end{definicion}

\begin{observacion}[Por qué $\chi>0$]\label{obs:chi-positiva}
Los dos mecanismos de la \cref{sec:microscopico} producen dipolos en el sentido del campo: el desplazamiento de
la nube electrónica y la orientación de los dipolos permanentes. Por eso $\vP$ tiene el sentido de $\vE$ y la
susceptibilidad es positiva. En un cristal no cúbico la respuesta depende de la dirección y $(\chi_{ij})$ no es
proporcional a la identidad, aunque sí simétrica y definida positiva~\cite{Landau}.
\end{observacion}

\subsection{Permitividad}\label{sec:permitividad}

En un medio \LHI{} el desplazamiento es proporcional al campo:
\[
  \vD=\eo\vE+\vP=\eo\vE+\eo\chi\vE=\eo(1+\chi)\,\vE .
\]

\begin{definicion}[Permitividad]\label{def:permitividad}
En un dieléctrico \LHI{} se definen la \emph{permitividad dieléctrica} $\varepsilon$ y la \emph{permitividad relativa}
(o constante dieléctrica) $\er$,
\begin{equation}\label{eq:epsilon}
  \varepsilon\defeq\eo(1+\chi),\qquad \er\defeq\frac{\varepsilon}{\eo}=1+\chi ,
\end{equation}
de modo que
\begin{equation}\label{eq:D-epsilon}
  \boxed{\;\vD=\varepsilon\,\vE,\qquad \vP=(\varepsilon-\eo)\,\vE\; .}
\end{equation}
Como $\chi>0$, $\varepsilon>\eo$ y $\er>1$. Las unidades de $\varepsilon$ son las de $\eo$, $\si{F/m}$; $\er$ es adimensional
(algunos textos la llaman $\kappa$ o $K$~\cite{Purcell}).
\end{definicion}

\begin{table}[t]
\centering
\caption{Permitividades relativas de algunos materiales a temperatura ambiente~\cite{Griffiths}.}
\label{tab:permitividades}
\begin{tabular}{@{}lc@{\qquad}lc@{}}
\toprule
Material & $\er$ & Material & $\er$\\
\midrule
Vacío & $1$ (exacto) & Diamante & $5{,}7$\\
Aire (seco) & $1{,}00054$ & Sal común (NaCl) & $5{,}9$\\
Benceno & $2{,}28$ & Silicio & $11{,}7$\\
Metanol & $33{,}0$ & Agua (\SI{20}{\celsius}) & $80{,}1$\\
\bottomrule
\end{tabular}
\end{table}

\subsection{Consecuencias en un medio lineal}\label{sec:consecuencias}

\begin{proposicion}[Carga ligada en un medio homogéneo]\label{prop:rhoP-rhof}
En el interior de un dieléctrico \LHI{},
\begin{equation}\label{eq:rhoP-rhof}
  \rhoP=-\frac{\chi}{1+\chi}\,\rhof=-\Big(1-\frac{\eo}{\varepsilon}\Big)\rhof .
\end{equation}
En particular, si no hay carga libre en su volumen tampoco hay carga ligada de volumen: la carga de
polarización se concentra en las superficies.
\end{proposicion}

\begin{proof}
Por \eqref{eq:chi}, \eqref{eq:D-epsilon} y \eqref{eq:gauss-D}, con $\chi$ y $\varepsilon$ constantes,
$\rhoP=-\divg\vP=-\eo\chi\,\divg\vE=-\dfrac{\eo\chi}{\varepsilon}\divg\vD=-\dfrac{\chi}{1+\chi}\,\rhof$.
\end{proof}

\begin{proposicion}[Un único medio que llena todo el espacio]\label{prop:un-medio}
Si todo el espacio está ocupado por un único dieléctrico \LHI{} de permitividad $\varepsilon$, entonces
\begin{equation}\label{eq:coulomb-epsilon}
  \divg\vE=\frac{\rhof}{\varepsilon},\qquad \rot\vE=\vcero,\qquad
  \vE(\vr)=\frac{1}{4\pi\varepsilon}\int\frac{\rhof(\vr\,')\,\uv R}{R^2}\,\dd V' :
\end{equation}
el campo es el del vacío con $\eo$ sustituida por $\varepsilon$, y por tanto $\er$ veces menor.
\end{proposicion}

\begin{proof}
Con $\varepsilon$ constante, $\divg\vE=\divg\vD/\varepsilon=\rhof/\varepsilon$. El campo $\varepsilon\vE=\vD$ tiene divergencia $\rhof$ y rotacional
nulo, igual que el campo de Coulomb de la carga libre en el vacío multiplicado por $\eo$; por el
\cref{lem:helmholtz} ambos coinciden.
\end{proof}

\begin{trampa}[¡Solo si hay un único medio!]
Las expresiones $\divg\vE=\rhof/\varepsilon$, $\oiint\vE\cdot\dS=\qfenc/\varepsilon$ y $\vE=\frac{1}{4\pi\varepsilon}\int\rhof\,\uv R/R^2\,\dd V'$ \emph{no son
ciertas en general}. Si $\varepsilon$ varía de un punto a otro,
\[
  \rhof=\divg(\varepsilon\vE)=\varepsilon\,\divg\vE+\vE\cdot\grad\varepsilon ,
\]
y en una interfase, donde $\varepsilon$ cambia bruscamente, $\grad\varepsilon$ es singular: ese término extra es precisamente la
carga de polarización acumulada en la frontera. Por eso en las fronteras se trabaja con $\vD$, cuyas
condiciones de contorno no dependen del material (\cref{sec:frontera}).
\end{trampa}

Los dos ejemplos canónicos son una carga libre en un dieléctrico infinito, que reduce su campo en el factor
$\er$ (\cref{prob:carga-en-dielectrico}), y el condensador plano relleno de dieléctrico, cuya capacidad
aumenta en el mismo factor, $C=\er C_0$ (\cref{prob:condensador-dielectrico}). De ahí procede el método
experimental clásico para medir $\er$: comparar la capacidad de un condensador lleno y vacío.

% =============================================================================
\section{Condiciones de frontera}\label{sec:frontera}
% =============================================================================

Las ecuaciones $\divg\vD=\rhof$ y $\rot\vE=\vcero$ valen en el interior de cada medio, pero en una interfase entre
dos dieléctricos, o entre un dieléctrico y el vacío o un conductor, los campos pueden ser discontinuos y
las derivadas dejan de tener sentido. Sus formas integrales, en cambio, siguen valiendo, y de ellas se
deducen las relaciones entre los campos a uno y otro lado. Repasamos la deducción del Tema~1~\cite{T1},
ahora con $\vD$.

Sea $\Sigma$ una superficie que separa los medios~1 y~2, con normal unitaria $\nn$ dirigida de~1 a~2 y
posiblemente con una densidad de carga libre $\sigmaf$. Para un campo $\vec F$, escribimos $F_n=\vec F\cdot\nn$ y
$\vec F_t=\vec F-F_n\nn$.

% Figura: deducción de las condiciones de frontera (cilindro de Gauss y circuito
% rectangular) e interpretación del salto de la componente normal de P.
\begin{figure}[t]
\centering
\begin{subfigure}[b]{0.33\textwidth}\centering
\begin{tikzpicture}[scale=0.95]
  \fill[dielec] (-1.8,-1.4) rectangle (1.8,0);
  \fill[dielecB] (-1.8,0) rectangle (1.8,1.4);
  \draw[texto, line width=0.6pt] (-1.8,0) -- (1.8,0);
  \node[etiq, anchor=west] at (-1.75,1.12) {$2$};
  \node[etiq, anchor=west] at (-1.75,-1.12) {$1$};
  % cilindro
  \fill[fondo!50!dielecB] (-0.7,-0.42) rectangle (0.7,0.42);
  \draw[gauss] (-0.7,-0.42) -- (-0.7,0.42) (0.7,-0.42) -- (0.7,0.42);
  \draw[gauss] (0,0.42) ellipse (0.7 and 0.16);
  \draw[gauss] (-0.7,-0.42) arc (180:360:0.7 and 0.16);
  \draw[gauss, densely dotted] (0.7,-0.42) arc (0:180:0.7 and 0.16);
  \draw[vecn] (0,0.42) -- (0,1.2) node[right, etiq] {$\nn$};
  \draw[vecn] (0,-0.58) -- (0,-1.3) node[right, etiq] {$-\nn$};
  \node[etiqt] at (-1.12,0.62) {$\Delta S$};
  \draw[cota] (1.0,-0.42) -- (1.0,0.42) node[midway, right, etiq] {$h\!\to\!0$};
\end{tikzpicture}
\caption{Normal de $\vD$}
\end{subfigure}\hfill
\begin{subfigure}[b]{0.33\textwidth}\centering
\begin{tikzpicture}[scale=0.95]
  \fill[dielec] (-1.8,-1.4) rectangle (1.8,0);
  \fill[dielecB] (-1.8,0) rectangle (1.8,1.4);
  \draw[texto, line width=0.6pt] (-1.8,0) -- (1.8,0);
  \node[etiq, anchor=west] at (-1.75,1.12) {$2$};
  \node[etiq, anchor=west] at (-1.75,-1.12) {$1$};
  \draw[acentoB, line width=0.7pt] (-0.8,-0.35) rectangle (0.8,0.35);
  \draw[acentoB, line width=0.7pt, flechamedio] (-0.8,-0.35) -- (0.8,-0.35);
  \draw[acentoB, line width=0.7pt, flechamedio] (0.8,0.35) -- (-0.8,0.35);
  \node[etiqt] at (0,0.62) {$\Delta\ell$};
  \draw[vecn] (-1.3,0) -- (-1.3,0.75) node[right, etiq] {$\nn$};
  \draw[vecn] (-1.3,0) -- (-0.55,0) node[pos=0.55, below=1pt, etiq] {$\tu$};
  \draw[cota] (1.1,-0.35) -- (1.1,0.35) node[midway, right, etiq] {$h\!\to\!0$};
\end{tikzpicture}
\caption{Tangencial de $\vE$}
\end{subfigure}\hfill
\begin{subfigure}[b]{0.32\textwidth}\centering
\begin{tikzpicture}[scale=0.95]
  \fill[dielec] (-1.7,-1.4) rectangle (1.7,0);
  \fill[dielecB] (-1.7,0) rectangle (1.7,1.4);
  \draw[texto, line width=0.6pt] (-1.7,0) -- (1.7,0);
  \node[etiq, anchor=west] at (-1.65,1.12) {$2$};
  \node[etiq, anchor=west] at (-1.65,-1.12) {$1$};
  \draw[vecP] (-0.45,-1.15) -- (-0.45,-0.05) node[pos=0.35, left, etiq] {$\vP_1$};
  \draw[vecP] (0.45,0.05) -- (0.45,0.85) node[pos=0.6, right, etiq] {$\vP_2$};
  \draw[vecn] (1.25,0) -- (1.25,0.65) node[right, etiq] {$\nn$};
  \node[cpos] at (-0.45,0.17) {$+P_{1n}$};
  \node[cneg] at (0.45,-0.2) {$-P_{2n}$};
\end{tikzpicture}
\caption{$\sigmaP=P_{1n}-P_{2n}$}
\end{subfigure}
\caption{Condiciones de frontera entre dos medios; $\nn$ es la normal unitaria del medio 1 al 2.
(a) Cilindro de Gauss de bases $\Delta S$ y altura $h\to0$: el flujo lateral desaparece y queda
$(D_{2n}-D_{1n})\Delta S=\sigmaf\,\Delta S$. (b) Circuito rectangular de lados $\Delta\ell$ y $h\to0$:
la circulación de $\vE$ da $(E_{2t}-E_{1t})\Delta\ell=0$. (c) Cada medio aporta a la interfase la carga
ligada $\vP\cdot\nn_{\mathrm{ext}}$ de su propia cara: $+P_{1n}$ el medio 1 (su normal exterior es $\nn$)
y $-P_{2n}$ el medio 2 (su normal exterior es $-\nn$).}
\label{fig:frontera}
\end{figure}


\subsection{Las dos condiciones fundamentales}\label{sec:condiciones}

\begin{teorema}[Condiciones de frontera]\label{teo:frontera}
En cada punto de $\Sigma$ en que los campos son acotados a ambos lados:
\begin{enumerate}
  \item la componente normal de $\vD$ salta en la densidad de carga libre,
        \begin{equation}\label{eq:frontera-D}
          D_{2n}-D_{1n}=\sigmaf
          \qquad\Longleftrightarrow\qquad
          \nn\cdot(\vD_2-\vD_1)=\sigmaf ;
        \end{equation}
  \item la componente tangencial de $\vE$ es continua,
        \begin{equation}\label{eq:frontera-E}
          \vE_{2t}=\vE_{1t}
          \qquad\Longleftrightarrow\qquad
          \nn\times(\vE_2-\vE_1)=\vcero ;
        \end{equation}
  \item el potencial es continuo, $\phi_1=\phi_2$ sobre $\Sigma$.
\end{enumerate}
\end{teorema}

\begin{proof}
(i) Aplicamos la forma integral de la ley de Gauss para $\vD$ a un cilindro de bases $\Delta S$, paralelas a $\Sigma$ a uno
y otro lado, y altura $h$ (\cref{fig:frontera}\,(a)); las bases son tan pequeñas que $\vD$ es prácticamente uniforme en
cada una. Al hacer $h\to0$, el área lateral tiende a cero y, como $\vD$ es acotado, su flujo desaparece. Solo quedan las
bases, con normales exteriores $\nn$ (en el medio~2) y $-\nn$ (en el medio~1):
\[
  \oiint\vD\cdot\dS\ \xrightarrow{\;h\to0\;}\ \vD_2\cdot\nn\,\Delta S+\vD_1\cdot(-\nn)\,\Delta S=(D_{2n}-D_{1n})\,\Delta S .
\]
En cuanto a la carga libre encerrada, la volúmica es $\rhof\,\Delta S\,h\to0$ y solo sobrevive la superficial, $\sigmaf\,\Delta S$.
Igualando y dividiendo por $\Delta S$ se obtiene \eqref{eq:frontera-D}.

(ii) Como el campo electrostático es conservativo, $\oint_C\vE\cdot\dl=0$ para cualquier curva cerrada. Tomamos un rectángulo
de lados $\Delta\ell$ paralelos a $\Sigma$, según una dirección tangente $\tu$, y lados $h$ perpendiculares (\cref{fig:frontera}\,(b)).
Al hacer $h\to0$ los lados perpendiculares no contribuyen y
\[
  \vE_2\cdot(-\tu)\,\Delta\ell+\vE_1\cdot\tu\,\Delta\ell=0
  \quad\Longrightarrow\quad (E_{2t}-E_{1t})\,\Delta\ell=0 ,
\]
para cualquier dirección tangente $\tu$, lo que equivale a \eqref{eq:frontera-E}.

(iii) Sean $P_1$ y $P_2$ dos puntos a uno y otro lado de $\Sigma$, separados una distancia $h$ según la normal. Entonces
$\phi(P_2)-\phi(P_1)=-\int_{P_1}^{P_2}\vE\cdot\dl$, que en módulo no supera $h\,\sup\abs{\vE}$ y tiende a cero con $h$ porque $\vE$ es acotado.
\end{proof}

\begin{observacion}[Continuidad de $\phi$ y de $\vE_t$ son equivalentes]\label{obs:phi-Et}
Si $\phi$ es continuo, $\phi_1=\phi_2$ en todos los puntos de $\Sigma$: ambos potenciales coinciden como funciones sobre la
superficie, y también sus derivadas a lo largo de ella, $E_{1t}=-\partial\phi_1/\partial t=-\partial\phi_2/\partial t=E_{2t}$. Recíprocamente, si
$\vE_t$ es continuo, para dos puntos $A$ y $B$ de la frontera, recorriéndola justo por encima y justo por debajo,
$\phi_2(B)-\phi_2(A)=-\int_A^BE_{2t}\,\dd\ell=-\int_A^BE_{1t}\,\dd\ell=\phi_1(B)-\phi_1(A)$, de modo que $\phi_2-\phi_1$ es constante a lo largo
de $\Sigma$; el argumento de (iii) muestra que esa constante es nula. Las dos condiciones son la misma
información; la (iii) es la que se usa al resolver la ecuación de Laplace.
\end{observacion}

\subsection{Las otras dos componentes}\label{sec:otras-componentes}

Las condiciones \eqref{eq:frontera-D} y \eqref{eq:frontera-E} no dicen nada directamente de $E_n$ ni de $\vD_t$. Se obtienen
de ellas despejando $\vE=(\vD-\vP)/\eo$.

\begin{corolario}[Saltos de $E_n$ y de $\vD_t$]\label{cor:otras}
En las mismas condiciones del \cref{teo:frontera}:
\begin{enumerate}[start=4]
  \item la componente normal de $\vE$ salta en la carga total, libre y ligada, dividida por $\eo$:
        \begin{equation}\label{eq:salto-En}
          E_{2n}-E_{1n}=\frac{1}{\eo}\big(\sigmaf+\sigmaP\big),
          \qquad \sigmaP=P_{1n}-P_{2n} ;
        \end{equation}
  \item la componente tangencial de $\vD$ salta como la de $\vP$:
        \begin{equation}\label{eq:salto-Dt}
          \vD_{2t}-\vD_{1t}=\vP_{2t}-\vP_{1t} .
        \end{equation}
\end{enumerate}
\end{corolario}

\begin{proof}
(iv) Sustituyendo $D_n=\eo E_n+P_n$ en \eqref{eq:frontera-D},
$(\eo E_{2n}+P_{2n})-(\eo E_{1n}+P_{1n})=\sigmaf$, es decir, $\eo(E_{2n}-E_{1n})=\sigmaf+P_{1n}-P_{2n}$. Que $P_{1n}-P_{2n}$ es la carga ligada neta de la
interfase se sigue del \cref{teo:densidades}: cada medio aporta $\vP\cdot\nn_{\mathrm{ext}}$ en su propia cara, y la normal exterior del
medio~1 es $\nn$ mientras que la del medio~2 es $-\nn$ (\cref{fig:frontera}\,(c)).
(v) Tomando componentes tangenciales en $\vD=\eo\vE+\vP$ y usando \eqref{eq:frontera-E},
$\vD_{2t}-\vD_{1t}=\eo(\vE_{2t}-\vE_{1t})+(\vP_{2t}-\vP_{1t})=\vP_{2t}-\vP_{1t}$.
\end{proof}

\begin{observacion}[Lectura de los saltos]\label{obs:lectura-saltos}
La ecuación \eqref{eq:salto-En} es la condición del vacío con la carga total: el campo $\vE$ «ve» todas las cargas. La
\eqref{eq:salto-Dt} es la versión superficial de $\rot\vD=\rot\vP$: como $\vD$ no es conservativo, su componente tangencial no tiene
por qué conservarse. Dos casos particulares: en la cara de un dieléctrico en el vacío ($\vP_2=\vcero$) se recupera
$\sigmaP=\vP\cdot\nn$; y junto a un conductor en equilibrio ($\vE_1=\vD_1=\vcero$ en el metal), $D_{2n}=\sigmaf$ es la carga libre del conductor y
$\eo E_{2n}=\sigmaf+\sigmaP$ incluye la carga ligada del dieléctrico adyacente.
\end{observacion}

\subsection{Medios lineales: ley de refracción}\label{sec:refraccion}

Si ambos medios son \LHI{}, $\vD_1=\varepsilon_1\vE_1$ y $\vD_2=\varepsilon_2\vE_2$, y las condiciones fundamentales toman la forma
\begin{equation}\label{eq:frontera-LHI}
  \varepsilon_2E_{2n}-\varepsilon_1E_{1n}=\sigmaf,\qquad
  \frac{\vD_{2t}}{\varepsilon_2}-\frac{\vD_{1t}}{\varepsilon_1}=\vcero .
\end{equation}
Si $\varepsilon_1\neq\varepsilon_2$, ni $E_n$ ni $\vD_t$ son continuos, aunque no haya carga libre en la interfase. Esta discontinuidad del
campo se traduce en una refracción de sus líneas.

\begin{proposicion}[Ley de refracción de las líneas de campo]\label{prop:refraccion}
En una interfase sin carga libre entre dos medios \LHI{}, si $\theta_1$ y $\theta_2$ son los ángulos que forman $\vE_1$ y $\vE_2$ con la
normal,
\begin{equation}\label{eq:refraccion}
  \boxed{\;\frac{\tan\theta_2}{\varepsilon_2}=\frac{\tan\theta_1}{\varepsilon_1}\; .}
\end{equation}
Las líneas de $\vD$ se refractan con los mismos ángulos, pues $\vD\parallel\vE$ en cada medio.
\end{proposicion}

\begin{proof}
Con $\sigmaf=0$, $E_{2t}=E_{1t}$ y $\varepsilon_2E_{2n}=\varepsilon_1E_{1n}$. Como $\tan\theta_i=E_{it}/E_{in}$ (\cref{fig:refraccion}\,(a)),
\[
  \frac{\tan\theta_2}{\tan\theta_1}=\frac{E_{2t}}{E_{1t}}\,\frac{E_{1n}}{E_{2n}}=1\cdot\frac{\varepsilon_2}{\varepsilon_1}. \qedhere
\]
\end{proof}

Al pasar a un medio de mayor permitividad, $\theta_2>\theta_1$: la línea se aleja de la normal y tiende a «tumbarse» a lo largo
de la interfase (\cref{fig:refraccion}\,(b)). En el límite $\varepsilon_2\to\infty$, las líneas del medio~1 llegan perpendiculares a la
interfase, como a la superficie de un conductor~\cite{T3}.

% Figura: ley de refracción de las líneas de campo y ángulo refractado.
\begin{figure}[t]
\centering
\begin{subfigure}[b]{0.46\textwidth}\centering
\begin{tikzpicture}[scale=1.0]
  % medio 1 abajo (eps1), medio 2 arriba (eps2 > eps1); tan t2/eps2 = tan t1/eps1
  \fill[dielec] (-2.4,-1.9) rectangle (2.4,0);
  \fill[dielecB] (-2.4,0) rectangle (2.4,1.9);
  \draw[texto, line width=0.6pt] (-2.4,0) -- (2.4,0);
  \node[etiq, anchor=west] at (-2.35,1.6) {$\varepsilon_2$};
  \node[etiq, anchor=west] at (-2.35,-1.6) {$\varepsilon_1$};
  \draw[texto!50!fondo, dashed, line width=0.35pt] (0,-1.9) -- (0,1.9);
  % E1: ángulo 30 grados con la normal; E2: atan(2.5 tan 30) = 55.3 grados
  \coordinate (E1) at ({-1.55*sin(30)},{-1.55*cos(30)});
  \coordinate (E2) at ({1.75*sin(55.3)},{1.75*cos(55.3)});
  \draw[vecE] (E1) -- (0,0);
  \draw[vecE] (0,0) -- (E2);
  \node[etiq, anchor=north east] at ($(E1)+(0.02,-0.02)$) {$\vE_1$};
  \node[etiq, anchor=west] at ($(E2)+(0.05,0.05)$) {$\vE_2$};
  % componentes
  \draw[campoE!60!fondo, dashed, line width=0.4pt] (E2) -- ({1.75*sin(55.3)},0)
     node[pos=0.5, right, etiq] {$E_{2n}$};
  \draw[campoE!60!fondo, dashed, line width=0.4pt] (E2) -- (0,{1.75*cos(55.3)})
     node[pos=0.5, above, etiq] {$E_{2t}$};
  \draw[campoE!60!fondo, dashed, line width=0.4pt] (E1) -- (0,{-1.55*cos(30)})
     node[pos=0.5, below, etiq] {$E_{1t}$};
  \draw[campoE!60!fondo, dashed, line width=0.4pt] (E1) -- ({-1.55*sin(30)},0)
     node[pos=0.72, left, etiq] {$E_{1n}$};
  % ángulos
  \draw[tenue, line width=0.4pt] (0,0.55) arc (90:34.7:0.55);
  \node[etiqt] at (62:0.78) {$\theta_2$};
  \draw[tenue, line width=0.4pt] (0,-0.55) arc (-90:-120:0.55);
  \node[etiqt] at (-105:0.82) {$\theta_1$};
  \draw[vecn] (-1.75,0) -- (-1.75,0.7) node[right, etiq] {$\nn$};
\end{tikzpicture}
\caption{Refracción de $\vE$ ($\varepsilon_2>\varepsilon_1$, $\sigmaf=0$)}
\end{subfigure}\hfill
\begin{subfigure}[b]{0.5\textwidth}\centering
\begin{tikzpicture}
\begin{axis}[width=6.9cm, height=5.6cm, axis lines=left, xmin=0, xmax=90, ymin=0, ymax=92,
  xlabel={$\theta_1$ (grados)}, ylabel={$\theta_2$ (grados)}, xtick={0,30,60,90}, ytick={0,30,60,90},
  grid=major, clip=false]
  \foreach \k/\c in {0.25/campoE!45!fondo, 0.5/campoE!70!fondo, 2/campoD!70!fondo, 4/campoD} {
    \edef\temp{\noexpand\addplot[draw=\c, line width=0.9pt, domain=0:89.5, samples=120]
      {atan(\k*tan(x))};}\temp
  }
  \addplot[draw=tenue, line width=0.6pt, dashed, domain=0:90] {x};
  \node[etiqf, font=\scriptsize, anchor=west] at (axis cs:5,82) {$\varepsilon_2/\varepsilon_1=4$};
  \node[etiqf, font=\scriptsize] at (axis cs:38,57.4) {$2$};
  \node[etiqf, font=\scriptsize] at (axis cs:50,50) {$1$};
  \node[etiqf, font=\scriptsize] at (axis cs:62,43.2) {$\tfrac12$};
  \node[etiqf, font=\scriptsize] at (axis cs:76,44.7) {$\tfrac14$};
\end{axis}
\end{tikzpicture}
\caption{$\theta_2=\arctan\!\big[(\varepsilon_2/\varepsilon_1)\tan\theta_1\big]$}
\end{subfigure}
\caption{Ley de refracción de las líneas de campo en una interfase sin carga libre. (a) La
componente tangencial de $\vE$ se conserva y la normal se reduce en el factor $\varepsilon_1/\varepsilon_2$: al
pasar al medio de mayor permitividad la línea se aleja de la normal. (b) Ángulo refractado para
varios cocientes de permitividades; la diagonal ($\varepsilon_1=\varepsilon_2$) no tiene refracción.}
\label{fig:refraccion}
\end{figure}


\begin{ejemplo}[Esfera dieléctrica en un campo uniforme]\label{ej:esfera-dielectrica}
Una esfera de radio $a$ y permitividad $\varepsilon=\er\eo$ se coloca en un campo uniforme $\vE_0=E_0\,\uv z$. Busquemos una solución en
la que la esfera se polariza \emph{uniformemente}. Por el \cref{prob:esfera-polarizada}, una esfera con polarización uniforme
$\vP$ crea en su interior el campo uniforme $-\vP/3\eo$, y fuera, el de un dipolo $\vp=\tfrac43\pi a^3\vP$. El campo interior total es
$\vE_{\mathrm{int}}=\vE_0-\vP/3\eo$, y la relación constitutiva exige $\vP=\eo\chi\vE_{\mathrm{int}}$. Resolviendo este sistema autoconsistente,
\[
  \vE_{\mathrm{int}}=\vE_0-\frac{\chi}{3}\vE_{\mathrm{int}}
  \quad\Longrightarrow\quad
  \vE_{\mathrm{int}}=\frac{3}{\er+2}\,\vE_0 ,\qquad
  \vP=3\eo\,\frac{\er-1}{\er+2}\,\vE_0 .
\]
En coordenadas esféricas, el potencial correspondiente es
\[
  \phi_{\mathrm{int}}=-\frac{3}{\er+2}\,E_0\,r\cos\theta,
  \qquad
  \phi_{\mathrm{ext}}=-E_0\,r\cos\theta+\frac{\er-1}{\er+2}\,\frac{E_0a^3\cos\theta}{r^2}.
\]
\emph{Comprobación.} Ambos son armónicos ($r\cos\theta=z$ y $\cos\theta/r^2$ es el potencial de un dipolo); $\phi_{\mathrm{ext}}\to-E_0z$ en el
infinito; en $r=a$ los dos valen $-3E_0a\cos\theta/(\er+2)$, y las componentes radiales de $\vD$ coinciden:
\[
  -\varepsilon\,\frac{\partial\phi_{\mathrm{int}}}{\partial r}=\frac{3\er\eo E_0\cos\theta}{\er+2}
  =-\eo\,\frac{\partial\phi_{\mathrm{ext}}}{\partial r}\bigg|_{r=a}=\eo E_0\cos\theta\Big(1+2\,\frac{\er-1}{\er+2}\Big).
\]
Por el teorema de unicidad~\cite{T3,Jackson}, que se extiende sin cambios a medios \LHI{} a trozos, esta es la solución.
El campo interior es uniforme y más débil que el aplicado; las líneas de $\vD$, continuas, se concentran en la esfera
(\cref{fig:esfera-dielectrica}). En el límite $\er\to\infty$, $\vE_{\mathrm{int}}\to\vcero$ y el dipolo inducido tiende a $4\pi\eo a^3\vE_0$: se
recupera la esfera conductora del Tema~3~\cite{T3}.
\end{ejemplo}

% Figura (simulación): esfera dieléctrica eps_r = 4 en un campo uniforme.
% Datos: scripts/lineas_campo.py -> datos/esfera-dielectrica-{E,D}.tex
\begin{figure}[t]
\centering
\pgfplotsset{esfera/.style={marco, width=7.6cm, xmin=-3, xmax=3, ymin=-2.2, ymax=2.2}}
\begin{subfigure}[b]{0.49\textwidth}\centering
\begin{tikzpicture}
\begin{axis}[esfera]
  \fill[dielec] (axis cs:0,0) circle[radius=1];
  %% Generado por scripts/lineas_campo.py -- no editar a mano
\addplot[equipot] coordinates {(2.781,-2.200) (2.782,-2.165) (2.782,-2.130) (2.783,-2.094) (2.784,-2.059) (2.784,-2.024) (2.785,-1.989) (2.785,-1.954) (2.786,-1.918) (2.787,-1.883) (2.787,-1.848) (2.788,-1.813) (2.789,-1.778) (2.789,-1.742) (2.790,-1.707) (2.791,-1.672) (2.791,-1.637) (2.792,-1.602) (2.793,-1.566) (2.793,-1.531) (2.794,-1.496) (2.795,-1.461) (2.795,-1.426) (2.796,-1.390) (2.797,-1.355) (2.797,-1.320) (2.798,-1.285) (2.799,-1.250) (2.799,-1.214) (2.800,-1.179) (2.801,-1.144) (2.801,-1.109) (2.802,-1.074) (2.802,-1.038) (2.803,-1.003) (2.804,-0.968) (2.804,-0.933) (2.805,-0.898) (2.805,-0.862) (2.806,-0.827) (2.807,-0.792) (2.807,-0.757) (2.808,-0.722) (2.808,-0.686) (2.809,-0.651) (2.809,-0.616) (2.809,-0.581) (2.810,-0.546) (2.810,-0.510) (2.811,-0.475) (2.811,-0.440) (2.811,-0.405) (2.812,-0.370) (2.812,-0.334) (2.812,-0.299) (2.812,-0.264) (2.813,-0.229) (2.813,-0.194) (2.813,-0.158) (2.813,-0.123) (2.813,-0.088) (2.813,-0.053) (2.813,-0.018) (2.813,0.018) (2.813,0.053) (2.813,0.088) (2.813,0.123) (2.813,0.158) (2.813,0.194) (2.813,0.229) (2.812,0.264) (2.812,0.299) (2.812,0.334) (2.812,0.370) (2.811,0.405) (2.811,0.440) (2.811,0.475) (2.810,0.510) (2.810,0.546) (2.809,0.581) (2.809,0.616) (2.809,0.651) (2.808,0.686) (2.808,0.722) (2.807,0.757) (2.807,0.792) (2.806,0.827) (2.805,0.862) (2.805,0.898) (2.804,0.933) (2.804,0.968) (2.803,1.003) (2.802,1.038) (2.802,1.074) (2.801,1.109) (2.801,1.144) (2.800,1.179) (2.799,1.214) (2.799,1.250) (2.798,1.285) (2.797,1.320) (2.797,1.355) (2.796,1.390) (2.795,1.426) (2.795,1.461) (2.794,1.496) (2.793,1.531) (2.793,1.566) (2.792,1.602) (2.791,1.637) (2.791,1.672) (2.790,1.707) (2.789,1.742) (2.789,1.778) (2.788,1.813) (2.787,1.848) (2.787,1.883) (2.786,1.918) (2.785,1.954) (2.785,1.989) (2.784,2.024) (2.784,2.059) (2.783,2.094) (2.782,2.130) (2.782,2.165) (2.781,2.200)};
\addplot[equipot] coordinates {(2.286,-2.200) (2.287,-2.165) (2.287,-2.130) (2.288,-2.094) (2.289,-2.059) (2.290,-2.024) (2.291,-1.989) (2.292,-1.954) (2.293,-1.918) (2.294,-1.883) (2.295,-1.848) (2.296,-1.813) (2.297,-1.778) (2.298,-1.742) (2.299,-1.707) (2.300,-1.672) (2.301,-1.637) (2.302,-1.602) (2.303,-1.566) (2.304,-1.531) (2.306,-1.496) (2.307,-1.461) (2.308,-1.426) (2.309,-1.390) (2.310,-1.355) (2.311,-1.320) (2.312,-1.285) (2.314,-1.250) (2.315,-1.214) (2.316,-1.179) (2.317,-1.144) (2.318,-1.109) (2.319,-1.074) (2.321,-1.039) (2.322,-1.003) (2.323,-0.968) (2.324,-0.933) (2.325,-0.898) (2.326,-0.863) (2.327,-0.827) (2.328,-0.792) (2.329,-0.757) (2.330,-0.722) (2.331,-0.687) (2.332,-0.651) (2.333,-0.616) (2.334,-0.581) (2.335,-0.546) (2.335,-0.511) (2.336,-0.475) (2.337,-0.440) (2.338,-0.405) (2.338,-0.370) (2.339,-0.335) (2.339,-0.299) (2.340,-0.264) (2.340,-0.229) (2.340,-0.194) (2.341,-0.158) (2.341,-0.123) (2.341,-0.088) (2.341,-0.053) (2.341,-0.018) (2.341,0.018) (2.341,0.053) (2.341,0.088) (2.341,0.123) (2.341,0.158) (2.340,0.194) (2.340,0.229) (2.340,0.264) (2.339,0.299) (2.339,0.335) (2.338,0.370) (2.338,0.405) (2.337,0.440) (2.336,0.475) (2.335,0.511) (2.335,0.546) (2.334,0.581) (2.333,0.616) (2.332,0.651) (2.331,0.687) (2.330,0.722) (2.329,0.757) (2.328,0.792) (2.327,0.827) (2.326,0.863) (2.325,0.898) (2.324,0.933) (2.323,0.968) (2.322,1.003) (2.321,1.039) (2.319,1.074) (2.318,1.109) (2.317,1.144) (2.316,1.179) (2.315,1.214) (2.314,1.250) (2.312,1.285) (2.311,1.320) (2.310,1.355) (2.309,1.390) (2.308,1.426) (2.307,1.461) (2.306,1.496) (2.304,1.531) (2.303,1.566) (2.302,1.602) (2.301,1.637) (2.300,1.672) (2.299,1.707) (2.298,1.742) (2.297,1.778) (2.296,1.813) (2.295,1.848) (2.294,1.883) (2.293,1.918) (2.292,1.954) (2.291,1.989) (2.290,2.024) (2.289,2.059) (2.288,2.094) (2.287,2.130) (2.287,2.165) (2.286,2.200)};
\addplot[equipot] coordinates {(1.789,-2.200) (1.790,-2.165) (1.792,-2.130) (1.793,-2.094) (1.794,-2.059) (1.795,-2.024) (1.797,-1.989) (1.798,-1.953) (1.799,-1.918) (1.801,-1.883) (1.802,-1.848) (1.804,-1.813) (1.806,-1.777) (1.807,-1.742) (1.809,-1.707) (1.810,-1.672) (1.812,-1.637) (1.814,-1.601) (1.816,-1.566) (1.818,-1.531) (1.820,-1.496) (1.822,-1.461) (1.824,-1.425) (1.826,-1.390) (1.828,-1.355) (1.830,-1.320) (1.832,-1.285) (1.834,-1.250) (1.836,-1.214) (1.838,-1.179) (1.840,-1.144) (1.843,-1.109) (1.845,-1.074) (1.847,-1.039) (1.849,-1.003) (1.851,-0.968) (1.854,-0.933) (1.856,-0.898) (1.858,-0.863) (1.860,-0.827) (1.862,-0.792) (1.864,-0.757) (1.866,-0.722) (1.868,-0.687) (1.870,-0.652) (1.872,-0.616) (1.874,-0.581) (1.876,-0.546) (1.877,-0.511) (1.879,-0.476) (1.881,-0.440) (1.882,-0.405) (1.883,-0.370) (1.884,-0.335) (1.885,-0.300) (1.886,-0.264) (1.887,-0.229) (1.888,-0.194) (1.889,-0.159) (1.889,-0.123) (1.890,-0.088) (1.890,-0.053) (1.890,-0.018) (1.890,0.018) (1.890,0.053) (1.890,0.088) (1.889,0.123) (1.889,0.159) (1.888,0.194) (1.887,0.229) (1.886,0.264) (1.885,0.300) (1.884,0.335) (1.883,0.370) (1.882,0.405) (1.881,0.440) (1.879,0.476) (1.877,0.511) (1.876,0.546) (1.874,0.581) (1.872,0.616) (1.870,0.652) (1.868,0.687) (1.866,0.722) (1.864,0.757) (1.862,0.792) (1.860,0.827) (1.858,0.863) (1.856,0.898) (1.854,0.933) (1.851,0.968) (1.849,1.003) (1.847,1.039) (1.845,1.074) (1.843,1.109) (1.840,1.144) (1.838,1.179) (1.836,1.214) (1.834,1.250) (1.832,1.285) (1.830,1.320) (1.828,1.355) (1.826,1.390) (1.824,1.425) (1.822,1.461) (1.820,1.496) (1.818,1.531) (1.816,1.566) (1.814,1.601) (1.812,1.637) (1.810,1.672) (1.809,1.707) (1.807,1.742) (1.806,1.777) (1.804,1.813) (1.802,1.848) (1.801,1.883) (1.799,1.918) (1.798,1.953) (1.797,1.989) (1.795,2.024) (1.794,2.059) (1.793,2.094) (1.792,2.130) (1.790,2.165) (1.789,2.200)};
\addplot[equipot] coordinates {(1.289,-2.200) (1.290,-2.165) (1.292,-2.130) (1.293,-2.095) (1.295,-2.060) (1.297,-2.025) (1.298,-1.990) (1.300,-1.955) (1.302,-1.920) (1.304,-1.885) (1.306,-1.850) (1.308,-1.815) (1.311,-1.780) (1.313,-1.745) (1.316,-1.710) (1.318,-1.675) (1.321,-1.640) (1.324,-1.605) (1.326,-1.570) (1.329,-1.535) (1.332,-1.500) (1.336,-1.465) (1.339,-1.430) (1.342,-1.395) (1.346,-1.360) (1.350,-1.326) (1.353,-1.291) (1.357,-1.256) (1.361,-1.221) (1.365,-1.186) (1.370,-1.151) (1.374,-1.116) (1.378,-1.082) (1.383,-1.047) (1.387,-1.012) (1.392,-0.977) (1.396,-0.943) (1.401,-0.908) (1.405,-0.873) (1.410,-0.838) (1.414,-0.803) (1.419,-0.769) (1.423,-0.734) (1.428,-0.699) (1.432,-0.664) (1.436,-0.629) (1.440,-0.595) (1.444,-0.560) (1.448,-0.525) (1.452,-0.490) (1.455,-0.455) (1.459,-0.420) (1.462,-0.385) (1.464,-0.350) (1.467,-0.315) (1.469,-0.280) (1.472,-0.245) (1.473,-0.210) (1.475,-0.175) (1.476,-0.140) (1.477,-0.105) (1.478,-0.070) (1.479,-0.035) (1.479,0.000) (1.479,0.035) (1.478,0.070) (1.477,0.105) (1.476,0.140) (1.475,0.175) (1.473,0.210) (1.472,0.245) (1.469,0.280) (1.467,0.315) (1.464,0.350) (1.462,0.385) (1.459,0.420) (1.455,0.455) (1.452,0.490) (1.448,0.525) (1.444,0.560) (1.440,0.595) (1.436,0.629) (1.432,0.664) (1.428,0.699) (1.423,0.734) (1.419,0.769) (1.414,0.803) (1.410,0.838) (1.405,0.873) (1.401,0.908) (1.396,0.943) (1.392,0.977) (1.387,1.012) (1.383,1.047) (1.378,1.082) (1.374,1.116) (1.370,1.151) (1.365,1.186) (1.361,1.221) (1.357,1.256) (1.353,1.291) (1.350,1.326) (1.346,1.360) (1.342,1.395) (1.339,1.430) (1.336,1.465) (1.332,1.500) (1.329,1.535) (1.326,1.570) (1.324,1.605) (1.321,1.640) (1.318,1.675) (1.316,1.710) (1.313,1.745) (1.311,1.780) (1.308,1.815) (1.306,1.850) (1.304,1.885) (1.302,1.920) (1.300,1.955) (1.298,1.990) (1.297,2.025) (1.295,2.060) (1.293,2.095) (1.292,2.130) (1.290,2.165) (1.289,2.200)};
\addplot[equipot] coordinates {(0.781,-2.200) (0.782,-2.165) (0.784,-2.130) (0.785,-2.094) (0.787,-2.059) (0.788,-2.024) (0.790,-1.989) (0.792,-1.954) (0.794,-1.918) (0.797,-1.883) (0.799,-1.848) (0.801,-1.813) (0.804,-1.778) (0.807,-1.743) (0.810,-1.708) (0.813,-1.673) (0.817,-1.638) (0.820,-1.602) (0.824,-1.567) (0.828,-1.532) (0.833,-1.498) (0.837,-1.463) (0.842,-1.428) (0.848,-1.393) (0.853,-1.358) (0.859,-1.323) (0.866,-1.289) (0.872,-1.254) (0.879,-1.220) (0.887,-1.185) (0.894,-1.151) (0.903,-1.117) (0.911,-1.082) (0.920,-1.048) (0.929,-1.014) (0.938,-0.980) (0.947,-0.946) (0.957,-0.912) (0.967,-0.878) (0.977,-0.845) (0.987,-0.811) (0.997,-0.777) (1.007,-0.743) (1.017,-0.710) (1.027,-0.676) (1.036,-0.642) (1.046,-0.608) (1.055,-0.574) (1.063,-0.540) (1.072,-0.505) (1.080,-0.471) (1.088,-0.437) (1.095,-0.402) (1.102,-0.368) (1.108,-0.333) (1.113,-0.298) (1.119,-0.263) (1.123,-0.228) (1.127,-0.193) (1.130,-0.158) (1.133,-0.123) (1.135,-0.088) (1.136,-0.053) (1.137,-0.018) (1.137,0.018) (1.136,0.053) (1.135,0.088) (1.133,0.123) (1.130,0.158) (1.127,0.193) (1.123,0.228) (1.119,0.263) (1.113,0.298) (1.108,0.333) (1.102,0.368) (1.095,0.402) (1.088,0.437) (1.080,0.471) (1.072,0.505) (1.063,0.540) (1.055,0.574) (1.046,0.608) (1.036,0.642) (1.027,0.676) (1.017,0.710) (1.007,0.743) (0.997,0.777) (0.987,0.811) (0.977,0.845) (0.967,0.878) (0.957,0.912) (0.947,0.946) (0.938,0.980) (0.929,1.014) (0.920,1.048) (0.911,1.082) (0.903,1.117) (0.894,1.151) (0.887,1.185) (0.879,1.220) (0.872,1.254) (0.866,1.289) (0.859,1.323) (0.853,1.358) (0.848,1.393) (0.842,1.428) (0.837,1.463) (0.833,1.498) (0.828,1.532) (0.824,1.567) (0.820,1.602) (0.817,1.638) (0.813,1.673) (0.810,1.708) (0.807,1.743) (0.804,1.778) (0.801,1.813) (0.799,1.848) (0.797,1.883) (0.794,1.918) (0.792,1.954) (0.790,1.989) (0.788,2.024) (0.787,2.059) (0.785,2.094) (0.784,2.130) (0.782,2.165) (0.781,2.200)};
\addplot[equipot] coordinates {(0.262,-2.200) (0.263,-2.165) (0.263,-2.130) (0.264,-2.095) (0.265,-2.060) (0.266,-2.025) (0.266,-1.990) (0.267,-1.955) (0.268,-1.919) (0.270,-1.884) (0.271,-1.849) (0.272,-1.814) (0.273,-1.779) (0.275,-1.744) (0.277,-1.709) (0.278,-1.674) (0.280,-1.639) (0.283,-1.604) (0.285,-1.569) (0.288,-1.534) (0.291,-1.499) (0.294,-1.464) (0.298,-1.429) (0.302,-1.395) (0.307,-1.360) (0.312,-1.325) (0.318,-1.291) (0.324,-1.256) (0.332,-1.222) (0.340,-1.188) (0.350,-1.154) (0.360,-1.121) (0.373,-1.088) (0.386,-1.055) (0.401,-1.024) (0.417,-0.993) (0.435,-0.962) (0.453,-0.933) (0.473,-0.904) (0.493,-0.875) (0.500,-0.842) (0.500,-0.807) (0.500,-0.772) (0.500,-0.737) (0.500,-0.701) (0.500,-0.666) (0.500,-0.631) (0.500,-0.596) (0.500,-0.561) (0.500,-0.526) (0.500,-0.491) (0.500,-0.456) (0.500,-0.421) (0.500,-0.386) (0.500,-0.351) (0.500,-0.316) (0.500,-0.281) (0.500,-0.246) (0.500,-0.210) (0.500,-0.175) (0.500,-0.140) (0.500,-0.105) (0.500,-0.070) (0.500,-0.035) (0.500,0.000) (0.500,0.035) (0.500,0.070) (0.500,0.105) (0.500,0.140) (0.500,0.175) (0.500,0.210) (0.500,0.246) (0.500,0.281) (0.500,0.316) (0.500,0.351) (0.500,0.386) (0.500,0.421) (0.500,0.456) (0.500,0.491) (0.500,0.526) (0.500,0.561) (0.500,0.596) (0.500,0.631) (0.500,0.666) (0.500,0.701) (0.500,0.737) (0.500,0.772) (0.500,0.807) (0.500,0.842) (0.493,0.875) (0.473,0.904) (0.453,0.933) (0.435,0.962) (0.417,0.993) (0.401,1.024) (0.386,1.055) (0.373,1.088) (0.360,1.121) (0.350,1.154) (0.340,1.188) (0.332,1.222) (0.324,1.256) (0.318,1.291) (0.312,1.325) (0.307,1.360) (0.302,1.395) (0.298,1.429) (0.294,1.464) (0.291,1.499) (0.288,1.534) (0.285,1.569) (0.283,1.604) (0.280,1.639) (0.278,1.674) (0.277,1.709) (0.275,1.744) (0.273,1.779) (0.272,1.814) (0.271,1.849) (0.270,1.884) (0.268,1.919) (0.267,1.955) (0.266,1.990) (0.266,2.025) (0.265,2.060) (0.264,2.095) (0.263,2.130) (0.263,2.165) (0.262,2.200)};
\addplot[equipot] coordinates {(-0.262,-2.200) (-0.263,-2.165) (-0.263,-2.130) (-0.264,-2.095) (-0.265,-2.060) (-0.266,-2.025) (-0.266,-1.990) (-0.267,-1.955) (-0.268,-1.919) (-0.270,-1.884) (-0.271,-1.849) (-0.272,-1.814) (-0.273,-1.779) (-0.275,-1.744) (-0.277,-1.709) (-0.278,-1.674) (-0.280,-1.639) (-0.283,-1.604) (-0.285,-1.569) (-0.288,-1.534) (-0.291,-1.499) (-0.294,-1.464) (-0.298,-1.429) (-0.302,-1.395) (-0.307,-1.360) (-0.312,-1.325) (-0.318,-1.291) (-0.324,-1.256) (-0.332,-1.222) (-0.340,-1.188) (-0.350,-1.154) (-0.360,-1.121) (-0.373,-1.088) (-0.386,-1.055) (-0.401,-1.024) (-0.417,-0.993) (-0.435,-0.962) (-0.453,-0.933) (-0.473,-0.904) (-0.493,-0.875) (-0.500,-0.842) (-0.500,-0.807) (-0.500,-0.772) (-0.500,-0.737) (-0.500,-0.701) (-0.500,-0.666) (-0.500,-0.631) (-0.500,-0.596) (-0.500,-0.561) (-0.500,-0.526) (-0.500,-0.491) (-0.500,-0.456) (-0.500,-0.421) (-0.500,-0.386) (-0.500,-0.351) (-0.500,-0.316) (-0.500,-0.281) (-0.500,-0.246) (-0.500,-0.210) (-0.500,-0.175) (-0.500,-0.140) (-0.500,-0.105) (-0.500,-0.070) (-0.500,-0.035) (-0.500,-0.000) (-0.500,0.035) (-0.500,0.070) (-0.500,0.105) (-0.500,0.140) (-0.500,0.175) (-0.500,0.210) (-0.500,0.246) (-0.500,0.281) (-0.500,0.316) (-0.500,0.351) (-0.500,0.386) (-0.500,0.421) (-0.500,0.456) (-0.500,0.491) (-0.500,0.526) (-0.500,0.561) (-0.500,0.596) (-0.500,0.631) (-0.500,0.666) (-0.500,0.701) (-0.500,0.737) (-0.500,0.772) (-0.500,0.807) (-0.500,0.842) (-0.493,0.875) (-0.473,0.904) (-0.453,0.933) (-0.435,0.962) (-0.417,0.993) (-0.401,1.024) (-0.386,1.055) (-0.373,1.088) (-0.360,1.121) (-0.350,1.154) (-0.340,1.188) (-0.332,1.222) (-0.324,1.256) (-0.318,1.291) (-0.312,1.325) (-0.307,1.360) (-0.302,1.395) (-0.298,1.429) (-0.294,1.464) (-0.291,1.499) (-0.288,1.534) (-0.285,1.569) (-0.283,1.604) (-0.280,1.639) (-0.278,1.674) (-0.277,1.709) (-0.275,1.744) (-0.273,1.779) (-0.272,1.814) (-0.271,1.849) (-0.270,1.884) (-0.268,1.919) (-0.267,1.955) (-0.266,1.990) (-0.266,2.025) (-0.265,2.060) (-0.264,2.095) (-0.263,2.130) (-0.263,2.165) (-0.262,2.200)};
\addplot[equipot] coordinates {(-0.781,-2.200) (-0.782,-2.165) (-0.784,-2.130) (-0.785,-2.094) (-0.787,-2.059) (-0.788,-2.024) (-0.790,-1.989) (-0.792,-1.954) (-0.794,-1.918) (-0.797,-1.883) (-0.799,-1.848) (-0.801,-1.813) (-0.804,-1.778) (-0.807,-1.743) (-0.810,-1.708) (-0.813,-1.673) (-0.817,-1.638) (-0.820,-1.602) (-0.824,-1.567) (-0.828,-1.532) (-0.833,-1.498) (-0.837,-1.463) (-0.842,-1.428) (-0.848,-1.393) (-0.853,-1.358) (-0.859,-1.323) (-0.866,-1.289) (-0.872,-1.254) (-0.879,-1.220) (-0.887,-1.185) (-0.894,-1.151) (-0.903,-1.117) (-0.911,-1.082) (-0.920,-1.048) (-0.929,-1.014) (-0.938,-0.980) (-0.947,-0.946) (-0.957,-0.912) (-0.967,-0.878) (-0.977,-0.845) (-0.987,-0.811) (-0.997,-0.777) (-1.007,-0.743) (-1.017,-0.710) (-1.027,-0.676) (-1.036,-0.642) (-1.046,-0.608) (-1.055,-0.574) (-1.063,-0.540) (-1.072,-0.505) (-1.080,-0.471) (-1.088,-0.437) (-1.095,-0.402) (-1.102,-0.368) (-1.108,-0.333) (-1.113,-0.298) (-1.119,-0.263) (-1.123,-0.228) (-1.127,-0.193) (-1.130,-0.158) (-1.133,-0.123) (-1.135,-0.088) (-1.136,-0.053) (-1.137,-0.018) (-1.137,0.018) (-1.136,0.053) (-1.135,0.088) (-1.133,0.123) (-1.130,0.158) (-1.127,0.193) (-1.123,0.228) (-1.119,0.263) (-1.113,0.298) (-1.108,0.333) (-1.102,0.368) (-1.095,0.402) (-1.088,0.437) (-1.080,0.471) (-1.072,0.505) (-1.063,0.540) (-1.055,0.574) (-1.046,0.608) (-1.036,0.642) (-1.027,0.676) (-1.017,0.710) (-1.007,0.743) (-0.997,0.777) (-0.987,0.811) (-0.977,0.845) (-0.967,0.878) (-0.957,0.912) (-0.947,0.946) (-0.938,0.980) (-0.929,1.014) (-0.920,1.048) (-0.911,1.082) (-0.903,1.117) (-0.894,1.151) (-0.887,1.185) (-0.879,1.220) (-0.872,1.254) (-0.866,1.289) (-0.859,1.323) (-0.853,1.358) (-0.848,1.393) (-0.842,1.428) (-0.837,1.463) (-0.833,1.498) (-0.828,1.532) (-0.824,1.567) (-0.820,1.602) (-0.817,1.638) (-0.813,1.673) (-0.810,1.708) (-0.807,1.743) (-0.804,1.778) (-0.801,1.813) (-0.799,1.848) (-0.797,1.883) (-0.794,1.918) (-0.792,1.954) (-0.790,1.989) (-0.788,2.024) (-0.787,2.059) (-0.785,2.094) (-0.784,2.130) (-0.782,2.165) (-0.781,2.200)};
\addplot[equipot] coordinates {(-1.289,-2.200) (-1.290,-2.165) (-1.292,-2.130) (-1.293,-2.095) (-1.295,-2.060) (-1.297,-2.025) (-1.298,-1.990) (-1.300,-1.955) (-1.302,-1.920) (-1.304,-1.885) (-1.306,-1.850) (-1.308,-1.815) (-1.311,-1.780) (-1.313,-1.745) (-1.316,-1.710) (-1.318,-1.675) (-1.321,-1.640) (-1.324,-1.605) (-1.326,-1.570) (-1.329,-1.535) (-1.332,-1.500) (-1.336,-1.465) (-1.339,-1.430) (-1.342,-1.395) (-1.346,-1.360) (-1.350,-1.326) (-1.353,-1.291) (-1.357,-1.256) (-1.361,-1.221) (-1.365,-1.186) (-1.370,-1.151) (-1.374,-1.116) (-1.378,-1.082) (-1.383,-1.047) (-1.387,-1.012) (-1.392,-0.977) (-1.396,-0.943) (-1.401,-0.908) (-1.405,-0.873) (-1.410,-0.838) (-1.414,-0.803) (-1.419,-0.769) (-1.423,-0.734) (-1.428,-0.699) (-1.432,-0.664) (-1.436,-0.629) (-1.440,-0.595) (-1.444,-0.560) (-1.448,-0.525) (-1.452,-0.490) (-1.455,-0.455) (-1.459,-0.420) (-1.462,-0.385) (-1.464,-0.350) (-1.467,-0.315) (-1.469,-0.280) (-1.472,-0.245) (-1.473,-0.210) (-1.475,-0.175) (-1.476,-0.140) (-1.477,-0.105) (-1.478,-0.070) (-1.479,-0.035) (-1.479,0.000) (-1.479,0.035) (-1.478,0.070) (-1.477,0.105) (-1.476,0.140) (-1.475,0.175) (-1.473,0.210) (-1.472,0.245) (-1.469,0.280) (-1.467,0.315) (-1.464,0.350) (-1.462,0.385) (-1.459,0.420) (-1.455,0.455) (-1.452,0.490) (-1.448,0.525) (-1.444,0.560) (-1.440,0.595) (-1.436,0.629) (-1.432,0.664) (-1.428,0.699) (-1.423,0.734) (-1.419,0.769) (-1.414,0.803) (-1.410,0.838) (-1.405,0.873) (-1.401,0.908) (-1.396,0.943) (-1.392,0.977) (-1.387,1.012) (-1.383,1.047) (-1.378,1.082) (-1.374,1.116) (-1.370,1.151) (-1.365,1.186) (-1.361,1.221) (-1.357,1.256) (-1.353,1.291) (-1.350,1.326) (-1.346,1.360) (-1.342,1.395) (-1.339,1.430) (-1.336,1.465) (-1.332,1.500) (-1.329,1.535) (-1.326,1.570) (-1.324,1.605) (-1.321,1.640) (-1.318,1.675) (-1.316,1.710) (-1.313,1.745) (-1.311,1.780) (-1.308,1.815) (-1.306,1.850) (-1.304,1.885) (-1.302,1.920) (-1.300,1.955) (-1.298,1.990) (-1.297,2.025) (-1.295,2.060) (-1.293,2.095) (-1.292,2.130) (-1.290,2.165) (-1.289,2.200)};
\addplot[equipot] coordinates {(-1.789,-2.200) (-1.790,-2.165) (-1.792,-2.130) (-1.793,-2.094) (-1.794,-2.059) (-1.795,-2.024) (-1.797,-1.989) (-1.798,-1.953) (-1.799,-1.918) (-1.801,-1.883) (-1.802,-1.848) (-1.804,-1.813) (-1.806,-1.777) (-1.807,-1.742) (-1.809,-1.707) (-1.810,-1.672) (-1.812,-1.637) (-1.814,-1.601) (-1.816,-1.566) (-1.818,-1.531) (-1.820,-1.496) (-1.822,-1.461) (-1.824,-1.425) (-1.826,-1.390) (-1.828,-1.355) (-1.830,-1.320) (-1.832,-1.285) (-1.834,-1.250) (-1.836,-1.214) (-1.838,-1.179) (-1.840,-1.144) (-1.843,-1.109) (-1.845,-1.074) (-1.847,-1.039) (-1.849,-1.003) (-1.851,-0.968) (-1.854,-0.933) (-1.856,-0.898) (-1.858,-0.863) (-1.860,-0.827) (-1.862,-0.792) (-1.864,-0.757) (-1.866,-0.722) (-1.868,-0.687) (-1.870,-0.652) (-1.872,-0.616) (-1.874,-0.581) (-1.876,-0.546) (-1.877,-0.511) (-1.879,-0.476) (-1.881,-0.440) (-1.882,-0.405) (-1.883,-0.370) (-1.884,-0.335) (-1.885,-0.300) (-1.886,-0.264) (-1.887,-0.229) (-1.888,-0.194) (-1.889,-0.159) (-1.889,-0.123) (-1.890,-0.088) (-1.890,-0.053) (-1.890,-0.018) (-1.890,0.018) (-1.890,0.053) (-1.890,0.088) (-1.889,0.123) (-1.889,0.159) (-1.888,0.194) (-1.887,0.229) (-1.886,0.264) (-1.885,0.300) (-1.884,0.335) (-1.883,0.370) (-1.882,0.405) (-1.881,0.440) (-1.879,0.476) (-1.877,0.511) (-1.876,0.546) (-1.874,0.581) (-1.872,0.616) (-1.870,0.652) (-1.868,0.687) (-1.866,0.722) (-1.864,0.757) (-1.862,0.792) (-1.860,0.827) (-1.858,0.863) (-1.856,0.898) (-1.854,0.933) (-1.851,0.968) (-1.849,1.003) (-1.847,1.039) (-1.845,1.074) (-1.843,1.109) (-1.840,1.144) (-1.838,1.179) (-1.836,1.214) (-1.834,1.250) (-1.832,1.285) (-1.830,1.320) (-1.828,1.355) (-1.826,1.390) (-1.824,1.425) (-1.822,1.461) (-1.820,1.496) (-1.818,1.531) (-1.816,1.566) (-1.814,1.601) (-1.812,1.637) (-1.810,1.672) (-1.809,1.707) (-1.807,1.742) (-1.806,1.777) (-1.804,1.813) (-1.802,1.848) (-1.801,1.883) (-1.799,1.918) (-1.798,1.953) (-1.797,1.989) (-1.795,2.024) (-1.794,2.059) (-1.793,2.094) (-1.792,2.130) (-1.790,2.165) (-1.789,2.200)};
\addplot[equipot] coordinates {(-2.286,-2.200) (-2.287,-2.165) (-2.287,-2.130) (-2.288,-2.094) (-2.289,-2.059) (-2.290,-2.024) (-2.291,-1.989) (-2.292,-1.954) (-2.293,-1.918) (-2.294,-1.883) (-2.295,-1.848) (-2.296,-1.813) (-2.297,-1.778) (-2.298,-1.742) (-2.299,-1.707) (-2.300,-1.672) (-2.301,-1.637) (-2.302,-1.602) (-2.303,-1.566) (-2.304,-1.531) (-2.306,-1.496) (-2.307,-1.461) (-2.308,-1.426) (-2.309,-1.390) (-2.310,-1.355) (-2.311,-1.320) (-2.312,-1.285) (-2.314,-1.250) (-2.315,-1.214) (-2.316,-1.179) (-2.317,-1.144) (-2.318,-1.109) (-2.319,-1.074) (-2.321,-1.039) (-2.322,-1.003) (-2.323,-0.968) (-2.324,-0.933) (-2.325,-0.898) (-2.326,-0.863) (-2.327,-0.827) (-2.328,-0.792) (-2.329,-0.757) (-2.330,-0.722) (-2.331,-0.687) (-2.332,-0.651) (-2.333,-0.616) (-2.334,-0.581) (-2.335,-0.546) (-2.335,-0.511) (-2.336,-0.475) (-2.337,-0.440) (-2.338,-0.405) (-2.338,-0.370) (-2.339,-0.335) (-2.339,-0.299) (-2.340,-0.264) (-2.340,-0.229) (-2.340,-0.194) (-2.341,-0.158) (-2.341,-0.123) (-2.341,-0.088) (-2.341,-0.053) (-2.341,-0.018) (-2.341,0.018) (-2.341,0.053) (-2.341,0.088) (-2.341,0.123) (-2.341,0.158) (-2.340,0.194) (-2.340,0.229) (-2.340,0.264) (-2.339,0.299) (-2.339,0.335) (-2.338,0.370) (-2.338,0.405) (-2.337,0.440) (-2.336,0.475) (-2.335,0.511) (-2.335,0.546) (-2.334,0.581) (-2.333,0.616) (-2.332,0.651) (-2.331,0.687) (-2.330,0.722) (-2.329,0.757) (-2.328,0.792) (-2.327,0.827) (-2.326,0.863) (-2.325,0.898) (-2.324,0.933) (-2.323,0.968) (-2.322,1.003) (-2.321,1.039) (-2.319,1.074) (-2.318,1.109) (-2.317,1.144) (-2.316,1.179) (-2.315,1.214) (-2.314,1.250) (-2.312,1.285) (-2.311,1.320) (-2.310,1.355) (-2.309,1.390) (-2.308,1.426) (-2.307,1.461) (-2.306,1.496) (-2.304,1.531) (-2.303,1.566) (-2.302,1.602) (-2.301,1.637) (-2.300,1.672) (-2.299,1.707) (-2.298,1.742) (-2.297,1.778) (-2.296,1.813) (-2.295,1.848) (-2.294,1.883) (-2.293,1.918) (-2.292,1.954) (-2.291,1.989) (-2.290,2.024) (-2.289,2.059) (-2.288,2.094) (-2.287,2.130) (-2.287,2.165) (-2.286,2.200)};
\addplot[equipot] coordinates {(-2.781,-2.200) (-2.782,-2.165) (-2.782,-2.130) (-2.783,-2.094) (-2.784,-2.059) (-2.784,-2.024) (-2.785,-1.989) (-2.785,-1.954) (-2.786,-1.918) (-2.787,-1.883) (-2.787,-1.848) (-2.788,-1.813) (-2.789,-1.778) (-2.789,-1.742) (-2.790,-1.707) (-2.791,-1.672) (-2.791,-1.637) (-2.792,-1.602) (-2.793,-1.566) (-2.793,-1.531) (-2.794,-1.496) (-2.795,-1.461) (-2.795,-1.426) (-2.796,-1.390) (-2.797,-1.355) (-2.797,-1.320) (-2.798,-1.285) (-2.799,-1.250) (-2.799,-1.214) (-2.800,-1.179) (-2.801,-1.144) (-2.801,-1.109) (-2.802,-1.074) (-2.802,-1.038) (-2.803,-1.003) (-2.804,-0.968) (-2.804,-0.933) (-2.805,-0.898) (-2.805,-0.862) (-2.806,-0.827) (-2.807,-0.792) (-2.807,-0.757) (-2.808,-0.722) (-2.808,-0.686) (-2.809,-0.651) (-2.809,-0.616) (-2.809,-0.581) (-2.810,-0.546) (-2.810,-0.510) (-2.811,-0.475) (-2.811,-0.440) (-2.811,-0.405) (-2.812,-0.370) (-2.812,-0.334) (-2.812,-0.299) (-2.812,-0.264) (-2.813,-0.229) (-2.813,-0.194) (-2.813,-0.158) (-2.813,-0.123) (-2.813,-0.088) (-2.813,-0.053) (-2.813,-0.018) (-2.813,0.018) (-2.813,0.053) (-2.813,0.088) (-2.813,0.123) (-2.813,0.158) (-2.813,0.194) (-2.813,0.229) (-2.812,0.264) (-2.812,0.299) (-2.812,0.334) (-2.812,0.370) (-2.811,0.405) (-2.811,0.440) (-2.811,0.475) (-2.810,0.510) (-2.810,0.546) (-2.809,0.581) (-2.809,0.616) (-2.809,0.651) (-2.808,0.686) (-2.808,0.722) (-2.807,0.757) (-2.807,0.792) (-2.806,0.827) (-2.805,0.862) (-2.805,0.898) (-2.804,0.933) (-2.804,0.968) (-2.803,1.003) (-2.802,1.038) (-2.802,1.074) (-2.801,1.109) (-2.801,1.144) (-2.800,1.179) (-2.799,1.214) (-2.799,1.250) (-2.798,1.285) (-2.797,1.320) (-2.797,1.355) (-2.796,1.390) (-2.795,1.426) (-2.795,1.461) (-2.794,1.496) (-2.793,1.531) (-2.793,1.566) (-2.792,1.602) (-2.791,1.637) (-2.791,1.672) (-2.790,1.707) (-2.789,1.742) (-2.789,1.778) (-2.788,1.813) (-2.787,1.848) (-2.787,1.883) (-2.786,1.918) (-2.785,1.954) (-2.785,1.989) (-2.784,2.024) (-2.784,2.059) (-2.783,2.094) (-2.782,2.130) (-2.782,2.165) (-2.781,2.200)};
\addplot[lineaE, flechaen=0.577] coordinates {(0.990,0.141) (1.025,0.145) (1.060,0.148) (1.095,0.152) (1.130,0.155) (1.166,0.157) (1.201,0.160) (1.236,0.162) (1.271,0.165) (1.307,0.167) (1.342,0.169) (1.377,0.171) (1.412,0.172) (1.447,0.174) (1.483,0.175) (1.518,0.177) (1.553,0.178) (1.589,0.179) (1.624,0.180) (1.659,0.181) (1.694,0.182) (1.730,0.183) (1.765,0.184) (1.800,0.185) (1.835,0.186) (1.871,0.186) (1.906,0.187) (1.941,0.188) (1.977,0.188) (2.012,0.189) (2.047,0.189) (2.082,0.190) (2.118,0.190) (2.153,0.191) (2.188,0.191) (2.224,0.192) (2.259,0.192) (2.294,0.192) (2.330,0.193) (2.365,0.193) (2.400,0.193) (2.435,0.193) (2.471,0.194) (2.506,0.194) (2.541,0.194) (2.577,0.194) (2.612,0.195) (2.647,0.195) (2.682,0.195) (2.718,0.195) (2.753,0.195) (2.788,0.196) (2.824,0.196) (2.859,0.196) (2.894,0.196) (2.929,0.196) (2.965,0.196) (3.000,0.196)};
\addplot[lineaE, flechaen=0.423] coordinates {(-3.000,0.196) (-2.965,0.196) (-2.929,0.196) (-2.894,0.196) (-2.859,0.196) (-2.824,0.196) (-2.788,0.196) (-2.753,0.195) (-2.718,0.195) (-2.682,0.195) (-2.647,0.195) (-2.612,0.195) (-2.577,0.194) (-2.541,0.194) (-2.506,0.194) (-2.471,0.194) (-2.435,0.193) (-2.400,0.193) (-2.365,0.193) (-2.330,0.193) (-2.294,0.192) (-2.259,0.192) (-2.224,0.192) (-2.188,0.191) (-2.153,0.191) (-2.118,0.190) (-2.082,0.190) (-2.047,0.189) (-2.012,0.189) (-1.977,0.188) (-1.941,0.188) (-1.906,0.187) (-1.871,0.186) (-1.835,0.186) (-1.800,0.185) (-1.765,0.184) (-1.730,0.183) (-1.694,0.182) (-1.659,0.181) (-1.624,0.180) (-1.589,0.179) (-1.553,0.178) (-1.518,0.177) (-1.483,0.175) (-1.447,0.174) (-1.412,0.172) (-1.377,0.171) (-1.342,0.169) (-1.307,0.167) (-1.271,0.165) (-1.236,0.162) (-1.201,0.160) (-1.166,0.157) (-1.130,0.155) (-1.095,0.152) (-1.060,0.148) (-1.025,0.145) (-0.990,0.141)};
\addplot[lineaE, flechaen=0.585] coordinates {(0.959,0.282) (0.994,0.290) (1.028,0.297) (1.063,0.303) (1.098,0.309) (1.133,0.315) (1.168,0.320) (1.203,0.325) (1.238,0.329) (1.273,0.333) (1.308,0.337) (1.343,0.341) (1.378,0.344) (1.413,0.348) (1.448,0.351) (1.483,0.353) (1.519,0.356) (1.554,0.358) (1.589,0.360) (1.624,0.363) (1.659,0.365) (1.695,0.366) (1.730,0.368) (1.765,0.370) (1.800,0.371) (1.836,0.373) (1.871,0.374) (1.906,0.375) (1.942,0.376) (1.977,0.378) (2.012,0.379) (2.047,0.380) (2.083,0.380) (2.118,0.381) (2.153,0.382) (2.188,0.383) (2.224,0.384) (2.259,0.384) (2.294,0.385) (2.330,0.386) (2.365,0.386) (2.400,0.387) (2.435,0.387) (2.471,0.388) (2.506,0.388) (2.541,0.389) (2.577,0.389) (2.612,0.390) (2.647,0.390) (2.682,0.390) (2.718,0.391) (2.753,0.391) (2.788,0.391) (2.824,0.392) (2.859,0.392) (2.894,0.392) (2.929,0.392) (2.965,0.393) (3.000,0.393)};
\addplot[lineaE, flechaen=0.415] coordinates {(-3.000,0.393) (-2.965,0.393) (-2.929,0.392) (-2.894,0.392) (-2.859,0.392) (-2.824,0.392) (-2.788,0.391) (-2.753,0.391) (-2.718,0.391) (-2.682,0.390) (-2.647,0.390) (-2.612,0.390) (-2.577,0.389) (-2.541,0.389) (-2.506,0.388) (-2.471,0.388) (-2.435,0.387) (-2.400,0.387) (-2.365,0.386) (-2.330,0.386) (-2.294,0.385) (-2.259,0.384) (-2.224,0.384) (-2.188,0.383) (-2.153,0.382) (-2.118,0.381) (-2.083,0.380) (-2.047,0.380) (-2.012,0.379) (-1.977,0.378) (-1.942,0.376) (-1.906,0.375) (-1.871,0.374) (-1.836,0.373) (-1.800,0.371) (-1.765,0.370) (-1.730,0.368) (-1.695,0.366) (-1.659,0.365) (-1.624,0.363) (-1.589,0.360) (-1.554,0.358) (-1.519,0.356) (-1.483,0.353) (-1.448,0.351) (-1.413,0.348) (-1.378,0.344) (-1.343,0.341) (-1.308,0.337) (-1.273,0.333) (-1.238,0.329) (-1.203,0.325) (-1.168,0.320) (-1.133,0.315) (-1.098,0.309) (-1.063,0.303) (-1.028,0.297) (-0.994,0.290) (-0.959,0.282)};
\addplot[lineaE, flechaen=0.597] coordinates {(0.906,0.424) (0.939,0.435) (0.973,0.445) (1.006,0.455) (1.040,0.463) (1.075,0.472) (1.109,0.479) (1.143,0.486) (1.178,0.493) (1.212,0.499) (1.247,0.505) (1.282,0.511) (1.316,0.516) (1.351,0.520) (1.386,0.525) (1.421,0.529) (1.456,0.533) (1.491,0.536) (1.526,0.540) (1.561,0.543) (1.596,0.546) (1.631,0.549) (1.666,0.551) (1.701,0.554) (1.736,0.556) (1.771,0.558) (1.806,0.560) (1.841,0.562) (1.876,0.564) (1.911,0.565) (1.946,0.567) (1.982,0.568) (2.017,0.570) (2.052,0.571) (2.087,0.572) (2.122,0.574) (2.157,0.575) (2.192,0.576) (2.227,0.577) (2.262,0.578) (2.298,0.579) (2.333,0.579) (2.368,0.580) (2.403,0.581) (2.438,0.582) (2.473,0.582) (2.508,0.583) (2.543,0.584) (2.579,0.584) (2.614,0.585) (2.649,0.586) (2.684,0.586) (2.719,0.587) (2.754,0.587) (2.789,0.587) (2.824,0.588) (2.859,0.588) (2.895,0.589) (2.930,0.589) (2.965,0.589) (3.000,0.590)};
\addplot[lineaE, flechaen=0.403] coordinates {(-3.000,0.590) (-2.965,0.589) (-2.930,0.589) (-2.895,0.589) (-2.859,0.588) (-2.824,0.588) (-2.789,0.587) (-2.754,0.587) (-2.719,0.587) (-2.684,0.586) (-2.649,0.586) (-2.614,0.585) (-2.579,0.584) (-2.543,0.584) (-2.508,0.583) (-2.473,0.582) (-2.438,0.582) (-2.403,0.581) (-2.368,0.580) (-2.333,0.579) (-2.298,0.579) (-2.262,0.578) (-2.227,0.577) (-2.192,0.576) (-2.157,0.575) (-2.122,0.574) (-2.087,0.572) (-2.052,0.571) (-2.017,0.570) (-1.982,0.568) (-1.946,0.567) (-1.911,0.565) (-1.876,0.564) (-1.841,0.562) (-1.806,0.560) (-1.771,0.558) (-1.736,0.556) (-1.701,0.554) (-1.666,0.551) (-1.631,0.549) (-1.596,0.546) (-1.561,0.543) (-1.526,0.540) (-1.491,0.536) (-1.456,0.533) (-1.421,0.529) (-1.386,0.525) (-1.351,0.520) (-1.316,0.516) (-1.282,0.511) (-1.247,0.505) (-1.212,0.499) (-1.178,0.493) (-1.143,0.486) (-1.109,0.479) (-1.075,0.472) (-1.040,0.463) (-1.006,0.455) (-0.973,0.445) (-0.939,0.435) (-0.906,0.424)};
\addplot[lineaE, flechaen=0.613] coordinates {(0.825,0.565) (0.857,0.580) (0.890,0.593) (0.923,0.606) (0.957,0.618) (0.990,0.629) (1.024,0.639) (1.058,0.648) (1.093,0.657) (1.127,0.665) (1.162,0.673) (1.197,0.680) (1.231,0.687) (1.266,0.693) (1.301,0.699) (1.336,0.704) (1.371,0.709) (1.407,0.714) (1.442,0.719) (1.477,0.723) (1.512,0.727) (1.548,0.730) (1.583,0.734) (1.618,0.737) (1.653,0.740) (1.689,0.743) (1.724,0.746) (1.760,0.748) (1.795,0.751) (1.830,0.753) (1.866,0.755) (1.901,0.757) (1.937,0.759) (1.972,0.761) (2.007,0.762) (2.043,0.764) (2.078,0.766) (2.114,0.767) (2.149,0.768) (2.185,0.770) (2.220,0.771) (2.255,0.772) (2.291,0.773) (2.326,0.774) (2.362,0.775) (2.397,0.776) (2.433,0.777) (2.468,0.778) (2.504,0.779) (2.539,0.779) (2.575,0.780) (2.610,0.781) (2.645,0.782) (2.681,0.782) (2.716,0.783) (2.752,0.783) (2.787,0.784) (2.823,0.785) (2.858,0.785) (2.894,0.786) (2.929,0.786) (2.965,0.786) (3.000,0.787)};
\addplot[lineaE, flechaen=0.387] coordinates {(-3.000,0.787) (-2.965,0.786) (-2.929,0.786) (-2.894,0.786) (-2.858,0.785) (-2.823,0.785) (-2.787,0.784) (-2.752,0.783) (-2.716,0.783) (-2.681,0.782) (-2.645,0.782) (-2.610,0.781) (-2.575,0.780) (-2.539,0.779) (-2.504,0.779) (-2.468,0.778) (-2.433,0.777) (-2.397,0.776) (-2.362,0.775) (-2.326,0.774) (-2.291,0.773) (-2.255,0.772) (-2.220,0.771) (-2.185,0.770) (-2.149,0.768) (-2.114,0.767) (-2.078,0.766) (-2.043,0.764) (-2.007,0.762) (-1.972,0.761) (-1.937,0.759) (-1.901,0.757) (-1.866,0.755) (-1.830,0.753) (-1.795,0.751) (-1.760,0.748) (-1.724,0.746) (-1.689,0.743) (-1.653,0.740) (-1.618,0.737) (-1.583,0.734) (-1.548,0.730) (-1.512,0.727) (-1.477,0.723) (-1.442,0.719) (-1.407,0.714) (-1.371,0.709) (-1.336,0.704) (-1.301,0.699) (-1.266,0.693) (-1.231,0.687) (-1.197,0.680) (-1.162,0.673) (-1.127,0.665) (-1.093,0.657) (-1.058,0.648) (-1.024,0.639) (-0.990,0.629) (-0.957,0.618) (-0.923,0.606) (-0.890,0.593) (-0.857,0.580) (-0.825,0.565)};
\addplot[lineaE, flechaen=0.635] coordinates {(0.708,0.707) (0.738,0.725) (0.769,0.741) (0.801,0.756) (0.833,0.771) (0.866,0.784) (0.899,0.796) (0.932,0.808) (0.966,0.818) (1.000,0.828) (1.034,0.838) (1.068,0.847) (1.102,0.855) (1.137,0.862) (1.171,0.870) (1.206,0.876) (1.240,0.883) (1.275,0.889) (1.310,0.894) (1.345,0.899) (1.380,0.904) (1.415,0.909) (1.450,0.913) (1.485,0.917) (1.520,0.921) (1.555,0.925) (1.590,0.928) (1.625,0.932) (1.661,0.935) (1.696,0.938) (1.731,0.940) (1.766,0.943) (1.801,0.945) (1.836,0.948) (1.872,0.950) (1.907,0.952) (1.942,0.954) (1.977,0.956) (2.013,0.958) (2.048,0.959) (2.083,0.961) (2.118,0.963) (2.154,0.964) (2.189,0.965) (2.224,0.967) (2.259,0.968) (2.295,0.969) (2.330,0.970) (2.365,0.971) (2.400,0.972) (2.436,0.973) (2.471,0.974) (2.506,0.975) (2.541,0.976) (2.577,0.977) (2.612,0.978) (2.647,0.978) (2.683,0.979) (2.718,0.980) (2.753,0.981) (2.788,0.981) (2.824,0.982) (2.859,0.982) (2.894,0.983) (2.929,0.983) (2.965,0.984) (3.000,0.984)};
\addplot[lineaE, flechaen=0.365] coordinates {(-3.000,0.984) (-2.965,0.984) (-2.929,0.983) (-2.894,0.983) (-2.859,0.982) (-2.824,0.982) (-2.788,0.981) (-2.753,0.981) (-2.718,0.980) (-2.683,0.979) (-2.647,0.978) (-2.612,0.978) (-2.577,0.977) (-2.541,0.976) (-2.506,0.975) (-2.471,0.974) (-2.436,0.973) (-2.400,0.972) (-2.365,0.971) (-2.330,0.970) (-2.295,0.969) (-2.259,0.968) (-2.224,0.967) (-2.189,0.965) (-2.154,0.964) (-2.118,0.963) (-2.083,0.961) (-2.048,0.959) (-2.013,0.958) (-1.977,0.956) (-1.942,0.954) (-1.907,0.952) (-1.872,0.950) (-1.836,0.948) (-1.801,0.945) (-1.766,0.943) (-1.731,0.940) (-1.696,0.938) (-1.661,0.935) (-1.625,0.932) (-1.590,0.928) (-1.555,0.925) (-1.520,0.921) (-1.485,0.917) (-1.450,0.913) (-1.415,0.909) (-1.380,0.904) (-1.345,0.899) (-1.310,0.894) (-1.275,0.889) (-1.240,0.883) (-1.206,0.876) (-1.171,0.870) (-1.137,0.862) (-1.102,0.855) (-1.068,0.847) (-1.034,0.838) (-1.000,0.828) (-0.966,0.818) (-0.932,0.808) (-0.899,0.796) (-0.866,0.784) (-0.833,0.771) (-0.801,0.756) (-0.769,0.741) (-0.738,0.725) (-0.708,0.707)};
\addplot[lineaE, flechaen=0.662] coordinates {(0.531,0.848) (0.559,0.869) (0.589,0.888) (0.620,0.905) (0.651,0.921) (0.683,0.936) (0.716,0.951) (0.749,0.964) (0.782,0.976) (0.816,0.988) (0.849,0.999) (0.883,1.009) (0.917,1.018) (0.952,1.027) (0.986,1.036) (1.021,1.044) (1.055,1.051) (1.090,1.059) (1.125,1.065) (1.160,1.072) (1.195,1.078) (1.230,1.083) (1.265,1.089) (1.300,1.094) (1.335,1.098) (1.371,1.103) (1.406,1.107) (1.441,1.111) (1.476,1.115) (1.512,1.119) (1.547,1.122) (1.582,1.126) (1.618,1.129) (1.653,1.132) (1.688,1.134) (1.724,1.137) (1.759,1.140) (1.794,1.142) (1.830,1.144) (1.865,1.147) (1.901,1.149) (1.936,1.151) (1.972,1.153) (2.007,1.154) (2.042,1.156) (2.078,1.158) (2.113,1.159) (2.149,1.161) (2.184,1.162) (2.220,1.164) (2.255,1.165) (2.291,1.166) (2.326,1.167) (2.361,1.168) (2.397,1.170) (2.432,1.171) (2.468,1.172) (2.503,1.173) (2.539,1.173) (2.574,1.174) (2.610,1.175) (2.645,1.176) (2.681,1.177) (2.716,1.178) (2.752,1.178) (2.787,1.179) (2.823,1.180) (2.858,1.180) (2.894,1.181) (2.929,1.181) (2.965,1.182) (3.000,1.182)};
\addplot[lineaE, flechaen=0.338] coordinates {(-3.000,1.182) (-2.965,1.182) (-2.929,1.181) (-2.894,1.181) (-2.858,1.180) (-2.823,1.180) (-2.787,1.179) (-2.752,1.178) (-2.716,1.178) (-2.681,1.177) (-2.645,1.176) (-2.610,1.175) (-2.574,1.174) (-2.539,1.173) (-2.503,1.173) (-2.468,1.172) (-2.432,1.171) (-2.397,1.170) (-2.361,1.168) (-2.326,1.167) (-2.291,1.166) (-2.255,1.165) (-2.220,1.164) (-2.184,1.162) (-2.149,1.161) (-2.113,1.159) (-2.078,1.158) (-2.042,1.156) (-2.007,1.154) (-1.972,1.153) (-1.936,1.151) (-1.901,1.149) (-1.865,1.147) (-1.830,1.144) (-1.794,1.142) (-1.759,1.140) (-1.724,1.137) (-1.688,1.134) (-1.653,1.132) (-1.618,1.129) (-1.582,1.126) (-1.547,1.122) (-1.512,1.119) (-1.476,1.115) (-1.441,1.111) (-1.406,1.107) (-1.371,1.103) (-1.335,1.098) (-1.300,1.094) (-1.265,1.089) (-1.230,1.083) (-1.195,1.078) (-1.160,1.072) (-1.125,1.065) (-1.090,1.059) (-1.055,1.051) (-1.021,1.044) (-0.986,1.036) (-0.952,1.027) (-0.917,1.018) (-0.883,1.009) (-0.849,0.999) (-0.816,0.988) (-0.782,0.976) (-0.749,0.964) (-0.716,0.951) (-0.683,0.936) (-0.651,0.921) (-0.620,0.905) (-0.589,0.888) (-0.559,0.869) (-0.531,0.848)};
\addplot[lineaE, flechaen=0.707] coordinates {(0.148,0.989) (0.178,1.005) (0.210,1.019) (0.243,1.033) (0.275,1.047) (0.308,1.061) (0.340,1.074) (0.373,1.087) (0.406,1.100) (0.440,1.112) (0.473,1.124) (0.507,1.135) (0.540,1.146) (0.574,1.157) (0.608,1.167) (0.642,1.177) (0.676,1.186) (0.710,1.195) (0.745,1.204) (0.779,1.212) (0.814,1.220) (0.848,1.227) (0.883,1.235) (0.918,1.241) (0.953,1.248) (0.987,1.254) (1.022,1.260) (1.057,1.266) (1.092,1.271) (1.127,1.277) (1.162,1.282) (1.197,1.287) (1.232,1.291) (1.268,1.295) (1.303,1.300) (1.338,1.304) (1.373,1.307) (1.408,1.311) (1.444,1.314) (1.479,1.318) (1.514,1.321) (1.549,1.324) (1.585,1.327) (1.620,1.330) (1.655,1.332) (1.691,1.335) (1.726,1.337) (1.761,1.340) (1.797,1.342) (1.832,1.344) (1.867,1.346) (1.903,1.348) (1.938,1.350) (1.974,1.351) (2.009,1.353) (2.044,1.355) (2.080,1.356) (2.115,1.358) (2.150,1.359) (2.186,1.361) (2.221,1.362) (2.257,1.363) (2.292,1.365) (2.327,1.366) (2.363,1.367) (2.398,1.368) (2.434,1.369) (2.469,1.370) (2.504,1.371) (2.540,1.372) (2.575,1.373) (2.611,1.373) (2.646,1.374) (2.681,1.375) (2.717,1.376) (2.752,1.377) (2.788,1.377) (2.823,1.378) (2.858,1.379) (2.894,1.379) (2.929,1.380) (2.965,1.380) (3.000,1.381)};
\addplot[lineaE, flechaen=0.293] coordinates {(-3.000,1.381) (-2.965,1.380) (-2.929,1.380) (-2.894,1.379) (-2.858,1.379) (-2.823,1.378) (-2.788,1.377) (-2.752,1.377) (-2.717,1.376) (-2.681,1.375) (-2.646,1.374) (-2.611,1.373) (-2.575,1.373) (-2.540,1.372) (-2.504,1.371) (-2.469,1.370) (-2.434,1.369) (-2.398,1.368) (-2.363,1.367) (-2.327,1.366) (-2.292,1.365) (-2.257,1.363) (-2.221,1.362) (-2.186,1.361) (-2.150,1.359) (-2.115,1.358) (-2.080,1.356) (-2.044,1.355) (-2.009,1.353) (-1.974,1.351) (-1.938,1.350) (-1.903,1.348) (-1.867,1.346) (-1.832,1.344) (-1.797,1.342) (-1.761,1.340) (-1.726,1.337) (-1.691,1.335) (-1.655,1.332) (-1.620,1.330) (-1.585,1.327) (-1.549,1.324) (-1.514,1.321) (-1.479,1.318) (-1.444,1.314) (-1.408,1.311) (-1.373,1.307) (-1.338,1.304) (-1.303,1.300) (-1.268,1.295) (-1.232,1.291) (-1.197,1.287) (-1.162,1.282) (-1.127,1.277) (-1.092,1.271) (-1.057,1.266) (-1.022,1.260) (-0.987,1.254) (-0.953,1.248) (-0.918,1.241) (-0.883,1.235) (-0.848,1.227) (-0.814,1.220) (-0.779,1.212) (-0.745,1.204) (-0.710,1.195) (-0.676,1.186) (-0.642,1.177) (-0.608,1.167) (-0.574,1.157) (-0.540,1.146) (-0.507,1.135) (-0.473,1.124) (-0.440,1.112) (-0.406,1.100) (-0.373,1.087) (-0.340,1.074) (-0.308,1.061) (-0.275,1.047) (-0.243,1.033) (-0.210,1.019) (-0.178,1.005) (-0.148,0.989)};
\addplot[lineaE, flechaen=0.141, flechaen=0.859] coordinates {(-3.000,1.580) (-2.965,1.579) (-2.930,1.579) (-2.895,1.578) (-2.860,1.578) (-2.825,1.577) (-2.790,1.576) (-2.755,1.576) (-2.720,1.575) (-2.685,1.574) (-2.650,1.573) (-2.615,1.573) (-2.580,1.572) (-2.545,1.571) (-2.510,1.570) (-2.475,1.569) (-2.440,1.568) (-2.405,1.567) (-2.370,1.566) (-2.335,1.565) (-2.300,1.564) (-2.265,1.563) (-2.230,1.562) (-2.195,1.561) (-2.160,1.559) (-2.125,1.558) (-2.090,1.557) (-2.055,1.555) (-2.020,1.554) (-1.985,1.552) (-1.950,1.551) (-1.915,1.549) (-1.880,1.547) (-1.845,1.546) (-1.810,1.544) (-1.775,1.542) (-1.740,1.540) (-1.705,1.538) (-1.670,1.536) (-1.635,1.533) (-1.600,1.531) (-1.565,1.529) (-1.530,1.526) (-1.495,1.524) (-1.460,1.521) (-1.425,1.518) (-1.390,1.515) (-1.355,1.512) (-1.320,1.509) (-1.286,1.506) (-1.251,1.503) (-1.216,1.499) (-1.181,1.496) (-1.146,1.492) (-1.111,1.488) (-1.077,1.484) (-1.042,1.480) (-1.007,1.476) (-0.972,1.471) (-0.937,1.467) (-0.903,1.462) (-0.868,1.458) (-0.833,1.453) (-0.799,1.448) (-0.764,1.443) (-0.729,1.438) (-0.695,1.432) (-0.660,1.427) (-0.626,1.421) (-0.591,1.416) (-0.556,1.410) (-0.522,1.405) (-0.487,1.399) (-0.453,1.394) (-0.418,1.388) (-0.383,1.383) (-0.349,1.378) (-0.314,1.373) (-0.279,1.368) (-0.245,1.364) (-0.210,1.360) (-0.175,1.357) (-0.140,1.354) (-0.105,1.352) (-0.070,1.350) (-0.035,1.349) (0.000,1.348) (0.035,1.349) (0.070,1.350) (0.105,1.352) (0.140,1.354) (0.175,1.357) (0.210,1.360) (0.245,1.364) (0.279,1.368) (0.314,1.373) (0.349,1.378) (0.383,1.383) (0.418,1.388) (0.453,1.394) (0.487,1.399) (0.522,1.405) (0.556,1.410) (0.591,1.416) (0.626,1.421) (0.660,1.427) (0.695,1.432) (0.729,1.438) (0.764,1.443) (0.799,1.448) (0.833,1.453) (0.868,1.458) (0.903,1.462) (0.937,1.467) (0.972,1.471) (1.007,1.476) (1.042,1.480) (1.077,1.484) (1.111,1.488) (1.146,1.492) (1.181,1.496) (1.216,1.499) (1.251,1.503) (1.286,1.506) (1.320,1.509) (1.355,1.512) (1.390,1.515) (1.425,1.518) (1.460,1.521) (1.495,1.524) (1.530,1.526) (1.565,1.529) (1.600,1.531) (1.635,1.533) (1.670,1.536) (1.705,1.538) (1.740,1.540) (1.775,1.542) (1.810,1.544) (1.845,1.546) (1.880,1.547) (1.915,1.549) (1.950,1.551) (1.985,1.552) (2.020,1.554) (2.055,1.555) (2.090,1.557) (2.125,1.558) (2.160,1.559) (2.195,1.561) (2.230,1.562) (2.265,1.563) (2.300,1.564) (2.335,1.565) (2.370,1.566) (2.405,1.567) (2.440,1.568) (2.475,1.569) (2.510,1.570) (2.545,1.571) (2.580,1.572) (2.615,1.573) (2.650,1.573) (2.685,1.574) (2.720,1.575) (2.755,1.576) (2.790,1.576) (2.825,1.577) (2.860,1.578) (2.895,1.578) (2.930,1.579) (2.965,1.579) (3.000,1.580)};
\addplot[lineaE, flechaen=0.141, flechaen=0.859] coordinates {(-3.000,1.779) (-2.965,1.779) (-2.930,1.778) (-2.895,1.777) (-2.859,1.777) (-2.824,1.776) (-2.789,1.776) (-2.754,1.775) (-2.719,1.774) (-2.684,1.774) (-2.649,1.773) (-2.613,1.772) (-2.578,1.771) (-2.543,1.770) (-2.508,1.770) (-2.473,1.769) (-2.438,1.768) (-2.403,1.767) (-2.367,1.766) (-2.332,1.765) (-2.297,1.764) (-2.262,1.763) (-2.227,1.762) (-2.192,1.761) (-2.157,1.760) (-2.121,1.758) (-2.086,1.757) (-2.051,1.756) (-2.016,1.755) (-1.981,1.753) (-1.946,1.752) (-1.911,1.750) (-1.876,1.749) (-1.841,1.747) (-1.805,1.746) (-1.770,1.744) (-1.735,1.742) (-1.700,1.741) (-1.665,1.739) (-1.630,1.737) (-1.595,1.735) (-1.560,1.733) (-1.525,1.731) (-1.489,1.729) (-1.454,1.727) (-1.419,1.724) (-1.384,1.722) (-1.349,1.720) (-1.314,1.717) (-1.279,1.715) (-1.244,1.712) (-1.209,1.709) (-1.174,1.707) (-1.139,1.704) (-1.104,1.701) (-1.069,1.698) (-1.034,1.695) (-0.999,1.692) (-0.964,1.689) (-0.929,1.686) (-0.894,1.682) (-0.859,1.679) (-0.824,1.676) (-0.789,1.673) (-0.754,1.669) (-0.719,1.666) (-0.684,1.663) (-0.649,1.659) (-0.614,1.656) (-0.579,1.653) (-0.544,1.649) (-0.509,1.646) (-0.474,1.643) (-0.439,1.640) (-0.404,1.637) (-0.369,1.635) (-0.334,1.632) (-0.299,1.630) (-0.263,1.628) (-0.228,1.626) (-0.193,1.624) (-0.158,1.622) (-0.123,1.621) (-0.088,1.620) (-0.053,1.620) (-0.018,1.619) (0.018,1.619) (0.053,1.620) (0.088,1.620) (0.123,1.621) (0.158,1.622) (0.193,1.624) (0.228,1.626) (0.263,1.628) (0.299,1.630) (0.334,1.632) (0.369,1.635) (0.404,1.637) (0.439,1.640) (0.474,1.643) (0.509,1.646) (0.544,1.649) (0.579,1.653) (0.614,1.656) (0.649,1.659) (0.684,1.663) (0.719,1.666) (0.754,1.669) (0.789,1.673) (0.824,1.676) (0.859,1.679) (0.894,1.682) (0.929,1.686) (0.964,1.689) (0.999,1.692) (1.034,1.695) (1.069,1.698) (1.104,1.701) (1.139,1.704) (1.174,1.707) (1.209,1.709) (1.244,1.712) (1.279,1.715) (1.314,1.717) (1.349,1.720) (1.384,1.722) (1.419,1.724) (1.454,1.727) (1.489,1.729) (1.525,1.731) (1.560,1.733) (1.595,1.735) (1.630,1.737) (1.665,1.739) (1.700,1.741) (1.735,1.742) (1.770,1.744) (1.805,1.746) (1.841,1.747) (1.876,1.749) (1.911,1.750) (1.946,1.752) (1.981,1.753) (2.016,1.755) (2.051,1.756) (2.086,1.757) (2.121,1.758) (2.157,1.760) (2.192,1.761) (2.227,1.762) (2.262,1.763) (2.297,1.764) (2.332,1.765) (2.367,1.766) (2.403,1.767) (2.438,1.768) (2.473,1.769) (2.508,1.770) (2.543,1.770) (2.578,1.771) (2.613,1.772) (2.649,1.773) (2.684,1.774) (2.719,1.774) (2.754,1.775) (2.789,1.776) (2.824,1.776) (2.859,1.777) (2.895,1.777) (2.930,1.778) (2.965,1.779) (3.000,1.779)};
\addplot[lineaE, flechaen=0.142, flechaen=0.858] coordinates {(-3.000,1.979) (-2.965,1.978) (-2.930,1.978) (-2.895,1.977) (-2.860,1.977) (-2.824,1.976) (-2.789,1.975) (-2.754,1.975) (-2.719,1.974) (-2.684,1.973) (-2.649,1.973) (-2.614,1.972) (-2.579,1.971) (-2.544,1.971) (-2.508,1.970) (-2.473,1.969) (-2.438,1.968) (-2.403,1.967) (-2.368,1.967) (-2.333,1.966) (-2.298,1.965) (-2.263,1.964) (-2.228,1.963) (-2.192,1.962) (-2.157,1.961) (-2.122,1.960) (-2.087,1.959) (-2.052,1.958) (-2.017,1.956) (-1.982,1.955) (-1.947,1.954) (-1.912,1.953) (-1.877,1.951) (-1.841,1.950) (-1.806,1.949) (-1.771,1.947) (-1.736,1.946) (-1.701,1.944) (-1.666,1.943) (-1.631,1.941) (-1.596,1.940) (-1.561,1.938) (-1.526,1.936) (-1.491,1.935) (-1.455,1.933) (-1.420,1.931) (-1.385,1.929) (-1.350,1.927) (-1.315,1.926) (-1.280,1.924) (-1.245,1.922) (-1.210,1.920) (-1.175,1.917) (-1.140,1.915) (-1.105,1.913) (-1.070,1.911) (-1.035,1.909) (-1.000,1.907) (-0.965,1.905) (-0.930,1.902) (-0.895,1.900) (-0.859,1.898) (-0.824,1.896) (-0.789,1.893) (-0.754,1.891) (-0.719,1.889) (-0.684,1.887) (-0.649,1.885) (-0.614,1.882) (-0.579,1.880) (-0.544,1.878) (-0.509,1.876) (-0.474,1.875) (-0.439,1.873) (-0.404,1.871) (-0.369,1.869) (-0.334,1.868) (-0.298,1.867) (-0.263,1.865) (-0.228,1.864) (-0.193,1.863) (-0.158,1.862) (-0.123,1.862) (-0.088,1.861) (-0.053,1.861) (-0.018,1.861) (0.018,1.861) (0.053,1.861) (0.088,1.861) (0.123,1.862) (0.158,1.862) (0.193,1.863) (0.228,1.864) (0.263,1.865) (0.298,1.867) (0.334,1.868) (0.369,1.869) (0.404,1.871) (0.439,1.873) (0.474,1.875) (0.509,1.876) (0.544,1.878) (0.579,1.880) (0.614,1.882) (0.649,1.885) (0.684,1.887) (0.719,1.889) (0.754,1.891) (0.789,1.893) (0.824,1.896) (0.859,1.898) (0.895,1.900) (0.930,1.902) (0.965,1.905) (1.000,1.907) (1.035,1.909) (1.070,1.911) (1.105,1.913) (1.140,1.915) (1.175,1.917) (1.210,1.920) (1.245,1.922) (1.280,1.924) (1.315,1.926) (1.350,1.927) (1.385,1.929) (1.420,1.931) (1.455,1.933) (1.491,1.935) (1.526,1.936) (1.561,1.938) (1.596,1.940) (1.631,1.941) (1.666,1.943) (1.701,1.944) (1.736,1.946) (1.771,1.947) (1.806,1.949) (1.841,1.950) (1.877,1.951) (1.912,1.953) (1.947,1.954) (1.982,1.955) (2.017,1.956) (2.052,1.958) (2.087,1.959) (2.122,1.960) (2.157,1.961) (2.192,1.962) (2.228,1.963) (2.263,1.964) (2.298,1.965) (2.333,1.966) (2.368,1.967) (2.403,1.967) (2.438,1.968) (2.473,1.969) (2.508,1.970) (2.544,1.971) (2.579,1.971) (2.614,1.972) (2.649,1.973) (2.684,1.973) (2.719,1.974) (2.754,1.975) (2.789,1.975) (2.824,1.976) (2.860,1.977) (2.895,1.977) (2.930,1.978) (2.965,1.978) (3.000,1.979)};
\addplot[lineaE, flechaen=0.142, flechaen=0.858] coordinates {(-3.000,2.179) (-2.965,2.178) (-2.930,2.178) (-2.895,2.177) (-2.860,2.177) (-2.824,2.176) (-2.789,2.176) (-2.754,2.175) (-2.719,2.174) (-2.684,2.174) (-2.649,2.173) (-2.614,2.173) (-2.579,2.172) (-2.544,2.171) (-2.509,2.170) (-2.473,2.170) (-2.438,2.169) (-2.403,2.168) (-2.368,2.167) (-2.333,2.167) (-2.298,2.166) (-2.263,2.165) (-2.228,2.164) (-2.193,2.163) (-2.158,2.162) (-2.123,2.161) (-2.087,2.161) (-2.052,2.160) (-2.017,2.159) (-1.982,2.158) (-1.947,2.156) (-1.912,2.155) (-1.877,2.154) (-1.842,2.153) (-1.807,2.152) (-1.772,2.151) (-1.737,2.150) (-1.701,2.148) (-1.666,2.147) (-1.631,2.146) (-1.596,2.145) (-1.561,2.143) (-1.526,2.142) (-1.491,2.141) (-1.456,2.139) (-1.421,2.138) (-1.386,2.136) (-1.351,2.135) (-1.316,2.133) (-1.281,2.132) (-1.245,2.130) (-1.210,2.129) (-1.175,2.127) (-1.140,2.126) (-1.105,2.124) (-1.070,2.122) (-1.035,2.121) (-1.000,2.119) (-0.965,2.118) (-0.930,2.116) (-0.895,2.114) (-0.860,2.113) (-0.825,2.111) (-0.790,2.110) (-0.754,2.108) (-0.719,2.107) (-0.684,2.105) (-0.649,2.104) (-0.614,2.102) (-0.579,2.101) (-0.544,2.100) (-0.509,2.098) (-0.474,2.097) (-0.439,2.096) (-0.404,2.095) (-0.369,2.094) (-0.333,2.093) (-0.298,2.092) (-0.263,2.091) (-0.228,2.090) (-0.193,2.090) (-0.158,2.089) (-0.123,2.089) (-0.088,2.089) (-0.053,2.088) (-0.018,2.088) (0.018,2.088) (0.053,2.088) (0.088,2.089) (0.123,2.089) (0.158,2.089) (0.193,2.090) (0.228,2.090) (0.263,2.091) (0.298,2.092) (0.333,2.093) (0.369,2.094) (0.404,2.095) (0.439,2.096) (0.474,2.097) (0.509,2.098) (0.544,2.100) (0.579,2.101) (0.614,2.102) (0.649,2.104) (0.684,2.105) (0.719,2.107) (0.754,2.108) (0.790,2.110) (0.825,2.111) (0.860,2.113) (0.895,2.114) (0.930,2.116) (0.965,2.118) (1.000,2.119) (1.035,2.121) (1.070,2.122) (1.105,2.124) (1.140,2.126) (1.175,2.127) (1.210,2.129) (1.245,2.130) (1.281,2.132) (1.316,2.133) (1.351,2.135) (1.386,2.136) (1.421,2.138) (1.456,2.139) (1.491,2.141) (1.526,2.142) (1.561,2.143) (1.596,2.145) (1.631,2.146) (1.666,2.147) (1.701,2.148) (1.737,2.150) (1.772,2.151) (1.807,2.152) (1.842,2.153) (1.877,2.154) (1.912,2.155) (1.947,2.156) (1.982,2.158) (2.017,2.159) (2.052,2.160) (2.087,2.161) (2.123,2.161) (2.158,2.162) (2.193,2.163) (2.228,2.164) (2.263,2.165) (2.298,2.166) (2.333,2.167) (2.368,2.167) (2.403,2.168) (2.438,2.169) (2.473,2.170) (2.509,2.170) (2.544,2.171) (2.579,2.172) (2.614,2.173) (2.649,2.173) (2.684,2.174) (2.719,2.174) (2.754,2.175) (2.789,2.176) (2.824,2.176) (2.860,2.177) (2.895,2.177) (2.930,2.178) (2.965,2.178) (3.000,2.179)};
\addplot[lineaE, flechaen=0.239] coordinates {(-0.959,0.283) (-0.924,0.283) (-0.888,0.283) (-0.853,0.283) (-0.817,0.283) (-0.781,0.283) (-0.746,0.283) (-0.710,0.283) (-0.675,0.283) (-0.639,0.283) (-0.604,0.283) (-0.568,0.283) (-0.533,0.283) (-0.497,0.283) (-0.462,0.283) (-0.426,0.283) (-0.391,0.283) (-0.355,0.283) (-0.320,0.283) (-0.284,0.283) (-0.249,0.283) (-0.213,0.283) (-0.178,0.283) (-0.142,0.283) (-0.107,0.283) (-0.071,0.283) (-0.036,0.283) (-0.000,0.283) (0.036,0.283) (0.071,0.283) (0.107,0.283) (0.142,0.283) (0.178,0.283) (0.213,0.283) (0.249,0.283) (0.284,0.283) (0.320,0.283) (0.355,0.283) (0.391,0.283) (0.426,0.283) (0.462,0.283) (0.497,0.283) (0.533,0.283) (0.568,0.283) (0.604,0.283) (0.639,0.283) (0.675,0.283) (0.710,0.283) (0.746,0.283) (0.781,0.283) (0.817,0.283) (0.853,0.283) (0.888,0.283) (0.924,0.283) (0.959,0.283)};
\addplot[lineaE, flechaen=0.804] coordinates {(-0.824,0.567) (-0.789,0.566) (-0.754,0.566) (-0.719,0.566) (-0.684,0.566) (-0.649,0.566) (-0.614,0.566) (-0.579,0.566) (-0.544,0.566) (-0.508,0.566) (-0.473,0.566) (-0.438,0.566) (-0.403,0.566) (-0.368,0.566) (-0.333,0.566) (-0.298,0.566) (-0.263,0.566) (-0.228,0.566) (-0.193,0.566) (-0.158,0.566) (-0.123,0.566) (-0.088,0.566) (-0.053,0.566) (-0.018,0.566) (0.018,0.566) (0.053,0.566) (0.088,0.566) (0.123,0.566) (0.158,0.566) (0.193,0.566) (0.228,0.566) (0.263,0.566) (0.298,0.566) (0.333,0.566) (0.368,0.566) (0.403,0.566) (0.438,0.566) (0.473,0.566) (0.508,0.566) (0.544,0.566) (0.579,0.566) (0.614,0.566) (0.649,0.566) (0.684,0.566) (0.719,0.566) (0.754,0.566) (0.789,0.566) (0.824,0.567)};
\addplot[lineaE, flechaen=0.500] coordinates {(-0.525,0.851) (-0.491,0.849) (-0.456,0.849) (-0.421,0.849) (-0.386,0.849) (-0.351,0.849) (-0.316,0.849) (-0.281,0.849) (-0.246,0.849) (-0.210,0.849) (-0.175,0.849) (-0.140,0.849) (-0.105,0.849) (-0.070,0.849) (-0.035,0.849) (0.000,0.849) (0.035,0.849) (0.070,0.849) (0.105,0.849) (0.140,0.849) (0.175,0.849) (0.210,0.849) (0.246,0.849) (0.281,0.849) (0.316,0.849) (0.351,0.849) (0.386,0.849) (0.421,0.849) (0.456,0.849) (0.491,0.849) (0.525,0.851)};
\addplot[lineaE, flechaen=0.577] coordinates {(0.990,-0.141) (1.025,-0.145) (1.060,-0.148) (1.095,-0.152) (1.130,-0.155) (1.166,-0.157) (1.201,-0.160) (1.236,-0.162) (1.271,-0.165) (1.307,-0.167) (1.342,-0.169) (1.377,-0.171) (1.412,-0.172) (1.447,-0.174) (1.483,-0.175) (1.518,-0.177) (1.553,-0.178) (1.589,-0.179) (1.624,-0.180) (1.659,-0.181) (1.694,-0.182) (1.730,-0.183) (1.765,-0.184) (1.800,-0.185) (1.835,-0.186) (1.871,-0.186) (1.906,-0.187) (1.941,-0.188) (1.977,-0.188) (2.012,-0.189) (2.047,-0.189) (2.082,-0.190) (2.118,-0.190) (2.153,-0.191) (2.188,-0.191) (2.224,-0.192) (2.259,-0.192) (2.294,-0.192) (2.330,-0.193) (2.365,-0.193) (2.400,-0.193) (2.435,-0.193) (2.471,-0.194) (2.506,-0.194) (2.541,-0.194) (2.577,-0.194) (2.612,-0.195) (2.647,-0.195) (2.682,-0.195) (2.718,-0.195) (2.753,-0.195) (2.788,-0.196) (2.824,-0.196) (2.859,-0.196) (2.894,-0.196) (2.929,-0.196) (2.965,-0.196) (3.000,-0.196)};
\addplot[lineaE, flechaen=0.423] coordinates {(-3.000,-0.196) (-2.965,-0.196) (-2.929,-0.196) (-2.894,-0.196) (-2.859,-0.196) (-2.824,-0.196) (-2.788,-0.196) (-2.753,-0.195) (-2.718,-0.195) (-2.682,-0.195) (-2.647,-0.195) (-2.612,-0.195) (-2.577,-0.194) (-2.541,-0.194) (-2.506,-0.194) (-2.471,-0.194) (-2.435,-0.193) (-2.400,-0.193) (-2.365,-0.193) (-2.330,-0.193) (-2.294,-0.192) (-2.259,-0.192) (-2.224,-0.192) (-2.188,-0.191) (-2.153,-0.191) (-2.118,-0.190) (-2.082,-0.190) (-2.047,-0.189) (-2.012,-0.189) (-1.977,-0.188) (-1.941,-0.188) (-1.906,-0.187) (-1.871,-0.186) (-1.835,-0.186) (-1.800,-0.185) (-1.765,-0.184) (-1.730,-0.183) (-1.694,-0.182) (-1.659,-0.181) (-1.624,-0.180) (-1.589,-0.179) (-1.553,-0.178) (-1.518,-0.177) (-1.483,-0.175) (-1.447,-0.174) (-1.412,-0.172) (-1.377,-0.171) (-1.342,-0.169) (-1.307,-0.167) (-1.271,-0.165) (-1.236,-0.162) (-1.201,-0.160) (-1.166,-0.157) (-1.130,-0.155) (-1.095,-0.152) (-1.060,-0.148) (-1.025,-0.145) (-0.990,-0.141)};
\addplot[lineaE, flechaen=0.585] coordinates {(0.959,-0.282) (0.994,-0.290) (1.028,-0.297) (1.063,-0.303) (1.098,-0.309) (1.133,-0.315) (1.168,-0.320) (1.203,-0.325) (1.238,-0.329) (1.273,-0.333) (1.308,-0.337) (1.343,-0.341) (1.378,-0.344) (1.413,-0.348) (1.448,-0.351) (1.483,-0.353) (1.519,-0.356) (1.554,-0.358) (1.589,-0.360) (1.624,-0.363) (1.659,-0.365) (1.695,-0.366) (1.730,-0.368) (1.765,-0.370) (1.800,-0.371) (1.836,-0.373) (1.871,-0.374) (1.906,-0.375) (1.942,-0.376) (1.977,-0.378) (2.012,-0.379) (2.047,-0.380) (2.083,-0.380) (2.118,-0.381) (2.153,-0.382) (2.188,-0.383) (2.224,-0.384) (2.259,-0.384) (2.294,-0.385) (2.330,-0.386) (2.365,-0.386) (2.400,-0.387) (2.435,-0.387) (2.471,-0.388) (2.506,-0.388) (2.541,-0.389) (2.577,-0.389) (2.612,-0.390) (2.647,-0.390) (2.682,-0.390) (2.718,-0.391) (2.753,-0.391) (2.788,-0.391) (2.824,-0.392) (2.859,-0.392) (2.894,-0.392) (2.929,-0.392) (2.965,-0.393) (3.000,-0.393)};
\addplot[lineaE, flechaen=0.415] coordinates {(-3.000,-0.393) (-2.965,-0.393) (-2.929,-0.392) (-2.894,-0.392) (-2.859,-0.392) (-2.824,-0.392) (-2.788,-0.391) (-2.753,-0.391) (-2.718,-0.391) (-2.682,-0.390) (-2.647,-0.390) (-2.612,-0.390) (-2.577,-0.389) (-2.541,-0.389) (-2.506,-0.388) (-2.471,-0.388) (-2.435,-0.387) (-2.400,-0.387) (-2.365,-0.386) (-2.330,-0.386) (-2.294,-0.385) (-2.259,-0.384) (-2.224,-0.384) (-2.188,-0.383) (-2.153,-0.382) (-2.118,-0.381) (-2.083,-0.380) (-2.047,-0.380) (-2.012,-0.379) (-1.977,-0.378) (-1.942,-0.376) (-1.906,-0.375) (-1.871,-0.374) (-1.836,-0.373) (-1.800,-0.371) (-1.765,-0.370) (-1.730,-0.368) (-1.695,-0.366) (-1.659,-0.365) (-1.624,-0.363) (-1.589,-0.360) (-1.554,-0.358) (-1.519,-0.356) (-1.483,-0.353) (-1.448,-0.351) (-1.413,-0.348) (-1.378,-0.344) (-1.343,-0.341) (-1.308,-0.337) (-1.273,-0.333) (-1.238,-0.329) (-1.203,-0.325) (-1.168,-0.320) (-1.133,-0.315) (-1.098,-0.309) (-1.063,-0.303) (-1.028,-0.297) (-0.994,-0.290) (-0.959,-0.282)};
\addplot[lineaE, flechaen=0.597] coordinates {(0.906,-0.424) (0.939,-0.435) (0.973,-0.445) (1.006,-0.455) (1.040,-0.463) (1.075,-0.472) (1.109,-0.479) (1.143,-0.486) (1.178,-0.493) (1.212,-0.499) (1.247,-0.505) (1.282,-0.511) (1.316,-0.516) (1.351,-0.520) (1.386,-0.525) (1.421,-0.529) (1.456,-0.533) (1.491,-0.536) (1.526,-0.540) (1.561,-0.543) (1.596,-0.546) (1.631,-0.549) (1.666,-0.551) (1.701,-0.554) (1.736,-0.556) (1.771,-0.558) (1.806,-0.560) (1.841,-0.562) (1.876,-0.564) (1.911,-0.565) (1.946,-0.567) (1.982,-0.568) (2.017,-0.570) (2.052,-0.571) (2.087,-0.572) (2.122,-0.574) (2.157,-0.575) (2.192,-0.576) (2.227,-0.577) (2.262,-0.578) (2.298,-0.579) (2.333,-0.579) (2.368,-0.580) (2.403,-0.581) (2.438,-0.582) (2.473,-0.582) (2.508,-0.583) (2.543,-0.584) (2.579,-0.584) (2.614,-0.585) (2.649,-0.586) (2.684,-0.586) (2.719,-0.587) (2.754,-0.587) (2.789,-0.587) (2.824,-0.588) (2.859,-0.588) (2.895,-0.589) (2.930,-0.589) (2.965,-0.589) (3.000,-0.590)};
\addplot[lineaE, flechaen=0.403] coordinates {(-3.000,-0.590) (-2.965,-0.589) (-2.930,-0.589) (-2.895,-0.589) (-2.859,-0.588) (-2.824,-0.588) (-2.789,-0.587) (-2.754,-0.587) (-2.719,-0.587) (-2.684,-0.586) (-2.649,-0.586) (-2.614,-0.585) (-2.579,-0.584) (-2.543,-0.584) (-2.508,-0.583) (-2.473,-0.582) (-2.438,-0.582) (-2.403,-0.581) (-2.368,-0.580) (-2.333,-0.579) (-2.298,-0.579) (-2.262,-0.578) (-2.227,-0.577) (-2.192,-0.576) (-2.157,-0.575) (-2.122,-0.574) (-2.087,-0.572) (-2.052,-0.571) (-2.017,-0.570) (-1.982,-0.568) (-1.946,-0.567) (-1.911,-0.565) (-1.876,-0.564) (-1.841,-0.562) (-1.806,-0.560) (-1.771,-0.558) (-1.736,-0.556) (-1.701,-0.554) (-1.666,-0.551) (-1.631,-0.549) (-1.596,-0.546) (-1.561,-0.543) (-1.526,-0.540) (-1.491,-0.536) (-1.456,-0.533) (-1.421,-0.529) (-1.386,-0.525) (-1.351,-0.520) (-1.316,-0.516) (-1.282,-0.511) (-1.247,-0.505) (-1.212,-0.499) (-1.178,-0.493) (-1.143,-0.486) (-1.109,-0.479) (-1.075,-0.472) (-1.040,-0.463) (-1.006,-0.455) (-0.973,-0.445) (-0.939,-0.435) (-0.906,-0.424)};
\addplot[lineaE, flechaen=0.613] coordinates {(0.825,-0.565) (0.857,-0.580) (0.890,-0.593) (0.923,-0.606) (0.957,-0.618) (0.990,-0.629) (1.024,-0.639) (1.058,-0.648) (1.093,-0.657) (1.127,-0.665) (1.162,-0.673) (1.197,-0.680) (1.231,-0.687) (1.266,-0.693) (1.301,-0.699) (1.336,-0.704) (1.371,-0.709) (1.407,-0.714) (1.442,-0.719) (1.477,-0.723) (1.512,-0.727) (1.548,-0.730) (1.583,-0.734) (1.618,-0.737) (1.653,-0.740) (1.689,-0.743) (1.724,-0.746) (1.760,-0.748) (1.795,-0.751) (1.830,-0.753) (1.866,-0.755) (1.901,-0.757) (1.937,-0.759) (1.972,-0.761) (2.007,-0.762) (2.043,-0.764) (2.078,-0.766) (2.114,-0.767) (2.149,-0.768) (2.185,-0.770) (2.220,-0.771) (2.255,-0.772) (2.291,-0.773) (2.326,-0.774) (2.362,-0.775) (2.397,-0.776) (2.433,-0.777) (2.468,-0.778) (2.504,-0.779) (2.539,-0.779) (2.575,-0.780) (2.610,-0.781) (2.645,-0.782) (2.681,-0.782) (2.716,-0.783) (2.752,-0.783) (2.787,-0.784) (2.823,-0.785) (2.858,-0.785) (2.894,-0.786) (2.929,-0.786) (2.965,-0.786) (3.000,-0.787)};
\addplot[lineaE, flechaen=0.387] coordinates {(-3.000,-0.787) (-2.965,-0.786) (-2.929,-0.786) (-2.894,-0.786) (-2.858,-0.785) (-2.823,-0.785) (-2.787,-0.784) (-2.752,-0.783) (-2.716,-0.783) (-2.681,-0.782) (-2.645,-0.782) (-2.610,-0.781) (-2.575,-0.780) (-2.539,-0.779) (-2.504,-0.779) (-2.468,-0.778) (-2.433,-0.777) (-2.397,-0.776) (-2.362,-0.775) (-2.326,-0.774) (-2.291,-0.773) (-2.255,-0.772) (-2.220,-0.771) (-2.185,-0.770) (-2.149,-0.768) (-2.114,-0.767) (-2.078,-0.766) (-2.043,-0.764) (-2.007,-0.762) (-1.972,-0.761) (-1.937,-0.759) (-1.901,-0.757) (-1.866,-0.755) (-1.830,-0.753) (-1.795,-0.751) (-1.760,-0.748) (-1.724,-0.746) (-1.689,-0.743) (-1.653,-0.740) (-1.618,-0.737) (-1.583,-0.734) (-1.548,-0.730) (-1.512,-0.727) (-1.477,-0.723) (-1.442,-0.719) (-1.407,-0.714) (-1.371,-0.709) (-1.336,-0.704) (-1.301,-0.699) (-1.266,-0.693) (-1.231,-0.687) (-1.197,-0.680) (-1.162,-0.673) (-1.127,-0.665) (-1.093,-0.657) (-1.058,-0.648) (-1.024,-0.639) (-0.990,-0.629) (-0.957,-0.618) (-0.923,-0.606) (-0.890,-0.593) (-0.857,-0.580) (-0.825,-0.565)};
\addplot[lineaE, flechaen=0.635] coordinates {(0.708,-0.707) (0.738,-0.725) (0.769,-0.741) (0.801,-0.756) (0.833,-0.771) (0.866,-0.784) (0.899,-0.796) (0.932,-0.808) (0.966,-0.818) (1.000,-0.828) (1.034,-0.838) (1.068,-0.847) (1.102,-0.855) (1.137,-0.862) (1.171,-0.870) (1.206,-0.876) (1.240,-0.883) (1.275,-0.889) (1.310,-0.894) (1.345,-0.899) (1.380,-0.904) (1.415,-0.909) (1.450,-0.913) (1.485,-0.917) (1.520,-0.921) (1.555,-0.925) (1.590,-0.928) (1.625,-0.932) (1.661,-0.935) (1.696,-0.938) (1.731,-0.940) (1.766,-0.943) (1.801,-0.945) (1.836,-0.948) (1.872,-0.950) (1.907,-0.952) (1.942,-0.954) (1.977,-0.956) (2.013,-0.958) (2.048,-0.959) (2.083,-0.961) (2.118,-0.963) (2.154,-0.964) (2.189,-0.965) (2.224,-0.967) (2.259,-0.968) (2.295,-0.969) (2.330,-0.970) (2.365,-0.971) (2.400,-0.972) (2.436,-0.973) (2.471,-0.974) (2.506,-0.975) (2.541,-0.976) (2.577,-0.977) (2.612,-0.978) (2.647,-0.978) (2.683,-0.979) (2.718,-0.980) (2.753,-0.981) (2.788,-0.981) (2.824,-0.982) (2.859,-0.982) (2.894,-0.983) (2.929,-0.983) (2.965,-0.984) (3.000,-0.984)};
\addplot[lineaE, flechaen=0.365] coordinates {(-3.000,-0.984) (-2.965,-0.984) (-2.929,-0.983) (-2.894,-0.983) (-2.859,-0.982) (-2.824,-0.982) (-2.788,-0.981) (-2.753,-0.981) (-2.718,-0.980) (-2.683,-0.979) (-2.647,-0.978) (-2.612,-0.978) (-2.577,-0.977) (-2.541,-0.976) (-2.506,-0.975) (-2.471,-0.974) (-2.436,-0.973) (-2.400,-0.972) (-2.365,-0.971) (-2.330,-0.970) (-2.295,-0.969) (-2.259,-0.968) (-2.224,-0.967) (-2.189,-0.965) (-2.154,-0.964) (-2.118,-0.963) (-2.083,-0.961) (-2.048,-0.959) (-2.013,-0.958) (-1.977,-0.956) (-1.942,-0.954) (-1.907,-0.952) (-1.872,-0.950) (-1.836,-0.948) (-1.801,-0.945) (-1.766,-0.943) (-1.731,-0.940) (-1.696,-0.938) (-1.661,-0.935) (-1.625,-0.932) (-1.590,-0.928) (-1.555,-0.925) (-1.520,-0.921) (-1.485,-0.917) (-1.450,-0.913) (-1.415,-0.909) (-1.380,-0.904) (-1.345,-0.899) (-1.310,-0.894) (-1.275,-0.889) (-1.240,-0.883) (-1.206,-0.876) (-1.171,-0.870) (-1.137,-0.862) (-1.102,-0.855) (-1.068,-0.847) (-1.034,-0.838) (-1.000,-0.828) (-0.966,-0.818) (-0.932,-0.808) (-0.899,-0.796) (-0.866,-0.784) (-0.833,-0.771) (-0.801,-0.756) (-0.769,-0.741) (-0.738,-0.725) (-0.708,-0.707)};
\addplot[lineaE, flechaen=0.662] coordinates {(0.531,-0.848) (0.559,-0.869) (0.589,-0.888) (0.620,-0.905) (0.651,-0.921) (0.683,-0.936) (0.716,-0.951) (0.749,-0.964) (0.782,-0.976) (0.816,-0.988) (0.849,-0.999) (0.883,-1.009) (0.917,-1.018) (0.952,-1.027) (0.986,-1.036) (1.021,-1.044) (1.055,-1.051) (1.090,-1.059) (1.125,-1.065) (1.160,-1.072) (1.195,-1.078) (1.230,-1.083) (1.265,-1.089) (1.300,-1.094) (1.335,-1.098) (1.371,-1.103) (1.406,-1.107) (1.441,-1.111) (1.476,-1.115) (1.512,-1.119) (1.547,-1.122) (1.582,-1.126) (1.618,-1.129) (1.653,-1.132) (1.688,-1.134) (1.724,-1.137) (1.759,-1.140) (1.794,-1.142) (1.830,-1.144) (1.865,-1.147) (1.901,-1.149) (1.936,-1.151) (1.972,-1.153) (2.007,-1.154) (2.042,-1.156) (2.078,-1.158) (2.113,-1.159) (2.149,-1.161) (2.184,-1.162) (2.220,-1.164) (2.255,-1.165) (2.291,-1.166) (2.326,-1.167) (2.361,-1.168) (2.397,-1.170) (2.432,-1.171) (2.468,-1.172) (2.503,-1.173) (2.539,-1.173) (2.574,-1.174) (2.610,-1.175) (2.645,-1.176) (2.681,-1.177) (2.716,-1.178) (2.752,-1.178) (2.787,-1.179) (2.823,-1.180) (2.858,-1.180) (2.894,-1.181) (2.929,-1.181) (2.965,-1.182) (3.000,-1.182)};
\addplot[lineaE, flechaen=0.338] coordinates {(-3.000,-1.182) (-2.965,-1.182) (-2.929,-1.181) (-2.894,-1.181) (-2.858,-1.180) (-2.823,-1.180) (-2.787,-1.179) (-2.752,-1.178) (-2.716,-1.178) (-2.681,-1.177) (-2.645,-1.176) (-2.610,-1.175) (-2.574,-1.174) (-2.539,-1.173) (-2.503,-1.173) (-2.468,-1.172) (-2.432,-1.171) (-2.397,-1.170) (-2.361,-1.168) (-2.326,-1.167) (-2.291,-1.166) (-2.255,-1.165) (-2.220,-1.164) (-2.184,-1.162) (-2.149,-1.161) (-2.113,-1.159) (-2.078,-1.158) (-2.042,-1.156) (-2.007,-1.154) (-1.972,-1.153) (-1.936,-1.151) (-1.901,-1.149) (-1.865,-1.147) (-1.830,-1.144) (-1.794,-1.142) (-1.759,-1.140) (-1.724,-1.137) (-1.688,-1.134) (-1.653,-1.132) (-1.618,-1.129) (-1.582,-1.126) (-1.547,-1.122) (-1.512,-1.119) (-1.476,-1.115) (-1.441,-1.111) (-1.406,-1.107) (-1.371,-1.103) (-1.335,-1.098) (-1.300,-1.094) (-1.265,-1.089) (-1.230,-1.083) (-1.195,-1.078) (-1.160,-1.072) (-1.125,-1.065) (-1.090,-1.059) (-1.055,-1.051) (-1.021,-1.044) (-0.986,-1.036) (-0.952,-1.027) (-0.917,-1.018) (-0.883,-1.009) (-0.849,-0.999) (-0.816,-0.988) (-0.782,-0.976) (-0.749,-0.964) (-0.716,-0.951) (-0.683,-0.936) (-0.651,-0.921) (-0.620,-0.905) (-0.589,-0.888) (-0.559,-0.869) (-0.531,-0.848)};
\addplot[lineaE, flechaen=0.707] coordinates {(0.148,-0.989) (0.178,-1.005) (0.210,-1.019) (0.243,-1.033) (0.275,-1.047) (0.308,-1.061) (0.340,-1.074) (0.373,-1.087) (0.406,-1.100) (0.440,-1.112) (0.473,-1.124) (0.507,-1.135) (0.540,-1.146) (0.574,-1.157) (0.608,-1.167) (0.642,-1.177) (0.676,-1.186) (0.710,-1.195) (0.745,-1.204) (0.779,-1.212) (0.814,-1.220) (0.848,-1.227) (0.883,-1.235) (0.918,-1.241) (0.953,-1.248) (0.987,-1.254) (1.022,-1.260) (1.057,-1.266) (1.092,-1.271) (1.127,-1.277) (1.162,-1.282) (1.197,-1.287) (1.232,-1.291) (1.268,-1.295) (1.303,-1.300) (1.338,-1.304) (1.373,-1.307) (1.408,-1.311) (1.444,-1.314) (1.479,-1.318) (1.514,-1.321) (1.549,-1.324) (1.585,-1.327) (1.620,-1.330) (1.655,-1.332) (1.691,-1.335) (1.726,-1.337) (1.761,-1.340) (1.797,-1.342) (1.832,-1.344) (1.867,-1.346) (1.903,-1.348) (1.938,-1.350) (1.974,-1.351) (2.009,-1.353) (2.044,-1.355) (2.080,-1.356) (2.115,-1.358) (2.150,-1.359) (2.186,-1.361) (2.221,-1.362) (2.257,-1.363) (2.292,-1.365) (2.327,-1.366) (2.363,-1.367) (2.398,-1.368) (2.434,-1.369) (2.469,-1.370) (2.504,-1.371) (2.540,-1.372) (2.575,-1.373) (2.611,-1.373) (2.646,-1.374) (2.681,-1.375) (2.717,-1.376) (2.752,-1.377) (2.788,-1.377) (2.823,-1.378) (2.858,-1.379) (2.894,-1.379) (2.929,-1.380) (2.965,-1.380) (3.000,-1.381)};
\addplot[lineaE, flechaen=0.293] coordinates {(-3.000,-1.381) (-2.965,-1.380) (-2.929,-1.380) (-2.894,-1.379) (-2.858,-1.379) (-2.823,-1.378) (-2.788,-1.377) (-2.752,-1.377) (-2.717,-1.376) (-2.681,-1.375) (-2.646,-1.374) (-2.611,-1.373) (-2.575,-1.373) (-2.540,-1.372) (-2.504,-1.371) (-2.469,-1.370) (-2.434,-1.369) (-2.398,-1.368) (-2.363,-1.367) (-2.327,-1.366) (-2.292,-1.365) (-2.257,-1.363) (-2.221,-1.362) (-2.186,-1.361) (-2.150,-1.359) (-2.115,-1.358) (-2.080,-1.356) (-2.044,-1.355) (-2.009,-1.353) (-1.974,-1.351) (-1.938,-1.350) (-1.903,-1.348) (-1.867,-1.346) (-1.832,-1.344) (-1.797,-1.342) (-1.761,-1.340) (-1.726,-1.337) (-1.691,-1.335) (-1.655,-1.332) (-1.620,-1.330) (-1.585,-1.327) (-1.549,-1.324) (-1.514,-1.321) (-1.479,-1.318) (-1.444,-1.314) (-1.408,-1.311) (-1.373,-1.307) (-1.338,-1.304) (-1.303,-1.300) (-1.268,-1.295) (-1.232,-1.291) (-1.197,-1.287) (-1.162,-1.282) (-1.127,-1.277) (-1.092,-1.271) (-1.057,-1.266) (-1.022,-1.260) (-0.987,-1.254) (-0.953,-1.248) (-0.918,-1.241) (-0.883,-1.235) (-0.848,-1.227) (-0.814,-1.220) (-0.779,-1.212) (-0.745,-1.204) (-0.710,-1.195) (-0.676,-1.186) (-0.642,-1.177) (-0.608,-1.167) (-0.574,-1.157) (-0.540,-1.146) (-0.507,-1.135) (-0.473,-1.124) (-0.440,-1.112) (-0.406,-1.100) (-0.373,-1.087) (-0.340,-1.074) (-0.308,-1.061) (-0.275,-1.047) (-0.243,-1.033) (-0.210,-1.019) (-0.178,-1.005) (-0.148,-0.989)};
\addplot[lineaE, flechaen=0.141, flechaen=0.859] coordinates {(-3.000,-1.580) (-2.965,-1.579) (-2.930,-1.579) (-2.895,-1.578) (-2.860,-1.578) (-2.825,-1.577) (-2.790,-1.576) (-2.755,-1.576) (-2.720,-1.575) (-2.685,-1.574) (-2.650,-1.573) (-2.615,-1.573) (-2.580,-1.572) (-2.545,-1.571) (-2.510,-1.570) (-2.475,-1.569) (-2.440,-1.568) (-2.405,-1.567) (-2.370,-1.566) (-2.335,-1.565) (-2.300,-1.564) (-2.265,-1.563) (-2.230,-1.562) (-2.195,-1.561) (-2.160,-1.559) (-2.125,-1.558) (-2.090,-1.557) (-2.055,-1.555) (-2.020,-1.554) (-1.985,-1.552) (-1.950,-1.551) (-1.915,-1.549) (-1.880,-1.547) (-1.845,-1.546) (-1.810,-1.544) (-1.775,-1.542) (-1.740,-1.540) (-1.705,-1.538) (-1.670,-1.536) (-1.635,-1.533) (-1.600,-1.531) (-1.565,-1.529) (-1.530,-1.526) (-1.495,-1.524) (-1.460,-1.521) (-1.425,-1.518) (-1.390,-1.515) (-1.355,-1.512) (-1.320,-1.509) (-1.286,-1.506) (-1.251,-1.503) (-1.216,-1.499) (-1.181,-1.496) (-1.146,-1.492) (-1.111,-1.488) (-1.077,-1.484) (-1.042,-1.480) (-1.007,-1.476) (-0.972,-1.471) (-0.937,-1.467) (-0.903,-1.462) (-0.868,-1.458) (-0.833,-1.453) (-0.799,-1.448) (-0.764,-1.443) (-0.729,-1.438) (-0.695,-1.432) (-0.660,-1.427) (-0.626,-1.421) (-0.591,-1.416) (-0.556,-1.410) (-0.522,-1.405) (-0.487,-1.399) (-0.453,-1.394) (-0.418,-1.388) (-0.383,-1.383) (-0.349,-1.378) (-0.314,-1.373) (-0.279,-1.368) (-0.245,-1.364) (-0.210,-1.360) (-0.175,-1.357) (-0.140,-1.354) (-0.105,-1.352) (-0.070,-1.350) (-0.035,-1.349) (0.000,-1.348) (0.035,-1.349) (0.070,-1.350) (0.105,-1.352) (0.140,-1.354) (0.175,-1.357) (0.210,-1.360) (0.245,-1.364) (0.279,-1.368) (0.314,-1.373) (0.349,-1.378) (0.383,-1.383) (0.418,-1.388) (0.453,-1.394) (0.487,-1.399) (0.522,-1.405) (0.556,-1.410) (0.591,-1.416) (0.626,-1.421) (0.660,-1.427) (0.695,-1.432) (0.729,-1.438) (0.764,-1.443) (0.799,-1.448) (0.833,-1.453) (0.868,-1.458) (0.903,-1.462) (0.937,-1.467) (0.972,-1.471) (1.007,-1.476) (1.042,-1.480) (1.077,-1.484) (1.111,-1.488) (1.146,-1.492) (1.181,-1.496) (1.216,-1.499) (1.251,-1.503) (1.286,-1.506) (1.320,-1.509) (1.355,-1.512) (1.390,-1.515) (1.425,-1.518) (1.460,-1.521) (1.495,-1.524) (1.530,-1.526) (1.565,-1.529) (1.600,-1.531) (1.635,-1.533) (1.670,-1.536) (1.705,-1.538) (1.740,-1.540) (1.775,-1.542) (1.810,-1.544) (1.845,-1.546) (1.880,-1.547) (1.915,-1.549) (1.950,-1.551) (1.985,-1.552) (2.020,-1.554) (2.055,-1.555) (2.090,-1.557) (2.125,-1.558) (2.160,-1.559) (2.195,-1.561) (2.230,-1.562) (2.265,-1.563) (2.300,-1.564) (2.335,-1.565) (2.370,-1.566) (2.405,-1.567) (2.440,-1.568) (2.475,-1.569) (2.510,-1.570) (2.545,-1.571) (2.580,-1.572) (2.615,-1.573) (2.650,-1.573) (2.685,-1.574) (2.720,-1.575) (2.755,-1.576) (2.790,-1.576) (2.825,-1.577) (2.860,-1.578) (2.895,-1.578) (2.930,-1.579) (2.965,-1.579) (3.000,-1.580)};
\addplot[lineaE, flechaen=0.141, flechaen=0.859] coordinates {(-3.000,-1.779) (-2.965,-1.779) (-2.930,-1.778) (-2.895,-1.777) (-2.859,-1.777) (-2.824,-1.776) (-2.789,-1.776) (-2.754,-1.775) (-2.719,-1.774) (-2.684,-1.774) (-2.649,-1.773) (-2.613,-1.772) (-2.578,-1.771) (-2.543,-1.770) (-2.508,-1.770) (-2.473,-1.769) (-2.438,-1.768) (-2.403,-1.767) (-2.367,-1.766) (-2.332,-1.765) (-2.297,-1.764) (-2.262,-1.763) (-2.227,-1.762) (-2.192,-1.761) (-2.157,-1.760) (-2.121,-1.758) (-2.086,-1.757) (-2.051,-1.756) (-2.016,-1.755) (-1.981,-1.753) (-1.946,-1.752) (-1.911,-1.750) (-1.876,-1.749) (-1.841,-1.747) (-1.805,-1.746) (-1.770,-1.744) (-1.735,-1.742) (-1.700,-1.741) (-1.665,-1.739) (-1.630,-1.737) (-1.595,-1.735) (-1.560,-1.733) (-1.525,-1.731) (-1.489,-1.729) (-1.454,-1.727) (-1.419,-1.724) (-1.384,-1.722) (-1.349,-1.720) (-1.314,-1.717) (-1.279,-1.715) (-1.244,-1.712) (-1.209,-1.709) (-1.174,-1.707) (-1.139,-1.704) (-1.104,-1.701) (-1.069,-1.698) (-1.034,-1.695) (-0.999,-1.692) (-0.964,-1.689) (-0.929,-1.686) (-0.894,-1.682) (-0.859,-1.679) (-0.824,-1.676) (-0.789,-1.673) (-0.754,-1.669) (-0.719,-1.666) (-0.684,-1.663) (-0.649,-1.659) (-0.614,-1.656) (-0.579,-1.653) (-0.544,-1.649) (-0.509,-1.646) (-0.474,-1.643) (-0.439,-1.640) (-0.404,-1.637) (-0.369,-1.635) (-0.334,-1.632) (-0.299,-1.630) (-0.263,-1.628) (-0.228,-1.626) (-0.193,-1.624) (-0.158,-1.622) (-0.123,-1.621) (-0.088,-1.620) (-0.053,-1.620) (-0.018,-1.619) (0.018,-1.619) (0.053,-1.620) (0.088,-1.620) (0.123,-1.621) (0.158,-1.622) (0.193,-1.624) (0.228,-1.626) (0.263,-1.628) (0.299,-1.630) (0.334,-1.632) (0.369,-1.635) (0.404,-1.637) (0.439,-1.640) (0.474,-1.643) (0.509,-1.646) (0.544,-1.649) (0.579,-1.653) (0.614,-1.656) (0.649,-1.659) (0.684,-1.663) (0.719,-1.666) (0.754,-1.669) (0.789,-1.673) (0.824,-1.676) (0.859,-1.679) (0.894,-1.682) (0.929,-1.686) (0.964,-1.689) (0.999,-1.692) (1.034,-1.695) (1.069,-1.698) (1.104,-1.701) (1.139,-1.704) (1.174,-1.707) (1.209,-1.709) (1.244,-1.712) (1.279,-1.715) (1.314,-1.717) (1.349,-1.720) (1.384,-1.722) (1.419,-1.724) (1.454,-1.727) (1.489,-1.729) (1.525,-1.731) (1.560,-1.733) (1.595,-1.735) (1.630,-1.737) (1.665,-1.739) (1.700,-1.741) (1.735,-1.742) (1.770,-1.744) (1.805,-1.746) (1.841,-1.747) (1.876,-1.749) (1.911,-1.750) (1.946,-1.752) (1.981,-1.753) (2.016,-1.755) (2.051,-1.756) (2.086,-1.757) (2.121,-1.758) (2.157,-1.760) (2.192,-1.761) (2.227,-1.762) (2.262,-1.763) (2.297,-1.764) (2.332,-1.765) (2.367,-1.766) (2.403,-1.767) (2.438,-1.768) (2.473,-1.769) (2.508,-1.770) (2.543,-1.770) (2.578,-1.771) (2.613,-1.772) (2.649,-1.773) (2.684,-1.774) (2.719,-1.774) (2.754,-1.775) (2.789,-1.776) (2.824,-1.776) (2.859,-1.777) (2.895,-1.777) (2.930,-1.778) (2.965,-1.779) (3.000,-1.779)};
\addplot[lineaE, flechaen=0.142, flechaen=0.858] coordinates {(-3.000,-1.979) (-2.965,-1.978) (-2.930,-1.978) (-2.895,-1.977) (-2.860,-1.977) (-2.824,-1.976) (-2.789,-1.975) (-2.754,-1.975) (-2.719,-1.974) (-2.684,-1.973) (-2.649,-1.973) (-2.614,-1.972) (-2.579,-1.971) (-2.544,-1.971) (-2.508,-1.970) (-2.473,-1.969) (-2.438,-1.968) (-2.403,-1.967) (-2.368,-1.967) (-2.333,-1.966) (-2.298,-1.965) (-2.263,-1.964) (-2.228,-1.963) (-2.192,-1.962) (-2.157,-1.961) (-2.122,-1.960) (-2.087,-1.959) (-2.052,-1.958) (-2.017,-1.956) (-1.982,-1.955) (-1.947,-1.954) (-1.912,-1.953) (-1.877,-1.951) (-1.841,-1.950) (-1.806,-1.949) (-1.771,-1.947) (-1.736,-1.946) (-1.701,-1.944) (-1.666,-1.943) (-1.631,-1.941) (-1.596,-1.940) (-1.561,-1.938) (-1.526,-1.936) (-1.491,-1.935) (-1.455,-1.933) (-1.420,-1.931) (-1.385,-1.929) (-1.350,-1.927) (-1.315,-1.926) (-1.280,-1.924) (-1.245,-1.922) (-1.210,-1.920) (-1.175,-1.917) (-1.140,-1.915) (-1.105,-1.913) (-1.070,-1.911) (-1.035,-1.909) (-1.000,-1.907) (-0.965,-1.905) (-0.930,-1.902) (-0.895,-1.900) (-0.859,-1.898) (-0.824,-1.896) (-0.789,-1.893) (-0.754,-1.891) (-0.719,-1.889) (-0.684,-1.887) (-0.649,-1.885) (-0.614,-1.882) (-0.579,-1.880) (-0.544,-1.878) (-0.509,-1.876) (-0.474,-1.875) (-0.439,-1.873) (-0.404,-1.871) (-0.369,-1.869) (-0.334,-1.868) (-0.298,-1.867) (-0.263,-1.865) (-0.228,-1.864) (-0.193,-1.863) (-0.158,-1.862) (-0.123,-1.862) (-0.088,-1.861) (-0.053,-1.861) (-0.018,-1.861) (0.018,-1.861) (0.053,-1.861) (0.088,-1.861) (0.123,-1.862) (0.158,-1.862) (0.193,-1.863) (0.228,-1.864) (0.263,-1.865) (0.298,-1.867) (0.334,-1.868) (0.369,-1.869) (0.404,-1.871) (0.439,-1.873) (0.474,-1.875) (0.509,-1.876) (0.544,-1.878) (0.579,-1.880) (0.614,-1.882) (0.649,-1.885) (0.684,-1.887) (0.719,-1.889) (0.754,-1.891) (0.789,-1.893) (0.824,-1.896) (0.859,-1.898) (0.895,-1.900) (0.930,-1.902) (0.965,-1.905) (1.000,-1.907) (1.035,-1.909) (1.070,-1.911) (1.105,-1.913) (1.140,-1.915) (1.175,-1.917) (1.210,-1.920) (1.245,-1.922) (1.280,-1.924) (1.315,-1.926) (1.350,-1.927) (1.385,-1.929) (1.420,-1.931) (1.455,-1.933) (1.491,-1.935) (1.526,-1.936) (1.561,-1.938) (1.596,-1.940) (1.631,-1.941) (1.666,-1.943) (1.701,-1.944) (1.736,-1.946) (1.771,-1.947) (1.806,-1.949) (1.841,-1.950) (1.877,-1.951) (1.912,-1.953) (1.947,-1.954) (1.982,-1.955) (2.017,-1.956) (2.052,-1.958) (2.087,-1.959) (2.122,-1.960) (2.157,-1.961) (2.192,-1.962) (2.228,-1.963) (2.263,-1.964) (2.298,-1.965) (2.333,-1.966) (2.368,-1.967) (2.403,-1.967) (2.438,-1.968) (2.473,-1.969) (2.508,-1.970) (2.544,-1.971) (2.579,-1.971) (2.614,-1.972) (2.649,-1.973) (2.684,-1.973) (2.719,-1.974) (2.754,-1.975) (2.789,-1.975) (2.824,-1.976) (2.860,-1.977) (2.895,-1.977) (2.930,-1.978) (2.965,-1.978) (3.000,-1.979)};
\addplot[lineaE, flechaen=0.142, flechaen=0.858] coordinates {(-3.000,-2.179) (-2.965,-2.178) (-2.930,-2.178) (-2.895,-2.177) (-2.860,-2.177) (-2.824,-2.176) (-2.789,-2.176) (-2.754,-2.175) (-2.719,-2.174) (-2.684,-2.174) (-2.649,-2.173) (-2.614,-2.173) (-2.579,-2.172) (-2.544,-2.171) (-2.509,-2.170) (-2.473,-2.170) (-2.438,-2.169) (-2.403,-2.168) (-2.368,-2.167) (-2.333,-2.167) (-2.298,-2.166) (-2.263,-2.165) (-2.228,-2.164) (-2.193,-2.163) (-2.158,-2.162) (-2.123,-2.161) (-2.087,-2.161) (-2.052,-2.160) (-2.017,-2.159) (-1.982,-2.158) (-1.947,-2.156) (-1.912,-2.155) (-1.877,-2.154) (-1.842,-2.153) (-1.807,-2.152) (-1.772,-2.151) (-1.737,-2.150) (-1.701,-2.148) (-1.666,-2.147) (-1.631,-2.146) (-1.596,-2.145) (-1.561,-2.143) (-1.526,-2.142) (-1.491,-2.141) (-1.456,-2.139) (-1.421,-2.138) (-1.386,-2.136) (-1.351,-2.135) (-1.316,-2.133) (-1.281,-2.132) (-1.245,-2.130) (-1.210,-2.129) (-1.175,-2.127) (-1.140,-2.126) (-1.105,-2.124) (-1.070,-2.122) (-1.035,-2.121) (-1.000,-2.119) (-0.965,-2.118) (-0.930,-2.116) (-0.895,-2.114) (-0.860,-2.113) (-0.825,-2.111) (-0.790,-2.110) (-0.754,-2.108) (-0.719,-2.107) (-0.684,-2.105) (-0.649,-2.104) (-0.614,-2.102) (-0.579,-2.101) (-0.544,-2.100) (-0.509,-2.098) (-0.474,-2.097) (-0.439,-2.096) (-0.404,-2.095) (-0.369,-2.094) (-0.333,-2.093) (-0.298,-2.092) (-0.263,-2.091) (-0.228,-2.090) (-0.193,-2.090) (-0.158,-2.089) (-0.123,-2.089) (-0.088,-2.089) (-0.053,-2.088) (-0.018,-2.088) (0.018,-2.088) (0.053,-2.088) (0.088,-2.089) (0.123,-2.089) (0.158,-2.089) (0.193,-2.090) (0.228,-2.090) (0.263,-2.091) (0.298,-2.092) (0.333,-2.093) (0.369,-2.094) (0.404,-2.095) (0.439,-2.096) (0.474,-2.097) (0.509,-2.098) (0.544,-2.100) (0.579,-2.101) (0.614,-2.102) (0.649,-2.104) (0.684,-2.105) (0.719,-2.107) (0.754,-2.108) (0.790,-2.110) (0.825,-2.111) (0.860,-2.113) (0.895,-2.114) (0.930,-2.116) (0.965,-2.118) (1.000,-2.119) (1.035,-2.121) (1.070,-2.122) (1.105,-2.124) (1.140,-2.126) (1.175,-2.127) (1.210,-2.129) (1.245,-2.130) (1.281,-2.132) (1.316,-2.133) (1.351,-2.135) (1.386,-2.136) (1.421,-2.138) (1.456,-2.139) (1.491,-2.141) (1.526,-2.142) (1.561,-2.143) (1.596,-2.145) (1.631,-2.146) (1.666,-2.147) (1.701,-2.148) (1.737,-2.150) (1.772,-2.151) (1.807,-2.152) (1.842,-2.153) (1.877,-2.154) (1.912,-2.155) (1.947,-2.156) (1.982,-2.158) (2.017,-2.159) (2.052,-2.160) (2.087,-2.161) (2.123,-2.161) (2.158,-2.162) (2.193,-2.163) (2.228,-2.164) (2.263,-2.165) (2.298,-2.166) (2.333,-2.167) (2.368,-2.167) (2.403,-2.168) (2.438,-2.169) (2.473,-2.170) (2.509,-2.170) (2.544,-2.171) (2.579,-2.172) (2.614,-2.173) (2.649,-2.173) (2.684,-2.174) (2.719,-2.174) (2.754,-2.175) (2.789,-2.176) (2.824,-2.176) (2.860,-2.177) (2.895,-2.177) (2.930,-2.178) (2.965,-2.178) (3.000,-2.179)};
\addplot[lineaE, flechaen=0.239] coordinates {(-0.959,-0.283) (-0.924,-0.283) (-0.888,-0.283) (-0.853,-0.283) (-0.817,-0.283) (-0.781,-0.283) (-0.746,-0.283) (-0.710,-0.283) (-0.675,-0.283) (-0.639,-0.283) (-0.604,-0.283) (-0.568,-0.283) (-0.533,-0.283) (-0.497,-0.283) (-0.462,-0.283) (-0.426,-0.283) (-0.391,-0.283) (-0.355,-0.283) (-0.320,-0.283) (-0.284,-0.283) (-0.249,-0.283) (-0.213,-0.283) (-0.178,-0.283) (-0.142,-0.283) (-0.107,-0.283) (-0.071,-0.283) (-0.036,-0.283) (-0.000,-0.283) (0.036,-0.283) (0.071,-0.283) (0.107,-0.283) (0.142,-0.283) (0.178,-0.283) (0.213,-0.283) (0.249,-0.283) (0.284,-0.283) (0.320,-0.283) (0.355,-0.283) (0.391,-0.283) (0.426,-0.283) (0.462,-0.283) (0.497,-0.283) (0.533,-0.283) (0.568,-0.283) (0.604,-0.283) (0.639,-0.283) (0.675,-0.283) (0.710,-0.283) (0.746,-0.283) (0.781,-0.283) (0.817,-0.283) (0.853,-0.283) (0.888,-0.283) (0.924,-0.283) (0.959,-0.283)};
\addplot[lineaE, flechaen=0.804] coordinates {(-0.824,-0.567) (-0.789,-0.566) (-0.754,-0.566) (-0.719,-0.566) (-0.684,-0.566) (-0.649,-0.566) (-0.614,-0.566) (-0.579,-0.566) (-0.544,-0.566) (-0.508,-0.566) (-0.473,-0.566) (-0.438,-0.566) (-0.403,-0.566) (-0.368,-0.566) (-0.333,-0.566) (-0.298,-0.566) (-0.263,-0.566) (-0.228,-0.566) (-0.193,-0.566) (-0.158,-0.566) (-0.123,-0.566) (-0.088,-0.566) (-0.053,-0.566) (-0.018,-0.566) (0.018,-0.566) (0.053,-0.566) (0.088,-0.566) (0.123,-0.566) (0.158,-0.566) (0.193,-0.566) (0.228,-0.566) (0.263,-0.566) (0.298,-0.566) (0.333,-0.566) (0.368,-0.566) (0.403,-0.566) (0.438,-0.566) (0.473,-0.566) (0.508,-0.566) (0.544,-0.566) (0.579,-0.566) (0.614,-0.566) (0.649,-0.566) (0.684,-0.566) (0.719,-0.566) (0.754,-0.566) (0.789,-0.566) (0.824,-0.567)};
\addplot[lineaE, flechaen=0.500] coordinates {(-0.525,-0.851) (-0.491,-0.849) (-0.456,-0.849) (-0.421,-0.849) (-0.386,-0.849) (-0.351,-0.849) (-0.316,-0.849) (-0.281,-0.849) (-0.246,-0.849) (-0.210,-0.849) (-0.175,-0.849) (-0.140,-0.849) (-0.105,-0.849) (-0.070,-0.849) (-0.035,-0.849) (0.000,-0.849) (0.035,-0.849) (0.070,-0.849) (0.105,-0.849) (0.140,-0.849) (0.175,-0.849) (0.210,-0.849) (0.246,-0.849) (0.281,-0.849) (0.316,-0.849) (0.351,-0.849) (0.386,-0.849) (0.421,-0.849) (0.456,-0.849) (0.491,-0.849) (0.525,-0.851)};
\addplot[lineaE, flechaen=0.142, flechaen=0.858] coordinates {(-3.000,0.000) (3.000,0.000)};

  \draw[texto, line width=0.6pt] (axis cs:0,0) circle[radius=1];
  \pgfplotsinvokeforeach{-60,-30,0,30,60}{
    \node[cpos, font=\scriptsize\bfseries, fill=fondo, circle, inner sep=0.2pt] at (axis cs:{1.16*cos(#1)},{1.16*sin(#1)}) {$+$};
    \node[cneg, font=\scriptsize\bfseries, fill=fondo, circle, inner sep=0.2pt] at (axis cs:{-1.16*cos(#1)},{1.16*sin(#1)}) {$-$};
  }
\end{axis}
\end{tikzpicture}
\caption{Líneas de $\vE$}
\end{subfigure}\hfill
\begin{subfigure}[b]{0.49\textwidth}\centering
\begin{tikzpicture}
\begin{axis}[esfera]
  \fill[dielec] (axis cs:0,0) circle[radius=1];
  %% Generado por scripts/lineas_campo.py -- no editar a mano
\addplot[equipot] coordinates {(2.781,-2.200) (2.782,-2.165) (2.782,-2.130) (2.783,-2.094) (2.784,-2.059) (2.784,-2.024) (2.785,-1.989) (2.785,-1.954) (2.786,-1.918) (2.787,-1.883) (2.787,-1.848) (2.788,-1.813) (2.789,-1.778) (2.789,-1.742) (2.790,-1.707) (2.791,-1.672) (2.791,-1.637) (2.792,-1.602) (2.793,-1.566) (2.793,-1.531) (2.794,-1.496) (2.795,-1.461) (2.795,-1.426) (2.796,-1.390) (2.797,-1.355) (2.797,-1.320) (2.798,-1.285) (2.799,-1.250) (2.799,-1.214) (2.800,-1.179) (2.801,-1.144) (2.801,-1.109) (2.802,-1.074) (2.802,-1.038) (2.803,-1.003) (2.804,-0.968) (2.804,-0.933) (2.805,-0.898) (2.805,-0.862) (2.806,-0.827) (2.807,-0.792) (2.807,-0.757) (2.808,-0.722) (2.808,-0.686) (2.809,-0.651) (2.809,-0.616) (2.809,-0.581) (2.810,-0.546) (2.810,-0.510) (2.811,-0.475) (2.811,-0.440) (2.811,-0.405) (2.812,-0.370) (2.812,-0.334) (2.812,-0.299) (2.812,-0.264) (2.813,-0.229) (2.813,-0.194) (2.813,-0.158) (2.813,-0.123) (2.813,-0.088) (2.813,-0.053) (2.813,-0.018) (2.813,0.018) (2.813,0.053) (2.813,0.088) (2.813,0.123) (2.813,0.158) (2.813,0.194) (2.813,0.229) (2.812,0.264) (2.812,0.299) (2.812,0.334) (2.812,0.370) (2.811,0.405) (2.811,0.440) (2.811,0.475) (2.810,0.510) (2.810,0.546) (2.809,0.581) (2.809,0.616) (2.809,0.651) (2.808,0.686) (2.808,0.722) (2.807,0.757) (2.807,0.792) (2.806,0.827) (2.805,0.862) (2.805,0.898) (2.804,0.933) (2.804,0.968) (2.803,1.003) (2.802,1.038) (2.802,1.074) (2.801,1.109) (2.801,1.144) (2.800,1.179) (2.799,1.214) (2.799,1.250) (2.798,1.285) (2.797,1.320) (2.797,1.355) (2.796,1.390) (2.795,1.426) (2.795,1.461) (2.794,1.496) (2.793,1.531) (2.793,1.566) (2.792,1.602) (2.791,1.637) (2.791,1.672) (2.790,1.707) (2.789,1.742) (2.789,1.778) (2.788,1.813) (2.787,1.848) (2.787,1.883) (2.786,1.918) (2.785,1.954) (2.785,1.989) (2.784,2.024) (2.784,2.059) (2.783,2.094) (2.782,2.130) (2.782,2.165) (2.781,2.200)};
\addplot[equipot] coordinates {(2.286,-2.200) (2.287,-2.165) (2.287,-2.130) (2.288,-2.094) (2.289,-2.059) (2.290,-2.024) (2.291,-1.989) (2.292,-1.954) (2.293,-1.918) (2.294,-1.883) (2.295,-1.848) (2.296,-1.813) (2.297,-1.778) (2.298,-1.742) (2.299,-1.707) (2.300,-1.672) (2.301,-1.637) (2.302,-1.602) (2.303,-1.566) (2.304,-1.531) (2.306,-1.496) (2.307,-1.461) (2.308,-1.426) (2.309,-1.390) (2.310,-1.355) (2.311,-1.320) (2.312,-1.285) (2.314,-1.250) (2.315,-1.214) (2.316,-1.179) (2.317,-1.144) (2.318,-1.109) (2.319,-1.074) (2.321,-1.039) (2.322,-1.003) (2.323,-0.968) (2.324,-0.933) (2.325,-0.898) (2.326,-0.863) (2.327,-0.827) (2.328,-0.792) (2.329,-0.757) (2.330,-0.722) (2.331,-0.687) (2.332,-0.651) (2.333,-0.616) (2.334,-0.581) (2.335,-0.546) (2.335,-0.511) (2.336,-0.475) (2.337,-0.440) (2.338,-0.405) (2.338,-0.370) (2.339,-0.335) (2.339,-0.299) (2.340,-0.264) (2.340,-0.229) (2.340,-0.194) (2.341,-0.158) (2.341,-0.123) (2.341,-0.088) (2.341,-0.053) (2.341,-0.018) (2.341,0.018) (2.341,0.053) (2.341,0.088) (2.341,0.123) (2.341,0.158) (2.340,0.194) (2.340,0.229) (2.340,0.264) (2.339,0.299) (2.339,0.335) (2.338,0.370) (2.338,0.405) (2.337,0.440) (2.336,0.475) (2.335,0.511) (2.335,0.546) (2.334,0.581) (2.333,0.616) (2.332,0.651) (2.331,0.687) (2.330,0.722) (2.329,0.757) (2.328,0.792) (2.327,0.827) (2.326,0.863) (2.325,0.898) (2.324,0.933) (2.323,0.968) (2.322,1.003) (2.321,1.039) (2.319,1.074) (2.318,1.109) (2.317,1.144) (2.316,1.179) (2.315,1.214) (2.314,1.250) (2.312,1.285) (2.311,1.320) (2.310,1.355) (2.309,1.390) (2.308,1.426) (2.307,1.461) (2.306,1.496) (2.304,1.531) (2.303,1.566) (2.302,1.602) (2.301,1.637) (2.300,1.672) (2.299,1.707) (2.298,1.742) (2.297,1.778) (2.296,1.813) (2.295,1.848) (2.294,1.883) (2.293,1.918) (2.292,1.954) (2.291,1.989) (2.290,2.024) (2.289,2.059) (2.288,2.094) (2.287,2.130) (2.287,2.165) (2.286,2.200)};
\addplot[equipot] coordinates {(1.789,-2.200) (1.790,-2.165) (1.792,-2.130) (1.793,-2.094) (1.794,-2.059) (1.795,-2.024) (1.797,-1.989) (1.798,-1.953) (1.799,-1.918) (1.801,-1.883) (1.802,-1.848) (1.804,-1.813) (1.806,-1.777) (1.807,-1.742) (1.809,-1.707) (1.810,-1.672) (1.812,-1.637) (1.814,-1.601) (1.816,-1.566) (1.818,-1.531) (1.820,-1.496) (1.822,-1.461) (1.824,-1.425) (1.826,-1.390) (1.828,-1.355) (1.830,-1.320) (1.832,-1.285) (1.834,-1.250) (1.836,-1.214) (1.838,-1.179) (1.840,-1.144) (1.843,-1.109) (1.845,-1.074) (1.847,-1.039) (1.849,-1.003) (1.851,-0.968) (1.854,-0.933) (1.856,-0.898) (1.858,-0.863) (1.860,-0.827) (1.862,-0.792) (1.864,-0.757) (1.866,-0.722) (1.868,-0.687) (1.870,-0.652) (1.872,-0.616) (1.874,-0.581) (1.876,-0.546) (1.877,-0.511) (1.879,-0.476) (1.881,-0.440) (1.882,-0.405) (1.883,-0.370) (1.884,-0.335) (1.885,-0.300) (1.886,-0.264) (1.887,-0.229) (1.888,-0.194) (1.889,-0.159) (1.889,-0.123) (1.890,-0.088) (1.890,-0.053) (1.890,-0.018) (1.890,0.018) (1.890,0.053) (1.890,0.088) (1.889,0.123) (1.889,0.159) (1.888,0.194) (1.887,0.229) (1.886,0.264) (1.885,0.300) (1.884,0.335) (1.883,0.370) (1.882,0.405) (1.881,0.440) (1.879,0.476) (1.877,0.511) (1.876,0.546) (1.874,0.581) (1.872,0.616) (1.870,0.652) (1.868,0.687) (1.866,0.722) (1.864,0.757) (1.862,0.792) (1.860,0.827) (1.858,0.863) (1.856,0.898) (1.854,0.933) (1.851,0.968) (1.849,1.003) (1.847,1.039) (1.845,1.074) (1.843,1.109) (1.840,1.144) (1.838,1.179) (1.836,1.214) (1.834,1.250) (1.832,1.285) (1.830,1.320) (1.828,1.355) (1.826,1.390) (1.824,1.425) (1.822,1.461) (1.820,1.496) (1.818,1.531) (1.816,1.566) (1.814,1.601) (1.812,1.637) (1.810,1.672) (1.809,1.707) (1.807,1.742) (1.806,1.777) (1.804,1.813) (1.802,1.848) (1.801,1.883) (1.799,1.918) (1.798,1.953) (1.797,1.989) (1.795,2.024) (1.794,2.059) (1.793,2.094) (1.792,2.130) (1.790,2.165) (1.789,2.200)};
\addplot[equipot] coordinates {(1.289,-2.200) (1.290,-2.165) (1.292,-2.130) (1.293,-2.095) (1.295,-2.060) (1.297,-2.025) (1.298,-1.990) (1.300,-1.955) (1.302,-1.920) (1.304,-1.885) (1.306,-1.850) (1.308,-1.815) (1.311,-1.780) (1.313,-1.745) (1.316,-1.710) (1.318,-1.675) (1.321,-1.640) (1.324,-1.605) (1.326,-1.570) (1.329,-1.535) (1.332,-1.500) (1.336,-1.465) (1.339,-1.430) (1.342,-1.395) (1.346,-1.360) (1.350,-1.326) (1.353,-1.291) (1.357,-1.256) (1.361,-1.221) (1.365,-1.186) (1.370,-1.151) (1.374,-1.116) (1.378,-1.082) (1.383,-1.047) (1.387,-1.012) (1.392,-0.977) (1.396,-0.943) (1.401,-0.908) (1.405,-0.873) (1.410,-0.838) (1.414,-0.803) (1.419,-0.769) (1.423,-0.734) (1.428,-0.699) (1.432,-0.664) (1.436,-0.629) (1.440,-0.595) (1.444,-0.560) (1.448,-0.525) (1.452,-0.490) (1.455,-0.455) (1.459,-0.420) (1.462,-0.385) (1.464,-0.350) (1.467,-0.315) (1.469,-0.280) (1.472,-0.245) (1.473,-0.210) (1.475,-0.175) (1.476,-0.140) (1.477,-0.105) (1.478,-0.070) (1.479,-0.035) (1.479,0.000) (1.479,0.035) (1.478,0.070) (1.477,0.105) (1.476,0.140) (1.475,0.175) (1.473,0.210) (1.472,0.245) (1.469,0.280) (1.467,0.315) (1.464,0.350) (1.462,0.385) (1.459,0.420) (1.455,0.455) (1.452,0.490) (1.448,0.525) (1.444,0.560) (1.440,0.595) (1.436,0.629) (1.432,0.664) (1.428,0.699) (1.423,0.734) (1.419,0.769) (1.414,0.803) (1.410,0.838) (1.405,0.873) (1.401,0.908) (1.396,0.943) (1.392,0.977) (1.387,1.012) (1.383,1.047) (1.378,1.082) (1.374,1.116) (1.370,1.151) (1.365,1.186) (1.361,1.221) (1.357,1.256) (1.353,1.291) (1.350,1.326) (1.346,1.360) (1.342,1.395) (1.339,1.430) (1.336,1.465) (1.332,1.500) (1.329,1.535) (1.326,1.570) (1.324,1.605) (1.321,1.640) (1.318,1.675) (1.316,1.710) (1.313,1.745) (1.311,1.780) (1.308,1.815) (1.306,1.850) (1.304,1.885) (1.302,1.920) (1.300,1.955) (1.298,1.990) (1.297,2.025) (1.295,2.060) (1.293,2.095) (1.292,2.130) (1.290,2.165) (1.289,2.200)};
\addplot[equipot] coordinates {(0.781,-2.200) (0.782,-2.165) (0.784,-2.130) (0.785,-2.094) (0.787,-2.059) (0.788,-2.024) (0.790,-1.989) (0.792,-1.954) (0.794,-1.918) (0.797,-1.883) (0.799,-1.848) (0.801,-1.813) (0.804,-1.778) (0.807,-1.743) (0.810,-1.708) (0.813,-1.673) (0.817,-1.638) (0.820,-1.602) (0.824,-1.567) (0.828,-1.532) (0.833,-1.498) (0.837,-1.463) (0.842,-1.428) (0.848,-1.393) (0.853,-1.358) (0.859,-1.323) (0.866,-1.289) (0.872,-1.254) (0.879,-1.220) (0.887,-1.185) (0.894,-1.151) (0.903,-1.117) (0.911,-1.082) (0.920,-1.048) (0.929,-1.014) (0.938,-0.980) (0.947,-0.946) (0.957,-0.912) (0.967,-0.878) (0.977,-0.845) (0.987,-0.811) (0.997,-0.777) (1.007,-0.743) (1.017,-0.710) (1.027,-0.676) (1.036,-0.642) (1.046,-0.608) (1.055,-0.574) (1.063,-0.540) (1.072,-0.505) (1.080,-0.471) (1.088,-0.437) (1.095,-0.402) (1.102,-0.368) (1.108,-0.333) (1.113,-0.298) (1.119,-0.263) (1.123,-0.228) (1.127,-0.193) (1.130,-0.158) (1.133,-0.123) (1.135,-0.088) (1.136,-0.053) (1.137,-0.018) (1.137,0.018) (1.136,0.053) (1.135,0.088) (1.133,0.123) (1.130,0.158) (1.127,0.193) (1.123,0.228) (1.119,0.263) (1.113,0.298) (1.108,0.333) (1.102,0.368) (1.095,0.402) (1.088,0.437) (1.080,0.471) (1.072,0.505) (1.063,0.540) (1.055,0.574) (1.046,0.608) (1.036,0.642) (1.027,0.676) (1.017,0.710) (1.007,0.743) (0.997,0.777) (0.987,0.811) (0.977,0.845) (0.967,0.878) (0.957,0.912) (0.947,0.946) (0.938,0.980) (0.929,1.014) (0.920,1.048) (0.911,1.082) (0.903,1.117) (0.894,1.151) (0.887,1.185) (0.879,1.220) (0.872,1.254) (0.866,1.289) (0.859,1.323) (0.853,1.358) (0.848,1.393) (0.842,1.428) (0.837,1.463) (0.833,1.498) (0.828,1.532) (0.824,1.567) (0.820,1.602) (0.817,1.638) (0.813,1.673) (0.810,1.708) (0.807,1.743) (0.804,1.778) (0.801,1.813) (0.799,1.848) (0.797,1.883) (0.794,1.918) (0.792,1.954) (0.790,1.989) (0.788,2.024) (0.787,2.059) (0.785,2.094) (0.784,2.130) (0.782,2.165) (0.781,2.200)};
\addplot[equipot] coordinates {(0.262,-2.200) (0.263,-2.165) (0.263,-2.130) (0.264,-2.095) (0.265,-2.060) (0.266,-2.025) (0.266,-1.990) (0.267,-1.955) (0.268,-1.919) (0.270,-1.884) (0.271,-1.849) (0.272,-1.814) (0.273,-1.779) (0.275,-1.744) (0.277,-1.709) (0.278,-1.674) (0.280,-1.639) (0.283,-1.604) (0.285,-1.569) (0.288,-1.534) (0.291,-1.499) (0.294,-1.464) (0.298,-1.429) (0.302,-1.395) (0.307,-1.360) (0.312,-1.325) (0.318,-1.291) (0.324,-1.256) (0.332,-1.222) (0.340,-1.188) (0.350,-1.154) (0.360,-1.121) (0.373,-1.088) (0.386,-1.055) (0.401,-1.024) (0.417,-0.993) (0.435,-0.962) (0.453,-0.933) (0.473,-0.904) (0.493,-0.875) (0.500,-0.842) (0.500,-0.807) (0.500,-0.772) (0.500,-0.737) (0.500,-0.701) (0.500,-0.666) (0.500,-0.631) (0.500,-0.596) (0.500,-0.561) (0.500,-0.526) (0.500,-0.491) (0.500,-0.456) (0.500,-0.421) (0.500,-0.386) (0.500,-0.351) (0.500,-0.316) (0.500,-0.281) (0.500,-0.246) (0.500,-0.210) (0.500,-0.175) (0.500,-0.140) (0.500,-0.105) (0.500,-0.070) (0.500,-0.035) (0.500,0.000) (0.500,0.035) (0.500,0.070) (0.500,0.105) (0.500,0.140) (0.500,0.175) (0.500,0.210) (0.500,0.246) (0.500,0.281) (0.500,0.316) (0.500,0.351) (0.500,0.386) (0.500,0.421) (0.500,0.456) (0.500,0.491) (0.500,0.526) (0.500,0.561) (0.500,0.596) (0.500,0.631) (0.500,0.666) (0.500,0.701) (0.500,0.737) (0.500,0.772) (0.500,0.807) (0.500,0.842) (0.493,0.875) (0.473,0.904) (0.453,0.933) (0.435,0.962) (0.417,0.993) (0.401,1.024) (0.386,1.055) (0.373,1.088) (0.360,1.121) (0.350,1.154) (0.340,1.188) (0.332,1.222) (0.324,1.256) (0.318,1.291) (0.312,1.325) (0.307,1.360) (0.302,1.395) (0.298,1.429) (0.294,1.464) (0.291,1.499) (0.288,1.534) (0.285,1.569) (0.283,1.604) (0.280,1.639) (0.278,1.674) (0.277,1.709) (0.275,1.744) (0.273,1.779) (0.272,1.814) (0.271,1.849) (0.270,1.884) (0.268,1.919) (0.267,1.955) (0.266,1.990) (0.266,2.025) (0.265,2.060) (0.264,2.095) (0.263,2.130) (0.263,2.165) (0.262,2.200)};
\addplot[equipot] coordinates {(-0.262,-2.200) (-0.263,-2.165) (-0.263,-2.130) (-0.264,-2.095) (-0.265,-2.060) (-0.266,-2.025) (-0.266,-1.990) (-0.267,-1.955) (-0.268,-1.919) (-0.270,-1.884) (-0.271,-1.849) (-0.272,-1.814) (-0.273,-1.779) (-0.275,-1.744) (-0.277,-1.709) (-0.278,-1.674) (-0.280,-1.639) (-0.283,-1.604) (-0.285,-1.569) (-0.288,-1.534) (-0.291,-1.499) (-0.294,-1.464) (-0.298,-1.429) (-0.302,-1.395) (-0.307,-1.360) (-0.312,-1.325) (-0.318,-1.291) (-0.324,-1.256) (-0.332,-1.222) (-0.340,-1.188) (-0.350,-1.154) (-0.360,-1.121) (-0.373,-1.088) (-0.386,-1.055) (-0.401,-1.024) (-0.417,-0.993) (-0.435,-0.962) (-0.453,-0.933) (-0.473,-0.904) (-0.493,-0.875) (-0.500,-0.842) (-0.500,-0.807) (-0.500,-0.772) (-0.500,-0.737) (-0.500,-0.701) (-0.500,-0.666) (-0.500,-0.631) (-0.500,-0.596) (-0.500,-0.561) (-0.500,-0.526) (-0.500,-0.491) (-0.500,-0.456) (-0.500,-0.421) (-0.500,-0.386) (-0.500,-0.351) (-0.500,-0.316) (-0.500,-0.281) (-0.500,-0.246) (-0.500,-0.210) (-0.500,-0.175) (-0.500,-0.140) (-0.500,-0.105) (-0.500,-0.070) (-0.500,-0.035) (-0.500,-0.000) (-0.500,0.035) (-0.500,0.070) (-0.500,0.105) (-0.500,0.140) (-0.500,0.175) (-0.500,0.210) (-0.500,0.246) (-0.500,0.281) (-0.500,0.316) (-0.500,0.351) (-0.500,0.386) (-0.500,0.421) (-0.500,0.456) (-0.500,0.491) (-0.500,0.526) (-0.500,0.561) (-0.500,0.596) (-0.500,0.631) (-0.500,0.666) (-0.500,0.701) (-0.500,0.737) (-0.500,0.772) (-0.500,0.807) (-0.500,0.842) (-0.493,0.875) (-0.473,0.904) (-0.453,0.933) (-0.435,0.962) (-0.417,0.993) (-0.401,1.024) (-0.386,1.055) (-0.373,1.088) (-0.360,1.121) (-0.350,1.154) (-0.340,1.188) (-0.332,1.222) (-0.324,1.256) (-0.318,1.291) (-0.312,1.325) (-0.307,1.360) (-0.302,1.395) (-0.298,1.429) (-0.294,1.464) (-0.291,1.499) (-0.288,1.534) (-0.285,1.569) (-0.283,1.604) (-0.280,1.639) (-0.278,1.674) (-0.277,1.709) (-0.275,1.744) (-0.273,1.779) (-0.272,1.814) (-0.271,1.849) (-0.270,1.884) (-0.268,1.919) (-0.267,1.955) (-0.266,1.990) (-0.266,2.025) (-0.265,2.060) (-0.264,2.095) (-0.263,2.130) (-0.263,2.165) (-0.262,2.200)};
\addplot[equipot] coordinates {(-0.781,-2.200) (-0.782,-2.165) (-0.784,-2.130) (-0.785,-2.094) (-0.787,-2.059) (-0.788,-2.024) (-0.790,-1.989) (-0.792,-1.954) (-0.794,-1.918) (-0.797,-1.883) (-0.799,-1.848) (-0.801,-1.813) (-0.804,-1.778) (-0.807,-1.743) (-0.810,-1.708) (-0.813,-1.673) (-0.817,-1.638) (-0.820,-1.602) (-0.824,-1.567) (-0.828,-1.532) (-0.833,-1.498) (-0.837,-1.463) (-0.842,-1.428) (-0.848,-1.393) (-0.853,-1.358) (-0.859,-1.323) (-0.866,-1.289) (-0.872,-1.254) (-0.879,-1.220) (-0.887,-1.185) (-0.894,-1.151) (-0.903,-1.117) (-0.911,-1.082) (-0.920,-1.048) (-0.929,-1.014) (-0.938,-0.980) (-0.947,-0.946) (-0.957,-0.912) (-0.967,-0.878) (-0.977,-0.845) (-0.987,-0.811) (-0.997,-0.777) (-1.007,-0.743) (-1.017,-0.710) (-1.027,-0.676) (-1.036,-0.642) (-1.046,-0.608) (-1.055,-0.574) (-1.063,-0.540) (-1.072,-0.505) (-1.080,-0.471) (-1.088,-0.437) (-1.095,-0.402) (-1.102,-0.368) (-1.108,-0.333) (-1.113,-0.298) (-1.119,-0.263) (-1.123,-0.228) (-1.127,-0.193) (-1.130,-0.158) (-1.133,-0.123) (-1.135,-0.088) (-1.136,-0.053) (-1.137,-0.018) (-1.137,0.018) (-1.136,0.053) (-1.135,0.088) (-1.133,0.123) (-1.130,0.158) (-1.127,0.193) (-1.123,0.228) (-1.119,0.263) (-1.113,0.298) (-1.108,0.333) (-1.102,0.368) (-1.095,0.402) (-1.088,0.437) (-1.080,0.471) (-1.072,0.505) (-1.063,0.540) (-1.055,0.574) (-1.046,0.608) (-1.036,0.642) (-1.027,0.676) (-1.017,0.710) (-1.007,0.743) (-0.997,0.777) (-0.987,0.811) (-0.977,0.845) (-0.967,0.878) (-0.957,0.912) (-0.947,0.946) (-0.938,0.980) (-0.929,1.014) (-0.920,1.048) (-0.911,1.082) (-0.903,1.117) (-0.894,1.151) (-0.887,1.185) (-0.879,1.220) (-0.872,1.254) (-0.866,1.289) (-0.859,1.323) (-0.853,1.358) (-0.848,1.393) (-0.842,1.428) (-0.837,1.463) (-0.833,1.498) (-0.828,1.532) (-0.824,1.567) (-0.820,1.602) (-0.817,1.638) (-0.813,1.673) (-0.810,1.708) (-0.807,1.743) (-0.804,1.778) (-0.801,1.813) (-0.799,1.848) (-0.797,1.883) (-0.794,1.918) (-0.792,1.954) (-0.790,1.989) (-0.788,2.024) (-0.787,2.059) (-0.785,2.094) (-0.784,2.130) (-0.782,2.165) (-0.781,2.200)};
\addplot[equipot] coordinates {(-1.289,-2.200) (-1.290,-2.165) (-1.292,-2.130) (-1.293,-2.095) (-1.295,-2.060) (-1.297,-2.025) (-1.298,-1.990) (-1.300,-1.955) (-1.302,-1.920) (-1.304,-1.885) (-1.306,-1.850) (-1.308,-1.815) (-1.311,-1.780) (-1.313,-1.745) (-1.316,-1.710) (-1.318,-1.675) (-1.321,-1.640) (-1.324,-1.605) (-1.326,-1.570) (-1.329,-1.535) (-1.332,-1.500) (-1.336,-1.465) (-1.339,-1.430) (-1.342,-1.395) (-1.346,-1.360) (-1.350,-1.326) (-1.353,-1.291) (-1.357,-1.256) (-1.361,-1.221) (-1.365,-1.186) (-1.370,-1.151) (-1.374,-1.116) (-1.378,-1.082) (-1.383,-1.047) (-1.387,-1.012) (-1.392,-0.977) (-1.396,-0.943) (-1.401,-0.908) (-1.405,-0.873) (-1.410,-0.838) (-1.414,-0.803) (-1.419,-0.769) (-1.423,-0.734) (-1.428,-0.699) (-1.432,-0.664) (-1.436,-0.629) (-1.440,-0.595) (-1.444,-0.560) (-1.448,-0.525) (-1.452,-0.490) (-1.455,-0.455) (-1.459,-0.420) (-1.462,-0.385) (-1.464,-0.350) (-1.467,-0.315) (-1.469,-0.280) (-1.472,-0.245) (-1.473,-0.210) (-1.475,-0.175) (-1.476,-0.140) (-1.477,-0.105) (-1.478,-0.070) (-1.479,-0.035) (-1.479,0.000) (-1.479,0.035) (-1.478,0.070) (-1.477,0.105) (-1.476,0.140) (-1.475,0.175) (-1.473,0.210) (-1.472,0.245) (-1.469,0.280) (-1.467,0.315) (-1.464,0.350) (-1.462,0.385) (-1.459,0.420) (-1.455,0.455) (-1.452,0.490) (-1.448,0.525) (-1.444,0.560) (-1.440,0.595) (-1.436,0.629) (-1.432,0.664) (-1.428,0.699) (-1.423,0.734) (-1.419,0.769) (-1.414,0.803) (-1.410,0.838) (-1.405,0.873) (-1.401,0.908) (-1.396,0.943) (-1.392,0.977) (-1.387,1.012) (-1.383,1.047) (-1.378,1.082) (-1.374,1.116) (-1.370,1.151) (-1.365,1.186) (-1.361,1.221) (-1.357,1.256) (-1.353,1.291) (-1.350,1.326) (-1.346,1.360) (-1.342,1.395) (-1.339,1.430) (-1.336,1.465) (-1.332,1.500) (-1.329,1.535) (-1.326,1.570) (-1.324,1.605) (-1.321,1.640) (-1.318,1.675) (-1.316,1.710) (-1.313,1.745) (-1.311,1.780) (-1.308,1.815) (-1.306,1.850) (-1.304,1.885) (-1.302,1.920) (-1.300,1.955) (-1.298,1.990) (-1.297,2.025) (-1.295,2.060) (-1.293,2.095) (-1.292,2.130) (-1.290,2.165) (-1.289,2.200)};
\addplot[equipot] coordinates {(-1.789,-2.200) (-1.790,-2.165) (-1.792,-2.130) (-1.793,-2.094) (-1.794,-2.059) (-1.795,-2.024) (-1.797,-1.989) (-1.798,-1.953) (-1.799,-1.918) (-1.801,-1.883) (-1.802,-1.848) (-1.804,-1.813) (-1.806,-1.777) (-1.807,-1.742) (-1.809,-1.707) (-1.810,-1.672) (-1.812,-1.637) (-1.814,-1.601) (-1.816,-1.566) (-1.818,-1.531) (-1.820,-1.496) (-1.822,-1.461) (-1.824,-1.425) (-1.826,-1.390) (-1.828,-1.355) (-1.830,-1.320) (-1.832,-1.285) (-1.834,-1.250) (-1.836,-1.214) (-1.838,-1.179) (-1.840,-1.144) (-1.843,-1.109) (-1.845,-1.074) (-1.847,-1.039) (-1.849,-1.003) (-1.851,-0.968) (-1.854,-0.933) (-1.856,-0.898) (-1.858,-0.863) (-1.860,-0.827) (-1.862,-0.792) (-1.864,-0.757) (-1.866,-0.722) (-1.868,-0.687) (-1.870,-0.652) (-1.872,-0.616) (-1.874,-0.581) (-1.876,-0.546) (-1.877,-0.511) (-1.879,-0.476) (-1.881,-0.440) (-1.882,-0.405) (-1.883,-0.370) (-1.884,-0.335) (-1.885,-0.300) (-1.886,-0.264) (-1.887,-0.229) (-1.888,-0.194) (-1.889,-0.159) (-1.889,-0.123) (-1.890,-0.088) (-1.890,-0.053) (-1.890,-0.018) (-1.890,0.018) (-1.890,0.053) (-1.890,0.088) (-1.889,0.123) (-1.889,0.159) (-1.888,0.194) (-1.887,0.229) (-1.886,0.264) (-1.885,0.300) (-1.884,0.335) (-1.883,0.370) (-1.882,0.405) (-1.881,0.440) (-1.879,0.476) (-1.877,0.511) (-1.876,0.546) (-1.874,0.581) (-1.872,0.616) (-1.870,0.652) (-1.868,0.687) (-1.866,0.722) (-1.864,0.757) (-1.862,0.792) (-1.860,0.827) (-1.858,0.863) (-1.856,0.898) (-1.854,0.933) (-1.851,0.968) (-1.849,1.003) (-1.847,1.039) (-1.845,1.074) (-1.843,1.109) (-1.840,1.144) (-1.838,1.179) (-1.836,1.214) (-1.834,1.250) (-1.832,1.285) (-1.830,1.320) (-1.828,1.355) (-1.826,1.390) (-1.824,1.425) (-1.822,1.461) (-1.820,1.496) (-1.818,1.531) (-1.816,1.566) (-1.814,1.601) (-1.812,1.637) (-1.810,1.672) (-1.809,1.707) (-1.807,1.742) (-1.806,1.777) (-1.804,1.813) (-1.802,1.848) (-1.801,1.883) (-1.799,1.918) (-1.798,1.953) (-1.797,1.989) (-1.795,2.024) (-1.794,2.059) (-1.793,2.094) (-1.792,2.130) (-1.790,2.165) (-1.789,2.200)};
\addplot[equipot] coordinates {(-2.286,-2.200) (-2.287,-2.165) (-2.287,-2.130) (-2.288,-2.094) (-2.289,-2.059) (-2.290,-2.024) (-2.291,-1.989) (-2.292,-1.954) (-2.293,-1.918) (-2.294,-1.883) (-2.295,-1.848) (-2.296,-1.813) (-2.297,-1.778) (-2.298,-1.742) (-2.299,-1.707) (-2.300,-1.672) (-2.301,-1.637) (-2.302,-1.602) (-2.303,-1.566) (-2.304,-1.531) (-2.306,-1.496) (-2.307,-1.461) (-2.308,-1.426) (-2.309,-1.390) (-2.310,-1.355) (-2.311,-1.320) (-2.312,-1.285) (-2.314,-1.250) (-2.315,-1.214) (-2.316,-1.179) (-2.317,-1.144) (-2.318,-1.109) (-2.319,-1.074) (-2.321,-1.039) (-2.322,-1.003) (-2.323,-0.968) (-2.324,-0.933) (-2.325,-0.898) (-2.326,-0.863) (-2.327,-0.827) (-2.328,-0.792) (-2.329,-0.757) (-2.330,-0.722) (-2.331,-0.687) (-2.332,-0.651) (-2.333,-0.616) (-2.334,-0.581) (-2.335,-0.546) (-2.335,-0.511) (-2.336,-0.475) (-2.337,-0.440) (-2.338,-0.405) (-2.338,-0.370) (-2.339,-0.335) (-2.339,-0.299) (-2.340,-0.264) (-2.340,-0.229) (-2.340,-0.194) (-2.341,-0.158) (-2.341,-0.123) (-2.341,-0.088) (-2.341,-0.053) (-2.341,-0.018) (-2.341,0.018) (-2.341,0.053) (-2.341,0.088) (-2.341,0.123) (-2.341,0.158) (-2.340,0.194) (-2.340,0.229) (-2.340,0.264) (-2.339,0.299) (-2.339,0.335) (-2.338,0.370) (-2.338,0.405) (-2.337,0.440) (-2.336,0.475) (-2.335,0.511) (-2.335,0.546) (-2.334,0.581) (-2.333,0.616) (-2.332,0.651) (-2.331,0.687) (-2.330,0.722) (-2.329,0.757) (-2.328,0.792) (-2.327,0.827) (-2.326,0.863) (-2.325,0.898) (-2.324,0.933) (-2.323,0.968) (-2.322,1.003) (-2.321,1.039) (-2.319,1.074) (-2.318,1.109) (-2.317,1.144) (-2.316,1.179) (-2.315,1.214) (-2.314,1.250) (-2.312,1.285) (-2.311,1.320) (-2.310,1.355) (-2.309,1.390) (-2.308,1.426) (-2.307,1.461) (-2.306,1.496) (-2.304,1.531) (-2.303,1.566) (-2.302,1.602) (-2.301,1.637) (-2.300,1.672) (-2.299,1.707) (-2.298,1.742) (-2.297,1.778) (-2.296,1.813) (-2.295,1.848) (-2.294,1.883) (-2.293,1.918) (-2.292,1.954) (-2.291,1.989) (-2.290,2.024) (-2.289,2.059) (-2.288,2.094) (-2.287,2.130) (-2.287,2.165) (-2.286,2.200)};
\addplot[equipot] coordinates {(-2.781,-2.200) (-2.782,-2.165) (-2.782,-2.130) (-2.783,-2.094) (-2.784,-2.059) (-2.784,-2.024) (-2.785,-1.989) (-2.785,-1.954) (-2.786,-1.918) (-2.787,-1.883) (-2.787,-1.848) (-2.788,-1.813) (-2.789,-1.778) (-2.789,-1.742) (-2.790,-1.707) (-2.791,-1.672) (-2.791,-1.637) (-2.792,-1.602) (-2.793,-1.566) (-2.793,-1.531) (-2.794,-1.496) (-2.795,-1.461) (-2.795,-1.426) (-2.796,-1.390) (-2.797,-1.355) (-2.797,-1.320) (-2.798,-1.285) (-2.799,-1.250) (-2.799,-1.214) (-2.800,-1.179) (-2.801,-1.144) (-2.801,-1.109) (-2.802,-1.074) (-2.802,-1.038) (-2.803,-1.003) (-2.804,-0.968) (-2.804,-0.933) (-2.805,-0.898) (-2.805,-0.862) (-2.806,-0.827) (-2.807,-0.792) (-2.807,-0.757) (-2.808,-0.722) (-2.808,-0.686) (-2.809,-0.651) (-2.809,-0.616) (-2.809,-0.581) (-2.810,-0.546) (-2.810,-0.510) (-2.811,-0.475) (-2.811,-0.440) (-2.811,-0.405) (-2.812,-0.370) (-2.812,-0.334) (-2.812,-0.299) (-2.812,-0.264) (-2.813,-0.229) (-2.813,-0.194) (-2.813,-0.158) (-2.813,-0.123) (-2.813,-0.088) (-2.813,-0.053) (-2.813,-0.018) (-2.813,0.018) (-2.813,0.053) (-2.813,0.088) (-2.813,0.123) (-2.813,0.158) (-2.813,0.194) (-2.813,0.229) (-2.812,0.264) (-2.812,0.299) (-2.812,0.334) (-2.812,0.370) (-2.811,0.405) (-2.811,0.440) (-2.811,0.475) (-2.810,0.510) (-2.810,0.546) (-2.809,0.581) (-2.809,0.616) (-2.809,0.651) (-2.808,0.686) (-2.808,0.722) (-2.807,0.757) (-2.807,0.792) (-2.806,0.827) (-2.805,0.862) (-2.805,0.898) (-2.804,0.933) (-2.804,0.968) (-2.803,1.003) (-2.802,1.038) (-2.802,1.074) (-2.801,1.109) (-2.801,1.144) (-2.800,1.179) (-2.799,1.214) (-2.799,1.250) (-2.798,1.285) (-2.797,1.320) (-2.797,1.355) (-2.796,1.390) (-2.795,1.426) (-2.795,1.461) (-2.794,1.496) (-2.793,1.531) (-2.793,1.566) (-2.792,1.602) (-2.791,1.637) (-2.791,1.672) (-2.790,1.707) (-2.789,1.742) (-2.789,1.778) (-2.788,1.813) (-2.787,1.848) (-2.787,1.883) (-2.786,1.918) (-2.785,1.954) (-2.785,1.989) (-2.784,2.024) (-2.784,2.059) (-2.783,2.094) (-2.782,2.130) (-2.782,2.165) (-2.781,2.200)};
\addplot[lineaD, flechaen=0.142, flechaen=0.858] coordinates {(-3.000,0.196) (-2.965,0.196) (-2.930,0.196) (-2.895,0.196) (-2.860,0.196) (-2.824,0.196) (-2.789,0.196) (-2.754,0.195) (-2.719,0.195) (-2.684,0.195) (-2.649,0.195) (-2.614,0.195) (-2.579,0.194) (-2.544,0.194) (-2.509,0.194) (-2.473,0.194) (-2.438,0.193) (-2.403,0.193) (-2.368,0.193) (-2.333,0.193) (-2.298,0.192) (-2.263,0.192) (-2.228,0.192) (-2.193,0.191) (-2.158,0.191) (-2.122,0.190) (-2.087,0.190) (-2.052,0.189) (-2.017,0.189) (-1.982,0.188) (-1.947,0.188) (-1.912,0.187) (-1.877,0.187) (-1.842,0.186) (-1.806,0.185) (-1.771,0.184) (-1.736,0.183) (-1.701,0.183) (-1.666,0.182) (-1.631,0.181) (-1.596,0.179) (-1.561,0.178) (-1.526,0.177) (-1.491,0.176) (-1.456,0.174) (-1.421,0.173) (-1.385,0.171) (-1.350,0.169) (-1.315,0.167) (-1.280,0.165) (-1.245,0.163) (-1.210,0.161) (-1.175,0.158) (-1.140,0.155) (-1.105,0.152) (-1.070,0.149) (-1.035,0.146) (-1.000,0.143) (-0.965,0.141) (-0.930,0.141) (-0.895,0.141) (-0.860,0.141) (-0.825,0.141) (-0.790,0.141) (-0.755,0.141) (-0.720,0.141) (-0.685,0.141) (-0.649,0.141) (-0.614,0.141) (-0.579,0.141) (-0.544,0.141) (-0.509,0.141) (-0.474,0.141) (-0.439,0.141) (-0.404,0.141) (-0.369,0.141) (-0.333,0.141) (-0.298,0.141) (-0.263,0.141) (-0.228,0.141) (-0.193,0.141) (-0.158,0.141) (-0.123,0.141) (-0.088,0.141) (-0.053,0.141) (-0.018,0.141) (0.018,0.141) (0.053,0.141) (0.088,0.141) (0.123,0.141) (0.158,0.141) (0.193,0.141) (0.228,0.141) (0.263,0.141) (0.298,0.141) (0.333,0.141) (0.369,0.141) (0.404,0.141) (0.439,0.141) (0.474,0.141) (0.509,0.141) (0.544,0.141) (0.579,0.141) (0.614,0.141) (0.649,0.141) (0.685,0.141) (0.720,0.141) (0.755,0.141) (0.790,0.141) (0.825,0.141) (0.860,0.141) (0.895,0.141) (0.930,0.141) (0.965,0.141) (1.000,0.143) (1.035,0.146) (1.070,0.149) (1.105,0.152) (1.140,0.155) (1.175,0.158) (1.210,0.161) (1.245,0.163) (1.280,0.165) (1.315,0.167) (1.350,0.169) (1.385,0.171) (1.421,0.173) (1.456,0.174) (1.491,0.176) (1.526,0.177) (1.561,0.178) (1.596,0.179) (1.631,0.181) (1.666,0.182) (1.701,0.183) (1.736,0.183) (1.771,0.184) (1.806,0.185) (1.842,0.186) (1.877,0.187) (1.912,0.187) (1.947,0.188) (1.982,0.188) (2.017,0.189) (2.052,0.189) (2.087,0.190) (2.122,0.190) (2.158,0.191) (2.193,0.191) (2.228,0.192) (2.263,0.192) (2.298,0.192) (2.333,0.193) (2.368,0.193) (2.403,0.193) (2.438,0.193) (2.473,0.194) (2.509,0.194) (2.544,0.194) (2.579,0.194) (2.614,0.195) (2.649,0.195) (2.684,0.195) (2.719,0.195) (2.754,0.195) (2.789,0.196) (2.824,0.196) (2.860,0.196) (2.895,0.196) (2.930,0.196) (2.965,0.196) (3.000,0.196)};
\addplot[lineaD, flechaen=0.141, flechaen=0.859] coordinates {(-3.000,0.393) (-2.965,0.393) (-2.930,0.392) (-2.895,0.392) (-2.859,0.392) (-2.824,0.392) (-2.789,0.391) (-2.754,0.391) (-2.719,0.391) (-2.684,0.390) (-2.648,0.390) (-2.613,0.390) (-2.578,0.389) (-2.543,0.389) (-2.508,0.388) (-2.473,0.388) (-2.438,0.387) (-2.402,0.387) (-2.367,0.386) (-2.332,0.386) (-2.297,0.385) (-2.262,0.384) (-2.227,0.384) (-2.191,0.383) (-2.156,0.382) (-2.121,0.381) (-2.086,0.381) (-2.051,0.380) (-2.016,0.379) (-1.981,0.378) (-1.945,0.377) (-1.910,0.375) (-1.875,0.374) (-1.840,0.373) (-1.805,0.371) (-1.770,0.370) (-1.735,0.368) (-1.700,0.367) (-1.664,0.365) (-1.629,0.363) (-1.594,0.361) (-1.559,0.359) (-1.524,0.356) (-1.489,0.354) (-1.454,0.351) (-1.419,0.348) (-1.384,0.345) (-1.349,0.342) (-1.314,0.338) (-1.279,0.334) (-1.244,0.330) (-1.209,0.326) (-1.174,0.321) (-1.140,0.316) (-1.105,0.310) (-1.070,0.304) (-1.036,0.298) (-1.001,0.292) (-0.967,0.284) (-0.932,0.283) (-0.896,0.283) (-0.861,0.283) (-0.826,0.283) (-0.791,0.283) (-0.756,0.283) (-0.721,0.283) (-0.686,0.283) (-0.650,0.283) (-0.615,0.283) (-0.580,0.283) (-0.545,0.283) (-0.510,0.283) (-0.475,0.283) (-0.439,0.283) (-0.404,0.283) (-0.369,0.283) (-0.334,0.283) (-0.299,0.283) (-0.264,0.283) (-0.229,0.283) (-0.193,0.283) (-0.158,0.283) (-0.123,0.283) (-0.088,0.283) (-0.053,0.283) (-0.018,0.283) (0.018,0.283) (0.053,0.283) (0.088,0.283) (0.123,0.283) (0.158,0.283) (0.193,0.283) (0.229,0.283) (0.264,0.283) (0.299,0.283) (0.334,0.283) (0.369,0.283) (0.404,0.283) (0.439,0.283) (0.475,0.283) (0.510,0.283) (0.545,0.283) (0.580,0.283) (0.615,0.283) (0.650,0.283) (0.686,0.283) (0.721,0.283) (0.756,0.283) (0.791,0.283) (0.826,0.283) (0.861,0.283) (0.896,0.283) (0.932,0.283) (0.967,0.284) (1.001,0.292) (1.036,0.298) (1.070,0.304) (1.105,0.310) (1.140,0.316) (1.174,0.321) (1.209,0.326) (1.244,0.330) (1.279,0.334) (1.314,0.338) (1.349,0.342) (1.384,0.345) (1.419,0.348) (1.454,0.351) (1.489,0.354) (1.524,0.356) (1.559,0.359) (1.594,0.361) (1.629,0.363) (1.664,0.365) (1.700,0.367) (1.735,0.368) (1.770,0.370) (1.805,0.371) (1.840,0.373) (1.875,0.374) (1.910,0.375) (1.945,0.377) (1.981,0.378) (2.016,0.379) (2.051,0.380) (2.086,0.381) (2.121,0.381) (2.156,0.382) (2.191,0.383) (2.227,0.384) (2.262,0.384) (2.297,0.385) (2.332,0.386) (2.367,0.386) (2.402,0.387) (2.438,0.387) (2.473,0.388) (2.508,0.388) (2.543,0.389) (2.578,0.389) (2.613,0.390) (2.648,0.390) (2.684,0.390) (2.719,0.391) (2.754,0.391) (2.789,0.391) (2.824,0.392) (2.859,0.392) (2.895,0.392) (2.930,0.392) (2.965,0.393) (3.000,0.393)};
\addplot[lineaD, flechaen=0.141, flechaen=0.859] coordinates {(-3.000,0.590) (-2.965,0.589) (-2.930,0.589) (-2.895,0.589) (-2.860,0.588) (-2.825,0.588) (-2.790,0.587) (-2.755,0.587) (-2.720,0.587) (-2.685,0.586) (-2.650,0.586) (-2.615,0.585) (-2.580,0.584) (-2.545,0.584) (-2.510,0.583) (-2.474,0.583) (-2.439,0.582) (-2.404,0.581) (-2.369,0.580) (-2.334,0.580) (-2.299,0.579) (-2.264,0.578) (-2.229,0.577) (-2.194,0.576) (-2.159,0.575) (-2.124,0.574) (-2.089,0.572) (-2.054,0.571) (-2.019,0.570) (-1.984,0.569) (-1.949,0.567) (-1.914,0.565) (-1.879,0.564) (-1.844,0.562) (-1.809,0.560) (-1.774,0.558) (-1.739,0.556) (-1.704,0.554) (-1.669,0.551) (-1.634,0.549) (-1.599,0.546) (-1.565,0.543) (-1.530,0.540) (-1.495,0.537) (-1.460,0.533) (-1.425,0.529) (-1.390,0.525) (-1.356,0.521) (-1.321,0.516) (-1.286,0.511) (-1.251,0.506) (-1.217,0.500) (-1.182,0.494) (-1.148,0.487) (-1.114,0.480) (-1.080,0.473) (-1.045,0.465) (-1.011,0.456) (-0.978,0.447) (-0.944,0.437) (-0.911,0.426) (-0.876,0.424) (-0.841,0.424) (-0.806,0.424) (-0.771,0.424) (-0.736,0.424) (-0.701,0.424) (-0.666,0.424) (-0.631,0.424) (-0.596,0.424) (-0.561,0.424) (-0.526,0.424) (-0.491,0.424) (-0.455,0.424) (-0.420,0.424) (-0.385,0.424) (-0.350,0.424) (-0.315,0.424) (-0.280,0.424) (-0.245,0.424) (-0.210,0.424) (-0.175,0.424) (-0.140,0.424) (-0.105,0.424) (-0.070,0.424) (-0.035,0.424) (0.000,0.424) (0.035,0.424) (0.070,0.424) (0.105,0.424) (0.140,0.424) (0.175,0.424) (0.210,0.424) (0.245,0.424) (0.280,0.424) (0.315,0.424) (0.350,0.424) (0.385,0.424) (0.420,0.424) (0.455,0.424) (0.491,0.424) (0.526,0.424) (0.561,0.424) (0.596,0.424) (0.631,0.424) (0.666,0.424) (0.701,0.424) (0.736,0.424) (0.771,0.424) (0.806,0.424) (0.841,0.424) (0.876,0.424) (0.911,0.426) (0.944,0.437) (0.978,0.447) (1.011,0.456) (1.045,0.465) (1.080,0.473) (1.114,0.480) (1.148,0.487) (1.182,0.494) (1.217,0.500) (1.251,0.506) (1.286,0.511) (1.321,0.516) (1.356,0.521) (1.390,0.525) (1.425,0.529) (1.460,0.533) (1.495,0.537) (1.530,0.540) (1.565,0.543) (1.599,0.546) (1.634,0.549) (1.669,0.551) (1.704,0.554) (1.739,0.556) (1.774,0.558) (1.809,0.560) (1.844,0.562) (1.879,0.564) (1.914,0.565) (1.949,0.567) (1.984,0.569) (2.019,0.570) (2.054,0.571) (2.089,0.572) (2.124,0.574) (2.159,0.575) (2.194,0.576) (2.229,0.577) (2.264,0.578) (2.299,0.579) (2.334,0.580) (2.369,0.580) (2.404,0.581) (2.439,0.582) (2.474,0.583) (2.510,0.583) (2.545,0.584) (2.580,0.584) (2.615,0.585) (2.650,0.586) (2.685,0.586) (2.720,0.587) (2.755,0.587) (2.790,0.587) (2.825,0.588) (2.860,0.588) (2.895,0.589) (2.930,0.589) (2.965,0.589) (3.000,0.590)};
\addplot[lineaD, flechaen=0.141, flechaen=0.859] coordinates {(-3.000,0.787) (-2.965,0.786) (-2.930,0.786) (-2.895,0.786) (-2.859,0.785) (-2.824,0.785) (-2.789,0.784) (-2.754,0.783) (-2.719,0.783) (-2.684,0.782) (-2.649,0.782) (-2.613,0.781) (-2.578,0.780) (-2.543,0.780) (-2.508,0.779) (-2.473,0.778) (-2.438,0.777) (-2.402,0.776) (-2.367,0.775) (-2.332,0.774) (-2.297,0.773) (-2.262,0.772) (-2.227,0.771) (-2.192,0.770) (-2.157,0.769) (-2.121,0.767) (-2.086,0.766) (-2.051,0.764) (-2.016,0.763) (-1.981,0.761) (-1.946,0.759) (-1.911,0.758) (-1.876,0.756) (-1.841,0.754) (-1.805,0.751) (-1.770,0.749) (-1.735,0.747) (-1.700,0.744) (-1.665,0.741) (-1.630,0.738) (-1.595,0.735) (-1.560,0.732) (-1.525,0.728) (-1.490,0.724) (-1.455,0.720) (-1.420,0.716) (-1.386,0.711) (-1.351,0.706) (-1.316,0.701) (-1.281,0.696) (-1.247,0.690) (-1.212,0.683) (-1.178,0.676) (-1.143,0.669) (-1.109,0.661) (-1.075,0.653) (-1.041,0.643) (-1.007,0.634) (-0.974,0.623) (-0.940,0.612) (-0.907,0.600) (-0.874,0.587) (-0.842,0.574) (-0.809,0.566) (-0.773,0.566) (-0.738,0.566) (-0.703,0.566) (-0.668,0.566) (-0.633,0.566) (-0.598,0.566) (-0.562,0.566) (-0.527,0.566) (-0.492,0.566) (-0.457,0.566) (-0.422,0.566) (-0.387,0.566) (-0.352,0.566) (-0.316,0.566) (-0.281,0.566) (-0.246,0.566) (-0.211,0.566) (-0.176,0.566) (-0.141,0.566) (-0.105,0.566) (-0.070,0.566) (-0.035,0.566) (0.000,0.566) (0.035,0.566) (0.070,0.566) (0.105,0.566) (0.141,0.566) (0.176,0.566) (0.211,0.566) (0.246,0.566) (0.281,0.566) (0.316,0.566) (0.352,0.566) (0.387,0.566) (0.422,0.566) (0.457,0.566) (0.492,0.566) (0.527,0.566) (0.562,0.566) (0.598,0.566) (0.633,0.566) (0.668,0.566) (0.703,0.566) (0.738,0.566) (0.773,0.566) (0.809,0.566) (0.842,0.574) (0.874,0.587) (0.907,0.600) (0.940,0.612) (0.974,0.623) (1.007,0.634) (1.041,0.643) (1.075,0.653) (1.109,0.661) (1.143,0.669) (1.178,0.676) (1.212,0.683) (1.247,0.690) (1.281,0.696) (1.316,0.701) (1.351,0.706) (1.386,0.711) (1.420,0.716) (1.455,0.720) (1.490,0.724) (1.525,0.728) (1.560,0.732) (1.595,0.735) (1.630,0.738) (1.665,0.741) (1.700,0.744) (1.735,0.747) (1.770,0.749) (1.805,0.751) (1.841,0.754) (1.876,0.756) (1.911,0.758) (1.946,0.759) (1.981,0.761) (2.016,0.763) (2.051,0.764) (2.086,0.766) (2.121,0.767) (2.157,0.769) (2.192,0.770) (2.227,0.771) (2.262,0.772) (2.297,0.773) (2.332,0.774) (2.367,0.775) (2.402,0.776) (2.438,0.777) (2.473,0.778) (2.508,0.779) (2.543,0.780) (2.578,0.780) (2.613,0.781) (2.649,0.782) (2.684,0.782) (2.719,0.783) (2.754,0.783) (2.789,0.784) (2.824,0.785) (2.859,0.785) (2.895,0.786) (2.930,0.786) (2.965,0.786) (3.000,0.787)};
\addplot[lineaD, flechaen=0.140, flechaen=0.860] coordinates {(-3.000,0.984) (-2.965,0.984) (-2.930,0.983) (-2.895,0.983) (-2.860,0.982) (-2.825,0.982) (-2.789,0.981) (-2.754,0.981) (-2.719,0.980) (-2.684,0.979) (-2.649,0.979) (-2.614,0.978) (-2.579,0.977) (-2.544,0.976) (-2.509,0.975) (-2.474,0.974) (-2.439,0.973) (-2.404,0.972) (-2.368,0.971) (-2.333,0.970) (-2.298,0.969) (-2.263,0.968) (-2.228,0.967) (-2.193,0.966) (-2.158,0.964) (-2.123,0.963) (-2.088,0.961) (-2.053,0.960) (-2.018,0.958) (-1.983,0.956) (-1.948,0.954) (-1.913,0.952) (-1.878,0.950) (-1.843,0.948) (-1.808,0.946) (-1.773,0.943) (-1.738,0.941) (-1.703,0.938) (-1.668,0.935) (-1.633,0.932) (-1.598,0.929) (-1.563,0.926) (-1.528,0.922) (-1.493,0.918) (-1.458,0.914) (-1.423,0.910) (-1.389,0.906) (-1.354,0.901) (-1.319,0.896) (-1.284,0.890) (-1.250,0.884) (-1.215,0.878) (-1.181,0.872) (-1.146,0.865) (-1.112,0.857) (-1.078,0.849) (-1.044,0.841) (-1.010,0.831) (-0.976,0.822) (-0.943,0.811) (-0.910,0.800) (-0.877,0.788) (-0.844,0.775) (-0.812,0.761) (-0.780,0.746) (-0.749,0.731) (-0.718,0.713) (-0.684,0.707) (-0.649,0.707) (-0.614,0.707) (-0.579,0.707) (-0.544,0.707) (-0.509,0.707) (-0.474,0.707) (-0.439,0.707) (-0.404,0.707) (-0.368,0.707) (-0.333,0.707) (-0.298,0.707) (-0.263,0.707) (-0.228,0.707) (-0.193,0.707) (-0.158,0.707) (-0.123,0.707) (-0.088,0.707) (-0.053,0.707) (-0.018,0.707) (0.018,0.707) (0.053,0.707) (0.088,0.707) (0.123,0.707) (0.158,0.707) (0.193,0.707) (0.228,0.707) (0.263,0.707) (0.298,0.707) (0.333,0.707) (0.368,0.707) (0.404,0.707) (0.439,0.707) (0.474,0.707) (0.509,0.707) (0.544,0.707) (0.579,0.707) (0.614,0.707) (0.649,0.707) (0.684,0.707) (0.718,0.713) (0.749,0.731) (0.780,0.746) (0.812,0.761) (0.844,0.775) (0.877,0.788) (0.910,0.800) (0.943,0.811) (0.976,0.822) (1.010,0.831) (1.044,0.841) (1.078,0.849) (1.112,0.857) (1.146,0.865) (1.181,0.872) (1.215,0.878) (1.250,0.884) (1.284,0.890) (1.319,0.896) (1.354,0.901) (1.389,0.906) (1.423,0.910) (1.458,0.914) (1.493,0.918) (1.528,0.922) (1.563,0.926) (1.598,0.929) (1.633,0.932) (1.668,0.935) (1.703,0.938) (1.738,0.941) (1.773,0.943) (1.808,0.946) (1.843,0.948) (1.878,0.950) (1.913,0.952) (1.948,0.954) (1.983,0.956) (2.018,0.958) (2.053,0.960) (2.088,0.961) (2.123,0.963) (2.158,0.964) (2.193,0.966) (2.228,0.967) (2.263,0.968) (2.298,0.969) (2.333,0.970) (2.368,0.971) (2.404,0.972) (2.439,0.973) (2.474,0.974) (2.509,0.975) (2.544,0.976) (2.579,0.977) (2.614,0.978) (2.649,0.979) (2.684,0.979) (2.719,0.980) (2.754,0.981) (2.789,0.981) (2.825,0.982) (2.860,0.982) (2.895,0.983) (2.930,0.983) (2.965,0.984) (3.000,0.984)};
\addplot[lineaD, flechaen=0.139, flechaen=0.861] coordinates {(-3.000,1.182) (-2.965,1.182) (-2.930,1.181) (-2.895,1.181) (-2.860,1.180) (-2.825,1.180) (-2.790,1.179) (-2.755,1.178) (-2.720,1.178) (-2.685,1.177) (-2.650,1.176) (-2.615,1.175) (-2.580,1.174) (-2.545,1.174) (-2.510,1.173) (-2.475,1.172) (-2.440,1.171) (-2.404,1.170) (-2.369,1.169) (-2.334,1.168) (-2.299,1.166) (-2.264,1.165) (-2.229,1.164) (-2.194,1.163) (-2.159,1.161) (-2.124,1.160) (-2.089,1.158) (-2.054,1.157) (-2.019,1.155) (-1.984,1.153) (-1.949,1.151) (-1.914,1.149) (-1.879,1.147) (-1.844,1.145) (-1.809,1.143) (-1.774,1.141) (-1.740,1.138) (-1.705,1.136) (-1.670,1.133) (-1.635,1.130) (-1.600,1.127) (-1.565,1.124) (-1.530,1.121) (-1.495,1.117) (-1.460,1.113) (-1.426,1.110) (-1.391,1.105) (-1.356,1.101) (-1.321,1.097) (-1.287,1.092) (-1.252,1.087) (-1.217,1.081) (-1.183,1.076) (-1.148,1.070) (-1.114,1.063) (-1.079,1.056) (-1.045,1.049) (-1.011,1.042) (-0.977,1.034) (-0.943,1.025) (-0.909,1.016) (-0.875,1.006) (-0.842,0.996) (-0.808,0.985) (-0.775,0.974) (-0.742,0.961) (-0.710,0.948) (-0.678,0.934) (-0.646,0.919) (-0.615,0.903) (-0.585,0.885) (-0.555,0.867) (-0.526,0.849) (-0.491,0.849) (-0.456,0.849) (-0.420,0.849) (-0.385,0.849) (-0.350,0.849) (-0.315,0.849) (-0.280,0.849) (-0.245,0.849) (-0.210,0.849) (-0.175,0.849) (-0.140,0.849) (-0.105,0.849) (-0.070,0.849) (-0.035,0.849) (0.000,0.849) (0.035,0.849) (0.070,0.849) (0.105,0.849) (0.140,0.849) (0.175,0.849) (0.210,0.849) (0.245,0.849) (0.280,0.849) (0.315,0.849) (0.350,0.849) (0.385,0.849) (0.420,0.849) (0.456,0.849) (0.491,0.849) (0.526,0.849) (0.555,0.867) (0.585,0.885) (0.615,0.903) (0.646,0.919) (0.678,0.934) (0.710,0.948) (0.742,0.961) (0.775,0.974) (0.808,0.985) (0.842,0.996) (0.875,1.006) (0.909,1.016) (0.943,1.025) (0.977,1.034) (1.011,1.042) (1.045,1.049) (1.079,1.056) (1.114,1.063) (1.148,1.070) (1.183,1.076) (1.217,1.081) (1.252,1.087) (1.287,1.092) (1.321,1.097) (1.356,1.101) (1.391,1.105) (1.426,1.110) (1.460,1.113) (1.495,1.117) (1.530,1.121) (1.565,1.124) (1.600,1.127) (1.635,1.130) (1.670,1.133) (1.705,1.136) (1.740,1.138) (1.774,1.141) (1.809,1.143) (1.844,1.145) (1.879,1.147) (1.914,1.149) (1.949,1.151) (1.984,1.153) (2.019,1.155) (2.054,1.157) (2.089,1.158) (2.124,1.160) (2.159,1.161) (2.194,1.163) (2.229,1.164) (2.264,1.165) (2.299,1.166) (2.334,1.168) (2.369,1.169) (2.404,1.170) (2.440,1.171) (2.475,1.172) (2.510,1.173) (2.545,1.174) (2.580,1.174) (2.615,1.175) (2.650,1.176) (2.685,1.177) (2.720,1.178) (2.755,1.178) (2.790,1.179) (2.825,1.180) (2.860,1.180) (2.895,1.181) (2.930,1.181) (2.965,1.182) (3.000,1.182)};
\addplot[lineaD, flechaen=0.139, flechaen=0.861] coordinates {(-3.000,1.381) (-2.965,1.380) (-2.930,1.380) (-2.895,1.379) (-2.860,1.379) (-2.825,1.378) (-2.790,1.377) (-2.755,1.377) (-2.720,1.376) (-2.685,1.375) (-2.650,1.374) (-2.615,1.374) (-2.580,1.373) (-2.545,1.372) (-2.510,1.371) (-2.475,1.370) (-2.439,1.369) (-2.404,1.368) (-2.369,1.367) (-2.334,1.366) (-2.299,1.365) (-2.264,1.364) (-2.229,1.362) (-2.194,1.361) (-2.159,1.360) (-2.124,1.358) (-2.089,1.357) (-2.054,1.355) (-2.019,1.354) (-1.984,1.352) (-1.949,1.350) (-1.914,1.348) (-1.879,1.347) (-1.844,1.345) (-1.809,1.343) (-1.774,1.340) (-1.739,1.338) (-1.704,1.336) (-1.669,1.333) (-1.635,1.331) (-1.600,1.328) (-1.565,1.325) (-1.530,1.322) (-1.495,1.319) (-1.460,1.316) (-1.425,1.313) (-1.390,1.309) (-1.355,1.305) (-1.321,1.302) (-1.286,1.298) (-1.251,1.293) (-1.216,1.289) (-1.181,1.284) (-1.147,1.280) (-1.112,1.274) (-1.077,1.269) (-1.043,1.264) (-1.008,1.258) (-0.974,1.252) (-0.939,1.246) (-0.905,1.239) (-0.871,1.232) (-0.836,1.225) (-0.802,1.217) (-0.768,1.209) (-0.734,1.201) (-0.700,1.193) (-0.666,1.184) (-0.632,1.174) (-0.599,1.164) (-0.565,1.154) (-0.532,1.144) (-0.498,1.133) (-0.465,1.121) (-0.432,1.109) (-0.399,1.097) (-0.367,1.085) (-0.334,1.071) (-0.302,1.058) (-0.270,1.044) (-0.237,1.031) (-0.205,1.017) (-0.173,1.003) (-0.140,0.991) (-0.105,0.990) (-0.070,0.990) (-0.035,0.990) (-0.000,0.990) (0.035,0.990) (0.070,0.990) (0.105,0.990) (0.140,0.991) (0.173,1.003) (0.205,1.017) (0.237,1.031) (0.270,1.044) (0.302,1.058) (0.334,1.071) (0.367,1.085) (0.399,1.097) (0.432,1.109) (0.465,1.121) (0.498,1.133) (0.532,1.144) (0.565,1.154) (0.599,1.164) (0.632,1.174) (0.666,1.184) (0.700,1.193) (0.734,1.201) (0.768,1.209) (0.802,1.217) (0.836,1.225) (0.871,1.232) (0.905,1.239) (0.939,1.246) (0.974,1.252) (1.008,1.258) (1.043,1.264) (1.077,1.269) (1.112,1.274) (1.147,1.280) (1.181,1.284) (1.216,1.289) (1.251,1.293) (1.286,1.298) (1.321,1.302) (1.355,1.305) (1.390,1.309) (1.425,1.313) (1.460,1.316) (1.495,1.319) (1.530,1.322) (1.565,1.325) (1.600,1.328) (1.635,1.331) (1.669,1.333) (1.704,1.336) (1.739,1.338) (1.774,1.340) (1.809,1.343) (1.844,1.345) (1.879,1.347) (1.914,1.348) (1.949,1.350) (1.984,1.352) (2.019,1.354) (2.054,1.355) (2.089,1.357) (2.124,1.358) (2.159,1.360) (2.194,1.361) (2.229,1.362) (2.264,1.364) (2.299,1.365) (2.334,1.366) (2.369,1.367) (2.404,1.368) (2.439,1.369) (2.475,1.370) (2.510,1.371) (2.545,1.372) (2.580,1.373) (2.615,1.374) (2.650,1.374) (2.685,1.375) (2.720,1.376) (2.755,1.377) (2.790,1.377) (2.825,1.378) (2.860,1.379) (2.895,1.379) (2.930,1.380) (2.965,1.380) (3.000,1.381)};
\addplot[lineaD, flechaen=0.141, flechaen=0.859] coordinates {(-3.000,1.580) (-2.965,1.579) (-2.930,1.579) (-2.895,1.578) (-2.860,1.578) (-2.825,1.577) (-2.790,1.576) (-2.755,1.576) (-2.720,1.575) (-2.685,1.574) (-2.650,1.573) (-2.615,1.573) (-2.580,1.572) (-2.545,1.571) (-2.510,1.570) (-2.475,1.569) (-2.440,1.568) (-2.405,1.567) (-2.370,1.566) (-2.335,1.565) (-2.300,1.564) (-2.265,1.563) (-2.230,1.562) (-2.195,1.561) (-2.160,1.559) (-2.125,1.558) (-2.090,1.557) (-2.055,1.555) (-2.020,1.554) (-1.985,1.552) (-1.950,1.551) (-1.915,1.549) (-1.880,1.547) (-1.845,1.546) (-1.810,1.544) (-1.775,1.542) (-1.740,1.540) (-1.705,1.538) (-1.670,1.536) (-1.635,1.533) (-1.600,1.531) (-1.565,1.529) (-1.530,1.526) (-1.495,1.524) (-1.460,1.521) (-1.425,1.518) (-1.390,1.515) (-1.355,1.512) (-1.320,1.509) (-1.286,1.506) (-1.251,1.503) (-1.216,1.499) (-1.181,1.496) (-1.146,1.492) (-1.111,1.488) (-1.077,1.484) (-1.042,1.480) (-1.007,1.476) (-0.972,1.471) (-0.937,1.467) (-0.903,1.462) (-0.868,1.458) (-0.833,1.453) (-0.799,1.448) (-0.764,1.443) (-0.729,1.438) (-0.695,1.432) (-0.660,1.427) (-0.626,1.421) (-0.591,1.416) (-0.556,1.410) (-0.522,1.405) (-0.487,1.399) (-0.453,1.394) (-0.418,1.388) (-0.383,1.383) (-0.349,1.378) (-0.314,1.373) (-0.279,1.368) (-0.245,1.364) (-0.210,1.360) (-0.175,1.357) (-0.140,1.354) (-0.105,1.352) (-0.070,1.350) (-0.035,1.349) (0.000,1.348) (0.035,1.349) (0.070,1.350) (0.105,1.352) (0.140,1.354) (0.175,1.357) (0.210,1.360) (0.245,1.364) (0.279,1.368) (0.314,1.373) (0.349,1.378) (0.383,1.383) (0.418,1.388) (0.453,1.394) (0.487,1.399) (0.522,1.405) (0.556,1.410) (0.591,1.416) (0.626,1.421) (0.660,1.427) (0.695,1.432) (0.729,1.438) (0.764,1.443) (0.799,1.448) (0.833,1.453) (0.868,1.458) (0.903,1.462) (0.937,1.467) (0.972,1.471) (1.007,1.476) (1.042,1.480) (1.077,1.484) (1.111,1.488) (1.146,1.492) (1.181,1.496) (1.216,1.499) (1.251,1.503) (1.286,1.506) (1.320,1.509) (1.355,1.512) (1.390,1.515) (1.425,1.518) (1.460,1.521) (1.495,1.524) (1.530,1.526) (1.565,1.529) (1.600,1.531) (1.635,1.533) (1.670,1.536) (1.705,1.538) (1.740,1.540) (1.775,1.542) (1.810,1.544) (1.845,1.546) (1.880,1.547) (1.915,1.549) (1.950,1.551) (1.985,1.552) (2.020,1.554) (2.055,1.555) (2.090,1.557) (2.125,1.558) (2.160,1.559) (2.195,1.561) (2.230,1.562) (2.265,1.563) (2.300,1.564) (2.335,1.565) (2.370,1.566) (2.405,1.567) (2.440,1.568) (2.475,1.569) (2.510,1.570) (2.545,1.571) (2.580,1.572) (2.615,1.573) (2.650,1.573) (2.685,1.574) (2.720,1.575) (2.755,1.576) (2.790,1.576) (2.825,1.577) (2.860,1.578) (2.895,1.578) (2.930,1.579) (2.965,1.579) (3.000,1.580)};
\addplot[lineaD, flechaen=0.141, flechaen=0.859] coordinates {(-3.000,1.779) (-2.965,1.779) (-2.930,1.778) (-2.895,1.777) (-2.859,1.777) (-2.824,1.776) (-2.789,1.776) (-2.754,1.775) (-2.719,1.774) (-2.684,1.774) (-2.649,1.773) (-2.613,1.772) (-2.578,1.771) (-2.543,1.770) (-2.508,1.770) (-2.473,1.769) (-2.438,1.768) (-2.403,1.767) (-2.367,1.766) (-2.332,1.765) (-2.297,1.764) (-2.262,1.763) (-2.227,1.762) (-2.192,1.761) (-2.157,1.760) (-2.121,1.758) (-2.086,1.757) (-2.051,1.756) (-2.016,1.755) (-1.981,1.753) (-1.946,1.752) (-1.911,1.750) (-1.876,1.749) (-1.841,1.747) (-1.805,1.746) (-1.770,1.744) (-1.735,1.742) (-1.700,1.741) (-1.665,1.739) (-1.630,1.737) (-1.595,1.735) (-1.560,1.733) (-1.525,1.731) (-1.489,1.729) (-1.454,1.727) (-1.419,1.724) (-1.384,1.722) (-1.349,1.720) (-1.314,1.717) (-1.279,1.715) (-1.244,1.712) (-1.209,1.709) (-1.174,1.707) (-1.139,1.704) (-1.104,1.701) (-1.069,1.698) (-1.034,1.695) (-0.999,1.692) (-0.964,1.689) (-0.929,1.686) (-0.894,1.682) (-0.859,1.679) (-0.824,1.676) (-0.789,1.673) (-0.754,1.669) (-0.719,1.666) (-0.684,1.663) (-0.649,1.659) (-0.614,1.656) (-0.579,1.653) (-0.544,1.649) (-0.509,1.646) (-0.474,1.643) (-0.439,1.640) (-0.404,1.637) (-0.369,1.635) (-0.334,1.632) (-0.299,1.630) (-0.263,1.628) (-0.228,1.626) (-0.193,1.624) (-0.158,1.622) (-0.123,1.621) (-0.088,1.620) (-0.053,1.620) (-0.018,1.619) (0.018,1.619) (0.053,1.620) (0.088,1.620) (0.123,1.621) (0.158,1.622) (0.193,1.624) (0.228,1.626) (0.263,1.628) (0.299,1.630) (0.334,1.632) (0.369,1.635) (0.404,1.637) (0.439,1.640) (0.474,1.643) (0.509,1.646) (0.544,1.649) (0.579,1.653) (0.614,1.656) (0.649,1.659) (0.684,1.663) (0.719,1.666) (0.754,1.669) (0.789,1.673) (0.824,1.676) (0.859,1.679) (0.894,1.682) (0.929,1.686) (0.964,1.689) (0.999,1.692) (1.034,1.695) (1.069,1.698) (1.104,1.701) (1.139,1.704) (1.174,1.707) (1.209,1.709) (1.244,1.712) (1.279,1.715) (1.314,1.717) (1.349,1.720) (1.384,1.722) (1.419,1.724) (1.454,1.727) (1.489,1.729) (1.525,1.731) (1.560,1.733) (1.595,1.735) (1.630,1.737) (1.665,1.739) (1.700,1.741) (1.735,1.742) (1.770,1.744) (1.805,1.746) (1.841,1.747) (1.876,1.749) (1.911,1.750) (1.946,1.752) (1.981,1.753) (2.016,1.755) (2.051,1.756) (2.086,1.757) (2.121,1.758) (2.157,1.760) (2.192,1.761) (2.227,1.762) (2.262,1.763) (2.297,1.764) (2.332,1.765) (2.367,1.766) (2.403,1.767) (2.438,1.768) (2.473,1.769) (2.508,1.770) (2.543,1.770) (2.578,1.771) (2.613,1.772) (2.649,1.773) (2.684,1.774) (2.719,1.774) (2.754,1.775) (2.789,1.776) (2.824,1.776) (2.859,1.777) (2.895,1.777) (2.930,1.778) (2.965,1.779) (3.000,1.779)};
\addplot[lineaD, flechaen=0.142, flechaen=0.858] coordinates {(-3.000,1.979) (-2.965,1.978) (-2.930,1.978) (-2.895,1.977) (-2.860,1.977) (-2.824,1.976) (-2.789,1.975) (-2.754,1.975) (-2.719,1.974) (-2.684,1.973) (-2.649,1.973) (-2.614,1.972) (-2.579,1.971) (-2.544,1.971) (-2.508,1.970) (-2.473,1.969) (-2.438,1.968) (-2.403,1.967) (-2.368,1.967) (-2.333,1.966) (-2.298,1.965) (-2.263,1.964) (-2.228,1.963) (-2.192,1.962) (-2.157,1.961) (-2.122,1.960) (-2.087,1.959) (-2.052,1.958) (-2.017,1.956) (-1.982,1.955) (-1.947,1.954) (-1.912,1.953) (-1.877,1.951) (-1.841,1.950) (-1.806,1.949) (-1.771,1.947) (-1.736,1.946) (-1.701,1.944) (-1.666,1.943) (-1.631,1.941) (-1.596,1.940) (-1.561,1.938) (-1.526,1.936) (-1.491,1.935) (-1.455,1.933) (-1.420,1.931) (-1.385,1.929) (-1.350,1.927) (-1.315,1.926) (-1.280,1.924) (-1.245,1.922) (-1.210,1.920) (-1.175,1.917) (-1.140,1.915) (-1.105,1.913) (-1.070,1.911) (-1.035,1.909) (-1.000,1.907) (-0.965,1.905) (-0.930,1.902) (-0.895,1.900) (-0.859,1.898) (-0.824,1.896) (-0.789,1.893) (-0.754,1.891) (-0.719,1.889) (-0.684,1.887) (-0.649,1.885) (-0.614,1.882) (-0.579,1.880) (-0.544,1.878) (-0.509,1.876) (-0.474,1.875) (-0.439,1.873) (-0.404,1.871) (-0.369,1.869) (-0.334,1.868) (-0.298,1.867) (-0.263,1.865) (-0.228,1.864) (-0.193,1.863) (-0.158,1.862) (-0.123,1.862) (-0.088,1.861) (-0.053,1.861) (-0.018,1.861) (0.018,1.861) (0.053,1.861) (0.088,1.861) (0.123,1.862) (0.158,1.862) (0.193,1.863) (0.228,1.864) (0.263,1.865) (0.298,1.867) (0.334,1.868) (0.369,1.869) (0.404,1.871) (0.439,1.873) (0.474,1.875) (0.509,1.876) (0.544,1.878) (0.579,1.880) (0.614,1.882) (0.649,1.885) (0.684,1.887) (0.719,1.889) (0.754,1.891) (0.789,1.893) (0.824,1.896) (0.859,1.898) (0.895,1.900) (0.930,1.902) (0.965,1.905) (1.000,1.907) (1.035,1.909) (1.070,1.911) (1.105,1.913) (1.140,1.915) (1.175,1.917) (1.210,1.920) (1.245,1.922) (1.280,1.924) (1.315,1.926) (1.350,1.927) (1.385,1.929) (1.420,1.931) (1.455,1.933) (1.491,1.935) (1.526,1.936) (1.561,1.938) (1.596,1.940) (1.631,1.941) (1.666,1.943) (1.701,1.944) (1.736,1.946) (1.771,1.947) (1.806,1.949) (1.841,1.950) (1.877,1.951) (1.912,1.953) (1.947,1.954) (1.982,1.955) (2.017,1.956) (2.052,1.958) (2.087,1.959) (2.122,1.960) (2.157,1.961) (2.192,1.962) (2.228,1.963) (2.263,1.964) (2.298,1.965) (2.333,1.966) (2.368,1.967) (2.403,1.967) (2.438,1.968) (2.473,1.969) (2.508,1.970) (2.544,1.971) (2.579,1.971) (2.614,1.972) (2.649,1.973) (2.684,1.973) (2.719,1.974) (2.754,1.975) (2.789,1.975) (2.824,1.976) (2.860,1.977) (2.895,1.977) (2.930,1.978) (2.965,1.978) (3.000,1.979)};
\addplot[lineaD, flechaen=0.142, flechaen=0.858] coordinates {(-3.000,2.179) (-2.965,2.178) (-2.930,2.178) (-2.895,2.177) (-2.860,2.177) (-2.824,2.176) (-2.789,2.176) (-2.754,2.175) (-2.719,2.174) (-2.684,2.174) (-2.649,2.173) (-2.614,2.173) (-2.579,2.172) (-2.544,2.171) (-2.509,2.170) (-2.473,2.170) (-2.438,2.169) (-2.403,2.168) (-2.368,2.167) (-2.333,2.167) (-2.298,2.166) (-2.263,2.165) (-2.228,2.164) (-2.193,2.163) (-2.158,2.162) (-2.123,2.161) (-2.087,2.161) (-2.052,2.160) (-2.017,2.159) (-1.982,2.158) (-1.947,2.156) (-1.912,2.155) (-1.877,2.154) (-1.842,2.153) (-1.807,2.152) (-1.772,2.151) (-1.737,2.150) (-1.701,2.148) (-1.666,2.147) (-1.631,2.146) (-1.596,2.145) (-1.561,2.143) (-1.526,2.142) (-1.491,2.141) (-1.456,2.139) (-1.421,2.138) (-1.386,2.136) (-1.351,2.135) (-1.316,2.133) (-1.281,2.132) (-1.245,2.130) (-1.210,2.129) (-1.175,2.127) (-1.140,2.126) (-1.105,2.124) (-1.070,2.122) (-1.035,2.121) (-1.000,2.119) (-0.965,2.118) (-0.930,2.116) (-0.895,2.114) (-0.860,2.113) (-0.825,2.111) (-0.790,2.110) (-0.754,2.108) (-0.719,2.107) (-0.684,2.105) (-0.649,2.104) (-0.614,2.102) (-0.579,2.101) (-0.544,2.100) (-0.509,2.098) (-0.474,2.097) (-0.439,2.096) (-0.404,2.095) (-0.369,2.094) (-0.333,2.093) (-0.298,2.092) (-0.263,2.091) (-0.228,2.090) (-0.193,2.090) (-0.158,2.089) (-0.123,2.089) (-0.088,2.089) (-0.053,2.088) (-0.018,2.088) (0.018,2.088) (0.053,2.088) (0.088,2.089) (0.123,2.089) (0.158,2.089) (0.193,2.090) (0.228,2.090) (0.263,2.091) (0.298,2.092) (0.333,2.093) (0.369,2.094) (0.404,2.095) (0.439,2.096) (0.474,2.097) (0.509,2.098) (0.544,2.100) (0.579,2.101) (0.614,2.102) (0.649,2.104) (0.684,2.105) (0.719,2.107) (0.754,2.108) (0.790,2.110) (0.825,2.111) (0.860,2.113) (0.895,2.114) (0.930,2.116) (0.965,2.118) (1.000,2.119) (1.035,2.121) (1.070,2.122) (1.105,2.124) (1.140,2.126) (1.175,2.127) (1.210,2.129) (1.245,2.130) (1.281,2.132) (1.316,2.133) (1.351,2.135) (1.386,2.136) (1.421,2.138) (1.456,2.139) (1.491,2.141) (1.526,2.142) (1.561,2.143) (1.596,2.145) (1.631,2.146) (1.666,2.147) (1.701,2.148) (1.737,2.150) (1.772,2.151) (1.807,2.152) (1.842,2.153) (1.877,2.154) (1.912,2.155) (1.947,2.156) (1.982,2.158) (2.017,2.159) (2.052,2.160) (2.087,2.161) (2.123,2.161) (2.158,2.162) (2.193,2.163) (2.228,2.164) (2.263,2.165) (2.298,2.166) (2.333,2.167) (2.368,2.167) (2.403,2.168) (2.438,2.169) (2.473,2.170) (2.509,2.170) (2.544,2.171) (2.579,2.172) (2.614,2.173) (2.649,2.173) (2.684,2.174) (2.719,2.174) (2.754,2.175) (2.789,2.176) (2.824,2.176) (2.860,2.177) (2.895,2.177) (2.930,2.178) (2.965,2.178) (3.000,2.179)};
\addplot[lineaD, flechaen=0.142, flechaen=0.858] coordinates {(-3.000,-0.196) (-2.965,-0.196) (-2.930,-0.196) (-2.895,-0.196) (-2.860,-0.196) (-2.824,-0.196) (-2.789,-0.196) (-2.754,-0.195) (-2.719,-0.195) (-2.684,-0.195) (-2.649,-0.195) (-2.614,-0.195) (-2.579,-0.194) (-2.544,-0.194) (-2.509,-0.194) (-2.473,-0.194) (-2.438,-0.193) (-2.403,-0.193) (-2.368,-0.193) (-2.333,-0.193) (-2.298,-0.192) (-2.263,-0.192) (-2.228,-0.192) (-2.193,-0.191) (-2.158,-0.191) (-2.122,-0.190) (-2.087,-0.190) (-2.052,-0.189) (-2.017,-0.189) (-1.982,-0.188) (-1.947,-0.188) (-1.912,-0.187) (-1.877,-0.187) (-1.842,-0.186) (-1.806,-0.185) (-1.771,-0.184) (-1.736,-0.183) (-1.701,-0.183) (-1.666,-0.182) (-1.631,-0.181) (-1.596,-0.179) (-1.561,-0.178) (-1.526,-0.177) (-1.491,-0.176) (-1.456,-0.174) (-1.421,-0.173) (-1.385,-0.171) (-1.350,-0.169) (-1.315,-0.167) (-1.280,-0.165) (-1.245,-0.163) (-1.210,-0.161) (-1.175,-0.158) (-1.140,-0.155) (-1.105,-0.152) (-1.070,-0.149) (-1.035,-0.146) (-1.000,-0.143) (-0.965,-0.141) (-0.930,-0.141) (-0.895,-0.141) (-0.860,-0.141) (-0.825,-0.141) (-0.790,-0.141) (-0.755,-0.141) (-0.720,-0.141) (-0.685,-0.141) (-0.649,-0.141) (-0.614,-0.141) (-0.579,-0.141) (-0.544,-0.141) (-0.509,-0.141) (-0.474,-0.141) (-0.439,-0.141) (-0.404,-0.141) (-0.369,-0.141) (-0.333,-0.141) (-0.298,-0.141) (-0.263,-0.141) (-0.228,-0.141) (-0.193,-0.141) (-0.158,-0.141) (-0.123,-0.141) (-0.088,-0.141) (-0.053,-0.141) (-0.018,-0.141) (0.018,-0.141) (0.053,-0.141) (0.088,-0.141) (0.123,-0.141) (0.158,-0.141) (0.193,-0.141) (0.228,-0.141) (0.263,-0.141) (0.298,-0.141) (0.333,-0.141) (0.369,-0.141) (0.404,-0.141) (0.439,-0.141) (0.474,-0.141) (0.509,-0.141) (0.544,-0.141) (0.579,-0.141) (0.614,-0.141) (0.649,-0.141) (0.685,-0.141) (0.720,-0.141) (0.755,-0.141) (0.790,-0.141) (0.825,-0.141) (0.860,-0.141) (0.895,-0.141) (0.930,-0.141) (0.965,-0.141) (1.000,-0.143) (1.035,-0.146) (1.070,-0.149) (1.105,-0.152) (1.140,-0.155) (1.175,-0.158) (1.210,-0.161) (1.245,-0.163) (1.280,-0.165) (1.315,-0.167) (1.350,-0.169) (1.385,-0.171) (1.421,-0.173) (1.456,-0.174) (1.491,-0.176) (1.526,-0.177) (1.561,-0.178) (1.596,-0.179) (1.631,-0.181) (1.666,-0.182) (1.701,-0.183) (1.736,-0.183) (1.771,-0.184) (1.806,-0.185) (1.842,-0.186) (1.877,-0.187) (1.912,-0.187) (1.947,-0.188) (1.982,-0.188) (2.017,-0.189) (2.052,-0.189) (2.087,-0.190) (2.122,-0.190) (2.158,-0.191) (2.193,-0.191) (2.228,-0.192) (2.263,-0.192) (2.298,-0.192) (2.333,-0.193) (2.368,-0.193) (2.403,-0.193) (2.438,-0.193) (2.473,-0.194) (2.509,-0.194) (2.544,-0.194) (2.579,-0.194) (2.614,-0.195) (2.649,-0.195) (2.684,-0.195) (2.719,-0.195) (2.754,-0.195) (2.789,-0.196) (2.824,-0.196) (2.860,-0.196) (2.895,-0.196) (2.930,-0.196) (2.965,-0.196) (3.000,-0.196)};
\addplot[lineaD, flechaen=0.141, flechaen=0.859] coordinates {(-3.000,-0.393) (-2.965,-0.393) (-2.930,-0.392) (-2.895,-0.392) (-2.859,-0.392) (-2.824,-0.392) (-2.789,-0.391) (-2.754,-0.391) (-2.719,-0.391) (-2.684,-0.390) (-2.648,-0.390) (-2.613,-0.390) (-2.578,-0.389) (-2.543,-0.389) (-2.508,-0.388) (-2.473,-0.388) (-2.438,-0.387) (-2.402,-0.387) (-2.367,-0.386) (-2.332,-0.386) (-2.297,-0.385) (-2.262,-0.384) (-2.227,-0.384) (-2.191,-0.383) (-2.156,-0.382) (-2.121,-0.381) (-2.086,-0.381) (-2.051,-0.380) (-2.016,-0.379) (-1.981,-0.378) (-1.945,-0.377) (-1.910,-0.375) (-1.875,-0.374) (-1.840,-0.373) (-1.805,-0.371) (-1.770,-0.370) (-1.735,-0.368) (-1.700,-0.367) (-1.664,-0.365) (-1.629,-0.363) (-1.594,-0.361) (-1.559,-0.359) (-1.524,-0.356) (-1.489,-0.354) (-1.454,-0.351) (-1.419,-0.348) (-1.384,-0.345) (-1.349,-0.342) (-1.314,-0.338) (-1.279,-0.334) (-1.244,-0.330) (-1.209,-0.326) (-1.174,-0.321) (-1.140,-0.316) (-1.105,-0.310) (-1.070,-0.304) (-1.036,-0.298) (-1.001,-0.292) (-0.967,-0.284) (-0.932,-0.283) (-0.896,-0.283) (-0.861,-0.283) (-0.826,-0.283) (-0.791,-0.283) (-0.756,-0.283) (-0.721,-0.283) (-0.686,-0.283) (-0.650,-0.283) (-0.615,-0.283) (-0.580,-0.283) (-0.545,-0.283) (-0.510,-0.283) (-0.475,-0.283) (-0.439,-0.283) (-0.404,-0.283) (-0.369,-0.283) (-0.334,-0.283) (-0.299,-0.283) (-0.264,-0.283) (-0.229,-0.283) (-0.193,-0.283) (-0.158,-0.283) (-0.123,-0.283) (-0.088,-0.283) (-0.053,-0.283) (-0.018,-0.283) (0.018,-0.283) (0.053,-0.283) (0.088,-0.283) (0.123,-0.283) (0.158,-0.283) (0.193,-0.283) (0.229,-0.283) (0.264,-0.283) (0.299,-0.283) (0.334,-0.283) (0.369,-0.283) (0.404,-0.283) (0.439,-0.283) (0.475,-0.283) (0.510,-0.283) (0.545,-0.283) (0.580,-0.283) (0.615,-0.283) (0.650,-0.283) (0.686,-0.283) (0.721,-0.283) (0.756,-0.283) (0.791,-0.283) (0.826,-0.283) (0.861,-0.283) (0.896,-0.283) (0.932,-0.283) (0.967,-0.284) (1.001,-0.292) (1.036,-0.298) (1.070,-0.304) (1.105,-0.310) (1.140,-0.316) (1.174,-0.321) (1.209,-0.326) (1.244,-0.330) (1.279,-0.334) (1.314,-0.338) (1.349,-0.342) (1.384,-0.345) (1.419,-0.348) (1.454,-0.351) (1.489,-0.354) (1.524,-0.356) (1.559,-0.359) (1.594,-0.361) (1.629,-0.363) (1.664,-0.365) (1.700,-0.367) (1.735,-0.368) (1.770,-0.370) (1.805,-0.371) (1.840,-0.373) (1.875,-0.374) (1.910,-0.375) (1.945,-0.377) (1.981,-0.378) (2.016,-0.379) (2.051,-0.380) (2.086,-0.381) (2.121,-0.381) (2.156,-0.382) (2.191,-0.383) (2.227,-0.384) (2.262,-0.384) (2.297,-0.385) (2.332,-0.386) (2.367,-0.386) (2.402,-0.387) (2.438,-0.387) (2.473,-0.388) (2.508,-0.388) (2.543,-0.389) (2.578,-0.389) (2.613,-0.390) (2.648,-0.390) (2.684,-0.390) (2.719,-0.391) (2.754,-0.391) (2.789,-0.391) (2.824,-0.392) (2.859,-0.392) (2.895,-0.392) (2.930,-0.392) (2.965,-0.393) (3.000,-0.393)};
\addplot[lineaD, flechaen=0.141, flechaen=0.859] coordinates {(-3.000,-0.590) (-2.965,-0.589) (-2.930,-0.589) (-2.895,-0.589) (-2.860,-0.588) (-2.825,-0.588) (-2.790,-0.587) (-2.755,-0.587) (-2.720,-0.587) (-2.685,-0.586) (-2.650,-0.586) (-2.615,-0.585) (-2.580,-0.584) (-2.545,-0.584) (-2.510,-0.583) (-2.474,-0.583) (-2.439,-0.582) (-2.404,-0.581) (-2.369,-0.580) (-2.334,-0.580) (-2.299,-0.579) (-2.264,-0.578) (-2.229,-0.577) (-2.194,-0.576) (-2.159,-0.575) (-2.124,-0.574) (-2.089,-0.572) (-2.054,-0.571) (-2.019,-0.570) (-1.984,-0.569) (-1.949,-0.567) (-1.914,-0.565) (-1.879,-0.564) (-1.844,-0.562) (-1.809,-0.560) (-1.774,-0.558) (-1.739,-0.556) (-1.704,-0.554) (-1.669,-0.551) (-1.634,-0.549) (-1.599,-0.546) (-1.565,-0.543) (-1.530,-0.540) (-1.495,-0.537) (-1.460,-0.533) (-1.425,-0.529) (-1.390,-0.525) (-1.356,-0.521) (-1.321,-0.516) (-1.286,-0.511) (-1.251,-0.506) (-1.217,-0.500) (-1.182,-0.494) (-1.148,-0.487) (-1.114,-0.480) (-1.080,-0.473) (-1.045,-0.465) (-1.011,-0.456) (-0.978,-0.447) (-0.944,-0.437) (-0.911,-0.426) (-0.876,-0.424) (-0.841,-0.424) (-0.806,-0.424) (-0.771,-0.424) (-0.736,-0.424) (-0.701,-0.424) (-0.666,-0.424) (-0.631,-0.424) (-0.596,-0.424) (-0.561,-0.424) (-0.526,-0.424) (-0.491,-0.424) (-0.455,-0.424) (-0.420,-0.424) (-0.385,-0.424) (-0.350,-0.424) (-0.315,-0.424) (-0.280,-0.424) (-0.245,-0.424) (-0.210,-0.424) (-0.175,-0.424) (-0.140,-0.424) (-0.105,-0.424) (-0.070,-0.424) (-0.035,-0.424) (0.000,-0.424) (0.035,-0.424) (0.070,-0.424) (0.105,-0.424) (0.140,-0.424) (0.175,-0.424) (0.210,-0.424) (0.245,-0.424) (0.280,-0.424) (0.315,-0.424) (0.350,-0.424) (0.385,-0.424) (0.420,-0.424) (0.455,-0.424) (0.491,-0.424) (0.526,-0.424) (0.561,-0.424) (0.596,-0.424) (0.631,-0.424) (0.666,-0.424) (0.701,-0.424) (0.736,-0.424) (0.771,-0.424) (0.806,-0.424) (0.841,-0.424) (0.876,-0.424) (0.911,-0.426) (0.944,-0.437) (0.978,-0.447) (1.011,-0.456) (1.045,-0.465) (1.080,-0.473) (1.114,-0.480) (1.148,-0.487) (1.182,-0.494) (1.217,-0.500) (1.251,-0.506) (1.286,-0.511) (1.321,-0.516) (1.356,-0.521) (1.390,-0.525) (1.425,-0.529) (1.460,-0.533) (1.495,-0.537) (1.530,-0.540) (1.565,-0.543) (1.599,-0.546) (1.634,-0.549) (1.669,-0.551) (1.704,-0.554) (1.739,-0.556) (1.774,-0.558) (1.809,-0.560) (1.844,-0.562) (1.879,-0.564) (1.914,-0.565) (1.949,-0.567) (1.984,-0.569) (2.019,-0.570) (2.054,-0.571) (2.089,-0.572) (2.124,-0.574) (2.159,-0.575) (2.194,-0.576) (2.229,-0.577) (2.264,-0.578) (2.299,-0.579) (2.334,-0.580) (2.369,-0.580) (2.404,-0.581) (2.439,-0.582) (2.474,-0.583) (2.510,-0.583) (2.545,-0.584) (2.580,-0.584) (2.615,-0.585) (2.650,-0.586) (2.685,-0.586) (2.720,-0.587) (2.755,-0.587) (2.790,-0.587) (2.825,-0.588) (2.860,-0.588) (2.895,-0.589) (2.930,-0.589) (2.965,-0.589) (3.000,-0.590)};
\addplot[lineaD, flechaen=0.141, flechaen=0.859] coordinates {(-3.000,-0.787) (-2.965,-0.786) (-2.930,-0.786) (-2.895,-0.786) (-2.859,-0.785) (-2.824,-0.785) (-2.789,-0.784) (-2.754,-0.783) (-2.719,-0.783) (-2.684,-0.782) (-2.649,-0.782) (-2.613,-0.781) (-2.578,-0.780) (-2.543,-0.780) (-2.508,-0.779) (-2.473,-0.778) (-2.438,-0.777) (-2.402,-0.776) (-2.367,-0.775) (-2.332,-0.774) (-2.297,-0.773) (-2.262,-0.772) (-2.227,-0.771) (-2.192,-0.770) (-2.157,-0.769) (-2.121,-0.767) (-2.086,-0.766) (-2.051,-0.764) (-2.016,-0.763) (-1.981,-0.761) (-1.946,-0.759) (-1.911,-0.758) (-1.876,-0.756) (-1.841,-0.754) (-1.805,-0.751) (-1.770,-0.749) (-1.735,-0.747) (-1.700,-0.744) (-1.665,-0.741) (-1.630,-0.738) (-1.595,-0.735) (-1.560,-0.732) (-1.525,-0.728) (-1.490,-0.724) (-1.455,-0.720) (-1.420,-0.716) (-1.386,-0.711) (-1.351,-0.706) (-1.316,-0.701) (-1.281,-0.696) (-1.247,-0.690) (-1.212,-0.683) (-1.178,-0.676) (-1.143,-0.669) (-1.109,-0.661) (-1.075,-0.653) (-1.041,-0.643) (-1.007,-0.634) (-0.974,-0.623) (-0.940,-0.612) (-0.907,-0.600) (-0.874,-0.587) (-0.842,-0.574) (-0.809,-0.566) (-0.773,-0.566) (-0.738,-0.566) (-0.703,-0.566) (-0.668,-0.566) (-0.633,-0.566) (-0.598,-0.566) (-0.562,-0.566) (-0.527,-0.566) (-0.492,-0.566) (-0.457,-0.566) (-0.422,-0.566) (-0.387,-0.566) (-0.352,-0.566) (-0.316,-0.566) (-0.281,-0.566) (-0.246,-0.566) (-0.211,-0.566) (-0.176,-0.566) (-0.141,-0.566) (-0.105,-0.566) (-0.070,-0.566) (-0.035,-0.566) (0.000,-0.566) (0.035,-0.566) (0.070,-0.566) (0.105,-0.566) (0.141,-0.566) (0.176,-0.566) (0.211,-0.566) (0.246,-0.566) (0.281,-0.566) (0.316,-0.566) (0.352,-0.566) (0.387,-0.566) (0.422,-0.566) (0.457,-0.566) (0.492,-0.566) (0.527,-0.566) (0.562,-0.566) (0.598,-0.566) (0.633,-0.566) (0.668,-0.566) (0.703,-0.566) (0.738,-0.566) (0.773,-0.566) (0.809,-0.566) (0.842,-0.574) (0.874,-0.587) (0.907,-0.600) (0.940,-0.612) (0.974,-0.623) (1.007,-0.634) (1.041,-0.643) (1.075,-0.653) (1.109,-0.661) (1.143,-0.669) (1.178,-0.676) (1.212,-0.683) (1.247,-0.690) (1.281,-0.696) (1.316,-0.701) (1.351,-0.706) (1.386,-0.711) (1.420,-0.716) (1.455,-0.720) (1.490,-0.724) (1.525,-0.728) (1.560,-0.732) (1.595,-0.735) (1.630,-0.738) (1.665,-0.741) (1.700,-0.744) (1.735,-0.747) (1.770,-0.749) (1.805,-0.751) (1.841,-0.754) (1.876,-0.756) (1.911,-0.758) (1.946,-0.759) (1.981,-0.761) (2.016,-0.763) (2.051,-0.764) (2.086,-0.766) (2.121,-0.767) (2.157,-0.769) (2.192,-0.770) (2.227,-0.771) (2.262,-0.772) (2.297,-0.773) (2.332,-0.774) (2.367,-0.775) (2.402,-0.776) (2.438,-0.777) (2.473,-0.778) (2.508,-0.779) (2.543,-0.780) (2.578,-0.780) (2.613,-0.781) (2.649,-0.782) (2.684,-0.782) (2.719,-0.783) (2.754,-0.783) (2.789,-0.784) (2.824,-0.785) (2.859,-0.785) (2.895,-0.786) (2.930,-0.786) (2.965,-0.786) (3.000,-0.787)};
\addplot[lineaD, flechaen=0.140, flechaen=0.860] coordinates {(-3.000,-0.984) (-2.965,-0.984) (-2.930,-0.983) (-2.895,-0.983) (-2.860,-0.982) (-2.825,-0.982) (-2.789,-0.981) (-2.754,-0.981) (-2.719,-0.980) (-2.684,-0.979) (-2.649,-0.979) (-2.614,-0.978) (-2.579,-0.977) (-2.544,-0.976) (-2.509,-0.975) (-2.474,-0.974) (-2.439,-0.973) (-2.404,-0.972) (-2.368,-0.971) (-2.333,-0.970) (-2.298,-0.969) (-2.263,-0.968) (-2.228,-0.967) (-2.193,-0.966) (-2.158,-0.964) (-2.123,-0.963) (-2.088,-0.961) (-2.053,-0.960) (-2.018,-0.958) (-1.983,-0.956) (-1.948,-0.954) (-1.913,-0.952) (-1.878,-0.950) (-1.843,-0.948) (-1.808,-0.946) (-1.773,-0.943) (-1.738,-0.941) (-1.703,-0.938) (-1.668,-0.935) (-1.633,-0.932) (-1.598,-0.929) (-1.563,-0.926) (-1.528,-0.922) (-1.493,-0.918) (-1.458,-0.914) (-1.423,-0.910) (-1.389,-0.906) (-1.354,-0.901) (-1.319,-0.896) (-1.284,-0.890) (-1.250,-0.884) (-1.215,-0.878) (-1.181,-0.872) (-1.146,-0.865) (-1.112,-0.857) (-1.078,-0.849) (-1.044,-0.841) (-1.010,-0.831) (-0.976,-0.822) (-0.943,-0.811) (-0.910,-0.800) (-0.877,-0.788) (-0.844,-0.775) (-0.812,-0.761) (-0.780,-0.746) (-0.749,-0.731) (-0.718,-0.713) (-0.684,-0.707) (-0.649,-0.707) (-0.614,-0.707) (-0.579,-0.707) (-0.544,-0.707) (-0.509,-0.707) (-0.474,-0.707) (-0.439,-0.707) (-0.404,-0.707) (-0.368,-0.707) (-0.333,-0.707) (-0.298,-0.707) (-0.263,-0.707) (-0.228,-0.707) (-0.193,-0.707) (-0.158,-0.707) (-0.123,-0.707) (-0.088,-0.707) (-0.053,-0.707) (-0.018,-0.707) (0.018,-0.707) (0.053,-0.707) (0.088,-0.707) (0.123,-0.707) (0.158,-0.707) (0.193,-0.707) (0.228,-0.707) (0.263,-0.707) (0.298,-0.707) (0.333,-0.707) (0.368,-0.707) (0.404,-0.707) (0.439,-0.707) (0.474,-0.707) (0.509,-0.707) (0.544,-0.707) (0.579,-0.707) (0.614,-0.707) (0.649,-0.707) (0.684,-0.707) (0.718,-0.713) (0.749,-0.731) (0.780,-0.746) (0.812,-0.761) (0.844,-0.775) (0.877,-0.788) (0.910,-0.800) (0.943,-0.811) (0.976,-0.822) (1.010,-0.831) (1.044,-0.841) (1.078,-0.849) (1.112,-0.857) (1.146,-0.865) (1.181,-0.872) (1.215,-0.878) (1.250,-0.884) (1.284,-0.890) (1.319,-0.896) (1.354,-0.901) (1.389,-0.906) (1.423,-0.910) (1.458,-0.914) (1.493,-0.918) (1.528,-0.922) (1.563,-0.926) (1.598,-0.929) (1.633,-0.932) (1.668,-0.935) (1.703,-0.938) (1.738,-0.941) (1.773,-0.943) (1.808,-0.946) (1.843,-0.948) (1.878,-0.950) (1.913,-0.952) (1.948,-0.954) (1.983,-0.956) (2.018,-0.958) (2.053,-0.960) (2.088,-0.961) (2.123,-0.963) (2.158,-0.964) (2.193,-0.966) (2.228,-0.967) (2.263,-0.968) (2.298,-0.969) (2.333,-0.970) (2.368,-0.971) (2.404,-0.972) (2.439,-0.973) (2.474,-0.974) (2.509,-0.975) (2.544,-0.976) (2.579,-0.977) (2.614,-0.978) (2.649,-0.979) (2.684,-0.979) (2.719,-0.980) (2.754,-0.981) (2.789,-0.981) (2.825,-0.982) (2.860,-0.982) (2.895,-0.983) (2.930,-0.983) (2.965,-0.984) (3.000,-0.984)};
\addplot[lineaD, flechaen=0.139, flechaen=0.861] coordinates {(-3.000,-1.182) (-2.965,-1.182) (-2.930,-1.181) (-2.895,-1.181) (-2.860,-1.180) (-2.825,-1.180) (-2.790,-1.179) (-2.755,-1.178) (-2.720,-1.178) (-2.685,-1.177) (-2.650,-1.176) (-2.615,-1.175) (-2.580,-1.174) (-2.545,-1.174) (-2.510,-1.173) (-2.475,-1.172) (-2.440,-1.171) (-2.404,-1.170) (-2.369,-1.169) (-2.334,-1.168) (-2.299,-1.166) (-2.264,-1.165) (-2.229,-1.164) (-2.194,-1.163) (-2.159,-1.161) (-2.124,-1.160) (-2.089,-1.158) (-2.054,-1.157) (-2.019,-1.155) (-1.984,-1.153) (-1.949,-1.151) (-1.914,-1.149) (-1.879,-1.147) (-1.844,-1.145) (-1.809,-1.143) (-1.774,-1.141) (-1.740,-1.138) (-1.705,-1.136) (-1.670,-1.133) (-1.635,-1.130) (-1.600,-1.127) (-1.565,-1.124) (-1.530,-1.121) (-1.495,-1.117) (-1.460,-1.113) (-1.426,-1.110) (-1.391,-1.105) (-1.356,-1.101) (-1.321,-1.097) (-1.287,-1.092) (-1.252,-1.087) (-1.217,-1.081) (-1.183,-1.076) (-1.148,-1.070) (-1.114,-1.063) (-1.079,-1.056) (-1.045,-1.049) (-1.011,-1.042) (-0.977,-1.034) (-0.943,-1.025) (-0.909,-1.016) (-0.875,-1.006) (-0.842,-0.996) (-0.808,-0.985) (-0.775,-0.974) (-0.742,-0.961) (-0.710,-0.948) (-0.678,-0.934) (-0.646,-0.919) (-0.615,-0.903) (-0.585,-0.885) (-0.555,-0.867) (-0.526,-0.849) (-0.491,-0.849) (-0.456,-0.849) (-0.420,-0.849) (-0.385,-0.849) (-0.350,-0.849) (-0.315,-0.849) (-0.280,-0.849) (-0.245,-0.849) (-0.210,-0.849) (-0.175,-0.849) (-0.140,-0.849) (-0.105,-0.849) (-0.070,-0.849) (-0.035,-0.849) (0.000,-0.849) (0.035,-0.849) (0.070,-0.849) (0.105,-0.849) (0.140,-0.849) (0.175,-0.849) (0.210,-0.849) (0.245,-0.849) (0.280,-0.849) (0.315,-0.849) (0.350,-0.849) (0.385,-0.849) (0.420,-0.849) (0.456,-0.849) (0.491,-0.849) (0.526,-0.849) (0.555,-0.867) (0.585,-0.885) (0.615,-0.903) (0.646,-0.919) (0.678,-0.934) (0.710,-0.948) (0.742,-0.961) (0.775,-0.974) (0.808,-0.985) (0.842,-0.996) (0.875,-1.006) (0.909,-1.016) (0.943,-1.025) (0.977,-1.034) (1.011,-1.042) (1.045,-1.049) (1.079,-1.056) (1.114,-1.063) (1.148,-1.070) (1.183,-1.076) (1.217,-1.081) (1.252,-1.087) (1.287,-1.092) (1.321,-1.097) (1.356,-1.101) (1.391,-1.105) (1.426,-1.110) (1.460,-1.113) (1.495,-1.117) (1.530,-1.121) (1.565,-1.124) (1.600,-1.127) (1.635,-1.130) (1.670,-1.133) (1.705,-1.136) (1.740,-1.138) (1.774,-1.141) (1.809,-1.143) (1.844,-1.145) (1.879,-1.147) (1.914,-1.149) (1.949,-1.151) (1.984,-1.153) (2.019,-1.155) (2.054,-1.157) (2.089,-1.158) (2.124,-1.160) (2.159,-1.161) (2.194,-1.163) (2.229,-1.164) (2.264,-1.165) (2.299,-1.166) (2.334,-1.168) (2.369,-1.169) (2.404,-1.170) (2.440,-1.171) (2.475,-1.172) (2.510,-1.173) (2.545,-1.174) (2.580,-1.174) (2.615,-1.175) (2.650,-1.176) (2.685,-1.177) (2.720,-1.178) (2.755,-1.178) (2.790,-1.179) (2.825,-1.180) (2.860,-1.180) (2.895,-1.181) (2.930,-1.181) (2.965,-1.182) (3.000,-1.182)};
\addplot[lineaD, flechaen=0.139, flechaen=0.861] coordinates {(-3.000,-1.381) (-2.965,-1.380) (-2.930,-1.380) (-2.895,-1.379) (-2.860,-1.379) (-2.825,-1.378) (-2.790,-1.377) (-2.755,-1.377) (-2.720,-1.376) (-2.685,-1.375) (-2.650,-1.374) (-2.615,-1.374) (-2.580,-1.373) (-2.545,-1.372) (-2.510,-1.371) (-2.475,-1.370) (-2.439,-1.369) (-2.404,-1.368) (-2.369,-1.367) (-2.334,-1.366) (-2.299,-1.365) (-2.264,-1.364) (-2.229,-1.362) (-2.194,-1.361) (-2.159,-1.360) (-2.124,-1.358) (-2.089,-1.357) (-2.054,-1.355) (-2.019,-1.354) (-1.984,-1.352) (-1.949,-1.350) (-1.914,-1.348) (-1.879,-1.347) (-1.844,-1.345) (-1.809,-1.343) (-1.774,-1.340) (-1.739,-1.338) (-1.704,-1.336) (-1.669,-1.333) (-1.635,-1.331) (-1.600,-1.328) (-1.565,-1.325) (-1.530,-1.322) (-1.495,-1.319) (-1.460,-1.316) (-1.425,-1.313) (-1.390,-1.309) (-1.355,-1.305) (-1.321,-1.302) (-1.286,-1.298) (-1.251,-1.293) (-1.216,-1.289) (-1.181,-1.284) (-1.147,-1.280) (-1.112,-1.274) (-1.077,-1.269) (-1.043,-1.264) (-1.008,-1.258) (-0.974,-1.252) (-0.939,-1.246) (-0.905,-1.239) (-0.871,-1.232) (-0.836,-1.225) (-0.802,-1.217) (-0.768,-1.209) (-0.734,-1.201) (-0.700,-1.193) (-0.666,-1.184) (-0.632,-1.174) (-0.599,-1.164) (-0.565,-1.154) (-0.532,-1.144) (-0.498,-1.133) (-0.465,-1.121) (-0.432,-1.109) (-0.399,-1.097) (-0.367,-1.085) (-0.334,-1.071) (-0.302,-1.058) (-0.270,-1.044) (-0.237,-1.031) (-0.205,-1.017) (-0.173,-1.003) (-0.140,-0.991) (-0.105,-0.990) (-0.070,-0.990) (-0.035,-0.990) (-0.000,-0.990) (0.035,-0.990) (0.070,-0.990) (0.105,-0.990) (0.140,-0.991) (0.173,-1.003) (0.205,-1.017) (0.237,-1.031) (0.270,-1.044) (0.302,-1.058) (0.334,-1.071) (0.367,-1.085) (0.399,-1.097) (0.432,-1.109) (0.465,-1.121) (0.498,-1.133) (0.532,-1.144) (0.565,-1.154) (0.599,-1.164) (0.632,-1.174) (0.666,-1.184) (0.700,-1.193) (0.734,-1.201) (0.768,-1.209) (0.802,-1.217) (0.836,-1.225) (0.871,-1.232) (0.905,-1.239) (0.939,-1.246) (0.974,-1.252) (1.008,-1.258) (1.043,-1.264) (1.077,-1.269) (1.112,-1.274) (1.147,-1.280) (1.181,-1.284) (1.216,-1.289) (1.251,-1.293) (1.286,-1.298) (1.321,-1.302) (1.355,-1.305) (1.390,-1.309) (1.425,-1.313) (1.460,-1.316) (1.495,-1.319) (1.530,-1.322) (1.565,-1.325) (1.600,-1.328) (1.635,-1.331) (1.669,-1.333) (1.704,-1.336) (1.739,-1.338) (1.774,-1.340) (1.809,-1.343) (1.844,-1.345) (1.879,-1.347) (1.914,-1.348) (1.949,-1.350) (1.984,-1.352) (2.019,-1.354) (2.054,-1.355) (2.089,-1.357) (2.124,-1.358) (2.159,-1.360) (2.194,-1.361) (2.229,-1.362) (2.264,-1.364) (2.299,-1.365) (2.334,-1.366) (2.369,-1.367) (2.404,-1.368) (2.439,-1.369) (2.475,-1.370) (2.510,-1.371) (2.545,-1.372) (2.580,-1.373) (2.615,-1.374) (2.650,-1.374) (2.685,-1.375) (2.720,-1.376) (2.755,-1.377) (2.790,-1.377) (2.825,-1.378) (2.860,-1.379) (2.895,-1.379) (2.930,-1.380) (2.965,-1.380) (3.000,-1.381)};
\addplot[lineaD, flechaen=0.141, flechaen=0.859] coordinates {(-3.000,-1.580) (-2.965,-1.579) (-2.930,-1.579) (-2.895,-1.578) (-2.860,-1.578) (-2.825,-1.577) (-2.790,-1.576) (-2.755,-1.576) (-2.720,-1.575) (-2.685,-1.574) (-2.650,-1.573) (-2.615,-1.573) (-2.580,-1.572) (-2.545,-1.571) (-2.510,-1.570) (-2.475,-1.569) (-2.440,-1.568) (-2.405,-1.567) (-2.370,-1.566) (-2.335,-1.565) (-2.300,-1.564) (-2.265,-1.563) (-2.230,-1.562) (-2.195,-1.561) (-2.160,-1.559) (-2.125,-1.558) (-2.090,-1.557) (-2.055,-1.555) (-2.020,-1.554) (-1.985,-1.552) (-1.950,-1.551) (-1.915,-1.549) (-1.880,-1.547) (-1.845,-1.546) (-1.810,-1.544) (-1.775,-1.542) (-1.740,-1.540) (-1.705,-1.538) (-1.670,-1.536) (-1.635,-1.533) (-1.600,-1.531) (-1.565,-1.529) (-1.530,-1.526) (-1.495,-1.524) (-1.460,-1.521) (-1.425,-1.518) (-1.390,-1.515) (-1.355,-1.512) (-1.320,-1.509) (-1.286,-1.506) (-1.251,-1.503) (-1.216,-1.499) (-1.181,-1.496) (-1.146,-1.492) (-1.111,-1.488) (-1.077,-1.484) (-1.042,-1.480) (-1.007,-1.476) (-0.972,-1.471) (-0.937,-1.467) (-0.903,-1.462) (-0.868,-1.458) (-0.833,-1.453) (-0.799,-1.448) (-0.764,-1.443) (-0.729,-1.438) (-0.695,-1.432) (-0.660,-1.427) (-0.626,-1.421) (-0.591,-1.416) (-0.556,-1.410) (-0.522,-1.405) (-0.487,-1.399) (-0.453,-1.394) (-0.418,-1.388) (-0.383,-1.383) (-0.349,-1.378) (-0.314,-1.373) (-0.279,-1.368) (-0.245,-1.364) (-0.210,-1.360) (-0.175,-1.357) (-0.140,-1.354) (-0.105,-1.352) (-0.070,-1.350) (-0.035,-1.349) (0.000,-1.348) (0.035,-1.349) (0.070,-1.350) (0.105,-1.352) (0.140,-1.354) (0.175,-1.357) (0.210,-1.360) (0.245,-1.364) (0.279,-1.368) (0.314,-1.373) (0.349,-1.378) (0.383,-1.383) (0.418,-1.388) (0.453,-1.394) (0.487,-1.399) (0.522,-1.405) (0.556,-1.410) (0.591,-1.416) (0.626,-1.421) (0.660,-1.427) (0.695,-1.432) (0.729,-1.438) (0.764,-1.443) (0.799,-1.448) (0.833,-1.453) (0.868,-1.458) (0.903,-1.462) (0.937,-1.467) (0.972,-1.471) (1.007,-1.476) (1.042,-1.480) (1.077,-1.484) (1.111,-1.488) (1.146,-1.492) (1.181,-1.496) (1.216,-1.499) (1.251,-1.503) (1.286,-1.506) (1.320,-1.509) (1.355,-1.512) (1.390,-1.515) (1.425,-1.518) (1.460,-1.521) (1.495,-1.524) (1.530,-1.526) (1.565,-1.529) (1.600,-1.531) (1.635,-1.533) (1.670,-1.536) (1.705,-1.538) (1.740,-1.540) (1.775,-1.542) (1.810,-1.544) (1.845,-1.546) (1.880,-1.547) (1.915,-1.549) (1.950,-1.551) (1.985,-1.552) (2.020,-1.554) (2.055,-1.555) (2.090,-1.557) (2.125,-1.558) (2.160,-1.559) (2.195,-1.561) (2.230,-1.562) (2.265,-1.563) (2.300,-1.564) (2.335,-1.565) (2.370,-1.566) (2.405,-1.567) (2.440,-1.568) (2.475,-1.569) (2.510,-1.570) (2.545,-1.571) (2.580,-1.572) (2.615,-1.573) (2.650,-1.573) (2.685,-1.574) (2.720,-1.575) (2.755,-1.576) (2.790,-1.576) (2.825,-1.577) (2.860,-1.578) (2.895,-1.578) (2.930,-1.579) (2.965,-1.579) (3.000,-1.580)};
\addplot[lineaD, flechaen=0.141, flechaen=0.859] coordinates {(-3.000,-1.779) (-2.965,-1.779) (-2.930,-1.778) (-2.895,-1.777) (-2.859,-1.777) (-2.824,-1.776) (-2.789,-1.776) (-2.754,-1.775) (-2.719,-1.774) (-2.684,-1.774) (-2.649,-1.773) (-2.613,-1.772) (-2.578,-1.771) (-2.543,-1.770) (-2.508,-1.770) (-2.473,-1.769) (-2.438,-1.768) (-2.403,-1.767) (-2.367,-1.766) (-2.332,-1.765) (-2.297,-1.764) (-2.262,-1.763) (-2.227,-1.762) (-2.192,-1.761) (-2.157,-1.760) (-2.121,-1.758) (-2.086,-1.757) (-2.051,-1.756) (-2.016,-1.755) (-1.981,-1.753) (-1.946,-1.752) (-1.911,-1.750) (-1.876,-1.749) (-1.841,-1.747) (-1.805,-1.746) (-1.770,-1.744) (-1.735,-1.742) (-1.700,-1.741) (-1.665,-1.739) (-1.630,-1.737) (-1.595,-1.735) (-1.560,-1.733) (-1.525,-1.731) (-1.489,-1.729) (-1.454,-1.727) (-1.419,-1.724) (-1.384,-1.722) (-1.349,-1.720) (-1.314,-1.717) (-1.279,-1.715) (-1.244,-1.712) (-1.209,-1.709) (-1.174,-1.707) (-1.139,-1.704) (-1.104,-1.701) (-1.069,-1.698) (-1.034,-1.695) (-0.999,-1.692) (-0.964,-1.689) (-0.929,-1.686) (-0.894,-1.682) (-0.859,-1.679) (-0.824,-1.676) (-0.789,-1.673) (-0.754,-1.669) (-0.719,-1.666) (-0.684,-1.663) (-0.649,-1.659) (-0.614,-1.656) (-0.579,-1.653) (-0.544,-1.649) (-0.509,-1.646) (-0.474,-1.643) (-0.439,-1.640) (-0.404,-1.637) (-0.369,-1.635) (-0.334,-1.632) (-0.299,-1.630) (-0.263,-1.628) (-0.228,-1.626) (-0.193,-1.624) (-0.158,-1.622) (-0.123,-1.621) (-0.088,-1.620) (-0.053,-1.620) (-0.018,-1.619) (0.018,-1.619) (0.053,-1.620) (0.088,-1.620) (0.123,-1.621) (0.158,-1.622) (0.193,-1.624) (0.228,-1.626) (0.263,-1.628) (0.299,-1.630) (0.334,-1.632) (0.369,-1.635) (0.404,-1.637) (0.439,-1.640) (0.474,-1.643) (0.509,-1.646) (0.544,-1.649) (0.579,-1.653) (0.614,-1.656) (0.649,-1.659) (0.684,-1.663) (0.719,-1.666) (0.754,-1.669) (0.789,-1.673) (0.824,-1.676) (0.859,-1.679) (0.894,-1.682) (0.929,-1.686) (0.964,-1.689) (0.999,-1.692) (1.034,-1.695) (1.069,-1.698) (1.104,-1.701) (1.139,-1.704) (1.174,-1.707) (1.209,-1.709) (1.244,-1.712) (1.279,-1.715) (1.314,-1.717) (1.349,-1.720) (1.384,-1.722) (1.419,-1.724) (1.454,-1.727) (1.489,-1.729) (1.525,-1.731) (1.560,-1.733) (1.595,-1.735) (1.630,-1.737) (1.665,-1.739) (1.700,-1.741) (1.735,-1.742) (1.770,-1.744) (1.805,-1.746) (1.841,-1.747) (1.876,-1.749) (1.911,-1.750) (1.946,-1.752) (1.981,-1.753) (2.016,-1.755) (2.051,-1.756) (2.086,-1.757) (2.121,-1.758) (2.157,-1.760) (2.192,-1.761) (2.227,-1.762) (2.262,-1.763) (2.297,-1.764) (2.332,-1.765) (2.367,-1.766) (2.403,-1.767) (2.438,-1.768) (2.473,-1.769) (2.508,-1.770) (2.543,-1.770) (2.578,-1.771) (2.613,-1.772) (2.649,-1.773) (2.684,-1.774) (2.719,-1.774) (2.754,-1.775) (2.789,-1.776) (2.824,-1.776) (2.859,-1.777) (2.895,-1.777) (2.930,-1.778) (2.965,-1.779) (3.000,-1.779)};
\addplot[lineaD, flechaen=0.142, flechaen=0.858] coordinates {(-3.000,-1.979) (-2.965,-1.978) (-2.930,-1.978) (-2.895,-1.977) (-2.860,-1.977) (-2.824,-1.976) (-2.789,-1.975) (-2.754,-1.975) (-2.719,-1.974) (-2.684,-1.973) (-2.649,-1.973) (-2.614,-1.972) (-2.579,-1.971) (-2.544,-1.971) (-2.508,-1.970) (-2.473,-1.969) (-2.438,-1.968) (-2.403,-1.967) (-2.368,-1.967) (-2.333,-1.966) (-2.298,-1.965) (-2.263,-1.964) (-2.228,-1.963) (-2.192,-1.962) (-2.157,-1.961) (-2.122,-1.960) (-2.087,-1.959) (-2.052,-1.958) (-2.017,-1.956) (-1.982,-1.955) (-1.947,-1.954) (-1.912,-1.953) (-1.877,-1.951) (-1.841,-1.950) (-1.806,-1.949) (-1.771,-1.947) (-1.736,-1.946) (-1.701,-1.944) (-1.666,-1.943) (-1.631,-1.941) (-1.596,-1.940) (-1.561,-1.938) (-1.526,-1.936) (-1.491,-1.935) (-1.455,-1.933) (-1.420,-1.931) (-1.385,-1.929) (-1.350,-1.927) (-1.315,-1.926) (-1.280,-1.924) (-1.245,-1.922) (-1.210,-1.920) (-1.175,-1.917) (-1.140,-1.915) (-1.105,-1.913) (-1.070,-1.911) (-1.035,-1.909) (-1.000,-1.907) (-0.965,-1.905) (-0.930,-1.902) (-0.895,-1.900) (-0.859,-1.898) (-0.824,-1.896) (-0.789,-1.893) (-0.754,-1.891) (-0.719,-1.889) (-0.684,-1.887) (-0.649,-1.885) (-0.614,-1.882) (-0.579,-1.880) (-0.544,-1.878) (-0.509,-1.876) (-0.474,-1.875) (-0.439,-1.873) (-0.404,-1.871) (-0.369,-1.869) (-0.334,-1.868) (-0.298,-1.867) (-0.263,-1.865) (-0.228,-1.864) (-0.193,-1.863) (-0.158,-1.862) (-0.123,-1.862) (-0.088,-1.861) (-0.053,-1.861) (-0.018,-1.861) (0.018,-1.861) (0.053,-1.861) (0.088,-1.861) (0.123,-1.862) (0.158,-1.862) (0.193,-1.863) (0.228,-1.864) (0.263,-1.865) (0.298,-1.867) (0.334,-1.868) (0.369,-1.869) (0.404,-1.871) (0.439,-1.873) (0.474,-1.875) (0.509,-1.876) (0.544,-1.878) (0.579,-1.880) (0.614,-1.882) (0.649,-1.885) (0.684,-1.887) (0.719,-1.889) (0.754,-1.891) (0.789,-1.893) (0.824,-1.896) (0.859,-1.898) (0.895,-1.900) (0.930,-1.902) (0.965,-1.905) (1.000,-1.907) (1.035,-1.909) (1.070,-1.911) (1.105,-1.913) (1.140,-1.915) (1.175,-1.917) (1.210,-1.920) (1.245,-1.922) (1.280,-1.924) (1.315,-1.926) (1.350,-1.927) (1.385,-1.929) (1.420,-1.931) (1.455,-1.933) (1.491,-1.935) (1.526,-1.936) (1.561,-1.938) (1.596,-1.940) (1.631,-1.941) (1.666,-1.943) (1.701,-1.944) (1.736,-1.946) (1.771,-1.947) (1.806,-1.949) (1.841,-1.950) (1.877,-1.951) (1.912,-1.953) (1.947,-1.954) (1.982,-1.955) (2.017,-1.956) (2.052,-1.958) (2.087,-1.959) (2.122,-1.960) (2.157,-1.961) (2.192,-1.962) (2.228,-1.963) (2.263,-1.964) (2.298,-1.965) (2.333,-1.966) (2.368,-1.967) (2.403,-1.967) (2.438,-1.968) (2.473,-1.969) (2.508,-1.970) (2.544,-1.971) (2.579,-1.971) (2.614,-1.972) (2.649,-1.973) (2.684,-1.973) (2.719,-1.974) (2.754,-1.975) (2.789,-1.975) (2.824,-1.976) (2.860,-1.977) (2.895,-1.977) (2.930,-1.978) (2.965,-1.978) (3.000,-1.979)};
\addplot[lineaD, flechaen=0.142, flechaen=0.858] coordinates {(-3.000,-2.179) (-2.965,-2.178) (-2.930,-2.178) (-2.895,-2.177) (-2.860,-2.177) (-2.824,-2.176) (-2.789,-2.176) (-2.754,-2.175) (-2.719,-2.174) (-2.684,-2.174) (-2.649,-2.173) (-2.614,-2.173) (-2.579,-2.172) (-2.544,-2.171) (-2.509,-2.170) (-2.473,-2.170) (-2.438,-2.169) (-2.403,-2.168) (-2.368,-2.167) (-2.333,-2.167) (-2.298,-2.166) (-2.263,-2.165) (-2.228,-2.164) (-2.193,-2.163) (-2.158,-2.162) (-2.123,-2.161) (-2.087,-2.161) (-2.052,-2.160) (-2.017,-2.159) (-1.982,-2.158) (-1.947,-2.156) (-1.912,-2.155) (-1.877,-2.154) (-1.842,-2.153) (-1.807,-2.152) (-1.772,-2.151) (-1.737,-2.150) (-1.701,-2.148) (-1.666,-2.147) (-1.631,-2.146) (-1.596,-2.145) (-1.561,-2.143) (-1.526,-2.142) (-1.491,-2.141) (-1.456,-2.139) (-1.421,-2.138) (-1.386,-2.136) (-1.351,-2.135) (-1.316,-2.133) (-1.281,-2.132) (-1.245,-2.130) (-1.210,-2.129) (-1.175,-2.127) (-1.140,-2.126) (-1.105,-2.124) (-1.070,-2.122) (-1.035,-2.121) (-1.000,-2.119) (-0.965,-2.118) (-0.930,-2.116) (-0.895,-2.114) (-0.860,-2.113) (-0.825,-2.111) (-0.790,-2.110) (-0.754,-2.108) (-0.719,-2.107) (-0.684,-2.105) (-0.649,-2.104) (-0.614,-2.102) (-0.579,-2.101) (-0.544,-2.100) (-0.509,-2.098) (-0.474,-2.097) (-0.439,-2.096) (-0.404,-2.095) (-0.369,-2.094) (-0.333,-2.093) (-0.298,-2.092) (-0.263,-2.091) (-0.228,-2.090) (-0.193,-2.090) (-0.158,-2.089) (-0.123,-2.089) (-0.088,-2.089) (-0.053,-2.088) (-0.018,-2.088) (0.018,-2.088) (0.053,-2.088) (0.088,-2.089) (0.123,-2.089) (0.158,-2.089) (0.193,-2.090) (0.228,-2.090) (0.263,-2.091) (0.298,-2.092) (0.333,-2.093) (0.369,-2.094) (0.404,-2.095) (0.439,-2.096) (0.474,-2.097) (0.509,-2.098) (0.544,-2.100) (0.579,-2.101) (0.614,-2.102) (0.649,-2.104) (0.684,-2.105) (0.719,-2.107) (0.754,-2.108) (0.790,-2.110) (0.825,-2.111) (0.860,-2.113) (0.895,-2.114) (0.930,-2.116) (0.965,-2.118) (1.000,-2.119) (1.035,-2.121) (1.070,-2.122) (1.105,-2.124) (1.140,-2.126) (1.175,-2.127) (1.210,-2.129) (1.245,-2.130) (1.281,-2.132) (1.316,-2.133) (1.351,-2.135) (1.386,-2.136) (1.421,-2.138) (1.456,-2.139) (1.491,-2.141) (1.526,-2.142) (1.561,-2.143) (1.596,-2.145) (1.631,-2.146) (1.666,-2.147) (1.701,-2.148) (1.737,-2.150) (1.772,-2.151) (1.807,-2.152) (1.842,-2.153) (1.877,-2.154) (1.912,-2.155) (1.947,-2.156) (1.982,-2.158) (2.017,-2.159) (2.052,-2.160) (2.087,-2.161) (2.123,-2.161) (2.158,-2.162) (2.193,-2.163) (2.228,-2.164) (2.263,-2.165) (2.298,-2.166) (2.333,-2.167) (2.368,-2.167) (2.403,-2.168) (2.438,-2.169) (2.473,-2.170) (2.509,-2.170) (2.544,-2.171) (2.579,-2.172) (2.614,-2.173) (2.649,-2.173) (2.684,-2.174) (2.719,-2.174) (2.754,-2.175) (2.789,-2.176) (2.824,-2.176) (2.860,-2.177) (2.895,-2.177) (2.930,-2.178) (2.965,-2.178) (3.000,-2.179)};
\addplot[lineaD, flechaen=0.142, flechaen=0.858] coordinates {(-3.000,0.000) (3.000,0.000)};

  \draw[texto, line width=0.6pt] (axis cs:0,0) circle[radius=1];
\end{axis}
\end{tikzpicture}
\caption{Líneas de $\vD$}
\end{subfigure}
\caption{Esfera dieléctrica ($\er=4$, sombreada) en un campo uniforme $\vE_0$ dirigido de izquierda a
derecha; trazos: equipotenciales. Solución exacta de la \cref{ej:esfera-dielectrica}. (a) El campo
interior es uniforme y más débil, $\vE_{\mathrm{int}}=3\vE_0/(\er+2)$: parte de las líneas de $\vE$ termina en la carga
ligada negativa del hemisferio izquierdo y renace en la positiva del derecho. (b) Las líneas de $\vD$, en
cambio, son continuas (no hay carga libre) y se concentran dentro de la esfera, donde
$D_{\mathrm{int}}=3\er\eo E_0/(\er+2)>\eo E_0$. Ambas familias se refractan en la superficie según la \cref{fig:refraccion}.}
\label{fig:esfera-dielectrica}
\end{figure}


% =============================================================================
\section{Cavidades en un dieléctrico}\label{sec:cavidades}
% =============================================================================

Dentro de un dieléctrico el campo macroscópico $\vE$ es un promedio: no podemos introducir una carga de
prueba entre las moléculas para medirlo. Lo que sí podemos hacer es practicar una pequeña cavidad y medir
el campo en su interior. La pregunta es si ese campo coincide con el del medio, y la respuesta es que
depende de la forma de la cavidad. Dos formas extremas, la aguja y la lenteja, dan respectivamente $\vE$ y
$\vD/\eo$, y con ello una definición operacional de ambos campos~\cite{DFA-teoria,Landau}.

\subsection{Planteamiento: el principio de superposición}\label{sec:cav-superposicion}

Sea un dieléctrico con campos $\vE_d$, $\vD_d$ y $\vP_d$, uniformes a la escala de la cavidad que vamos a practicar, y sea
$V_c$ la región de la cavidad, muy pequeña frente a las distancias en que varían los campos.

\begin{proposicion}[Superposición para cavidades]\label{prop:cav-superposicion}
Si la polarización del medio fuera de la cavidad no se altera al practicarla, el campo en la cavidad es
\begin{equation}\label{eq:cav-superposicion}
  \vEcav=\vE_d+\vE_{-\vP_d},
\end{equation}
donde $\vE_{-\vP_d}$ es el campo que crearía un cuerpo con la forma de la cavidad y polarización uniforme $-\vP_d$.
\end{proposicion}

\begin{proof}
La polarización del medio con cavidad es $\vP=\vP_d-\vP_d\,\mathbb 1_{V_c}$, donde $\mathbb 1_{V_c}$ vale uno en la cavidad y cero fuera.
El potencial \eqref{eq:phi-P} depende linealmente de $\vP$, así que el campo es la suma del creado por el medio macizo,
$\vE_d$, y del creado por la polarización $-\vP_d$ restringida a $V_c$ (\cref{fig:cavidades}\,(a)). Las cargas libres, si las hay,
no cambian.
\end{proof}

\begin{observacion}[Sobre la hipótesis]\label{obs:cav-hipotesis}
En un electrete la hipótesis es exacta: su polarización no depende del campo. En un medio \LHI{} la cavidad modifica
ligeramente el campo y la polarización en sus proximidades, y la \cref{prop:cav-superposicion} es solo aproximada. Los
resultados de las cavidades acicular y lenticular, sin embargo, son exactos en el límite de forma extrema, porque se
deducen directamente de las condiciones de frontera, como veremos.
\end{observacion}

% Figura: cavidades en un dieléctrico. Principio de superposición, cavidad
% acicular y cavidad lenticular, con los perfiles axiales que las justifican.
\begin{figure}[t]
\centering
\begin{subfigure}[b]{\textwidth}\centering
\begin{tikzpicture}[scale=0.92]
  \tikzset{bloque/.style={fill=dielec, draw=texto!45!fondo, line width=0.4pt, rounded corners=2pt}}
  \tikzset{aguja/.style={fill=fondo, draw=texto, line width=0.5pt, rounded corners=1.4pt}}
  % dieléctrico con cavidad
  \draw[bloque] (0,0) rectangle (3.0,1.9);
  \draw[aguja] (0.95,0.7) rectangle (2.05,0.86);
  \draw[vecP] (1.05,1.38) -- (1.95,1.38) node[above, midway, etiq] {$\vP_d$};
  \node[etiq] at (3.55,0.95) {$=$};
  % dieléctrico macizo
  \draw[bloque] (4.1,0) rectangle (7.1,1.9);
  \draw[vecP] (5.15,1.38) -- (6.05,1.38) node[above, midway, etiq] {$\vP_d$};
  \node[etiq] at (7.65,0.95) {$+$};
  % «anticavidad»: cilindro con -P_d
  \draw[aguja, fill=dielecB] (8.25,0.7) rectangle (9.75,0.86);
  \draw[vecP] (9.55,0.78) -- (8.45,0.78);
  \node[etiq, anchor=north] at (9.0,0.6) {$-\vP_d$};
\end{tikzpicture}
\caption{Superposición: dieléctrico con cavidad $=$ dieléctrico macizo $+$ cuerpo con la forma de la cavidad y polarización $-\vP_d$}
\end{subfigure}\\[1.0em]
\begin{subfigure}[b]{0.235\textwidth}\centering
\begin{tikzpicture}[scale=0.86]
  \fill[dielec, rounded corners=2pt] (0,0) rectangle (3.0,2.4);
  \draw[vecE] (0.25,2.0) -- (1.15,2.0) node[above, midway, etiq] {$\vE_d$};
  \draw[vecP] (1.8,2.0) -- (2.7,2.0) node[above, midway, etiq] {$\vP_d$};
  \draw[fill=fondo, draw=texto, line width=0.5pt, rounded corners=1.4pt] (0.55,0.62) rectangle (2.45,0.8);
  \draw[vecn] (1.5,0.8) -- (1.5,1.42) node[right, etiq] {$\nn$};
  \draw[vecn] (1.5,0.8) -- (2.1,0.8);
  \node[etiq, anchor=south] at (2.05,0.8) {$\tu$};
  \node[etiq, anchor=north] at (1.5,0.55) {$\vEcav=\vE_d$};
\end{tikzpicture}
\caption{Acicular}
\end{subfigure}\hfill
\begin{subfigure}[b]{0.24\textwidth}\centering
\begin{tikzpicture}
\begin{axis}[perfil, width=4.25cm, height=3.6cm, ymin=-0.75, ymax=1.15, ytick={-0.5,0.5,1}]
  \curvaPeje \curvaEeje{0.1}
  \node[etiqf, font=\tiny, anchor=north west] at (rel axis cs:0.02,0.98) {$a/\ell=0{,}1$};
\end{axis}
\end{tikzpicture}
\caption{Cilindro largo}
\end{subfigure}\hfill
\begin{subfigure}[b]{0.235\textwidth}\centering
\begin{tikzpicture}[scale=0.86]
  \fill[dielec, rounded corners=2pt] (0,0) rectangle (3.0,2.4);
  \draw[vecE] (0.25,2.0) -- (1.15,2.0) node[above, midway, etiq] {$\vE_d$};
  \draw[vecP] (1.8,2.0) -- (2.7,2.0) node[above, midway, etiq] {$\vP_d$};
  \draw[fill=fondo, draw=texto, line width=0.5pt] (1.1,0.95) ellipse (0.09 and 0.75);
  \draw[vecn] (1.19,0.95) -- (1.8,0.95) node[above, etiq] {$\nn$};
  \draw[vecn] (1.19,0.95) -- (1.19,1.5) node[right, etiq] {$\tu$};
  \node[etiq, anchor=west] at (1.3,0.42) {$\vDcav=\vD_d$};
\end{tikzpicture}
\caption{Lenticular}
\end{subfigure}\hfill
\begin{subfigure}[b]{0.24\textwidth}\centering
\begin{tikzpicture}
\begin{axis}[perfil, width=4.25cm, height=3.6cm, ymin=-1.15, ymax=1.15, ytick={-1,1}]
  \curvaPeje \curvaEeje{50}
  \node[etiqf, font=\tiny, anchor=north west] at (rel axis cs:0.02,0.98) {$a/\ell=50$};
\end{axis}
\end{tikzpicture}
\caption{Disco delgado}
\end{subfigure}
\caption{Cavidades en un dieléctrico polarizado ($\vE_d$, $\vP_d$, $\vD_d$ son los campos del medio lejos de la
cavidad). (a) Por superposición, basta sumar al campo del medio el de un cuerpo con la forma de la
cavidad y polarización $-\vP_d$. (b) En una cavidad acicular paralela al campo domina la superficie
lateral y se conserva la componente tangencial: $\vEcav=\vE_d$. (c) Perfil axial de $E$ (línea continua,
unidades $P/\eo$) y de $P$ (trazos) de un cilindro polarizado largo: dentro, lejos de los extremos, $E\simeq0$.
(d) En una cavidad lenticular perpendicular al campo dominan las caras y se conserva la componente
normal de $\vD$: $\vDcav=\vD_d$. (e) En un disco delgado, $E\simeq-P/\eo$ dentro; con $-\vP_d$ esto aporta $+\vP_d/\eo$.}
\label{fig:cavidades}
\end{figure}


\subsection{Cavidad acicular}\label{sec:acicular}

\begin{definicion}[Cavidad acicular]\label{def:acicular}
Una \emph{cavidad acicular} es un cilindro de radio $a$ y longitud $2\ell$ con $a\ll\ell$ (una aguja), cuyo eje es paralelo a $\vE_d$.
\end{definicion}

\begin{proposicion}[Campo en una cavidad acicular]\label{prop:acicular}
En el centro de una cavidad acicular, en el límite $a/\ell\to0$,
\begin{equation}\label{eq:acicular}
  \boxed{\;\vEcav=\vE_d\;},\qquad \vDcav=\eo\vE_d=\vD_d-\vP_d .
\end{equation}
\end{proposicion}

\begin{proof}
\emph{Por condiciones de frontera.} La superficie de la aguja es casi toda lateral, y en ella $\vE_d$ es tangente. La componente
tangencial de $\vE$ es continua (\cref{teo:frontera}), luego justo dentro de la pared lateral $\vEcav=\vE_d$. Las caras de los
extremos tienen área $\pi a^2$, despreciable, y su influencia en el centro desaparece en el límite.

\emph{Por superposición.} El cuerpo de la \cref{prop:cav-superposicion} es un cilindro con polarización axial $-\vP_d$, que equivale a dos
discos de radio $a$ con densidades $\pm P_d$ en $z=\mp\ell$. En el centro, por el resultado del \cref{prob:cilindro-polarizado} con
$P\to-P_d$ y $z=0$,
\[
  \vE_{-\vP_d}(0)=\frac{\vP_d}{\eo}\left(1-\frac{\ell}{\sqrt{a^2+\ell^2}}\right)=\frac{\vP_d}{\eo}\left(\frac{a^2}{2\ell^2}+\mathcal O\big(a^4/\ell^4\big)\right)
  \xrightarrow{\;a/\ell\to0\;}\vcero ,
\]
como muestra el perfil del cilindro largo, en el que $E\simeq0$ lejos de los extremos (\cref{fig:cavidades}\,(c)).
\end{proof}

\subsection{Cavidad lenticular}\label{sec:lenticular}

\begin{definicion}[Cavidad lenticular]\label{def:lenticular}
Una \emph{cavidad lenticular} es un cilindro de radio $a$ y espesor $2\ell$ con $\ell\ll a$ (una lenteja o disco), cuyas caras son
perpendiculares a $\vE_d$.
\end{definicion}

\begin{proposicion}[Campo en una cavidad lenticular]\label{prop:lenticular}
En el centro de una cavidad lenticular, en el límite $\ell/a\to0$,
\begin{equation}\label{eq:lenticular}
  \boxed{\;\vDcav=\vD_d\;},\qquad \vEcav=\frac{\vD_d}{\eo}=\vE_d+\frac{\vP_d}{\eo} .
\end{equation}
\end{proposicion}

\begin{proof}
\emph{Por condiciones de frontera.} Ahora la superficie es casi toda cara, perpendicular a $\vE_d$ y por tanto a $\vD_d$. No hay carga libre
en ella, y la componente normal de $\vD$ es continua (\cref{teo:frontera}): justo dentro de la cara, $\vDcav=\vD_d$. El borde lateral, de
área $4\pi a\ell$, es despreciable frente a las caras en el límite.

\emph{Por superposición.} El cuerpo es ahora un disco delgado con $-\vP_d$. En su centro, con la misma fórmula y $\ell/a\to0$,
\[
  \vE_{-\vP_d}(0)=\frac{\vP_d}{\eo}\left(1-\frac{\ell}{\sqrt{a^2+\ell^2}}\right)\xrightarrow{\;\ell/a\to0\;}\frac{\vP_d}{\eo},
\]
el campo de un condensador plano con densidades $\pm P_d$, como en el perfil del disco delgado (\cref{fig:cavidades}\,(e)). Sumando
$\vE_d$ se obtiene \eqref{eq:lenticular}.
\end{proof}

\subsection{Interpretación: definición operacional de $\vE$ y $\vD$}\label{sec:cav-interpretacion}

Las dos proposiciones anteriores, debidas esencialmente a Kelvin, permiten definir los campos macroscópicos mediante
medidas: \emph{$\vE$ es la fuerza por unidad de carga que actúa sobre una carga de prueba en una cavidad acicular
paralela al campo, y $\vD/\eo$ es la que actúa en una cavidad lenticular perpendicular a él}. La diferencia entre ambas
es $\vP/\eo$: es el campo de las cargas ligadas de las caras del disco, que la aguja evita y la lenteja no.

\begin{observacion}[La cavidad esférica y el campo de Lorentz]\label{obs:lorentz}
Para una cavidad esférica, el cuerpo de la \cref{prop:cav-superposicion} es una esfera con $-\vP_d$, que crea en su interior el campo
uniforme $+\vP_d/3\eo$ (\cref{prob:esfera-polarizada}). Por tanto,
\[
  \vEcav=\vE_d+\frac{\vP_d}{3\eo},
\]
a medio camino entre la aguja y la lenteja. Es el llamado \emph{campo local de Lorentz}: una estimación del campo que actúa
sobre una molécula, que «ve» a sus vecinas como un medio continuo a partir de cierta distancia. Si cada molécula
responde con $\vp=\alpha\vE_{\mathrm{loc}}$ y hay $N$ por unidad de volumen, $\vP=N\alpha\,(\vE+\vP/3\eo)$, y despejando $\chi=P/(\eo E)$ se obtiene
la relación de Clausius--Mossotti~\cite{Griffiths,Jackson},
\[
  \frac{\er-1}{\er+2}=\frac{N\alpha}{3\eo},
\]
que liga la polarizabilidad microscópica de la \cref{sec:microscopico} con la permitividad macroscópica de la
\cref{sec:lineales}.
\end{observacion}

\begin{observacion}[Cavidad elipsoidal]\label{obs:elipsoide}
Más en general, una cavidad con forma de elipsoide de revolución cuyo eje es paralelo a $\vE_d$ tiene en su interior un campo
uniforme $\vEcav=\vE_d+N\,\vP_d/\eo$, donde el \emph{factor de despolarización} $N\in[0,1]$ solo depende de la relación de aspecto~\cite{Osborn,Landau}.
Vale $N\to0$ para la aguja, $N=1/3$ para la esfera y $N\to1$ para la lenteja: las tres cavidades de esta sección son puntos de
una misma curva (\cref{fig:despolarizacion}).
\end{observacion}

% Figura: factor de despolarización N de un esferoide de revolución (eje paralelo
% al campo) en función de su relación de aspecto m = c/a.  Datos: scripts/lineas_campo.py
\begin{figure}[t]
\centering
\begin{tikzpicture}
\begin{axis}[width=10.5cm, height=5.4cm, xmode=log, axis lines=left, xmin=0.01, xmax=100,
  ymin=0, ymax=1.08, grid=major, xlabel={relación de aspecto $m=c/a$ (eje paralelo / eje transversal)},
  ylabel={$N(m)$}, ytick={0,0.3333,1}, yticklabels={$0$,$\tfrac13$,$1$},
  xtick={0.01,0.1,1,10,100}, xticklabels={$10^{-2}$,$10^{-1}$,$1$,$10$,$10^{2}$}, clip=false]
  \addplot[curvaE, line width=1.1pt] table[x=m, y=N] {datos/despolarizacion.dat};
  \addplot[only marks, mark=*, mark size=1.8pt, color=acentoC] coordinates {(1,0.3333)};
  % iconos de forma
  \draw[texto, fill=dielecB, line width=0.4pt] (axis cs:0.016,0.42) ellipse [x radius=0.05cm, y radius=0.36cm];
  \node[etiqt, anchor=west] at (axis cs:0.022,0.42) {lenteja: $N\to1$,\; $\vDcav=\vD_d$};
  \draw[texto, fill=dielecB, line width=0.4pt] (axis cs:1,0.58) circle (0.17cm);
  \node[etiqt, anchor=west] at (axis cs:1.35,0.58) {esfera: $N=\tfrac13$ (campo de Lorentz)};
  \draw[texto, fill=dielecB, line width=0.4pt] (axis cs:30,0.25) ellipse [x radius=0.42cm, y radius=0.05cm];
  \node[etiqt, anchor=south] at (axis cs:30,0.29) {aguja: $N\to0$, $\vEcav=\vE_d$};
  \draw[vecE] (axis cs:0.013,0.15) -- (axis cs:0.04,0.15);
  \node[etiqt, anchor=west] at (axis cs:0.042,0.15) {$\vE_d$};
\end{axis}
\end{tikzpicture}
\caption{Factor de despolarización $N$ de una cavidad con forma de esferoide de revolución cuyo eje,
de semieje $c$, es paralelo a $\vE_d$ (el otro semieje es $a$). El campo en la cavidad es
$\vEcav=\vE_d+N\,\vP_d/\eo$, que interpola entre los dos casos límite del texto: la aguja ($m\to\infty$) y la
lenteja ($m\to0$). La esfera da el campo local de Lorentz, $\vE_d+\vP_d/3\eo$.}
\label{fig:despolarizacion}
\end{figure}


% =============================================================================
\section{Seminarios prácticos}\label{sec:seminarios}
% =============================================================================

Esta sección resuelve con detalle los doce problemas de la relación del tema~\cite{DFA-problemas}. Cada
problema sigue la estructura de los temas anteriores: el enunciado, la idea que guía la resolución, la
resolución paso a paso y un comentario sobre la interpretación del resultado y sus comprobaciones. La guía
de métodos que acompaña a estos apuntes~\cite{Guia-T4} recoge los mismos problemas en formato de ficha.

\subsection{Cómo plantear un problema de dieléctricos}\label{sec:plantear}

Todo problema de dieléctricos empieza por identificar qué es dato: la polarización (electretes), la carga
libre (conductores cargados en medios lineales) o los potenciales (baterías). La \cref{tab:traduccion}
traduce cada situación a la herramienta adecuada, y la \cref{tab:decision}, el método que conviene.

\begin{table}[H]
\centering
\caption{Traducción de la situación física a ecuaciones.}\label{tab:traduccion}
\small
\begin{tabularx}{\textwidth}{@{}>{\raggedright\arraybackslash}p{4.3cm}>{\raggedright\arraybackslash}X@{}}
\toprule
Situación & Qué se sabe y qué se usa\\
\midrule
Electrete ($\vP$ dada) & $\rhoP=-\divg\vP$, $\sigmaP=\vP\cdot\nn$; si $\rhof=0$ y $\vP$ es irrotacional, $\vD=\vcero$\\
Carga libre en medio \LHI{} & Gauss para $\vD$ con la carga libre; después $\vE=\vD/\varepsilon$ y $\vP=(\varepsilon-\eo)\vE$\\
Interfase sin carga libre & $D_n$ continua y $\vE_t$ continuo\\
Conductor en equilibrio & $\vE=\vD=\vcero$ dentro; $D_n=\sigmaf$ en su superficie; volumen equipotencial\\
Batería $V_0$ entre placas & $\int\vE\cdot\dl=V_0$ entre placas a lo largo de cualquier camino\\
Placas unidas por un cable & $\phi_1=\phi_2$, es decir, $\int\vE\cdot\dl=0$ de una placa a otra\\
Condensador aislado & la carga libre de cada placa se conserva\\
\bottomrule
\end{tabularx}
\end{table}

\begin{table}[H]
\centering
\caption{Mapa de decisión: qué método conviene según lo que pide el enunciado.}\label{tab:decision}
\small
\begin{tabularx}{\textwidth}{@{}>{\raggedright\arraybackslash}X>{\raggedright\arraybackslash}p{4.6cm}l@{}}
\toprule
Lo que pide el enunciado & Método & Problemas\\
\midrule
$\vP$ dada, sin simetría que permita usar Gauss & M1. Cargas ligadas y superposición & 4.1, 4.4, 4.5\\
$\vP$ dada con simetría esférica o cilíndrica & M2. $\vD$ por Gauss, $\vE=(\vD-\vP)/\eo$ & 4.6, 4.7\\
Carga libre en medios \LHI{} con simetría & M3. Gauss para $\vD$ & 4.2, 4.3, 4.8\\
Condensadores parcialmente llenos & M4. Serie y paralelo & 4.8, 4.10--4.12\\
Electretes entre placas conectadas & M5. $D$ continuo y $\oint\vE\cdot\dl=0$ & 4.9\\
\bottomrule
\end{tabularx}
\end{table}

\begin{metodo}{Los cinco métodos}
\begin{description}[leftmargin=0pt, itemsep=0.35em, font=\normalfont\bfseries\color{acentoA}]
  \item[M1. Cargas ligadas y superposición.] Calcular $\rhoP$ y $\sigmaP$ con el \cref{teo:densidades}; comprobar $\QP=0$; obtener $\phi$ o $\vE$
        integrando como para cargas libres (anillos, discos, capas). Útil para campos sobre un eje.
  \item[M2. $\vD$ por Gauss con $\vP$ dada.] Si hay simetría, Gauss para $\vD$ con la carga libre encerrada da $\vD$ directamente
        (a menudo $\vD=\vcero$, \cref{prop:D-nulo}); luego $\vE=(\vD-\vP)/\eo$. Comprobar con Gauss para $\vE$ y las cargas ligadas.
  \item[M3. Gauss para $\vD$ en medios lineales.] (1) $\vD$ con la carga libre; (2) $\vE=\vD/\varepsilon$ en cada medio; (3) $\vP=\vD-\eo\vE$;
        (4) $\rhoP$, $\sigmaP$; (5) potencial y capacidad, $C=Q/\Delta\phi$.
  \item[M4. Serie y paralelo.] Interfase paralela a las placas: $\vD$ común, los campos se suman en $\Delta\phi$, capacidades
        en serie. Interfase perpendicular: $\vE$ común, las cargas se suman, capacidades en paralelo.
  \item[M5. Electrete entre placas conectadas.] $D$ continuo en las interfases sin carga libre, $\vE=(\vD-\vP)/\eo$ en cada capa, y la
        condición del circuito ($\int\vE\cdot\dl=0$ o $V_0$) fija $D$, que es la carga inducida en las placas.
\end{description}
\end{metodo}

% -----------------------------------------------------------------------------
\subsection{Problema 4.1. Esfera uniformemente polarizada}\label{sec:p41}

\begin{problema}[Esfera uniformemente polarizada]\label{prob:esfera-polarizada}
Sea una esfera de radio $a$ uniformemente polarizada con polarización $\vP=P\,\uv z$. Calcula el potencial y el campo eléctrico a lo
largo del eje $z$.
\end{problema}

\begin{resolucion}
\idea No hay simetría que permita usar Gauss ($\vD$ no es nulo: $\vP$ tiene componente tangencial en la superficie). Calculamos
las cargas ligadas e integramos su potencial sobre el eje (método M1).

\paso{Paso 1: densidades ligadas} Como $\vP$ es uniforme, $\rhoP=-\divg\vP=0$. En la superficie $\nn=\uv r$, y con $\theta'$ el ángulo polar del
punto fuente,
\[
  \sigmaP(\theta')=\vP\cdot\uv r=P\cos\theta' .
\]
Es positiva en el hemisferio $z>0$ y negativa en el otro (\cref{fig:p41}\,(a)); su integral es nula, como exige la \cref{prop:QP-nula}.

\paso{Paso 2: potencial en el eje} Para un punto $(0,0,z)$ con $z>0$, la distancia a un punto de la superficie es
$R=\sqrt{a^2+z^2-2az\cos\theta'}$ y $\dd S'=a^2\sen\theta'\,\dd\theta'\,\dd\varphi'$. Con $u=\cos\theta'$,
\[
  \phi(z)=\frac{1}{4\pi\eo}\oiint\frac{\sigmaP}{R}\,\dd S'=\frac{Pa^2}{2\eo}\int_{-1}^{1}\frac{u\,\dd u}{\sqrt{a^2+z^2-2azu}}\equiv\frac{Pa^2}{2\eo}\,I(z).
\]
El cambio $s=\sqrt{a^2+z^2-2azu}$, con $u=(a^2+z^2-s^2)/2az$ y $\dd u=-s\,\dd s/(az)$, da
\[
  I(z)=\frac{1}{2a^2z^2}\int_{\abs{a-z}}^{a+z}\big(a^2+z^2-s^2\big)\,\dd s
      =\frac{1}{2a^2z^2}\Big[(a^2+z^2)\,s-\frac{s^3}{3}\Big]_{\abs{a-z}}^{a+z}.
\]
Si $z>a$, el corchete vale $2a(a^2+z^2)-\frac13\big[(z+a)^3-(z-a)^3\big]=2a(a^2+z^2)-\frac13(6az^2+2a^3)=\frac43a^3$; si $0<z<a$, vale
$2z(a^2+z^2)-\frac13\big[(a+z)^3-(a-z)^3\big]=\frac43z^3$. Por tanto,
\[
  I(z)=\begin{cases}\dfrac{2a}{3z^2}, & z\geq a,\\[2mm] \dfrac{2z}{3a^2}, & 0\leq z\leq a.\end{cases}
\]
Como $\sigmaP$ es impar bajo la reflexión $z\to-z$, el potencial también lo es, $\phi(-z)=-\phi(z)$. En conclusión,
\begin{equation}\label{eq:p41-phi}
  \phi(z)=\begin{cases}\dfrac{Pa^3}{3\eo}\,\dfrac{z}{\abs{z}^3}, & \abs z\geq a,\\[3mm] \dfrac{P}{3\eo}\,z, & \abs z\leq a.\end{cases}
\end{equation}

\paso{Paso 3: campo en el eje} Por simetría de revolución, en el eje $\vE=E_z\,\uv z$ con $E_z=-\dd\phi/\dd z$:
\begin{equation}\label{eq:p41-E}
  \recuadro{E_z(z)=\frac{2Pa^3}{3\eo\abs{z}^3}\quad(\abs z>a),\qquad E_z(z)=-\frac{P}{3\eo}\quad(\abs z<a).}
\end{equation}
\end{resolucion}

\begin{observacion}[Interpretación y comprobaciones]\label{obs:p41}
(a) Fuera, \eqref{eq:p41-phi} es el potencial en el eje de un dipolo $\vp=\frac43\pi a^3\vP$ situado en el centro, $p/(4\pi\eo z^2)$: el dipolo
total de la esfera, como predice la \cref{obs:momento-ligado}. (b) Dentro, $\phi=Pz/3\eo$; de hecho, el potencial interior es
$\phi=\vP\cdot\vr/3\eo$ en todos los puntos, y el campo es \emph{uniforme} y opuesto a la polarización, $\vE=-\vP/3\eo$: es el \emph{campo
despolarizante} (\cref{fig:esfera-polarizada}). (c) En los polos, $E_z$ salta de $-P/3\eo$ a $2P/3\eo$, es decir, $P/\eo=\sigmaP(0)/\eo$, como exige
\eqref{eq:salto-En}. (d) Dentro, $\vD=\eo\vE+\vP=\frac23\vP\neq\vcero$ aunque no haya carga libre: $\vP$ tiene componente tangencial en la superficie.
Este resultado es la pieza clave de la esfera dieléctrica (\cref{ej:esfera-dielectrica}) y de la cavidad esférica (\cref{obs:lorentz}).
\end{observacion}

% Figura: problema 4.1. Esfera uniformemente polarizada: geometría y potencial y
% campo sobre el eje (unidades P a/eps0 y P/eps0).
\begin{figure}[t]
\centering
\begin{subfigure}[b]{0.4\textwidth}\centering
\begin{tikzpicture}[scale=1.05]
  \fill[dielec] (0,0) circle (1.2);
  \draw[borde] (0,0) circle (1.2);
  \draw[texto!70!fondo, line width=0.5pt, -{Stealth[length=2mm]}] (-1.75,0) -- (3.3,0) node[below, etiq] {$z$};
  \foreach \t in {-70,-35,0,35,70} {
    \node[cpos, fill=fondo, circle, inner sep=0.2pt] at (\t:1.42) {$+$};
    \node[cneg, fill=fondo, circle, inner sep=0.2pt] at (180-\t:1.42) {$-$};
  }
  \draw[vecP] (-0.55,-0.55) -- (0.55,-0.55) node[right, etiq] {$\vP$};
  % punto del eje y elemento de superficie
  \coordinate (Z) at (2.75,0);
  \fill[texto] (Z) circle (0.04);
  \node[etiq, anchor=south] at (Z) {$z$};
  \coordinate (Q) at (55:1.2);
  \fill[acentoB] (Q) circle (0.05);
  \draw[acentoB, line width=0.5pt] (0,0) -- (Q) node[pos=0.55, left, etiq] {$a$};
  \draw[acentoB, line width=0.5pt, dashed] (Q) -- (Z) node[pos=0.5, above right=-1pt, etiq] {$R$};
  \draw[tenue, line width=0.4pt] (0.42,0) arc (0:55:0.42);
  \node[etiqt] at (27:0.62) {$\theta'$};
\end{tikzpicture}
\caption{$\sigmaP=P\cos\theta'$, $\rhoP=0$}
\end{subfigure}\hfill
\begin{subfigure}[b]{0.56\textwidth}\centering
\begin{tikzpicture}
\begin{axis}[width=8.0cm, height=5.0cm, xmin=-3.2, xmax=3.2, ymin=-0.62, ymax=0.82,
  xlabel={$z/a$}, xtick={-3,-2,-1,0,1,2,3}, ytick={-0.3333,0,0.3333,0.6667}, yticklabels={$-\frac13$,$0$,$\frac13$,$\frac23$},
  grid=major, xlabel style={yshift=3pt}, clip=false,
  legend style={at={(0.02,0.98)}, anchor=north west}]
  \fill[dielec, opacity=0.6] (axis cs:-1,-0.62) rectangle (axis cs:1,0.82);
  \addplot[curvaV, domain=-3.2:-1, samples=80] {-1/(3*x^2)}; \addlegendentry{$\phi\,/\,(Pa/\eo)$}
  \addplot[curvaE, domain=-3.2:-1, samples=80] {2/(3*abs(x)^3)}; \addlegendentry{$E_z\,/\,(P/\eo)$}
  \addplot[curvaV, domain=-1:1, samples=2] {x/3};
  \addplot[curvaV, domain=1:3.2, samples=80] {1/(3*x^2)};
  \addplot[curvaE, domain=-1:1, samples=2] {-1/3};
  \addplot[curvaE, domain=1:3.2, samples=80] {2/(3*x^3)};
  \draw[equipot] (axis cs:-1,-0.62) -- (axis cs:-1,0.82) (axis cs:1,-0.62) -- (axis cs:1,0.82);
\end{axis}
\end{tikzpicture}
\caption{Potencial y campo sobre el eje}
\end{subfigure}
\caption{\Cref{prob:esfera-polarizada}. (a) La polarización uniforme no deja carga de volumen; en la
superficie, $\sigmaP=P\cos\theta'$ es positiva en el hemisferio hacia el que apunta $\vP$ y negativa en el opuesto.
(b) Sobre el eje (zona sombreada: interior), el potencial es continuo, lineal dentro y dipolar fuera; el
campo es uniforme dentro, $E_z=-P/3\eo$, y salta $\sigmaP/\eo=P/\eo$ en los polos.}
\label{fig:p41}
\end{figure}

% Figura (simulación): esfera uniformemente polarizada (problema 4.1).
% Datos: scripts/lineas_campo.py -> datos/esfera-polarizada-{E,D}.tex
\begin{figure}[t]
\centering
\pgfplotsset{esferap/.style={marco, width=7.4cm, xmin=-2.6, xmax=2.6, ymin=-2.0, ymax=2.0}}
\begin{subfigure}[b]{0.49\textwidth}\centering
\begin{tikzpicture}
\begin{axis}[esferap]
  \fill[dielec] (axis cs:0,0) circle[radius=1];
  %% Generado por scripts/lineas_campo.py -- no editar a mano
\addplot[equipot] coordinates {(-0.905,-0.429) (-0.926,-0.401) (-0.946,-0.372) (-0.964,-0.342) (-0.981,-0.311) (-0.996,-0.280) (-1.009,-0.247) (-1.021,-0.214) (-1.031,-0.180) (-1.039,-0.146) (-1.045,-0.112) (-1.050,-0.077) (-1.053,-0.042) (-1.054,-0.007) (-1.054,0.028) (-1.051,0.063) (-1.047,0.098) (-1.041,0.133) (-1.034,0.167) (-1.025,0.201) (-1.014,0.234) (-1.001,0.267) (-0.987,0.299) (-0.971,0.330) (-0.953,0.361) (-0.934,0.390) (-0.914,0.419) (-0.900,0.411) (-0.900,0.375) (-0.900,0.340) (-0.900,0.305) (-0.900,0.270) (-0.900,0.235) (-0.900,0.200) (-0.900,0.165) (-0.900,0.130) (-0.900,0.095) (-0.900,0.059) (-0.900,0.024) (-0.900,-0.011) (-0.900,-0.046) (-0.900,-0.081) (-0.900,-0.116) (-0.900,-0.151) (-0.900,-0.186) (-0.900,-0.221) (-0.900,-0.257) (-0.900,-0.292) (-0.900,-0.327) (-0.900,-0.362) (-0.900,-0.397) (-0.905,-0.429)};
\addplot[equipot] coordinates {(-0.726,-0.692) (-0.759,-0.678) (-0.790,-0.663) (-0.821,-0.646) (-0.852,-0.627) (-0.881,-0.607) (-0.908,-0.585) (-0.935,-0.562) (-0.961,-0.538) (-0.985,-0.512) (-1.008,-0.485) (-1.030,-0.457) (-1.050,-0.428) (-1.069,-0.398) (-1.087,-0.367) (-1.103,-0.336) (-1.117,-0.304) (-1.130,-0.271) (-1.142,-0.237) (-1.152,-0.203) (-1.160,-0.169) (-1.167,-0.134) (-1.172,-0.099) (-1.176,-0.064) (-1.178,-0.029) (-1.178,0.007) (-1.177,0.042) (-1.175,0.077) (-1.170,0.112) (-1.165,0.147) (-1.157,0.182) (-1.148,0.216) (-1.138,0.250) (-1.126,0.283) (-1.112,0.316) (-1.097,0.348) (-1.080,0.379) (-1.062,0.409) (-1.043,0.439) (-1.022,0.468) (-1.000,0.495) (-0.976,0.522) (-0.951,0.547) (-0.925,0.571) (-0.898,0.593) (-0.870,0.615) (-0.840,0.634) (-0.810,0.652) (-0.779,0.669) (-0.746,0.683) (-0.720,0.683) (-0.720,0.648) (-0.720,0.613) (-0.720,0.577) (-0.720,0.542) (-0.720,0.506) (-0.720,0.471) (-0.720,0.436) (-0.720,0.400) (-0.720,0.365) (-0.720,0.330) (-0.720,0.294) (-0.720,0.259) (-0.720,0.223) (-0.720,0.188) (-0.720,0.153) (-0.720,0.117) (-0.720,0.082) (-0.720,0.046) (-0.720,0.011) (-0.720,-0.024) (-0.720,-0.060) (-0.720,-0.095) (-0.720,-0.130) (-0.720,-0.166) (-0.720,-0.201) (-0.720,-0.237) (-0.720,-0.272) (-0.720,-0.307) (-0.720,-0.343) (-0.720,-0.378) (-0.720,-0.414) (-0.720,-0.449) (-0.720,-0.484) (-0.720,-0.520) (-0.720,-0.555) (-0.720,-0.590) (-0.720,-0.626) (-0.720,-0.661) (-0.726,-0.692)};
\addplot[equipot] coordinates {(-0.662,-0.841) (-0.697,-0.836) (-0.732,-0.830) (-0.766,-0.822) (-0.799,-0.812) (-0.833,-0.801) (-0.865,-0.788) (-0.897,-0.773) (-0.928,-0.757) (-0.959,-0.740) (-0.988,-0.721) (-1.017,-0.700) (-1.045,-0.679) (-1.072,-0.656) (-1.097,-0.633) (-1.122,-0.608) (-1.146,-0.582) (-1.168,-0.555) (-1.190,-0.527) (-1.210,-0.498) (-1.229,-0.469) (-1.247,-0.439) (-1.263,-0.408) (-1.279,-0.376) (-1.293,-0.344) (-1.306,-0.311) (-1.317,-0.278) (-1.327,-0.244) (-1.336,-0.210) (-1.344,-0.176) (-1.350,-0.142) (-1.355,-0.107) (-1.358,-0.072) (-1.360,-0.037) (-1.361,-0.002) (-1.360,0.033) (-1.358,0.068) (-1.355,0.103) (-1.350,0.138) (-1.344,0.173) (-1.337,0.207) (-1.328,0.241) (-1.318,0.275) (-1.307,0.308) (-1.294,0.341) (-1.280,0.373) (-1.265,0.404) (-1.249,0.435) (-1.231,0.466) (-1.212,0.495) (-1.192,0.524) (-1.171,0.552) (-1.148,0.579) (-1.125,0.605) (-1.100,0.630) (-1.074,0.654) (-1.048,0.677) (-1.020,0.698) (-0.991,0.719) (-0.962,0.738) (-0.932,0.755) (-0.900,0.772) (-0.869,0.786) (-0.836,0.800) (-0.803,0.811) (-0.769,0.821) (-0.735,0.829) (-0.701,0.836) (-0.666,0.841) (-0.631,0.843) (-0.596,0.844) (-0.561,0.843) (-0.540,0.827) (-0.540,0.792) (-0.540,0.756) (-0.540,0.721) (-0.540,0.686) (-0.540,0.651) (-0.540,0.616) (-0.540,0.581) (-0.540,0.546) (-0.540,0.511) (-0.540,0.476) (-0.540,0.441) (-0.540,0.405) (-0.540,0.370) (-0.540,0.335) (-0.540,0.300) (-0.540,0.265) (-0.540,0.230) (-0.540,0.195) (-0.540,0.160) (-0.540,0.125) (-0.540,0.090) (-0.540,0.054) (-0.540,0.019) (-0.540,-0.016) (-0.540,-0.051) (-0.540,-0.086) (-0.540,-0.121) (-0.540,-0.156) (-0.540,-0.191) (-0.540,-0.226) (-0.540,-0.261) (-0.540,-0.297) (-0.540,-0.332) (-0.540,-0.367) (-0.540,-0.402) (-0.540,-0.437) (-0.540,-0.472) (-0.540,-0.507) (-0.540,-0.542) (-0.540,-0.577) (-0.540,-0.613) (-0.540,-0.648) (-0.540,-0.683) (-0.540,-0.718) (-0.540,-0.753) (-0.540,-0.788) (-0.540,-0.823) (-0.557,-0.843) (-0.592,-0.844) (-0.627,-0.844) (-0.662,-0.841)};
\addplot[equipot] coordinates {(-0.732,-1.034) (-0.767,-1.033) (-0.802,-1.031) (-0.837,-1.027) (-0.872,-1.021) (-0.906,-1.015) (-0.940,-1.006) (-0.974,-0.996) (-1.007,-0.985) (-1.040,-0.973) (-1.073,-0.959) (-1.104,-0.944) (-1.136,-0.928) (-1.166,-0.911) (-1.196,-0.892) (-1.225,-0.873) (-1.254,-0.852) (-1.281,-0.830) (-1.308,-0.808) (-1.334,-0.784) (-1.360,-0.759) (-1.384,-0.734) (-1.407,-0.708) (-1.430,-0.681) (-1.451,-0.653) (-1.472,-0.625) (-1.492,-0.595) (-1.510,-0.566) (-1.528,-0.535) (-1.545,-0.504) (-1.560,-0.473) (-1.575,-0.441) (-1.589,-0.408) (-1.601,-0.375) (-1.612,-0.342) (-1.623,-0.309) (-1.632,-0.275) (-1.640,-0.241) (-1.647,-0.206) (-1.653,-0.171) (-1.658,-0.137) (-1.662,-0.102) (-1.665,-0.067) (-1.666,-0.031) (-1.667,0.004) (-1.666,0.039) (-1.664,0.074) (-1.661,0.109) (-1.657,0.144) (-1.652,0.179) (-1.646,0.213) (-1.639,0.248) (-1.630,0.282) (-1.621,0.316) (-1.610,0.349) (-1.599,0.382) (-1.586,0.415) (-1.572,0.448) (-1.557,0.479) (-1.541,0.511) (-1.524,0.542) (-1.507,0.572) (-1.488,0.602) (-1.468,0.631) (-1.447,0.659) (-1.425,0.687) (-1.402,0.713) (-1.379,0.739) (-1.354,0.765) (-1.329,0.789) (-1.303,0.812) (-1.276,0.835) (-1.248,0.856) (-1.219,0.877) (-1.190,0.896) (-1.160,0.914) (-1.129,0.932) (-1.098,0.947) (-1.066,0.962) (-1.033,0.976) (-1.000,0.988) (-0.967,0.999) (-0.933,1.008) (-0.899,1.016) (-0.864,1.023) (-0.830,1.028) (-0.795,1.031) (-0.759,1.033) (-0.724,1.034) (-0.689,1.033) (-0.654,1.030) (-0.619,1.026) (-0.585,1.020) (-0.550,1.012) (-0.516,1.003) (-0.483,0.992) (-0.450,0.979) (-0.418,0.964) (-0.387,0.948) (-0.360,0.928) (-0.360,0.893) (-0.360,0.858) (-0.360,0.823) (-0.360,0.787) (-0.360,0.752) (-0.360,0.717) (-0.360,0.682) (-0.360,0.647) (-0.360,0.612) (-0.360,0.576) (-0.360,0.541) (-0.360,0.506) (-0.360,0.471) (-0.360,0.436) (-0.360,0.401) (-0.360,0.366) (-0.360,0.330) (-0.360,0.295) (-0.360,0.260) (-0.360,0.225) (-0.360,0.190) (-0.360,0.155) (-0.360,0.119) (-0.360,0.084) (-0.360,0.049) (-0.360,0.014) (-0.360,-0.021) (-0.360,-0.056) (-0.360,-0.092) (-0.360,-0.127) (-0.360,-0.162) (-0.360,-0.197) (-0.360,-0.232) (-0.360,-0.267) (-0.360,-0.303) (-0.360,-0.338) (-0.360,-0.373) (-0.360,-0.408) (-0.360,-0.443) (-0.360,-0.478) (-0.360,-0.514) (-0.360,-0.549) (-0.360,-0.584) (-0.360,-0.619) (-0.360,-0.654) (-0.360,-0.689) (-0.360,-0.725) (-0.360,-0.760) (-0.360,-0.795) (-0.360,-0.830) (-0.360,-0.865) (-0.360,-0.900) (-0.363,-0.935) (-0.393,-0.952) (-0.425,-0.968) (-0.457,-0.982) (-0.490,-0.994) (-0.523,-1.005) (-0.557,-1.014) (-0.592,-1.021) (-0.627,-1.027) (-0.661,-1.031) (-0.697,-1.033) (-0.732,-1.034)};
\addplot[equipot] coordinates {(-1.148,-1.456) (-1.183,-1.452) (-1.218,-1.447) (-1.252,-1.441) (-1.287,-1.434) (-1.321,-1.425) (-1.355,-1.416) (-1.388,-1.406) (-1.421,-1.395) (-1.454,-1.383) (-1.487,-1.370) (-1.519,-1.356) (-1.551,-1.341) (-1.582,-1.325) (-1.613,-1.309) (-1.644,-1.291) (-1.674,-1.273) (-1.703,-1.254) (-1.732,-1.234) (-1.761,-1.214) (-1.789,-1.193) (-1.816,-1.171) (-1.843,-1.148) (-1.869,-1.125) (-1.895,-1.101) (-1.920,-1.077) (-1.945,-1.051) (-1.968,-1.026) (-1.992,-0.999) (-2.014,-0.973) (-2.036,-0.945) (-2.057,-0.917) (-2.078,-0.889) (-2.098,-0.860) (-2.117,-0.831) (-2.136,-0.801) (-2.153,-0.771) (-2.170,-0.740) (-2.187,-0.709) (-2.203,-0.677) (-2.217,-0.646) (-2.232,-0.614) (-2.245,-0.581) (-2.258,-0.549) (-2.270,-0.516) (-2.281,-0.482) (-2.291,-0.449) (-2.301,-0.415) (-2.310,-0.381) (-2.318,-0.347) (-2.326,-0.313) (-2.332,-0.278) (-2.338,-0.244) (-2.343,-0.209) (-2.347,-0.174) (-2.351,-0.139) (-2.354,-0.104) (-2.356,-0.069) (-2.357,-0.034) (-2.357,0.001) (-2.357,0.036) (-2.355,0.071) (-2.353,0.106) (-2.351,0.141) (-2.347,0.176) (-2.343,0.211) (-2.338,0.246) (-2.332,0.280) (-2.325,0.315) (-2.318,0.349) (-2.310,0.383) (-2.301,0.417) (-2.291,0.451) (-2.280,0.484) (-2.269,0.517) (-2.257,0.550) (-2.244,0.583) (-2.231,0.615) (-2.217,0.647) (-2.202,0.679) (-2.186,0.711) (-2.170,0.742) (-2.152,0.772) (-2.135,0.802) (-2.116,0.832) (-2.097,0.862) (-2.077,0.890) (-2.056,0.919) (-2.035,0.947) (-2.013,0.974) (-1.990,1.001) (-1.967,1.027) (-1.943,1.053) (-1.919,1.078) (-1.894,1.102) (-1.868,1.126) (-1.842,1.150) (-1.815,1.172) (-1.787,1.194) (-1.759,1.215) (-1.731,1.235) (-1.702,1.255) (-1.672,1.274) (-1.642,1.292) (-1.611,1.309) (-1.580,1.326) (-1.549,1.342) (-1.517,1.356) (-1.485,1.370) (-1.452,1.383) (-1.420,1.395) (-1.386,1.406) (-1.353,1.417) (-1.319,1.426) (-1.285,1.434) (-1.250,1.441) (-1.216,1.447) (-1.181,1.453) (-1.146,1.457) (-1.111,1.460) (-1.076,1.461) (-1.041,1.462) (-1.006,1.462) (-0.971,1.460) (-0.936,1.458) (-0.901,1.454) (-0.866,1.449) (-0.832,1.443) (-0.798,1.436) (-0.763,1.427) (-0.730,1.418) (-0.696,1.407) (-0.663,1.395) (-0.631,1.382) (-0.599,1.367) (-0.567,1.352) (-0.536,1.335) (-0.506,1.318) (-0.477,1.299) (-0.448,1.279) (-0.420,1.258) (-0.392,1.236) (-0.366,1.213) (-0.340,1.189) (-0.316,1.164) (-0.292,1.138) (-0.269,1.111) (-0.247,1.084) (-0.226,1.056) (-0.207,1.027) (-0.188,0.997) (-0.180,0.964) (-0.180,0.929) (-0.180,0.894) (-0.180,0.859) (-0.180,0.824) (-0.180,0.789) (-0.180,0.753) (-0.180,0.718) (-0.180,0.683) (-0.180,0.648) (-0.180,0.613) (-0.180,0.578) (-0.180,0.543) (-0.180,0.508) (-0.180,0.473) (-0.180,0.438) (-0.180,0.403) (-0.180,0.367) (-0.180,0.332) (-0.180,0.297) (-0.180,0.262) (-0.180,0.227) (-0.180,0.192) (-0.180,0.157) (-0.180,0.122) (-0.180,0.087) (-0.180,0.052) (-0.180,0.017) (-0.180,-0.019) (-0.180,-0.054) (-0.180,-0.089) (-0.180,-0.124) (-0.180,-0.159) (-0.180,-0.194) (-0.180,-0.229) (-0.180,-0.264) (-0.180,-0.299) (-0.180,-0.334) (-0.180,-0.369) (-0.180,-0.404) (-0.180,-0.440) (-0.180,-0.475) (-0.180,-0.510) (-0.180,-0.545) (-0.180,-0.580) (-0.180,-0.615) (-0.180,-0.650) (-0.180,-0.685) (-0.180,-0.720) (-0.180,-0.755) (-0.180,-0.790) (-0.180,-0.826) (-0.180,-0.861) (-0.180,-0.896) (-0.180,-0.931) (-0.180,-0.966) (-0.189,-0.999) (-0.208,-1.028) (-0.228,-1.057) (-0.248,-1.085) (-0.270,-1.113) (-0.293,-1.139) (-0.317,-1.165) (-0.342,-1.190) (-0.367,-1.214) (-0.394,-1.237) (-0.421,-1.259) (-0.449,-1.280) (-0.478,-1.300) (-0.508,-1.319) (-0.538,-1.336) (-0.569,-1.353) (-0.601,-1.368) (-0.633,-1.383) (-0.665,-1.396) (-0.698,-1.408) (-0.732,-1.418) (-0.765,-1.428) (-0.799,-1.436) (-0.834,-1.443) (-0.868,-1.449) (-0.903,-1.454) (-0.938,-1.458) (-0.973,-1.461) (-1.008,-1.462) (-1.043,-1.462) (-1.078,-1.461) (-1.113,-1.459) (-1.148,-1.456)};
\addplot[equipot] coordinates {(0.923,-1.457) (0.958,-1.460) (0.993,-1.462) (1.028,-1.462) (1.063,-1.462) (1.098,-1.460) (1.133,-1.458) (1.168,-1.454) (1.202,-1.450) (1.237,-1.444) (1.271,-1.437) (1.306,-1.429) (1.340,-1.420) (1.373,-1.411) (1.407,-1.400) (1.440,-1.388) (1.473,-1.375) (1.505,-1.362) (1.537,-1.347) (1.568,-1.332) (1.600,-1.316) (1.630,-1.299) (1.661,-1.281) (1.690,-1.263) (1.720,-1.243) (1.748,-1.223) (1.776,-1.202) (1.804,-1.181) (1.831,-1.158) (1.858,-1.135) (1.884,-1.112) (1.909,-1.087) (1.934,-1.063) (1.958,-1.037) (1.981,-1.011) (2.004,-0.984) (2.027,-0.957) (2.048,-0.930) (2.069,-0.901) (2.089,-0.873) (2.109,-0.844) (2.127,-0.814) (2.146,-0.784) (2.163,-0.753) (2.180,-0.723) (2.196,-0.691) (2.211,-0.660) (2.225,-0.628) (2.239,-0.596) (2.252,-0.563) (2.265,-0.530) (2.276,-0.497) (2.287,-0.464) (2.297,-0.430) (2.306,-0.396) (2.315,-0.362) (2.322,-0.328) (2.329,-0.293) (2.336,-0.259) (2.341,-0.224) (2.346,-0.189) (2.349,-0.155) (2.352,-0.120) (2.355,-0.085) (2.356,-0.050) (2.357,-0.014) (2.357,0.021) (2.356,0.056) (2.354,0.091) (2.352,0.126) (2.349,0.161) (2.345,0.196) (2.340,0.230) (2.335,0.265) (2.328,0.299) (2.321,0.334) (2.313,0.368) (2.305,0.402) (2.295,0.436) (2.285,0.469) (2.274,0.503) (2.262,0.536) (2.250,0.569) (2.237,0.601) (2.223,0.633) (2.208,0.665) (2.193,0.697) (2.177,0.728) (2.160,0.759) (2.143,0.789) (2.124,0.819) (2.105,0.849) (2.086,0.878) (2.065,0.906) (2.044,0.934) (2.023,0.962) (2.000,0.989) (1.977,1.016) (1.954,1.042) (1.930,1.067) (1.905,1.092) (1.879,1.116) (1.853,1.139) (1.827,1.162) (1.799,1.184) (1.772,1.206) (1.743,1.227) (1.714,1.247) (1.685,1.266) (1.655,1.284) (1.625,1.302) (1.594,1.319) (1.563,1.335) (1.531,1.350) (1.499,1.364) (1.467,1.378) (1.434,1.390) (1.401,1.402) (1.367,1.412) (1.334,1.422) (1.300,1.431) (1.266,1.438) (1.231,1.445) (1.196,1.450) (1.162,1.455) (1.127,1.458) (1.092,1.461) (1.057,1.462) (1.022,1.462) (0.986,1.461) (0.951,1.459) (0.917,1.456) (0.882,1.452) (0.847,1.446) (0.813,1.439) (0.778,1.431) (0.745,1.422) (0.711,1.412) (0.678,1.400) (0.645,1.388) (0.613,1.374) (0.581,1.359) (0.550,1.343) (0.519,1.326) (0.490,1.307) (0.460,1.288) (0.432,1.267) (0.404,1.246) (0.377,1.223) (0.351,1.199) (0.326,1.175) (0.302,1.149) (0.279,1.123) (0.257,1.096) (0.236,1.068) (0.215,1.039) (0.196,1.010) (0.180,0.979) (0.180,0.944) (0.180,0.909) (0.180,0.874) (0.180,0.839) (0.180,0.804) (0.180,0.769) (0.180,0.734) (0.180,0.699) (0.180,0.664) (0.180,0.629) (0.180,0.593) (0.180,0.558) (0.180,0.523) (0.180,0.488) (0.180,0.453) (0.180,0.418) (0.180,0.383) (0.180,0.348) (0.180,0.313) (0.180,0.278) (0.180,0.243) (0.180,0.207) (0.180,0.172) (0.180,0.137) (0.180,0.102) (0.180,0.067) (0.180,0.032) (0.180,-0.003) (0.180,-0.038) (0.180,-0.073) (0.180,-0.108) (0.180,-0.143) (0.180,-0.178) (0.180,-0.214) (0.180,-0.249) (0.180,-0.284) (0.180,-0.319) (0.180,-0.354) (0.180,-0.389) (0.180,-0.424) (0.180,-0.459) (0.180,-0.494) (0.180,-0.529) (0.180,-0.564) (0.180,-0.600) (0.180,-0.635) (0.180,-0.670) (0.180,-0.705) (0.180,-0.740) (0.180,-0.775) (0.180,-0.810) (0.180,-0.845) (0.180,-0.880) (0.180,-0.915) (0.180,-0.950) (0.181,-0.985) (0.199,-1.015) (0.219,-1.044) (0.239,-1.073) (0.261,-1.101) (0.283,-1.128) (0.306,-1.154) (0.331,-1.179) (0.356,-1.204) (0.382,-1.227) (0.409,-1.249) (0.437,-1.271) (0.465,-1.291) (0.495,-1.310) (0.525,-1.329) (0.555,-1.346) (0.587,-1.362) (0.618,-1.376) (0.651,-1.390) (0.684,-1.402) (0.717,-1.414) (0.750,-1.424) (0.784,-1.433) (0.819,-1.440) (0.853,-1.447) (0.888,-1.452) (0.923,-1.457)};
\addplot[equipot] coordinates {(0.732,-1.034) (0.767,-1.033) (0.802,-1.031) (0.837,-1.027) (0.872,-1.021) (0.906,-1.015) (0.940,-1.006) (0.974,-0.996) (1.007,-0.985) (1.040,-0.973) (1.073,-0.959) (1.104,-0.944) (1.136,-0.928) (1.166,-0.911) (1.196,-0.892) (1.225,-0.873) (1.254,-0.852) (1.281,-0.830) (1.308,-0.808) (1.334,-0.784) (1.360,-0.759) (1.384,-0.734) (1.407,-0.708) (1.430,-0.681) (1.451,-0.653) (1.472,-0.625) (1.492,-0.595) (1.510,-0.566) (1.528,-0.535) (1.545,-0.504) (1.560,-0.473) (1.575,-0.441) (1.589,-0.408) (1.601,-0.375) (1.612,-0.342) (1.623,-0.309) (1.632,-0.275) (1.640,-0.241) (1.647,-0.206) (1.653,-0.171) (1.658,-0.137) (1.662,-0.102) (1.665,-0.067) (1.666,-0.031) (1.667,0.004) (1.666,0.039) (1.664,0.074) (1.661,0.109) (1.657,0.144) (1.652,0.179) (1.646,0.213) (1.639,0.248) (1.630,0.282) (1.621,0.316) (1.610,0.349) (1.599,0.382) (1.586,0.415) (1.572,0.448) (1.557,0.479) (1.541,0.511) (1.524,0.542) (1.507,0.572) (1.488,0.602) (1.468,0.631) (1.447,0.659) (1.425,0.687) (1.402,0.713) (1.379,0.739) (1.354,0.765) (1.329,0.789) (1.303,0.812) (1.276,0.835) (1.248,0.856) (1.219,0.877) (1.190,0.896) (1.160,0.914) (1.129,0.932) (1.098,0.947) (1.066,0.962) (1.033,0.976) (1.000,0.988) (0.967,0.999) (0.933,1.008) (0.899,1.016) (0.864,1.023) (0.830,1.028) (0.795,1.031) (0.759,1.033) (0.724,1.034) (0.689,1.033) (0.654,1.030) (0.619,1.026) (0.585,1.020) (0.550,1.012) (0.516,1.003) (0.483,0.992) (0.450,0.979) (0.418,0.964) (0.387,0.948) (0.360,0.928) (0.360,0.893) (0.360,0.858) (0.360,0.823) (0.360,0.787) (0.360,0.752) (0.360,0.717) (0.360,0.682) (0.360,0.647) (0.360,0.612) (0.360,0.576) (0.360,0.541) (0.360,0.506) (0.360,0.471) (0.360,0.436) (0.360,0.401) (0.360,0.366) (0.360,0.330) (0.360,0.295) (0.360,0.260) (0.360,0.225) (0.360,0.190) (0.360,0.155) (0.360,0.119) (0.360,0.084) (0.360,0.049) (0.360,0.014) (0.360,-0.021) (0.360,-0.056) (0.360,-0.092) (0.360,-0.127) (0.360,-0.162) (0.360,-0.197) (0.360,-0.232) (0.360,-0.267) (0.360,-0.303) (0.360,-0.338) (0.360,-0.373) (0.360,-0.408) (0.360,-0.443) (0.360,-0.478) (0.360,-0.514) (0.360,-0.549) (0.360,-0.584) (0.360,-0.619) (0.360,-0.654) (0.360,-0.689) (0.360,-0.725) (0.360,-0.760) (0.360,-0.795) (0.360,-0.830) (0.360,-0.865) (0.360,-0.900) (0.363,-0.935) (0.393,-0.952) (0.425,-0.968) (0.457,-0.982) (0.490,-0.994) (0.523,-1.005) (0.557,-1.014) (0.592,-1.021) (0.627,-1.027) (0.661,-1.031) (0.697,-1.033) (0.732,-1.034)};
\addplot[equipot] coordinates {(0.541,-0.841) (0.576,-0.844) (0.611,-0.844) (0.646,-0.842) (0.681,-0.839) (0.716,-0.833) (0.750,-0.826) (0.784,-0.817) (0.817,-0.806) (0.850,-0.794) (0.882,-0.780) (0.914,-0.765) (0.945,-0.748) (0.975,-0.730) (1.004,-0.710) (1.032,-0.689) (1.059,-0.667) (1.086,-0.644) (1.111,-0.619) (1.135,-0.594) (1.158,-0.567) (1.180,-0.540) (1.201,-0.512) (1.220,-0.483) (1.239,-0.453) (1.256,-0.422) (1.272,-0.391) (1.286,-0.359) (1.300,-0.327) (1.312,-0.294) (1.323,-0.260) (1.332,-0.226) (1.340,-0.192) (1.347,-0.158) (1.352,-0.123) (1.357,-0.088) (1.359,-0.053) (1.361,-0.018) (1.361,0.017) (1.359,0.052) (1.357,0.087) (1.353,0.122) (1.347,0.157) (1.340,0.191) (1.332,0.225) (1.323,0.259) (1.312,0.292) (1.300,0.325) (1.287,0.358) (1.272,0.390) (1.256,0.421) (1.239,0.452) (1.221,0.482) (1.201,0.511) (1.181,0.539) (1.159,0.567) (1.136,0.593) (1.112,0.619) (1.086,0.643) (1.060,0.666) (1.033,0.688) (1.005,0.709) (0.976,0.729) (0.946,0.747) (0.915,0.764) (0.884,0.780) (0.851,0.794) (0.818,0.806) (0.785,0.817) (0.751,0.826) (0.717,0.833) (0.682,0.839) (0.647,0.842) (0.612,0.844) (0.577,0.844) (0.542,0.842) (0.540,0.808) (0.540,0.773) (0.540,0.738) (0.540,0.703) (0.540,0.668) (0.540,0.632) (0.540,0.597) (0.540,0.562) (0.540,0.527) (0.540,0.492) (0.540,0.457) (0.540,0.422) (0.540,0.387) (0.540,0.352) (0.540,0.317) (0.540,0.281) (0.540,0.246) (0.540,0.211) (0.540,0.176) (0.540,0.141) (0.540,0.106) (0.540,0.071) (0.540,0.036) (0.540,0.001) (0.540,-0.034) (0.540,-0.070) (0.540,-0.105) (0.540,-0.140) (0.540,-0.175) (0.540,-0.210) (0.540,-0.245) (0.540,-0.280) (0.540,-0.315) (0.540,-0.350) (0.540,-0.385) (0.540,-0.421) (0.540,-0.456) (0.540,-0.491) (0.540,-0.526) (0.540,-0.561) (0.540,-0.596) (0.540,-0.631) (0.540,-0.666) (0.540,-0.701) (0.540,-0.737) (0.540,-0.772) (0.540,-0.807) (0.541,-0.841)};
\addplot[equipot] coordinates {(0.720,-0.691) (0.753,-0.681) (0.785,-0.666) (0.816,-0.649) (0.847,-0.630) (0.876,-0.610) (0.904,-0.589) (0.931,-0.566) (0.957,-0.542) (0.981,-0.516) (1.004,-0.490) (1.026,-0.462) (1.047,-0.433) (1.066,-0.403) (1.084,-0.373) (1.100,-0.341) (1.115,-0.309) (1.128,-0.276) (1.140,-0.243) (1.150,-0.209) (1.159,-0.175) (1.166,-0.140) (1.171,-0.105) (1.175,-0.070) (1.178,-0.035) (1.179,0.001) (1.178,0.036) (1.175,0.071) (1.171,0.107) (1.166,0.141) (1.158,0.176) (1.150,0.210) (1.139,0.244) (1.128,0.278) (1.114,0.310) (1.099,0.342) (1.083,0.374) (1.065,0.404) (1.046,0.434) (1.025,0.463) (1.003,0.491) (0.980,0.517) (0.956,0.543) (0.930,0.567) (0.903,0.590) (0.875,0.611) (0.845,0.631) (0.815,0.649) (0.784,0.666) (0.752,0.681) (0.720,0.689) (0.720,0.654) (0.720,0.619) (0.720,0.583) (0.720,0.548) (0.720,0.512) (0.720,0.477) (0.720,0.442) (0.720,0.406) (0.720,0.371) (0.720,0.335) (0.720,0.300) (0.720,0.265) (0.720,0.229) (0.720,0.194) (0.720,0.159) (0.720,0.123) (0.720,0.088) (0.720,0.052) (0.720,0.017) (0.720,-0.018) (0.720,-0.054) (0.720,-0.089) (0.720,-0.125) (0.720,-0.160) (0.720,-0.195) (0.720,-0.231) (0.720,-0.266) (0.720,-0.301) (0.720,-0.337) (0.720,-0.372) (0.720,-0.408) (0.720,-0.443) (0.720,-0.478) (0.720,-0.514) (0.720,-0.549) (0.720,-0.585) (0.720,-0.620) (0.720,-0.655) (0.720,-0.691)};
\addplot[equipot] coordinates {(0.905,-0.429) (0.926,-0.401) (0.946,-0.372) (0.964,-0.342) (0.981,-0.311) (0.996,-0.280) (1.009,-0.247) (1.021,-0.214) (1.031,-0.180) (1.039,-0.146) (1.045,-0.112) (1.050,-0.077) (1.053,-0.042) (1.054,-0.007) (1.054,0.028) (1.051,0.063) (1.047,0.098) (1.041,0.133) (1.034,0.167) (1.025,0.201) (1.014,0.234) (1.001,0.267) (0.987,0.299) (0.971,0.330) (0.953,0.361) (0.934,0.390) (0.914,0.419) (0.900,0.411) (0.900,0.375) (0.900,0.340) (0.900,0.305) (0.900,0.270) (0.900,0.235) (0.900,0.200) (0.900,0.165) (0.900,0.130) (0.900,0.095) (0.900,0.059) (0.900,0.024) (0.900,-0.011) (0.900,-0.046) (0.900,-0.081) (0.900,-0.116) (0.900,-0.151) (0.900,-0.186) (0.900,-0.221) (0.900,-0.257) (0.900,-0.292) (0.900,-0.327) (0.900,-0.362) (0.900,-0.397) (0.905,-0.429)};
\addplot[lineaE, flechaen=0.711] coordinates {(0.975,0.221) (1.009,0.232) (1.043,0.244) (1.076,0.256) (1.110,0.268) (1.143,0.281) (1.176,0.294) (1.210,0.307) (1.243,0.320) (1.276,0.333) (1.309,0.347) (1.342,0.360) (1.375,0.374) (1.407,0.388) (1.440,0.402) (1.473,0.417) (1.505,0.431) (1.538,0.446) (1.570,0.461) (1.603,0.476) (1.635,0.491) (1.667,0.506) (1.699,0.521) (1.732,0.537) (1.764,0.553) (1.795,0.569) (1.827,0.585) (1.859,0.601) (1.891,0.617) (1.922,0.634) (1.954,0.650) (1.986,0.667) (2.017,0.684) (2.048,0.701) (2.080,0.718) (2.111,0.735) (2.142,0.753) (2.173,0.770) (2.204,0.788) (2.235,0.806) (2.266,0.823) (2.296,0.841) (2.327,0.860) (2.358,0.878) (2.388,0.896) (2.419,0.915) (2.449,0.934) (2.479,0.952) (2.510,0.971) (2.540,0.990) (2.570,1.009) (2.600,1.029)};
\addplot[lineaE, flechaen=0.289] coordinates {(-2.600,1.029) (-2.570,1.009) (-2.540,0.990) (-2.510,0.971) (-2.479,0.952) (-2.449,0.934) (-2.419,0.915) (-2.388,0.896) (-2.358,0.878) (-2.327,0.860) (-2.296,0.841) (-2.266,0.823) (-2.235,0.806) (-2.204,0.788) (-2.173,0.770) (-2.142,0.753) (-2.111,0.735) (-2.080,0.718) (-2.048,0.701) (-2.017,0.684) (-1.986,0.667) (-1.954,0.650) (-1.922,0.634) (-1.891,0.617) (-1.859,0.601) (-1.827,0.585) (-1.795,0.569) (-1.764,0.553) (-1.732,0.537) (-1.699,0.521) (-1.667,0.506) (-1.635,0.491) (-1.603,0.476) (-1.570,0.461) (-1.538,0.446) (-1.505,0.431) (-1.473,0.417) (-1.440,0.402) (-1.407,0.388) (-1.375,0.374) (-1.342,0.360) (-1.309,0.347) (-1.276,0.333) (-1.243,0.320) (-1.210,0.307) (-1.176,0.294) (-1.143,0.281) (-1.110,0.268) (-1.076,0.256) (-1.043,0.244) (-1.009,0.232) (-0.975,0.221)};
\addplot[lineaE, flechaen=0.874] coordinates {(0.925,0.381) (0.954,0.400) (0.984,0.420) (1.012,0.441) (1.041,0.462) (1.070,0.483) (1.098,0.505) (1.126,0.527) (1.154,0.549) (1.181,0.571) (1.209,0.594) (1.236,0.617) (1.263,0.640) (1.289,0.663) (1.316,0.687) (1.342,0.711) (1.368,0.735) (1.394,0.760) (1.419,0.785) (1.444,0.810) (1.469,0.835) (1.494,0.860) (1.519,0.886) (1.543,0.912) (1.567,0.938) (1.591,0.964) (1.615,0.991) (1.638,1.018) (1.661,1.044) (1.684,1.072) (1.707,1.099) (1.729,1.127) (1.751,1.154) (1.773,1.182) (1.795,1.210) (1.816,1.239) (1.838,1.267) (1.858,1.296) (1.879,1.325) (1.900,1.354) (1.920,1.383) (1.940,1.412) (1.959,1.442) (1.979,1.472) (1.998,1.502) (2.017,1.532) (2.036,1.562) (2.054,1.592) (2.072,1.623) (2.090,1.653) (2.108,1.684) (2.125,1.715) (2.143,1.746) (2.159,1.778) (2.176,1.809) (2.193,1.840) (2.209,1.872) (2.224,1.904) (2.240,1.936) (2.255,1.968) (2.270,2.000)};
\addplot[lineaE, flechaen=0.126] coordinates {(-2.270,2.000) (-2.255,1.968) (-2.240,1.936) (-2.224,1.904) (-2.209,1.872) (-2.193,1.840) (-2.176,1.809) (-2.159,1.778) (-2.143,1.746) (-2.125,1.715) (-2.108,1.684) (-2.090,1.653) (-2.072,1.623) (-2.054,1.592) (-2.036,1.562) (-2.017,1.532) (-1.998,1.502) (-1.979,1.472) (-1.959,1.442) (-1.940,1.412) (-1.920,1.383) (-1.900,1.354) (-1.879,1.325) (-1.858,1.296) (-1.838,1.267) (-1.816,1.239) (-1.795,1.210) (-1.773,1.182) (-1.751,1.154) (-1.729,1.127) (-1.707,1.099) (-1.684,1.072) (-1.661,1.044) (-1.638,1.018) (-1.615,0.991) (-1.591,0.964) (-1.567,0.938) (-1.543,0.912) (-1.519,0.886) (-1.494,0.860) (-1.469,0.835) (-1.444,0.810) (-1.419,0.785) (-1.394,0.760) (-1.368,0.735) (-1.342,0.711) (-1.316,0.687) (-1.289,0.663) (-1.263,0.640) (-1.236,0.617) (-1.209,0.594) (-1.181,0.571) (-1.154,0.549) (-1.126,0.527) (-1.098,0.505) (-1.070,0.483) (-1.041,0.462) (-1.012,0.441) (-0.984,0.420) (-0.954,0.400) (-0.925,0.381)};
\addplot[lineaE, flechaen=0.500] coordinates {(0.854,0.520) (0.877,0.547) (0.900,0.573) (0.923,0.601) (0.944,0.628) (0.966,0.656) (0.987,0.684) (1.007,0.713) (1.028,0.742) (1.047,0.771) (1.066,0.801) (1.085,0.831) (1.103,0.861) (1.121,0.891) (1.138,0.922) (1.155,0.953) (1.171,0.984) (1.187,1.015) (1.203,1.047) (1.218,1.079) (1.232,1.111) (1.246,1.143) (1.259,1.176) (1.272,1.209) (1.284,1.242) (1.296,1.275) (1.308,1.308) (1.319,1.342) (1.329,1.375) (1.339,1.409) (1.348,1.443) (1.357,1.477) (1.365,1.511) (1.373,1.546) (1.380,1.580) (1.387,1.615) (1.393,1.650) (1.398,1.684) (1.403,1.719) (1.408,1.754) (1.412,1.789) (1.415,1.824) (1.418,1.859) (1.420,1.894) (1.422,1.930) (1.423,1.965) (1.423,2.000)};
\addplot[lineaE, flechaen=0.500] coordinates {(-1.423,2.000) (-1.423,1.965) (-1.422,1.930) (-1.420,1.894) (-1.418,1.859) (-1.415,1.824) (-1.412,1.789) (-1.408,1.754) (-1.403,1.719) (-1.398,1.684) (-1.393,1.650) (-1.387,1.615) (-1.380,1.580) (-1.373,1.546) (-1.365,1.511) (-1.357,1.477) (-1.348,1.443) (-1.339,1.409) (-1.329,1.375) (-1.319,1.342) (-1.308,1.308) (-1.296,1.275) (-1.284,1.242) (-1.272,1.209) (-1.259,1.176) (-1.246,1.143) (-1.232,1.111) (-1.218,1.079) (-1.203,1.047) (-1.187,1.015) (-1.171,0.984) (-1.155,0.953) (-1.138,0.922) (-1.121,0.891) (-1.103,0.861) (-1.085,0.831) (-1.066,0.801) (-1.047,0.771) (-1.028,0.742) (-1.007,0.713) (-0.987,0.684) (-0.966,0.656) (-0.944,0.628) (-0.923,0.601) (-0.900,0.573) (-0.877,0.547) (-0.854,0.520)};
\addplot[lineaE, flechaen=0.500] coordinates {(0.768,0.641) (0.785,0.672) (0.800,0.704) (0.815,0.736) (0.828,0.769) (0.841,0.802) (0.854,0.836) (0.865,0.869) (0.876,0.903) (0.886,0.937) (0.895,0.972) (0.903,1.006) (0.911,1.041) (0.917,1.076) (0.923,1.111) (0.928,1.146) (0.932,1.181) (0.935,1.217) (0.938,1.252) (0.939,1.288) (0.940,1.323) (0.939,1.359) (0.938,1.395) (0.936,1.430) (0.933,1.465) (0.929,1.501) (0.924,1.536) (0.918,1.571) (0.911,1.606) (0.904,1.641) (0.895,1.675) (0.885,1.709) (0.875,1.743) (0.863,1.777) (0.850,1.810) (0.837,1.843) (0.822,1.875) (0.807,1.907) (0.790,1.939) (0.773,1.970) (0.754,2.000)};
\addplot[lineaE, flechaen=0.500] coordinates {(-0.754,2.000) (-0.773,1.970) (-0.790,1.939) (-0.807,1.907) (-0.822,1.875) (-0.837,1.843) (-0.850,1.810) (-0.863,1.777) (-0.875,1.743) (-0.885,1.709) (-0.895,1.675) (-0.904,1.641) (-0.911,1.606) (-0.918,1.571) (-0.924,1.536) (-0.929,1.501) (-0.933,1.465) (-0.936,1.430) (-0.938,1.395) (-0.939,1.359) (-0.940,1.323) (-0.939,1.288) (-0.938,1.252) (-0.935,1.217) (-0.932,1.181) (-0.928,1.146) (-0.923,1.111) (-0.917,1.076) (-0.911,1.041) (-0.903,1.006) (-0.895,0.972) (-0.886,0.937) (-0.876,0.903) (-0.865,0.869) (-0.854,0.836) (-0.841,0.802) (-0.828,0.769) (-0.815,0.736) (-0.800,0.704) (-0.785,0.672) (-0.768,0.641)};
\addplot[lineaE, flechaen=0.707] coordinates {(0.661,0.751) (0.668,0.785) (0.674,0.820) (0.678,0.855) (0.681,0.890) (0.683,0.926) (0.684,0.961) (0.684,0.996) (0.682,1.032) (0.679,1.067) (0.675,1.102) (0.670,1.137) (0.663,1.172) (0.655,1.206) (0.646,1.240) (0.635,1.274) (0.623,1.307) (0.610,1.340) (0.596,1.372) (0.580,1.404) (0.562,1.435) (0.544,1.465) (0.524,1.494) (0.503,1.522) (0.480,1.550) (0.456,1.576) (0.431,1.601) (0.405,1.624) (0.377,1.646) (0.349,1.667) (0.319,1.686) (0.288,1.704) (0.257,1.720) (0.224,1.734) (0.191,1.746) (0.157,1.757) (0.123,1.765) (0.088,1.771) (0.053,1.775) (0.018,1.778) (-0.018,1.778) (-0.053,1.775) (-0.088,1.771) (-0.123,1.765) (-0.157,1.757) (-0.191,1.746) (-0.224,1.734) (-0.257,1.720) (-0.288,1.704) (-0.319,1.686) (-0.349,1.667) (-0.377,1.646) (-0.405,1.624) (-0.431,1.601) (-0.456,1.576) (-0.480,1.550) (-0.503,1.522) (-0.524,1.494) (-0.544,1.465) (-0.562,1.435) (-0.580,1.404) (-0.596,1.372) (-0.610,1.340) (-0.623,1.307) (-0.635,1.274) (-0.646,1.240) (-0.655,1.206) (-0.663,1.172) (-0.670,1.137) (-0.675,1.102) (-0.679,1.067) (-0.682,1.032) (-0.684,0.996) (-0.684,0.961) (-0.683,0.926) (-0.681,0.890) (-0.678,0.855) (-0.674,0.820) (-0.668,0.785) (-0.661,0.751)};
\addplot[lineaE, flechaen=0.077] coordinates {(0.525,0.851) (0.521,0.886) (0.515,0.921) (0.506,0.955) (0.495,0.989) (0.483,1.023) (0.470,1.056) (0.454,1.088) (0.437,1.119) (0.418,1.149) (0.397,1.178) (0.374,1.206) (0.350,1.232) (0.324,1.256) (0.297,1.279) (0.268,1.300) (0.238,1.319) (0.206,1.336) (0.174,1.350) (0.140,1.362) (0.106,1.372) (0.071,1.379) (0.036,1.383) (-0.000,1.384) (-0.036,1.383) (-0.071,1.379) (-0.106,1.372) (-0.140,1.362) (-0.174,1.350) (-0.206,1.336) (-0.238,1.319) (-0.268,1.300) (-0.297,1.279) (-0.324,1.256) (-0.350,1.232) (-0.374,1.206) (-0.397,1.178) (-0.418,1.149) (-0.437,1.119) (-0.454,1.088) (-0.470,1.056) (-0.483,1.023) (-0.495,0.989) (-0.506,0.955) (-0.515,0.921) (-0.521,0.886) (-0.525,0.851)};
\addplot[lineaE, flechaen=0.500] coordinates {(0.340,0.941) (0.320,0.969) (0.297,0.996) (0.272,1.021) (0.245,1.044) (0.217,1.065) (0.186,1.083) (0.155,1.099) (0.122,1.112) (0.088,1.121) (0.053,1.128) (0.018,1.131) (-0.018,1.131) (-0.053,1.128) (-0.088,1.121) (-0.122,1.112) (-0.155,1.099) (-0.186,1.083) (-0.217,1.065) (-0.245,1.044) (-0.272,1.021) (-0.297,0.996) (-0.320,0.969) (-0.340,0.941)};
\addplot[lineaE, flechaen=0.500] coordinates {(0.335,0.942) (0.301,0.940) (0.266,0.940) (0.230,0.940) (0.195,0.940) (0.159,0.940) (0.124,0.940) (0.089,0.940) (0.053,0.940) (0.018,0.940) (-0.018,0.940) (-0.053,0.940) (-0.089,0.940) (-0.124,0.940) (-0.159,0.940) (-0.195,0.940) (-0.230,0.940) (-0.266,0.940) (-0.301,0.940) (-0.335,0.942)};
\addplot[lineaE, flechaen=0.500] coordinates {(0.522,0.853) (0.488,0.850) (0.451,0.850) (0.415,0.850) (0.379,0.850) (0.343,0.850) (0.307,0.850) (0.271,0.850) (0.235,0.850) (0.199,0.850) (0.163,0.850) (0.126,0.850) (0.090,0.850) (0.054,0.850) (0.018,0.850) (-0.018,0.850) (-0.054,0.850) (-0.090,0.850) (-0.126,0.850) (-0.163,0.850) (-0.199,0.850) (-0.235,0.850) (-0.271,0.850) (-0.307,0.850) (-0.343,0.850) (-0.379,0.850) (-0.415,0.850) (-0.451,0.850) (-0.488,0.850) (-0.522,0.853)};
\addplot[lineaE, flechaen=0.121] coordinates {(0.660,0.751) (0.625,0.750) (0.589,0.750) (0.553,0.750) (0.518,0.750) (0.482,0.750) (0.446,0.750) (0.411,0.750) (0.375,0.750) (0.339,0.750) (0.303,0.750) (0.268,0.750) (0.232,0.750) (0.196,0.750) (0.161,0.750) (0.125,0.750) (0.089,0.750) (0.054,0.750) (0.018,0.750) (-0.018,0.750) (-0.054,0.750) (-0.089,0.750) (-0.125,0.750) (-0.161,0.750) (-0.196,0.750) (-0.232,0.750) (-0.268,0.750) (-0.303,0.750) (-0.339,0.750) (-0.375,0.750) (-0.411,0.750) (-0.446,0.750) (-0.482,0.750) (-0.518,0.750) (-0.553,0.750) (-0.589,0.750) (-0.625,0.750) (-0.660,0.751)};
\addplot[lineaE, flechaen=0.826] coordinates {(0.768,0.641) (0.732,0.640) (0.697,0.640) (0.661,0.640) (0.625,0.640) (0.589,0.640) (0.554,0.640) (0.518,0.640) (0.482,0.640) (0.447,0.640) (0.411,0.640) (0.375,0.640) (0.339,0.640) (0.304,0.640) (0.268,0.640) (0.232,0.640) (0.196,0.640) (0.161,0.640) (0.125,0.640) (0.089,0.640) (0.054,0.640) (0.018,0.640) (-0.018,0.640) (-0.054,0.640) (-0.089,0.640) (-0.125,0.640) (-0.161,0.640) (-0.196,0.640) (-0.232,0.640) (-0.268,0.640) (-0.304,0.640) (-0.339,0.640) (-0.375,0.640) (-0.411,0.640) (-0.447,0.640) (-0.482,0.640) (-0.518,0.640) (-0.554,0.640) (-0.589,0.640) (-0.625,0.640) (-0.661,0.640) (-0.697,0.640) (-0.732,0.640) (-0.768,0.641)};
\addplot[lineaE, flechaen=0.207] coordinates {(0.853,0.521) (0.818,0.520) (0.783,0.520) (0.747,0.520) (0.711,0.520) (0.676,0.520) (0.640,0.520) (0.605,0.520) (0.569,0.520) (0.534,0.520) (0.498,0.520) (0.462,0.520) (0.427,0.520) (0.391,0.520) (0.356,0.520) (0.320,0.520) (0.285,0.520) (0.249,0.520) (0.213,0.520) (0.178,0.520) (0.142,0.520) (0.107,0.520) (0.071,0.520) (0.036,0.520) (0.000,0.520) (-0.036,0.520) (-0.071,0.520) (-0.107,0.520) (-0.142,0.520) (-0.178,0.520) (-0.213,0.520) (-0.249,0.520) (-0.285,0.520) (-0.320,0.520) (-0.356,0.520) (-0.391,0.520) (-0.427,0.520) (-0.462,0.520) (-0.498,0.520) (-0.534,0.520) (-0.569,0.520) (-0.605,0.520) (-0.640,0.520) (-0.676,0.520) (-0.711,0.520) (-0.747,0.520) (-0.783,0.520) (-0.818,0.520) (-0.853,0.521)};
\addplot[lineaE, flechaen=0.770] coordinates {(0.925,0.381) (0.889,0.380) (0.854,0.380) (0.818,0.380) (0.783,0.380) (0.747,0.380) (0.711,0.380) (0.676,0.380) (0.640,0.380) (0.605,0.380) (0.569,0.380) (0.534,0.380) (0.498,0.380) (0.462,0.380) (0.427,0.380) (0.391,0.380) (0.356,0.380) (0.320,0.380) (0.285,0.380) (0.249,0.380) (0.213,0.380) (0.178,0.380) (0.142,0.380) (0.107,0.380) (0.071,0.380) (0.036,0.380) (-0.000,0.380) (-0.036,0.380) (-0.071,0.380) (-0.107,0.380) (-0.142,0.380) (-0.178,0.380) (-0.213,0.380) (-0.249,0.380) (-0.285,0.380) (-0.320,0.380) (-0.356,0.380) (-0.391,0.380) (-0.427,0.380) (-0.462,0.380) (-0.498,0.380) (-0.534,0.380) (-0.569,0.380) (-0.605,0.380) (-0.640,0.380) (-0.676,0.380) (-0.711,0.380) (-0.747,0.380) (-0.783,0.380) (-0.818,0.380) (-0.854,0.380) (-0.889,0.380) (-0.925,0.381)};
\addplot[lineaE, flechaen=0.244] coordinates {(0.975,0.220) (0.940,0.220) (0.905,0.220) (0.869,0.220) (0.834,0.220) (0.798,0.220) (0.763,0.220) (0.727,0.220) (0.692,0.220) (0.656,0.220) (0.621,0.220) (0.585,0.220) (0.550,0.220) (0.514,0.220) (0.479,0.220) (0.443,0.220) (0.408,0.220) (0.372,0.220) (0.337,0.220) (0.302,0.220) (0.266,0.220) (0.231,0.220) (0.195,0.220) (0.160,0.220) (0.124,0.220) (0.089,0.220) (0.053,0.220) (0.018,0.220) (-0.018,0.220) (-0.053,0.220) (-0.089,0.220) (-0.124,0.220) (-0.160,0.220) (-0.195,0.220) (-0.231,0.220) (-0.266,0.220) (-0.302,0.220) (-0.337,0.220) (-0.372,0.220) (-0.408,0.220) (-0.443,0.220) (-0.479,0.220) (-0.514,0.220) (-0.550,0.220) (-0.585,0.220) (-0.621,0.220) (-0.656,0.220) (-0.692,0.220) (-0.727,0.220) (-0.763,0.220) (-0.798,0.220) (-0.834,0.220) (-0.869,0.220) (-0.905,0.220) (-0.940,0.220) (-0.975,0.220)};
\addplot[lineaE, flechaen=0.711] coordinates {(0.975,-0.221) (1.009,-0.232) (1.043,-0.244) (1.076,-0.256) (1.110,-0.268) (1.143,-0.281) (1.176,-0.294) (1.210,-0.307) (1.243,-0.320) (1.276,-0.333) (1.309,-0.347) (1.342,-0.360) (1.375,-0.374) (1.407,-0.388) (1.440,-0.402) (1.473,-0.417) (1.505,-0.431) (1.538,-0.446) (1.570,-0.461) (1.603,-0.476) (1.635,-0.491) (1.667,-0.506) (1.699,-0.521) (1.732,-0.537) (1.764,-0.553) (1.795,-0.569) (1.827,-0.585) (1.859,-0.601) (1.891,-0.617) (1.922,-0.634) (1.954,-0.650) (1.986,-0.667) (2.017,-0.684) (2.048,-0.701) (2.080,-0.718) (2.111,-0.735) (2.142,-0.753) (2.173,-0.770) (2.204,-0.788) (2.235,-0.806) (2.266,-0.823) (2.296,-0.841) (2.327,-0.860) (2.358,-0.878) (2.388,-0.896) (2.419,-0.915) (2.449,-0.934) (2.479,-0.952) (2.510,-0.971) (2.540,-0.990) (2.570,-1.009) (2.600,-1.029)};
\addplot[lineaE, flechaen=0.289] coordinates {(-2.600,-1.029) (-2.570,-1.009) (-2.540,-0.990) (-2.510,-0.971) (-2.479,-0.952) (-2.449,-0.934) (-2.419,-0.915) (-2.388,-0.896) (-2.358,-0.878) (-2.327,-0.860) (-2.296,-0.841) (-2.266,-0.823) (-2.235,-0.806) (-2.204,-0.788) (-2.173,-0.770) (-2.142,-0.753) (-2.111,-0.735) (-2.080,-0.718) (-2.048,-0.701) (-2.017,-0.684) (-1.986,-0.667) (-1.954,-0.650) (-1.922,-0.634) (-1.891,-0.617) (-1.859,-0.601) (-1.827,-0.585) (-1.795,-0.569) (-1.764,-0.553) (-1.732,-0.537) (-1.699,-0.521) (-1.667,-0.506) (-1.635,-0.491) (-1.603,-0.476) (-1.570,-0.461) (-1.538,-0.446) (-1.505,-0.431) (-1.473,-0.417) (-1.440,-0.402) (-1.407,-0.388) (-1.375,-0.374) (-1.342,-0.360) (-1.309,-0.347) (-1.276,-0.333) (-1.243,-0.320) (-1.210,-0.307) (-1.176,-0.294) (-1.143,-0.281) (-1.110,-0.268) (-1.076,-0.256) (-1.043,-0.244) (-1.009,-0.232) (-0.975,-0.221)};
\addplot[lineaE, flechaen=0.874] coordinates {(0.925,-0.381) (0.954,-0.400) (0.984,-0.420) (1.012,-0.441) (1.041,-0.462) (1.070,-0.483) (1.098,-0.505) (1.126,-0.527) (1.154,-0.549) (1.181,-0.571) (1.209,-0.594) (1.236,-0.617) (1.263,-0.640) (1.289,-0.663) (1.316,-0.687) (1.342,-0.711) (1.368,-0.735) (1.394,-0.760) (1.419,-0.785) (1.444,-0.810) (1.469,-0.835) (1.494,-0.860) (1.519,-0.886) (1.543,-0.912) (1.567,-0.938) (1.591,-0.964) (1.615,-0.991) (1.638,-1.018) (1.661,-1.044) (1.684,-1.072) (1.707,-1.099) (1.729,-1.127) (1.751,-1.154) (1.773,-1.182) (1.795,-1.210) (1.816,-1.239) (1.838,-1.267) (1.858,-1.296) (1.879,-1.325) (1.900,-1.354) (1.920,-1.383) (1.940,-1.412) (1.959,-1.442) (1.979,-1.472) (1.998,-1.502) (2.017,-1.532) (2.036,-1.562) (2.054,-1.592) (2.072,-1.623) (2.090,-1.653) (2.108,-1.684) (2.125,-1.715) (2.143,-1.746) (2.159,-1.778) (2.176,-1.809) (2.193,-1.840) (2.209,-1.872) (2.224,-1.904) (2.240,-1.936) (2.255,-1.968) (2.270,-2.000)};
\addplot[lineaE, flechaen=0.126] coordinates {(-2.270,-2.000) (-2.255,-1.968) (-2.240,-1.936) (-2.224,-1.904) (-2.209,-1.872) (-2.193,-1.840) (-2.176,-1.809) (-2.159,-1.778) (-2.143,-1.746) (-2.125,-1.715) (-2.108,-1.684) (-2.090,-1.653) (-2.072,-1.623) (-2.054,-1.592) (-2.036,-1.562) (-2.017,-1.532) (-1.998,-1.502) (-1.979,-1.472) (-1.959,-1.442) (-1.940,-1.412) (-1.920,-1.383) (-1.900,-1.354) (-1.879,-1.325) (-1.858,-1.296) (-1.838,-1.267) (-1.816,-1.239) (-1.795,-1.210) (-1.773,-1.182) (-1.751,-1.154) (-1.729,-1.127) (-1.707,-1.099) (-1.684,-1.072) (-1.661,-1.044) (-1.638,-1.018) (-1.615,-0.991) (-1.591,-0.964) (-1.567,-0.938) (-1.543,-0.912) (-1.519,-0.886) (-1.494,-0.860) (-1.469,-0.835) (-1.444,-0.810) (-1.419,-0.785) (-1.394,-0.760) (-1.368,-0.735) (-1.342,-0.711) (-1.316,-0.687) (-1.289,-0.663) (-1.263,-0.640) (-1.236,-0.617) (-1.209,-0.594) (-1.181,-0.571) (-1.154,-0.549) (-1.126,-0.527) (-1.098,-0.505) (-1.070,-0.483) (-1.041,-0.462) (-1.012,-0.441) (-0.984,-0.420) (-0.954,-0.400) (-0.925,-0.381)};
\addplot[lineaE, flechaen=0.500] coordinates {(0.854,-0.520) (0.877,-0.547) (0.900,-0.573) (0.923,-0.601) (0.944,-0.628) (0.966,-0.656) (0.987,-0.684) (1.007,-0.713) (1.028,-0.742) (1.047,-0.771) (1.066,-0.801) (1.085,-0.831) (1.103,-0.861) (1.121,-0.891) (1.138,-0.922) (1.155,-0.953) (1.171,-0.984) (1.187,-1.015) (1.203,-1.047) (1.218,-1.079) (1.232,-1.111) (1.246,-1.143) (1.259,-1.176) (1.272,-1.209) (1.284,-1.242) (1.296,-1.275) (1.308,-1.308) (1.319,-1.342) (1.329,-1.375) (1.339,-1.409) (1.348,-1.443) (1.357,-1.477) (1.365,-1.511) (1.373,-1.546) (1.380,-1.580) (1.387,-1.615) (1.393,-1.650) (1.398,-1.684) (1.403,-1.719) (1.408,-1.754) (1.412,-1.789) (1.415,-1.824) (1.418,-1.859) (1.420,-1.894) (1.422,-1.930) (1.423,-1.965) (1.423,-2.000)};
\addplot[lineaE, flechaen=0.500] coordinates {(-1.423,-2.000) (-1.423,-1.965) (-1.422,-1.930) (-1.420,-1.894) (-1.418,-1.859) (-1.415,-1.824) (-1.412,-1.789) (-1.408,-1.754) (-1.403,-1.719) (-1.398,-1.684) (-1.393,-1.650) (-1.387,-1.615) (-1.380,-1.580) (-1.373,-1.546) (-1.365,-1.511) (-1.357,-1.477) (-1.348,-1.443) (-1.339,-1.409) (-1.329,-1.375) (-1.319,-1.342) (-1.308,-1.308) (-1.296,-1.275) (-1.284,-1.242) (-1.272,-1.209) (-1.259,-1.176) (-1.246,-1.143) (-1.232,-1.111) (-1.218,-1.079) (-1.203,-1.047) (-1.187,-1.015) (-1.171,-0.984) (-1.155,-0.953) (-1.138,-0.922) (-1.121,-0.891) (-1.103,-0.861) (-1.085,-0.831) (-1.066,-0.801) (-1.047,-0.771) (-1.028,-0.742) (-1.007,-0.713) (-0.987,-0.684) (-0.966,-0.656) (-0.944,-0.628) (-0.923,-0.601) (-0.900,-0.573) (-0.877,-0.547) (-0.854,-0.520)};
\addplot[lineaE, flechaen=0.500] coordinates {(0.768,-0.641) (0.785,-0.672) (0.800,-0.704) (0.815,-0.736) (0.828,-0.769) (0.841,-0.802) (0.854,-0.836) (0.865,-0.869) (0.876,-0.903) (0.886,-0.937) (0.895,-0.972) (0.903,-1.006) (0.911,-1.041) (0.917,-1.076) (0.923,-1.111) (0.928,-1.146) (0.932,-1.181) (0.935,-1.217) (0.938,-1.252) (0.939,-1.288) (0.940,-1.323) (0.939,-1.359) (0.938,-1.395) (0.936,-1.430) (0.933,-1.465) (0.929,-1.501) (0.924,-1.536) (0.918,-1.571) (0.911,-1.606) (0.904,-1.641) (0.895,-1.675) (0.885,-1.709) (0.875,-1.743) (0.863,-1.777) (0.850,-1.810) (0.837,-1.843) (0.822,-1.875) (0.807,-1.907) (0.790,-1.939) (0.773,-1.970) (0.754,-2.000)};
\addplot[lineaE, flechaen=0.500] coordinates {(-0.754,-2.000) (-0.773,-1.970) (-0.790,-1.939) (-0.807,-1.907) (-0.822,-1.875) (-0.837,-1.843) (-0.850,-1.810) (-0.863,-1.777) (-0.875,-1.743) (-0.885,-1.709) (-0.895,-1.675) (-0.904,-1.641) (-0.911,-1.606) (-0.918,-1.571) (-0.924,-1.536) (-0.929,-1.501) (-0.933,-1.465) (-0.936,-1.430) (-0.938,-1.395) (-0.939,-1.359) (-0.940,-1.323) (-0.939,-1.288) (-0.938,-1.252) (-0.935,-1.217) (-0.932,-1.181) (-0.928,-1.146) (-0.923,-1.111) (-0.917,-1.076) (-0.911,-1.041) (-0.903,-1.006) (-0.895,-0.972) (-0.886,-0.937) (-0.876,-0.903) (-0.865,-0.869) (-0.854,-0.836) (-0.841,-0.802) (-0.828,-0.769) (-0.815,-0.736) (-0.800,-0.704) (-0.785,-0.672) (-0.768,-0.641)};
\addplot[lineaE, flechaen=0.293] coordinates {(0.661,-0.751) (0.668,-0.785) (0.674,-0.820) (0.678,-0.855) (0.681,-0.890) (0.683,-0.926) (0.684,-0.961) (0.684,-0.996) (0.682,-1.032) (0.679,-1.067) (0.675,-1.102) (0.670,-1.137) (0.663,-1.172) (0.655,-1.206) (0.646,-1.240) (0.635,-1.274) (0.623,-1.307) (0.610,-1.340) (0.596,-1.372) (0.580,-1.404) (0.562,-1.435) (0.544,-1.465) (0.524,-1.494) (0.503,-1.522) (0.480,-1.550) (0.456,-1.576) (0.431,-1.601) (0.405,-1.624) (0.377,-1.646) (0.349,-1.667) (0.319,-1.686) (0.288,-1.704) (0.257,-1.720) (0.224,-1.734) (0.191,-1.746) (0.157,-1.757) (0.123,-1.765) (0.088,-1.771) (0.053,-1.775) (0.018,-1.778) (-0.018,-1.778) (-0.053,-1.775) (-0.088,-1.771) (-0.123,-1.765) (-0.157,-1.757) (-0.191,-1.746) (-0.224,-1.734) (-0.257,-1.720) (-0.288,-1.704) (-0.319,-1.686) (-0.349,-1.667) (-0.377,-1.646) (-0.405,-1.624) (-0.431,-1.601) (-0.456,-1.576) (-0.480,-1.550) (-0.503,-1.522) (-0.524,-1.494) (-0.544,-1.465) (-0.562,-1.435) (-0.580,-1.404) (-0.596,-1.372) (-0.610,-1.340) (-0.623,-1.307) (-0.635,-1.274) (-0.646,-1.240) (-0.655,-1.206) (-0.663,-1.172) (-0.670,-1.137) (-0.675,-1.102) (-0.679,-1.067) (-0.682,-1.032) (-0.684,-0.996) (-0.684,-0.961) (-0.683,-0.926) (-0.681,-0.890) (-0.678,-0.855) (-0.674,-0.820) (-0.668,-0.785) (-0.661,-0.751)};
\addplot[lineaE, flechaen=0.923] coordinates {(0.525,-0.851) (0.521,-0.886) (0.515,-0.921) (0.506,-0.955) (0.495,-0.989) (0.483,-1.023) (0.470,-1.056) (0.454,-1.088) (0.437,-1.119) (0.418,-1.149) (0.397,-1.178) (0.374,-1.206) (0.350,-1.232) (0.324,-1.256) (0.297,-1.279) (0.268,-1.300) (0.238,-1.319) (0.206,-1.336) (0.174,-1.350) (0.140,-1.362) (0.106,-1.372) (0.071,-1.379) (0.036,-1.383) (-0.000,-1.384) (-0.036,-1.383) (-0.071,-1.379) (-0.106,-1.372) (-0.140,-1.362) (-0.174,-1.350) (-0.206,-1.336) (-0.238,-1.319) (-0.268,-1.300) (-0.297,-1.279) (-0.324,-1.256) (-0.350,-1.232) (-0.374,-1.206) (-0.397,-1.178) (-0.418,-1.149) (-0.437,-1.119) (-0.454,-1.088) (-0.470,-1.056) (-0.483,-1.023) (-0.495,-0.989) (-0.506,-0.955) (-0.515,-0.921) (-0.521,-0.886) (-0.525,-0.851)};
\addplot[lineaE, flechaen=0.500] coordinates {(0.340,-0.941) (0.320,-0.969) (0.297,-0.996) (0.272,-1.021) (0.245,-1.044) (0.217,-1.065) (0.186,-1.083) (0.155,-1.099) (0.122,-1.112) (0.088,-1.121) (0.053,-1.128) (0.018,-1.131) (-0.018,-1.131) (-0.053,-1.128) (-0.088,-1.121) (-0.122,-1.112) (-0.155,-1.099) (-0.186,-1.083) (-0.217,-1.065) (-0.245,-1.044) (-0.272,-1.021) (-0.297,-0.996) (-0.320,-0.969) (-0.340,-0.941)};
\addplot[lineaE, flechaen=0.500] coordinates {(0.335,-0.942) (0.301,-0.940) (0.266,-0.940) (0.230,-0.940) (0.195,-0.940) (0.159,-0.940) (0.124,-0.940) (0.089,-0.940) (0.053,-0.940) (0.018,-0.940) (-0.018,-0.940) (-0.053,-0.940) (-0.089,-0.940) (-0.124,-0.940) (-0.159,-0.940) (-0.195,-0.940) (-0.230,-0.940) (-0.266,-0.940) (-0.301,-0.940) (-0.335,-0.942)};
\addplot[lineaE, flechaen=0.500] coordinates {(0.522,-0.853) (0.488,-0.850) (0.451,-0.850) (0.415,-0.850) (0.379,-0.850) (0.343,-0.850) (0.307,-0.850) (0.271,-0.850) (0.235,-0.850) (0.199,-0.850) (0.163,-0.850) (0.126,-0.850) (0.090,-0.850) (0.054,-0.850) (0.018,-0.850) (-0.018,-0.850) (-0.054,-0.850) (-0.090,-0.850) (-0.126,-0.850) (-0.163,-0.850) (-0.199,-0.850) (-0.235,-0.850) (-0.271,-0.850) (-0.307,-0.850) (-0.343,-0.850) (-0.379,-0.850) (-0.415,-0.850) (-0.451,-0.850) (-0.488,-0.850) (-0.522,-0.853)};
\addplot[lineaE, flechaen=0.879] coordinates {(0.660,-0.751) (0.625,-0.750) (0.589,-0.750) (0.553,-0.750) (0.518,-0.750) (0.482,-0.750) (0.446,-0.750) (0.411,-0.750) (0.375,-0.750) (0.339,-0.750) (0.303,-0.750) (0.268,-0.750) (0.232,-0.750) (0.196,-0.750) (0.161,-0.750) (0.125,-0.750) (0.089,-0.750) (0.054,-0.750) (0.018,-0.750) (-0.018,-0.750) (-0.054,-0.750) (-0.089,-0.750) (-0.125,-0.750) (-0.161,-0.750) (-0.196,-0.750) (-0.232,-0.750) (-0.268,-0.750) (-0.303,-0.750) (-0.339,-0.750) (-0.375,-0.750) (-0.411,-0.750) (-0.446,-0.750) (-0.482,-0.750) (-0.518,-0.750) (-0.553,-0.750) (-0.589,-0.750) (-0.625,-0.750) (-0.660,-0.751)};
\addplot[lineaE, flechaen=0.174] coordinates {(0.768,-0.641) (0.732,-0.640) (0.697,-0.640) (0.661,-0.640) (0.625,-0.640) (0.589,-0.640) (0.554,-0.640) (0.518,-0.640) (0.482,-0.640) (0.447,-0.640) (0.411,-0.640) (0.375,-0.640) (0.339,-0.640) (0.304,-0.640) (0.268,-0.640) (0.232,-0.640) (0.196,-0.640) (0.161,-0.640) (0.125,-0.640) (0.089,-0.640) (0.054,-0.640) (0.018,-0.640) (-0.018,-0.640) (-0.054,-0.640) (-0.089,-0.640) (-0.125,-0.640) (-0.161,-0.640) (-0.196,-0.640) (-0.232,-0.640) (-0.268,-0.640) (-0.304,-0.640) (-0.339,-0.640) (-0.375,-0.640) (-0.411,-0.640) (-0.447,-0.640) (-0.482,-0.640) (-0.518,-0.640) (-0.554,-0.640) (-0.589,-0.640) (-0.625,-0.640) (-0.661,-0.640) (-0.697,-0.640) (-0.732,-0.640) (-0.768,-0.641)};
\addplot[lineaE, flechaen=0.793] coordinates {(0.853,-0.521) (0.818,-0.520) (0.783,-0.520) (0.747,-0.520) (0.711,-0.520) (0.676,-0.520) (0.640,-0.520) (0.605,-0.520) (0.569,-0.520) (0.534,-0.520) (0.498,-0.520) (0.462,-0.520) (0.427,-0.520) (0.391,-0.520) (0.356,-0.520) (0.320,-0.520) (0.285,-0.520) (0.249,-0.520) (0.213,-0.520) (0.178,-0.520) (0.142,-0.520) (0.107,-0.520) (0.071,-0.520) (0.036,-0.520) (0.000,-0.520) (-0.036,-0.520) (-0.071,-0.520) (-0.107,-0.520) (-0.142,-0.520) (-0.178,-0.520) (-0.213,-0.520) (-0.249,-0.520) (-0.285,-0.520) (-0.320,-0.520) (-0.356,-0.520) (-0.391,-0.520) (-0.427,-0.520) (-0.462,-0.520) (-0.498,-0.520) (-0.534,-0.520) (-0.569,-0.520) (-0.605,-0.520) (-0.640,-0.520) (-0.676,-0.520) (-0.711,-0.520) (-0.747,-0.520) (-0.783,-0.520) (-0.818,-0.520) (-0.853,-0.521)};
\addplot[lineaE, flechaen=0.230] coordinates {(0.925,-0.381) (0.889,-0.380) (0.854,-0.380) (0.818,-0.380) (0.783,-0.380) (0.747,-0.380) (0.711,-0.380) (0.676,-0.380) (0.640,-0.380) (0.605,-0.380) (0.569,-0.380) (0.534,-0.380) (0.498,-0.380) (0.462,-0.380) (0.427,-0.380) (0.391,-0.380) (0.356,-0.380) (0.320,-0.380) (0.285,-0.380) (0.249,-0.380) (0.213,-0.380) (0.178,-0.380) (0.142,-0.380) (0.107,-0.380) (0.071,-0.380) (0.036,-0.380) (-0.000,-0.380) (-0.036,-0.380) (-0.071,-0.380) (-0.107,-0.380) (-0.142,-0.380) (-0.178,-0.380) (-0.213,-0.380) (-0.249,-0.380) (-0.285,-0.380) (-0.320,-0.380) (-0.356,-0.380) (-0.391,-0.380) (-0.427,-0.380) (-0.462,-0.380) (-0.498,-0.380) (-0.534,-0.380) (-0.569,-0.380) (-0.605,-0.380) (-0.640,-0.380) (-0.676,-0.380) (-0.711,-0.380) (-0.747,-0.380) (-0.783,-0.380) (-0.818,-0.380) (-0.854,-0.380) (-0.889,-0.380) (-0.925,-0.381)};
\addplot[lineaE, flechaen=0.756] coordinates {(0.975,-0.220) (0.940,-0.220) (0.905,-0.220) (0.869,-0.220) (0.834,-0.220) (0.798,-0.220) (0.763,-0.220) (0.727,-0.220) (0.692,-0.220) (0.656,-0.220) (0.621,-0.220) (0.585,-0.220) (0.550,-0.220) (0.514,-0.220) (0.479,-0.220) (0.443,-0.220) (0.408,-0.220) (0.372,-0.220) (0.337,-0.220) (0.302,-0.220) (0.266,-0.220) (0.231,-0.220) (0.195,-0.220) (0.160,-0.220) (0.124,-0.220) (0.089,-0.220) (0.053,-0.220) (0.018,-0.220) (-0.018,-0.220) (-0.053,-0.220) (-0.089,-0.220) (-0.124,-0.220) (-0.160,-0.220) (-0.195,-0.220) (-0.231,-0.220) (-0.266,-0.220) (-0.302,-0.220) (-0.337,-0.220) (-0.372,-0.220) (-0.408,-0.220) (-0.443,-0.220) (-0.479,-0.220) (-0.514,-0.220) (-0.550,-0.220) (-0.585,-0.220) (-0.621,-0.220) (-0.656,-0.220) (-0.692,-0.220) (-0.727,-0.220) (-0.763,-0.220) (-0.798,-0.220) (-0.834,-0.220) (-0.869,-0.220) (-0.905,-0.220) (-0.940,-0.220) (-0.975,-0.220)};
\addplot[lineaE, flechaen=0.281] coordinates {(-2.600,0.000) (-1.000,0.000)};
\addplot[lineaE, flechaen=0.719] coordinates {(1.000,0.000) (2.600,0.000)};
\addplot[lineaE, flechaen=0.750] coordinates {(1.000,0.000) (-1.000,0.000)};

  \draw[texto, line width=0.6pt] (axis cs:0,0) circle[radius=1];
  \pgfplotsinvokeforeach{-60,-30,0,30,60}{
    \node[cpos, font=\scriptsize\bfseries, fill=fondo, circle, inner sep=0.2pt] at (axis cs:{1.16*cos(#1)},{1.16*sin(#1)}) {$+$};
    \node[cneg, font=\scriptsize\bfseries, fill=fondo, circle, inner sep=0.2pt] at (axis cs:{-1.16*cos(#1)},{1.16*sin(#1)}) {$-$};
  }
\end{axis}
\end{tikzpicture}
\caption{Líneas de $\vE$}
\end{subfigure}\hfill
\begin{subfigure}[b]{0.49\textwidth}\centering
\begin{tikzpicture}
\begin{axis}[esferap]
  \fill[dielec] (axis cs:0,0) circle[radius=1];
  %% Generado por scripts/lineas_campo.py -- no editar a mano
\addplot[equipot] coordinates {(-0.905,-0.429) (-0.926,-0.401) (-0.946,-0.372) (-0.964,-0.342) (-0.981,-0.311) (-0.996,-0.280) (-1.009,-0.247) (-1.021,-0.214) (-1.031,-0.180) (-1.039,-0.146) (-1.045,-0.112) (-1.050,-0.077) (-1.053,-0.042) (-1.054,-0.007) (-1.054,0.028) (-1.051,0.063) (-1.047,0.098) (-1.041,0.133) (-1.034,0.167) (-1.025,0.201) (-1.014,0.234) (-1.001,0.267) (-0.987,0.299) (-0.971,0.330) (-0.953,0.361) (-0.934,0.390) (-0.914,0.419) (-0.900,0.411) (-0.900,0.375) (-0.900,0.340) (-0.900,0.305) (-0.900,0.270) (-0.900,0.235) (-0.900,0.200) (-0.900,0.165) (-0.900,0.130) (-0.900,0.095) (-0.900,0.059) (-0.900,0.024) (-0.900,-0.011) (-0.900,-0.046) (-0.900,-0.081) (-0.900,-0.116) (-0.900,-0.151) (-0.900,-0.186) (-0.900,-0.221) (-0.900,-0.257) (-0.900,-0.292) (-0.900,-0.327) (-0.900,-0.362) (-0.900,-0.397) (-0.905,-0.429)};
\addplot[equipot] coordinates {(-0.726,-0.692) (-0.759,-0.678) (-0.790,-0.663) (-0.821,-0.646) (-0.852,-0.627) (-0.881,-0.607) (-0.908,-0.585) (-0.935,-0.562) (-0.961,-0.538) (-0.985,-0.512) (-1.008,-0.485) (-1.030,-0.457) (-1.050,-0.428) (-1.069,-0.398) (-1.087,-0.367) (-1.103,-0.336) (-1.117,-0.304) (-1.130,-0.271) (-1.142,-0.237) (-1.152,-0.203) (-1.160,-0.169) (-1.167,-0.134) (-1.172,-0.099) (-1.176,-0.064) (-1.178,-0.029) (-1.178,0.007) (-1.177,0.042) (-1.175,0.077) (-1.170,0.112) (-1.165,0.147) (-1.157,0.182) (-1.148,0.216) (-1.138,0.250) (-1.126,0.283) (-1.112,0.316) (-1.097,0.348) (-1.080,0.379) (-1.062,0.409) (-1.043,0.439) (-1.022,0.468) (-1.000,0.495) (-0.976,0.522) (-0.951,0.547) (-0.925,0.571) (-0.898,0.593) (-0.870,0.615) (-0.840,0.634) (-0.810,0.652) (-0.779,0.669) (-0.746,0.683) (-0.720,0.683) (-0.720,0.648) (-0.720,0.613) (-0.720,0.577) (-0.720,0.542) (-0.720,0.506) (-0.720,0.471) (-0.720,0.436) (-0.720,0.400) (-0.720,0.365) (-0.720,0.330) (-0.720,0.294) (-0.720,0.259) (-0.720,0.223) (-0.720,0.188) (-0.720,0.153) (-0.720,0.117) (-0.720,0.082) (-0.720,0.046) (-0.720,0.011) (-0.720,-0.024) (-0.720,-0.060) (-0.720,-0.095) (-0.720,-0.130) (-0.720,-0.166) (-0.720,-0.201) (-0.720,-0.237) (-0.720,-0.272) (-0.720,-0.307) (-0.720,-0.343) (-0.720,-0.378) (-0.720,-0.414) (-0.720,-0.449) (-0.720,-0.484) (-0.720,-0.520) (-0.720,-0.555) (-0.720,-0.590) (-0.720,-0.626) (-0.720,-0.661) (-0.726,-0.692)};
\addplot[equipot] coordinates {(-0.662,-0.841) (-0.697,-0.836) (-0.732,-0.830) (-0.766,-0.822) (-0.799,-0.812) (-0.833,-0.801) (-0.865,-0.788) (-0.897,-0.773) (-0.928,-0.757) (-0.959,-0.740) (-0.988,-0.721) (-1.017,-0.700) (-1.045,-0.679) (-1.072,-0.656) (-1.097,-0.633) (-1.122,-0.608) (-1.146,-0.582) (-1.168,-0.555) (-1.190,-0.527) (-1.210,-0.498) (-1.229,-0.469) (-1.247,-0.439) (-1.263,-0.408) (-1.279,-0.376) (-1.293,-0.344) (-1.306,-0.311) (-1.317,-0.278) (-1.327,-0.244) (-1.336,-0.210) (-1.344,-0.176) (-1.350,-0.142) (-1.355,-0.107) (-1.358,-0.072) (-1.360,-0.037) (-1.361,-0.002) (-1.360,0.033) (-1.358,0.068) (-1.355,0.103) (-1.350,0.138) (-1.344,0.173) (-1.337,0.207) (-1.328,0.241) (-1.318,0.275) (-1.307,0.308) (-1.294,0.341) (-1.280,0.373) (-1.265,0.404) (-1.249,0.435) (-1.231,0.466) (-1.212,0.495) (-1.192,0.524) (-1.171,0.552) (-1.148,0.579) (-1.125,0.605) (-1.100,0.630) (-1.074,0.654) (-1.048,0.677) (-1.020,0.698) (-0.991,0.719) (-0.962,0.738) (-0.932,0.755) (-0.900,0.772) (-0.869,0.786) (-0.836,0.800) (-0.803,0.811) (-0.769,0.821) (-0.735,0.829) (-0.701,0.836) (-0.666,0.841) (-0.631,0.843) (-0.596,0.844) (-0.561,0.843) (-0.540,0.827) (-0.540,0.792) (-0.540,0.756) (-0.540,0.721) (-0.540,0.686) (-0.540,0.651) (-0.540,0.616) (-0.540,0.581) (-0.540,0.546) (-0.540,0.511) (-0.540,0.476) (-0.540,0.441) (-0.540,0.405) (-0.540,0.370) (-0.540,0.335) (-0.540,0.300) (-0.540,0.265) (-0.540,0.230) (-0.540,0.195) (-0.540,0.160) (-0.540,0.125) (-0.540,0.090) (-0.540,0.054) (-0.540,0.019) (-0.540,-0.016) (-0.540,-0.051) (-0.540,-0.086) (-0.540,-0.121) (-0.540,-0.156) (-0.540,-0.191) (-0.540,-0.226) (-0.540,-0.261) (-0.540,-0.297) (-0.540,-0.332) (-0.540,-0.367) (-0.540,-0.402) (-0.540,-0.437) (-0.540,-0.472) (-0.540,-0.507) (-0.540,-0.542) (-0.540,-0.577) (-0.540,-0.613) (-0.540,-0.648) (-0.540,-0.683) (-0.540,-0.718) (-0.540,-0.753) (-0.540,-0.788) (-0.540,-0.823) (-0.557,-0.843) (-0.592,-0.844) (-0.627,-0.844) (-0.662,-0.841)};
\addplot[equipot] coordinates {(-0.732,-1.034) (-0.767,-1.033) (-0.802,-1.031) (-0.837,-1.027) (-0.872,-1.021) (-0.906,-1.015) (-0.940,-1.006) (-0.974,-0.996) (-1.007,-0.985) (-1.040,-0.973) (-1.073,-0.959) (-1.104,-0.944) (-1.136,-0.928) (-1.166,-0.911) (-1.196,-0.892) (-1.225,-0.873) (-1.254,-0.852) (-1.281,-0.830) (-1.308,-0.808) (-1.334,-0.784) (-1.360,-0.759) (-1.384,-0.734) (-1.407,-0.708) (-1.430,-0.681) (-1.451,-0.653) (-1.472,-0.625) (-1.492,-0.595) (-1.510,-0.566) (-1.528,-0.535) (-1.545,-0.504) (-1.560,-0.473) (-1.575,-0.441) (-1.589,-0.408) (-1.601,-0.375) (-1.612,-0.342) (-1.623,-0.309) (-1.632,-0.275) (-1.640,-0.241) (-1.647,-0.206) (-1.653,-0.171) (-1.658,-0.137) (-1.662,-0.102) (-1.665,-0.067) (-1.666,-0.031) (-1.667,0.004) (-1.666,0.039) (-1.664,0.074) (-1.661,0.109) (-1.657,0.144) (-1.652,0.179) (-1.646,0.213) (-1.639,0.248) (-1.630,0.282) (-1.621,0.316) (-1.610,0.349) (-1.599,0.382) (-1.586,0.415) (-1.572,0.448) (-1.557,0.479) (-1.541,0.511) (-1.524,0.542) (-1.507,0.572) (-1.488,0.602) (-1.468,0.631) (-1.447,0.659) (-1.425,0.687) (-1.402,0.713) (-1.379,0.739) (-1.354,0.765) (-1.329,0.789) (-1.303,0.812) (-1.276,0.835) (-1.248,0.856) (-1.219,0.877) (-1.190,0.896) (-1.160,0.914) (-1.129,0.932) (-1.098,0.947) (-1.066,0.962) (-1.033,0.976) (-1.000,0.988) (-0.967,0.999) (-0.933,1.008) (-0.899,1.016) (-0.864,1.023) (-0.830,1.028) (-0.795,1.031) (-0.759,1.033) (-0.724,1.034) (-0.689,1.033) (-0.654,1.030) (-0.619,1.026) (-0.585,1.020) (-0.550,1.012) (-0.516,1.003) (-0.483,0.992) (-0.450,0.979) (-0.418,0.964) (-0.387,0.948) (-0.360,0.928) (-0.360,0.893) (-0.360,0.858) (-0.360,0.823) (-0.360,0.787) (-0.360,0.752) (-0.360,0.717) (-0.360,0.682) (-0.360,0.647) (-0.360,0.612) (-0.360,0.576) (-0.360,0.541) (-0.360,0.506) (-0.360,0.471) (-0.360,0.436) (-0.360,0.401) (-0.360,0.366) (-0.360,0.330) (-0.360,0.295) (-0.360,0.260) (-0.360,0.225) (-0.360,0.190) (-0.360,0.155) (-0.360,0.119) (-0.360,0.084) (-0.360,0.049) (-0.360,0.014) (-0.360,-0.021) (-0.360,-0.056) (-0.360,-0.092) (-0.360,-0.127) (-0.360,-0.162) (-0.360,-0.197) (-0.360,-0.232) (-0.360,-0.267) (-0.360,-0.303) (-0.360,-0.338) (-0.360,-0.373) (-0.360,-0.408) (-0.360,-0.443) (-0.360,-0.478) (-0.360,-0.514) (-0.360,-0.549) (-0.360,-0.584) (-0.360,-0.619) (-0.360,-0.654) (-0.360,-0.689) (-0.360,-0.725) (-0.360,-0.760) (-0.360,-0.795) (-0.360,-0.830) (-0.360,-0.865) (-0.360,-0.900) (-0.363,-0.935) (-0.393,-0.952) (-0.425,-0.968) (-0.457,-0.982) (-0.490,-0.994) (-0.523,-1.005) (-0.557,-1.014) (-0.592,-1.021) (-0.627,-1.027) (-0.661,-1.031) (-0.697,-1.033) (-0.732,-1.034)};
\addplot[equipot] coordinates {(-1.148,-1.456) (-1.183,-1.452) (-1.218,-1.447) (-1.252,-1.441) (-1.287,-1.434) (-1.321,-1.425) (-1.355,-1.416) (-1.388,-1.406) (-1.421,-1.395) (-1.454,-1.383) (-1.487,-1.370) (-1.519,-1.356) (-1.551,-1.341) (-1.582,-1.325) (-1.613,-1.309) (-1.644,-1.291) (-1.674,-1.273) (-1.703,-1.254) (-1.732,-1.234) (-1.761,-1.214) (-1.789,-1.193) (-1.816,-1.171) (-1.843,-1.148) (-1.869,-1.125) (-1.895,-1.101) (-1.920,-1.077) (-1.945,-1.051) (-1.968,-1.026) (-1.992,-0.999) (-2.014,-0.973) (-2.036,-0.945) (-2.057,-0.917) (-2.078,-0.889) (-2.098,-0.860) (-2.117,-0.831) (-2.136,-0.801) (-2.153,-0.771) (-2.170,-0.740) (-2.187,-0.709) (-2.203,-0.677) (-2.217,-0.646) (-2.232,-0.614) (-2.245,-0.581) (-2.258,-0.549) (-2.270,-0.516) (-2.281,-0.482) (-2.291,-0.449) (-2.301,-0.415) (-2.310,-0.381) (-2.318,-0.347) (-2.326,-0.313) (-2.332,-0.278) (-2.338,-0.244) (-2.343,-0.209) (-2.347,-0.174) (-2.351,-0.139) (-2.354,-0.104) (-2.356,-0.069) (-2.357,-0.034) (-2.357,0.001) (-2.357,0.036) (-2.355,0.071) (-2.353,0.106) (-2.351,0.141) (-2.347,0.176) (-2.343,0.211) (-2.338,0.246) (-2.332,0.280) (-2.325,0.315) (-2.318,0.349) (-2.310,0.383) (-2.301,0.417) (-2.291,0.451) (-2.280,0.484) (-2.269,0.517) (-2.257,0.550) (-2.244,0.583) (-2.231,0.615) (-2.217,0.647) (-2.202,0.679) (-2.186,0.711) (-2.170,0.742) (-2.152,0.772) (-2.135,0.802) (-2.116,0.832) (-2.097,0.862) (-2.077,0.890) (-2.056,0.919) (-2.035,0.947) (-2.013,0.974) (-1.990,1.001) (-1.967,1.027) (-1.943,1.053) (-1.919,1.078) (-1.894,1.102) (-1.868,1.126) (-1.842,1.150) (-1.815,1.172) (-1.787,1.194) (-1.759,1.215) (-1.731,1.235) (-1.702,1.255) (-1.672,1.274) (-1.642,1.292) (-1.611,1.309) (-1.580,1.326) (-1.549,1.342) (-1.517,1.356) (-1.485,1.370) (-1.452,1.383) (-1.420,1.395) (-1.386,1.406) (-1.353,1.417) (-1.319,1.426) (-1.285,1.434) (-1.250,1.441) (-1.216,1.447) (-1.181,1.453) (-1.146,1.457) (-1.111,1.460) (-1.076,1.461) (-1.041,1.462) (-1.006,1.462) (-0.971,1.460) (-0.936,1.458) (-0.901,1.454) (-0.866,1.449) (-0.832,1.443) (-0.798,1.436) (-0.763,1.427) (-0.730,1.418) (-0.696,1.407) (-0.663,1.395) (-0.631,1.382) (-0.599,1.367) (-0.567,1.352) (-0.536,1.335) (-0.506,1.318) (-0.477,1.299) (-0.448,1.279) (-0.420,1.258) (-0.392,1.236) (-0.366,1.213) (-0.340,1.189) (-0.316,1.164) (-0.292,1.138) (-0.269,1.111) (-0.247,1.084) (-0.226,1.056) (-0.207,1.027) (-0.188,0.997) (-0.180,0.964) (-0.180,0.929) (-0.180,0.894) (-0.180,0.859) (-0.180,0.824) (-0.180,0.789) (-0.180,0.753) (-0.180,0.718) (-0.180,0.683) (-0.180,0.648) (-0.180,0.613) (-0.180,0.578) (-0.180,0.543) (-0.180,0.508) (-0.180,0.473) (-0.180,0.438) (-0.180,0.403) (-0.180,0.367) (-0.180,0.332) (-0.180,0.297) (-0.180,0.262) (-0.180,0.227) (-0.180,0.192) (-0.180,0.157) (-0.180,0.122) (-0.180,0.087) (-0.180,0.052) (-0.180,0.017) (-0.180,-0.019) (-0.180,-0.054) (-0.180,-0.089) (-0.180,-0.124) (-0.180,-0.159) (-0.180,-0.194) (-0.180,-0.229) (-0.180,-0.264) (-0.180,-0.299) (-0.180,-0.334) (-0.180,-0.369) (-0.180,-0.404) (-0.180,-0.440) (-0.180,-0.475) (-0.180,-0.510) (-0.180,-0.545) (-0.180,-0.580) (-0.180,-0.615) (-0.180,-0.650) (-0.180,-0.685) (-0.180,-0.720) (-0.180,-0.755) (-0.180,-0.790) (-0.180,-0.826) (-0.180,-0.861) (-0.180,-0.896) (-0.180,-0.931) (-0.180,-0.966) (-0.189,-0.999) (-0.208,-1.028) (-0.228,-1.057) (-0.248,-1.085) (-0.270,-1.113) (-0.293,-1.139) (-0.317,-1.165) (-0.342,-1.190) (-0.367,-1.214) (-0.394,-1.237) (-0.421,-1.259) (-0.449,-1.280) (-0.478,-1.300) (-0.508,-1.319) (-0.538,-1.336) (-0.569,-1.353) (-0.601,-1.368) (-0.633,-1.383) (-0.665,-1.396) (-0.698,-1.408) (-0.732,-1.418) (-0.765,-1.428) (-0.799,-1.436) (-0.834,-1.443) (-0.868,-1.449) (-0.903,-1.454) (-0.938,-1.458) (-0.973,-1.461) (-1.008,-1.462) (-1.043,-1.462) (-1.078,-1.461) (-1.113,-1.459) (-1.148,-1.456)};
\addplot[equipot] coordinates {(0.923,-1.457) (0.958,-1.460) (0.993,-1.462) (1.028,-1.462) (1.063,-1.462) (1.098,-1.460) (1.133,-1.458) (1.168,-1.454) (1.202,-1.450) (1.237,-1.444) (1.271,-1.437) (1.306,-1.429) (1.340,-1.420) (1.373,-1.411) (1.407,-1.400) (1.440,-1.388) (1.473,-1.375) (1.505,-1.362) (1.537,-1.347) (1.568,-1.332) (1.600,-1.316) (1.630,-1.299) (1.661,-1.281) (1.690,-1.263) (1.720,-1.243) (1.748,-1.223) (1.776,-1.202) (1.804,-1.181) (1.831,-1.158) (1.858,-1.135) (1.884,-1.112) (1.909,-1.087) (1.934,-1.063) (1.958,-1.037) (1.981,-1.011) (2.004,-0.984) (2.027,-0.957) (2.048,-0.930) (2.069,-0.901) (2.089,-0.873) (2.109,-0.844) (2.127,-0.814) (2.146,-0.784) (2.163,-0.753) (2.180,-0.723) (2.196,-0.691) (2.211,-0.660) (2.225,-0.628) (2.239,-0.596) (2.252,-0.563) (2.265,-0.530) (2.276,-0.497) (2.287,-0.464) (2.297,-0.430) (2.306,-0.396) (2.315,-0.362) (2.322,-0.328) (2.329,-0.293) (2.336,-0.259) (2.341,-0.224) (2.346,-0.189) (2.349,-0.155) (2.352,-0.120) (2.355,-0.085) (2.356,-0.050) (2.357,-0.014) (2.357,0.021) (2.356,0.056) (2.354,0.091) (2.352,0.126) (2.349,0.161) (2.345,0.196) (2.340,0.230) (2.335,0.265) (2.328,0.299) (2.321,0.334) (2.313,0.368) (2.305,0.402) (2.295,0.436) (2.285,0.469) (2.274,0.503) (2.262,0.536) (2.250,0.569) (2.237,0.601) (2.223,0.633) (2.208,0.665) (2.193,0.697) (2.177,0.728) (2.160,0.759) (2.143,0.789) (2.124,0.819) (2.105,0.849) (2.086,0.878) (2.065,0.906) (2.044,0.934) (2.023,0.962) (2.000,0.989) (1.977,1.016) (1.954,1.042) (1.930,1.067) (1.905,1.092) (1.879,1.116) (1.853,1.139) (1.827,1.162) (1.799,1.184) (1.772,1.206) (1.743,1.227) (1.714,1.247) (1.685,1.266) (1.655,1.284) (1.625,1.302) (1.594,1.319) (1.563,1.335) (1.531,1.350) (1.499,1.364) (1.467,1.378) (1.434,1.390) (1.401,1.402) (1.367,1.412) (1.334,1.422) (1.300,1.431) (1.266,1.438) (1.231,1.445) (1.196,1.450) (1.162,1.455) (1.127,1.458) (1.092,1.461) (1.057,1.462) (1.022,1.462) (0.986,1.461) (0.951,1.459) (0.917,1.456) (0.882,1.452) (0.847,1.446) (0.813,1.439) (0.778,1.431) (0.745,1.422) (0.711,1.412) (0.678,1.400) (0.645,1.388) (0.613,1.374) (0.581,1.359) (0.550,1.343) (0.519,1.326) (0.490,1.307) (0.460,1.288) (0.432,1.267) (0.404,1.246) (0.377,1.223) (0.351,1.199) (0.326,1.175) (0.302,1.149) (0.279,1.123) (0.257,1.096) (0.236,1.068) (0.215,1.039) (0.196,1.010) (0.180,0.979) (0.180,0.944) (0.180,0.909) (0.180,0.874) (0.180,0.839) (0.180,0.804) (0.180,0.769) (0.180,0.734) (0.180,0.699) (0.180,0.664) (0.180,0.629) (0.180,0.593) (0.180,0.558) (0.180,0.523) (0.180,0.488) (0.180,0.453) (0.180,0.418) (0.180,0.383) (0.180,0.348) (0.180,0.313) (0.180,0.278) (0.180,0.243) (0.180,0.207) (0.180,0.172) (0.180,0.137) (0.180,0.102) (0.180,0.067) (0.180,0.032) (0.180,-0.003) (0.180,-0.038) (0.180,-0.073) (0.180,-0.108) (0.180,-0.143) (0.180,-0.178) (0.180,-0.214) (0.180,-0.249) (0.180,-0.284) (0.180,-0.319) (0.180,-0.354) (0.180,-0.389) (0.180,-0.424) (0.180,-0.459) (0.180,-0.494) (0.180,-0.529) (0.180,-0.564) (0.180,-0.600) (0.180,-0.635) (0.180,-0.670) (0.180,-0.705) (0.180,-0.740) (0.180,-0.775) (0.180,-0.810) (0.180,-0.845) (0.180,-0.880) (0.180,-0.915) (0.180,-0.950) (0.181,-0.985) (0.199,-1.015) (0.219,-1.044) (0.239,-1.073) (0.261,-1.101) (0.283,-1.128) (0.306,-1.154) (0.331,-1.179) (0.356,-1.204) (0.382,-1.227) (0.409,-1.249) (0.437,-1.271) (0.465,-1.291) (0.495,-1.310) (0.525,-1.329) (0.555,-1.346) (0.587,-1.362) (0.618,-1.376) (0.651,-1.390) (0.684,-1.402) (0.717,-1.414) (0.750,-1.424) (0.784,-1.433) (0.819,-1.440) (0.853,-1.447) (0.888,-1.452) (0.923,-1.457)};
\addplot[equipot] coordinates {(0.732,-1.034) (0.767,-1.033) (0.802,-1.031) (0.837,-1.027) (0.872,-1.021) (0.906,-1.015) (0.940,-1.006) (0.974,-0.996) (1.007,-0.985) (1.040,-0.973) (1.073,-0.959) (1.104,-0.944) (1.136,-0.928) (1.166,-0.911) (1.196,-0.892) (1.225,-0.873) (1.254,-0.852) (1.281,-0.830) (1.308,-0.808) (1.334,-0.784) (1.360,-0.759) (1.384,-0.734) (1.407,-0.708) (1.430,-0.681) (1.451,-0.653) (1.472,-0.625) (1.492,-0.595) (1.510,-0.566) (1.528,-0.535) (1.545,-0.504) (1.560,-0.473) (1.575,-0.441) (1.589,-0.408) (1.601,-0.375) (1.612,-0.342) (1.623,-0.309) (1.632,-0.275) (1.640,-0.241) (1.647,-0.206) (1.653,-0.171) (1.658,-0.137) (1.662,-0.102) (1.665,-0.067) (1.666,-0.031) (1.667,0.004) (1.666,0.039) (1.664,0.074) (1.661,0.109) (1.657,0.144) (1.652,0.179) (1.646,0.213) (1.639,0.248) (1.630,0.282) (1.621,0.316) (1.610,0.349) (1.599,0.382) (1.586,0.415) (1.572,0.448) (1.557,0.479) (1.541,0.511) (1.524,0.542) (1.507,0.572) (1.488,0.602) (1.468,0.631) (1.447,0.659) (1.425,0.687) (1.402,0.713) (1.379,0.739) (1.354,0.765) (1.329,0.789) (1.303,0.812) (1.276,0.835) (1.248,0.856) (1.219,0.877) (1.190,0.896) (1.160,0.914) (1.129,0.932) (1.098,0.947) (1.066,0.962) (1.033,0.976) (1.000,0.988) (0.967,0.999) (0.933,1.008) (0.899,1.016) (0.864,1.023) (0.830,1.028) (0.795,1.031) (0.759,1.033) (0.724,1.034) (0.689,1.033) (0.654,1.030) (0.619,1.026) (0.585,1.020) (0.550,1.012) (0.516,1.003) (0.483,0.992) (0.450,0.979) (0.418,0.964) (0.387,0.948) (0.360,0.928) (0.360,0.893) (0.360,0.858) (0.360,0.823) (0.360,0.787) (0.360,0.752) (0.360,0.717) (0.360,0.682) (0.360,0.647) (0.360,0.612) (0.360,0.576) (0.360,0.541) (0.360,0.506) (0.360,0.471) (0.360,0.436) (0.360,0.401) (0.360,0.366) (0.360,0.330) (0.360,0.295) (0.360,0.260) (0.360,0.225) (0.360,0.190) (0.360,0.155) (0.360,0.119) (0.360,0.084) (0.360,0.049) (0.360,0.014) (0.360,-0.021) (0.360,-0.056) (0.360,-0.092) (0.360,-0.127) (0.360,-0.162) (0.360,-0.197) (0.360,-0.232) (0.360,-0.267) (0.360,-0.303) (0.360,-0.338) (0.360,-0.373) (0.360,-0.408) (0.360,-0.443) (0.360,-0.478) (0.360,-0.514) (0.360,-0.549) (0.360,-0.584) (0.360,-0.619) (0.360,-0.654) (0.360,-0.689) (0.360,-0.725) (0.360,-0.760) (0.360,-0.795) (0.360,-0.830) (0.360,-0.865) (0.360,-0.900) (0.363,-0.935) (0.393,-0.952) (0.425,-0.968) (0.457,-0.982) (0.490,-0.994) (0.523,-1.005) (0.557,-1.014) (0.592,-1.021) (0.627,-1.027) (0.661,-1.031) (0.697,-1.033) (0.732,-1.034)};
\addplot[equipot] coordinates {(0.541,-0.841) (0.576,-0.844) (0.611,-0.844) (0.646,-0.842) (0.681,-0.839) (0.716,-0.833) (0.750,-0.826) (0.784,-0.817) (0.817,-0.806) (0.850,-0.794) (0.882,-0.780) (0.914,-0.765) (0.945,-0.748) (0.975,-0.730) (1.004,-0.710) (1.032,-0.689) (1.059,-0.667) (1.086,-0.644) (1.111,-0.619) (1.135,-0.594) (1.158,-0.567) (1.180,-0.540) (1.201,-0.512) (1.220,-0.483) (1.239,-0.453) (1.256,-0.422) (1.272,-0.391) (1.286,-0.359) (1.300,-0.327) (1.312,-0.294) (1.323,-0.260) (1.332,-0.226) (1.340,-0.192) (1.347,-0.158) (1.352,-0.123) (1.357,-0.088) (1.359,-0.053) (1.361,-0.018) (1.361,0.017) (1.359,0.052) (1.357,0.087) (1.353,0.122) (1.347,0.157) (1.340,0.191) (1.332,0.225) (1.323,0.259) (1.312,0.292) (1.300,0.325) (1.287,0.358) (1.272,0.390) (1.256,0.421) (1.239,0.452) (1.221,0.482) (1.201,0.511) (1.181,0.539) (1.159,0.567) (1.136,0.593) (1.112,0.619) (1.086,0.643) (1.060,0.666) (1.033,0.688) (1.005,0.709) (0.976,0.729) (0.946,0.747) (0.915,0.764) (0.884,0.780) (0.851,0.794) (0.818,0.806) (0.785,0.817) (0.751,0.826) (0.717,0.833) (0.682,0.839) (0.647,0.842) (0.612,0.844) (0.577,0.844) (0.542,0.842) (0.540,0.808) (0.540,0.773) (0.540,0.738) (0.540,0.703) (0.540,0.668) (0.540,0.632) (0.540,0.597) (0.540,0.562) (0.540,0.527) (0.540,0.492) (0.540,0.457) (0.540,0.422) (0.540,0.387) (0.540,0.352) (0.540,0.317) (0.540,0.281) (0.540,0.246) (0.540,0.211) (0.540,0.176) (0.540,0.141) (0.540,0.106) (0.540,0.071) (0.540,0.036) (0.540,0.001) (0.540,-0.034) (0.540,-0.070) (0.540,-0.105) (0.540,-0.140) (0.540,-0.175) (0.540,-0.210) (0.540,-0.245) (0.540,-0.280) (0.540,-0.315) (0.540,-0.350) (0.540,-0.385) (0.540,-0.421) (0.540,-0.456) (0.540,-0.491) (0.540,-0.526) (0.540,-0.561) (0.540,-0.596) (0.540,-0.631) (0.540,-0.666) (0.540,-0.701) (0.540,-0.737) (0.540,-0.772) (0.540,-0.807) (0.541,-0.841)};
\addplot[equipot] coordinates {(0.720,-0.691) (0.753,-0.681) (0.785,-0.666) (0.816,-0.649) (0.847,-0.630) (0.876,-0.610) (0.904,-0.589) (0.931,-0.566) (0.957,-0.542) (0.981,-0.516) (1.004,-0.490) (1.026,-0.462) (1.047,-0.433) (1.066,-0.403) (1.084,-0.373) (1.100,-0.341) (1.115,-0.309) (1.128,-0.276) (1.140,-0.243) (1.150,-0.209) (1.159,-0.175) (1.166,-0.140) (1.171,-0.105) (1.175,-0.070) (1.178,-0.035) (1.179,0.001) (1.178,0.036) (1.175,0.071) (1.171,0.107) (1.166,0.141) (1.158,0.176) (1.150,0.210) (1.139,0.244) (1.128,0.278) (1.114,0.310) (1.099,0.342) (1.083,0.374) (1.065,0.404) (1.046,0.434) (1.025,0.463) (1.003,0.491) (0.980,0.517) (0.956,0.543) (0.930,0.567) (0.903,0.590) (0.875,0.611) (0.845,0.631) (0.815,0.649) (0.784,0.666) (0.752,0.681) (0.720,0.689) (0.720,0.654) (0.720,0.619) (0.720,0.583) (0.720,0.548) (0.720,0.512) (0.720,0.477) (0.720,0.442) (0.720,0.406) (0.720,0.371) (0.720,0.335) (0.720,0.300) (0.720,0.265) (0.720,0.229) (0.720,0.194) (0.720,0.159) (0.720,0.123) (0.720,0.088) (0.720,0.052) (0.720,0.017) (0.720,-0.018) (0.720,-0.054) (0.720,-0.089) (0.720,-0.125) (0.720,-0.160) (0.720,-0.195) (0.720,-0.231) (0.720,-0.266) (0.720,-0.301) (0.720,-0.337) (0.720,-0.372) (0.720,-0.408) (0.720,-0.443) (0.720,-0.478) (0.720,-0.514) (0.720,-0.549) (0.720,-0.585) (0.720,-0.620) (0.720,-0.655) (0.720,-0.691)};
\addplot[equipot] coordinates {(0.905,-0.429) (0.926,-0.401) (0.946,-0.372) (0.964,-0.342) (0.981,-0.311) (0.996,-0.280) (1.009,-0.247) (1.021,-0.214) (1.031,-0.180) (1.039,-0.146) (1.045,-0.112) (1.050,-0.077) (1.053,-0.042) (1.054,-0.007) (1.054,0.028) (1.051,0.063) (1.047,0.098) (1.041,0.133) (1.034,0.167) (1.025,0.201) (1.014,0.234) (1.001,0.267) (0.987,0.299) (0.971,0.330) (0.953,0.361) (0.934,0.390) (0.914,0.419) (0.900,0.411) (0.900,0.375) (0.900,0.340) (0.900,0.305) (0.900,0.270) (0.900,0.235) (0.900,0.200) (0.900,0.165) (0.900,0.130) (0.900,0.095) (0.900,0.059) (0.900,0.024) (0.900,-0.011) (0.900,-0.046) (0.900,-0.081) (0.900,-0.116) (0.900,-0.151) (0.900,-0.186) (0.900,-0.221) (0.900,-0.257) (0.900,-0.292) (0.900,-0.327) (0.900,-0.362) (0.900,-0.397) (0.905,-0.429)};
\addplot[lineaD, flechaen=0.094, flechaen=0.906] coordinates {(-2.600,1.029) (-2.570,1.010) (-2.541,0.991) (-2.511,0.972) (-2.481,0.953) (-2.451,0.935) (-2.421,0.916) (-2.391,0.898) (-2.361,0.880) (-2.331,0.862) (-2.301,0.844) (-2.271,0.826) (-2.240,0.809) (-2.210,0.791) (-2.179,0.774) (-2.149,0.756) (-2.118,0.739) (-2.087,0.722) (-2.056,0.705) (-2.026,0.688) (-1.995,0.672) (-1.964,0.655) (-1.932,0.639) (-1.901,0.623) (-1.870,0.606) (-1.839,0.590) (-1.807,0.575) (-1.776,0.559) (-1.744,0.543) (-1.713,0.528) (-1.681,0.513) (-1.649,0.498) (-1.618,0.483) (-1.586,0.468) (-1.554,0.453) (-1.522,0.438) (-1.490,0.424) (-1.458,0.410) (-1.425,0.396) (-1.393,0.382) (-1.361,0.368) (-1.328,0.355) (-1.296,0.341) (-1.263,0.328) (-1.231,0.315) (-1.198,0.302) (-1.165,0.289) (-1.132,0.277) (-1.099,0.265) (-1.066,0.252) (-1.033,0.240) (-1.000,0.229) (-0.967,0.220) (-0.931,0.220) (-0.896,0.220) (-0.861,0.220) (-0.826,0.220) (-0.791,0.220) (-0.756,0.220) (-0.721,0.220) (-0.685,0.220) (-0.650,0.220) (-0.615,0.220) (-0.580,0.220) (-0.545,0.220) (-0.510,0.220) (-0.475,0.220) (-0.439,0.220) (-0.404,0.220) (-0.369,0.220) (-0.334,0.220) (-0.299,0.220) (-0.264,0.220) (-0.228,0.220) (-0.193,0.220) (-0.158,0.220) (-0.123,0.220) (-0.088,0.220) (-0.053,0.220) (-0.018,0.220) (0.018,0.220) (0.053,0.220) (0.088,0.220) (0.123,0.220) (0.158,0.220) (0.193,0.220) (0.228,0.220) (0.264,0.220) (0.299,0.220) (0.334,0.220) (0.369,0.220) (0.404,0.220) (0.439,0.220) (0.475,0.220) (0.510,0.220) (0.545,0.220) (0.580,0.220) (0.615,0.220) (0.650,0.220) (0.685,0.220) (0.721,0.220) (0.756,0.220) (0.791,0.220) (0.826,0.220) (0.861,0.220) (0.896,0.220) (0.931,0.220) (0.967,0.220) (1.000,0.229) (1.033,0.240) (1.066,0.252) (1.099,0.265) (1.132,0.277) (1.165,0.289) (1.198,0.302) (1.231,0.315) (1.263,0.328) (1.296,0.341) (1.328,0.355) (1.361,0.368) (1.393,0.382) (1.425,0.396) (1.458,0.410) (1.490,0.424) (1.522,0.438) (1.554,0.453) (1.586,0.468) (1.618,0.483) (1.649,0.498) (1.681,0.513) (1.713,0.528) (1.744,0.543) (1.776,0.559) (1.807,0.575) (1.839,0.590) (1.870,0.606) (1.901,0.623) (1.932,0.639) (1.964,0.655) (1.995,0.672) (2.026,0.688) (2.056,0.705) (2.087,0.722) (2.118,0.739) (2.149,0.756) (2.179,0.774) (2.210,0.791) (2.240,0.809) (2.271,0.826) (2.301,0.844) (2.331,0.862) (2.361,0.880) (2.391,0.898) (2.421,0.916) (2.451,0.935) (2.481,0.953) (2.511,0.972) (2.541,0.991) (2.570,1.010) (2.600,1.029)};
\addplot[lineaD, flechaen=0.044, flechaen=0.956] coordinates {(-2.270,2.000) (-2.256,1.968) (-2.240,1.936) (-2.225,1.905) (-2.209,1.873) (-2.193,1.842) (-2.177,1.811) (-2.161,1.780) (-2.144,1.749) (-2.127,1.718) (-2.110,1.688) (-2.093,1.657) (-2.075,1.627) (-2.057,1.597) (-2.039,1.567) (-2.020,1.537) (-2.002,1.507) (-1.983,1.477) (-1.963,1.448) (-1.944,1.419) (-1.924,1.390) (-1.904,1.361) (-1.884,1.332) (-1.864,1.303) (-1.843,1.275) (-1.822,1.247) (-1.801,1.219) (-1.780,1.191) (-1.758,1.163) (-1.736,1.135) (-1.714,1.108) (-1.692,1.081) (-1.669,1.054) (-1.646,1.027) (-1.623,1.001) (-1.600,0.975) (-1.577,0.948) (-1.553,0.923) (-1.529,0.897) (-1.505,0.871) (-1.480,0.846) (-1.456,0.821) (-1.431,0.796) (-1.406,0.772) (-1.380,0.747) (-1.355,0.723) (-1.329,0.699) (-1.303,0.676) (-1.277,0.652) (-1.250,0.629) (-1.224,0.607) (-1.197,0.584) (-1.170,0.562) (-1.142,0.540) (-1.115,0.518) (-1.087,0.496) (-1.059,0.475) (-1.031,0.454) (-1.002,0.434) (-0.974,0.413) (-0.945,0.393) (-0.913,0.380) (-0.878,0.380) (-0.843,0.380) (-0.808,0.380) (-0.773,0.380) (-0.738,0.380) (-0.703,0.380) (-0.667,0.380) (-0.632,0.380) (-0.597,0.380) (-0.562,0.380) (-0.527,0.380) (-0.492,0.380) (-0.457,0.380) (-0.422,0.380) (-0.386,0.380) (-0.351,0.380) (-0.316,0.380) (-0.281,0.380) (-0.246,0.380) (-0.211,0.380) (-0.176,0.380) (-0.141,0.380) (-0.105,0.380) (-0.070,0.380) (-0.035,0.380) (-0.000,0.380) (0.035,0.380) (0.070,0.380) (0.105,0.380) (0.141,0.380) (0.176,0.380) (0.211,0.380) (0.246,0.380) (0.281,0.380) (0.316,0.380) (0.351,0.380) (0.386,0.380) (0.422,0.380) (0.457,0.380) (0.492,0.380) (0.527,0.380) (0.562,0.380) (0.597,0.380) (0.632,0.380) (0.667,0.380) (0.703,0.380) (0.738,0.380) (0.773,0.380) (0.808,0.380) (0.843,0.380) (0.878,0.380) (0.913,0.380) (0.945,0.393) (0.974,0.413) (1.002,0.434) (1.031,0.454) (1.059,0.475) (1.087,0.496) (1.115,0.518) (1.142,0.540) (1.170,0.562) (1.197,0.584) (1.224,0.607) (1.250,0.629) (1.277,0.652) (1.303,0.676) (1.329,0.699) (1.355,0.723) (1.380,0.747) (1.406,0.772) (1.431,0.796) (1.456,0.821) (1.480,0.846) (1.505,0.871) (1.529,0.897) (1.553,0.923) (1.577,0.948) (1.600,0.975) (1.623,1.001) (1.646,1.027) (1.669,1.054) (1.692,1.081) (1.714,1.108) (1.736,1.135) (1.758,1.163) (1.780,1.191) (1.801,1.219) (1.822,1.247) (1.843,1.275) (1.864,1.303) (1.884,1.332) (1.904,1.361) (1.924,1.390) (1.944,1.419) (1.963,1.448) (1.983,1.477) (2.002,1.507) (2.020,1.537) (2.039,1.567) (2.057,1.597) (2.075,1.627) (2.093,1.657) (2.110,1.688) (2.127,1.718) (2.144,1.749) (2.161,1.780) (2.177,1.811) (2.193,1.842) (2.209,1.873) (2.225,1.905) (2.240,1.936) (2.256,1.968) (2.270,2.000)};
\addplot[lineaD, flechaen=0.399] coordinates {(-1.423,2.000) (-1.423,1.965) (-1.422,1.930) (-1.420,1.895) (-1.418,1.860) (-1.415,1.825) (-1.412,1.790) (-1.408,1.755) (-1.404,1.720) (-1.399,1.685) (-1.393,1.651) (-1.387,1.616) (-1.380,1.582) (-1.373,1.547) (-1.366,1.513) (-1.357,1.479) (-1.349,1.445) (-1.339,1.411) (-1.330,1.378) (-1.319,1.344) (-1.309,1.311) (-1.297,1.277) (-1.285,1.244) (-1.273,1.212) (-1.260,1.179) (-1.247,1.146) (-1.233,1.114) (-1.219,1.082) (-1.204,1.050) (-1.189,1.019) (-1.173,0.987) (-1.157,0.956) (-1.140,0.925) (-1.123,0.895) (-1.105,0.864) (-1.087,0.834) (-1.069,0.805) (-1.050,0.775) (-1.030,0.746) (-1.010,0.717) (-0.990,0.689) (-0.969,0.660) (-0.948,0.632) (-0.926,0.605) (-0.904,0.578) (-0.881,0.551) (-0.858,0.524) (-0.825,0.520) (-0.789,0.520) (-0.754,0.520) (-0.719,0.520) (-0.684,0.520) (-0.649,0.520) (-0.614,0.520) (-0.579,0.520) (-0.544,0.520) (-0.509,0.520) (-0.474,0.520) (-0.439,0.520) (-0.404,0.520) (-0.368,0.520) (-0.333,0.520) (-0.298,0.520) (-0.263,0.520) (-0.228,0.520) (-0.193,0.520) (-0.158,0.520) (-0.123,0.520) (-0.088,0.520) (-0.053,0.520) (-0.018,0.520) (0.018,0.520) (0.053,0.520) (0.088,0.520) (0.123,0.520) (0.158,0.520) (0.193,0.520) (0.228,0.520) (0.263,0.520) (0.298,0.520) (0.333,0.520) (0.368,0.520) (0.404,0.520) (0.439,0.520) (0.474,0.520) (0.509,0.520) (0.544,0.520) (0.579,0.520) (0.614,0.520) (0.649,0.520) (0.684,0.520) (0.719,0.520) (0.754,0.520) (0.789,0.520) (0.825,0.520) (0.858,0.524) (0.881,0.551) (0.904,0.578) (0.926,0.605) (0.948,0.632) (0.969,0.660) (0.990,0.689) (1.010,0.717) (1.030,0.746) (1.050,0.775) (1.069,0.805) (1.087,0.834) (1.105,0.864) (1.123,0.895) (1.140,0.925) (1.157,0.956) (1.173,0.987) (1.189,1.019) (1.204,1.050) (1.219,1.082) (1.233,1.114) (1.247,1.146) (1.260,1.179) (1.273,1.212) (1.285,1.244) (1.297,1.277) (1.309,1.311) (1.319,1.344) (1.330,1.378) (1.339,1.411) (1.349,1.445) (1.357,1.479) (1.366,1.513) (1.373,1.547) (1.380,1.582) (1.387,1.616) (1.393,1.651) (1.399,1.685) (1.404,1.720) (1.408,1.755) (1.412,1.790) (1.415,1.825) (1.418,1.860) (1.420,1.895) (1.422,1.930) (1.423,1.965) (1.423,2.000)};
\addplot[lineaD, flechaen=0.614] coordinates {(-0.754,2.000) (-0.772,1.970) (-0.790,1.940) (-0.806,1.909) (-0.821,1.877) (-0.836,1.845) (-0.849,1.813) (-0.862,1.780) (-0.873,1.747) (-0.884,1.714) (-0.894,1.680) (-0.902,1.646) (-0.910,1.612) (-0.917,1.578) (-0.923,1.543) (-0.928,1.508) (-0.932,1.474) (-0.935,1.439) (-0.938,1.404) (-0.939,1.369) (-0.940,1.334) (-0.939,1.299) (-0.938,1.264) (-0.936,1.229) (-0.933,1.194) (-0.930,1.159) (-0.925,1.124) (-0.920,1.090) (-0.913,1.055) (-0.906,1.021) (-0.899,0.987) (-0.890,0.953) (-0.881,0.919) (-0.870,0.885) (-0.859,0.852) (-0.848,0.819) (-0.835,0.786) (-0.822,0.754) (-0.808,0.722) (-0.793,0.690) (-0.778,0.658) (-0.753,0.640) (-0.718,0.640) (-0.683,0.640) (-0.648,0.640) (-0.613,0.640) (-0.578,0.640) (-0.543,0.640) (-0.508,0.640) (-0.473,0.640) (-0.438,0.640) (-0.403,0.640) (-0.368,0.640) (-0.333,0.640) (-0.298,0.640) (-0.263,0.640) (-0.228,0.640) (-0.193,0.640) (-0.158,0.640) (-0.123,0.640) (-0.088,0.640) (-0.053,0.640) (-0.018,0.640) (0.018,0.640) (0.053,0.640) (0.088,0.640) (0.123,0.640) (0.158,0.640) (0.193,0.640) (0.228,0.640) (0.263,0.640) (0.298,0.640) (0.333,0.640) (0.368,0.640) (0.403,0.640) (0.438,0.640) (0.473,0.640) (0.508,0.640) (0.543,0.640) (0.578,0.640) (0.613,0.640) (0.648,0.640) (0.683,0.640) (0.718,0.640) (0.753,0.640) (0.778,0.658) (0.793,0.690) (0.808,0.722) (0.822,0.754) (0.835,0.786) (0.848,0.819) (0.859,0.852) (0.870,0.885) (0.881,0.919) (0.890,0.953) (0.899,0.987) (0.906,1.021) (0.913,1.055) (0.920,1.090) (0.925,1.124) (0.930,1.159) (0.933,1.194) (0.936,1.229) (0.938,1.264) (0.939,1.299) (0.940,1.334) (0.939,1.369) (0.938,1.404) (0.935,1.439) (0.932,1.474) (0.928,1.508) (0.923,1.543) (0.917,1.578) (0.910,1.612) (0.902,1.646) (0.894,1.680) (0.884,1.714) (0.873,1.747) (0.862,1.780) (0.849,1.813) (0.836,1.845) (0.821,1.877) (0.806,1.909) (0.790,1.940) (0.772,1.970) (0.754,2.000)};
\addplot[lineaD, flechaen=0.500] coordinates {(-0.657,0.750) (-0.621,0.750) (-0.586,0.750) (-0.551,0.750) (-0.516,0.750) (-0.481,0.750) (-0.446,0.750) (-0.410,0.750) (-0.375,0.750) (-0.340,0.750) (-0.305,0.750) (-0.270,0.750) (-0.234,0.750) (-0.199,0.750) (-0.164,0.750) (-0.129,0.750) (-0.094,0.750) (-0.059,0.750) (-0.023,0.750) (0.012,0.750) (0.047,0.750) (0.082,0.750) (0.117,0.750) (0.152,0.750) (0.188,0.750) (0.223,0.750) (0.258,0.750) (0.293,0.750) (0.328,0.750) (0.363,0.750) (0.399,0.750) (0.434,0.750) (0.469,0.750) (0.504,0.750) (0.539,0.750) (0.574,0.750) (0.610,0.750) (0.645,0.750) (0.665,0.770) (0.671,0.804) (0.676,0.839) (0.680,0.874) (0.683,0.909) (0.684,0.944) (0.684,0.979) (0.683,1.015) (0.681,1.050) (0.677,1.085) (0.673,1.119) (0.667,1.154) (0.660,1.189) (0.651,1.223) (0.641,1.256) (0.630,1.290) (0.617,1.323) (0.604,1.355) (0.589,1.387) (0.572,1.418) (0.554,1.448) (0.535,1.478) (0.515,1.506) (0.493,1.534) (0.470,1.561) (0.446,1.586) (0.420,1.610) (0.394,1.633) (0.366,1.655) (0.337,1.675) (0.307,1.693) (0.276,1.710) (0.245,1.725) (0.212,1.739) (0.179,1.750) (0.145,1.760) (0.111,1.767) (0.076,1.773) (0.041,1.776) (0.006,1.778) (-0.029,1.777) (-0.064,1.774) (-0.099,1.769) (-0.134,1.763) (-0.168,1.754) (-0.201,1.743) (-0.234,1.730) (-0.266,1.716) (-0.297,1.699) (-0.327,1.681) (-0.356,1.662) (-0.385,1.641) (-0.412,1.618) (-0.437,1.594) (-0.462,1.569) (-0.486,1.543) (-0.508,1.516) (-0.529,1.487) (-0.548,1.458) (-0.566,1.428) (-0.583,1.397) (-0.599,1.366) (-0.613,1.334) (-0.626,1.301) (-0.638,1.268) (-0.648,1.234) (-0.657,1.200) (-0.664,1.166) (-0.671,1.131) (-0.676,1.096) (-0.680,1.061) (-0.683,1.026) (-0.684,0.991) (-0.684,0.956) (-0.683,0.921) (-0.681,0.886) (-0.678,0.851) (-0.673,0.816) (-0.668,0.781) (-0.657,0.750)};
\addplot[lineaD, flechaen=0.382] coordinates {(-0.523,0.850) (-0.488,0.850) (-0.453,0.850) (-0.417,0.850) (-0.382,0.850) (-0.346,0.850) (-0.311,0.850) (-0.276,0.850) (-0.240,0.850) (-0.205,0.850) (-0.170,0.850) (-0.134,0.850) (-0.099,0.850) (-0.063,0.850) (-0.028,0.850) (0.007,0.850) (0.043,0.850) (0.078,0.850) (0.114,0.850) (0.149,0.850) (0.184,0.850) (0.220,0.850) (0.255,0.850) (0.291,0.850) (0.326,0.850) (0.361,0.850) (0.397,0.850) (0.432,0.850) (0.468,0.850) (0.503,0.850) (0.525,0.864) (0.519,0.898) (0.512,0.933) (0.502,0.967) (0.492,1.001) (0.479,1.034) (0.465,1.066) (0.449,1.098) (0.431,1.129) (0.411,1.158) (0.390,1.186) (0.367,1.213) (0.343,1.239) (0.317,1.263) (0.289,1.285) (0.260,1.305) (0.230,1.324) (0.198,1.340) (0.166,1.354) (0.132,1.365) (0.098,1.374) (0.063,1.380) (0.028,1.383) (-0.007,1.384) (-0.043,1.382) (-0.078,1.377) (-0.113,1.370) (-0.147,1.360) (-0.180,1.348) (-0.212,1.333) (-0.243,1.316) (-0.273,1.297) (-0.301,1.276) (-0.328,1.253) (-0.353,1.228) (-0.377,1.202) (-0.399,1.175) (-0.420,1.146) (-0.439,1.116) (-0.456,1.085) (-0.471,1.053) (-0.485,1.020) (-0.496,0.987) (-0.507,0.953) (-0.515,0.918) (-0.522,0.884) (-0.523,0.850)};
\addplot[lineaD, flechaen=0.500] coordinates {(-0.338,0.941) (-0.303,0.940) (-0.267,0.940) (-0.232,0.940) (-0.196,0.940) (-0.161,0.940) (-0.125,0.940) (-0.090,0.940) (-0.054,0.940) (-0.019,0.940) (0.017,0.940) (0.052,0.940) (0.088,0.940) (0.123,0.940) (0.159,0.940) (0.194,0.940) (0.230,0.940) (0.266,0.940) (0.301,0.940) (0.337,0.940) (0.321,0.967) (0.299,0.994) (0.274,1.020) (0.247,1.043) (0.218,1.064) (0.188,1.082) (0.156,1.098) (0.123,1.111) (0.089,1.121) (0.054,1.128) (0.019,1.131) (-0.017,1.131) (-0.052,1.128) (-0.087,1.122) (-0.121,1.112) (-0.154,1.099) (-0.186,1.083) (-0.217,1.065) (-0.245,1.044) (-0.272,1.021) (-0.297,0.996) (-0.320,0.969) (-0.338,0.941)};
\addplot[lineaD, flechaen=0.094, flechaen=0.906] coordinates {(-2.600,-1.029) (-2.570,-1.010) (-2.541,-0.991) (-2.511,-0.972) (-2.481,-0.953) (-2.451,-0.935) (-2.421,-0.916) (-2.391,-0.898) (-2.361,-0.880) (-2.331,-0.862) (-2.301,-0.844) (-2.271,-0.826) (-2.240,-0.809) (-2.210,-0.791) (-2.179,-0.774) (-2.149,-0.756) (-2.118,-0.739) (-2.087,-0.722) (-2.056,-0.705) (-2.026,-0.688) (-1.995,-0.672) (-1.964,-0.655) (-1.932,-0.639) (-1.901,-0.623) (-1.870,-0.606) (-1.839,-0.590) (-1.807,-0.575) (-1.776,-0.559) (-1.744,-0.543) (-1.713,-0.528) (-1.681,-0.513) (-1.649,-0.498) (-1.618,-0.483) (-1.586,-0.468) (-1.554,-0.453) (-1.522,-0.438) (-1.490,-0.424) (-1.458,-0.410) (-1.425,-0.396) (-1.393,-0.382) (-1.361,-0.368) (-1.328,-0.355) (-1.296,-0.341) (-1.263,-0.328) (-1.231,-0.315) (-1.198,-0.302) (-1.165,-0.289) (-1.132,-0.277) (-1.099,-0.265) (-1.066,-0.252) (-1.033,-0.240) (-1.000,-0.229) (-0.967,-0.220) (-0.931,-0.220) (-0.896,-0.220) (-0.861,-0.220) (-0.826,-0.220) (-0.791,-0.220) (-0.756,-0.220) (-0.721,-0.220) (-0.685,-0.220) (-0.650,-0.220) (-0.615,-0.220) (-0.580,-0.220) (-0.545,-0.220) (-0.510,-0.220) (-0.475,-0.220) (-0.439,-0.220) (-0.404,-0.220) (-0.369,-0.220) (-0.334,-0.220) (-0.299,-0.220) (-0.264,-0.220) (-0.228,-0.220) (-0.193,-0.220) (-0.158,-0.220) (-0.123,-0.220) (-0.088,-0.220) (-0.053,-0.220) (-0.018,-0.220) (0.018,-0.220) (0.053,-0.220) (0.088,-0.220) (0.123,-0.220) (0.158,-0.220) (0.193,-0.220) (0.228,-0.220) (0.264,-0.220) (0.299,-0.220) (0.334,-0.220) (0.369,-0.220) (0.404,-0.220) (0.439,-0.220) (0.475,-0.220) (0.510,-0.220) (0.545,-0.220) (0.580,-0.220) (0.615,-0.220) (0.650,-0.220) (0.685,-0.220) (0.721,-0.220) (0.756,-0.220) (0.791,-0.220) (0.826,-0.220) (0.861,-0.220) (0.896,-0.220) (0.931,-0.220) (0.967,-0.220) (1.000,-0.229) (1.033,-0.240) (1.066,-0.252) (1.099,-0.265) (1.132,-0.277) (1.165,-0.289) (1.198,-0.302) (1.231,-0.315) (1.263,-0.328) (1.296,-0.341) (1.328,-0.355) (1.361,-0.368) (1.393,-0.382) (1.425,-0.396) (1.458,-0.410) (1.490,-0.424) (1.522,-0.438) (1.554,-0.453) (1.586,-0.468) (1.618,-0.483) (1.649,-0.498) (1.681,-0.513) (1.713,-0.528) (1.744,-0.543) (1.776,-0.559) (1.807,-0.575) (1.839,-0.590) (1.870,-0.606) (1.901,-0.623) (1.932,-0.639) (1.964,-0.655) (1.995,-0.672) (2.026,-0.688) (2.056,-0.705) (2.087,-0.722) (2.118,-0.739) (2.149,-0.756) (2.179,-0.774) (2.210,-0.791) (2.240,-0.809) (2.271,-0.826) (2.301,-0.844) (2.331,-0.862) (2.361,-0.880) (2.391,-0.898) (2.421,-0.916) (2.451,-0.935) (2.481,-0.953) (2.511,-0.972) (2.541,-0.991) (2.570,-1.010) (2.600,-1.029)};
\addplot[lineaD, flechaen=0.044, flechaen=0.956] coordinates {(-2.270,-2.000) (-2.256,-1.968) (-2.240,-1.936) (-2.225,-1.905) (-2.209,-1.873) (-2.193,-1.842) (-2.177,-1.811) (-2.161,-1.780) (-2.144,-1.749) (-2.127,-1.718) (-2.110,-1.688) (-2.093,-1.657) (-2.075,-1.627) (-2.057,-1.597) (-2.039,-1.567) (-2.020,-1.537) (-2.002,-1.507) (-1.983,-1.477) (-1.963,-1.448) (-1.944,-1.419) (-1.924,-1.390) (-1.904,-1.361) (-1.884,-1.332) (-1.864,-1.303) (-1.843,-1.275) (-1.822,-1.247) (-1.801,-1.219) (-1.780,-1.191) (-1.758,-1.163) (-1.736,-1.135) (-1.714,-1.108) (-1.692,-1.081) (-1.669,-1.054) (-1.646,-1.027) (-1.623,-1.001) (-1.600,-0.975) (-1.577,-0.948) (-1.553,-0.923) (-1.529,-0.897) (-1.505,-0.871) (-1.480,-0.846) (-1.456,-0.821) (-1.431,-0.796) (-1.406,-0.772) (-1.380,-0.747) (-1.355,-0.723) (-1.329,-0.699) (-1.303,-0.676) (-1.277,-0.652) (-1.250,-0.629) (-1.224,-0.607) (-1.197,-0.584) (-1.170,-0.562) (-1.142,-0.540) (-1.115,-0.518) (-1.087,-0.496) (-1.059,-0.475) (-1.031,-0.454) (-1.002,-0.434) (-0.974,-0.413) (-0.945,-0.393) (-0.913,-0.380) (-0.878,-0.380) (-0.843,-0.380) (-0.808,-0.380) (-0.773,-0.380) (-0.738,-0.380) (-0.703,-0.380) (-0.667,-0.380) (-0.632,-0.380) (-0.597,-0.380) (-0.562,-0.380) (-0.527,-0.380) (-0.492,-0.380) (-0.457,-0.380) (-0.422,-0.380) (-0.386,-0.380) (-0.351,-0.380) (-0.316,-0.380) (-0.281,-0.380) (-0.246,-0.380) (-0.211,-0.380) (-0.176,-0.380) (-0.141,-0.380) (-0.105,-0.380) (-0.070,-0.380) (-0.035,-0.380) (-0.000,-0.380) (0.035,-0.380) (0.070,-0.380) (0.105,-0.380) (0.141,-0.380) (0.176,-0.380) (0.211,-0.380) (0.246,-0.380) (0.281,-0.380) (0.316,-0.380) (0.351,-0.380) (0.386,-0.380) (0.422,-0.380) (0.457,-0.380) (0.492,-0.380) (0.527,-0.380) (0.562,-0.380) (0.597,-0.380) (0.632,-0.380) (0.667,-0.380) (0.703,-0.380) (0.738,-0.380) (0.773,-0.380) (0.808,-0.380) (0.843,-0.380) (0.878,-0.380) (0.913,-0.380) (0.945,-0.393) (0.974,-0.413) (1.002,-0.434) (1.031,-0.454) (1.059,-0.475) (1.087,-0.496) (1.115,-0.518) (1.142,-0.540) (1.170,-0.562) (1.197,-0.584) (1.224,-0.607) (1.250,-0.629) (1.277,-0.652) (1.303,-0.676) (1.329,-0.699) (1.355,-0.723) (1.380,-0.747) (1.406,-0.772) (1.431,-0.796) (1.456,-0.821) (1.480,-0.846) (1.505,-0.871) (1.529,-0.897) (1.553,-0.923) (1.577,-0.948) (1.600,-0.975) (1.623,-1.001) (1.646,-1.027) (1.669,-1.054) (1.692,-1.081) (1.714,-1.108) (1.736,-1.135) (1.758,-1.163) (1.780,-1.191) (1.801,-1.219) (1.822,-1.247) (1.843,-1.275) (1.864,-1.303) (1.884,-1.332) (1.904,-1.361) (1.924,-1.390) (1.944,-1.419) (1.963,-1.448) (1.983,-1.477) (2.002,-1.507) (2.020,-1.537) (2.039,-1.567) (2.057,-1.597) (2.075,-1.627) (2.093,-1.657) (2.110,-1.688) (2.127,-1.718) (2.144,-1.749) (2.161,-1.780) (2.177,-1.811) (2.193,-1.842) (2.209,-1.873) (2.225,-1.905) (2.240,-1.936) (2.256,-1.968) (2.270,-2.000)};
\addplot[lineaD, flechaen=0.399] coordinates {(-1.423,-2.000) (-1.423,-1.965) (-1.422,-1.930) (-1.420,-1.895) (-1.418,-1.860) (-1.415,-1.825) (-1.412,-1.790) (-1.408,-1.755) (-1.404,-1.720) (-1.399,-1.685) (-1.393,-1.651) (-1.387,-1.616) (-1.380,-1.582) (-1.373,-1.547) (-1.366,-1.513) (-1.357,-1.479) (-1.349,-1.445) (-1.339,-1.411) (-1.330,-1.378) (-1.319,-1.344) (-1.309,-1.311) (-1.297,-1.277) (-1.285,-1.244) (-1.273,-1.212) (-1.260,-1.179) (-1.247,-1.146) (-1.233,-1.114) (-1.219,-1.082) (-1.204,-1.050) (-1.189,-1.019) (-1.173,-0.987) (-1.157,-0.956) (-1.140,-0.925) (-1.123,-0.895) (-1.105,-0.864) (-1.087,-0.834) (-1.069,-0.805) (-1.050,-0.775) (-1.030,-0.746) (-1.010,-0.717) (-0.990,-0.689) (-0.969,-0.660) (-0.948,-0.632) (-0.926,-0.605) (-0.904,-0.578) (-0.881,-0.551) (-0.858,-0.524) (-0.825,-0.520) (-0.789,-0.520) (-0.754,-0.520) (-0.719,-0.520) (-0.684,-0.520) (-0.649,-0.520) (-0.614,-0.520) (-0.579,-0.520) (-0.544,-0.520) (-0.509,-0.520) (-0.474,-0.520) (-0.439,-0.520) (-0.404,-0.520) (-0.368,-0.520) (-0.333,-0.520) (-0.298,-0.520) (-0.263,-0.520) (-0.228,-0.520) (-0.193,-0.520) (-0.158,-0.520) (-0.123,-0.520) (-0.088,-0.520) (-0.053,-0.520) (-0.018,-0.520) (0.018,-0.520) (0.053,-0.520) (0.088,-0.520) (0.123,-0.520) (0.158,-0.520) (0.193,-0.520) (0.228,-0.520) (0.263,-0.520) (0.298,-0.520) (0.333,-0.520) (0.368,-0.520) (0.404,-0.520) (0.439,-0.520) (0.474,-0.520) (0.509,-0.520) (0.544,-0.520) (0.579,-0.520) (0.614,-0.520) (0.649,-0.520) (0.684,-0.520) (0.719,-0.520) (0.754,-0.520) (0.789,-0.520) (0.825,-0.520) (0.858,-0.524) (0.881,-0.551) (0.904,-0.578) (0.926,-0.605) (0.948,-0.632) (0.969,-0.660) (0.990,-0.689) (1.010,-0.717) (1.030,-0.746) (1.050,-0.775) (1.069,-0.805) (1.087,-0.834) (1.105,-0.864) (1.123,-0.895) (1.140,-0.925) (1.157,-0.956) (1.173,-0.987) (1.189,-1.019) (1.204,-1.050) (1.219,-1.082) (1.233,-1.114) (1.247,-1.146) (1.260,-1.179) (1.273,-1.212) (1.285,-1.244) (1.297,-1.277) (1.309,-1.311) (1.319,-1.344) (1.330,-1.378) (1.339,-1.411) (1.349,-1.445) (1.357,-1.479) (1.366,-1.513) (1.373,-1.547) (1.380,-1.582) (1.387,-1.616) (1.393,-1.651) (1.399,-1.685) (1.404,-1.720) (1.408,-1.755) (1.412,-1.790) (1.415,-1.825) (1.418,-1.860) (1.420,-1.895) (1.422,-1.930) (1.423,-1.965) (1.423,-2.000)};
\addplot[lineaD, flechaen=0.614] coordinates {(-0.754,-2.000) (-0.772,-1.970) (-0.790,-1.940) (-0.806,-1.909) (-0.821,-1.877) (-0.836,-1.845) (-0.849,-1.813) (-0.862,-1.780) (-0.873,-1.747) (-0.884,-1.714) (-0.894,-1.680) (-0.902,-1.646) (-0.910,-1.612) (-0.917,-1.578) (-0.923,-1.543) (-0.928,-1.508) (-0.932,-1.474) (-0.935,-1.439) (-0.938,-1.404) (-0.939,-1.369) (-0.940,-1.334) (-0.939,-1.299) (-0.938,-1.264) (-0.936,-1.229) (-0.933,-1.194) (-0.930,-1.159) (-0.925,-1.124) (-0.920,-1.090) (-0.913,-1.055) (-0.906,-1.021) (-0.899,-0.987) (-0.890,-0.953) (-0.881,-0.919) (-0.870,-0.885) (-0.859,-0.852) (-0.848,-0.819) (-0.835,-0.786) (-0.822,-0.754) (-0.808,-0.722) (-0.793,-0.690) (-0.778,-0.658) (-0.753,-0.640) (-0.718,-0.640) (-0.683,-0.640) (-0.648,-0.640) (-0.613,-0.640) (-0.578,-0.640) (-0.543,-0.640) (-0.508,-0.640) (-0.473,-0.640) (-0.438,-0.640) (-0.403,-0.640) (-0.368,-0.640) (-0.333,-0.640) (-0.298,-0.640) (-0.263,-0.640) (-0.228,-0.640) (-0.193,-0.640) (-0.158,-0.640) (-0.123,-0.640) (-0.088,-0.640) (-0.053,-0.640) (-0.018,-0.640) (0.018,-0.640) (0.053,-0.640) (0.088,-0.640) (0.123,-0.640) (0.158,-0.640) (0.193,-0.640) (0.228,-0.640) (0.263,-0.640) (0.298,-0.640) (0.333,-0.640) (0.368,-0.640) (0.403,-0.640) (0.438,-0.640) (0.473,-0.640) (0.508,-0.640) (0.543,-0.640) (0.578,-0.640) (0.613,-0.640) (0.648,-0.640) (0.683,-0.640) (0.718,-0.640) (0.753,-0.640) (0.778,-0.658) (0.793,-0.690) (0.808,-0.722) (0.822,-0.754) (0.835,-0.786) (0.848,-0.819) (0.859,-0.852) (0.870,-0.885) (0.881,-0.919) (0.890,-0.953) (0.899,-0.987) (0.906,-1.021) (0.913,-1.055) (0.920,-1.090) (0.925,-1.124) (0.930,-1.159) (0.933,-1.194) (0.936,-1.229) (0.938,-1.264) (0.939,-1.299) (0.940,-1.334) (0.939,-1.369) (0.938,-1.404) (0.935,-1.439) (0.932,-1.474) (0.928,-1.508) (0.923,-1.543) (0.917,-1.578) (0.910,-1.612) (0.902,-1.646) (0.894,-1.680) (0.884,-1.714) (0.873,-1.747) (0.862,-1.780) (0.849,-1.813) (0.836,-1.845) (0.821,-1.877) (0.806,-1.909) (0.790,-1.940) (0.772,-1.970) (0.754,-2.000)};
\addplot[lineaD, flechaen=0.500] coordinates {(-0.657,-0.750) (-0.621,-0.750) (-0.586,-0.750) (-0.551,-0.750) (-0.516,-0.750) (-0.481,-0.750) (-0.446,-0.750) (-0.410,-0.750) (-0.375,-0.750) (-0.340,-0.750) (-0.305,-0.750) (-0.270,-0.750) (-0.234,-0.750) (-0.199,-0.750) (-0.164,-0.750) (-0.129,-0.750) (-0.094,-0.750) (-0.059,-0.750) (-0.023,-0.750) (0.012,-0.750) (0.047,-0.750) (0.082,-0.750) (0.117,-0.750) (0.152,-0.750) (0.188,-0.750) (0.223,-0.750) (0.258,-0.750) (0.293,-0.750) (0.328,-0.750) (0.363,-0.750) (0.399,-0.750) (0.434,-0.750) (0.469,-0.750) (0.504,-0.750) (0.539,-0.750) (0.574,-0.750) (0.610,-0.750) (0.645,-0.750) (0.665,-0.770) (0.671,-0.804) (0.676,-0.839) (0.680,-0.874) (0.683,-0.909) (0.684,-0.944) (0.684,-0.979) (0.683,-1.015) (0.681,-1.050) (0.677,-1.085) (0.673,-1.119) (0.667,-1.154) (0.660,-1.189) (0.651,-1.223) (0.641,-1.256) (0.630,-1.290) (0.617,-1.323) (0.604,-1.355) (0.589,-1.387) (0.572,-1.418) (0.554,-1.448) (0.535,-1.478) (0.515,-1.506) (0.493,-1.534) (0.470,-1.561) (0.446,-1.586) (0.420,-1.610) (0.394,-1.633) (0.366,-1.655) (0.337,-1.675) (0.307,-1.693) (0.276,-1.710) (0.245,-1.725) (0.212,-1.739) (0.179,-1.750) (0.145,-1.760) (0.111,-1.767) (0.076,-1.773) (0.041,-1.776) (0.006,-1.778) (-0.029,-1.777) (-0.064,-1.774) (-0.099,-1.769) (-0.134,-1.763) (-0.168,-1.754) (-0.201,-1.743) (-0.234,-1.730) (-0.266,-1.716) (-0.297,-1.699) (-0.327,-1.681) (-0.356,-1.662) (-0.385,-1.641) (-0.412,-1.618) (-0.437,-1.594) (-0.462,-1.569) (-0.486,-1.543) (-0.508,-1.516) (-0.529,-1.487) (-0.548,-1.458) (-0.566,-1.428) (-0.583,-1.397) (-0.599,-1.366) (-0.613,-1.334) (-0.626,-1.301) (-0.638,-1.268) (-0.648,-1.234) (-0.657,-1.200) (-0.664,-1.166) (-0.671,-1.131) (-0.676,-1.096) (-0.680,-1.061) (-0.683,-1.026) (-0.684,-0.991) (-0.684,-0.956) (-0.683,-0.921) (-0.681,-0.886) (-0.678,-0.851) (-0.673,-0.816) (-0.668,-0.781) (-0.657,-0.750)};
\addplot[lineaD, flechaen=0.382] coordinates {(-0.523,-0.850) (-0.488,-0.850) (-0.453,-0.850) (-0.417,-0.850) (-0.382,-0.850) (-0.346,-0.850) (-0.311,-0.850) (-0.276,-0.850) (-0.240,-0.850) (-0.205,-0.850) (-0.170,-0.850) (-0.134,-0.850) (-0.099,-0.850) (-0.063,-0.850) (-0.028,-0.850) (0.007,-0.850) (0.043,-0.850) (0.078,-0.850) (0.114,-0.850) (0.149,-0.850) (0.184,-0.850) (0.220,-0.850) (0.255,-0.850) (0.291,-0.850) (0.326,-0.850) (0.361,-0.850) (0.397,-0.850) (0.432,-0.850) (0.468,-0.850) (0.503,-0.850) (0.525,-0.864) (0.519,-0.898) (0.512,-0.933) (0.502,-0.967) (0.492,-1.001) (0.479,-1.034) (0.465,-1.066) (0.449,-1.098) (0.431,-1.129) (0.411,-1.158) (0.390,-1.186) (0.367,-1.213) (0.343,-1.239) (0.317,-1.263) (0.289,-1.285) (0.260,-1.305) (0.230,-1.324) (0.198,-1.340) (0.166,-1.354) (0.132,-1.365) (0.098,-1.374) (0.063,-1.380) (0.028,-1.383) (-0.007,-1.384) (-0.043,-1.382) (-0.078,-1.377) (-0.113,-1.370) (-0.147,-1.360) (-0.180,-1.348) (-0.212,-1.333) (-0.243,-1.316) (-0.273,-1.297) (-0.301,-1.276) (-0.328,-1.253) (-0.353,-1.228) (-0.377,-1.202) (-0.399,-1.175) (-0.420,-1.146) (-0.439,-1.116) (-0.456,-1.085) (-0.471,-1.053) (-0.485,-1.020) (-0.496,-0.987) (-0.507,-0.953) (-0.515,-0.918) (-0.522,-0.884) (-0.523,-0.850)};
\addplot[lineaD, flechaen=0.500] coordinates {(-0.338,-0.941) (-0.303,-0.940) (-0.267,-0.940) (-0.232,-0.940) (-0.196,-0.940) (-0.161,-0.940) (-0.125,-0.940) (-0.090,-0.940) (-0.054,-0.940) (-0.019,-0.940) (0.017,-0.940) (0.052,-0.940) (0.088,-0.940) (0.123,-0.940) (0.159,-0.940) (0.194,-0.940) (0.230,-0.940) (0.266,-0.940) (0.301,-0.940) (0.337,-0.940) (0.321,-0.967) (0.299,-0.994) (0.274,-1.020) (0.247,-1.043) (0.218,-1.064) (0.188,-1.082) (0.156,-1.098) (0.123,-1.111) (0.089,-1.121) (0.054,-1.128) (0.019,-1.131) (-0.017,-1.131) (-0.052,-1.128) (-0.087,-1.122) (-0.121,-1.112) (-0.154,-1.099) (-0.186,-1.083) (-0.217,-1.065) (-0.245,-1.044) (-0.272,-1.021) (-0.297,-0.996) (-0.320,-0.969) (-0.338,-0.941)};
\addplot[lineaD, flechaen=0.087, flechaen=0.913] coordinates {(-2.600,0.000) (2.600,0.000)};

  \draw[texto, line width=0.6pt] (axis cs:0,0) circle[radius=1];
  \draw[vecP] (axis cs:-0.45,-0.55) -- (axis cs:0.45,-0.55);
  \node[etiqm, font=\footnotesize] at (axis cs:0,-0.78) {$\vP$};
\end{axis}
\end{tikzpicture}
\caption{Líneas de $\vD$}
\end{subfigure}
\caption{Esfera uniformemente polarizada (\cref{prob:esfera-polarizada}); trazos: equipotenciales. (a) Fuera,
el campo es exactamente el de un dipolo $\vp=\tfrac43\pi a^3\vP$ en el centro; dentro es uniforme y
opuesto a la polarización, $\vE=-\vP/3\eo$ (campo despolarizante): las líneas nacen en la carga ligada
positiva $\sigmaP=P\cos\theta$ y mueren en la negativa. (b) Sin carga libre, $\vD$ es solenoidal y sus
líneas son cerradas; dentro, $\vD=\eo\vE+\vP=\tfrac23\vP$ tiene el sentido de $\vP$.}
\label{fig:esfera-polarizada}
\end{figure}


% -----------------------------------------------------------------------------
\subsection{Problema 4.2. Carga libre en un dieléctrico}\label{sec:p42}

\begin{problema}[Esfera cargada en un dieléctrico]\label{prob:carga-en-dielectrico}
Una esfera de radio $a$ y carga $Q$ está situada en un dieléctrico lineal, isótropo y homogéneo de permitividad $\varepsilon$. Calcula
$\vE$, $\vD$ y $\vP$ en el medio dieléctrico.
\end{problema}

\begin{resolucion}
\idea Un material que se polariza en respuesta a un campo tiene que tener cerca alguna carga libre que lo cree: aquí, la de
la esfera, que supondremos conductora. Conocemos la carga libre, no la ligada, así que usamos Gauss para $\vD$ (método M3).

\paso{Paso 1: $\vD$ por Gauss} Por la simetría esférica (dieléctrico \LHI{} que llena todo el espacio exterior), $\vD=D(r)\,\uv r$.
Con una esfera gaussiana de radio $r>a$,
\[
  \oiint_S\vD\cdot\dS=D(r)\,4\pi r^2=Q
  \quad\Longrightarrow\quad
  \vD=\frac{Q}{4\pi r^2}\,\uv r .
\]
\paso{Paso 2: campo y polarización} En el medio \LHI{},
\[
  \recuadro{\vE=\frac{\vD}{\varepsilon}=\frac{Q}{4\pi\varepsilon r^2}\,\uv r,\qquad
  \vP=\vD-\eo\vE=(\varepsilon-\eo)\,\frac{Q}{4\pi\varepsilon r^2}\,\uv r .}
\]
\paso{Paso 3: cargas ligadas} Como $\divg(\uv r/r^2)=0$ para $r>0$, $\rhoP=0$, en acuerdo con la \cref{prop:rhoP-rhof}. En la superficie
del dieléctrico en contacto con la esfera, la normal exterior al dieléctrico es $\nn=-\uv r$:
\[
  \sigmaP(a)=\vP\cdot\nn=-\frac{Q}{4\pi a^2}\,\frac{\varepsilon-\eo}{\varepsilon}.
\]
\paso{Paso 4: comprobación con Gauss para $\vE$} La carga total encerrada por una esfera de radio $r>a$ es
\[
  Q+4\pi a^2\sigmaP=Q-Q\,\frac{\varepsilon-\eo}{\varepsilon}=\frac{\eo}{\varepsilon}\,Q
  \quad\Longrightarrow\quad
  E\,4\pi r^2=\frac{Q}{\varepsilon},\quad E=\frac{Q}{4\pi\varepsilon r^2}. \qedhere
\]
\end{resolucion}

\begin{observacion}[Apantallamiento]\label{obs:p42}
La carga ligada negativa que se pega a la esfera apantalla parcialmente su carga, y el campo es $\er$ veces menor que en el vacío
(\cref{fig:p42}). En el límite $\varepsilon\to\infty$ el apantallamiento es total y $\vE\to\vcero$: un dieléctrico de permitividad infinita se comporta como un
conductor. El potencial de la esfera es $Q/4\pi\varepsilon a$, y su capacidad, $C=4\pi\varepsilon a$, $\er$ veces la del vacío. Si el dieléctrico fuera
finito, en su superficie exterior aparecería la carga ligada opuesta, $+Q(\varepsilon-\eo)/\varepsilon$, que garantiza $\QP=0$; el resultado anterior solo es
exacto si esa superficie es una esfera concéntrica o está en el infinito.
\end{observacion}

% Figura: problema 4.2. Esfera conductora cargada en un dieléctrico LHI infinito.
\begin{figure}[t]
\centering
\begin{subfigure}[b]{0.42\textwidth}\centering
\begin{tikzpicture}[scale=1.0]
  \fill[dielec] plot[smooth cycle, tension=0.8] coordinates
     {(-2.3,0.2) (-1.6,1.6) (0.2,1.9) (1.9,1.5) (2.3,-0.1) (1.6,-1.6) (-0.3,-1.85) (-1.9,-1.3)};
  \node[etiq] at (1.55,1.25) {$\varepsilon$};
  \draw[gauss] (0,0) circle (1.25);
  \node[etiq, text=acentoB] at (-50:1.55) {$S$};
  \shade[ball color=conductorluz] (0,0) circle (0.55);
  \draw[texto, line width=0.5pt] (0,0) circle (0.55);
  \foreach \t in {0,40,...,320} \node[cneg, font=\tiny\bfseries] at (\t:0.7) {$-$};
  \node[font=\small, text=black] at (0,0.05) {$Q$};
  \draw[texto, line width=0.4pt, -{Stealth[length=1.6mm]}] (0,0) -- (-140:0.55);
  \draw[acentoB, line width=0.4pt, -{Stealth[length=1.6mm]}] (0,0) -- (115:1.25)
    node[pos=0.72, left, etiq] {$r$};
  \node[etiqt] at (-118:0.92) {$a$};
\end{tikzpicture}
\caption{Geometría}
\end{subfigure}\hfill
\begin{subfigure}[b]{0.54\textwidth}\centering
\begin{tikzpicture}
\begin{axis}[width=7.6cm, height=4.8cm, axis lines=left, xmin=0, xmax=4, ymin=0, ymax=1.1,
  xlabel={$r/a$}, ylabel={unidades de $Q/4\pi a^{2}$}, xtick={0,1,2,3,4}, ytick={0,0.3333,1},
  yticklabels={$0$,$\frac1{\er}$,$1$}, legend style={at={(0.98,0.98)}, anchor=north east}]
  \addplot[curvaD, dashed, domain=1:4, samples=80] {1/x^2}; \addlegendentry{$D$ \;($=\eo E$ en vacío)}
  \addplot[curvaE, domain=1:4, samples=80] {1/(3*x^2)}; \addlegendentry{$\eo E$ \;($\er=3$)}
  \addplot[curvaP, domain=1:4, samples=80] {2/(3*x^2)}; \addlegendentry{$P$}
  \fill[conductor, opacity=0.35] (axis cs:0,0) rectangle (axis cs:1,1.1);
  \node[etiqt] at (axis cs:0.5,0.55) {metal};
\end{axis}
\end{tikzpicture}
\caption{$D$, $\eo E$ y $P$ frente a $r$}
\end{subfigure}
\caption{\Cref{prob:carga-en-dielectrico}. (a) Conductor de radio $a$ y carga libre $Q$ sumergido en un dieléctrico
\LHI{} de permitividad $\varepsilon$; la superficie gaussiana $S$ es una esfera concéntrica. Junto al metal aparece
una carga ligada negativa que apantalla parcialmente a $Q$. (b) $D$ no depende del medio; el campo se
reduce en el factor $\er$ y la diferencia, $D-\eo E=P$, es la polarización.}
\label{fig:p42}
\end{figure}


% -----------------------------------------------------------------------------
\subsection{Problema 4.3. Condensador plano con dieléctrico}\label{sec:p43}

\begin{problema}[Condensador plano con dieléctrico]\label{prob:condensador-dielectrico}
Analiza los campos en el interior de un condensador plano cuando las densidades de carga de las placas son $\pm\sigma$ y entre ellas
se introduce un dieléctrico de permitividad $\varepsilon$. Las placas tienen área $A$ y están separadas una distancia $d$. Calcula la capacidad.
\end{problema}

\begin{resolucion}
\idea Despreciando los efectos de borde el problema tiene simetría plana; conocemos la carga libre de las placas y usamos Gauss
para $\vD$ (método M3). Tomamos $\uv z$ de la placa positiva a la negativa (\cref{fig:p43}).

\paso{Paso 0: sin dieléctrico} El campo de dos planos con $\pm\sigma$ es $\sigma/\eo$ entre ellos, y $C_0=\dfrac{Q}{\phi_+-\phi_-}=\dfrac{\sigma A}{(\sigma/\eo)\,d}=\eo\dfrac{A}{d}$.

\paso{Paso 1: $\vD$ por Gauss} Con la caja $S'$ de la \cref{fig:p43}, con una cara dentro de la placa inferior (donde $\vD=\vcero$) y la otra
en el dieléctrico, de área $S'$: $D\,S'=\sigma S'$. Con la caja $S$ en la placa superior, $-D\,S=-\sigma S$. En ambos casos
\[
  \vD=\sigma\,\uv z .
\]
\paso{Paso 2: campo, polarización y cargas ligadas}
\[
  \vE=\frac{\vD}{\varepsilon}=\frac{\sigma}{\varepsilon}\,\uv z,\qquad
  \vP=\vD-\eo\vE=\frac{\varepsilon-\eo}{\varepsilon}\,\sigma\,\uv z,\qquad
  \sigmaP=\vP\cdot\nn=\pm\frac{\varepsilon-\eo}{\varepsilon}\,\sigma ,
\]
con el signo $+$ en la cara adyacente a la placa negativa ($\nn=\uv z$) y $-$ junto a la positiva.

\paso{Paso 3: capacidad} $\phi_+-\phi_-=E\,d=\sigma d/\varepsilon$ y
\[
  \recuadro{C=\frac{\sigma A}{\sigma d/\varepsilon}=\varepsilon\,\frac{A}{d},\qquad \frac{C}{C_0}=\frac{\varepsilon}{\eo}=\er .}
\]
\paso{Paso 4: comprobación con Gauss para $\vE$} En la caja $S'$, la carga total encerrada es $(\sigma-\sigmaP)S'$:
$E\,S'=\dfrac{1}{\eo}\Big(\sigma-\sigma\dfrac{\varepsilon-\eo}{\varepsilon}\Big)S'=\dfrac{\sigma}{\varepsilon}\,S'$. Es una comprobación, no otra forma de resolverlo: para
usarla hay que conocer antes $\sigmaP$.
\end{resolucion}

\begin{observacion}[Carga fija o potencial fijo]\label{obs:p43}
Si el condensador está aislado (carga fija), introducir el dieléctrico reduce la diferencia de potencial en el factor $\er$; si está
conectado a una batería (potencial fijo), la batería suministra carga hasta multiplicarla por $\er$. En ambos casos $C=\er C_0$: la
capacidad solo depende de la geometría y del material.
\end{observacion}

% Figura: problema 4.3. Condensador plano con dieléctrico: cargas libres, ligadas,
% campos y superficies gaussianas.
\begin{figure}[t]
\centering
\begin{tikzpicture}[scale=1.0]
  % dieléctrico y placas
  \fill[dielec] (0,0) rectangle (7.0,2.6);
  \fill[conductor] (0,2.6) rectangle (7.0,2.82);
  \fill[conductor] (0,-0.22) rectangle (7.0,0);
  \draw[texto, line width=0.5pt] (0,2.6) rectangle (7.0,2.82) (0,-0.22) rectangle (7.0,0);
  \foreach \x in {0.3,0.85,...,6.9} {
    \node[cneg, font=\scriptsize\bfseries] at (\x,2.71) {$-$};
    \node[cpos, font=\scriptsize\bfseries] at (\x,-0.11) {$+$};
    \node[cpos, font=\tiny\bfseries] at (\x,2.47) {$+$};
    \node[cneg, font=\tiny\bfseries] at (\x,0.13) {$-$};
  }
  \node[etiq, anchor=west] at (7.1,2.71) {$-\sigma$};
  \node[etiq, anchor=west] at (7.1,-0.11) {$+\sigma$};
  \node[etiq, anchor=west, text=acentoC] at (7.1,2.42) {$+\sigmaP$};
  \node[etiq, anchor=west, text=acentoC] at (7.1,0.18) {$-\sigmaP$};
  % campos
  \draw[vecD] (0.65,0.35) -- (0.65,2.2) node[pos=0.5, left, etiq] {$\vD$};
  \draw[vecE] (1.35,0.84) -- (1.35,1.76) node[pos=0.5, left, etiq] {$\vE$};
  \draw[vecP] (2.0,0.84) -- (2.0,1.76) node[pos=0.5, left, etiq] {$\vP$};
  \node[etiq] at (2.85,1.3) {$\varepsilon$};
  % superficie gaussiana S (encierra la placa superior)
  \draw[gauss] (3.5,1.5) rectangle (4.7,3.2);
  \node[etiq, text=acentoB, anchor=west] at (4.75,3.05) {$S$};
  \draw[vecn] (4.1,1.5) -- (4.1,0.95) node[right, etiq] {$\nn$};
  % superficie gaussiana S' (solo la cara del dieléctrico y la placa inferior)
  \draw[gauss] (5.3,-0.5) rectangle (6.5,1.05);
  \node[etiq, text=acentoB, anchor=west] at (6.55,0.95) {$S'$};
  \draw[vecn] (5.9,1.05) -- (5.9,1.6) node[right, etiq] {$\nn'$};
  % cotas
  \draw[cota] (-0.4,0) -- (-0.4,2.6) node[midway, left, etiq] {$d$};
  \node[etiq, anchor=south] at (1.6,2.86) {$A$};
\end{tikzpicture}
\caption{\Cref{prob:condensador-dielectrico}. Condensador plano con cargas libres $\pm\sigma$ en las placas y un dieléctrico de
permitividad $\varepsilon$. Junto a cada placa aparece una carga ligada de signo opuesto, $\pm\sigmaP$ con
$\sigmaP=\sigma(\varepsilon-\eo)/\varepsilon<\sigma$. Las tres flechas indican el sentido común de $\vD$, $\vE$ y $\vP$ (de la placa positiva a
la negativa); sus longitudes son proporcionales a $D=\sigma$, $\eo E=\sigma/\er$ y $P=\sigmaP$ (dibujadas con $\er=2$). Las superficies
$S$ y $S'$ sirven para aplicar la ley de Gauss para $\vD$ (solo cuenta la carga libre).}
\label{fig:p43}
\end{figure}


% -----------------------------------------------------------------------------
\subsection{Problema 4.4. Electrete cilíndrico}\label{sec:p44}

\begin{problema}[Cilindro uniformemente polarizado]\label{prob:cilindro-polarizado}
Sea un cilindro de radio $a$ y longitud $2\ell$ de material dieléctrico uniformemente polarizado en la dirección de su eje. Calcula los
campos $\vE$ y $\vD$ en los puntos del eje y obtén expresiones aproximadas si consideramos un cilindro muy largo ($a\ll\ell$) o muy corto
($a\gg\ell$). Representa las líneas de $\vE$, $\vD$ y $\vP$.
\end{problema}

\begin{resolucion}
\idea $\vD\neq\vcero$ aunque no haya carga libre, porque $\rot\vP\neq\vcero$ en la pared lateral (\cref{ej:stokes}), y sin simetría no podemos usar Gauss.
Sustituimos el cilindro por sus cargas ligadas, que son dos discos, y superponemos sus campos (método M1).

\paso{Paso 1: cargas ligadas} Con $\vP=P\,\uv z$ uniforme, $\rhoP=0$; en las caras $z=\pm\ell$, $\nn=\pm\uv z$ y $\sigmaP=\pm P$; en la pared lateral,
$\vP\perp\nn$ y $\sigmaP=0$ (\cref{fig:p44}\,(a)).

\paso{Paso 2: campo de un disco} El campo en el eje de un disco de radio $a$ y densidad $\sigma$, a distancia $z$ de su centro, es~\cite{T1}
\[
  \vE_{\text{disco}}(z)=\frac{\sigma}{2\eo}\left(\frac{z}{\abs z}-\frac{z}{\sqrt{a^2+z^2}}\right)\uv z .
\]
\paso{Paso 3: superposición} Sumamos el disco $+P$ en $z=\ell$ y el disco $-P$ en $z=-\ell$, y usamos los ángulos de la \cref{fig:p44}\,(b),(c),
\[
  \cos\theta_1=\frac{z-\ell}{\sqrt{a^2+(z-\ell)^2}},\qquad \cos\theta_2=\frac{z+\ell}{\sqrt{a^2+(z+\ell)^2}} :
\]
\[
  E_z(z)=\frac{P}{2\eo}\Big[\sgn(z-\ell)-\sgn(z+\ell)+\cos\theta_2-\cos\theta_1\Big].
\]
Fuera ($\abs z>\ell$) los signos se cancelan; dentro ($\abs z<\ell$) su diferencia vale $-2$. Así,
\begin{equation}\label{eq:p44-E}
  \recuadro{\begin{aligned}
    \vE&=\frac{P}{2\eo}\big(\cos\theta_2-\cos\theta_1\big)\,\uv z, &&\abs z>\ell,\\
    \vE&=-\frac{P}{2\eo}\big(2+\cos\theta_1-\cos\theta_2\big)\,\uv z, &&\abs z<\ell.
  \end{aligned}}
\end{equation}
\paso{Paso 4: desplazamiento} Fuera, $\vD=\eo\vE$; dentro, $\vD=\eo\vE+P\,\uv z$. En ambos casos
\begin{equation}\label{eq:p44-D}
  \recuadro{\vD=\frac{P}{2}\big(\cos\theta_2-\cos\theta_1\big)\,\uv z\qquad\text{en todo el eje}.}
\end{equation}
\end{resolucion}

\begin{observacion}[Comprobaciones]\label{obs:p44-comprobaciones}
$D_z$ es continua en las caras, como debe ser sin carga libre (\cref{teo:frontera}), mientras que $E_z$ salta en $P/\eo=\sigmaP/\eo$ (\cref{cor:otras}). En el
centro, $z=0$, $\cos\theta_2=-\cos\theta_1=\ell/\sqrt{a^2+\ell^2}$ y $E_z(0)=-\dfrac{P}{\eo}\Big(1-\dfrac{\ell}{\sqrt{a^2+\ell^2}}\Big)$: es la fórmula que usamos para las
cavidades (\cref{sec:cavidades}).
\end{observacion}

\paso{Límites} \emph{Cilindro largo} ($a\ll\ell$). Para puntos interiores lejos de los extremos, $\theta_2\simeq0$ y $\theta_1\simeq\pi$, de modo que
$\cos\theta_2\simeq1$ y $\cos\theta_1\simeq-1$:
\[
  \vEint\simeq\frac{\vP}{2\eo}\big[-2+1-(-1)\big]\simeq\vcero,\qquad \vD_{\mathrm{int}}\simeq\frac{\vP}{2}\big[1-(-1)\big]=\vP ;
\]
fuera, cerca de una cara, solo sobrevive el disco próximo, $\vEext\simeq\frac{\vP}{2\eo}(1-\cos\theta_1)$ y $\vD_{\mathrm{ext}}\simeq\frac{\vP}{2}(1-\cos\theta_1)$, que decaen
rápidamente: el campo se concentra en los extremos (\cref{fig:p44-perfiles}, tercera columna).

\emph{Cilindro corto} ($a\gg\ell$). Para $\abs z\lesssim\ell\ll a$, ambos ángulos son próximos a $\pi/2$ y $\cos\theta_{1,2}\simeq0$:
\[
  \vEext\simeq\vcero,\quad \vD_{\mathrm{ext}}\simeq\vcero,\qquad \vEint\simeq-\frac{\vP}{\eo},\quad \vD_{\mathrm{int}}\simeq\vcero :
\]
el electrete se comporta como un condensador plano con densidades $\pm P$ en las caras (\cref{fig:p44-perfiles}, segunda columna).

\begin{observacion}[Líneas de campo]\label{obs:p44-lineas}
La \cref{fig:cilindro-lineas} muestra las tres familias. $\vP$ es uniforme y solo existe dentro. $\vE$ nace en la cara positiva y muere en la
negativa: fuera es el campo de dos discos; dentro se opone a $\vP$. $\vD$ forma líneas cerradas que atraviesan el cilindro en el sentido de $\vP$
y se cierran por fuera como las de un imán de barra: la analogía es exacta, $\vD$ es formalmente el campo de un solenoide con «corriente»
superficial $\vP\times\nn$.
\end{observacion}

% Figura: problema 4.4. Electrete cilíndrico: circuito de Stokes, discos
% equivalentes y ángulos theta_1, theta_2 para puntos del eje fuera y dentro.
\begin{figure}[t]
\centering
\begin{subfigure}[b]{\textwidth}\centering
\begin{tikzpicture}[scale=0.95]
  \fill[dielec] (0,-0.8) -- (3.6,-0.8) arc (-90:90:0.28 and 0.8) -- (0,0.8) arc (90:270:0.28 and 0.8);
  \draw[borde] (0,-0.8) -- (3.6,-0.8) (0,0.8) -- (3.6,0.8);
  \draw[borde] (0,0) ellipse (0.28 and 0.8);
  \fill[dielecB] (3.6,0) ellipse (0.28 and 0.8);
  \draw[borde] (3.6,0) ellipse (0.28 and 0.8);
  \draw[vecP] (0.9,0.1) -- (2.6,0.1) node[above, midway, etiq] {$\vP$};
  \draw[texto, line width=0.4pt, -{Stealth[length=1.6mm]}] (3.6,0) -- (3.75,0.72);
  \node[etiq] at (4.12,0.42) {$a$};
  \draw[cota] (0,1.15) -- (3.6,1.15) node[midway, above, etiq] {$2\ell$};
  % circuito de Stokes a caballo de la pared lateral
  \draw[cargapos, line width=0.7pt] (1.3,-1.05) rectangle (2.3,-0.55);
  \draw[cargapos, line width=0.7pt, flechamedio] (1.3,-0.55) -- (2.3,-0.55);
  \draw[cargapos, line width=0.7pt, flechamedio] (2.3,-1.05) -- (1.3,-1.05);
  \node[etiq, anchor=north, text=cargapos] at (1.8,-1.12) {$\oint\vP\cdot\dl=P\,\Delta\ell\neq0$};
  % equivalente
  \node[etiq] at (5.1,0) {$\equiv$};
  \draw[borde, fill=dielecB] (6.3,0) ellipse (0.28 and 0.8);
  \draw[borde, fill=dielecB] (9.9,0) ellipse (0.28 and 0.8);
  \draw[texto!70!fondo, line width=0.5pt] (6.3,0) -- (9.9,0);
  \node[cneg, font=\small] at (6.3,1.1) {$-\sigmaP$};
  \node[cpos, font=\small] at (9.9,1.1) {$+\sigmaP$};
  \draw[cota] (6.3,-1.1) -- (9.9,-1.1) node[midway, below, etiq] {$2\ell$};
\end{tikzpicture}
\caption{Circuito de Stokes y discos equivalentes ($\rhoP=0$, $\sigmaP=\pm P$)}
\end{subfigure}\\[0.9em]
\begin{subfigure}[b]{0.49\textwidth}\centering
\begin{tikzpicture}[scale=0.9]
  \draw[texto!70!fondo, line width=0.5pt, -{Stealth[length=2mm]}] (-2.6,0) -- (4.5,0) node[below, etiq] {$z$};
  \fill[dielecB, opacity=0.8] (-1.6,0) ellipse (0.3 and 1.3);
  \fill[dielecB, opacity=0.8] (1.6,0) ellipse (0.3 and 1.3);
  \draw[borde] (-1.6,0) ellipse (0.3 and 1.3) (1.6,0) ellipse (0.3 and 1.3);
  \node[cneg, font=\small] at (-1.6,1.6) {$-P$};
  \node[cpos, font=\small] at (1.6,1.6) {$+P$};
  \draw[texto, line width=0.4pt] (-1.6,0) -- (-1.6,1.3);
  \node[etiq, anchor=east] at (-1.95,0.65) {$a$};
  \coordinate (Z) at (3.6,0);
  \fill[texto] (Z) circle (0.04);
  \draw[campoE, line width=0.55pt] (-1.6,1.3) -- (Z);
  \draw[campoD, line width=0.55pt] (1.6,1.3) -- (Z);
  \draw[tenue, line width=0.4pt] ($(Z)+(-0.95,0)$) arc (180:{180-atan(1.3/5.2)}:0.95);
  \node[etiqt] at ($(Z)+(-1.2,0.13)$) {$\theta_2$};
  \draw[tenue, line width=0.4pt] ($(Z)+(-0.5,0)$) arc (180:{180-atan(1.3/2.0)}:0.5);
  \node[etiqt] at ($(Z)+(-0.4,0.42)$) {$\theta_1$};
  \foreach \x/\t in {-1.6/{-\ell},0/0,1.6/\ell} \node[etiqt, anchor=north] at (\x,-0.05) {$\t$};
  \node[etiqt, anchor=north] at (Z) {$z$};
\end{tikzpicture}
\caption{Punto exterior, $z>\ell$}
\end{subfigure}\hfill
\begin{subfigure}[b]{0.49\textwidth}\centering
\begin{tikzpicture}[scale=0.9]
  \draw[texto!70!fondo, line width=0.5pt, -{Stealth[length=2mm]}] (-2.6,0) -- (3.2,0) node[below, etiq] {$z$};
  \fill[dielecB, opacity=0.8] (-1.6,0) ellipse (0.3 and 1.3);
  \fill[dielecB, opacity=0.8] (1.6,0) ellipse (0.3 and 1.3);
  \draw[borde] (-1.6,0) ellipse (0.3 and 1.3) (1.6,0) ellipse (0.3 and 1.3);
  \node[cneg, font=\small] at (-1.6,1.6) {$-P$};
  \node[cpos, font=\small] at (1.6,1.6) {$+P$};
  \coordinate (Z) at (0.55,0);
  \fill[texto] (Z) circle (0.04);
  \draw[campoE, line width=0.55pt] (-1.6,1.3) -- (Z);
  \draw[campoD, line width=0.55pt] (1.6,1.3) -- (Z);
  \draw[tenue, line width=0.4pt] ($(Z)+(-0.7,0)$) arc (180:{180-atan(1.3/2.15)}:0.7);
  \node[etiqt] at ($(Z)+(-0.95,0.2)$) {$\theta_2$};
  \draw[tenue, line width=0.4pt] ($(Z)+(0.42,0)$) arc (0:{atan(1.3/1.05)}:0.42);
  \node[etiqt] at ($(Z)+(0.75,0.28)$) {$\pi-\theta_1$};
  \draw[tenue, line width=0.4pt, densely dotted] ($(Z)+(-0.3,0)$) arc (180:{atan(1.3/1.05)}:0.3);
  \node[etiqt] at ($(Z)+(-0.05,0.5)$) {$\theta_1$};
  \foreach \x/\t in {-1.6/{-\ell},1.6/\ell} \node[etiqt, anchor=north] at (\x,-0.05) {$\t$};
  \node[etiqt, anchor=north] at (0,-0.05) {$0$};
  \node[etiqt, anchor=north] at (Z) {$z$};
\end{tikzpicture}
\caption{Punto interior, $|z|<\ell$}
\end{subfigure}
\caption{\Cref{prob:cilindro-polarizado}. (a) En la pared lateral $\vP$ salta de $P\,\uv z$ a $\vcero$: la circulación sobre un
rectángulo que la atraviesa no se anula, luego $\rot\vP\neq\vcero$ y $\vD$ no tiene por qué anularse. Como $\rhoP=0$,
el cilindro equivale a dos discos con densidades $\mp P$. (b), (c) Ángulos que el eje forma con las
rectas que unen el punto con el borde de cada disco: $\cos\theta_1=(z-\ell)/\sqrt{a^2+(z-\ell)^2}$ y
$\cos\theta_2=(z+\ell)/\sqrt{a^2+(z+\ell)^2}$; para $|z|<\ell$, $\theta_1>\pi/2$.}
\label{fig:p44}
\end{figure}

% Figura: problema 4.4. Perfiles axiales de E (unid. P/eps0), D y P (unid. P) para
% tres relaciones de aspecto a/l.  Funciones analíticas: \curvaEeje, \curvaDeje (preámbulo).
\begin{figure}[t]
\centering
\newcommand{\icono}[2]{\draw[texto, fill=dielecB, line width=0.35pt]
   (rel axis cs:0.86,0.82) ++(-#1/2,-#2/2) rectangle ++(#1,#2);}
\begin{tikzpicture}
\begin{groupplot}[group style={group size=3 by 2, horizontal sep=0.95cm, vertical sep=1.05cm},
   perfil, width=5.55cm, height=4.0cm]
  % --- fila superior: E y P
  \nextgroupplot[ymin=-1.15, ymax=1.15, ytick={-1,-0.5,0,0.5,1}, ylabel={$E,\ P$}]
    \curvaPeje \curvaEeje{1}
    \node[etiqf, font=\tiny, anchor=north west] at (rel axis cs:0.02,0.98) {$a/\ell=1$};
    \icono{0.30cm}{0.30cm}
  \nextgroupplot[ymin=-1.15, ymax=1.15, ytick={-1,-0.5,0,0.5,1}]
    \curvaPeje \curvaEeje{50}
    \node[etiqf, font=\tiny, anchor=north west] at (rel axis cs:0.02,0.98) {$a/\ell=50$};
    \icono{0.05cm}{0.5cm}
  \nextgroupplot[ymin=-1.15, ymax=1.15, ytick={-1,-0.5,0,0.5,1}]
    \curvaPeje \curvaEeje{0.1}
    \node[etiqf, font=\tiny, anchor=north west] at (rel axis cs:0.02,0.98) {$a/\ell=0{,}1$};
    \icono{0.55cm}{0.06cm}
  % --- fila inferior: D y P
  \nextgroupplot[ymin=-0.08, ymax=1.15, ytick={0,0.5,1}, ylabel={$D,\ P$}]
    \curvaPeje \curvaDeje{1}
  \nextgroupplot[ymin=-0.08, ymax=1.15, ytick={0,0.5,1}]
    \curvaPeje \curvaDeje{50}
  \nextgroupplot[ymin=-0.08, ymax=1.15, ytick={0,0.5,1}]
    \curvaPeje \curvaDeje{0.1}
\end{groupplot}
\end{tikzpicture}
\caption{\Cref{prob:cilindro-polarizado}. Campos sobre el eje del electrete cilíndrico: $E_z$ en unidades de $P/\eo$ (continua, cian),
$D_z$ en unidades de $P$ (continua, violeta) y $P$ (trazos). Arriba: $E_z$ salta $P/\eo$ en cada cara y dentro se
opone a $\vP$. Abajo: $D_z$ es continuo (no hay carga libre) y coincide con la polarización solo en el límite del
cilindro largo. Columnas: cilindro equilátero, disco delgado ($E\simeq-P/\eo$, $D\simeq0$) y cilindro largo
($E\simeq0$, $D\simeq P$ lejos de los extremos).}
\label{fig:p44-perfiles}
\end{figure}

% Figura (simulación): electrete cilíndrico uniformemente polarizado (problema 4.4),
% líneas de P, E y D.  Datos: scripts/lineas_campo.py -> datos/cilindro-polarizado-{E,D}.tex
\begin{figure}[t]
\centering
\pgfplotsset{cil/.style={marco, width=5.55cm, xmin=-3.2, xmax=3.2, ymin=-1.9, ymax=1.9}}
\begin{subfigure}[b]{0.325\textwidth}\centering
\begin{tikzpicture}
\begin{axis}[cil]
  \path (rel axis cs:0,0) rectangle (rel axis cs:1,1);
  \fill[dielec] (axis cs:-1.5,-0.5) rectangle (axis cs:1.5,0.5);
  \pgfplotsinvokeforeach{-0.4,-0.2,0,0.2,0.4}{
    \addplot[draw=campoP, line width=0.6pt, flechaen=0.55] coordinates {(-1.5,#1) (1.5,#1)};}
  \draw[texto, line width=0.6pt] (axis cs:-1.5,-0.5) rectangle (axis cs:1.5,0.5);
\end{axis}
\end{tikzpicture}
\caption{$\vP(\vr)$}
\end{subfigure}\hfill
\begin{subfigure}[b]{0.325\textwidth}\centering
\begin{tikzpicture}
\begin{axis}[cil]
  \fill[dielec] (axis cs:-1.5,-0.5) rectangle (axis cs:1.5,0.5);
  %% Generado por scripts/lineas_campo.py -- no editar a mano
\addplot[equipot] coordinates {(-1.500,-0.395) (-1.516,-0.364) (-1.531,-0.332) (-1.545,-0.300) (-1.557,-0.267) (-1.568,-0.234) (-1.577,-0.200) (-1.585,-0.166) (-1.591,-0.131) (-1.596,-0.097) (-1.599,-0.062) (-1.601,-0.027) (-1.601,0.008) (-1.600,0.043) (-1.598,0.078) (-1.594,0.113) (-1.588,0.147) (-1.581,0.182) (-1.573,0.216) (-1.563,0.249) (-1.552,0.282) (-1.539,0.315) (-1.524,0.347) (-1.508,0.378) (-1.494,0.380) (-1.477,0.349) (-1.463,0.317) (-1.451,0.284) (-1.439,0.251) (-1.430,0.217) (-1.422,0.183) (-1.415,0.149) (-1.410,0.114) (-1.406,0.079) (-1.403,0.044) (-1.402,0.009) (-1.402,-0.026) (-1.404,-0.061) (-1.407,-0.095) (-1.412,-0.130) (-1.418,-0.165) (-1.425,-0.199) (-1.434,-0.233) (-1.444,-0.266) (-1.456,-0.299) (-1.470,-0.332) (-1.484,-0.363) (-1.500,-0.395)};
\addplot[equipot] coordinates {(-1.500,-0.482) (-1.526,-0.459) (-1.550,-0.434) (-1.572,-0.406) (-1.592,-0.378) (-1.610,-0.347) (-1.627,-0.316) (-1.641,-0.284) (-1.654,-0.252) (-1.665,-0.218) (-1.674,-0.185) (-1.682,-0.150) (-1.688,-0.116) (-1.692,-0.081) (-1.695,-0.046) (-1.697,-0.011) (-1.696,0.024) (-1.694,0.059) (-1.691,0.094) (-1.686,0.129) (-1.679,0.163) (-1.671,0.197) (-1.661,0.231) (-1.649,0.264) (-1.636,0.297) (-1.621,0.328) (-1.604,0.359) (-1.585,0.389) (-1.564,0.417) (-1.541,0.444) (-1.516,0.468) (-1.490,0.474) (-1.465,0.449) (-1.442,0.422) (-1.422,0.394) (-1.403,0.364) (-1.386,0.333) (-1.372,0.301) (-1.358,0.269) (-1.347,0.235) (-1.337,0.202) (-1.329,0.168) (-1.323,0.133) (-1.317,0.098) (-1.314,0.063) (-1.312,0.028) (-1.312,-0.007) (-1.313,-0.042) (-1.315,-0.077) (-1.319,-0.112) (-1.325,-0.146) (-1.332,-0.181) (-1.341,-0.215) (-1.351,-0.248) (-1.363,-0.281) (-1.377,-0.313) (-1.392,-0.345) (-1.410,-0.375) (-1.429,-0.405) (-1.451,-0.432) (-1.474,-0.458) (-1.500,-0.482)};
\addplot[equipot] coordinates {(-1.548,-0.528) (-1.580,-0.514) (-1.610,-0.495) (-1.637,-0.472) (-1.662,-0.448) (-1.685,-0.421) (-1.707,-0.393) (-1.726,-0.363) (-1.743,-0.333) (-1.759,-0.301) (-1.773,-0.269) (-1.785,-0.236) (-1.796,-0.202) (-1.804,-0.168) (-1.812,-0.133) (-1.817,-0.098) (-1.821,-0.063) (-1.823,-0.028) (-1.823,0.007) (-1.822,0.043) (-1.819,0.078) (-1.815,0.113) (-1.809,0.148) (-1.801,0.182) (-1.791,0.216) (-1.780,0.250) (-1.767,0.282) (-1.753,0.315) (-1.736,0.346) (-1.718,0.376) (-1.698,0.405) (-1.676,0.432) (-1.652,0.458) (-1.626,0.482) (-1.597,0.503) (-1.567,0.520) (-1.534,0.533) (-1.499,0.537) (-1.464,0.531) (-1.432,0.517) (-1.402,0.498) (-1.375,0.475) (-1.350,0.450) (-1.327,0.423) (-1.307,0.395) (-1.288,0.365) (-1.271,0.334) (-1.256,0.302) (-1.243,0.269) (-1.231,0.236) (-1.221,0.202) (-1.213,0.168) (-1.206,0.133) (-1.201,0.098) (-1.197,0.063) (-1.196,0.028) (-1.195,-0.007) (-1.196,-0.043) (-1.199,-0.078) (-1.203,-0.113) (-1.209,-0.148) (-1.216,-0.182) (-1.225,-0.216) (-1.236,-0.250) (-1.248,-0.283) (-1.262,-0.316) (-1.278,-0.347) (-1.295,-0.378) (-1.315,-0.407) (-1.337,-0.435) (-1.360,-0.461) (-1.386,-0.485) (-1.414,-0.506) (-1.445,-0.523) (-1.478,-0.535) (-1.514,-0.536) (-1.548,-0.528)};
\addplot[equipot] coordinates {(-1.548,-0.656) (-1.583,-0.651) (-1.617,-0.642) (-1.650,-0.630) (-1.682,-0.616) (-1.712,-0.599) (-1.742,-0.579) (-1.770,-0.558) (-1.796,-0.535) (-1.821,-0.510) (-1.845,-0.484) (-1.867,-0.457) (-1.887,-0.428) (-1.906,-0.399) (-1.924,-0.368) (-1.940,-0.337) (-1.954,-0.305) (-1.967,-0.272) (-1.978,-0.239) (-1.988,-0.205) (-1.996,-0.171) (-2.003,-0.137) (-2.008,-0.102) (-2.011,-0.067) (-2.014,-0.032) (-2.014,0.003) (-2.013,0.039) (-2.011,0.074) (-2.007,0.108) (-2.001,0.143) (-1.994,0.178) (-1.986,0.212) (-1.976,0.245) (-1.964,0.279) (-1.951,0.311) (-1.937,0.343) (-1.920,0.374) (-1.903,0.405) (-1.883,0.434) (-1.863,0.462) (-1.840,0.490) (-1.817,0.515) (-1.791,0.540) (-1.764,0.562) (-1.736,0.583) (-1.707,0.602) (-1.676,0.619) (-1.643,0.633) (-1.610,0.644) (-1.576,0.652) (-1.541,0.657) (-1.506,0.658) (-1.471,0.656) (-1.436,0.650) (-1.403,0.640) (-1.370,0.628) (-1.338,0.612) (-1.308,0.594) (-1.279,0.574) (-1.252,0.552) (-1.227,0.527) (-1.203,0.501) (-1.181,0.474) (-1.160,0.446) (-1.141,0.416) (-1.124,0.386) (-1.108,0.354) (-1.094,0.322) (-1.081,0.289) (-1.070,0.256) (-1.060,0.222) (-1.052,0.188) (-1.045,0.154) (-1.039,0.119) (-1.035,0.084) (-1.032,0.049) (-1.031,0.014) (-1.031,-0.021) (-1.033,-0.056) (-1.036,-0.091) (-1.040,-0.126) (-1.046,-0.161) (-1.053,-0.195) (-1.062,-0.229) (-1.072,-0.263) (-1.084,-0.296) (-1.097,-0.328) (-1.111,-0.360) (-1.127,-0.392) (-1.145,-0.422) (-1.164,-0.451) (-1.185,-0.480) (-1.208,-0.507) (-1.232,-0.532) (-1.257,-0.556) (-1.285,-0.578) (-1.314,-0.598) (-1.344,-0.616) (-1.376,-0.631) (-1.409,-0.642) (-1.443,-0.651) (-1.478,-0.656) (-1.513,-0.658) (-1.548,-0.656)};
\addplot[equipot] coordinates {(-1.660,-0.922) (-1.695,-0.915) (-1.729,-0.907) (-1.763,-0.898) (-1.796,-0.887) (-1.829,-0.874) (-1.861,-0.860) (-1.893,-0.844) (-1.924,-0.827) (-1.954,-0.809) (-1.983,-0.789) (-2.011,-0.769) (-2.039,-0.747) (-2.065,-0.724) (-2.091,-0.699) (-2.115,-0.674) (-2.139,-0.648) (-2.162,-0.621) (-2.183,-0.593) (-2.204,-0.565) (-2.223,-0.535) (-2.241,-0.505) (-2.258,-0.474) (-2.274,-0.443) (-2.289,-0.411) (-2.303,-0.379) (-2.315,-0.346) (-2.327,-0.313) (-2.337,-0.279) (-2.346,-0.245) (-2.354,-0.211) (-2.360,-0.176) (-2.366,-0.141) (-2.370,-0.106) (-2.373,-0.071) (-2.375,-0.036) (-2.376,-0.001) (-2.375,0.034) (-2.373,0.069) (-2.370,0.104) (-2.366,0.139) (-2.361,0.174) (-2.354,0.209) (-2.346,0.243) (-2.337,0.277) (-2.327,0.311) (-2.316,0.344) (-2.304,0.377) (-2.290,0.409) (-2.275,0.441) (-2.259,0.473) (-2.242,0.503) (-2.224,0.534) (-2.205,0.563) (-2.184,0.592) (-2.163,0.620) (-2.140,0.647) (-2.117,0.673) (-2.092,0.698) (-2.067,0.722) (-2.040,0.745) (-2.013,0.767) (-1.985,0.788) (-1.955,0.808) (-1.925,0.826) (-1.895,0.843) (-1.863,0.859) (-1.831,0.873) (-1.798,0.886) (-1.765,0.897) (-1.731,0.907) (-1.697,0.915) (-1.662,0.921) (-1.627,0.926) (-1.592,0.929) (-1.557,0.930) (-1.522,0.930) (-1.487,0.927) (-1.452,0.923) (-1.417,0.917) (-1.383,0.909) (-1.349,0.900) (-1.316,0.889) (-1.283,0.876) (-1.251,0.861) (-1.219,0.846) (-1.189,0.828) (-1.159,0.809) (-1.130,0.789) (-1.103,0.767) (-1.076,0.744) (-1.050,0.720) (-1.026,0.695) (-1.002,0.669) (-0.980,0.642) (-0.959,0.613) (-0.939,0.584) (-0.920,0.555) (-0.902,0.524) (-0.886,0.493) (-0.870,0.462) (-0.856,0.430) (-0.843,0.397) (-0.831,0.364) (-0.820,0.330) (-0.810,0.297) (-0.801,0.263) (-0.794,0.228) (-0.787,0.194) (-0.782,0.159) (-0.778,0.124) (-0.774,0.089) (-0.772,0.054) (-0.771,0.019) (-0.771,-0.017) (-0.772,-0.052) (-0.774,-0.087) (-0.777,-0.122) (-0.782,-0.157) (-0.787,-0.192) (-0.793,-0.226) (-0.801,-0.261) (-0.809,-0.295) (-0.819,-0.329) (-0.830,-0.362) (-0.842,-0.395) (-0.855,-0.428) (-0.869,-0.460) (-0.885,-0.492) (-0.901,-0.523) (-0.919,-0.553) (-0.938,-0.583) (-0.958,-0.612) (-0.979,-0.640) (-1.001,-0.667) (-1.024,-0.693) (-1.049,-0.719) (-1.075,-0.743) (-1.101,-0.766) (-1.129,-0.788) (-1.158,-0.808) (-1.187,-0.827) (-1.218,-0.845) (-1.249,-0.861) (-1.281,-0.875) (-1.314,-0.888) (-1.347,-0.899) (-1.381,-0.909) (-1.415,-0.917) (-1.450,-0.923) (-1.485,-0.927) (-1.520,-0.929) (-1.555,-0.930) (-1.590,-0.929) (-1.625,-0.926) (-1.660,-0.922)};
\addplot[equipot] coordinates {(-3.200,1.227) (-3.178,1.254) (-3.155,1.281) (-3.131,1.307) (-3.107,1.333) (-3.083,1.358) (-3.058,1.382) (-3.033,1.407) (-3.007,1.430) (-2.981,1.454) (-2.954,1.476) (-2.927,1.499) (-2.899,1.520) (-2.871,1.541) (-2.843,1.562) (-2.814,1.582) (-2.785,1.601) (-2.755,1.620) (-2.725,1.639) (-2.695,1.656) (-2.665,1.673) (-2.634,1.690) (-2.602,1.706) (-2.571,1.721) (-2.539,1.735) (-2.507,1.749) (-2.474,1.762) (-2.441,1.775) (-2.408,1.787) (-2.375,1.798) (-2.342,1.808) (-2.308,1.818) (-2.274,1.827) (-2.240,1.835) (-2.206,1.843) (-2.171,1.850) (-2.137,1.856) (-2.102,1.861) (-2.067,1.866) (-2.033,1.869) (-1.998,1.872) (-1.963,1.875) (-1.928,1.876) (-1.893,1.877) (-1.857,1.877) (-1.822,1.876) (-1.787,1.874) (-1.752,1.872) (-1.718,1.868) (-1.683,1.864) (-1.648,1.859) (-1.613,1.854) (-1.579,1.847) (-1.545,1.840) (-1.511,1.832) (-1.477,1.823) (-1.443,1.813) (-1.409,1.803) (-1.376,1.791) (-1.343,1.779) (-1.311,1.766) (-1.278,1.753) (-1.247,1.738) (-1.215,1.723) (-1.184,1.707) (-1.153,1.690) (-1.122,1.673) (-1.093,1.655) (-1.063,1.636) (-1.034,1.616) (-1.005,1.596) (-0.977,1.575) (-0.950,1.553) (-0.923,1.531) (-0.896,1.508) (-0.870,1.484) (-0.845,1.460) (-0.820,1.435) (-0.796,1.410) (-0.772,1.384) (-0.749,1.358) (-0.727,1.331) (-0.705,1.303) (-0.684,1.275) (-0.663,1.247) (-0.643,1.218) (-0.624,1.189) (-0.605,1.159) (-0.588,1.129) (-0.570,1.098) (-0.553,1.067) (-0.537,1.036) (-0.522,1.005) (-0.507,0.973) (-0.493,0.941) (-0.479,0.909) (-0.466,0.876) (-0.454,0.843) (-0.442,0.810) (-0.431,0.777) (-0.420,0.744) (-0.410,0.710) (-0.401,0.676) (-0.392,0.642) (-0.383,0.608) (-0.375,0.574) (-0.368,0.540) (-0.361,0.506) (-0.355,0.471) (-0.349,0.436) (-0.344,0.402) (-0.339,0.367) (-0.334,0.332) (-0.330,0.297) (-0.327,0.263) (-0.324,0.228) (-0.321,0.193) (-0.319,0.158) (-0.317,0.123) (-0.316,0.088) (-0.315,0.053) (-0.315,0.018) (-0.315,-0.018) (-0.315,-0.053) (-0.316,-0.088) (-0.317,-0.123) (-0.319,-0.158) (-0.321,-0.193) (-0.324,-0.228) (-0.327,-0.263) (-0.330,-0.297) (-0.334,-0.332) (-0.339,-0.367) (-0.344,-0.402) (-0.349,-0.436) (-0.355,-0.471) (-0.361,-0.506) (-0.368,-0.540) (-0.375,-0.574) (-0.383,-0.608) (-0.392,-0.642) (-0.401,-0.676) (-0.410,-0.710) (-0.420,-0.744) (-0.431,-0.777) (-0.442,-0.810) (-0.454,-0.843) (-0.466,-0.876) (-0.479,-0.909) (-0.493,-0.941) (-0.507,-0.973) (-0.522,-1.005) (-0.537,-1.036) (-0.553,-1.067) (-0.570,-1.098) (-0.588,-1.129) (-0.605,-1.159) (-0.624,-1.189) (-0.643,-1.218) (-0.663,-1.247) (-0.684,-1.275) (-0.705,-1.303) (-0.727,-1.331) (-0.749,-1.358) (-0.772,-1.384) (-0.796,-1.410) (-0.820,-1.435) (-0.845,-1.460) (-0.870,-1.484) (-0.896,-1.508) (-0.923,-1.531) (-0.950,-1.553) (-0.977,-1.575) (-1.005,-1.596) (-1.034,-1.616) (-1.063,-1.636) (-1.093,-1.655) (-1.122,-1.673) (-1.153,-1.690) (-1.184,-1.707) (-1.215,-1.723) (-1.247,-1.738) (-1.278,-1.753) (-1.311,-1.766) (-1.343,-1.779) (-1.376,-1.791) (-1.409,-1.803) (-1.443,-1.813) (-1.477,-1.823) (-1.511,-1.832) (-1.545,-1.840) (-1.579,-1.847) (-1.613,-1.854) (-1.648,-1.859) (-1.683,-1.864) (-1.718,-1.868) (-1.752,-1.872) (-1.787,-1.874) (-1.822,-1.876) (-1.857,-1.877) (-1.893,-1.877) (-1.928,-1.876) (-1.963,-1.875) (-1.998,-1.872) (-2.033,-1.869) (-2.067,-1.866) (-2.102,-1.861) (-2.137,-1.856) (-2.171,-1.850) (-2.206,-1.843) (-2.240,-1.835) (-2.274,-1.827) (-2.308,-1.818) (-2.342,-1.808) (-2.375,-1.798) (-2.408,-1.787) (-2.441,-1.775) (-2.474,-1.762) (-2.507,-1.749) (-2.539,-1.735) (-2.571,-1.721) (-2.602,-1.706) (-2.634,-1.690) (-2.665,-1.673) (-2.695,-1.656) (-2.725,-1.639) (-2.755,-1.620) (-2.785,-1.601) (-2.814,-1.582) (-2.843,-1.562) (-2.871,-1.541) (-2.899,-1.520) (-2.927,-1.499) (-2.954,-1.476) (-2.981,-1.454) (-3.007,-1.430) (-3.033,-1.407) (-3.058,-1.382) (-3.083,-1.358) (-3.107,-1.333) (-3.131,-1.307) (-3.155,-1.281) (-3.178,-1.254) (-3.200,-1.227)};
\addplot[equipot] coordinates {(3.200,1.227) (3.178,1.254) (3.155,1.281) (3.131,1.307) (3.107,1.333) (3.083,1.358) (3.058,1.382) (3.033,1.407) (3.007,1.430) (2.981,1.454) (2.954,1.476) (2.927,1.499) (2.899,1.520) (2.871,1.541) (2.843,1.562) (2.814,1.582) (2.785,1.601) (2.755,1.620) (2.725,1.639) (2.695,1.656) (2.665,1.673) (2.634,1.690) (2.602,1.706) (2.571,1.721) (2.539,1.735) (2.507,1.749) (2.474,1.762) (2.441,1.775) (2.408,1.787) (2.375,1.798) (2.342,1.808) (2.308,1.818) (2.274,1.827) (2.240,1.835) (2.206,1.843) (2.171,1.850) (2.137,1.856) (2.102,1.861) (2.067,1.866) (2.033,1.869) (1.998,1.872) (1.963,1.875) (1.928,1.876) (1.893,1.877) (1.857,1.877) (1.822,1.876) (1.787,1.874) (1.752,1.872) (1.718,1.868) (1.683,1.864) (1.648,1.859) (1.613,1.854) (1.579,1.847) (1.545,1.840) (1.511,1.832) (1.477,1.823) (1.443,1.813) (1.409,1.803) (1.376,1.791) (1.343,1.779) (1.311,1.766) (1.278,1.753) (1.247,1.738) (1.215,1.723) (1.184,1.707) (1.153,1.690) (1.122,1.673) (1.093,1.655) (1.063,1.636) (1.034,1.616) (1.005,1.596) (0.977,1.575) (0.950,1.553) (0.923,1.531) (0.896,1.508) (0.870,1.484) (0.845,1.460) (0.820,1.435) (0.796,1.410) (0.772,1.384) (0.749,1.358) (0.727,1.331) (0.705,1.303) (0.684,1.275) (0.663,1.247) (0.643,1.218) (0.624,1.189) (0.605,1.159) (0.588,1.129) (0.570,1.098) (0.553,1.067) (0.537,1.036) (0.522,1.005) (0.507,0.973) (0.493,0.941) (0.479,0.909) (0.466,0.876) (0.454,0.843) (0.442,0.810) (0.431,0.777) (0.420,0.744) (0.410,0.710) (0.401,0.676) (0.392,0.642) (0.383,0.608) (0.375,0.574) (0.368,0.540) (0.361,0.506) (0.355,0.471) (0.349,0.436) (0.344,0.402) (0.339,0.367) (0.334,0.332) (0.330,0.297) (0.327,0.263) (0.324,0.228) (0.321,0.193) (0.319,0.158) (0.317,0.123) (0.316,0.088) (0.315,0.053) (0.315,0.018) (0.315,-0.018) (0.315,-0.053) (0.316,-0.088) (0.317,-0.123) (0.319,-0.158) (0.321,-0.193) (0.324,-0.228) (0.327,-0.263) (0.330,-0.297) (0.334,-0.332) (0.339,-0.367) (0.344,-0.402) (0.349,-0.436) (0.355,-0.471) (0.361,-0.506) (0.368,-0.540) (0.375,-0.574) (0.383,-0.608) (0.392,-0.642) (0.401,-0.676) (0.410,-0.710) (0.420,-0.744) (0.431,-0.777) (0.442,-0.810) (0.454,-0.843) (0.466,-0.876) (0.479,-0.909) (0.493,-0.941) (0.507,-0.973) (0.522,-1.005) (0.537,-1.036) (0.553,-1.067) (0.570,-1.098) (0.588,-1.129) (0.605,-1.159) (0.624,-1.189) (0.643,-1.218) (0.663,-1.247) (0.684,-1.275) (0.705,-1.303) (0.727,-1.331) (0.749,-1.358) (0.772,-1.384) (0.796,-1.410) (0.820,-1.435) (0.845,-1.460) (0.870,-1.484) (0.896,-1.508) (0.923,-1.531) (0.950,-1.553) (0.977,-1.575) (1.005,-1.596) (1.034,-1.616) (1.063,-1.636) (1.093,-1.655) (1.122,-1.673) (1.153,-1.690) (1.184,-1.707) (1.215,-1.723) (1.247,-1.738) (1.278,-1.753) (1.311,-1.766) (1.343,-1.779) (1.376,-1.791) (1.409,-1.803) (1.443,-1.813) (1.477,-1.823) (1.511,-1.832) (1.545,-1.840) (1.579,-1.847) (1.613,-1.854) (1.648,-1.859) (1.683,-1.864) (1.718,-1.868) (1.752,-1.872) (1.787,-1.874) (1.822,-1.876) (1.857,-1.877) (1.893,-1.877) (1.928,-1.876) (1.963,-1.875) (1.998,-1.872) (2.033,-1.869) (2.067,-1.866) (2.102,-1.861) (2.137,-1.856) (2.171,-1.850) (2.206,-1.843) (2.240,-1.835) (2.274,-1.827) (2.308,-1.818) (2.342,-1.808) (2.375,-1.798) (2.408,-1.787) (2.441,-1.775) (2.474,-1.762) (2.507,-1.749) (2.539,-1.735) (2.571,-1.721) (2.602,-1.706) (2.634,-1.690) (2.665,-1.673) (2.695,-1.656) (2.725,-1.639) (2.755,-1.620) (2.785,-1.601) (2.814,-1.582) (2.843,-1.562) (2.871,-1.541) (2.899,-1.520) (2.927,-1.499) (2.954,-1.476) (2.981,-1.454) (3.007,-1.430) (3.033,-1.407) (3.058,-1.382) (3.083,-1.358) (3.107,-1.333) (3.131,-1.307) (3.155,-1.281) (3.178,-1.254) (3.200,-1.227)};
\addplot[equipot] coordinates {(1.452,-0.923) (1.487,-0.927) (1.522,-0.930) (1.557,-0.930) (1.592,-0.929) (1.627,-0.926) (1.662,-0.921) (1.697,-0.915) (1.731,-0.907) (1.765,-0.897) (1.798,-0.886) (1.831,-0.873) (1.863,-0.859) (1.894,-0.844) (1.925,-0.826) (1.955,-0.808) (1.984,-0.788) (2.013,-0.768) (2.040,-0.745) (2.067,-0.722) (2.092,-0.698) (2.117,-0.673) (2.140,-0.647) (2.163,-0.620) (2.184,-0.592) (2.205,-0.563) (2.224,-0.534) (2.242,-0.504) (2.259,-0.473) (2.275,-0.441) (2.290,-0.410) (2.304,-0.377) (2.316,-0.344) (2.327,-0.311) (2.337,-0.277) (2.346,-0.243) (2.354,-0.209) (2.361,-0.174) (2.366,-0.139) (2.370,-0.105) (2.373,-0.070) (2.375,-0.034) (2.376,0.001) (2.375,0.036) (2.373,0.071) (2.370,0.106) (2.366,0.141) (2.360,0.176) (2.354,0.210) (2.346,0.245) (2.337,0.279) (2.327,0.312) (2.315,0.346) (2.303,0.379) (2.289,0.411) (2.274,0.443) (2.258,0.474) (2.241,0.505) (2.223,0.535) (2.204,0.565) (2.183,0.593) (2.162,0.621) (2.139,0.648) (2.116,0.674) (2.091,0.699) (2.065,0.723) (2.039,0.747) (2.011,0.769) (1.983,0.789) (1.954,0.809) (1.924,0.827) (1.893,0.844) (1.861,0.860) (1.829,0.874) (1.796,0.887) (1.763,0.898) (1.729,0.907) (1.695,0.915) (1.660,0.922) (1.625,0.926) (1.590,0.929) (1.555,0.930) (1.520,0.929) (1.485,0.927) (1.450,0.923) (1.415,0.917) (1.381,0.909) (1.347,0.899) (1.314,0.888) (1.281,0.875) (1.249,0.861) (1.218,0.845) (1.187,0.827) (1.158,0.808) (1.129,0.788) (1.101,0.766) (1.075,0.743) (1.049,0.719) (1.025,0.694) (1.001,0.667) (0.979,0.640) (0.958,0.612) (0.938,0.583) (0.919,0.553) (0.901,0.523) (0.885,0.492) (0.869,0.460) (0.855,0.428) (0.842,0.395) (0.830,0.362) (0.819,0.329) (0.809,0.295) (0.801,0.261) (0.793,0.226) (0.787,0.192) (0.782,0.157) (0.777,0.122) (0.774,0.087) (0.772,0.052) (0.771,0.017) (0.771,-0.018) (0.772,-0.054) (0.774,-0.089) (0.778,-0.124) (0.782,-0.159) (0.787,-0.193) (0.794,-0.228) (0.801,-0.262) (0.810,-0.296) (0.820,-0.330) (0.831,-0.364) (0.843,-0.397) (0.856,-0.429) (0.870,-0.462) (0.885,-0.493) (0.902,-0.524) (0.920,-0.555) (0.939,-0.584) (0.959,-0.613) (0.980,-0.641) (1.002,-0.669) (1.026,-0.695) (1.050,-0.720) (1.076,-0.744) (1.103,-0.767) (1.130,-0.789) (1.159,-0.809) (1.189,-0.828) (1.219,-0.845) (1.251,-0.861) (1.283,-0.876) (1.315,-0.889) (1.349,-0.900) (1.383,-0.909) (1.417,-0.917) (1.452,-0.923)};
\addplot[equipot] coordinates {(1.468,-0.655) (1.503,-0.658) (1.538,-0.657) (1.573,-0.653) (1.607,-0.645) (1.640,-0.634) (1.673,-0.620) (1.704,-0.604) (1.733,-0.585) (1.762,-0.564) (1.789,-0.542) (1.814,-0.518) (1.838,-0.492) (1.861,-0.465) (1.882,-0.437) (1.901,-0.407) (1.919,-0.377) (1.935,-0.346) (1.950,-0.314) (1.963,-0.282) (1.975,-0.249) (1.985,-0.215) (1.994,-0.181) (2.001,-0.146) (2.006,-0.112) (2.011,-0.077) (2.013,-0.042) (2.014,-0.007) (2.014,0.028) (2.012,0.063) (2.008,0.098) (2.003,0.133) (1.997,0.168) (1.989,0.202) (1.979,0.236) (1.968,0.269) (1.955,0.302) (1.941,0.334) (1.925,0.365) (1.908,0.396) (1.889,0.426) (1.869,0.454) (1.847,0.482) (1.824,0.508) (1.799,0.533) (1.772,0.556) (1.745,0.577) (1.715,0.597) (1.685,0.614) (1.653,0.629) (1.620,0.641) (1.586,0.650) (1.551,0.656) (1.516,0.658) (1.481,0.657) (1.446,0.652) (1.412,0.643) (1.379,0.632) (1.347,0.617) (1.317,0.600) (1.288,0.580) (1.260,0.558) (1.234,0.534) (1.210,0.509) (1.187,0.482) (1.166,0.454) (1.147,0.425) (1.129,0.395) (1.113,0.363) (1.098,0.331) (1.085,0.299) (1.073,0.266) (1.063,0.232) (1.054,0.198) (1.047,0.164) (1.041,0.129) (1.036,0.094) (1.033,0.059) (1.031,0.024) (1.031,-0.011) (1.032,-0.046) (1.035,-0.081) (1.039,-0.116) (1.044,-0.151) (1.051,-0.185) (1.059,-0.219) (1.069,-0.253) (1.080,-0.286) (1.093,-0.319) (1.107,-0.351) (1.123,-0.383) (1.140,-0.413) (1.159,-0.443) (1.179,-0.472) (1.201,-0.499) (1.225,-0.525) (1.250,-0.549) (1.277,-0.572) (1.305,-0.593) (1.335,-0.611) (1.367,-0.627) (1.399,-0.639) (1.433,-0.649) (1.468,-0.655)};
\addplot[equipot] coordinates {(1.468,-0.532) (1.503,-0.537) (1.537,-0.532) (1.570,-0.519) (1.600,-0.501) (1.629,-0.480) (1.654,-0.456) (1.678,-0.430) (1.700,-0.402) (1.720,-0.373) (1.738,-0.342) (1.754,-0.311) (1.769,-0.279) (1.782,-0.246) (1.793,-0.212) (1.802,-0.178) (1.809,-0.144) (1.815,-0.109) (1.820,-0.074) (1.822,-0.039) (1.823,-0.004) (1.823,0.032) (1.820,0.067) (1.816,0.102) (1.811,0.137) (1.804,0.171) (1.795,0.205) (1.784,0.239) (1.772,0.272) (1.757,0.304) (1.742,0.336) (1.724,0.367) (1.704,0.396) (1.683,0.424) (1.660,0.450) (1.634,0.475) (1.606,0.497) (1.577,0.515) (1.544,0.529) (1.510,0.537) (1.475,0.534) (1.442,0.522) (1.411,0.504) (1.383,0.483) (1.357,0.459) (1.334,0.432) (1.313,0.404) (1.293,0.374) (1.276,0.344) (1.260,0.312) (1.247,0.280) (1.234,0.246) (1.224,0.213) (1.215,0.179) (1.208,0.144) (1.202,0.109) (1.198,0.074) (1.196,0.039) (1.195,0.004) (1.196,-0.032) (1.198,-0.067) (1.201,-0.102) (1.207,-0.137) (1.214,-0.171) (1.222,-0.206) (1.232,-0.240) (1.244,-0.273) (1.258,-0.305) (1.273,-0.337) (1.290,-0.368) (1.309,-0.398) (1.330,-0.426) (1.352,-0.453) (1.378,-0.478) (1.405,-0.500) (1.435,-0.519) (1.468,-0.532)};
\addplot[equipot] coordinates {(1.500,-0.482) (1.526,-0.459) (1.550,-0.434) (1.572,-0.406) (1.592,-0.378) (1.610,-0.347) (1.627,-0.316) (1.641,-0.284) (1.654,-0.252) (1.665,-0.218) (1.674,-0.185) (1.682,-0.150) (1.688,-0.116) (1.692,-0.081) (1.695,-0.046) (1.697,-0.011) (1.696,0.024) (1.694,0.059) (1.691,0.094) (1.686,0.129) (1.679,0.163) (1.671,0.197) (1.661,0.231) (1.649,0.264) (1.636,0.297) (1.621,0.328) (1.604,0.359) (1.585,0.389) (1.564,0.417) (1.541,0.444) (1.516,0.468) (1.490,0.474) (1.465,0.449) (1.442,0.422) (1.422,0.394) (1.403,0.364) (1.386,0.333) (1.372,0.301) (1.358,0.269) (1.347,0.235) (1.337,0.202) (1.329,0.168) (1.323,0.133) (1.317,0.098) (1.314,0.063) (1.312,0.028) (1.312,-0.007) (1.313,-0.042) (1.315,-0.077) (1.319,-0.112) (1.325,-0.146) (1.332,-0.181) (1.341,-0.215) (1.351,-0.248) (1.363,-0.281) (1.377,-0.313) (1.392,-0.345) (1.410,-0.375) (1.429,-0.405) (1.451,-0.432) (1.474,-0.458) (1.500,-0.482)};
\addplot[equipot] coordinates {(1.500,-0.395) (1.516,-0.364) (1.531,-0.332) (1.545,-0.300) (1.557,-0.267) (1.568,-0.234) (1.577,-0.200) (1.585,-0.166) (1.591,-0.131) (1.596,-0.097) (1.599,-0.062) (1.601,-0.027) (1.601,0.008) (1.600,0.043) (1.598,0.078) (1.594,0.113) (1.588,0.147) (1.581,0.182) (1.573,0.216) (1.563,0.249) (1.552,0.282) (1.539,0.315) (1.524,0.347) (1.508,0.378) (1.494,0.380) (1.477,0.349) (1.463,0.317) (1.451,0.284) (1.439,0.251) (1.430,0.217) (1.422,0.183) (1.415,0.149) (1.410,0.114) (1.406,0.079) (1.403,0.044) (1.402,0.009) (1.402,-0.026) (1.404,-0.061) (1.407,-0.095) (1.412,-0.130) (1.418,-0.165) (1.425,-0.199) (1.434,-0.233) (1.444,-0.266) (1.456,-0.299) (1.470,-0.332) (1.484,-0.363) (1.500,-0.395)};
\addplot[lineaE, flechaen=0.375] coordinates {(1.500,0.112) (1.535,0.115) (1.571,0.120) (1.606,0.125) (1.641,0.131) (1.676,0.137) (1.711,0.143) (1.746,0.150) (1.780,0.158) (1.815,0.165) (1.850,0.173) (1.884,0.182) (1.919,0.190) (1.953,0.199) (1.988,0.208) (2.022,0.218) (2.056,0.227) (2.090,0.237) (2.125,0.247) (2.159,0.257) (2.193,0.267) (2.227,0.277) (2.261,0.288) (2.295,0.298) (2.329,0.309) (2.362,0.320) (2.396,0.331) (2.430,0.342) (2.464,0.353) (2.498,0.364) (2.531,0.376) (2.565,0.387) (2.599,0.398) (2.632,0.410) (2.666,0.422) (2.699,0.433) (2.733,0.445) (2.766,0.457) (2.800,0.469) (2.833,0.481) (2.867,0.493) (2.900,0.505) (2.934,0.517) (2.967,0.529) (3.000,0.542) (3.034,0.554) (3.067,0.567) (3.100,0.579) (3.134,0.592) (3.167,0.604) (3.200,0.617)};
\addplot[lineaE, flechaen=0.625] coordinates {(-3.200,0.617) (-3.167,0.604) (-3.134,0.592) (-3.100,0.579) (-3.067,0.567) (-3.034,0.554) (-3.000,0.542) (-2.967,0.529) (-2.934,0.517) (-2.900,0.505) (-2.867,0.493) (-2.833,0.481) (-2.800,0.469) (-2.766,0.457) (-2.733,0.445) (-2.699,0.433) (-2.666,0.422) (-2.632,0.410) (-2.599,0.398) (-2.565,0.387) (-2.531,0.376) (-2.498,0.364) (-2.464,0.353) (-2.430,0.342) (-2.396,0.331) (-2.362,0.320) (-2.329,0.309) (-2.295,0.298) (-2.261,0.288) (-2.227,0.277) (-2.193,0.267) (-2.159,0.257) (-2.125,0.247) (-2.090,0.237) (-2.056,0.227) (-2.022,0.218) (-1.988,0.208) (-1.953,0.199) (-1.919,0.190) (-1.884,0.182) (-1.850,0.173) (-1.815,0.165) (-1.780,0.158) (-1.746,0.150) (-1.711,0.143) (-1.676,0.137) (-1.641,0.131) (-1.606,0.125) (-1.571,0.120) (-1.535,0.115) (-1.500,0.112)};
\addplot[lineaE, flechaen=0.348] coordinates {(1.500,0.218) (1.534,0.225) (1.568,0.234) (1.601,0.245) (1.635,0.256) (1.667,0.269) (1.700,0.282) (1.732,0.296) (1.764,0.311) (1.795,0.327) (1.826,0.343) (1.857,0.360) (1.888,0.377) (1.918,0.394) (1.948,0.412) (1.978,0.431) (2.008,0.449) (2.037,0.468) (2.066,0.487) (2.096,0.507) (2.125,0.526) (2.154,0.546) (2.182,0.566) (2.211,0.587) (2.239,0.607) (2.268,0.628) (2.296,0.648) (2.324,0.669) (2.352,0.690) (2.380,0.711) (2.408,0.733) (2.436,0.754) (2.463,0.776) (2.491,0.797) (2.518,0.819) (2.546,0.841) (2.573,0.863) (2.600,0.885) (2.627,0.908) (2.654,0.930) (2.681,0.953) (2.708,0.975) (2.734,0.998) (2.761,1.021) (2.788,1.044) (2.814,1.067) (2.840,1.090) (2.867,1.113) (2.893,1.136) (2.919,1.160) (2.945,1.183) (2.971,1.207) (2.997,1.231) (3.022,1.254) (3.048,1.278) (3.074,1.302) (3.099,1.326) (3.124,1.351) (3.150,1.375) (3.175,1.399) (3.200,1.424)};
\addplot[lineaE, flechaen=0.652] coordinates {(-3.200,1.424) (-3.175,1.399) (-3.150,1.375) (-3.124,1.351) (-3.099,1.326) (-3.074,1.302) (-3.048,1.278) (-3.022,1.254) (-2.997,1.231) (-2.971,1.207) (-2.945,1.183) (-2.919,1.160) (-2.893,1.136) (-2.867,1.113) (-2.840,1.090) (-2.814,1.067) (-2.788,1.044) (-2.761,1.021) (-2.734,0.998) (-2.708,0.975) (-2.681,0.953) (-2.654,0.930) (-2.627,0.908) (-2.600,0.885) (-2.573,0.863) (-2.546,0.841) (-2.518,0.819) (-2.491,0.797) (-2.463,0.776) (-2.436,0.754) (-2.408,0.733) (-2.380,0.711) (-2.352,0.690) (-2.324,0.669) (-2.296,0.648) (-2.268,0.628) (-2.239,0.607) (-2.211,0.587) (-2.182,0.566) (-2.154,0.546) (-2.125,0.526) (-2.096,0.507) (-2.066,0.487) (-2.037,0.468) (-2.008,0.449) (-1.978,0.431) (-1.948,0.412) (-1.918,0.394) (-1.888,0.377) (-1.857,0.360) (-1.826,0.343) (-1.795,0.327) (-1.764,0.311) (-1.732,0.296) (-1.700,0.282) (-1.667,0.269) (-1.635,0.256) (-1.601,0.245) (-1.568,0.234) (-1.534,0.225) (-1.500,0.218)};
\addplot[lineaE, flechaen=0.508] coordinates {(1.500,0.325) (1.533,0.336) (1.565,0.352) (1.595,0.369) (1.624,0.389) (1.653,0.410) (1.680,0.432) (1.706,0.456) (1.732,0.480) (1.757,0.505) (1.781,0.531) (1.804,0.557) (1.827,0.584) (1.850,0.611) (1.872,0.639) (1.894,0.666) (1.915,0.694) (1.936,0.723) (1.957,0.751) (1.978,0.780) (1.998,0.809) (2.018,0.838) (2.038,0.867) (2.057,0.896) (2.076,0.926) (2.095,0.956) (2.114,0.986) (2.133,1.016) (2.151,1.046) (2.169,1.076) (2.187,1.106) (2.205,1.137) (2.223,1.167) (2.240,1.198) (2.257,1.229) (2.274,1.260) (2.291,1.291) (2.308,1.322) (2.324,1.353) (2.340,1.385) (2.356,1.416) (2.372,1.448) (2.387,1.479) (2.402,1.511) (2.418,1.543) (2.432,1.575) (2.447,1.607) (2.462,1.639) (2.476,1.672) (2.490,1.704) (2.504,1.736) (2.517,1.769) (2.531,1.802) (2.544,1.834) (2.557,1.867) (2.570,1.900)};
\addplot[lineaE, flechaen=0.492] coordinates {(-2.570,1.900) (-2.557,1.867) (-2.544,1.834) (-2.531,1.802) (-2.517,1.769) (-2.504,1.736) (-2.490,1.704) (-2.476,1.672) (-2.462,1.639) (-2.447,1.607) (-2.432,1.575) (-2.418,1.543) (-2.402,1.511) (-2.387,1.479) (-2.372,1.448) (-2.356,1.416) (-2.340,1.385) (-2.324,1.353) (-2.308,1.322) (-2.291,1.291) (-2.274,1.260) (-2.257,1.229) (-2.240,1.198) (-2.223,1.167) (-2.205,1.137) (-2.187,1.106) (-2.169,1.076) (-2.151,1.046) (-2.133,1.016) (-2.114,0.986) (-2.095,0.956) (-2.076,0.926) (-2.057,0.896) (-2.038,0.867) (-2.018,0.838) (-1.998,0.809) (-1.978,0.780) (-1.957,0.751) (-1.936,0.723) (-1.915,0.694) (-1.894,0.666) (-1.872,0.639) (-1.850,0.611) (-1.827,0.584) (-1.804,0.557) (-1.781,0.531) (-1.757,0.505) (-1.732,0.480) (-1.706,0.456) (-1.680,0.432) (-1.653,0.410) (-1.624,0.389) (-1.595,0.369) (-1.565,0.352) (-1.533,0.336) (-1.500,0.325)};
\addplot[lineaE, flechaen=0.500] coordinates {(1.500,0.433) (1.529,0.451) (1.553,0.477) (1.572,0.507) (1.588,0.538) (1.602,0.571) (1.615,0.604) (1.627,0.638) (1.638,0.671) (1.649,0.705) (1.658,0.739) (1.668,0.774) (1.677,0.808) (1.685,0.842) (1.693,0.877) (1.701,0.912) (1.708,0.946) (1.715,0.981) (1.722,1.016) (1.728,1.051) (1.734,1.086) (1.739,1.121) (1.745,1.156) (1.750,1.191) (1.754,1.227) (1.758,1.262) (1.762,1.297) (1.766,1.333) (1.769,1.368) (1.772,1.403) (1.774,1.439) (1.776,1.474) (1.778,1.510) (1.779,1.545) (1.780,1.581) (1.781,1.616) (1.781,1.652) (1.781,1.687) (1.780,1.723) (1.779,1.758) (1.778,1.794) (1.776,1.829) (1.774,1.865) (1.772,1.900)};
\addplot[lineaE, flechaen=0.500] coordinates {(-1.772,1.900) (-1.774,1.865) (-1.776,1.829) (-1.778,1.794) (-1.779,1.758) (-1.780,1.723) (-1.781,1.687) (-1.781,1.652) (-1.781,1.616) (-1.780,1.581) (-1.779,1.545) (-1.778,1.510) (-1.776,1.474) (-1.774,1.439) (-1.772,1.403) (-1.769,1.368) (-1.766,1.333) (-1.762,1.297) (-1.758,1.262) (-1.754,1.227) (-1.750,1.191) (-1.745,1.156) (-1.739,1.121) (-1.734,1.086) (-1.728,1.051) (-1.722,1.016) (-1.715,0.981) (-1.708,0.946) (-1.701,0.912) (-1.693,0.877) (-1.685,0.842) (-1.677,0.808) (-1.668,0.774) (-1.658,0.739) (-1.649,0.705) (-1.638,0.671) (-1.627,0.638) (-1.615,0.604) (-1.602,0.571) (-1.588,0.538) (-1.572,0.507) (-1.553,0.477) (-1.529,0.451) (-1.500,0.433)};
\addplot[lineaE, flechaen=0.500] coordinates {(1.472,0.500) (1.461,0.534) (1.451,0.568) (1.442,0.602) (1.433,0.637) (1.425,0.672) (1.416,0.706) (1.408,0.741) (1.399,0.775) (1.391,0.810) (1.382,0.845) (1.372,0.879) (1.363,0.913) (1.353,0.948) (1.343,0.982) (1.332,1.016) (1.322,1.050) (1.310,1.084) (1.299,1.117) (1.287,1.151) (1.274,1.184) (1.261,1.218) (1.248,1.251) (1.234,1.283) (1.220,1.316) (1.205,1.349) (1.189,1.381) (1.174,1.413) (1.157,1.444) (1.140,1.476) (1.123,1.507) (1.105,1.538) (1.086,1.568) (1.067,1.598) (1.047,1.628) (1.027,1.657) (1.006,1.686) (0.985,1.715) (0.963,1.743) (0.940,1.770) (0.917,1.797) (0.893,1.824) (0.869,1.850) (0.844,1.875) (0.818,1.900)};
\addplot[lineaE, flechaen=0.500] coordinates {(-0.818,1.900) (-0.844,1.875) (-0.869,1.850) (-0.893,1.824) (-0.917,1.797) (-0.940,1.770) (-0.963,1.743) (-0.985,1.715) (-1.006,1.686) (-1.027,1.657) (-1.047,1.628) (-1.067,1.598) (-1.086,1.568) (-1.105,1.538) (-1.123,1.507) (-1.140,1.476) (-1.157,1.444) (-1.174,1.413) (-1.189,1.381) (-1.205,1.349) (-1.220,1.316) (-1.234,1.283) (-1.248,1.251) (-1.261,1.218) (-1.274,1.184) (-1.287,1.151) (-1.299,1.117) (-1.310,1.084) (-1.322,1.050) (-1.332,1.016) (-1.343,0.982) (-1.353,0.948) (-1.363,0.913) (-1.372,0.879) (-1.382,0.845) (-1.391,0.810) (-1.399,0.775) (-1.408,0.741) (-1.416,0.706) (-1.425,0.672) (-1.433,0.637) (-1.442,0.602) (-1.451,0.568) (-1.461,0.534) (-1.472,0.500)};
\addplot[lineaE, flechaen=0.686] coordinates {(1.196,0.500) (1.173,0.525) (1.147,0.549) (1.121,0.572) (1.094,0.596) (1.068,0.619) (1.042,0.642) (1.015,0.666) (0.989,0.689) (0.962,0.712) (0.935,0.734) (0.908,0.757) (0.880,0.779) (0.853,0.800) (0.825,0.822) (0.796,0.843) (0.768,0.863) (0.739,0.883) (0.709,0.903) (0.680,0.922) (0.650,0.940) (0.619,0.958) (0.589,0.975) (0.558,0.991) (0.526,1.007) (0.494,1.022) (0.462,1.036) (0.429,1.050) (0.396,1.062) (0.363,1.073) (0.330,1.084) (0.296,1.094) (0.262,1.102) (0.227,1.110) (0.193,1.117) (0.158,1.122) (0.123,1.127) (0.088,1.130) (0.053,1.132) (0.018,1.133) (-0.018,1.133) (-0.053,1.132) (-0.088,1.130) (-0.123,1.127) (-0.158,1.122) (-0.193,1.117) (-0.227,1.110) (-0.262,1.102) (-0.296,1.094) (-0.330,1.084) (-0.363,1.073) (-0.396,1.062) (-0.429,1.050) (-0.462,1.036) (-0.494,1.022) (-0.526,1.007) (-0.558,0.991) (-0.589,0.975) (-0.619,0.958) (-0.650,0.940) (-0.680,0.922) (-0.709,0.903) (-0.739,0.883) (-0.768,0.863) (-0.796,0.843) (-0.825,0.822) (-0.853,0.800) (-0.880,0.779) (-0.908,0.757) (-0.935,0.734) (-0.962,0.712) (-0.989,0.689) (-1.015,0.666) (-1.042,0.642) (-1.068,0.619) (-1.094,0.596) (-1.121,0.572) (-1.147,0.549) (-1.173,0.525) (-1.196,0.500)};
\addplot[lineaE, flechaen=0.336] coordinates {(1.500,0.131) (1.465,0.135) (1.430,0.141) (1.396,0.147) (1.361,0.153) (1.326,0.160) (1.292,0.167) (1.258,0.175) (1.223,0.183) (1.189,0.191) (1.155,0.200) (1.121,0.208) (1.087,0.217) (1.053,0.227) (1.019,0.236) (0.985,0.245) (0.951,0.255) (0.917,0.264) (0.883,0.273) (0.849,0.283) (0.815,0.292) (0.781,0.301) (0.747,0.310) (0.713,0.319) (0.679,0.327) (0.644,0.335) (0.610,0.344) (0.576,0.351) (0.541,0.359) (0.507,0.366) (0.472,0.373) (0.438,0.379) (0.403,0.385) (0.368,0.391) (0.333,0.396) (0.299,0.401) (0.264,0.405) (0.229,0.409) (0.193,0.412) (0.158,0.414) (0.123,0.417) (0.088,0.418) (0.053,0.419) (0.018,0.420) (-0.018,0.420) (-0.053,0.419) (-0.088,0.418) (-0.123,0.417) (-0.158,0.414) (-0.193,0.412) (-0.229,0.409) (-0.264,0.405) (-0.299,0.401) (-0.333,0.396) (-0.368,0.391) (-0.403,0.385) (-0.438,0.379) (-0.472,0.373) (-0.507,0.366) (-0.541,0.359) (-0.576,0.351) (-0.610,0.344) (-0.644,0.335) (-0.679,0.327) (-0.713,0.319) (-0.747,0.310) (-0.781,0.301) (-0.815,0.292) (-0.849,0.283) (-0.883,0.273) (-0.917,0.264) (-0.951,0.255) (-0.985,0.245) (-1.019,0.236) (-1.053,0.227) (-1.087,0.217) (-1.121,0.208) (-1.155,0.200) (-1.189,0.191) (-1.223,0.183) (-1.258,0.175) (-1.292,0.167) (-1.326,0.160) (-1.361,0.153) (-1.396,0.147) (-1.430,0.141) (-1.465,0.135) (-1.500,0.131)};
\addplot[lineaE, flechaen=0.665] coordinates {(1.500,0.101) (1.465,0.104) (1.430,0.108) (1.395,0.113) (1.360,0.117) (1.325,0.123) (1.290,0.128) (1.255,0.134) (1.220,0.140) (1.186,0.147) (1.151,0.153) (1.116,0.160) (1.082,0.167) (1.047,0.174) (1.012,0.181) (0.978,0.188) (0.943,0.195) (0.909,0.203) (0.874,0.210) (0.840,0.217) (0.805,0.224) (0.770,0.231) (0.736,0.238) (0.701,0.244) (0.666,0.251) (0.632,0.257) (0.597,0.263) (0.562,0.269) (0.527,0.275) (0.492,0.280) (0.457,0.286) (0.422,0.290) (0.387,0.295) (0.352,0.299) (0.317,0.303) (0.282,0.307) (0.247,0.310) (0.212,0.312) (0.176,0.315) (0.141,0.317) (0.106,0.318) (0.071,0.319) (0.035,0.320) (0.000,0.320) (-0.035,0.320) (-0.071,0.319) (-0.106,0.318) (-0.141,0.317) (-0.176,0.315) (-0.212,0.312) (-0.247,0.310) (-0.282,0.307) (-0.317,0.303) (-0.352,0.299) (-0.387,0.295) (-0.422,0.290) (-0.457,0.286) (-0.492,0.280) (-0.527,0.275) (-0.562,0.269) (-0.597,0.263) (-0.632,0.257) (-0.666,0.251) (-0.701,0.244) (-0.736,0.238) (-0.770,0.231) (-0.805,0.224) (-0.840,0.217) (-0.874,0.210) (-0.909,0.203) (-0.943,0.195) (-0.978,0.188) (-1.012,0.181) (-1.047,0.174) (-1.082,0.167) (-1.116,0.160) (-1.151,0.153) (-1.186,0.147) (-1.220,0.140) (-1.255,0.134) (-1.290,0.128) (-1.325,0.123) (-1.360,0.117) (-1.395,0.113) (-1.430,0.108) (-1.465,0.104) (-1.500,0.101)};
\addplot[lineaE, flechaen=0.334] coordinates {(1.500,0.070) (1.465,0.072) (1.430,0.075) (1.395,0.078) (1.360,0.081) (1.325,0.085) (1.290,0.088) (1.255,0.092) (1.221,0.097) (1.186,0.101) (1.151,0.105) (1.116,0.110) (1.082,0.115) (1.047,0.120) (1.012,0.124) (0.977,0.129) (0.943,0.134) (0.908,0.139) (0.873,0.144) (0.838,0.149) (0.804,0.154) (0.769,0.159) (0.734,0.163) (0.699,0.168) (0.664,0.172) (0.630,0.177) (0.595,0.181) (0.560,0.185) (0.525,0.189) (0.490,0.193) (0.455,0.196) (0.420,0.200) (0.385,0.203) (0.350,0.206) (0.315,0.208) (0.280,0.211) (0.245,0.213) (0.210,0.215) (0.175,0.216) (0.140,0.218) (0.105,0.219) (0.070,0.219) (0.035,0.220) (-0.000,0.220) (-0.035,0.220) (-0.070,0.219) (-0.105,0.219) (-0.140,0.218) (-0.175,0.216) (-0.210,0.215) (-0.245,0.213) (-0.280,0.211) (-0.315,0.208) (-0.350,0.206) (-0.385,0.203) (-0.420,0.200) (-0.455,0.196) (-0.490,0.193) (-0.525,0.189) (-0.560,0.185) (-0.595,0.181) (-0.630,0.177) (-0.664,0.172) (-0.699,0.168) (-0.734,0.163) (-0.769,0.159) (-0.804,0.154) (-0.838,0.149) (-0.873,0.144) (-0.908,0.139) (-0.943,0.134) (-0.977,0.129) (-1.012,0.124) (-1.047,0.120) (-1.082,0.115) (-1.116,0.110) (-1.151,0.105) (-1.186,0.101) (-1.221,0.097) (-1.255,0.092) (-1.290,0.088) (-1.325,0.085) (-1.360,0.081) (-1.395,0.078) (-1.430,0.075) (-1.465,0.072) (-1.500,0.070)};
\addplot[lineaE, flechaen=0.666] coordinates {(1.500,0.038) (1.465,0.039) (1.429,0.041) (1.394,0.043) (1.359,0.044) (1.323,0.046) (1.288,0.048) (1.253,0.051) (1.218,0.053) (1.182,0.055) (1.147,0.058) (1.112,0.060) (1.077,0.063) (1.041,0.066) (1.006,0.068) (0.971,0.071) (0.935,0.074) (0.900,0.076) (0.865,0.079) (0.830,0.082) (0.794,0.084) (0.759,0.087) (0.724,0.090) (0.689,0.092) (0.653,0.095) (0.618,0.097) (0.583,0.099) (0.548,0.102) (0.512,0.104) (0.477,0.106) (0.442,0.108) (0.406,0.110) (0.371,0.111) (0.336,0.113) (0.300,0.114) (0.265,0.115) (0.230,0.117) (0.194,0.117) (0.159,0.118) (0.124,0.119) (0.088,0.119) (0.053,0.120) (0.018,0.120) (-0.018,0.120) (-0.053,0.120) (-0.088,0.119) (-0.124,0.119) (-0.159,0.118) (-0.194,0.117) (-0.230,0.117) (-0.265,0.115) (-0.300,0.114) (-0.336,0.113) (-0.371,0.111) (-0.406,0.110) (-0.442,0.108) (-0.477,0.106) (-0.512,0.104) (-0.548,0.102) (-0.583,0.099) (-0.618,0.097) (-0.653,0.095) (-0.689,0.092) (-0.724,0.090) (-0.759,0.087) (-0.794,0.084) (-0.830,0.082) (-0.865,0.079) (-0.900,0.076) (-0.935,0.074) (-0.971,0.071) (-1.006,0.068) (-1.041,0.066) (-1.077,0.063) (-1.112,0.060) (-1.147,0.058) (-1.182,0.055) (-1.218,0.053) (-1.253,0.051) (-1.288,0.048) (-1.323,0.046) (-1.359,0.044) (-1.394,0.043) (-1.429,0.041) (-1.465,0.039) (-1.500,0.038)};
\addplot[lineaE, flechaen=0.375] coordinates {(1.500,-0.112) (1.535,-0.115) (1.571,-0.120) (1.606,-0.125) (1.641,-0.131) (1.676,-0.137) (1.711,-0.143) (1.746,-0.150) (1.780,-0.158) (1.815,-0.165) (1.850,-0.173) (1.884,-0.182) (1.919,-0.190) (1.953,-0.199) (1.988,-0.208) (2.022,-0.218) (2.056,-0.227) (2.090,-0.237) (2.125,-0.247) (2.159,-0.257) (2.193,-0.267) (2.227,-0.277) (2.261,-0.288) (2.295,-0.298) (2.329,-0.309) (2.362,-0.320) (2.396,-0.331) (2.430,-0.342) (2.464,-0.353) (2.498,-0.364) (2.531,-0.376) (2.565,-0.387) (2.599,-0.398) (2.632,-0.410) (2.666,-0.422) (2.699,-0.433) (2.733,-0.445) (2.766,-0.457) (2.800,-0.469) (2.833,-0.481) (2.867,-0.493) (2.900,-0.505) (2.934,-0.517) (2.967,-0.529) (3.000,-0.542) (3.034,-0.554) (3.067,-0.567) (3.100,-0.579) (3.134,-0.592) (3.167,-0.604) (3.200,-0.617)};
\addplot[lineaE, flechaen=0.625] coordinates {(-3.200,-0.617) (-3.167,-0.604) (-3.134,-0.592) (-3.100,-0.579) (-3.067,-0.567) (-3.034,-0.554) (-3.000,-0.542) (-2.967,-0.529) (-2.934,-0.517) (-2.900,-0.505) (-2.867,-0.493) (-2.833,-0.481) (-2.800,-0.469) (-2.766,-0.457) (-2.733,-0.445) (-2.699,-0.433) (-2.666,-0.422) (-2.632,-0.410) (-2.599,-0.398) (-2.565,-0.387) (-2.531,-0.376) (-2.498,-0.364) (-2.464,-0.353) (-2.430,-0.342) (-2.396,-0.331) (-2.362,-0.320) (-2.329,-0.309) (-2.295,-0.298) (-2.261,-0.288) (-2.227,-0.277) (-2.193,-0.267) (-2.159,-0.257) (-2.125,-0.247) (-2.090,-0.237) (-2.056,-0.227) (-2.022,-0.218) (-1.988,-0.208) (-1.953,-0.199) (-1.919,-0.190) (-1.884,-0.182) (-1.850,-0.173) (-1.815,-0.165) (-1.780,-0.158) (-1.746,-0.150) (-1.711,-0.143) (-1.676,-0.137) (-1.641,-0.131) (-1.606,-0.125) (-1.571,-0.120) (-1.535,-0.115) (-1.500,-0.112)};
\addplot[lineaE, flechaen=0.348] coordinates {(1.500,-0.218) (1.534,-0.225) (1.568,-0.234) (1.601,-0.245) (1.635,-0.256) (1.667,-0.269) (1.700,-0.282) (1.732,-0.296) (1.764,-0.311) (1.795,-0.327) (1.826,-0.343) (1.857,-0.360) (1.888,-0.377) (1.918,-0.394) (1.948,-0.412) (1.978,-0.431) (2.008,-0.449) (2.037,-0.468) (2.066,-0.487) (2.096,-0.507) (2.125,-0.526) (2.154,-0.546) (2.182,-0.566) (2.211,-0.587) (2.239,-0.607) (2.268,-0.628) (2.296,-0.648) (2.324,-0.669) (2.352,-0.690) (2.380,-0.711) (2.408,-0.733) (2.436,-0.754) (2.463,-0.776) (2.491,-0.797) (2.518,-0.819) (2.546,-0.841) (2.573,-0.863) (2.600,-0.885) (2.627,-0.908) (2.654,-0.930) (2.681,-0.953) (2.708,-0.975) (2.734,-0.998) (2.761,-1.021) (2.788,-1.044) (2.814,-1.067) (2.840,-1.090) (2.867,-1.113) (2.893,-1.136) (2.919,-1.160) (2.945,-1.183) (2.971,-1.207) (2.997,-1.231) (3.022,-1.254) (3.048,-1.278) (3.074,-1.302) (3.099,-1.326) (3.124,-1.351) (3.150,-1.375) (3.175,-1.399) (3.200,-1.424)};
\addplot[lineaE, flechaen=0.652] coordinates {(-3.200,-1.424) (-3.175,-1.399) (-3.150,-1.375) (-3.124,-1.351) (-3.099,-1.326) (-3.074,-1.302) (-3.048,-1.278) (-3.022,-1.254) (-2.997,-1.231) (-2.971,-1.207) (-2.945,-1.183) (-2.919,-1.160) (-2.893,-1.136) (-2.867,-1.113) (-2.840,-1.090) (-2.814,-1.067) (-2.788,-1.044) (-2.761,-1.021) (-2.734,-0.998) (-2.708,-0.975) (-2.681,-0.953) (-2.654,-0.930) (-2.627,-0.908) (-2.600,-0.885) (-2.573,-0.863) (-2.546,-0.841) (-2.518,-0.819) (-2.491,-0.797) (-2.463,-0.776) (-2.436,-0.754) (-2.408,-0.733) (-2.380,-0.711) (-2.352,-0.690) (-2.324,-0.669) (-2.296,-0.648) (-2.268,-0.628) (-2.239,-0.607) (-2.211,-0.587) (-2.182,-0.566) (-2.154,-0.546) (-2.125,-0.526) (-2.096,-0.507) (-2.066,-0.487) (-2.037,-0.468) (-2.008,-0.449) (-1.978,-0.431) (-1.948,-0.412) (-1.918,-0.394) (-1.888,-0.377) (-1.857,-0.360) (-1.826,-0.343) (-1.795,-0.327) (-1.764,-0.311) (-1.732,-0.296) (-1.700,-0.282) (-1.667,-0.269) (-1.635,-0.256) (-1.601,-0.245) (-1.568,-0.234) (-1.534,-0.225) (-1.500,-0.218)};
\addplot[lineaE, flechaen=0.508] coordinates {(1.500,-0.325) (1.533,-0.336) (1.565,-0.352) (1.595,-0.369) (1.624,-0.389) (1.653,-0.410) (1.680,-0.432) (1.706,-0.456) (1.732,-0.480) (1.757,-0.505) (1.781,-0.531) (1.804,-0.557) (1.827,-0.584) (1.850,-0.611) (1.872,-0.639) (1.894,-0.666) (1.915,-0.694) (1.936,-0.723) (1.957,-0.751) (1.978,-0.780) (1.998,-0.809) (2.018,-0.838) (2.038,-0.867) (2.057,-0.896) (2.076,-0.926) (2.095,-0.956) (2.114,-0.986) (2.133,-1.016) (2.151,-1.046) (2.169,-1.076) (2.187,-1.106) (2.205,-1.137) (2.223,-1.167) (2.240,-1.198) (2.257,-1.229) (2.274,-1.260) (2.291,-1.291) (2.308,-1.322) (2.324,-1.353) (2.340,-1.385) (2.356,-1.416) (2.372,-1.448) (2.387,-1.479) (2.402,-1.511) (2.418,-1.543) (2.432,-1.575) (2.447,-1.607) (2.462,-1.639) (2.476,-1.672) (2.490,-1.704) (2.504,-1.736) (2.517,-1.769) (2.531,-1.802) (2.544,-1.834) (2.557,-1.867) (2.570,-1.900)};
\addplot[lineaE, flechaen=0.492] coordinates {(-2.570,-1.900) (-2.557,-1.867) (-2.544,-1.834) (-2.531,-1.802) (-2.517,-1.769) (-2.504,-1.736) (-2.490,-1.704) (-2.476,-1.672) (-2.462,-1.639) (-2.447,-1.607) (-2.432,-1.575) (-2.418,-1.543) (-2.402,-1.511) (-2.387,-1.479) (-2.372,-1.448) (-2.356,-1.416) (-2.340,-1.385) (-2.324,-1.353) (-2.308,-1.322) (-2.291,-1.291) (-2.274,-1.260) (-2.257,-1.229) (-2.240,-1.198) (-2.223,-1.167) (-2.205,-1.137) (-2.187,-1.106) (-2.169,-1.076) (-2.151,-1.046) (-2.133,-1.016) (-2.114,-0.986) (-2.095,-0.956) (-2.076,-0.926) (-2.057,-0.896) (-2.038,-0.867) (-2.018,-0.838) (-1.998,-0.809) (-1.978,-0.780) (-1.957,-0.751) (-1.936,-0.723) (-1.915,-0.694) (-1.894,-0.666) (-1.872,-0.639) (-1.850,-0.611) (-1.827,-0.584) (-1.804,-0.557) (-1.781,-0.531) (-1.757,-0.505) (-1.732,-0.480) (-1.706,-0.456) (-1.680,-0.432) (-1.653,-0.410) (-1.624,-0.389) (-1.595,-0.369) (-1.565,-0.352) (-1.533,-0.336) (-1.500,-0.325)};
\addplot[lineaE, flechaen=0.500] coordinates {(1.500,-0.433) (1.529,-0.451) (1.553,-0.477) (1.572,-0.507) (1.588,-0.538) (1.602,-0.571) (1.615,-0.604) (1.627,-0.638) (1.638,-0.671) (1.649,-0.705) (1.658,-0.739) (1.668,-0.774) (1.677,-0.808) (1.685,-0.842) (1.693,-0.877) (1.701,-0.912) (1.708,-0.946) (1.715,-0.981) (1.722,-1.016) (1.728,-1.051) (1.734,-1.086) (1.739,-1.121) (1.745,-1.156) (1.750,-1.191) (1.754,-1.227) (1.758,-1.262) (1.762,-1.297) (1.766,-1.333) (1.769,-1.368) (1.772,-1.403) (1.774,-1.439) (1.776,-1.474) (1.778,-1.510) (1.779,-1.545) (1.780,-1.581) (1.781,-1.616) (1.781,-1.652) (1.781,-1.687) (1.780,-1.723) (1.779,-1.758) (1.778,-1.794) (1.776,-1.829) (1.774,-1.865) (1.772,-1.900)};
\addplot[lineaE, flechaen=0.500] coordinates {(-1.772,-1.900) (-1.774,-1.865) (-1.776,-1.829) (-1.778,-1.794) (-1.779,-1.758) (-1.780,-1.723) (-1.781,-1.687) (-1.781,-1.652) (-1.781,-1.616) (-1.780,-1.581) (-1.779,-1.545) (-1.778,-1.510) (-1.776,-1.474) (-1.774,-1.439) (-1.772,-1.403) (-1.769,-1.368) (-1.766,-1.333) (-1.762,-1.297) (-1.758,-1.262) (-1.754,-1.227) (-1.750,-1.191) (-1.745,-1.156) (-1.739,-1.121) (-1.734,-1.086) (-1.728,-1.051) (-1.722,-1.016) (-1.715,-0.981) (-1.708,-0.946) (-1.701,-0.912) (-1.693,-0.877) (-1.685,-0.842) (-1.677,-0.808) (-1.668,-0.774) (-1.658,-0.739) (-1.649,-0.705) (-1.638,-0.671) (-1.627,-0.638) (-1.615,-0.604) (-1.602,-0.571) (-1.588,-0.538) (-1.572,-0.507) (-1.553,-0.477) (-1.529,-0.451) (-1.500,-0.433)};
\addplot[lineaE, flechaen=0.500] coordinates {(1.472,-0.500) (1.461,-0.534) (1.451,-0.568) (1.442,-0.602) (1.433,-0.637) (1.425,-0.672) (1.416,-0.706) (1.408,-0.741) (1.399,-0.775) (1.391,-0.810) (1.382,-0.845) (1.372,-0.879) (1.363,-0.913) (1.353,-0.948) (1.343,-0.982) (1.332,-1.016) (1.322,-1.050) (1.310,-1.084) (1.299,-1.117) (1.287,-1.151) (1.274,-1.184) (1.261,-1.218) (1.248,-1.251) (1.234,-1.283) (1.220,-1.316) (1.205,-1.349) (1.189,-1.381) (1.174,-1.413) (1.157,-1.444) (1.140,-1.476) (1.123,-1.507) (1.105,-1.538) (1.086,-1.568) (1.067,-1.598) (1.047,-1.628) (1.027,-1.657) (1.006,-1.686) (0.985,-1.715) (0.963,-1.743) (0.940,-1.770) (0.917,-1.797) (0.893,-1.824) (0.869,-1.850) (0.844,-1.875) (0.818,-1.900)};
\addplot[lineaE, flechaen=0.500] coordinates {(-0.818,-1.900) (-0.844,-1.875) (-0.869,-1.850) (-0.893,-1.824) (-0.917,-1.797) (-0.940,-1.770) (-0.963,-1.743) (-0.985,-1.715) (-1.006,-1.686) (-1.027,-1.657) (-1.047,-1.628) (-1.067,-1.598) (-1.086,-1.568) (-1.105,-1.538) (-1.123,-1.507) (-1.140,-1.476) (-1.157,-1.444) (-1.174,-1.413) (-1.189,-1.381) (-1.205,-1.349) (-1.220,-1.316) (-1.234,-1.283) (-1.248,-1.251) (-1.261,-1.218) (-1.274,-1.184) (-1.287,-1.151) (-1.299,-1.117) (-1.310,-1.084) (-1.322,-1.050) (-1.332,-1.016) (-1.343,-0.982) (-1.353,-0.948) (-1.363,-0.913) (-1.372,-0.879) (-1.382,-0.845) (-1.391,-0.810) (-1.399,-0.775) (-1.408,-0.741) (-1.416,-0.706) (-1.425,-0.672) (-1.433,-0.637) (-1.442,-0.602) (-1.451,-0.568) (-1.461,-0.534) (-1.472,-0.500)};
\addplot[lineaE, flechaen=0.686] coordinates {(1.196,-0.500) (1.173,-0.525) (1.147,-0.549) (1.121,-0.572) (1.094,-0.596) (1.068,-0.619) (1.042,-0.642) (1.015,-0.666) (0.989,-0.689) (0.962,-0.712) (0.935,-0.734) (0.908,-0.757) (0.880,-0.779) (0.853,-0.800) (0.825,-0.822) (0.796,-0.843) (0.768,-0.863) (0.739,-0.883) (0.709,-0.903) (0.680,-0.922) (0.650,-0.940) (0.619,-0.958) (0.589,-0.975) (0.558,-0.991) (0.526,-1.007) (0.494,-1.022) (0.462,-1.036) (0.429,-1.050) (0.396,-1.062) (0.363,-1.073) (0.330,-1.084) (0.296,-1.094) (0.262,-1.102) (0.227,-1.110) (0.193,-1.117) (0.158,-1.122) (0.123,-1.127) (0.088,-1.130) (0.053,-1.132) (0.018,-1.133) (-0.018,-1.133) (-0.053,-1.132) (-0.088,-1.130) (-0.123,-1.127) (-0.158,-1.122) (-0.193,-1.117) (-0.227,-1.110) (-0.262,-1.102) (-0.296,-1.094) (-0.330,-1.084) (-0.363,-1.073) (-0.396,-1.062) (-0.429,-1.050) (-0.462,-1.036) (-0.494,-1.022) (-0.526,-1.007) (-0.558,-0.991) (-0.589,-0.975) (-0.619,-0.958) (-0.650,-0.940) (-0.680,-0.922) (-0.709,-0.903) (-0.739,-0.883) (-0.768,-0.863) (-0.796,-0.843) (-0.825,-0.822) (-0.853,-0.800) (-0.880,-0.779) (-0.908,-0.757) (-0.935,-0.734) (-0.962,-0.712) (-0.989,-0.689) (-1.015,-0.666) (-1.042,-0.642) (-1.068,-0.619) (-1.094,-0.596) (-1.121,-0.572) (-1.147,-0.549) (-1.173,-0.525) (-1.196,-0.500)};
\addplot[lineaE, flechaen=0.336] coordinates {(1.500,-0.131) (1.465,-0.135) (1.430,-0.141) (1.396,-0.147) (1.361,-0.153) (1.326,-0.160) (1.292,-0.167) (1.258,-0.175) (1.223,-0.183) (1.189,-0.191) (1.155,-0.200) (1.121,-0.208) (1.087,-0.217) (1.053,-0.227) (1.019,-0.236) (0.985,-0.245) (0.951,-0.255) (0.917,-0.264) (0.883,-0.273) (0.849,-0.283) (0.815,-0.292) (0.781,-0.301) (0.747,-0.310) (0.713,-0.319) (0.679,-0.327) (0.644,-0.335) (0.610,-0.344) (0.576,-0.351) (0.541,-0.359) (0.507,-0.366) (0.472,-0.373) (0.438,-0.379) (0.403,-0.385) (0.368,-0.391) (0.333,-0.396) (0.299,-0.401) (0.264,-0.405) (0.229,-0.409) (0.193,-0.412) (0.158,-0.414) (0.123,-0.417) (0.088,-0.418) (0.053,-0.419) (0.018,-0.420) (-0.018,-0.420) (-0.053,-0.419) (-0.088,-0.418) (-0.123,-0.417) (-0.158,-0.414) (-0.193,-0.412) (-0.229,-0.409) (-0.264,-0.405) (-0.299,-0.401) (-0.333,-0.396) (-0.368,-0.391) (-0.403,-0.385) (-0.438,-0.379) (-0.472,-0.373) (-0.507,-0.366) (-0.541,-0.359) (-0.576,-0.351) (-0.610,-0.344) (-0.644,-0.335) (-0.679,-0.327) (-0.713,-0.319) (-0.747,-0.310) (-0.781,-0.301) (-0.815,-0.292) (-0.849,-0.283) (-0.883,-0.273) (-0.917,-0.264) (-0.951,-0.255) (-0.985,-0.245) (-1.019,-0.236) (-1.053,-0.227) (-1.087,-0.217) (-1.121,-0.208) (-1.155,-0.200) (-1.189,-0.191) (-1.223,-0.183) (-1.258,-0.175) (-1.292,-0.167) (-1.326,-0.160) (-1.361,-0.153) (-1.396,-0.147) (-1.430,-0.141) (-1.465,-0.135) (-1.500,-0.131)};
\addplot[lineaE, flechaen=0.665] coordinates {(1.500,-0.101) (1.465,-0.104) (1.430,-0.108) (1.395,-0.113) (1.360,-0.117) (1.325,-0.123) (1.290,-0.128) (1.255,-0.134) (1.220,-0.140) (1.186,-0.147) (1.151,-0.153) (1.116,-0.160) (1.082,-0.167) (1.047,-0.174) (1.012,-0.181) (0.978,-0.188) (0.943,-0.195) (0.909,-0.203) (0.874,-0.210) (0.840,-0.217) (0.805,-0.224) (0.770,-0.231) (0.736,-0.238) (0.701,-0.244) (0.666,-0.251) (0.632,-0.257) (0.597,-0.263) (0.562,-0.269) (0.527,-0.275) (0.492,-0.280) (0.457,-0.286) (0.422,-0.290) (0.387,-0.295) (0.352,-0.299) (0.317,-0.303) (0.282,-0.307) (0.247,-0.310) (0.212,-0.312) (0.176,-0.315) (0.141,-0.317) (0.106,-0.318) (0.071,-0.319) (0.035,-0.320) (0.000,-0.320) (-0.035,-0.320) (-0.071,-0.319) (-0.106,-0.318) (-0.141,-0.317) (-0.176,-0.315) (-0.212,-0.312) (-0.247,-0.310) (-0.282,-0.307) (-0.317,-0.303) (-0.352,-0.299) (-0.387,-0.295) (-0.422,-0.290) (-0.457,-0.286) (-0.492,-0.280) (-0.527,-0.275) (-0.562,-0.269) (-0.597,-0.263) (-0.632,-0.257) (-0.666,-0.251) (-0.701,-0.244) (-0.736,-0.238) (-0.770,-0.231) (-0.805,-0.224) (-0.840,-0.217) (-0.874,-0.210) (-0.909,-0.203) (-0.943,-0.195) (-0.978,-0.188) (-1.012,-0.181) (-1.047,-0.174) (-1.082,-0.167) (-1.116,-0.160) (-1.151,-0.153) (-1.186,-0.147) (-1.220,-0.140) (-1.255,-0.134) (-1.290,-0.128) (-1.325,-0.123) (-1.360,-0.117) (-1.395,-0.113) (-1.430,-0.108) (-1.465,-0.104) (-1.500,-0.101)};
\addplot[lineaE, flechaen=0.334] coordinates {(1.500,-0.070) (1.465,-0.072) (1.430,-0.075) (1.395,-0.078) (1.360,-0.081) (1.325,-0.085) (1.290,-0.088) (1.255,-0.092) (1.221,-0.097) (1.186,-0.101) (1.151,-0.105) (1.116,-0.110) (1.082,-0.115) (1.047,-0.120) (1.012,-0.124) (0.977,-0.129) (0.943,-0.134) (0.908,-0.139) (0.873,-0.144) (0.838,-0.149) (0.804,-0.154) (0.769,-0.159) (0.734,-0.163) (0.699,-0.168) (0.664,-0.172) (0.630,-0.177) (0.595,-0.181) (0.560,-0.185) (0.525,-0.189) (0.490,-0.193) (0.455,-0.196) (0.420,-0.200) (0.385,-0.203) (0.350,-0.206) (0.315,-0.208) (0.280,-0.211) (0.245,-0.213) (0.210,-0.215) (0.175,-0.216) (0.140,-0.218) (0.105,-0.219) (0.070,-0.219) (0.035,-0.220) (-0.000,-0.220) (-0.035,-0.220) (-0.070,-0.219) (-0.105,-0.219) (-0.140,-0.218) (-0.175,-0.216) (-0.210,-0.215) (-0.245,-0.213) (-0.280,-0.211) (-0.315,-0.208) (-0.350,-0.206) (-0.385,-0.203) (-0.420,-0.200) (-0.455,-0.196) (-0.490,-0.193) (-0.525,-0.189) (-0.560,-0.185) (-0.595,-0.181) (-0.630,-0.177) (-0.664,-0.172) (-0.699,-0.168) (-0.734,-0.163) (-0.769,-0.159) (-0.804,-0.154) (-0.838,-0.149) (-0.873,-0.144) (-0.908,-0.139) (-0.943,-0.134) (-0.977,-0.129) (-1.012,-0.124) (-1.047,-0.120) (-1.082,-0.115) (-1.116,-0.110) (-1.151,-0.105) (-1.186,-0.101) (-1.221,-0.097) (-1.255,-0.092) (-1.290,-0.088) (-1.325,-0.085) (-1.360,-0.081) (-1.395,-0.078) (-1.430,-0.075) (-1.465,-0.072) (-1.500,-0.070)};
\addplot[lineaE, flechaen=0.666] coordinates {(1.500,-0.038) (1.465,-0.039) (1.429,-0.041) (1.394,-0.043) (1.359,-0.044) (1.323,-0.046) (1.288,-0.048) (1.253,-0.051) (1.218,-0.053) (1.182,-0.055) (1.147,-0.058) (1.112,-0.060) (1.077,-0.063) (1.041,-0.066) (1.006,-0.068) (0.971,-0.071) (0.935,-0.074) (0.900,-0.076) (0.865,-0.079) (0.830,-0.082) (0.794,-0.084) (0.759,-0.087) (0.724,-0.090) (0.689,-0.092) (0.653,-0.095) (0.618,-0.097) (0.583,-0.099) (0.548,-0.102) (0.512,-0.104) (0.477,-0.106) (0.442,-0.108) (0.406,-0.110) (0.371,-0.111) (0.336,-0.113) (0.300,-0.114) (0.265,-0.115) (0.230,-0.117) (0.194,-0.117) (0.159,-0.118) (0.124,-0.119) (0.088,-0.119) (0.053,-0.120) (0.018,-0.120) (-0.018,-0.120) (-0.053,-0.120) (-0.088,-0.119) (-0.124,-0.119) (-0.159,-0.118) (-0.194,-0.117) (-0.230,-0.117) (-0.265,-0.115) (-0.300,-0.114) (-0.336,-0.113) (-0.371,-0.111) (-0.406,-0.110) (-0.442,-0.108) (-0.477,-0.106) (-0.512,-0.104) (-0.548,-0.102) (-0.583,-0.099) (-0.618,-0.097) (-0.653,-0.095) (-0.689,-0.092) (-0.724,-0.090) (-0.759,-0.087) (-0.794,-0.084) (-0.830,-0.082) (-0.865,-0.079) (-0.900,-0.076) (-0.935,-0.074) (-0.971,-0.071) (-1.006,-0.068) (-1.041,-0.066) (-1.077,-0.063) (-1.112,-0.060) (-1.147,-0.058) (-1.182,-0.055) (-1.218,-0.053) (-1.253,-0.051) (-1.288,-0.048) (-1.323,-0.046) (-1.359,-0.044) (-1.394,-0.043) (-1.429,-0.041) (-1.465,-0.039) (-1.500,-0.038)};
\addplot[lineaE, flechaen=0.382] coordinates {(1.500,0.000) (3.200,0.000)};
\addplot[lineaE, flechaen=0.618] coordinates {(-3.200,0.000) (-1.500,0.000)};
\addplot[lineaE, flechaen=0.667] coordinates {(1.500,0.000) (-1.500,0.000)};

  \draw[texto, line width=0.6pt] (axis cs:-1.5,-0.5) rectangle (axis cs:1.5,0.5);
\end{axis}
\end{tikzpicture}
\caption{$\vE(\vr)$}
\end{subfigure}\hfill
\begin{subfigure}[b]{0.325\textwidth}\centering
\begin{tikzpicture}
\begin{axis}[cil]
  \fill[dielec] (axis cs:-1.5,-0.5) rectangle (axis cs:1.5,0.5);
  %% Generado por scripts/lineas_campo.py -- no editar a mano
\addplot[equipot] coordinates {(-1.500,-0.395) (-1.516,-0.364) (-1.531,-0.332) (-1.545,-0.300) (-1.557,-0.267) (-1.568,-0.234) (-1.577,-0.200) (-1.585,-0.166) (-1.591,-0.131) (-1.596,-0.097) (-1.599,-0.062) (-1.601,-0.027) (-1.601,0.008) (-1.600,0.043) (-1.598,0.078) (-1.594,0.113) (-1.588,0.147) (-1.581,0.182) (-1.573,0.216) (-1.563,0.249) (-1.552,0.282) (-1.539,0.315) (-1.524,0.347) (-1.508,0.378) (-1.494,0.380) (-1.477,0.349) (-1.463,0.317) (-1.451,0.284) (-1.439,0.251) (-1.430,0.217) (-1.422,0.183) (-1.415,0.149) (-1.410,0.114) (-1.406,0.079) (-1.403,0.044) (-1.402,0.009) (-1.402,-0.026) (-1.404,-0.061) (-1.407,-0.095) (-1.412,-0.130) (-1.418,-0.165) (-1.425,-0.199) (-1.434,-0.233) (-1.444,-0.266) (-1.456,-0.299) (-1.470,-0.332) (-1.484,-0.363) (-1.500,-0.395)};
\addplot[equipot] coordinates {(-1.500,-0.482) (-1.526,-0.459) (-1.550,-0.434) (-1.572,-0.406) (-1.592,-0.378) (-1.610,-0.347) (-1.627,-0.316) (-1.641,-0.284) (-1.654,-0.252) (-1.665,-0.218) (-1.674,-0.185) (-1.682,-0.150) (-1.688,-0.116) (-1.692,-0.081) (-1.695,-0.046) (-1.697,-0.011) (-1.696,0.024) (-1.694,0.059) (-1.691,0.094) (-1.686,0.129) (-1.679,0.163) (-1.671,0.197) (-1.661,0.231) (-1.649,0.264) (-1.636,0.297) (-1.621,0.328) (-1.604,0.359) (-1.585,0.389) (-1.564,0.417) (-1.541,0.444) (-1.516,0.468) (-1.490,0.474) (-1.465,0.449) (-1.442,0.422) (-1.422,0.394) (-1.403,0.364) (-1.386,0.333) (-1.372,0.301) (-1.358,0.269) (-1.347,0.235) (-1.337,0.202) (-1.329,0.168) (-1.323,0.133) (-1.317,0.098) (-1.314,0.063) (-1.312,0.028) (-1.312,-0.007) (-1.313,-0.042) (-1.315,-0.077) (-1.319,-0.112) (-1.325,-0.146) (-1.332,-0.181) (-1.341,-0.215) (-1.351,-0.248) (-1.363,-0.281) (-1.377,-0.313) (-1.392,-0.345) (-1.410,-0.375) (-1.429,-0.405) (-1.451,-0.432) (-1.474,-0.458) (-1.500,-0.482)};
\addplot[equipot] coordinates {(-1.548,-0.528) (-1.580,-0.514) (-1.610,-0.495) (-1.637,-0.472) (-1.662,-0.448) (-1.685,-0.421) (-1.707,-0.393) (-1.726,-0.363) (-1.743,-0.333) (-1.759,-0.301) (-1.773,-0.269) (-1.785,-0.236) (-1.796,-0.202) (-1.804,-0.168) (-1.812,-0.133) (-1.817,-0.098) (-1.821,-0.063) (-1.823,-0.028) (-1.823,0.007) (-1.822,0.043) (-1.819,0.078) (-1.815,0.113) (-1.809,0.148) (-1.801,0.182) (-1.791,0.216) (-1.780,0.250) (-1.767,0.282) (-1.753,0.315) (-1.736,0.346) (-1.718,0.376) (-1.698,0.405) (-1.676,0.432) (-1.652,0.458) (-1.626,0.482) (-1.597,0.503) (-1.567,0.520) (-1.534,0.533) (-1.499,0.537) (-1.464,0.531) (-1.432,0.517) (-1.402,0.498) (-1.375,0.475) (-1.350,0.450) (-1.327,0.423) (-1.307,0.395) (-1.288,0.365) (-1.271,0.334) (-1.256,0.302) (-1.243,0.269) (-1.231,0.236) (-1.221,0.202) (-1.213,0.168) (-1.206,0.133) (-1.201,0.098) (-1.197,0.063) (-1.196,0.028) (-1.195,-0.007) (-1.196,-0.043) (-1.199,-0.078) (-1.203,-0.113) (-1.209,-0.148) (-1.216,-0.182) (-1.225,-0.216) (-1.236,-0.250) (-1.248,-0.283) (-1.262,-0.316) (-1.278,-0.347) (-1.295,-0.378) (-1.315,-0.407) (-1.337,-0.435) (-1.360,-0.461) (-1.386,-0.485) (-1.414,-0.506) (-1.445,-0.523) (-1.478,-0.535) (-1.514,-0.536) (-1.548,-0.528)};
\addplot[equipot] coordinates {(-1.548,-0.656) (-1.583,-0.651) (-1.617,-0.642) (-1.650,-0.630) (-1.682,-0.616) (-1.712,-0.599) (-1.742,-0.579) (-1.770,-0.558) (-1.796,-0.535) (-1.821,-0.510) (-1.845,-0.484) (-1.867,-0.457) (-1.887,-0.428) (-1.906,-0.399) (-1.924,-0.368) (-1.940,-0.337) (-1.954,-0.305) (-1.967,-0.272) (-1.978,-0.239) (-1.988,-0.205) (-1.996,-0.171) (-2.003,-0.137) (-2.008,-0.102) (-2.011,-0.067) (-2.014,-0.032) (-2.014,0.003) (-2.013,0.039) (-2.011,0.074) (-2.007,0.108) (-2.001,0.143) (-1.994,0.178) (-1.986,0.212) (-1.976,0.245) (-1.964,0.279) (-1.951,0.311) (-1.937,0.343) (-1.920,0.374) (-1.903,0.405) (-1.883,0.434) (-1.863,0.462) (-1.840,0.490) (-1.817,0.515) (-1.791,0.540) (-1.764,0.562) (-1.736,0.583) (-1.707,0.602) (-1.676,0.619) (-1.643,0.633) (-1.610,0.644) (-1.576,0.652) (-1.541,0.657) (-1.506,0.658) (-1.471,0.656) (-1.436,0.650) (-1.403,0.640) (-1.370,0.628) (-1.338,0.612) (-1.308,0.594) (-1.279,0.574) (-1.252,0.552) (-1.227,0.527) (-1.203,0.501) (-1.181,0.474) (-1.160,0.446) (-1.141,0.416) (-1.124,0.386) (-1.108,0.354) (-1.094,0.322) (-1.081,0.289) (-1.070,0.256) (-1.060,0.222) (-1.052,0.188) (-1.045,0.154) (-1.039,0.119) (-1.035,0.084) (-1.032,0.049) (-1.031,0.014) (-1.031,-0.021) (-1.033,-0.056) (-1.036,-0.091) (-1.040,-0.126) (-1.046,-0.161) (-1.053,-0.195) (-1.062,-0.229) (-1.072,-0.263) (-1.084,-0.296) (-1.097,-0.328) (-1.111,-0.360) (-1.127,-0.392) (-1.145,-0.422) (-1.164,-0.451) (-1.185,-0.480) (-1.208,-0.507) (-1.232,-0.532) (-1.257,-0.556) (-1.285,-0.578) (-1.314,-0.598) (-1.344,-0.616) (-1.376,-0.631) (-1.409,-0.642) (-1.443,-0.651) (-1.478,-0.656) (-1.513,-0.658) (-1.548,-0.656)};
\addplot[equipot] coordinates {(-1.660,-0.922) (-1.695,-0.915) (-1.729,-0.907) (-1.763,-0.898) (-1.796,-0.887) (-1.829,-0.874) (-1.861,-0.860) (-1.893,-0.844) (-1.924,-0.827) (-1.954,-0.809) (-1.983,-0.789) (-2.011,-0.769) (-2.039,-0.747) (-2.065,-0.724) (-2.091,-0.699) (-2.115,-0.674) (-2.139,-0.648) (-2.162,-0.621) (-2.183,-0.593) (-2.204,-0.565) (-2.223,-0.535) (-2.241,-0.505) (-2.258,-0.474) (-2.274,-0.443) (-2.289,-0.411) (-2.303,-0.379) (-2.315,-0.346) (-2.327,-0.313) (-2.337,-0.279) (-2.346,-0.245) (-2.354,-0.211) (-2.360,-0.176) (-2.366,-0.141) (-2.370,-0.106) (-2.373,-0.071) (-2.375,-0.036) (-2.376,-0.001) (-2.375,0.034) (-2.373,0.069) (-2.370,0.104) (-2.366,0.139) (-2.361,0.174) (-2.354,0.209) (-2.346,0.243) (-2.337,0.277) (-2.327,0.311) (-2.316,0.344) (-2.304,0.377) (-2.290,0.409) (-2.275,0.441) (-2.259,0.473) (-2.242,0.503) (-2.224,0.534) (-2.205,0.563) (-2.184,0.592) (-2.163,0.620) (-2.140,0.647) (-2.117,0.673) (-2.092,0.698) (-2.067,0.722) (-2.040,0.745) (-2.013,0.767) (-1.985,0.788) (-1.955,0.808) (-1.925,0.826) (-1.895,0.843) (-1.863,0.859) (-1.831,0.873) (-1.798,0.886) (-1.765,0.897) (-1.731,0.907) (-1.697,0.915) (-1.662,0.921) (-1.627,0.926) (-1.592,0.929) (-1.557,0.930) (-1.522,0.930) (-1.487,0.927) (-1.452,0.923) (-1.417,0.917) (-1.383,0.909) (-1.349,0.900) (-1.316,0.889) (-1.283,0.876) (-1.251,0.861) (-1.219,0.846) (-1.189,0.828) (-1.159,0.809) (-1.130,0.789) (-1.103,0.767) (-1.076,0.744) (-1.050,0.720) (-1.026,0.695) (-1.002,0.669) (-0.980,0.642) (-0.959,0.613) (-0.939,0.584) (-0.920,0.555) (-0.902,0.524) (-0.886,0.493) (-0.870,0.462) (-0.856,0.430) (-0.843,0.397) (-0.831,0.364) (-0.820,0.330) (-0.810,0.297) (-0.801,0.263) (-0.794,0.228) (-0.787,0.194) (-0.782,0.159) (-0.778,0.124) (-0.774,0.089) (-0.772,0.054) (-0.771,0.019) (-0.771,-0.017) (-0.772,-0.052) (-0.774,-0.087) (-0.777,-0.122) (-0.782,-0.157) (-0.787,-0.192) (-0.793,-0.226) (-0.801,-0.261) (-0.809,-0.295) (-0.819,-0.329) (-0.830,-0.362) (-0.842,-0.395) (-0.855,-0.428) (-0.869,-0.460) (-0.885,-0.492) (-0.901,-0.523) (-0.919,-0.553) (-0.938,-0.583) (-0.958,-0.612) (-0.979,-0.640) (-1.001,-0.667) (-1.024,-0.693) (-1.049,-0.719) (-1.075,-0.743) (-1.101,-0.766) (-1.129,-0.788) (-1.158,-0.808) (-1.187,-0.827) (-1.218,-0.845) (-1.249,-0.861) (-1.281,-0.875) (-1.314,-0.888) (-1.347,-0.899) (-1.381,-0.909) (-1.415,-0.917) (-1.450,-0.923) (-1.485,-0.927) (-1.520,-0.929) (-1.555,-0.930) (-1.590,-0.929) (-1.625,-0.926) (-1.660,-0.922)};
\addplot[equipot] coordinates {(-3.200,1.227) (-3.178,1.254) (-3.155,1.281) (-3.131,1.307) (-3.107,1.333) (-3.083,1.358) (-3.058,1.382) (-3.033,1.407) (-3.007,1.430) (-2.981,1.454) (-2.954,1.476) (-2.927,1.499) (-2.899,1.520) (-2.871,1.541) (-2.843,1.562) (-2.814,1.582) (-2.785,1.601) (-2.755,1.620) (-2.725,1.639) (-2.695,1.656) (-2.665,1.673) (-2.634,1.690) (-2.602,1.706) (-2.571,1.721) (-2.539,1.735) (-2.507,1.749) (-2.474,1.762) (-2.441,1.775) (-2.408,1.787) (-2.375,1.798) (-2.342,1.808) (-2.308,1.818) (-2.274,1.827) (-2.240,1.835) (-2.206,1.843) (-2.171,1.850) (-2.137,1.856) (-2.102,1.861) (-2.067,1.866) (-2.033,1.869) (-1.998,1.872) (-1.963,1.875) (-1.928,1.876) (-1.893,1.877) (-1.857,1.877) (-1.822,1.876) (-1.787,1.874) (-1.752,1.872) (-1.718,1.868) (-1.683,1.864) (-1.648,1.859) (-1.613,1.854) (-1.579,1.847) (-1.545,1.840) (-1.511,1.832) (-1.477,1.823) (-1.443,1.813) (-1.409,1.803) (-1.376,1.791) (-1.343,1.779) (-1.311,1.766) (-1.278,1.753) (-1.247,1.738) (-1.215,1.723) (-1.184,1.707) (-1.153,1.690) (-1.122,1.673) (-1.093,1.655) (-1.063,1.636) (-1.034,1.616) (-1.005,1.596) (-0.977,1.575) (-0.950,1.553) (-0.923,1.531) (-0.896,1.508) (-0.870,1.484) (-0.845,1.460) (-0.820,1.435) (-0.796,1.410) (-0.772,1.384) (-0.749,1.358) (-0.727,1.331) (-0.705,1.303) (-0.684,1.275) (-0.663,1.247) (-0.643,1.218) (-0.624,1.189) (-0.605,1.159) (-0.588,1.129) (-0.570,1.098) (-0.553,1.067) (-0.537,1.036) (-0.522,1.005) (-0.507,0.973) (-0.493,0.941) (-0.479,0.909) (-0.466,0.876) (-0.454,0.843) (-0.442,0.810) (-0.431,0.777) (-0.420,0.744) (-0.410,0.710) (-0.401,0.676) (-0.392,0.642) (-0.383,0.608) (-0.375,0.574) (-0.368,0.540) (-0.361,0.506) (-0.355,0.471) (-0.349,0.436) (-0.344,0.402) (-0.339,0.367) (-0.334,0.332) (-0.330,0.297) (-0.327,0.263) (-0.324,0.228) (-0.321,0.193) (-0.319,0.158) (-0.317,0.123) (-0.316,0.088) (-0.315,0.053) (-0.315,0.018) (-0.315,-0.018) (-0.315,-0.053) (-0.316,-0.088) (-0.317,-0.123) (-0.319,-0.158) (-0.321,-0.193) (-0.324,-0.228) (-0.327,-0.263) (-0.330,-0.297) (-0.334,-0.332) (-0.339,-0.367) (-0.344,-0.402) (-0.349,-0.436) (-0.355,-0.471) (-0.361,-0.506) (-0.368,-0.540) (-0.375,-0.574) (-0.383,-0.608) (-0.392,-0.642) (-0.401,-0.676) (-0.410,-0.710) (-0.420,-0.744) (-0.431,-0.777) (-0.442,-0.810) (-0.454,-0.843) (-0.466,-0.876) (-0.479,-0.909) (-0.493,-0.941) (-0.507,-0.973) (-0.522,-1.005) (-0.537,-1.036) (-0.553,-1.067) (-0.570,-1.098) (-0.588,-1.129) (-0.605,-1.159) (-0.624,-1.189) (-0.643,-1.218) (-0.663,-1.247) (-0.684,-1.275) (-0.705,-1.303) (-0.727,-1.331) (-0.749,-1.358) (-0.772,-1.384) (-0.796,-1.410) (-0.820,-1.435) (-0.845,-1.460) (-0.870,-1.484) (-0.896,-1.508) (-0.923,-1.531) (-0.950,-1.553) (-0.977,-1.575) (-1.005,-1.596) (-1.034,-1.616) (-1.063,-1.636) (-1.093,-1.655) (-1.122,-1.673) (-1.153,-1.690) (-1.184,-1.707) (-1.215,-1.723) (-1.247,-1.738) (-1.278,-1.753) (-1.311,-1.766) (-1.343,-1.779) (-1.376,-1.791) (-1.409,-1.803) (-1.443,-1.813) (-1.477,-1.823) (-1.511,-1.832) (-1.545,-1.840) (-1.579,-1.847) (-1.613,-1.854) (-1.648,-1.859) (-1.683,-1.864) (-1.718,-1.868) (-1.752,-1.872) (-1.787,-1.874) (-1.822,-1.876) (-1.857,-1.877) (-1.893,-1.877) (-1.928,-1.876) (-1.963,-1.875) (-1.998,-1.872) (-2.033,-1.869) (-2.067,-1.866) (-2.102,-1.861) (-2.137,-1.856) (-2.171,-1.850) (-2.206,-1.843) (-2.240,-1.835) (-2.274,-1.827) (-2.308,-1.818) (-2.342,-1.808) (-2.375,-1.798) (-2.408,-1.787) (-2.441,-1.775) (-2.474,-1.762) (-2.507,-1.749) (-2.539,-1.735) (-2.571,-1.721) (-2.602,-1.706) (-2.634,-1.690) (-2.665,-1.673) (-2.695,-1.656) (-2.725,-1.639) (-2.755,-1.620) (-2.785,-1.601) (-2.814,-1.582) (-2.843,-1.562) (-2.871,-1.541) (-2.899,-1.520) (-2.927,-1.499) (-2.954,-1.476) (-2.981,-1.454) (-3.007,-1.430) (-3.033,-1.407) (-3.058,-1.382) (-3.083,-1.358) (-3.107,-1.333) (-3.131,-1.307) (-3.155,-1.281) (-3.178,-1.254) (-3.200,-1.227)};
\addplot[equipot] coordinates {(3.200,1.227) (3.178,1.254) (3.155,1.281) (3.131,1.307) (3.107,1.333) (3.083,1.358) (3.058,1.382) (3.033,1.407) (3.007,1.430) (2.981,1.454) (2.954,1.476) (2.927,1.499) (2.899,1.520) (2.871,1.541) (2.843,1.562) (2.814,1.582) (2.785,1.601) (2.755,1.620) (2.725,1.639) (2.695,1.656) (2.665,1.673) (2.634,1.690) (2.602,1.706) (2.571,1.721) (2.539,1.735) (2.507,1.749) (2.474,1.762) (2.441,1.775) (2.408,1.787) (2.375,1.798) (2.342,1.808) (2.308,1.818) (2.274,1.827) (2.240,1.835) (2.206,1.843) (2.171,1.850) (2.137,1.856) (2.102,1.861) (2.067,1.866) (2.033,1.869) (1.998,1.872) (1.963,1.875) (1.928,1.876) (1.893,1.877) (1.857,1.877) (1.822,1.876) (1.787,1.874) (1.752,1.872) (1.718,1.868) (1.683,1.864) (1.648,1.859) (1.613,1.854) (1.579,1.847) (1.545,1.840) (1.511,1.832) (1.477,1.823) (1.443,1.813) (1.409,1.803) (1.376,1.791) (1.343,1.779) (1.311,1.766) (1.278,1.753) (1.247,1.738) (1.215,1.723) (1.184,1.707) (1.153,1.690) (1.122,1.673) (1.093,1.655) (1.063,1.636) (1.034,1.616) (1.005,1.596) (0.977,1.575) (0.950,1.553) (0.923,1.531) (0.896,1.508) (0.870,1.484) (0.845,1.460) (0.820,1.435) (0.796,1.410) (0.772,1.384) (0.749,1.358) (0.727,1.331) (0.705,1.303) (0.684,1.275) (0.663,1.247) (0.643,1.218) (0.624,1.189) (0.605,1.159) (0.588,1.129) (0.570,1.098) (0.553,1.067) (0.537,1.036) (0.522,1.005) (0.507,0.973) (0.493,0.941) (0.479,0.909) (0.466,0.876) (0.454,0.843) (0.442,0.810) (0.431,0.777) (0.420,0.744) (0.410,0.710) (0.401,0.676) (0.392,0.642) (0.383,0.608) (0.375,0.574) (0.368,0.540) (0.361,0.506) (0.355,0.471) (0.349,0.436) (0.344,0.402) (0.339,0.367) (0.334,0.332) (0.330,0.297) (0.327,0.263) (0.324,0.228) (0.321,0.193) (0.319,0.158) (0.317,0.123) (0.316,0.088) (0.315,0.053) (0.315,0.018) (0.315,-0.018) (0.315,-0.053) (0.316,-0.088) (0.317,-0.123) (0.319,-0.158) (0.321,-0.193) (0.324,-0.228) (0.327,-0.263) (0.330,-0.297) (0.334,-0.332) (0.339,-0.367) (0.344,-0.402) (0.349,-0.436) (0.355,-0.471) (0.361,-0.506) (0.368,-0.540) (0.375,-0.574) (0.383,-0.608) (0.392,-0.642) (0.401,-0.676) (0.410,-0.710) (0.420,-0.744) (0.431,-0.777) (0.442,-0.810) (0.454,-0.843) (0.466,-0.876) (0.479,-0.909) (0.493,-0.941) (0.507,-0.973) (0.522,-1.005) (0.537,-1.036) (0.553,-1.067) (0.570,-1.098) (0.588,-1.129) (0.605,-1.159) (0.624,-1.189) (0.643,-1.218) (0.663,-1.247) (0.684,-1.275) (0.705,-1.303) (0.727,-1.331) (0.749,-1.358) (0.772,-1.384) (0.796,-1.410) (0.820,-1.435) (0.845,-1.460) (0.870,-1.484) (0.896,-1.508) (0.923,-1.531) (0.950,-1.553) (0.977,-1.575) (1.005,-1.596) (1.034,-1.616) (1.063,-1.636) (1.093,-1.655) (1.122,-1.673) (1.153,-1.690) (1.184,-1.707) (1.215,-1.723) (1.247,-1.738) (1.278,-1.753) (1.311,-1.766) (1.343,-1.779) (1.376,-1.791) (1.409,-1.803) (1.443,-1.813) (1.477,-1.823) (1.511,-1.832) (1.545,-1.840) (1.579,-1.847) (1.613,-1.854) (1.648,-1.859) (1.683,-1.864) (1.718,-1.868) (1.752,-1.872) (1.787,-1.874) (1.822,-1.876) (1.857,-1.877) (1.893,-1.877) (1.928,-1.876) (1.963,-1.875) (1.998,-1.872) (2.033,-1.869) (2.067,-1.866) (2.102,-1.861) (2.137,-1.856) (2.171,-1.850) (2.206,-1.843) (2.240,-1.835) (2.274,-1.827) (2.308,-1.818) (2.342,-1.808) (2.375,-1.798) (2.408,-1.787) (2.441,-1.775) (2.474,-1.762) (2.507,-1.749) (2.539,-1.735) (2.571,-1.721) (2.602,-1.706) (2.634,-1.690) (2.665,-1.673) (2.695,-1.656) (2.725,-1.639) (2.755,-1.620) (2.785,-1.601) (2.814,-1.582) (2.843,-1.562) (2.871,-1.541) (2.899,-1.520) (2.927,-1.499) (2.954,-1.476) (2.981,-1.454) (3.007,-1.430) (3.033,-1.407) (3.058,-1.382) (3.083,-1.358) (3.107,-1.333) (3.131,-1.307) (3.155,-1.281) (3.178,-1.254) (3.200,-1.227)};
\addplot[equipot] coordinates {(1.452,-0.923) (1.487,-0.927) (1.522,-0.930) (1.557,-0.930) (1.592,-0.929) (1.627,-0.926) (1.662,-0.921) (1.697,-0.915) (1.731,-0.907) (1.765,-0.897) (1.798,-0.886) (1.831,-0.873) (1.863,-0.859) (1.894,-0.844) (1.925,-0.826) (1.955,-0.808) (1.984,-0.788) (2.013,-0.768) (2.040,-0.745) (2.067,-0.722) (2.092,-0.698) (2.117,-0.673) (2.140,-0.647) (2.163,-0.620) (2.184,-0.592) (2.205,-0.563) (2.224,-0.534) (2.242,-0.504) (2.259,-0.473) (2.275,-0.441) (2.290,-0.410) (2.304,-0.377) (2.316,-0.344) (2.327,-0.311) (2.337,-0.277) (2.346,-0.243) (2.354,-0.209) (2.361,-0.174) (2.366,-0.139) (2.370,-0.105) (2.373,-0.070) (2.375,-0.034) (2.376,0.001) (2.375,0.036) (2.373,0.071) (2.370,0.106) (2.366,0.141) (2.360,0.176) (2.354,0.210) (2.346,0.245) (2.337,0.279) (2.327,0.312) (2.315,0.346) (2.303,0.379) (2.289,0.411) (2.274,0.443) (2.258,0.474) (2.241,0.505) (2.223,0.535) (2.204,0.565) (2.183,0.593) (2.162,0.621) (2.139,0.648) (2.116,0.674) (2.091,0.699) (2.065,0.723) (2.039,0.747) (2.011,0.769) (1.983,0.789) (1.954,0.809) (1.924,0.827) (1.893,0.844) (1.861,0.860) (1.829,0.874) (1.796,0.887) (1.763,0.898) (1.729,0.907) (1.695,0.915) (1.660,0.922) (1.625,0.926) (1.590,0.929) (1.555,0.930) (1.520,0.929) (1.485,0.927) (1.450,0.923) (1.415,0.917) (1.381,0.909) (1.347,0.899) (1.314,0.888) (1.281,0.875) (1.249,0.861) (1.218,0.845) (1.187,0.827) (1.158,0.808) (1.129,0.788) (1.101,0.766) (1.075,0.743) (1.049,0.719) (1.025,0.694) (1.001,0.667) (0.979,0.640) (0.958,0.612) (0.938,0.583) (0.919,0.553) (0.901,0.523) (0.885,0.492) (0.869,0.460) (0.855,0.428) (0.842,0.395) (0.830,0.362) (0.819,0.329) (0.809,0.295) (0.801,0.261) (0.793,0.226) (0.787,0.192) (0.782,0.157) (0.777,0.122) (0.774,0.087) (0.772,0.052) (0.771,0.017) (0.771,-0.018) (0.772,-0.054) (0.774,-0.089) (0.778,-0.124) (0.782,-0.159) (0.787,-0.193) (0.794,-0.228) (0.801,-0.262) (0.810,-0.296) (0.820,-0.330) (0.831,-0.364) (0.843,-0.397) (0.856,-0.429) (0.870,-0.462) (0.885,-0.493) (0.902,-0.524) (0.920,-0.555) (0.939,-0.584) (0.959,-0.613) (0.980,-0.641) (1.002,-0.669) (1.026,-0.695) (1.050,-0.720) (1.076,-0.744) (1.103,-0.767) (1.130,-0.789) (1.159,-0.809) (1.189,-0.828) (1.219,-0.845) (1.251,-0.861) (1.283,-0.876) (1.315,-0.889) (1.349,-0.900) (1.383,-0.909) (1.417,-0.917) (1.452,-0.923)};
\addplot[equipot] coordinates {(1.468,-0.655) (1.503,-0.658) (1.538,-0.657) (1.573,-0.653) (1.607,-0.645) (1.640,-0.634) (1.673,-0.620) (1.704,-0.604) (1.733,-0.585) (1.762,-0.564) (1.789,-0.542) (1.814,-0.518) (1.838,-0.492) (1.861,-0.465) (1.882,-0.437) (1.901,-0.407) (1.919,-0.377) (1.935,-0.346) (1.950,-0.314) (1.963,-0.282) (1.975,-0.249) (1.985,-0.215) (1.994,-0.181) (2.001,-0.146) (2.006,-0.112) (2.011,-0.077) (2.013,-0.042) (2.014,-0.007) (2.014,0.028) (2.012,0.063) (2.008,0.098) (2.003,0.133) (1.997,0.168) (1.989,0.202) (1.979,0.236) (1.968,0.269) (1.955,0.302) (1.941,0.334) (1.925,0.365) (1.908,0.396) (1.889,0.426) (1.869,0.454) (1.847,0.482) (1.824,0.508) (1.799,0.533) (1.772,0.556) (1.745,0.577) (1.715,0.597) (1.685,0.614) (1.653,0.629) (1.620,0.641) (1.586,0.650) (1.551,0.656) (1.516,0.658) (1.481,0.657) (1.446,0.652) (1.412,0.643) (1.379,0.632) (1.347,0.617) (1.317,0.600) (1.288,0.580) (1.260,0.558) (1.234,0.534) (1.210,0.509) (1.187,0.482) (1.166,0.454) (1.147,0.425) (1.129,0.395) (1.113,0.363) (1.098,0.331) (1.085,0.299) (1.073,0.266) (1.063,0.232) (1.054,0.198) (1.047,0.164) (1.041,0.129) (1.036,0.094) (1.033,0.059) (1.031,0.024) (1.031,-0.011) (1.032,-0.046) (1.035,-0.081) (1.039,-0.116) (1.044,-0.151) (1.051,-0.185) (1.059,-0.219) (1.069,-0.253) (1.080,-0.286) (1.093,-0.319) (1.107,-0.351) (1.123,-0.383) (1.140,-0.413) (1.159,-0.443) (1.179,-0.472) (1.201,-0.499) (1.225,-0.525) (1.250,-0.549) (1.277,-0.572) (1.305,-0.593) (1.335,-0.611) (1.367,-0.627) (1.399,-0.639) (1.433,-0.649) (1.468,-0.655)};
\addplot[equipot] coordinates {(1.468,-0.532) (1.503,-0.537) (1.537,-0.532) (1.570,-0.519) (1.600,-0.501) (1.629,-0.480) (1.654,-0.456) (1.678,-0.430) (1.700,-0.402) (1.720,-0.373) (1.738,-0.342) (1.754,-0.311) (1.769,-0.279) (1.782,-0.246) (1.793,-0.212) (1.802,-0.178) (1.809,-0.144) (1.815,-0.109) (1.820,-0.074) (1.822,-0.039) (1.823,-0.004) (1.823,0.032) (1.820,0.067) (1.816,0.102) (1.811,0.137) (1.804,0.171) (1.795,0.205) (1.784,0.239) (1.772,0.272) (1.757,0.304) (1.742,0.336) (1.724,0.367) (1.704,0.396) (1.683,0.424) (1.660,0.450) (1.634,0.475) (1.606,0.497) (1.577,0.515) (1.544,0.529) (1.510,0.537) (1.475,0.534) (1.442,0.522) (1.411,0.504) (1.383,0.483) (1.357,0.459) (1.334,0.432) (1.313,0.404) (1.293,0.374) (1.276,0.344) (1.260,0.312) (1.247,0.280) (1.234,0.246) (1.224,0.213) (1.215,0.179) (1.208,0.144) (1.202,0.109) (1.198,0.074) (1.196,0.039) (1.195,0.004) (1.196,-0.032) (1.198,-0.067) (1.201,-0.102) (1.207,-0.137) (1.214,-0.171) (1.222,-0.206) (1.232,-0.240) (1.244,-0.273) (1.258,-0.305) (1.273,-0.337) (1.290,-0.368) (1.309,-0.398) (1.330,-0.426) (1.352,-0.453) (1.378,-0.478) (1.405,-0.500) (1.435,-0.519) (1.468,-0.532)};
\addplot[equipot] coordinates {(1.500,-0.482) (1.526,-0.459) (1.550,-0.434) (1.572,-0.406) (1.592,-0.378) (1.610,-0.347) (1.627,-0.316) (1.641,-0.284) (1.654,-0.252) (1.665,-0.218) (1.674,-0.185) (1.682,-0.150) (1.688,-0.116) (1.692,-0.081) (1.695,-0.046) (1.697,-0.011) (1.696,0.024) (1.694,0.059) (1.691,0.094) (1.686,0.129) (1.679,0.163) (1.671,0.197) (1.661,0.231) (1.649,0.264) (1.636,0.297) (1.621,0.328) (1.604,0.359) (1.585,0.389) (1.564,0.417) (1.541,0.444) (1.516,0.468) (1.490,0.474) (1.465,0.449) (1.442,0.422) (1.422,0.394) (1.403,0.364) (1.386,0.333) (1.372,0.301) (1.358,0.269) (1.347,0.235) (1.337,0.202) (1.329,0.168) (1.323,0.133) (1.317,0.098) (1.314,0.063) (1.312,0.028) (1.312,-0.007) (1.313,-0.042) (1.315,-0.077) (1.319,-0.112) (1.325,-0.146) (1.332,-0.181) (1.341,-0.215) (1.351,-0.248) (1.363,-0.281) (1.377,-0.313) (1.392,-0.345) (1.410,-0.375) (1.429,-0.405) (1.451,-0.432) (1.474,-0.458) (1.500,-0.482)};
\addplot[equipot] coordinates {(1.500,-0.395) (1.516,-0.364) (1.531,-0.332) (1.545,-0.300) (1.557,-0.267) (1.568,-0.234) (1.577,-0.200) (1.585,-0.166) (1.591,-0.131) (1.596,-0.097) (1.599,-0.062) (1.601,-0.027) (1.601,0.008) (1.600,0.043) (1.598,0.078) (1.594,0.113) (1.588,0.147) (1.581,0.182) (1.573,0.216) (1.563,0.249) (1.552,0.282) (1.539,0.315) (1.524,0.347) (1.508,0.378) (1.494,0.380) (1.477,0.349) (1.463,0.317) (1.451,0.284) (1.439,0.251) (1.430,0.217) (1.422,0.183) (1.415,0.149) (1.410,0.114) (1.406,0.079) (1.403,0.044) (1.402,0.009) (1.402,-0.026) (1.404,-0.061) (1.407,-0.095) (1.412,-0.130) (1.418,-0.165) (1.425,-0.199) (1.434,-0.233) (1.444,-0.266) (1.456,-0.299) (1.470,-0.332) (1.484,-0.363) (1.500,-0.395)};
\addplot[lineaD, flechaen=0.169, flechaen=0.831] coordinates {(-3.200,0.617) (-3.167,0.604) (-3.134,0.592) (-3.102,0.580) (-3.069,0.567) (-3.036,0.555) (-3.003,0.543) (-2.970,0.531) (-2.937,0.518) (-2.904,0.506) (-2.871,0.494) (-2.838,0.483) (-2.805,0.471) (-2.772,0.459) (-2.739,0.447) (-2.706,0.436) (-2.673,0.424) (-2.640,0.413) (-2.607,0.401) (-2.574,0.390) (-2.541,0.379) (-2.507,0.367) (-2.474,0.356) (-2.441,0.345) (-2.407,0.334) (-2.374,0.324) (-2.341,0.313) (-2.307,0.302) (-2.274,0.292) (-2.240,0.282) (-2.207,0.271) (-2.173,0.261) (-2.140,0.251) (-2.106,0.241) (-2.072,0.232) (-2.038,0.222) (-2.005,0.213) (-1.971,0.204) (-1.937,0.195) (-1.903,0.186) (-1.869,0.178) (-1.835,0.170) (-1.800,0.162) (-1.766,0.155) (-1.732,0.148) (-1.697,0.141) (-1.663,0.135) (-1.628,0.129) (-1.594,0.123) (-1.559,0.118) (-1.524,0.114) (-1.489,0.110) (-1.455,0.106) (-1.420,0.103) (-1.385,0.100) (-1.350,0.098) (-1.315,0.095) (-1.280,0.093) (-1.245,0.092) (-1.210,0.090) (-1.175,0.089) (-1.140,0.088) (-1.104,0.087) (-1.069,0.086) (-1.034,0.085) (-0.999,0.085) (-0.964,0.084) (-0.929,0.084) (-0.894,0.083) (-0.859,0.083) (-0.824,0.083) (-0.789,0.082) (-0.754,0.082) (-0.719,0.082) (-0.684,0.081) (-0.649,0.081) (-0.614,0.081) (-0.579,0.081) (-0.544,0.081) (-0.508,0.081) (-0.473,0.081) (-0.438,0.080) (-0.403,0.080) (-0.368,0.080) (-0.333,0.080) (-0.298,0.080) (-0.263,0.080) (-0.228,0.080) (-0.193,0.080) (-0.158,0.080) (-0.123,0.080) (-0.088,0.080) (-0.053,0.080) (-0.018,0.080) (0.018,0.080) (0.053,0.080) (0.088,0.080) (0.123,0.080) (0.158,0.080) (0.193,0.080) (0.228,0.080) (0.263,0.080) (0.298,0.080) (0.333,0.080) (0.368,0.080) (0.403,0.080) (0.438,0.080) (0.473,0.081) (0.508,0.081) (0.544,0.081) (0.579,0.081) (0.614,0.081) (0.649,0.081) (0.684,0.081) (0.719,0.082) (0.754,0.082) (0.789,0.082) (0.824,0.083) (0.859,0.083) (0.894,0.083) (0.929,0.084) (0.964,0.084) (0.999,0.085) (1.034,0.085) (1.069,0.086) (1.104,0.087) (1.140,0.088) (1.175,0.089) (1.210,0.090) (1.245,0.092) (1.280,0.093) (1.315,0.095) (1.350,0.098) (1.385,0.100) (1.420,0.103) (1.455,0.106) (1.489,0.110) (1.524,0.114) (1.559,0.118) (1.594,0.123) (1.628,0.129) (1.663,0.135) (1.697,0.141) (1.732,0.148) (1.766,0.155) (1.800,0.162) (1.835,0.170) (1.869,0.178) (1.903,0.186) (1.937,0.195) (1.971,0.204) (2.005,0.213) (2.038,0.222) (2.072,0.232) (2.106,0.241) (2.140,0.251) (2.173,0.261) (2.207,0.271) (2.240,0.282) (2.274,0.292) (2.307,0.302) (2.341,0.313) (2.374,0.324) (2.407,0.334) (2.441,0.345) (2.474,0.356) (2.507,0.367) (2.541,0.379) (2.574,0.390) (2.607,0.401) (2.640,0.413) (2.673,0.424) (2.706,0.436) (2.739,0.447) (2.772,0.459) (2.805,0.471) (2.838,0.483) (2.871,0.494) (2.904,0.506) (2.937,0.518) (2.970,0.531) (3.003,0.543) (3.036,0.555) (3.069,0.567) (3.102,0.580) (3.134,0.592) (3.167,0.604) (3.200,0.617)};
\addplot[lineaD, flechaen=0.190, flechaen=0.810] coordinates {(-3.200,1.424) (-3.175,1.399) (-3.150,1.375) (-3.124,1.351) (-3.099,1.327) (-3.074,1.302) (-3.048,1.279) (-3.023,1.255) (-2.997,1.231) (-2.971,1.207) (-2.945,1.184) (-2.919,1.160) (-2.893,1.137) (-2.867,1.113) (-2.841,1.090) (-2.814,1.067) (-2.788,1.044) (-2.762,1.021) (-2.735,0.998) (-2.708,0.976) (-2.682,0.953) (-2.655,0.931) (-2.628,0.908) (-2.601,0.886) (-2.574,0.864) (-2.546,0.842) (-2.519,0.820) (-2.492,0.798) (-2.464,0.776) (-2.437,0.755) (-2.409,0.733) (-2.381,0.712) (-2.353,0.691) (-2.325,0.670) (-2.297,0.649) (-2.269,0.628) (-2.240,0.608) (-2.212,0.587) (-2.183,0.567) (-2.155,0.547) (-2.126,0.527) (-2.097,0.508) (-2.068,0.488) (-2.038,0.469) (-2.009,0.450) (-1.979,0.431) (-1.949,0.413) (-1.919,0.395) (-1.889,0.378) (-1.859,0.360) (-1.828,0.344) (-1.797,0.328) (-1.765,0.312) (-1.734,0.297) (-1.702,0.283) (-1.669,0.270) (-1.636,0.257) (-1.603,0.246) (-1.570,0.235) (-1.536,0.226) (-1.502,0.217) (-1.468,0.209) (-1.434,0.203) (-1.399,0.197) (-1.365,0.191) (-1.330,0.187) (-1.295,0.183) (-1.260,0.180) (-1.225,0.177) (-1.190,0.174) (-1.155,0.172) (-1.121,0.170) (-1.086,0.168) (-1.051,0.167) (-1.016,0.166) (-0.981,0.164) (-0.946,0.163) (-0.911,0.163) (-0.875,0.162) (-0.840,0.161) (-0.805,0.161) (-0.770,0.160) (-0.735,0.159) (-0.700,0.159) (-0.665,0.159) (-0.630,0.158) (-0.595,0.158) (-0.560,0.158) (-0.525,0.157) (-0.490,0.157) (-0.455,0.157) (-0.420,0.157) (-0.385,0.157) (-0.350,0.157) (-0.315,0.156) (-0.280,0.156) (-0.245,0.156) (-0.210,0.156) (-0.175,0.156) (-0.140,0.156) (-0.105,0.156) (-0.070,0.156) (-0.035,0.156) (0.000,0.156) (0.035,0.156) (0.070,0.156) (0.105,0.156) (0.140,0.156) (0.175,0.156) (0.210,0.156) (0.245,0.156) (0.280,0.156) (0.315,0.156) (0.350,0.157) (0.385,0.157) (0.420,0.157) (0.455,0.157) (0.490,0.157) (0.525,0.157) (0.560,0.158) (0.595,0.158) (0.630,0.158) (0.665,0.159) (0.700,0.159) (0.735,0.159) (0.770,0.160) (0.805,0.161) (0.840,0.161) (0.875,0.162) (0.911,0.163) (0.946,0.163) (0.981,0.164) (1.016,0.166) (1.051,0.167) (1.086,0.168) (1.121,0.170) (1.155,0.172) (1.190,0.174) (1.225,0.177) (1.260,0.180) (1.295,0.183) (1.330,0.187) (1.365,0.191) (1.399,0.197) (1.434,0.203) (1.468,0.209) (1.502,0.217) (1.536,0.226) (1.570,0.235) (1.603,0.246) (1.636,0.257) (1.669,0.270) (1.702,0.283) (1.734,0.297) (1.765,0.312) (1.797,0.328) (1.828,0.344) (1.859,0.360) (1.889,0.378) (1.919,0.395) (1.949,0.413) (1.979,0.431) (2.009,0.450) (2.038,0.469) (2.068,0.488) (2.097,0.508) (2.126,0.527) (2.155,0.547) (2.183,0.567) (2.212,0.587) (2.240,0.608) (2.269,0.628) (2.297,0.649) (2.325,0.670) (2.353,0.691) (2.381,0.712) (2.409,0.733) (2.437,0.755) (2.464,0.776) (2.492,0.798) (2.519,0.820) (2.546,0.842) (2.574,0.864) (2.601,0.886) (2.628,0.908) (2.655,0.931) (2.682,0.953) (2.708,0.976) (2.735,0.998) (2.762,1.021) (2.788,1.044) (2.814,1.067) (2.841,1.090) (2.867,1.113) (2.893,1.137) (2.919,1.160) (2.945,1.184) (2.971,1.207) (2.997,1.231) (3.023,1.255) (3.048,1.279) (3.074,1.302) (3.099,1.327) (3.124,1.351) (3.150,1.375) (3.175,1.399) (3.200,1.424)};
\addplot[lineaD, flechaen=0.138, flechaen=0.862] coordinates {(-2.570,1.900) (-2.557,1.867) (-2.544,1.835) (-2.531,1.802) (-2.518,1.770) (-2.504,1.738) (-2.490,1.705) (-2.477,1.673) (-2.462,1.641) (-2.448,1.609) (-2.434,1.577) (-2.419,1.546) (-2.404,1.514) (-2.389,1.483) (-2.373,1.451) (-2.358,1.420) (-2.342,1.388) (-2.326,1.357) (-2.310,1.326) (-2.293,1.295) (-2.277,1.264) (-2.260,1.234) (-2.243,1.203) (-2.226,1.173) (-2.208,1.142) (-2.191,1.112) (-2.173,1.082) (-2.155,1.052) (-2.137,1.022) (-2.118,0.992) (-2.100,0.962) (-2.081,0.933) (-2.062,0.904) (-2.042,0.874) (-2.023,0.845) (-2.003,0.816) (-1.983,0.788) (-1.963,0.759) (-1.942,0.731) (-1.921,0.703) (-1.900,0.675) (-1.879,0.647) (-1.857,0.620) (-1.835,0.593) (-1.812,0.566) (-1.789,0.540) (-1.765,0.514) (-1.741,0.489) (-1.716,0.464) (-1.690,0.441) (-1.663,0.418) (-1.635,0.397) (-1.607,0.377) (-1.577,0.359) (-1.546,0.342) (-1.514,0.328) (-1.481,0.315) (-1.448,0.304) (-1.414,0.294) (-1.380,0.286) (-1.346,0.279) (-1.312,0.273) (-1.277,0.268) (-1.242,0.263) (-1.207,0.259) (-1.173,0.256) (-1.138,0.253) (-1.103,0.251) (-1.068,0.248) (-1.033,0.247) (-0.998,0.245) (-0.963,0.243) (-0.928,0.242) (-0.893,0.241) (-0.858,0.240) (-0.823,0.239) (-0.788,0.238) (-0.753,0.237) (-0.718,0.237) (-0.683,0.236) (-0.648,0.236) (-0.613,0.235) (-0.578,0.235) (-0.543,0.234) (-0.508,0.234) (-0.473,0.234) (-0.438,0.233) (-0.403,0.233) (-0.368,0.233) (-0.333,0.233) (-0.298,0.233) (-0.263,0.232) (-0.228,0.232) (-0.193,0.232) (-0.158,0.232) (-0.123,0.232) (-0.088,0.232) (-0.053,0.232) (-0.018,0.232) (0.018,0.232) (0.053,0.232) (0.088,0.232) (0.123,0.232) (0.158,0.232) (0.193,0.232) (0.228,0.232) (0.263,0.232) (0.298,0.233) (0.333,0.233) (0.368,0.233) (0.403,0.233) (0.438,0.233) (0.473,0.234) (0.508,0.234) (0.543,0.234) (0.578,0.235) (0.613,0.235) (0.648,0.236) (0.683,0.236) (0.718,0.237) (0.753,0.237) (0.788,0.238) (0.823,0.239) (0.858,0.240) (0.893,0.241) (0.928,0.242) (0.963,0.243) (0.998,0.245) (1.033,0.247) (1.068,0.248) (1.103,0.251) (1.138,0.253) (1.173,0.256) (1.207,0.259) (1.242,0.263) (1.277,0.268) (1.312,0.273) (1.346,0.279) (1.380,0.286) (1.414,0.294) (1.448,0.304) (1.481,0.315) (1.514,0.328) (1.546,0.342) (1.577,0.359) (1.607,0.377) (1.635,0.397) (1.663,0.418) (1.690,0.441) (1.716,0.464) (1.741,0.489) (1.765,0.514) (1.789,0.540) (1.812,0.566) (1.835,0.593) (1.857,0.620) (1.879,0.647) (1.900,0.675) (1.921,0.703) (1.942,0.731) (1.963,0.759) (1.983,0.788) (2.003,0.816) (2.023,0.845) (2.042,0.874) (2.062,0.904) (2.081,0.933) (2.100,0.962) (2.118,0.992) (2.137,1.022) (2.155,1.052) (2.173,1.082) (2.191,1.112) (2.208,1.142) (2.226,1.173) (2.243,1.203) (2.260,1.234) (2.277,1.264) (2.293,1.295) (2.310,1.326) (2.326,1.357) (2.342,1.388) (2.358,1.420) (2.373,1.451) (2.389,1.483) (2.404,1.514) (2.419,1.546) (2.434,1.577) (2.448,1.609) (2.462,1.641) (2.477,1.673) (2.490,1.705) (2.504,1.738) (2.518,1.770) (2.531,1.802) (2.544,1.835) (2.557,1.867) (2.570,1.900)};
\addplot[lineaD, flechaen=0.582] coordinates {(-1.772,1.900) (-1.774,1.865) (-1.776,1.830) (-1.778,1.795) (-1.779,1.759) (-1.780,1.724) (-1.781,1.689) (-1.781,1.654) (-1.781,1.619) (-1.780,1.583) (-1.779,1.548) (-1.778,1.513) (-1.776,1.478) (-1.774,1.443) (-1.772,1.408) (-1.769,1.373) (-1.766,1.338) (-1.763,1.303) (-1.759,1.268) (-1.755,1.233) (-1.750,1.198) (-1.746,1.163) (-1.741,1.128) (-1.735,1.093) (-1.729,1.059) (-1.723,1.024) (-1.717,0.989) (-1.710,0.955) (-1.703,0.920) (-1.695,0.886) (-1.687,0.852) (-1.679,0.817) (-1.670,0.783) (-1.661,0.749) (-1.652,0.715) (-1.641,0.682) (-1.631,0.648) (-1.619,0.615) (-1.607,0.582) (-1.593,0.550) (-1.578,0.518) (-1.560,0.488) (-1.538,0.460) (-1.512,0.436) (-1.483,0.417) (-1.452,0.401) (-1.419,0.388) (-1.386,0.377) (-1.352,0.368) (-1.317,0.360) (-1.283,0.353) (-1.248,0.348) (-1.213,0.343) (-1.178,0.338) (-1.143,0.335) (-1.108,0.332) (-1.073,0.329) (-1.038,0.326) (-1.003,0.324) (-0.968,0.322) (-0.933,0.321) (-0.897,0.319) (-0.862,0.318) (-0.827,0.317) (-0.792,0.316) (-0.757,0.315) (-0.721,0.314) (-0.686,0.313) (-0.651,0.313) (-0.616,0.312) (-0.581,0.311) (-0.546,0.311) (-0.510,0.311) (-0.475,0.310) (-0.440,0.310) (-0.405,0.310) (-0.370,0.309) (-0.334,0.309) (-0.299,0.309) (-0.264,0.309) (-0.229,0.308) (-0.194,0.308) (-0.158,0.308) (-0.123,0.308) (-0.088,0.308) (-0.053,0.308) (-0.018,0.308) (0.018,0.308) (0.053,0.308) (0.088,0.308) (0.123,0.308) (0.158,0.308) (0.194,0.308) (0.229,0.308) (0.264,0.309) (0.299,0.309) (0.334,0.309) (0.370,0.309) (0.405,0.310) (0.440,0.310) (0.475,0.310) (0.510,0.311) (0.546,0.311) (0.581,0.311) (0.616,0.312) (0.651,0.313) (0.686,0.313) (0.721,0.314) (0.757,0.315) (0.792,0.316) (0.827,0.317) (0.862,0.318) (0.897,0.319) (0.933,0.321) (0.968,0.322) (1.003,0.324) (1.038,0.326) (1.073,0.329) (1.108,0.332) (1.143,0.335) (1.178,0.338) (1.213,0.343) (1.248,0.348) (1.283,0.353) (1.317,0.360) (1.352,0.368) (1.386,0.377) (1.419,0.388) (1.452,0.401) (1.483,0.417) (1.512,0.436) (1.538,0.460) (1.560,0.488) (1.578,0.518) (1.593,0.550) (1.607,0.582) (1.619,0.615) (1.631,0.648) (1.641,0.682) (1.652,0.715) (1.661,0.749) (1.670,0.783) (1.679,0.817) (1.687,0.852) (1.695,0.886) (1.703,0.920) (1.710,0.955) (1.717,0.989) (1.723,1.024) (1.729,1.059) (1.735,1.093) (1.741,1.128) (1.746,1.163) (1.750,1.198) (1.755,1.233) (1.759,1.268) (1.763,1.303) (1.766,1.338) (1.769,1.373) (1.772,1.408) (1.774,1.443) (1.776,1.478) (1.778,1.513) (1.779,1.548) (1.780,1.583) (1.781,1.619) (1.781,1.654) (1.781,1.689) (1.780,1.724) (1.779,1.759) (1.778,1.795) (1.776,1.830) (1.774,1.865) (1.772,1.900)};
\addplot[lineaD, flechaen=0.418] coordinates {(-0.818,1.900) (-0.843,1.876) (-0.868,1.851) (-0.892,1.825) (-0.916,1.799) (-0.939,1.772) (-0.961,1.745) (-0.983,1.718) (-1.004,1.690) (-1.024,1.661) (-1.044,1.632) (-1.064,1.603) (-1.083,1.574) (-1.101,1.544) (-1.119,1.513) (-1.136,1.483) (-1.153,1.452) (-1.169,1.421) (-1.185,1.389) (-1.200,1.358) (-1.215,1.326) (-1.229,1.294) (-1.243,1.262) (-1.256,1.229) (-1.269,1.196) (-1.282,1.164) (-1.294,1.131) (-1.306,1.098) (-1.317,1.064) (-1.328,1.031) (-1.338,0.997) (-1.348,0.964) (-1.358,0.930) (-1.368,0.896) (-1.377,0.862) (-1.386,0.828) (-1.395,0.794) (-1.403,0.760) (-1.411,0.726) (-1.420,0.692) (-1.428,0.658) (-1.436,0.624) (-1.445,0.590) (-1.454,0.556) (-1.464,0.522) (-1.462,0.493) (-1.431,0.478) (-1.398,0.465) (-1.365,0.454) (-1.331,0.445) (-1.297,0.437) (-1.262,0.431) (-1.228,0.425) (-1.193,0.420) (-1.158,0.416) (-1.123,0.412) (-1.088,0.409) (-1.053,0.406) (-1.018,0.404) (-0.983,0.402) (-0.948,0.400) (-0.913,0.398) (-0.878,0.396) (-0.843,0.395) (-0.808,0.394) (-0.772,0.393) (-0.737,0.392) (-0.702,0.391) (-0.667,0.390) (-0.632,0.389) (-0.597,0.388) (-0.562,0.388) (-0.527,0.387) (-0.492,0.387) (-0.456,0.386) (-0.421,0.386) (-0.386,0.386) (-0.351,0.385) (-0.316,0.385) (-0.281,0.385) (-0.246,0.385) (-0.211,0.384) (-0.176,0.384) (-0.140,0.384) (-0.105,0.384) (-0.070,0.384) (-0.035,0.384) (0.000,0.384) (0.035,0.384) (0.070,0.384) (0.105,0.384) (0.140,0.384) (0.176,0.384) (0.211,0.384) (0.246,0.385) (0.281,0.385) (0.316,0.385) (0.351,0.385) (0.386,0.386) (0.421,0.386) (0.456,0.386) (0.492,0.387) (0.527,0.387) (0.562,0.388) (0.597,0.388) (0.632,0.389) (0.667,0.390) (0.702,0.391) (0.737,0.392) (0.772,0.393) (0.808,0.394) (0.843,0.395) (0.878,0.396) (0.913,0.398) (0.948,0.400) (0.983,0.402) (1.018,0.404) (1.053,0.406) (1.088,0.409) (1.123,0.412) (1.158,0.416) (1.193,0.420) (1.228,0.425) (1.262,0.431) (1.297,0.437) (1.331,0.445) (1.365,0.454) (1.398,0.465) (1.431,0.478) (1.462,0.493) (1.464,0.522) (1.454,0.556) (1.445,0.590) (1.436,0.624) (1.428,0.658) (1.420,0.692) (1.411,0.726) (1.403,0.760) (1.395,0.794) (1.386,0.828) (1.377,0.862) (1.368,0.896) (1.358,0.930) (1.348,0.964) (1.338,0.997) (1.328,1.031) (1.317,1.064) (1.306,1.098) (1.294,1.131) (1.282,1.164) (1.269,1.196) (1.256,1.229) (1.243,1.262) (1.229,1.294) (1.215,1.326) (1.200,1.358) (1.185,1.389) (1.169,1.421) (1.153,1.452) (1.136,1.483) (1.119,1.513) (1.101,1.544) (1.083,1.574) (1.064,1.603) (1.044,1.632) (1.024,1.661) (1.004,1.690) (0.983,1.718) (0.961,1.745) (0.939,1.772) (0.916,1.799) (0.892,1.825) (0.868,1.851) (0.843,1.876) (0.818,1.900)};
\addplot[lineaD, flechaen=0.209] coordinates {(-0.576,0.465) (-0.541,0.464) (-0.506,0.463) (-0.470,0.463) (-0.435,0.462) (-0.400,0.462) (-0.365,0.462) (-0.330,0.461) (-0.294,0.461) (-0.259,0.461) (-0.224,0.461) (-0.189,0.460) (-0.154,0.460) (-0.118,0.460) (-0.083,0.460) (-0.048,0.460) (-0.013,0.460) (0.022,0.460) (0.058,0.460) (0.093,0.460) (0.128,0.460) (0.163,0.460) (0.198,0.460) (0.234,0.461) (0.269,0.461) (0.304,0.461) (0.339,0.461) (0.374,0.462) (0.410,0.462) (0.445,0.463) (0.480,0.463) (0.515,0.464) (0.550,0.464) (0.586,0.465) (0.621,0.465) (0.656,0.466) (0.691,0.467) (0.726,0.468) (0.762,0.469) (0.797,0.470) (0.832,0.472) (0.867,0.473) (0.902,0.475) (0.937,0.477) (0.973,0.479) (1.008,0.481) (1.043,0.483) (1.078,0.486) (1.113,0.490) (1.148,0.493) (1.182,0.498) (1.183,0.517) (1.156,0.540) (1.130,0.564) (1.104,0.587) (1.077,0.611) (1.051,0.634) (1.025,0.657) (0.998,0.680) (0.971,0.703) (0.945,0.726) (0.918,0.749) (0.890,0.771) (0.863,0.793) (0.835,0.814) (0.807,0.835) (0.778,0.856) (0.749,0.876) (0.720,0.896) (0.691,0.915) (0.661,0.934) (0.630,0.952) (0.600,0.969) (0.569,0.986) (0.537,1.002) (0.506,1.017) (0.474,1.031) (0.441,1.045) (0.408,1.058) (0.375,1.069) (0.342,1.080) (0.308,1.090) (0.274,1.099) (0.240,1.107) (0.205,1.114) (0.170,1.120) (0.136,1.125) (0.101,1.129) (0.066,1.131) (0.030,1.133) (-0.005,1.133) (-0.040,1.133) (-0.075,1.131) (-0.110,1.128) (-0.145,1.124) (-0.180,1.119) (-0.215,1.113) (-0.249,1.105) (-0.283,1.097) (-0.317,1.088) (-0.351,1.077) (-0.384,1.066) (-0.417,1.054) (-0.450,1.041) (-0.482,1.027) (-0.514,1.013) (-0.546,0.997) (-0.577,0.981) (-0.608,0.964) (-0.639,0.947) (-0.669,0.928) (-0.699,0.910) (-0.728,0.890) (-0.757,0.871) (-0.786,0.850) (-0.814,0.829) (-0.842,0.808) (-0.870,0.787) (-0.898,0.765) (-0.925,0.742) (-0.952,0.720) (-0.979,0.697) (-1.005,0.674) (-1.032,0.651) (-1.058,0.628) (-1.085,0.604) (-1.111,0.581) (-1.137,0.557) (-1.164,0.534) (-1.190,0.511) (-1.173,0.496) (-1.138,0.492) (-1.103,0.489) (-1.068,0.486) (-1.033,0.483) (-0.998,0.480) (-0.963,0.478) (-0.928,0.476) (-0.893,0.474) (-0.857,0.473) (-0.822,0.471) (-0.787,0.470) (-0.752,0.469) (-0.717,0.468) (-0.681,0.467) (-0.646,0.466) (-0.611,0.465) (-0.576,0.465)};
\addplot[lineaD, flechaen=0.169, flechaen=0.831] coordinates {(-3.200,-0.617) (-3.167,-0.604) (-3.134,-0.592) (-3.102,-0.580) (-3.069,-0.567) (-3.036,-0.555) (-3.003,-0.543) (-2.970,-0.531) (-2.937,-0.518) (-2.904,-0.506) (-2.871,-0.494) (-2.838,-0.483) (-2.805,-0.471) (-2.772,-0.459) (-2.739,-0.447) (-2.706,-0.436) (-2.673,-0.424) (-2.640,-0.413) (-2.607,-0.401) (-2.574,-0.390) (-2.541,-0.379) (-2.507,-0.367) (-2.474,-0.356) (-2.441,-0.345) (-2.407,-0.334) (-2.374,-0.324) (-2.341,-0.313) (-2.307,-0.302) (-2.274,-0.292) (-2.240,-0.282) (-2.207,-0.271) (-2.173,-0.261) (-2.140,-0.251) (-2.106,-0.241) (-2.072,-0.232) (-2.038,-0.222) (-2.005,-0.213) (-1.971,-0.204) (-1.937,-0.195) (-1.903,-0.186) (-1.869,-0.178) (-1.835,-0.170) (-1.800,-0.162) (-1.766,-0.155) (-1.732,-0.148) (-1.697,-0.141) (-1.663,-0.135) (-1.628,-0.129) (-1.594,-0.123) (-1.559,-0.118) (-1.524,-0.114) (-1.489,-0.110) (-1.455,-0.106) (-1.420,-0.103) (-1.385,-0.100) (-1.350,-0.098) (-1.315,-0.095) (-1.280,-0.093) (-1.245,-0.092) (-1.210,-0.090) (-1.175,-0.089) (-1.140,-0.088) (-1.104,-0.087) (-1.069,-0.086) (-1.034,-0.085) (-0.999,-0.085) (-0.964,-0.084) (-0.929,-0.084) (-0.894,-0.083) (-0.859,-0.083) (-0.824,-0.083) (-0.789,-0.082) (-0.754,-0.082) (-0.719,-0.082) (-0.684,-0.081) (-0.649,-0.081) (-0.614,-0.081) (-0.579,-0.081) (-0.544,-0.081) (-0.508,-0.081) (-0.473,-0.081) (-0.438,-0.080) (-0.403,-0.080) (-0.368,-0.080) (-0.333,-0.080) (-0.298,-0.080) (-0.263,-0.080) (-0.228,-0.080) (-0.193,-0.080) (-0.158,-0.080) (-0.123,-0.080) (-0.088,-0.080) (-0.053,-0.080) (-0.018,-0.080) (0.018,-0.080) (0.053,-0.080) (0.088,-0.080) (0.123,-0.080) (0.158,-0.080) (0.193,-0.080) (0.228,-0.080) (0.263,-0.080) (0.298,-0.080) (0.333,-0.080) (0.368,-0.080) (0.403,-0.080) (0.438,-0.080) (0.473,-0.081) (0.508,-0.081) (0.544,-0.081) (0.579,-0.081) (0.614,-0.081) (0.649,-0.081) (0.684,-0.081) (0.719,-0.082) (0.754,-0.082) (0.789,-0.082) (0.824,-0.083) (0.859,-0.083) (0.894,-0.083) (0.929,-0.084) (0.964,-0.084) (0.999,-0.085) (1.034,-0.085) (1.069,-0.086) (1.104,-0.087) (1.140,-0.088) (1.175,-0.089) (1.210,-0.090) (1.245,-0.092) (1.280,-0.093) (1.315,-0.095) (1.350,-0.098) (1.385,-0.100) (1.420,-0.103) (1.455,-0.106) (1.489,-0.110) (1.524,-0.114) (1.559,-0.118) (1.594,-0.123) (1.628,-0.129) (1.663,-0.135) (1.697,-0.141) (1.732,-0.148) (1.766,-0.155) (1.800,-0.162) (1.835,-0.170) (1.869,-0.178) (1.903,-0.186) (1.937,-0.195) (1.971,-0.204) (2.005,-0.213) (2.038,-0.222) (2.072,-0.232) (2.106,-0.241) (2.140,-0.251) (2.173,-0.261) (2.207,-0.271) (2.240,-0.282) (2.274,-0.292) (2.307,-0.302) (2.341,-0.313) (2.374,-0.324) (2.407,-0.334) (2.441,-0.345) (2.474,-0.356) (2.507,-0.367) (2.541,-0.379) (2.574,-0.390) (2.607,-0.401) (2.640,-0.413) (2.673,-0.424) (2.706,-0.436) (2.739,-0.447) (2.772,-0.459) (2.805,-0.471) (2.838,-0.483) (2.871,-0.494) (2.904,-0.506) (2.937,-0.518) (2.970,-0.531) (3.003,-0.543) (3.036,-0.555) (3.069,-0.567) (3.102,-0.580) (3.134,-0.592) (3.167,-0.604) (3.200,-0.617)};
\addplot[lineaD, flechaen=0.190, flechaen=0.810] coordinates {(-3.200,-1.424) (-3.175,-1.399) (-3.150,-1.375) (-3.124,-1.351) (-3.099,-1.327) (-3.074,-1.302) (-3.048,-1.279) (-3.023,-1.255) (-2.997,-1.231) (-2.971,-1.207) (-2.945,-1.184) (-2.919,-1.160) (-2.893,-1.137) (-2.867,-1.113) (-2.841,-1.090) (-2.814,-1.067) (-2.788,-1.044) (-2.762,-1.021) (-2.735,-0.998) (-2.708,-0.976) (-2.682,-0.953) (-2.655,-0.931) (-2.628,-0.908) (-2.601,-0.886) (-2.574,-0.864) (-2.546,-0.842) (-2.519,-0.820) (-2.492,-0.798) (-2.464,-0.776) (-2.437,-0.755) (-2.409,-0.733) (-2.381,-0.712) (-2.353,-0.691) (-2.325,-0.670) (-2.297,-0.649) (-2.269,-0.628) (-2.240,-0.608) (-2.212,-0.587) (-2.183,-0.567) (-2.155,-0.547) (-2.126,-0.527) (-2.097,-0.508) (-2.068,-0.488) (-2.038,-0.469) (-2.009,-0.450) (-1.979,-0.431) (-1.949,-0.413) (-1.919,-0.395) (-1.889,-0.378) (-1.859,-0.360) (-1.828,-0.344) (-1.797,-0.328) (-1.765,-0.312) (-1.734,-0.297) (-1.702,-0.283) (-1.669,-0.270) (-1.636,-0.257) (-1.603,-0.246) (-1.570,-0.235) (-1.536,-0.226) (-1.502,-0.217) (-1.468,-0.209) (-1.434,-0.203) (-1.399,-0.197) (-1.365,-0.191) (-1.330,-0.187) (-1.295,-0.183) (-1.260,-0.180) (-1.225,-0.177) (-1.190,-0.174) (-1.155,-0.172) (-1.121,-0.170) (-1.086,-0.168) (-1.051,-0.167) (-1.016,-0.166) (-0.981,-0.164) (-0.946,-0.163) (-0.911,-0.163) (-0.875,-0.162) (-0.840,-0.161) (-0.805,-0.161) (-0.770,-0.160) (-0.735,-0.159) (-0.700,-0.159) (-0.665,-0.159) (-0.630,-0.158) (-0.595,-0.158) (-0.560,-0.158) (-0.525,-0.157) (-0.490,-0.157) (-0.455,-0.157) (-0.420,-0.157) (-0.385,-0.157) (-0.350,-0.157) (-0.315,-0.156) (-0.280,-0.156) (-0.245,-0.156) (-0.210,-0.156) (-0.175,-0.156) (-0.140,-0.156) (-0.105,-0.156) (-0.070,-0.156) (-0.035,-0.156) (0.000,-0.156) (0.035,-0.156) (0.070,-0.156) (0.105,-0.156) (0.140,-0.156) (0.175,-0.156) (0.210,-0.156) (0.245,-0.156) (0.280,-0.156) (0.315,-0.156) (0.350,-0.157) (0.385,-0.157) (0.420,-0.157) (0.455,-0.157) (0.490,-0.157) (0.525,-0.157) (0.560,-0.158) (0.595,-0.158) (0.630,-0.158) (0.665,-0.159) (0.700,-0.159) (0.735,-0.159) (0.770,-0.160) (0.805,-0.161) (0.840,-0.161) (0.875,-0.162) (0.911,-0.163) (0.946,-0.163) (0.981,-0.164) (1.016,-0.166) (1.051,-0.167) (1.086,-0.168) (1.121,-0.170) (1.155,-0.172) (1.190,-0.174) (1.225,-0.177) (1.260,-0.180) (1.295,-0.183) (1.330,-0.187) (1.365,-0.191) (1.399,-0.197) (1.434,-0.203) (1.468,-0.209) (1.502,-0.217) (1.536,-0.226) (1.570,-0.235) (1.603,-0.246) (1.636,-0.257) (1.669,-0.270) (1.702,-0.283) (1.734,-0.297) (1.765,-0.312) (1.797,-0.328) (1.828,-0.344) (1.859,-0.360) (1.889,-0.378) (1.919,-0.395) (1.949,-0.413) (1.979,-0.431) (2.009,-0.450) (2.038,-0.469) (2.068,-0.488) (2.097,-0.508) (2.126,-0.527) (2.155,-0.547) (2.183,-0.567) (2.212,-0.587) (2.240,-0.608) (2.269,-0.628) (2.297,-0.649) (2.325,-0.670) (2.353,-0.691) (2.381,-0.712) (2.409,-0.733) (2.437,-0.755) (2.464,-0.776) (2.492,-0.798) (2.519,-0.820) (2.546,-0.842) (2.574,-0.864) (2.601,-0.886) (2.628,-0.908) (2.655,-0.931) (2.682,-0.953) (2.708,-0.976) (2.735,-0.998) (2.762,-1.021) (2.788,-1.044) (2.814,-1.067) (2.841,-1.090) (2.867,-1.113) (2.893,-1.137) (2.919,-1.160) (2.945,-1.184) (2.971,-1.207) (2.997,-1.231) (3.023,-1.255) (3.048,-1.279) (3.074,-1.302) (3.099,-1.327) (3.124,-1.351) (3.150,-1.375) (3.175,-1.399) (3.200,-1.424)};
\addplot[lineaD, flechaen=0.138, flechaen=0.862] coordinates {(-2.570,-1.900) (-2.557,-1.867) (-2.544,-1.835) (-2.531,-1.802) (-2.518,-1.770) (-2.504,-1.738) (-2.490,-1.705) (-2.477,-1.673) (-2.462,-1.641) (-2.448,-1.609) (-2.434,-1.577) (-2.419,-1.546) (-2.404,-1.514) (-2.389,-1.483) (-2.373,-1.451) (-2.358,-1.420) (-2.342,-1.388) (-2.326,-1.357) (-2.310,-1.326) (-2.293,-1.295) (-2.277,-1.264) (-2.260,-1.234) (-2.243,-1.203) (-2.226,-1.173) (-2.208,-1.142) (-2.191,-1.112) (-2.173,-1.082) (-2.155,-1.052) (-2.137,-1.022) (-2.118,-0.992) (-2.100,-0.962) (-2.081,-0.933) (-2.062,-0.904) (-2.042,-0.874) (-2.023,-0.845) (-2.003,-0.816) (-1.983,-0.788) (-1.963,-0.759) (-1.942,-0.731) (-1.921,-0.703) (-1.900,-0.675) (-1.879,-0.647) (-1.857,-0.620) (-1.835,-0.593) (-1.812,-0.566) (-1.789,-0.540) (-1.765,-0.514) (-1.741,-0.489) (-1.716,-0.464) (-1.690,-0.441) (-1.663,-0.418) (-1.635,-0.397) (-1.607,-0.377) (-1.577,-0.359) (-1.546,-0.342) (-1.514,-0.328) (-1.481,-0.315) (-1.448,-0.304) (-1.414,-0.294) (-1.380,-0.286) (-1.346,-0.279) (-1.312,-0.273) (-1.277,-0.268) (-1.242,-0.263) (-1.207,-0.259) (-1.173,-0.256) (-1.138,-0.253) (-1.103,-0.251) (-1.068,-0.248) (-1.033,-0.247) (-0.998,-0.245) (-0.963,-0.243) (-0.928,-0.242) (-0.893,-0.241) (-0.858,-0.240) (-0.823,-0.239) (-0.788,-0.238) (-0.753,-0.237) (-0.718,-0.237) (-0.683,-0.236) (-0.648,-0.236) (-0.613,-0.235) (-0.578,-0.235) (-0.543,-0.234) (-0.508,-0.234) (-0.473,-0.234) (-0.438,-0.233) (-0.403,-0.233) (-0.368,-0.233) (-0.333,-0.233) (-0.298,-0.233) (-0.263,-0.232) (-0.228,-0.232) (-0.193,-0.232) (-0.158,-0.232) (-0.123,-0.232) (-0.088,-0.232) (-0.053,-0.232) (-0.018,-0.232) (0.018,-0.232) (0.053,-0.232) (0.088,-0.232) (0.123,-0.232) (0.158,-0.232) (0.193,-0.232) (0.228,-0.232) (0.263,-0.232) (0.298,-0.233) (0.333,-0.233) (0.368,-0.233) (0.403,-0.233) (0.438,-0.233) (0.473,-0.234) (0.508,-0.234) (0.543,-0.234) (0.578,-0.235) (0.613,-0.235) (0.648,-0.236) (0.683,-0.236) (0.718,-0.237) (0.753,-0.237) (0.788,-0.238) (0.823,-0.239) (0.858,-0.240) (0.893,-0.241) (0.928,-0.242) (0.963,-0.243) (0.998,-0.245) (1.033,-0.247) (1.068,-0.248) (1.103,-0.251) (1.138,-0.253) (1.173,-0.256) (1.207,-0.259) (1.242,-0.263) (1.277,-0.268) (1.312,-0.273) (1.346,-0.279) (1.380,-0.286) (1.414,-0.294) (1.448,-0.304) (1.481,-0.315) (1.514,-0.328) (1.546,-0.342) (1.577,-0.359) (1.607,-0.377) (1.635,-0.397) (1.663,-0.418) (1.690,-0.441) (1.716,-0.464) (1.741,-0.489) (1.765,-0.514) (1.789,-0.540) (1.812,-0.566) (1.835,-0.593) (1.857,-0.620) (1.879,-0.647) (1.900,-0.675) (1.921,-0.703) (1.942,-0.731) (1.963,-0.759) (1.983,-0.788) (2.003,-0.816) (2.023,-0.845) (2.042,-0.874) (2.062,-0.904) (2.081,-0.933) (2.100,-0.962) (2.118,-0.992) (2.137,-1.022) (2.155,-1.052) (2.173,-1.082) (2.191,-1.112) (2.208,-1.142) (2.226,-1.173) (2.243,-1.203) (2.260,-1.234) (2.277,-1.264) (2.293,-1.295) (2.310,-1.326) (2.326,-1.357) (2.342,-1.388) (2.358,-1.420) (2.373,-1.451) (2.389,-1.483) (2.404,-1.514) (2.419,-1.546) (2.434,-1.577) (2.448,-1.609) (2.462,-1.641) (2.477,-1.673) (2.490,-1.705) (2.504,-1.738) (2.518,-1.770) (2.531,-1.802) (2.544,-1.835) (2.557,-1.867) (2.570,-1.900)};
\addplot[lineaD, flechaen=0.418] coordinates {(-1.772,-1.900) (-1.774,-1.865) (-1.776,-1.830) (-1.778,-1.795) (-1.779,-1.759) (-1.780,-1.724) (-1.781,-1.689) (-1.781,-1.654) (-1.781,-1.619) (-1.780,-1.583) (-1.779,-1.548) (-1.778,-1.513) (-1.776,-1.478) (-1.774,-1.443) (-1.772,-1.408) (-1.769,-1.373) (-1.766,-1.338) (-1.763,-1.303) (-1.759,-1.268) (-1.755,-1.233) (-1.750,-1.198) (-1.746,-1.163) (-1.741,-1.128) (-1.735,-1.093) (-1.729,-1.059) (-1.723,-1.024) (-1.717,-0.989) (-1.710,-0.955) (-1.703,-0.920) (-1.695,-0.886) (-1.687,-0.852) (-1.679,-0.817) (-1.670,-0.783) (-1.661,-0.749) (-1.652,-0.715) (-1.641,-0.682) (-1.631,-0.648) (-1.619,-0.615) (-1.607,-0.582) (-1.593,-0.550) (-1.578,-0.518) (-1.560,-0.488) (-1.538,-0.460) (-1.512,-0.436) (-1.483,-0.417) (-1.452,-0.401) (-1.419,-0.388) (-1.386,-0.377) (-1.352,-0.368) (-1.317,-0.360) (-1.283,-0.353) (-1.248,-0.348) (-1.213,-0.343) (-1.178,-0.338) (-1.143,-0.335) (-1.108,-0.332) (-1.073,-0.329) (-1.038,-0.326) (-1.003,-0.324) (-0.968,-0.322) (-0.933,-0.321) (-0.897,-0.319) (-0.862,-0.318) (-0.827,-0.317) (-0.792,-0.316) (-0.757,-0.315) (-0.721,-0.314) (-0.686,-0.313) (-0.651,-0.313) (-0.616,-0.312) (-0.581,-0.311) (-0.546,-0.311) (-0.510,-0.311) (-0.475,-0.310) (-0.440,-0.310) (-0.405,-0.310) (-0.370,-0.309) (-0.334,-0.309) (-0.299,-0.309) (-0.264,-0.309) (-0.229,-0.308) (-0.194,-0.308) (-0.158,-0.308) (-0.123,-0.308) (-0.088,-0.308) (-0.053,-0.308) (-0.018,-0.308) (0.018,-0.308) (0.053,-0.308) (0.088,-0.308) (0.123,-0.308) (0.158,-0.308) (0.194,-0.308) (0.229,-0.308) (0.264,-0.309) (0.299,-0.309) (0.334,-0.309) (0.370,-0.309) (0.405,-0.310) (0.440,-0.310) (0.475,-0.310) (0.510,-0.311) (0.546,-0.311) (0.581,-0.311) (0.616,-0.312) (0.651,-0.313) (0.686,-0.313) (0.721,-0.314) (0.757,-0.315) (0.792,-0.316) (0.827,-0.317) (0.862,-0.318) (0.897,-0.319) (0.933,-0.321) (0.968,-0.322) (1.003,-0.324) (1.038,-0.326) (1.073,-0.329) (1.108,-0.332) (1.143,-0.335) (1.178,-0.338) (1.213,-0.343) (1.248,-0.348) (1.283,-0.353) (1.317,-0.360) (1.352,-0.368) (1.386,-0.377) (1.419,-0.388) (1.452,-0.401) (1.483,-0.417) (1.512,-0.436) (1.538,-0.460) (1.560,-0.488) (1.578,-0.518) (1.593,-0.550) (1.607,-0.582) (1.619,-0.615) (1.631,-0.648) (1.641,-0.682) (1.652,-0.715) (1.661,-0.749) (1.670,-0.783) (1.679,-0.817) (1.687,-0.852) (1.695,-0.886) (1.703,-0.920) (1.710,-0.955) (1.717,-0.989) (1.723,-1.024) (1.729,-1.059) (1.735,-1.093) (1.741,-1.128) (1.746,-1.163) (1.750,-1.198) (1.755,-1.233) (1.759,-1.268) (1.763,-1.303) (1.766,-1.338) (1.769,-1.373) (1.772,-1.408) (1.774,-1.443) (1.776,-1.478) (1.778,-1.513) (1.779,-1.548) (1.780,-1.583) (1.781,-1.619) (1.781,-1.654) (1.781,-1.689) (1.780,-1.724) (1.779,-1.759) (1.778,-1.795) (1.776,-1.830) (1.774,-1.865) (1.772,-1.900)};
\addplot[lineaD, flechaen=0.582] coordinates {(-0.818,-1.900) (-0.843,-1.876) (-0.868,-1.851) (-0.892,-1.825) (-0.916,-1.799) (-0.939,-1.772) (-0.961,-1.745) (-0.983,-1.718) (-1.004,-1.690) (-1.024,-1.661) (-1.044,-1.632) (-1.064,-1.603) (-1.083,-1.574) (-1.101,-1.544) (-1.119,-1.513) (-1.136,-1.483) (-1.153,-1.452) (-1.169,-1.421) (-1.185,-1.389) (-1.200,-1.358) (-1.215,-1.326) (-1.229,-1.294) (-1.243,-1.262) (-1.256,-1.229) (-1.269,-1.196) (-1.282,-1.164) (-1.294,-1.131) (-1.306,-1.098) (-1.317,-1.064) (-1.328,-1.031) (-1.338,-0.997) (-1.348,-0.964) (-1.358,-0.930) (-1.368,-0.896) (-1.377,-0.862) (-1.386,-0.828) (-1.395,-0.794) (-1.403,-0.760) (-1.411,-0.726) (-1.420,-0.692) (-1.428,-0.658) (-1.436,-0.624) (-1.445,-0.590) (-1.454,-0.556) (-1.464,-0.522) (-1.462,-0.493) (-1.431,-0.478) (-1.398,-0.465) (-1.365,-0.454) (-1.331,-0.445) (-1.297,-0.437) (-1.262,-0.431) (-1.228,-0.425) (-1.193,-0.420) (-1.158,-0.416) (-1.123,-0.412) (-1.088,-0.409) (-1.053,-0.406) (-1.018,-0.404) (-0.983,-0.402) (-0.948,-0.400) (-0.913,-0.398) (-0.878,-0.396) (-0.843,-0.395) (-0.808,-0.394) (-0.772,-0.393) (-0.737,-0.392) (-0.702,-0.391) (-0.667,-0.390) (-0.632,-0.389) (-0.597,-0.388) (-0.562,-0.388) (-0.527,-0.387) (-0.492,-0.387) (-0.456,-0.386) (-0.421,-0.386) (-0.386,-0.386) (-0.351,-0.385) (-0.316,-0.385) (-0.281,-0.385) (-0.246,-0.385) (-0.211,-0.384) (-0.176,-0.384) (-0.140,-0.384) (-0.105,-0.384) (-0.070,-0.384) (-0.035,-0.384) (0.000,-0.384) (0.035,-0.384) (0.070,-0.384) (0.105,-0.384) (0.140,-0.384) (0.176,-0.384) (0.211,-0.384) (0.246,-0.385) (0.281,-0.385) (0.316,-0.385) (0.351,-0.385) (0.386,-0.386) (0.421,-0.386) (0.456,-0.386) (0.492,-0.387) (0.527,-0.387) (0.562,-0.388) (0.597,-0.388) (0.632,-0.389) (0.667,-0.390) (0.702,-0.391) (0.737,-0.392) (0.772,-0.393) (0.808,-0.394) (0.843,-0.395) (0.878,-0.396) (0.913,-0.398) (0.948,-0.400) (0.983,-0.402) (1.018,-0.404) (1.053,-0.406) (1.088,-0.409) (1.123,-0.412) (1.158,-0.416) (1.193,-0.420) (1.228,-0.425) (1.262,-0.431) (1.297,-0.437) (1.331,-0.445) (1.365,-0.454) (1.398,-0.465) (1.431,-0.478) (1.462,-0.493) (1.464,-0.522) (1.454,-0.556) (1.445,-0.590) (1.436,-0.624) (1.428,-0.658) (1.420,-0.692) (1.411,-0.726) (1.403,-0.760) (1.395,-0.794) (1.386,-0.828) (1.377,-0.862) (1.368,-0.896) (1.358,-0.930) (1.348,-0.964) (1.338,-0.997) (1.328,-1.031) (1.317,-1.064) (1.306,-1.098) (1.294,-1.131) (1.282,-1.164) (1.269,-1.196) (1.256,-1.229) (1.243,-1.262) (1.229,-1.294) (1.215,-1.326) (1.200,-1.358) (1.185,-1.389) (1.169,-1.421) (1.153,-1.452) (1.136,-1.483) (1.119,-1.513) (1.101,-1.544) (1.083,-1.574) (1.064,-1.603) (1.044,-1.632) (1.024,-1.661) (1.004,-1.690) (0.983,-1.718) (0.961,-1.745) (0.939,-1.772) (0.916,-1.799) (0.892,-1.825) (0.868,-1.851) (0.843,-1.876) (0.818,-1.900)};
\addplot[lineaD, flechaen=0.500] coordinates {(-0.576,-0.465) (-0.541,-0.464) (-0.506,-0.463) (-0.470,-0.463) (-0.435,-0.462) (-0.400,-0.462) (-0.365,-0.462) (-0.330,-0.461) (-0.294,-0.461) (-0.259,-0.461) (-0.224,-0.461) (-0.189,-0.460) (-0.154,-0.460) (-0.118,-0.460) (-0.083,-0.460) (-0.048,-0.460) (-0.013,-0.460) (0.022,-0.460) (0.058,-0.460) (0.093,-0.460) (0.128,-0.460) (0.163,-0.460) (0.198,-0.460) (0.234,-0.461) (0.269,-0.461) (0.304,-0.461) (0.339,-0.461) (0.374,-0.462) (0.410,-0.462) (0.445,-0.463) (0.480,-0.463) (0.515,-0.464) (0.550,-0.464) (0.586,-0.465) (0.621,-0.465) (0.656,-0.466) (0.691,-0.467) (0.726,-0.468) (0.762,-0.469) (0.797,-0.470) (0.832,-0.472) (0.867,-0.473) (0.902,-0.475) (0.937,-0.477) (0.973,-0.479) (1.008,-0.481) (1.043,-0.483) (1.078,-0.486) (1.113,-0.490) (1.148,-0.493) (1.182,-0.498) (1.183,-0.517) (1.156,-0.540) (1.130,-0.564) (1.104,-0.587) (1.077,-0.611) (1.051,-0.634) (1.025,-0.657) (0.998,-0.680) (0.971,-0.703) (0.945,-0.726) (0.918,-0.749) (0.890,-0.771) (0.863,-0.793) (0.835,-0.814) (0.807,-0.835) (0.778,-0.856) (0.749,-0.876) (0.720,-0.896) (0.691,-0.915) (0.661,-0.934) (0.630,-0.952) (0.600,-0.969) (0.569,-0.986) (0.537,-1.002) (0.506,-1.017) (0.474,-1.031) (0.441,-1.045) (0.408,-1.058) (0.375,-1.069) (0.342,-1.080) (0.308,-1.090) (0.274,-1.099) (0.240,-1.107) (0.205,-1.114) (0.170,-1.120) (0.136,-1.125) (0.101,-1.129) (0.066,-1.131) (0.030,-1.133) (-0.005,-1.133) (-0.040,-1.133) (-0.075,-1.131) (-0.110,-1.128) (-0.145,-1.124) (-0.180,-1.119) (-0.215,-1.113) (-0.249,-1.105) (-0.283,-1.097) (-0.317,-1.088) (-0.351,-1.077) (-0.384,-1.066) (-0.417,-1.054) (-0.450,-1.041) (-0.482,-1.027) (-0.514,-1.013) (-0.546,-0.997) (-0.577,-0.981) (-0.608,-0.964) (-0.639,-0.947) (-0.669,-0.928) (-0.699,-0.910) (-0.728,-0.890) (-0.757,-0.871) (-0.786,-0.850) (-0.814,-0.829) (-0.842,-0.808) (-0.870,-0.787) (-0.898,-0.765) (-0.925,-0.742) (-0.952,-0.720) (-0.979,-0.697) (-1.005,-0.674) (-1.032,-0.651) (-1.058,-0.628) (-1.085,-0.604) (-1.111,-0.581) (-1.137,-0.557) (-1.164,-0.534) (-1.190,-0.511) (-1.173,-0.496) (-1.138,-0.492) (-1.103,-0.489) (-1.068,-0.486) (-1.033,-0.483) (-0.998,-0.480) (-0.963,-0.478) (-0.928,-0.476) (-0.893,-0.474) (-0.857,-0.473) (-0.822,-0.471) (-0.787,-0.470) (-0.752,-0.469) (-0.717,-0.468) (-0.681,-0.467) (-0.646,-0.466) (-0.611,-0.465) (-0.576,-0.465)};
\addplot[lineaD, flechaen=0.164, flechaen=0.836] coordinates {(-3.200,0.000) (3.200,0.000)};

  \draw[texto, line width=0.6pt] (axis cs:-1.5,-0.5) rectangle (axis cs:1.5,0.5);
\end{axis}
\end{tikzpicture}
\caption{$\vD(\vr)$}
\end{subfigure}
\caption{Electrete cilíndrico uniformemente polarizado ($a/\ell=1/3$), calculado con las funciones de
corriente exactas; trazos: equipotenciales. (a) La polarización es uniforme y axial. (b) $\vE$ nace en
la cara positiva y muere en la negativa, de modo que dentro se opone a $\vP$; fuera es el campo de dos
discos con cargas $\pm\sigmaP$. (c) $\vD$ forma líneas cerradas, como el campo de un solenoide: dentro tiene
el sentido de $\vP$ y se curva junto a la pared lateral, donde $\rot\vP\neq\vcero$.}
\label{fig:cilindro-lineas}
\end{figure}


% -----------------------------------------------------------------------------
\subsection{Problema 4.5. Cubo con polarización radial}\label{sec:p45}

\begin{problema}[Cubo con polarización radial]\label{prob:cubo}
Sea un cubo de material dieléctrico de lado $\ell$ centrado en el origen que presenta una polarización radial $\vP=A\vr$, donde $A$ es una
constante y $\vr$ el vector de posición de cada punto. Obtén las densidades ligadas de carga y demuestra que la carga ligada total es cero.
\end{problema}

\begin{resolucion}
\idea Aplicación directa del \cref{teo:densidades}, en cartesianas porque la frontera es un cubo.

\paso{Paso 1: densidad volúmica} Con $\vP=A(x\,\uv x+y\,\uv y+z\,\uv z)$ y $\partial x_i/\partial x_i=1$,
\[
  \rhoP=-\divg\vP=-A\Big(\frac{\partial x}{\partial x}+\frac{\partial y}{\partial y}+\frac{\partial z}{\partial z}\Big)=-3A .
\]
\paso{Paso 2: densidades superficiales} En la cara $y=\ell/2$, $\nn=\uv y$ y $\sigmaP=\vP\cdot\uv y=Ay=A\ell/2$; en la cara $y=-\ell/2$, $\nn=-\uv y$ y
$\sigmaP=-Ay=A\ell/2$. Por simetría, lo mismo en las otras cuatro caras:
\[
  \recuadro{\rhoP=-3A,\qquad \sigmaP=\frac{A\ell}{2}\ \text{(uniforme en cada cara)}.}
\]
\paso{Paso 3: carga total}
\[
  \QP=\int_V\rhoP\,\dd V+\oiint_S\sigmaP\,\dd S=-3A\,\ell^3+6\cdot\frac{A\ell}{2}\,\ell^2=0 . \qedhere
\]
\end{resolucion}

\begin{observacion}[Más allá de la neutralidad]\label{obs:p45}
También el momento dipolar total se anula, $\int\vP\,\dd V=A\int\vr\,\dd V=\vcero$ por simetría (\cref{obs:momento-ligado}): lejos del cubo, el campo
decae más deprisa que el de un dipolo. Obsérvese que $\sigmaP$ es uniforme en cada cara aunque $\vP$ no lo sea, porque en una cara solo cuenta la
componente normal, que es constante (\cref{fig:p45}).
\end{observacion}

% Figura: problema 4.5. Cubo con polarización radial P = A r.
\begin{figure}[t]
\centering
\begin{tikzpicture}[scale=1.0, x={(1cm,0cm)}, y={(0.45cm,0.32cm)}, z={(0cm,1cm)}]
  \def\h{1.2}
  % caras traseras
  \draw[texto!45!fondo, line width=0.4pt, dashed] (-\h,\h,-\h) -- (\h,\h,-\h) (-\h,\h,-\h) -- (-\h,\h,\h) (-\h,\h,-\h) -- (-\h,-\h,-\h);
  % flechas radiales de P desde el centro
  \foreach \v in {{1,0,0},{-1,0,0},{0,0,1},{0,0,-1},{0.7,0,0.7},{-0.7,0,0.7},{0.7,0,-0.7},{-0.7,0,-0.7},{0,1,0},{0,-1,0}} {
    \draw[campoP, line width=0.8pt, -{Stealth[length=1.6mm]}] (0,0,0) -- ($0.95*\h*({\v})$);
  }
  % caras delanteras
  \fill[dielec, opacity=0.55] (-\h,-\h,-\h) -- (\h,-\h,-\h) -- (\h,-\h,\h) -- (-\h,-\h,\h) -- cycle;
  \fill[dielecB, opacity=0.5] (\h,-\h,-\h) -- (\h,\h,-\h) -- (\h,\h,\h) -- (\h,-\h,\h) -- cycle;
  \fill[dielecB, opacity=0.35] (-\h,-\h,\h) -- (\h,-\h,\h) -- (\h,\h,\h) -- (-\h,\h,\h) -- cycle;
  \draw[borde] (-\h,-\h,-\h) -- (\h,-\h,-\h) -- (\h,-\h,\h) -- (-\h,-\h,\h) -- cycle;
  \draw[borde] (\h,-\h,-\h) -- (\h,\h,-\h) -- (\h,\h,\h) -- (\h,-\h,\h);
  \draw[borde] (-\h,-\h,\h) -- (-\h,\h,\h) -- (\h,\h,\h);
  % normal en la cara x = l/2
  \draw[vecn] (\h,0,0) -- (\h+0.9,0,0) node[right, etiq] {$\nn=\uv x$};
  \node[etiq] at (0,-\h,-0.62*\h) {$\rhoP=-3A$};
  \node[etiq, anchor=west] at (\h+0.15,0,0.85*\h) {$\sigmaP=\tfrac12A\ell$};
  % ejes
  \draw[tenue, line width=0.4pt, -{Stealth[length=1.6mm]}] (\h,\h,-\h) ++(0.2,0,0) -- ++(0.7,0,0) node[right, etiqt] {$x$};
  \draw[tenue, line width=0.4pt, -{Stealth[length=1.6mm]}] (\h,\h,-\h) ++(0,0.3,0) -- ++(0,0.9,0) node[right, etiqt] {$y$};
  \draw[cota] (-\h,-\h,-\h-0.35) -- (\h,-\h,-\h-0.35) node[midway, below, etiq] {$\ell$};
\end{tikzpicture}
\caption{\Cref{prob:cubo}. Cubo de lado $\ell$ centrado en el origen con $\vP=A\vr$ ($A>0$). La polarización diverge
uniformemente, $\rhoP=-\divg\vP=-3A$, y en cada cara la componente normal de $\vP$ es constante: por ejemplo,
en $x=\ell/2$, $\sigmaP=\vP\cdot\uv x=A\ell/2$. La carga negativa del volumen compensa exactamente la positiva de las seis caras.}
\label{fig:p45}
\end{figure}


% -----------------------------------------------------------------------------
\subsection{Problema 4.6. Electrete esférico entre conductores}\label{sec:p46}

\begin{problema}[Capa esférica de electrete]\label{prob:electrete-esferico}
El espacio entre dos superficies esféricas conductoras concéntricas de espesor despreciable y radios $a$ y $4a$, respectivamente, está
parcialmente lleno con una capa esférica de electrete de radio interno $a$ y externo $2a$. Dicho electrete presenta una polarización dada
por $\vP=\dfrac{k}{r^2}\,\uv r$, donde $k$ es una constante y $r$ la distancia al centro del sistema. Calcula $\vE$ y $\vD$ para $0<r<\infty$ y
representa gráficamente.
\end{problema}

\begin{resolucion}
\idea Simetría esférica y $\vP$ radial: sin carga libre, $\vD=\vcero$ (\cref{ej:radial}) y $\vE=-\vP/\eo$ (método M2). Lo comprobamos con
Gauss para $\vE$ y las cargas ligadas. Los conductores, neutros, no adquieren carga.

\paso{Paso 1: cargas ligadas} En $a<r<2a$, $\rhoP=-\dfrac{1}{r^2}\dfrac{\dd}{\dd r}\Big(r^2\dfrac{k}{r^2}\Big)=0$. En las superficies del electrete,
\[
  \sigmaP(a)=\vP\cdot(-\uv r)=-\frac{k}{a^2},\qquad \sigmaP(2a)=\vP\cdot\uv r=\frac{k}{4a^2},
\]
con carga total $-\dfrac{k}{a^2}\,4\pi a^2+\dfrac{k}{4a^2}\,4\pi(2a)^2=-4\pi k+4\pi k=0$.

\paso{Paso 2: $\vD$} Por simetría, $\vD=D(r)\,\uv r$, y no hay carga libre en ningún sitio (los conductores son neutros y, como veremos, no se induce
carga en ellos): $D\,4\pi r^2=0$ y $\vD=\vcero$ para todo $r$.

\paso{Paso 3: campo} $\vE=(\vD-\vP)/\eo=-\vP/\eo$:
\begin{equation}\label{eq:p46}
  \recuadro{\vE=\vcero\ (r<a),\qquad \vE=-\frac{k}{\eo r^2}\,\uv r\ \ (a<r<2a),\qquad \vE=\vcero\ (r>2a);\qquad \vD=\vcero .}
\end{equation}
\paso{Paso 4: comprobación con Gauss para $\vE$} Para $a<r<2a$ solo se encierra $\sigmaP(a)$: $E\,4\pi r^2=-4\pi k/\eo$, de donde $E=-k/\eo r^2$. Para $r>2a$ la
carga ligada encerrada es nula y $E=0$.
\end{resolucion}

\begin{observacion}[Potencial y papel de los conductores]\label{obs:p46}
Con $\phi(\infty)=0$, el potencial es nulo para $r\geq2a$ y, dentro del electrete,
$\phi(r)=-\int_{2a}^{r}E_r\,\dd r'=\dfrac{k}{\eo}\Big(\dfrac{1}{2a}-\dfrac1r\Big)$, que en $r=a$ vale $-k/2\eo a$, el potencial de la esfera interior
(\cref{fig:p46}). Como $D_n=\sigmaf=0$ en ambos conductores, en ellos no se induce carga: son irrelevantes. El campo, dirigido hacia dentro (de la
carga ligada positiva exterior a la negativa interior), solo existe dentro del electrete.
\end{observacion}

% Figura: problema 4.6. Capa esférica de electrete P = k/r^2 entre conductores.
\begin{figure}[t]
\centering
\begin{subfigure}[b]{0.36\textwidth}\centering
\begin{tikzpicture}[scale=0.62]
  \draw[texto, line width=0.9pt] (0,0) circle (4.0);
  \fill[dielec] (0,0) circle (2.0);
  \fill[fondo] (0,0) circle (1.0);
  \draw[borde] (0,0) circle (2.0);
  \draw[texto, line width=0.9pt] (0,0) circle (1.0);
  \foreach \t in {0,45,...,315}
    \draw[campoP, line width=0.85pt, -{Stealth[length=1.6mm]}] (\t:1.12) -- (\t:1.88);
  \draw[texto, line width=0.4pt, -{Stealth[length=1.5mm]}] (0,0) -- (115:1.0);
  \node[etiq] at (150:0.45) {$a$};
  \draw[texto, line width=0.4pt, -{Stealth[length=1.5mm]}] (0,0) -- (-65:2.0);
  \node[etiq, fill=dielec, inner sep=1pt] at (-80:1.5) {$2a$};
  \draw[texto, line width=0.4pt, -{Stealth[length=1.5mm]}] (0,0) -- (200:4.0);
  \node[etiq] at (200:3.0) [above=1pt] {$4a$};
  \node[etiq, text=campoP] at (28:2.45) {$\vP$};
\end{tikzpicture}
\caption{Geometría}
\end{subfigure}\hfill
\begin{subfigure}[b]{0.62\textwidth}\centering
\begin{tikzpicture}
\begin{groupplot}[group style={group size=2 by 2, horizontal sep=1.1cm, vertical sep=0.95cm},
   width=5.0cm, height=3.25cm, axis lines=left, xmin=0, xmax=5, xtick={1,2,4}, xticklabels={$a$,$2a$,$4a$},
   tick label style={font=\tiny}, label style={font=\footnotesize}, clip=false, xlabel={$r$}]
  \nextgroupplot[ymin=0, ymax=1.15, ytick={0.25,1}, yticklabels={$\frac{k}{4a^2}$,$\frac{k}{a^2}$}, ylabel={$P$}]
    \addplot[curvaP, solid, domain=1:2, samples=40] {1/x^2};
    \addplot[curvaP, solid] coordinates {(0,0) (1,0)}; \addplot[curvaP, solid] coordinates {(2,0) (5,0)};
    \draw[equipot] (axis cs:1,0) -- (axis cs:1,1.15) (axis cs:2,0) -- (axis cs:2,1.15) (axis cs:4,0) -- (axis cs:4,1.15);
  \nextgroupplot[ymin=-1.15, ymax=0.15, ytick={-1,-0.25,0}, yticklabels={$-\frac{k}{\eo a^2}$,$-\frac{k}{4\eo a^2}$,$0$}, ylabel={$E_r$}]
    \addplot[curvaE, domain=1:2, samples=40] {-1/x^2};
    \addplot[curvaE] coordinates {(0,0) (1,0)}; \addplot[curvaE] coordinates {(2,0) (5,0)};
    \draw[equipot] (axis cs:1,-1.15) -- (axis cs:1,0.15) (axis cs:2,-1.15) -- (axis cs:2,0.15) (axis cs:4,-1.15) -- (axis cs:4,0.15);
  \nextgroupplot[ymin=-0.6, ymax=0.6, ytick={0}, ylabel={$D_r$}]
    \addplot[curvaD] coordinates {(0,0) (5,0)};
    \draw[equipot] (axis cs:1,-0.6) -- (axis cs:1,0.6) (axis cs:2,-0.6) -- (axis cs:2,0.6) (axis cs:4,-0.6) -- (axis cs:4,0.6);
    \node[etiqt, anchor=north, fill=fondo, inner sep=1pt] at (axis cs:3,-0.12) {$\vD=\vcero$ en todo punto};
  \nextgroupplot[ymin=-0.6, ymax=0.1, ytick={-0.5,0}, yticklabels={$-\frac{k}{2\eo a}$,$0$}, ylabel={$\phi$}]
    \addplot[curvaV] coordinates {(0,-0.5) (1,-0.5)};
    \addplot[curvaV, domain=1:2, samples=40] {0.5-1/x};
    \addplot[curvaV] coordinates {(2,0) (5,0)};
    \draw[equipot] (axis cs:1,-0.6) -- (axis cs:1,0.1) (axis cs:2,-0.6) -- (axis cs:2,0.1) (axis cs:4,-0.6) -- (axis cs:4,0.1);
\end{groupplot}
\end{tikzpicture}
\caption{$P$, $E_r$, $D_r$ y $\phi$ frente a $r$}
\end{subfigure}
\caption{\Cref{prob:electrete-esferico}. (a) Capa de electrete ($a<r<2a$, sombreada) con $\vP=(k/r^2)\,\uv r$ entre dos
esferas conductoras neutras de radios $a$ y $4a$. (b) Sin carga libre y con $\vP$ irrotacional, $\vD=\vcero$ en
todo el espacio; el campo solo existe dentro del electrete, $\vE=-\vP/\eo$, y apunta hacia dentro. El
potencial ($\phi(\infty)=0$) es nulo fuera de $r=2a$ y constante dentro de $r=a$.}
\label{fig:p46}
\end{figure}


% -----------------------------------------------------------------------------
\subsection{Problema 4.7. Electrete cilíndrico coaxial}\label{sec:p47}

\begin{problema}[Electrete coaxial]\label{prob:electrete-coaxial}
Sea un electrete limitado por dos superficies cilíndricas coaxiales infinitamente largas de radios $a$ y $b$ ($a<b$). El electrete está
polarizado con un momento dipolar por unidad de volumen $\vP=\dfrac{kr}{2}\,\uv r$, siendo $k$ una constante, $r$ la distancia al eje de
simetría ($a<r<b$) y $\uv r$ el vector unitario en la dirección radial. Calcula el campo electrostático $\vE$ en todo el espacio y el vector
desplazamiento $\vD$ en el material.
\end{problema}

\begin{resolucion}
\idea Mismo esquema que el \cref{prob:electrete-esferico}, con simetría cilíndrica (método M2). Ahora la polarización no es solenoidal y sí hay
carga ligada de volumen. No hay carga libre y $\vP$ es irrotacional (no varía en la dirección perpendicular a $\uv r$), luego el problema es
trivial: $\vD=\vcero$ para todo $\vr$.

\paso{Paso 1: campo} $\vE=-\vP/\eo$ en el material y nulo fuera:
\begin{equation}\label{eq:p47}
  \recuadro{\vE=\vcero\ (r<a),\qquad \vE=-\frac{kr}{2\eo}\,\uv r\ \ (a<r<b),\qquad \vE=\vcero\ (r>b);\qquad \vD=\vcero .}
\end{equation}
\paso{Paso 2: cargas ligadas} En coordenadas cilíndricas, $\rhoP=-\dfrac1r\dfrac{\partial}{\partial r}\Big(r\,\dfrac{kr}{2}\Big)=-k$. En las superficies, con la
normal exterior al electrete,
\[
  \sigmaP(a)=\vP\cdot(-\uv r)=-\frac{ka}{2}\quad(\text{apunta hacia el eje}),\qquad \sigmaP(b)=\vP\cdot\uv r=\frac{kb}{2}.
\]
\paso{Paso 3: comprobación con Gauss para $\vE$} Para un cilindro gaussiano de radio $a<r<b$ y longitud $\ell$,
\[
  \qPenc=-\frac{ka}{2}\,2\pi a\ell-k\,\pi(r^2-a^2)\,\ell=-k\pi r^2\ell
  \quad\Longrightarrow\quad
  E\,2\pi r\ell=-\frac{k\pi r^2\ell}{\eo},\quad \vE=-\frac{kr}{2\eo}\,\uv r .
\]
Para $r>b$, $\qPenc=-k\pi a^2\ell+kb\,\pi b\ell-k\pi(b^2-a^2)\ell=0$: la carga de polarización neta es siempre cero.
\end{resolucion}

\begin{observacion}[Diferencia de potencial]\label{obs:p47}
Aunque fuera del electrete no haya campo, entre sus caras sí hay diferencia de potencial:
$\phi(a)-\phi(b)=\int_a^bE_r\,\dd r=-\dfrac{k}{4\eo}(b^2-a^2)$. La superficie interior está a menor potencial, en acuerdo con la carga ligada
negativa que la recubre (\cref{fig:p47}).
\end{observacion}

% Figura: problema 4.7. Electrete cilíndrico coaxial con P = (k r/2) u_r.
\begin{figure}[t]
\centering
\begin{subfigure}[b]{0.27\textwidth}\centering
\begin{tikzpicture}[scale=0.8]
  % cilindro hueco en perspectiva
  \fill[dielec] (-1.3,0) -- (-1.3,-3.0) arc (180:360:1.3 and 0.38) -- (1.3,0) arc (0:180:1.3 and 0.38);
  \fill[dielec] (-1.3,0) arc (180:360:1.3 and 0.38) -- (1.3,0) arc (0:180:1.3 and 0.38);
  \fill[dielecB] (0,0) ellipse (1.3 and 0.38);
  \fill[fondo] (0,0) ellipse (0.55 and 0.16);
  \draw[borde] (0,0) ellipse (1.3 and 0.38) (0,0) ellipse (0.55 and 0.16);
  \draw[borde] (-1.3,0) -- (-1.3,-3.0) arc (180:360:1.3 and 0.38) -- (1.3,0);
  \draw[texto!50!fondo, dashed, line width=0.4pt] (-0.55,0) -- (-0.55,-3.0) (0.55,0) -- (0.55,-3.0);
  \draw[texto!50!fondo, dashed, line width=0.4pt] (-0.55,-3.0) arc (180:360:0.55 and 0.16);
  \foreach \t in {20,70,...,340}
    \draw[campoP, line width=0.7pt, -{Stealth[length=1.3mm]}] ({0.6*cos(\t)},{0.175*sin(\t)}) -- ({1.22*cos(\t)},{0.355*sin(\t)});
  \draw[cota] (1.65,0) -- (1.65,-3.0) node[midway, right, etiq] {$\ell$};
  \node[etiq] at (-0.25,0.52) {$a$};
  \node[etiq] at (1.0,0.62) {$b$};
\end{tikzpicture}
\caption{Electrete coaxial}
\end{subfigure}\hfill
\begin{subfigure}[b]{0.27\textwidth}\centering
\begin{tikzpicture}[scale=0.8]
  \fill[dielec] (0,0) circle (1.7);
  \fill[fondo] (0,0) circle (0.8);
  \draw[borde] (0,0) circle (1.7) (0,0) circle (0.8);
  \foreach \t in {0,45,...,315}
    \draw[campoP, line width=0.8pt, -{Stealth[length=1.5mm]}] (\t:0.9) -- (\t:{0.9+0.62*1.6/1.7});
  \draw[vecn] (135:0.8) -- (135:0.3) node[right=-1pt, etiq] {$\nn$};
  \draw[vecn] (1.7,0) -- (2.3,0) node[above, etiq] {$\nn$};
  \node[etiq, text=campoP] at (62:2.0) {$\vP$};
  \node[etiq, font=\footnotesize] at (0,-0.38) {$\vE=\vcero$};
\end{tikzpicture}
\caption{Corte transversal}
\end{subfigure}\hfill
\begin{subfigure}[b]{0.42\textwidth}\centering
\begin{tikzpicture}
\begin{axis}[width=5.8cm, height=4.3cm, axis lines=middle, xmin=0, xmax=3.1, ymin=-1.15, ymax=1.15,
  xtick={1,2}, xticklabels={$a$,$b$}, ytick=\empty, xlabel={$r$},
  xlabel style={at={(ticklabel* cs:1)}, anchor=north east}, clip=false,
  legend style={at={(0.5,-0.08)}, anchor=north, legend columns=2, font=\tiny, /tikz/every even column/.append style={column sep=4pt}}]
  \addplot[curvaP, domain=1:2, samples=2] {x/2}; \addlegendentry{$P=kr/2$}
  \addplot[curvaE, domain=1:2, samples=2] {-x/2}; \addlegendentry{$\eo E_r=-kr/2$}
  \addplot[curvaD] coordinates {(0,0.02) (3.1,0.02)}; \addlegendentry{$D_r=0$}
  \addplot[curvaP, forget plot] coordinates {(1,0) (1,0.5)};
  \addplot[curvaP, forget plot] coordinates {(2,0) (2,1)};
  \draw[equipot] (axis cs:1,-1.15) -- (axis cs:1,1.15) (axis cs:2,-1.15) -- (axis cs:2,1.15);
  \node[etiqt, anchor=west] at (axis cs:2.05,1.0) {$\sigmaP(b)=\tfrac{kb}{2}$};
  \node[etiqt, anchor=west] at (axis cs:1.03,-0.22) {$\sigmaP(a)=-\tfrac{ka}{2}$};
  \node[etiqt] at (axis cs:1.32,-1.05) {$\rhoP=-k$};
\end{axis}
\end{tikzpicture}
\caption{Perfiles radiales}
\end{subfigure}
\caption{\Cref{prob:electrete-coaxial}. (a), (b) Electrete limitado por dos cilindros coaxiales de radios $a<b$, con polarización
radial creciente. (c) Hay carga ligada de volumen uniforme, $\rhoP=-k$, y de superficie en ambas caras;
su total por unidad de longitud es nulo. Como no hay carga libre y $\vP$ es irrotacional, $\vD=\vcero$ y
$\vE=-\vP/\eo$ dentro del material; fuera de él, el campo es nulo.}
\label{fig:p47}
\end{figure}


% -----------------------------------------------------------------------------
\subsection{Problema 4.8. Condensador esférico con dos dieléctricos}\label{sec:p48}

\begin{problema}[Dos dieléctricos hemisféricos]\label{prob:dos-hemisferios}
Sean dos esferas conductoras concéntricas de radios $a$ y $b$ separadas por dos dieléctricos de permitividades $\varepsilon_1$ y $\varepsilon_2$ cuya superficie
de contacto es un plano ecuatorial (\cref{fig:p48}). Se le proporciona una carga $Q$ a la esfera interior y se conecta la esfera exterior a potencial
nulo. Demuestra que la expresión para el potencial es la misma en las dos zonas del recinto entre las dos esferas. Calcula dicha solución y la
capacidad del sistema. ¿Cómo se distribuye la carga $Q$ en la esfera interior?
\end{problema}

\begin{resolucion}
\idea Probamos un potencial que solo dependa de $r$ en ambas zonas y comprobamos que satisface todas las condiciones; por unicidad, es la
solución. Después, Gauss para $\vD$ (métodos M3 y M4).

\paso{Paso 1: el mismo potencial en ambas zonas} Busquemos $\phi=\phi(r)$ en $a<r<b$. Entonces $\vE=E(r)\,\uv r$ es radial y, por tanto, \emph{tangente} al
plano ecuatorial que separa los medios. Las condiciones de frontera en ese plano son:
\begin{itemize}
  \item componente tangencial de $\vE$ continua: $E_1(r)=E_2(r)\equiv E(r)$ en cada punto del plano, es decir, la misma función en ambos medios;
  \item componente normal de $\vD$ continua: $\vD_i=\varepsilon_iE\,\uv r$ no tiene componente normal al plano, y $0=0$ se cumple trivialmente.
\end{itemize}
Además, en cada medio $\phi$ debe ser armónico, y las soluciones radiales de $\lap\phi=0$ son $\phi=A/r+B$; las dos esferas son equipotenciales y
$\phi(b)=0$. Una única función $\phi(r)=A/r+B$ satisface todo: por el teorema de unicidad~\cite{T3}, la expresión del potencial es la misma en las
dos zonas.

\paso{Paso 2: $\vD$ y $\vE$ por Gauss} Con una esfera de radio $a<r<b$, la mitad en cada medio,
\[
  \oiint\vD\cdot\dS=\varepsilon_1E\,2\pi r^2+\varepsilon_2E\,2\pi r^2=Q
  \quad\Longrightarrow\quad
  \vE=\frac{Q}{2\pi(\varepsilon_1+\varepsilon_2)\,r^2}\,\uv r .
\]
\paso{Paso 3: potencial y capacidad} Con $\phi(b)=0$,
\[
  \recuadro{\phi(r)=-\int_b^r E\,\dd r'=\frac{Q}{2\pi(\varepsilon_1+\varepsilon_2)}\Big(\frac1r-\frac1b\Big),\qquad
  C=\frac{Q}{\phi(a)}=2\pi(\varepsilon_1+\varepsilon_2)\,\frac{ab}{b-a}.}
\]
\paso{Paso 4: reparto de la carga} En la superficie del conductor interior, $\sigmaf=D_n=\varepsilon_iE(a)$ en cada hemisferio:
\[
  \recuadro{\sigma_1=\frac{\varepsilon_1\,Q}{2\pi a^2(\varepsilon_1+\varepsilon_2)},\qquad \sigma_2=\frac{\varepsilon_2\,Q}{2\pi a^2(\varepsilon_1+\varepsilon_2)} .}
\]
Comprobación: $2\pi a^2(\sigma_1+\sigma_2)=Q$.
\end{resolucion}

\begin{observacion}[Dos condensadores en paralelo y la carga total]\label{obs:p48}
La capacidad es la suma $C=C_1+C_2$, con $C_i=2\pi\varepsilon_i\,ab/(b-a)$ la de medio condensador esférico lleno del dieléctrico $i$: los dos hemisferios
están al mismo potencial y se reparten la carga, como dos condensadores en paralelo. La carga libre se concentra en el hemisferio de mayor
permitividad, pero la carga ligada adyacente, $\sigma_{\mathrm{P},i}=-(\varepsilon_i-\eo)E(a)$, compensa justo esa diferencia: la carga total,
$\sigma_i+\sigma_{\mathrm{P},i}=\eo E(a)=Q\eo/[2\pi a^2(\varepsilon_1+\varepsilon_2)]$, es uniforme. Por eso el campo es radial y con simetría esférica, aunque el sistema no la tenga
(\cref{fig:p48}\,(b)). El problema se resolvió en los seminarios con exactamente estos pasos.
\end{observacion}

% Figura: problema 4.8. Condensador esférico con dos dieléctricos hemisféricos.
\begin{figure}[t]
\centering
\begin{subfigure}[b]{0.5\textwidth}\centering
\begin{tikzpicture}[scale=1.0]
  \fill[dielec] (0,2.1) arc (90:270:2.1) -- cycle;
  \fill[dielecB] (0,-2.1) arc (-90:90:2.1) -- cycle;
  \draw[texto!70!fondo, line width=0.5pt] (0,-2.1) -- (0,2.1);
  \draw[texto, line width=1.1pt] (0,0) circle (2.1);
  \shade[ball color=conductorluz] (0,0) circle (0.85);
  \draw[texto, line width=0.6pt] (0,0) circle (0.85);
  \draw[gauss] (0,0) circle (1.5);
  \node[etiq] at (-0.9,1.62) {$\varepsilon_1$};
  \node[etiq] at (0.9,1.62) {$\varepsilon_2$};
  \draw[texto, line width=0.4pt, -{Stealth[length=1.6mm]}] (0,0) -- (40:0.85);
  \node[font=\small, text=black] at (0.12,0.12) [above left=-1pt] {$a$};
  \draw[acentoB, line width=0.5pt, -{Stealth[length=1.6mm]}] (0,0) -- (-25:1.5);
  \node[etiq, text=acentoB] at (-20:1.25) [below=1pt] {$r$};
  \draw[texto, line width=0.4pt, -{Stealth[length=1.6mm]}] (0,0) -- (-125:2.1);
  \node[etiq] at (-125:1.75) [left=1pt] {$b$};
  \node[etiq] at (-1.55,0.25) {$Q$};
  \draw[texto, line width=0.4pt, -{Stealth[length=1.5mm]}] (-1.38,0.18) -- (165:0.9);
  % tierra
  \draw[texto, line width=0.6pt] (2.1,0) -- (2.75,0) -- (2.75,-0.25);
  \pic at (2.75,-0.25) {tierra};
  \node[etiqt, anchor=west] at (2.82,0.15) {$\phi=0$};
\end{tikzpicture}
\caption{Geometría y superficie gaussiana}
\end{subfigure}\hfill
\begin{subfigure}[b]{0.47\textwidth}\centering
\begin{tikzpicture}[scale=0.42]
  \newcommand{\mini}[1]{\begin{scope}[shift={(0,#1)}]
     \fill[dielec] (0,2.1) arc (90:270:2.1) -- cycle;
     \fill[dielecB] (0,-2.1) arc (-90:90:2.1) -- cycle;
     \draw[texto, line width=0.6pt] (0,0) circle (2.1);
     \fill[fondo] (0,0) circle (0.95);\end{scope}}
  % (i) interfase: sin carga
  \mini{5.6}
  \draw[acentoC, line width=1.6pt] (0,0.95+5.6) -- (0,2.1+5.6) (0,-0.95+5.6) -- (0,-2.1+5.6);
  \draw[texto, line width=0.5pt] (0,5.6) circle (0.95);
  \node[etiq, anchor=west, align=left, font=\footnotesize] at (2.5,5.6) {interfase plana: $E_n=D_n=0$,\\ sin carga libre ni ligada};
  % (ii) carga total uniforme
  \mini{0}
  \draw[acentoA, line width=2.0pt] (0,0) circle (0.95);
  \node[etiq, anchor=west, align=left, font=\footnotesize] at (2.5,0) {carga total en $r=a$:\\ $\sigmaf+\sigmaP=\eo E(a)$, uniforme};
  % (iii) carga libre mayor en el medio de mayor permitividad
  \mini{-5.6}
  \draw[cargapos, line width=1.0pt] (0,0.95-5.6) arc (90:270:0.95);
  \draw[cargapos, line width=2.6pt] (0,-0.95-5.6) arc (-90:90:0.95);
  \node[etiq, anchor=west, align=left, font=\footnotesize] at (2.5,-5.6) {carga libre: $\sigma_{\mathrm{f},i}=\varepsilon_i E(a)$,\\ mayor donde $\varepsilon$ es mayor};
\end{tikzpicture}
\caption{Reparto de la carga en $r=a$ ($\varepsilon_2>\varepsilon_1$)}
\end{subfigure}
\caption{\Cref{prob:dos-hemisferios}. (a) Esfera interior de radio $a$ y carga $Q$; esfera exterior de radio $b$ a tierra; el hueco
está lleno por dos dieléctricos separados por un plano ecuatorial. (b) El campo es radial e igual en ambos
medios, pero la carga libre de la esfera interior no es uniforme: se concentra en el hemisferio del
dieléctrico de mayor permitividad, de modo que la carga total (libre más ligada) sí resulta uniforme.}
\label{fig:p48}
\end{figure}


% -----------------------------------------------------------------------------
\subsection{Problema 4.9. Electrete en un condensador cortocircuitado}\label{sec:p49}

\begin{problema}[Electrete entre armaduras unidas]\label{prob:electrete-condensador}
Una plancha de electrete con polarización $\vP$ uniforme se coloca entre las armaduras de un condensador plano como se muestra en la
\cref{fig:p49}. Calcula los campos dentro y fuera del electrete cuando las armaduras del condensador se unen mediante un cable conductor.
Desprecia el efecto de bordes.
\end{problema}

\begin{resolucion}
\idea La polarización es uniforme e independiente de cualquier campo. El electrete tiene cargas ligadas $\pm P$ en sus caras, que inducen
carga en las armaduras; al estar unidas por el cable, sus potenciales son iguales. Sea $a$ el espesor del electrete y $b$ el del hueco de aire, y
$\uv z$ el sentido de $\vP$, normal a las placas (método M5).

\paso{Paso 1: $\vD$ es el mismo en ambas capas} Por la simetría plana, los campos son normales a las placas y solo dependen de $z$. En cada capa
no hay carga libre y $\dd D/\dd z=0$; en la interfase tampoco, y $D_n$ es continua. Luego $\vD=D\,\uv z$ con $D$ constante en todo el hueco.

\paso{Paso 2: los campos en función de $D$}
\[
  \eo\vEext=\vD\quad(\text{aire}),\qquad \eo\vEint=\vD-\vP\quad(\text{electrete}).
\]
\paso{Paso 3: la condición del cable} Las armaduras están al mismo potencial:
\[
  \phi_{\text{sup}}-\phi_{\text{inf}}=-\int_0^{a+b}E_z\,\dd z=-\big(\Eint\,a+\Eext\,b\big)=0 ,
\]
y sustituyendo los campos del paso~2,
\[
  (D-P)\,a+D\,b=0\quad\Longrightarrow\quad D=\frac{a}{a+b}\,P .
\]
\paso{Paso 4: resultado}
\begin{equation}\label{eq:p49}
  \recuadro{\vEint=-\frac{b}{\eo(a+b)}\,\vP,\qquad \vEext=\frac{a}{\eo(a+b)}\,\vP,\qquad \vD=\frac{a}{a+b}\,\vP .}
\end{equation}
Las cargas inducidas en las armaduras son $\sigma_{\text{inf}}=+D=\frac{a}{a+b}P$ en la inferior (junto a la cara negativa del electrete) y
$\sigma_{\text{sup}}=-D$ en la superior.
\end{resolucion}

\begin{observacion}[Comprobación por superposición de planos]\label{obs:p49}
Hay cuatro planos cargados: $+s$ (armadura inferior), $-P$ y $+P$ (caras del electrete) y $-s$ (armadura superior), con $s=\frac{a}{a+b}P$. Cada uno crea
$\sigma/2\eo$ a cada lado, alejándose si es positivo. En el aire, los tres planos inferiores suman $s/2\eo$ hacia arriba y la armadura superior, que
atrae, otros $s/2\eo$: $\Eext=s/\eo=aP/\eo(a+b)$. En el electrete, $\frac{s-P}{2\eo}-\frac{P}{2\eo}+\frac{s}{2\eo}=\frac{s-P}{\eo}=-\frac{bP}{\eo(a+b)}$. Coincide con \eqref{eq:p49}, como en
la resolución de los seminarios. Sin el cable ($s=0$) se recupera la lámina polarizada aislada, $\Eext=0$ y $\Eint=-P/\eo$; con el hueco cerrado ($b\to0$),
$\Eint\to0$: el electrete queda cortocircuitado por las placas.
\end{observacion}

% Figura: problema 4.9. Lámina de electrete entre las armaduras de un condensador
% unidas por un cable; perfiles de E, D, P y phi (a = 1, b = 2).
\begin{figure}[t]
\centering
\begin{subfigure}[b]{0.52\textwidth}\centering
\begin{tikzpicture}[scale=1.0]
  \def\a{1.0} \def\b{2.0} \def\L{5.0}
  \fill[dielec] (0,0) rectangle (\L,\a);
  \draw[borde] (0,0) rectangle (\L,\a);
  \fill[conductor] (0,-0.14) rectangle (\L,0);
  \fill[conductor] (0,\a+\b) rectangle (\L,\a+\b+0.14);
  \draw[texto, line width=0.5pt] (0,-0.14) rectangle (\L,0) (0,\a+\b) rectangle (\L,\a+\b+0.14);
  % cable
  \draw[texto, line width=0.7pt, rounded corners=6pt] (\L,\a+\b+0.07) -- (\L+0.7,\a+\b+0.07) -- (\L+0.7,-0.07) -- (\L,-0.07);
  % cargas
  \foreach \x in {0.3,0.95,...,4.4} {
    \node[cpos, font=\tiny\bfseries] at (\x,-0.07) {$+$};
    \node[cneg, font=\tiny\bfseries] at (\x,\a+\b+0.07) {$-$};
    \node[cpos, font=\tiny\bfseries] at (\x+0.32,\a-0.13) {$+$};
    \node[cneg, font=\tiny\bfseries] at (\x+0.32,0.13) {$-$};
  }
  % vectores
  \draw[vecP] (0.65,0.22) -- (0.65,0.8); \node[etiqm, font=\footnotesize, anchor=west] at (0.7,0.5) {$\vP$};
  \draw[vecE] (2.2,0.8) -- (2.2,0.32); \node[etiqm, font=\footnotesize, anchor=west] at (2.25,0.52) {$\vEint$};
  \draw[vecE] (2.2,\a+0.4) -- (2.2,\a+0.4+0.6); \node[etiq, font=\footnotesize, anchor=west] at (2.25,\a+0.75) {$\vEext$};
  \draw[vecD] (3.7,\a+0.4) -- (3.7,\a+1.0); \node[etiq, font=\footnotesize, anchor=west] at (3.75,\a+0.75) {$\vD$};
  \draw[vecD] (3.7,0.25) -- (3.7,0.85); \node[etiqm, font=\footnotesize, anchor=west] at (3.75,0.55) {$\vD$};
  % cotas
  \draw[cota] (-0.35,0) -- (-0.35,\a) node[midway, left, etiq] {$a$};
  \draw[cota] (-0.35,\a) -- (-0.35,\a+\b) node[midway, left, etiq] {$b$};
  \node[etiqt, anchor=south] at (\L/2,\a+\b+0.16) {$\sigma_{\text{sup}}=-\tfrac{a}{a+b}P$};
  \node[etiqt, anchor=north] at (\L/2,-0.17) {$\sigma_{\text{inf}}=+\tfrac{a}{a+b}P$};
\end{tikzpicture}
\caption{Geometría y cargas}
\end{subfigure}\hfill
\begin{subfigure}[b]{0.45\textwidth}\centering
\begin{tikzpicture}
\begin{axis}[width=6.3cm, height=5.6cm, xmin=-0.9, xmax=1.15, ymin=0, ymax=3,
  ytick={0,1,3}, yticklabels={$0$,$a$,$a+b$}, ylabel={altura}, xtick={-0.6667,0,0.3333,0.6667,1},
  xticklabels={$-\frac23$,$0$,$\frac13$,$\frac23$,$1$}, grid=major,
  legend style={at={(0.5,1.03)}, anchor=south, legend columns=2, font=\tiny, /tikz/every even column/.append style={column sep=3pt}}]
  \fill[dielec, opacity=0.6] (axis cs:-0.9,0) rectangle (axis cs:1.15,1);
  \addplot[curvaP] coordinates {(1,0) (1,1) (0,1) (0,3)}; \addlegendentry{$P$}
  \addplot[curvaD, line width=2.8pt, opacity=0.45] coordinates {(0.3333,0) (0.3333,3)}; \addlegendentry{$D$}
  \addplot[curvaE] coordinates {(-0.6667,0) (-0.6667,1)}; \addlegendentry{$\eo E$}
  \addplot[curvaV] coordinates {(0,0) (0.6667,1) (0,3)}; \addlegendentry{$\phi\,\eo/(Pa)$}
  \addplot[curvaE, forget plot] coordinates {(0.3333,1) (0.3333,3)};
  \addplot[curvaE, forget plot, dotted, line width=0.5pt] coordinates {(-0.6667,1) (0.3333,1)};
\end{axis}
\end{tikzpicture}
\caption{Perfiles ($b=2a$, unidades de $P$)}
\end{subfigure}
\caption{\Cref{prob:electrete-condensador}. (a) Electrete de espesor $a$ y polarización uniforme $\vP$ sobre la armadura inferior; las
armaduras están unidas por un cable, de modo que $\phi_{\text{sup}}=\phi_{\text{inf}}$. Sin carga libre en la interfase, $\vD$ es
el mismo en el electrete y en el aire, y la condición $\oint\vE\cdot\dl=0$ obliga a que $\vE$ cambie de sentido.
(b) Con $b=2a$: $\eo\Eint=-\tfrac23P$, $\eo\Eext=\tfrac13P$, $D=\tfrac13P$; el potencial sube en el electrete y baja en el aire hasta
volver al valor inicial. En el aire $D=\eo\Eext$, y ambas curvas coinciden.}
\label{fig:p49}
\end{figure}


% -----------------------------------------------------------------------------
\subsection{Problema 4.10. Dieléctricos en paralelo}\label{sec:p410}

\begin{problema}[Dos dieléctricos en paralelo]\label{prob:paralelo}
Entre las placas paralelas de un condensador plano al que aplicamos una diferencia de potencial $V_0$ se introducen dos dieléctricos de
permitividades relativas $\varepsilon_{r1}$ y $\varepsilon_{r2}$ como se indica en la \cref{fig:p410-412}\,(a). Calcula la carga en cada placa y la capacidad del
sistema. Desprecia el efecto de bordes.
\end{problema}

\begin{resolucion}
\idea La interfase entre los dieléctricos es perpendicular a las placas, de modo que los campos son tangentes a ella (método M4, paralelo).

\paso{Paso 1: el campo es común} La componente tangencial de $\vE$ es continua, $E_1=E_2\equiv E$. Como las placas son equipotenciales y separadas $d$,
\[
  E_1=E_2=E=\frac{V_0}{d}.
\]
\paso{Paso 2: carga en cada zona} En la superficie de la placa, $\sigmaf=D_n$: $\sigma_i=D_i=\varepsilon_{ri}\eo E=\varepsilon_{ri}\eo V_0/d$, y con las áreas $A_1$ y $A_2$,
\[
  \recuadro{Q_1=V_0\,\frac{A_1\varepsilon_{r1}\eo}{d},\qquad Q_2=V_0\,\frac{A_2\varepsilon_{r2}\eo}{d},}
\]
y la capacidad es la suma de las de ambos tramos,
\[
  \recuadro{C=\frac{Q_1+Q_2}{V_0}=\frac{A_1\varepsilon_{r1}\eo}{d}+\frac{A_2\varepsilon_{r2}\eo}{d}=C_1+C_2 .}
\]
\end{resolucion}

\begin{observacion}[Paralelo]\label{obs:p410}
Los dos tramos comparten la diferencia de potencial y se reparten la carga: son dos condensadores en paralelo. La solución de campo
uniforme cumple todas las condiciones (placas equipotenciales, $\vE_t$ continuo y $D_n$ continua, trivialmente, en la interfase), y es única. El
mismo argumento resolvió el \cref{prob:dos-hemisferios}.
\end{observacion}

% -----------------------------------------------------------------------------
\subsection{Problema 4.11. Dieléctrico en serie con aire}\label{sec:p411}

\begin{problema}[Condensador parcialmente lleno]\label{prob:serie}
Considera un condensador de placas paralelas de área $A$ y separación $d$. Entre sus placas se aplica una diferencia de potencial $V_0$ y se
introduce un dieléctrico de permitividad $\varepsilon$ que llena parcialmente el espacio entre placas, como se indica en la \cref{fig:p410-412}\,(b).
Calcula los campos dentro y fuera del dieléctrico y la capacidad del sistema. Desprecia el efecto de bordes.
\end{problema}

\begin{resolucion}
\idea Ahora la interfase es paralela a las placas y los campos son normales a ella (método M4, serie). Sean $a$ el espesor del dieléctrico y $b=d-a$ el
del aire.

\paso{Paso 1: $\vD$ es común} Sin carga libre en la interfase, $D_n$ es continua: $\eo\Eext=\varepsilon\Eint\equiv D$.

\paso{Paso 2: la batería} $V_0=\Eint\,a+\Eext\,b=\Eint\Big(a+\dfrac{\varepsilon}{\eo}\,b\Big)=\dfrac{\Eint}{\eo}\big(\eo a+\varepsilon b\big)$, de donde
\[
  \recuadro{\Eint=V_0\,\frac{\eo}{a\eo+b\varepsilon},\qquad \Eext=V_0\,\frac{\varepsilon}{a\eo+b\varepsilon},\qquad
  C=\frac{DA}{V_0}=A\,\frac{\varepsilon\eo}{a\eo+b\varepsilon} .}
\]
\paso{Paso 3: comprobación como serie} Con $C_1=\eo A/b$ (aire) y $C_2=\varepsilon A/a$ (dieléctrico),
\[
  \frac{C_1C_2}{C_1+C_2}=\frac{A^2\eo\varepsilon/(ab)}{A(\eo a+\varepsilon b)/(ab)}=A\,\frac{\varepsilon\eo}{a\eo+b\varepsilon}. \qedhere
\]
\end{resolucion}

\begin{observacion}[Límites]\label{obs:p411}
Si $\varepsilon\to\eo$, $C\to\eo A/(a+b)=\eo A/d$. Si $\varepsilon\to\infty$, $\Eint\to0$ y $C\to\eo A/b$: el dieléctrico se comporta como un conductor y la separación efectiva se
reduce al hueco de aire, que es el \cref{prob:lamina}. Casi toda la caída de potencial se produce en el aire, donde el campo es $\er$ veces mayor
(\cref{fig:p411-perfil}\,(a)): por eso una capa de aire en serie con un aislante es el punto débil frente a la ruptura dieléctrica.
\end{observacion}

% -----------------------------------------------------------------------------
\subsection{Problema 4.12. Lámina conductora intercalada}\label{sec:p412}

\begin{problema}[Lámina conductora]\label{prob:lamina}
Entre las placas de un condensador plano-paralelo de área $A$ y separación $d$ se introduce una lámina conductora de grosor $e$ como se indica en la
\cref{fig:p410-412}\,(c). ¿Cómo varía la capacidad del sistema?
\end{problema}

\begin{resolucion}
\idea Dentro del conductor $\vE=\vcero$; en sus caras se inducen cargas $\mp Q$ que dejan inalterado el campo de los huecos (método M4, serie).

\paso{Paso 1: campos} Con cargas $\pm Q$ en las placas, la lámina adquiere $-Q$ en la cara frente a la placa positiva y $+Q$ en la otra (regla de las caras del
Tema~3~\cite{T3}). En cada hueco, $D=\sigma=Q/A$ y $E=\sigma/\eo$; en el metal, $E=0$.

\paso{Paso 2: diferencia de potencial} Con $a$ y $b$ los espesores de los huecos, $a+b=d-e$:
\[
  \Delta\phi=\frac{\sigma}{\eo}\,a+0\cdot e+\frac{\sigma}{\eo}\,b=\frac{\sigma}{\eo}(d-e)
  \quad\Longrightarrow\quad
  \recuadro{C'=\frac{Q}{\Delta\phi}=\frac{\eo A}{d-e}>\frac{\eo A}{d}=C .}
\]
\end{resolucion}

\begin{observacion}[Independencia de la posición]\label{obs:p412}
La capacidad aumenta y no depende de dónde se coloque la lámina, solo de su espesor: equivale a dos condensadores en serie, $\eo A/a$ y $\eo A/b$, cuya
serie solo depende de $a+b$. Es el límite $\varepsilon\to\infty$ del \cref{prob:serie} (\cref{fig:p411-perfil}\,(b)).
\end{observacion}

% Figura: problemas 4.10, 4.11 y 4.12. Condensadores planos con dieléctricos
% «en paralelo», «en serie» y con una lámina conductora intercalada.
\begin{figure}[t]
\centering
\begin{subfigure}[b]{0.3\textwidth}\centering
\begin{tikzpicture}[scale=0.8]
  \fill[dielecB] (0,1.7) rectangle (2.0,3.4);
  \fill[dielec] (0,0) rectangle (2.0,1.7);
  \draw[texto!70!fondo, line width=0.5pt] (0,1.7) -- (2.0,1.7);
  \draw[placa] (0,-0.1) -- (0,3.5);
  \draw[placa] (2.0,-0.1) -- (2.0,3.5);
  \draw[cargapos, line width=1.6pt] (-0.06,0) -- (-0.06,3.4);
  \node[etiq] at (1.0,3.05) {$\varepsilon_{r1}$};
  \node[etiq] at (1.0,1.35) {$\varepsilon_{r2}$};
  \draw[vecE] (0.45,2.35) -- (1.55,2.35);
  \draw[vecE] (0.45,0.65) -- (1.55,0.65);
  \node[etiq, anchor=east] at (-0.15,2.95) {$A_1$};
  \node[etiq, anchor=east] at (-0.15,2.3) {$Q_1$};
  \node[etiq, anchor=east] at (-0.15,1.25) {$A_2$};
  \node[etiq, anchor=east] at (-0.15,0.6) {$Q_2$};
  \draw[cota] (0,3.8) -- (2.0,3.8) node[midway, above, etiq] {$d$};
  % batería
  \draw[texto, line width=0.6pt] (0,-0.1) -- (0,-0.95) -- (0.8,-0.95);
  \draw[texto, line width=0.6pt] (1.2,-0.95) -- (2.0,-0.95) -- (2.0,-0.1);
  \draw[texto, line width=1.1pt] (0.85,-0.72) -- (0.85,-1.18);
  \draw[texto, line width=0.6pt] (1.12,-0.82) -- (1.12,-1.08);
  \node[etiq] at (1.0,-1.45) {$V_0$};
\end{tikzpicture}
\caption{Problema 4.10}
\end{subfigure}\hfill
\begin{subfigure}[b]{0.34\textwidth}\centering
\begin{tikzpicture}[scale=0.8]
  \fill[dielec] (0,0) rectangle (3.2,1.2);
  \draw[texto!70!fondo, line width=0.5pt] (0,1.2) -- (3.2,1.2);
  \draw[placa] (-0.1,0) -- (3.3,0);
  \draw[placa] (-0.1,2.6) -- (3.3,2.6);
  \node[etiq] at (2.95,0.6) {$\varepsilon$};
  \node[etiq] at (2.95,1.9) {$\eo$};
  \draw[vecD] (0.6,0.15) -- (0.6,1.05); \draw[vecD] (0.6,1.45) -- (0.6,2.35);
  \node[etiq, anchor=west, font=\footnotesize] at (0.68,1.9) {$\vD$};
  \draw[vecE] (1.6,0.4) -- (1.6,0.8); \draw[vecE] (1.6,1.35) -- (1.6,2.45);
  \node[etiq, anchor=west, font=\footnotesize] at (1.68,1.9) {$\vEext$};
  \node[etiqm, anchor=west, font=\footnotesize] at (1.68,0.6) {$\vEint$};
  \draw[cota] (-0.35,0) -- (-0.35,1.2) node[midway, left, etiq] {$a$};
  \draw[cota] (-0.35,1.2) -- (-0.35,2.6) node[midway, left, etiq] {$b$};
  % batería a la derecha
  \draw[texto, line width=0.6pt] (3.3,2.6) -- (3.95,2.6) -- (3.95,1.5);
  \draw[texto, line width=0.6pt] (3.95,1.1) -- (3.95,0) -- (3.3,0);
  \draw[texto, line width=1.1pt] (3.72,1.45) -- (4.18,1.45);
  \draw[texto, line width=0.6pt] (3.82,1.18) -- (4.08,1.18);
  \node[etiq, anchor=west] at (4.2,1.3) {$V_0$};
\end{tikzpicture}
\caption{Problema 4.11}
\end{subfigure}\hfill
\begin{subfigure}[b]{0.32\textwidth}\centering
\begin{tikzpicture}[scale=0.8]
  \fill[conductor] (0,1.2) rectangle (3.2,1.65);
  \draw[texto, line width=0.5pt] (0,1.2) rectangle (3.2,1.65);
  \draw[placa] (-0.1,0) -- (3.3,0);
  \draw[placa] (-0.1,3.0) -- (3.3,3.0);
  \node[etiq] at (0.45,0.6) {$\eo$};
  \node[etiq] at (0.45,2.33) {$\eo$};
  \node[etiq, anchor=west] at (3.35,3.0) {$+Q$};
  \node[etiq, anchor=west] at (3.35,1.73) {$-Q$};
  \node[etiq, anchor=west] at (3.35,1.12) {$+Q$};
  \node[etiq, anchor=west] at (3.35,0) {$-Q$};
  \draw[vecE] (1.2,2.85) -- (1.2,1.8); \draw[vecE] (1.2,1.05) -- (1.2,0.15);
  \node[etiq, anchor=west, font=\footnotesize] at (1.28,2.33) {$\sigma/\eo$};
  \node[etiq, anchor=west, font=\footnotesize] at (1.28,0.6) {$\sigma/\eo$};
  \node[font=\footnotesize, text=black] at (1.9,1.425) {$\vE=\vcero$};
  \draw[cota] (-0.35,0) -- (-0.35,1.2) node[midway, left, etiq] {$a$};
  \draw[cota] (-0.35,1.2) -- (-0.35,1.65) node[midway, left, etiq] {$e$};
  \draw[cota] (-0.35,1.65) -- (-0.35,3.0) node[midway, left, etiq] {$b$};
\end{tikzpicture}
\caption{Problema 4.12}
\end{subfigure}
\caption{Condensadores planos de área $A$ y separación $d$ parcialmente llenos. (a) Dos dieléctricos con la interfase
perpendicular a las placas: el campo es el mismo en ambos ($V_0/d$), pero la carga libre se reparte en
proporción a $\varepsilon_{ri}A_i$; el sistema equivale a dos condensadores en paralelo. (b) Interfase paralela a las
placas: $\vD$ es el mismo en ambas capas y el campo es menor en el dieléctrico; equivale a dos condensadores en
serie. (c) Una lámina conductora de espesor $e$ anula el campo en su interior y reduce la separación efectiva
a $d-e$, cualquiera que sea su posición.}
\label{fig:p410-412}
\end{figure}

% Figura: problemas 4.11 y 4.12. Perfiles de campo y potencial frente a la altura.
\begin{figure}[t]
\centering
\begin{subfigure}[b]{0.48\textwidth}\centering
\begin{tikzpicture}
\begin{axis}[width=6.6cm, height=5.4cm, xmin=0, xmax=1.65, ymin=0, ymax=2,
  ytick={0,1,2}, yticklabels={$0$,$a$,$d$}, ylabel={altura}, xtick={0,0.5,1,1.5},
  xticklabels={$0$,$\frac12$,$1$,$\frac32$}, grid=major,
  legend style={at={(0.5,1.03)}, anchor=south, legend columns=3, font=\tiny, /tikz/every even column/.append style={column sep=3pt}}]
  \fill[dielec, opacity=0.6] (axis cs:0,0) rectangle (axis cs:1.65,1);
  \addplot[curvaD, line width=2.8pt, opacity=0.45] coordinates {(1.5,0) (1.5,2)}; \addlegendentry{$D\,d/(\eo V_0)$}
  \addplot[curvaE] coordinates {(0.5,0) (0.5,1)}; \addlegendentry{$E\,d/V_0$}
  \addplot[curvaV] coordinates {(0,0) (0.25,1) (1,2)}; \addlegendentry{$\phi/V_0$}
  \addplot[curvaE, forget plot] coordinates {(1.5,1) (1.5,2)};
  \addplot[curvaE, forget plot, dotted, line width=0.5pt] coordinates {(0.5,1) (1.5,1)};
  \node[etiqt, anchor=north east] at (axis cs:1.35,0.95) {$\varepsilon=3\eo$};
\end{axis}
\end{tikzpicture}
\caption{Problema 4.11 ($a=b$, $\er=3$)}
\end{subfigure}\hfill
\begin{subfigure}[b]{0.48\textwidth}\centering
\begin{tikzpicture}
\begin{axis}[width=6.6cm, height=5.4cm, xmin=0, xmax=1.4, ymin=0, ymax=3,
  ytick={0,1,1.6,3}, yticklabels={$0$,$a$,$a+e$,$d$}, ylabel={altura}, xtick={0,0.5,1},
  xticklabels={$0$,$\frac12$,$1$}, grid=major,
  legend style={at={(0.5,1.03)}, anchor=south, legend columns=3, font=\tiny, /tikz/every even column/.append style={column sep=3pt}}]
  \fill[conductor, opacity=0.45] (axis cs:0,1) rectangle (axis cs:1.4,1.6);
  \addplot[curvaE] coordinates {(1,0) (1,1) (0,1) (0,1.6) (1,1.6) (1,3)}; \addlegendentry{$E\,(d-e)/V_0$}
  \addplot[curvaV] coordinates {(0,0) (0.4167,1) (0.4167,1.6) (1,3)}; \addlegendentry{$\phi/V_0$}
  \node[etiqt] at (axis cs:0.95,1.3) {metal};
\end{axis}
\end{tikzpicture}
\caption{Problema 4.12 ($a=1$, $e=0{,}6$, $b=1{,}4$)}
\end{subfigure}
\caption{Perfiles en función de la altura sobre la placa inferior (que se toma a potencial nulo). (a) Capa dieléctrica de
espesor $a$ bajo una capa de aire de espesor $b$: $D$ es continuo, el campo salta en el factor $\er$ y el potencial
cae sobre todo en el aire. (b) Lámina conductora: el campo es nulo dentro del metal y el potencial es allí
constante; fuera, el campo es el de un condensador de separación $d-e$.}
\label{fig:p411-perfil}
\end{figure}


% -----------------------------------------------------------------------------
\subsection{Resumen de resultados y lista de control}\label{sec:resumen-seminarios}

\begin{table}[H]
\centering
\caption{Resultados de los seminarios del Tema 4.}\label{tab:resultados}
\small
\renewcommand{\arraystretch}{1.55}
\begin{tabularx}{\textwidth}{@{}c>{\raggedright\arraybackslash}Xc@{}}
\toprule
Problema & Resultado & Método\\
\midrule
4.1 & $\phi=\dfrac{Pa^3}{3\eo}\dfrac{z}{\abs z^3}$, $E_z=\dfrac{2Pa^3}{3\eo\abs z^3}$ ($\abs z>a$);\quad $\phi=\dfrac{Pz}{3\eo}$, $E_z=-\dfrac{P}{3\eo}$ ($\abs z<a$) & M1\\
4.2 & $\vD=\dfrac{Q}{4\pi r^2}\uv r$,\quad $\vE=\dfrac{Q}{4\pi\varepsilon r^2}\uv r$,\quad $\vP=(\varepsilon-\eo)\dfrac{Q}{4\pi\varepsilon r^2}\uv r$ & M3\\
4.3 & $\vD=\sigma\uv z$,\quad $\vE=\dfrac{\sigma}{\varepsilon}\uv z$,\quad $C=\dfrac{\varepsilon A}{d}=\er C_0$ & M3\\
4.4 & $\vE=\dfrac{\vP}{2\eo}(\cos\theta_2-\cos\theta_1)$ fuera, $-\dfrac{\vP}{2\eo}(2+\cos\theta_1-\cos\theta_2)$ dentro;\quad $\vD=\dfrac{\vP}{2}(\cos\theta_2-\cos\theta_1)$ & M1\\
4.5 & $\rhoP=-3A$,\quad $\sigmaP=A\ell/2$,\quad $\QP=0$ & M1\\
4.6 & $\vD=\vcero$;\quad $\vE=-\dfrac{k}{\eo r^2}\uv r$ en $a<r<2a$, nulo fuera & M2\\
4.7 & $\vD=\vcero$;\quad $\vE=-\dfrac{kr}{2\eo}\uv r$ en $a<r<b$, nulo fuera;\quad $\rhoP=-k$ & M2\\
4.8 & $\phi=\dfrac{Q}{2\pi(\varepsilon_1+\varepsilon_2)}\Big(\dfrac1r-\dfrac1b\Big)$,\quad $C=2\pi(\varepsilon_1+\varepsilon_2)\dfrac{ab}{b-a}$,\quad $\sigma_i=\dfrac{\varepsilon_iQ}{2\pi a^2(\varepsilon_1+\varepsilon_2)}$ & M3, M4\\
4.9 & $\vEint=-\dfrac{b\,\vP}{\eo(a+b)}$,\quad $\vEext=\dfrac{a\,\vP}{\eo(a+b)}$,\quad $\vD=\dfrac{a\,\vP}{a+b}$ & M5\\
4.10 & $Q_i=V_0\dfrac{A_i\varepsilon_{ri}\eo}{d}$,\quad $C=\dfrac{A_1\varepsilon_{r1}\eo}{d}+\dfrac{A_2\varepsilon_{r2}\eo}{d}$ & M4\\
4.11 & $\Eint=\dfrac{V_0\eo}{a\eo+b\varepsilon}$,\quad $\Eext=\dfrac{V_0\varepsilon}{a\eo+b\varepsilon}$,\quad $C=\dfrac{A\varepsilon\eo}{a\eo+b\varepsilon}$ & M4\\
4.12 & $C'=\dfrac{\eo A}{d-e}$ (independiente de la posición) & M4\\
\bottomrule
\end{tabularx}
\end{table}

\begin{metodo}{Lista de control}
\begin{itemize}
  \item ¿He distinguido carga libre (la que controlo) de carga ligada (la que dispone el material)? ¿Uso Gauss para $\vD$ solo con la libre?
  \item Antes de afirmar $\vD=\vcero$, ¿he comprobado que no hay carga libre \emph{y} que $\vP$ es irrotacional, también en las superficies?
  \item ¿He tomado en cada superficie la normal \emph{exterior al dieléctrico} al calcular $\sigmaP=\vP\cdot\nn$?
  \item ¿Se anula la carga ligada total? ¿Coincide el campo con el que da Gauss para $\vE$ con todas las cargas?
  \item ¿Uso $\oiint\vE\cdot\dS=\qfenc/\varepsilon$ solo si un único medio \LHI{} llena todo el espacio relevante?
  \item En las interfases sin carga libre, ¿he impuesto $D_n$ continua y $\vE_t$ continuo (y no al revés)?
  \item Condensadores: interfase paralela a las placas, $\vD$ común y serie; perpendicular, $\vE$ común y paralelo.
  \item ¿Recupera el resultado los límites conocidos: $\varepsilon\to\eo$ (vacío), $\varepsilon\to\infty$ (conductor), $b\to0$, $a\gg\ell$, $a\ll\ell$?
\end{itemize}
\end{metodo}


% =============================================================================
\clearpage
\appendix
\renewcommand{\thesection}{\Alph{section}}
\section{Formulario del Tema 4}\label{sec:formulario}
% =============================================================================

La \cref{tab:formulario} reúne las expresiones centrales del tema, con la referencia al enunciado en que se demuestran.

\begin{table}[H]
\centering
\caption{Formulario del Tema 4. Materiales dieléctricos ($\nn$: normal exterior al dieléctrico, o del medio 1 al 2 en una interfase).}
\label{tab:formulario}
\small
\renewcommand{\arraystretch}{1.6}
\begin{tabularx}{\textwidth}{@{}>{\raggedright\arraybackslash}p{4.1cm}>{\raggedright\arraybackslash}X>{\raggedright\arraybackslash}p{2.7cm}@{}}
\toprule
Magnitud o resultado & Expresión & Referencia\\
\midrule
\multicolumn{3}{@{}l}{\textit{Modelo microscópico}}\\
Polarizabilidad & $\vp=\alpha\vE$;\quad nube uniforme: $\alpha=4\pi\eo a^3$ & \cref{prop:polarizabilidad}\\
Dipolo en campo uniforme & $\vF=\vcero$,\quad $\vtau=\vp\times\vE$,\quad $U=-\vp\cdot\vE$ & \cref{prop:par}\\
\multicolumn{3}{@{}l}{\textit{Polarización y carga ligada}}\\
Vector polarización & $\vP=\dd\vp/\dd V$ & \cref{def:polarizacion}\\
Densidades de polarización & $\rhoP=-\divg\vP$,\quad $\sigmaP=\vP\cdot\nn$ & \cref{teo:densidades}\\
Potencial & $\phi=\dfrac{1}{4\pi\eo}\Big[\displaystyle\int\frac{\rhoP}{R}\dd V'+\oiint\frac{\sigmaP}{R}\dd S'\Big]$ & \cref{teo:densidades}\\
Neutralidad & $\QP=0$;\quad $\qPenc=-\oiint_S\vP\cdot\dS$ & \cref{prop:QP-nula}, \cref{lem:qP-enc}\\
\multicolumn{3}{@{}l}{\textit{Vector desplazamiento}}\\
Definición & $\vD=\eo\vE+\vP$ & \cref{def:D}\\
Ley de Gauss & $\divg\vD=\rhof$,\quad $\oiint_S\vD\cdot\dS=\qfenc$ & \cref{teo:gauss-D}\\
Rotacional & $\rot\vD=\rot\vP$;\quad $\rhof=0$ y $\vP$ irrotacional $\Leftrightarrow$ $\vD\equiv\vcero$ & \cref{prop:D-nulo}\\
\multicolumn{3}{@{}l}{\textit{Medios lineales, homogéneos e isótropos}}\\
Relación constitutiva & $\vP=\eo\chi\vE$,\quad $\vD=\varepsilon\vE$,\quad $\varepsilon=\eo(1+\chi)=\er\eo$ & \cref{def:LHI}, \cref{def:permitividad}\\
Carga ligada de volumen & $\rhoP=-(1-\eo/\varepsilon)\,\rhof$ & \cref{prop:rhoP-rhof}\\
Un único medio & $\vE=\dfrac{1}{4\pi\varepsilon}\displaystyle\int\frac{\rhof\,\uv R}{R^2}\dd V'$ & \cref{prop:un-medio}\\
\multicolumn{3}{@{}l}{\textit{Condiciones de frontera}}\\
Fundamentales & $D_{2n}-D_{1n}=\sigmaf$,\quad $\vE_{2t}=\vE_{1t}$,\quad $\phi_1=\phi_2$ & \cref{teo:frontera}\\
Derivadas & $\eo(E_{2n}-E_{1n})=\sigmaf+\sigmaP$,\ \ $\sigmaP=P_{1n}-P_{2n}$;\quad $\vD_{2t}-\vD_{1t}=\vP_{2t}-\vP_{1t}$ & \cref{cor:otras}\\
Medios \LHI{} & $\varepsilon_2E_{2n}-\varepsilon_1E_{1n}=\sigmaf$;\quad $\tan\theta_2/\varepsilon_2=\tan\theta_1/\varepsilon_1$ si $\sigmaf=0$ & \cref{prop:refraccion}\\
\multicolumn{3}{@{}l}{\textit{Resultados de referencia}}\\
Esfera uniformemente polarizada & dentro $\vE=-\vP/3\eo$, $\vD=\tfrac23\vP$; fuera, dipolo $\tfrac43\pi a^3\vP$ & \cref{prob:esfera-polarizada}\\
Esfera dieléctrica en $\vE_0$ & $\vEint=\dfrac{3}{\er+2}\vE_0$ & \cref{ej:esfera-dielectrica}\\
Condensador plano lleno & $C=\er C_0=\varepsilon A/d$ & \cref{prob:condensador-dielectrico}\\
Cavidades & aguja: $\vEcav=\vE_d$;\quad lenteja: $\vDcav=\vD_d$;\quad esfera: $\vEcav=\vE_d+\vP_d/3\eo$ & \cref{prop:acicular}, \cref{prop:lenticular}, \cref{obs:lorentz}\\
\bottomrule
\end{tabularx}
\end{table}

% =============================================================================
\begin{thebibliography}{99}
\setlength{\itemsep}{0.25em}
\small
\bibitem{T1} D.~Martín Guerrero. \emph{Leyes básicas de la electrostática en el vacío} (Electromagnetismo I, Tema 1). Apuntes formalizados con
  seminarios prácticos. Grado en Física, Universidad de Salamanca, 2026.
\bibitem{T2} D.~Martín Guerrero. \emph{Desarrollo multipolar} (Electromagnetismo I, Tema 2). Apuntes formalizados con seminarios prácticos. Grado
  en Física, Universidad de Salamanca, 2026.
\bibitem{T3} D.~Martín Guerrero. \emph{Conductores en equilibrio electrostático} (Electromagnetismo I, Tema 3). Apuntes formalizados con seminarios
  prácticos. Grado en Física, Universidad de Salamanca, 2026.
\bibitem{T0} D.~Martín Guerrero. \emph{Herramientas matemáticas para el electromagnetismo} (Electromagnetismo I, Tema 0). Apuntes formalizados con
  seminarios prácticos. Grado en Física, Universidad de Salamanca, 2026.
\bibitem{DFA-teoria} Departamento de Física Aplicada. \emph{Electromagnetismo I. Tema IV: Materiales dieléctricos}. Apuntes de teoría. Material de la
  asignatura. Grado en Física, Universidad de Salamanca, 2025.
\bibitem{DFA-figuras} Departamento de Física Aplicada. \emph{Electromagnetismo I. Tema 4: Materiales dieléctricos (figuras)}. Material de la asignatura
  (figuras elaboradas con Mathematica). Grado en Física, Universidad de Salamanca, 2025.
\bibitem{DFA-problemas} Departamento de Física Aplicada. \emph{Electromagnetismo I. Tema 4: Materiales dieléctricos (relación de problemas)}. Material
  de la asignatura. Grado en Física, Universidad de Salamanca, curso 2025-2026.
\bibitem{Guia-T4} D.~Martín Guerrero. \emph{Tema 4. Dieléctricos: formulario, métodos y fichas de los doce problemas}. Guías de estudio de
  Electromagnetismo I. Grado en Física, Universidad de Salamanca, 2026.
\bibitem{Griffiths} D.~J. Griffiths. \emph{Introduction to Electrodynamics}. 4.ª ed. Cap.~4: «Electric fields in matter» (§4.1--§4.4; tablas 4.1 y 4.2).
  Cambridge University Press, 2017.
\bibitem{Wangsness} R.~K. Wangsness. \emph{Electromagnetic Fields}. 2.ª ed. Cap.~10: «Electrostatics in the presence of matter». John Wiley \& Sons, 1986.
\bibitem{Purcell} E.~M. Purcell y D.~J. Morin. \emph{Electricity and Magnetism}. 3.ª ed. Cap.~10: «Electric fields in matter». Cambridge University
  Press, 2013.
\bibitem{Feynman} R.~P. Feynman, R.~B. Leighton y M.~Sands. \emph{The Feynman Lectures on Physics. Volume II}. Caps.~10 y 11: «Dielectrics» e «Inside
  dielectrics». Addison-Wesley, 1964.
\bibitem{Jackson} J.~D. Jackson. \emph{Classical Electrodynamics}. 3.ª ed. §4.3--§4.6: polarización, desplazamiento, condiciones de contorno con
  dieléctricos, relación de Clausius--Mossotti. John Wiley \& Sons, 1999.
\bibitem{Landau} L.~D. Landau, E.~M. Lifshitz y L.~P. Pitaevskii. \emph{Electrodynamics of Continuous Media}. 2.ª ed. §6--§8: dieléctricos, campo en
  cavidades y elipsoides en un campo uniforme. Pergamon Press, 1984.
\bibitem{Osborn} J.~A. Osborn. «Demagnetizing factors of the general ellipsoid». \emph{Physical Review} 67 (1945), 351--357.
  doi:\href{https://doi.org/10.1103/PhysRev.67.351}{10.1103/PhysRev.67.351}.
\bibitem{CODATA} E.~Tiesinga, P.~J. Mohr, D.~B. Newell y B.~N. Taylor. «CODATA recommended values of the fundamental physical constants: 2018».
  \emph{Reviews of Modern Physics} 93 (2021), 025010. doi:\href{https://doi.org/10.1103/RevModPhys.93.025010}{10.1103/RevModPhys.93.025010}.
\end{thebibliography}

